\chapter{动态规划进阶、容器原理和源码阅读}

刷题到中后期，很多错误不再来自不会写循环，而来自状态定义不稳定、容器成本估计错误、以及对标准库行为理解不够。动态规划进阶要把问题拆成“选择、容量、阶段、状态”；C++ 容器原理要知道数据如何存储、什么时候扩容、迭代器何时失效；源码阅读要能从接口行为反推实现约束。面试中如果能同时讲算法和容器成本，答案会明显更可靠。

0-1 背包是 DP 进阶的基础模型。给定若干物品，每个物品只能选一次，有重量和价值，背包容量有限，求最大价值。暴力方法是枚举每个物品选或不选，复杂度 \complexity{O(2^n)}。状态可以定义为 \code{dp[i][cap]}：只考虑前 \code{i} 个物品、容量为 \code{cap} 时的最大价值。第 \code{i} 个物品要么不选，要么在容量足够时选择。

\begin{lstlisting}[language=C++,caption={0-1 背包：一维滚动数组}]
int knapsack_01(const std::vector<int>& weights,
                const std::vector<int>& values,
                int capacity)
{
    std::vector<int> dp(capacity + 1, 0);

    for (int i = 0; i < static_cast<int>(weights.size()); ++i) {
        for (int cap = capacity; cap >= weights[i]; --cap) {
            dp[cap] = std::max(dp[cap], dp[cap - weights[i]] + values[i]);
        }
    }

    return dp[capacity];
}
\end{lstlisting}

一维滚动数组必须倒序遍历容量，因为每个物品只能使用一次。若正序遍历，\code{dp[cap - weights[i]]} 可能已经包含当前物品，等价于重复选择。这个遍历方向就是 0-1 背包和完全背包最重要的区别。面试中不要只背公式，要说清楚倒序是为了让当前物品的转移读取上一轮状态。

分割等和子集是 0-1 背包的判定版本。题目问数组能否分成两个和相等的子集。总和若为奇数，直接不可能；否则问题变成能否从数组中选一些数，使和为 \code{sum / 2}。每个数只能选一次，所以是 0-1 背包。状态 \code{dp[cap]} 表示是否能凑出容量 \code{cap}。

\begin{lstlisting}[language=C++,caption={LeetCode 416 分割等和子集：0-1 背包可达性}]
bool can_partition(const std::vector<int>& nums)
{
    const int total = std::accumulate(nums.begin(), nums.end(), 0);
    if (total % 2 != 0) {
        return false;
    }

    const int target = total / 2;
    std::vector<char> dp(target + 1, false);
    dp[0] = true;

    for (const int x : nums) {
        for (int cap = target; cap >= x; --cap) {
            dp[cap] = dp[cap] || dp[cap - x];
        }
    }

    return dp[target];
}
\end{lstlisting}

这题的关键不是代码，而是把“分成两个相等子集”转成“选择若干元素凑出总和一半”。如果数组元素都是正数，背包容量有限，DP 很自然；如果存在负数，容量模型就不直接适用，需要偏移或换思路。复杂度是 \complexity{O(n \cdot target)}，其中 \code{target} 与数值大小有关，所以这是伪多项式复杂度，不是严格输入长度多项式。

目标和可以转成背包计数。给每个数加正号或负号，使表达式等于 target。设正数集合和为 P，负数集合和为 N，总和 S。则 P - N = target，P + N = S，推出 P = (S + target) / 2。问题变成从数组中选若干数，使和为 P 的方案数。暴力枚举符号是 \complexity{O(2^n)}；DP 保存每个和的方案数。

\begin{lstlisting}[language=C++,caption={LeetCode 494 目标和：转换为 0-1 背包计数}]
int find_target_sum_ways(const std::vector<int>& nums, int target)
{
    const int total = std::accumulate(nums.begin(), nums.end(), 0);
    if (std::abs(target) > total || (total + target) % 2 != 0) {
        return 0;
    }

    const int positive_sum = (total + target) / 2;
    std::vector<int> dp(positive_sum + 1, 0);
    dp[0] = 1;

    for (const int x : nums) {
        for (int sum = positive_sum; sum >= x; --sum) {
            dp[sum] += dp[sum - x];
        }
    }

    return dp[positive_sum];
}
\end{lstlisting}

这题要特别注意 0。若数组里有 0，给它加正号或负号都不改变和，但方案数应翻倍。上述 DP 能自然处理：当 \code{x == 0} 时，\code{dp[sum] += dp[sum]}，所有已有方案翻倍。很多手写特殊逻辑反而容易错。计数型 DP 要关注整型范围，力扣原题结果通常能放进 int，但工程中应考虑 \code{long long}。

零钱兑换 II 是完全背包计数。硬币可以无限使用，问凑成金额的组合数。注意是组合数，不是排列数。若外层遍历金额、内层遍历硬币，会把不同顺序当作不同方案；若外层遍历硬币、内层正序遍历金额，每种硬币按固定顺序加入，就只统计组合。

\begin{lstlisting}[language=C++,caption={LeetCode 518 零钱兑换 II：完全背包组合计数}]
int change(int amount, const std::vector<int>& coins)
{
    std::vector<int> dp(amount + 1, 0);
    dp[0] = 1;

    for (const int coin : coins) {
        for (int value = coin; value <= amount; ++value) {
            dp[value] += dp[value - coin];
        }
    }

    return dp[amount];
}
\end{lstlisting}

为什么容量正序？因为当前硬币可以重复使用，\code{dp[value - coin]} 应该允许已经使用当前硬币。为什么硬币在外层？因为组合不关心顺序，我们规定按硬币种类逐步开放选择。这个题非常适合区分 0-1 背包、完全背包、组合计数和排列计数。只要循环顺序解释不清，背包题就还没有真正掌握。

股票题是状态机 DP。买卖股票的最佳时机 I 只允许一次交易，最简单做法是扫描时维护历史最低买入价和当前最大利润。暴力枚举买入日和卖出日是 \complexity{O(n^2)}；优化来自“卖出日固定时，最优买入日一定是它之前价格最低的一天”。

\begin{lstlisting}[language=C++,caption={LeetCode 121 买卖股票 I：维护历史最低价格}]
int max_profit_one_transaction(const std::vector<int>& prices)
{
    int min_price = std::numeric_limits<int>::max();
    int answer = 0;

    for (const int price : prices) {
        min_price = std::min(min_price, price);
        answer = std::max(answer, price - min_price);
    }

    return answer;
}
\end{lstlisting}

含冷冻期的股票题不能只维护最低价，因为今天能不能买取决于前一天是否卖出。状态机更清楚：\code{hold} 表示今天结束时持有股票的最大收益；\code{sold} 表示今天刚卖出；\code{rest} 表示今天结束时不持有且不在刚卖出的状态。转移时，买入只能从 \code{rest} 来，卖出只能从 \code{hold} 来，休息可以从昨天的 \code{rest} 或 \code{sold} 来。

\begin{lstlisting}[language=C++,caption={LeetCode 309 含冷冻期股票：三状态 DP}]
int max_profit_with_cooldown(const std::vector<int>& prices)
{
    if (prices.empty()) {
        return 0;
    }

    int hold = -prices[0];
    int sold = 0;
    int rest = 0;

    for (int i = 1; i < static_cast<int>(prices.size()); ++i) {
        const int previous_hold = hold;
        const int previous_sold = sold;
        const int previous_rest = rest;

        hold = std::max(previous_hold, previous_rest - prices[i]);
        sold = previous_hold + prices[i];
        rest = std::max(previous_rest, previous_sold);
    }

    return std::max(sold, rest);
}
\end{lstlisting}

这段代码的关键是先保存昨天状态，避免同一天内状态互相污染。返回值不能是 \code{hold}，因为结束时仍持有股票不代表已经实现收益。股票 DP 的复盘模板是：每天结束时有哪些状态，状态之间允许哪些动作，动作是否受冷冻期、手续费、交易次数限制。把限制放进状态机，公式就会自然出现。

动态规划最容易被学成公式库，但公式不是起点。起点永远是暴力搜索：有哪些选择，搜索树的节点表示什么，哪些节点会重复出现。以最长递增子序列为例，暴力做法是对每个元素选择“选或不选”，同时维护上一个被选元素，复杂度指数级。重复子问题来自“处理到某个位置、上一个选择约束相同”时，后续选择空间完全一样。第一种 DP 定义可以是 \code{dp[i]}：以 \code{nums[i]} 结尾的最长递增子序列长度。这样每个状态都问一个局部稳定的问题：最后一个元素固定为 \code{i}，前一个元素只能从 \code{0..i-1} 中找，并且值要更小。

\begin{lstlisting}[language=C++,caption={LeetCode 300 最长递增子序列：二次 DP 和路径还原}]
std::vector<int> reconstruct_lis(const std::vector<int>& nums)
{
    const int n = static_cast<int>(nums.size());
    if (n == 0) {
        return {};
    }

    std::vector<int> dp(n, 1);
    std::vector<int> parent(n, -1);
    int best_index = 0;

    for (int i = 0; i < n; ++i) {
        for (int j = 0; j < i; ++j) {
            if (nums[j] >= nums[i] || dp[j] + 1 <= dp[i]) {
                continue;
            }
            dp[i] = dp[j] + 1;
            parent[i] = j;
        }
        if (dp[i] > dp[best_index]) {
            best_index = i;
        }
    }

    std::vector<int> sequence;
    for (int index = best_index; index != -1; index = parent[index]) {
        sequence.push_back(nums[index]);
    }
    std::reverse(sequence.begin(), sequence.end());
    return sequence;
}
\end{lstlisting}

这段代码不是最优复杂度，但它最适合讲清楚状态。 \code{parent[i]} 记录让 \code{dp[i]} 变大的前驱下标，所以最后可以从最优结尾倒着还原路径。正确性来自一个简单不变量：外层处理到 \code{i} 时，所有 \code{j < i} 的 \code{dp[j]} 已经是以 \code{j} 结尾的最优答案；若 \code{nums[j] < nums[i]}，把 \code{i} 接到以 \code{j} 结尾的最优序列后面，一定得到一个合法递增序列。枚举所有可能前驱后取最大值，就不会漏掉以 \code{i} 结尾的最优解。复杂度 \complexity{O(n^2)}，空间 \complexity{O(n)}。

最长递增子序列还可以优化到 \complexity{O(n\log n)}。优化依据不是简单二分，而是维护一个更抽象的数组 \code{tails}：长度为 \code{len + 1} 的递增子序列，结尾最小值是多少。结尾越小，未来接上新元素的机会越大。扫描每个数 \code{x}，用二分找到第一个大于等于 \code{x} 的位置并替换它；如果找不到，就把 \code{x} 追加到末尾。 \code{tails} 不一定是一条真实子序列，但它的长度等于 LIS 长度。这个区别必须讲清楚，否则读者会误以为 \code{tails} 本身就是答案路径。

\begin{lstlisting}[language=C++,caption={LeetCode 300 最长递增子序列：tails 数组和二分优化}]
int length_of_lis(const std::vector<int>& nums)
{
    std::vector<int> tails;
    tails.reserve(nums.size());

    for (const int x : nums) {
        auto it = std::lower_bound(tails.begin(), tails.end(), x);
        if (it == tails.end()) {
            tails.push_back(x);
        } else {
            *it = x;
        }
    }

    return static_cast<int>(tails.size());
}
\end{lstlisting}

为什么替换不会破坏答案？因为替换位置 \code{pos} 的含义是：存在长度为 \code{pos + 1} 的递增子序列以 \code{x} 结尾，并且 \code{x} 不大于原来的结尾。更小的结尾不会让未来变差，只会让未来更容易扩展。若题目要求严格递增，使用 \code{lower_bound}；若要求非递减，使用 \code{upper_bound}。这是面试里常见追问：二分函数的选择取决于相等元素是否允许接在后面。

网格路径 DP 是二维状态的入门。不同路径问从左上到右下只能向右或向下走有多少条路径。暴力递归会从每个格子继续向右和向下，重复计算大量子矩形。状态 \code{dp[row][col]} 表示到达当前格子的路径数，它只依赖上方和左方。由于每个状态只依赖上一行和当前行左侧，可以压缩成一维数组。这里的遍历顺序不是随便的：从左到右更新时，\code{dp[col]} 仍表示上一行同列，\code{dp[col - 1]} 已经表示当前行左侧。

\begin{lstlisting}[language=C++,caption={LeetCode 62 不同路径：一维滚动数组}]
int unique_paths(int rows, int cols)
{
    std::vector<int> dp(cols, 1);
    for (int row = 1; row < rows; ++row) {
        for (int col = 1; col < cols; ++col) {
            dp[col] += dp[col - 1];
        }
    }
    return dp[cols - 1];
}
\end{lstlisting}

最小路径和与不同路径结构相同，只是“方案数相加”变成“代价取最小”。状态 \code{dp[col]} 表示处理到当前行时，到达该列的最小路径和。第一行只能从左边来，第一列只能从上边来，其余位置取上方和左方较小者。边界处理是这类题的主要错误来源。一个稳妥写法是先初始化第一行，再逐行处理第一列和内部格子；不要在双循环里塞很多特殊判断，让主转移被边界噪声淹没。

\begin{lstlisting}[language=C++,caption={LeetCode 64 最小路径和：边界初始化后主转移}]
int min_path_sum(const std::vector<std::vector<int>>& grid)
{
    const int rows = static_cast<int>(grid.size());
    const int cols = static_cast<int>(grid[0].size());
    std::vector<int> dp(cols, 0);

    dp[0] = grid[0][0];
    for (int col = 1; col < cols; ++col) {
        dp[col] = dp[col - 1] + grid[0][col];
    }

    for (int row = 1; row < rows; ++row) {
        dp[0] += grid[row][0];
        for (int col = 1; col < cols; ++col) {
            dp[col] = std::min(dp[col], dp[col - 1]) + grid[row][col];
        }
    }

    return dp[cols - 1];
}
\end{lstlisting}

编辑距离是字符串 DP 的核心题。给定两个单词，允许插入、删除、替换，求把一个单词变成另一个单词的最少操作数。暴力递归从两个字符串末尾比较：若字符相同，两个指针都左移；若不同，可以删除一个字符、插入一个字符、替换一个字符。暴力慢是因为同一对前缀长度会被多条路径反复访问。状态 \code{dp[i][j]} 表示 \code{word1} 的前 \code{i} 个字符变成 \code{word2} 的前 \code{j} 个字符的最少操作数。空串边界很重要：\code{dp[i][0] = i}，\code{dp[0][j] = j}。

\begin{lstlisting}[language=C++,caption={LeetCode 72 编辑距离：两行滚动数组}]
int min_distance(const std::string& word1, const std::string& word2)
{
    const int n = static_cast<int>(word1.size());
    const int m = static_cast<int>(word2.size());

    std::vector<int> previous(m + 1, 0);
    std::vector<int> current(m + 1, 0);
    for (int j = 0; j <= m; ++j) {
        previous[j] = j;
    }

    for (int i = 1; i <= n; ++i) {
        current[0] = i;
        for (int j = 1; j <= m; ++j) {
            if (word1[i - 1] == word2[j - 1]) {
                current[j] = previous[j - 1];
                continue;
            }
            const int remove_cost = previous[j] + 1;
            const int insert_cost = current[j - 1] + 1;
            const int replace_cost = previous[j - 1] + 1;
            current[j] = std::min({remove_cost, insert_cost, replace_cost});
        }
        std::swap(previous, current);
    }

    return previous[m];
}
\end{lstlisting}

三个操作的含义要能对上表格位置。删除 \code{word1[i-1]} 后，问题变成 \code{word1} 前 \code{i-1} 个字符到 \code{word2} 前 \code{j} 个字符，所以来自 \code{previous[j]}；插入 \code{word2[j-1]} 后，等价于先把 \code{word1} 前 \code{i} 个字符变成 \code{word2} 前 \code{j-1} 个字符，所以来自 \code{current[j-1]}；替换则来自 \code{previous[j-1]}。两行滚动数组的空间是 \complexity{O(m)}，时间是 \complexity{O(nm)}。如果题目要求还原操作序列，就不能只保留两行，需要记录决策或重新回溯。

最长公共子序列和编辑距离很接近。状态 \code{dp[i][j]} 表示两个前缀的 LCS 长度。若末尾字符相同，就从 \code{dp[i-1][j-1] + 1} 转移；否则取删掉任一末尾后的最大值。它和最长公共子串不同：子序列允许不连续，子串必须连续。这个概念差异决定了转移差异。最长公共子串在字符不同时必须归零，LCS 在字符不同时取上方或左方最大值。面试里一旦题目出现“可以删除若干字符但保持相对顺序”，就要想到子序列。

\begin{lstlisting}[language=C++,caption={LeetCode 1143 最长公共子序列：一维数组和左上角缓存}]
int longest_common_subsequence(const std::string& text1,
                               const std::string& text2)
{
    const int n = static_cast<int>(text1.size());
    const int m = static_cast<int>(text2.size());
    std::vector<int> dp(m + 1, 0);

    for (int i = 1; i <= n; ++i) {
        int diagonal = 0;
        for (int j = 1; j <= m; ++j) {
            const int old_up = dp[j];
            if (text1[i - 1] == text2[j - 1]) {
                dp[j] = diagonal + 1;
            } else {
                dp[j] = std::max(dp[j], dp[j - 1]);
            }
            diagonal = old_up;
        }
    }

    return dp[m];
}
\end{lstlisting}

一维 LCS 的难点在于保存左上角。更新前的 \code{dp[j]} 是上一行同列，\code{dp[j-1]} 是当前行左侧，\code{diagonal} 是上一行左上角。每次循环结束前把旧的 \code{dp[j]} 放进 \code{diagonal}，供下一列使用。若你觉得这个写法难以保证正确，可以先写二维表；面试时清楚解释二维转移比写一个容易错的一维优化更重要。空间优化应该建立在状态依赖已经完全理解之后。

单词拆分是“前缀可达性”DP。题目问字符串能否被字典中的单词拼接出来。暴力方法枚举第一段单词，然后递归处理剩余后缀；同一个起点会被反复递归。状态 \code{dp[i]} 表示前 \code{i} 个字符能否被拆分。转移枚举最后一个单词的起点 \code{j}：若 \code{dp[j]} 为真，且 \code{s[j..i)} 在字典中，则 \code{dp[i]} 为真。这里要注意 \code{substr} 的成本，它会创建新字符串。数据量不大时可接受；数据量大时，应限制最大单词长度，或用字典树、字符串视图、滚动哈希减少复制。

\begin{lstlisting}[language=C++,caption={LeetCode 139 单词拆分：限制最大单词长度降低无效枚举}]
bool word_break(const std::string& s,
                const std::vector<std::string>& word_dict)
{
    std::unordered_set<std::string> words(word_dict.begin(), word_dict.end());
    int max_word_length = 0;
    for (const std::string& word : word_dict) {
        max_word_length = std::max(max_word_length,
                                   static_cast<int>(word.size()));
    }

    const int n = static_cast<int>(s.size());
    std::vector<char> dp(n + 1, false);
    dp[0] = true;

    for (int end = 1; end <= n; ++end) {
        const int begin_limit = std::max(0, end - max_word_length);
        for (int begin = end - 1; begin >= begin_limit; --begin) {
            if (!dp[begin]) {
                continue;
            }
            if (words.contains(s.substr(begin, end - begin))) {
                dp[end] = true;
                break;
            }
        }
    }

    return dp[n];
}
\end{lstlisting}

这个优化不是改变问题模型，而是减少不可能成为单词的长度。若字典中最长单词长度为 \code{L}，枚举起点时只需要看最近 \code{L} 个位置。复杂度粗略为 \complexity{O(nL)} 次查询，但每次 \code{substr} 仍有拷贝成本。面试表达可以分层：先给清楚的 DP，再说明字符串切片成本，最后提出可选优化。不要一上来写复杂字典树，把主线埋掉。

区间 DP 的状态通常是“处理一个连续区间的最优值”。它与普通线性 DP 的区别是：最后一步可能把区间切成两半，或者选择区间内某个点作为最后处理对象。戳气球是代表题。直接思考“先戳哪个气球”很混乱，因为戳掉一个气球会改变邻居；换成“最后戳哪个气球”就稳定了。设 \code{dp[left][right]} 表示开区间 \code{(left, right)} 内的气球全部戳掉能获得的最大硬币数。若最后戳 \code{mid}，此时它左右邻居一定是 \code{left} 和 \code{right}，收益为 \code{nums[left] * nums[mid] * nums[right]}，左右两边互不影响。

\begin{lstlisting}[language=C++,caption={LeetCode 312 戳气球：选择最后一个被戳的气球}]
int max_coins(const std::vector<int>& nums)
{
    std::vector<int> values;
    values.reserve(nums.size() + 2);
    values.push_back(1);
    values.insert(values.end(), nums.begin(), nums.end());
    values.push_back(1);

    const int n = static_cast<int>(values.size());
    std::vector<std::vector<int>> dp(n, std::vector<int>(n, 0));

    for (int length = 2; length < n; ++length) {
        for (int left = 0; left + length < n; ++left) {
            const int right = left + length;
            for (int mid = left + 1; mid < right; ++mid) {
                const int score = dp[left][mid] + dp[mid][right] +
                    values[left] * values[mid] * values[right];
                dp[left][right] = std::max(dp[left][right], score);
            }
        }
    }

    return dp[0][n - 1];
}
\end{lstlisting}

区间 DP 的遍历顺序按区间长度从小到大，因为大区间依赖小区间。这里的边界 \code{1} 是虚拟气球，不会被戳，只是让左右边界收益统一。正确性要抓住“最后一步”：如果 \code{mid} 是最后被戳的气球，那么 \code{left..mid} 和 \code{mid..right} 两个子问题已经在它之前完成，且互相独立；枚举所有可能的最后气球，取最大值。复杂度是 \complexity{O(n^3)}，空间 \complexity{O(n^2)}。若面试官追问能否优化，要先说明一般区间 DP 不一定有四边形不等式或决策单调性，不能乱套优化。

状态压缩 DP 把集合放进整数二进制位。适用条件通常是元素数量不大，例如 \code{n <= 20}。每一位表示某个任务、节点或物品是否已经被选择。经典问题是旅行商或访问所有节点的最短路径。若图无权，访问所有节点的最短路径可以用 BFS，把状态定义为 \code{(node, mask)}：当前在 \code{node}，已经访问的节点集合是 \code{mask}。这其实是状态压缩和图 BFS 的结合，而不是传统表格 DP。

\begin{lstlisting}[language=C++,caption={LeetCode 847 访问所有节点的最短路径：节点加集合状态 BFS}]
int shortest_path_length(const std::vector<std::vector<int>>& graph)
{
    const int n = static_cast<int>(graph.size());
    const int full_mask = (1 << n) - 1;

    std::queue<std::pair<int, int>> queue;
    std::vector<std::vector<int>> distance(n, std::vector<int>(1 << n, -1));

    for (int node = 0; node < n; ++node) {
        const int mask = 1 << node;
        distance[node][mask] = 0;
        queue.emplace(node, mask);
    }

    while (!queue.empty()) {
        const auto [node, mask] = queue.front();
        queue.pop();
        if (mask == full_mask) {
            return distance[node][mask];
        }

        for (const int next : graph[node]) {
            const int next_mask = mask | (1 << next);
            if (distance[next][next_mask] != -1) {
                continue;
            }
            distance[next][next_mask] = distance[node][mask] + 1;
            queue.emplace(next, next_mask);
        }
    }

    return 0;
}
\end{lstlisting}

这题如果只把节点作为访问标记，会出错。到达同一个节点时，已访问集合不同，未来能力不同。 \code{(node, mask)} 才是完整状态。多源初始化也很关键：路径可以从任意节点开始，所以所有单节点状态都在第 0 层。复杂度是 \complexity{O(n2^n + m2^n)}，只适合小 \code{n}。状态压缩题要先估算 \code{2^n} 是否可承受，再决定是否使用；看到集合并不意味着一定能压缩，如果 \code{n} 是一百，二进制集合就不现实。

另一类状态压缩是子集枚举。比如把数组分成若干组、每组和相等。可以用 \code{dp[mask]} 表示选择集合 \code{mask} 后，当前桶中已经放入的和对目标桶容量取模是多少；若 \code{dp[mask] == -1} 表示不可达。每次尝试加入一个未选元素，只要当前桶不超容量，就更新新集合。这个写法比回溯更像 DP，因为每个集合状态只计算一次。

\begin{lstlisting}[language=C++,caption={LeetCode 698 划分为 K 个等和子集：集合状态 DP}]
bool can_partition_k_subsets(std::vector<int> nums, int k)
{
    const int total = std::accumulate(nums.begin(), nums.end(), 0);
    if (total % k != 0) {
        return false;
    }
    const int target = total / k;
    std::sort(nums.begin(), nums.end(), std::greater<int>());
    if (!nums.empty() && nums.front() > target) {
        return false;
    }

    const int n = static_cast<int>(nums.size());
    const int state_count = 1 << n;
    std::vector<int> dp(state_count, -1);
    dp[0] = 0;

    for (int mask = 0; mask < state_count; ++mask) {
        if (dp[mask] == -1) {
            continue;
        }
        for (int i = 0; i < n; ++i) {
            if ((mask & (1 << i)) != 0) {
                continue;
            }
            const int next_sum = dp[mask] + nums[i];
            if (next_sum > target) {
                continue;
            }
            const int next_mask = mask | (1 << i);
            dp[next_mask] = next_sum % target;
        }
    }

    return dp[state_count - 1] == 0;
}
\end{lstlisting}

这个 DP 的状态值不是总和，而是当前桶的剩余进度。每当一个桶刚好装满，对 \code{target} 取模后的进度回到 0，表示开始装下一个桶。排序不是正确性必需，但大的数先处理能更早剪掉不可能状态。复杂度是 \complexity{O(n2^n)}。与回溯相比，状态 DP 的优势是避免不同选择顺序到达同一集合时重复搜索；缺点是内存随 \code{2^n} 增长。

单调队列优化 DP 是从“转移枚举范围”中挤出复杂度。以带限制的最大子序列和为例，题目要求选择一个非空子序列，相邻被选元素下标差最多为 \code{k}，求最大和。朴素 DP 定义 \code{dp[i]} 为必须选择 \code{nums[i]} 且以它结尾的最大和，则
\[
dp[i] = nums[i] + \max(0, dp[j]) \quad i-k \le j < i.
\]
如果每个 \code{i} 都扫描前 \code{k} 个位置，复杂度 \complexity{O(nk)}。优化来自窗口最大值：我们只需要最近 \code{k} 个 \code{dp[j]} 的最大值，用单调队列维护即可。

\begin{lstlisting}[language=C++,caption={LeetCode 1425 受限子序列和：单调队列优化转移最大值}]
int constrained_subset_sum(const std::vector<int>& nums, int k)
{
    std::deque<int> deque;
    std::vector<int> dp(nums.size(), 0);
    int answer = std::numeric_limits<int>::min();

    for (int i = 0; i < static_cast<int>(nums.size()); ++i) {
        while (!deque.empty() && deque.front() < i - k) {
            deque.pop_front();
        }

        dp[i] = nums[i];
        if (!deque.empty()) {
            dp[i] = std::max(dp[i], nums[i] + dp[deque.front()]);
        }
        answer = std::max(answer, dp[i]);

        while (!deque.empty() && dp[deque.back()] <= dp[i]) {
            deque.pop_back();
        }
        if (dp[i] > 0) {
            deque.push_back(i);
        }
    }

    return answer;
}
\end{lstlisting}

队列中保存下标，而不是保存值，因为要判断是否过期。队首永远是当前窗口内 \code{dp} 最大的位置。只把正的 \code{dp[i]} 放入队列，是因为负贡献不如从当前元素重新开始。正确性来自两个不变量：队列下标从前到后递增，保证能判断窗口；对应 \code{dp} 值从前到后递减，保证队首最大。每个下标最多进队一次、出队一次，所以总复杂度 \complexity{O(n)}。

动态规划和最短路之间有很深的联系。许多 DP 可以看成在有向无环图上找路径：状态是节点，转移是边，遍历顺序是拓扑序。若状态图有环但边权非负，可能变成 Dijkstra；若边权全为 1，可能变成 BFS；若存在负边但没有负环，可能变成 Bellman-Ford。这样理解后，题型边界会更清楚：不是所有“求最优”都叫 DP，也不是所有“状态转移”都要填表。你要先问状态图是否有自然阶段顺序，是否允许回头，转移代价是否单调。

\begin{warningbox}
DP 常见误区有四个。第一，只背公式，不说明状态含义，导致换一个题面就失效。第二，滚动数组时忘记转移读的是上一阶段还是当前阶段。第三，把组合数和排列数混在一起，循环顺序写反。第四，看到最值就写 DP，却没有检查状态数量是否可承受。每次写完 DP，都要用一句话解释 \code{dp} 的含义，用一句话解释遍历顺序，用一句话解释为什么不会漏解或重复计数。
\end{warningbox}

C++ 容器原理是刷题代码质量的底座。 \code{std::vector} 是连续数组，支持随机访问，扩容时会搬迁元素，导致指针、引用和迭代器失效。 \code{std::deque} 分段连续，支持两端高效插入删除，但随机访问局部性通常不如 \code{vector}。 \code{std::list} 节点稳定，但缓存局部性差，刷题中很少是首选。 \code{std::unordered_map} 平均常数查询，但 rehash 会让迭代器失效，且不维护顺序。 \code{std::map} 有序但每次操作是 \complexity{O(\log n)}。

阅读标准库源码时，不需要一开始钻进所有模板细节。先看三层：公开接口、核心数据成员、关键操作路径。以 \code{vector} 为例，大多数实现内部会维护三个指针或等价状态：开始位置、当前结束位置、容量结束位置。 \code{size()} 是当前结束减开始，\code{capacity()} 是容量结束减开始，\code{push_back} 在容量足够时原地构造，在容量不足时申请更大内存并移动旧元素。理解这三根指针，就能解释很多迭代器失效规则。

\begin{lstlisting}[language=C++,caption={vector 使用习惯：reserve、resize 和引用失效}]
void collect_values(const std::vector<int>& input)
{
    std::vector<int> values;
    values.reserve(input.size());

    for (const int x : input) {
        values.push_back(x * 2);
    }

    const int* data = values.data();
    values.push_back(42);
    // 如果 push_back 触发扩容，data 已经失效，不能继续解引用。
}
\end{lstlisting}

\code{reserve} 和 \code{resize} 的区别必须非常熟。 \code{reserve} 只保证容量，不改变大小；\code{resize} 改变大小并构造元素。DP 数组需要可下标访问，应使用 \code{vector<int> dp(n + 1)} 或 \code{resize}；收集答案时知道最多数量，可以先 \code{reserve} 再 \code{push_back}。把二者混用，是 C++ 刷题里非常常见的越界来源。

\code{unordered_map} 的源码阅读重点是桶数组、哈希函数、负载因子和 rehash。插入元素时，容器通过哈希值定位桶，再在桶内查找等价 key。平均常数复杂度依赖哈希分布；当元素数量相对桶数过多，会触发 rehash，重新分配桶并把元素放到新桶中。刷题时使用 \code{reserve} 可以减少 rehash 次数，尤其是两数之和、前缀和计数、频次统计这类一次性插入较多 key 的题。

\begin{lstlisting}[language=C++,caption={unordered\_map 查询：contains 与 find 的取舍}]
int count_frequency(const std::vector<int>& nums, int target)
{
    std::unordered_map<int, int> frequency;
    frequency.reserve(nums.size());

    for (const int x : nums) {
        ++frequency[x];
    }

    const auto found = frequency.find(target);
    if (found == frequency.end()) {
        return 0;
    }
    return found->second;
}
\end{lstlisting}

如果只判断存在，C++20 的 \code{contains} 很清楚；如果判断后还要取值，用 \code{find} 避免查两次。不要用 \code{operator[]} 做纯查询，因为 key 不存在时会插入默认值，改变容器状态。这个细节在前缀和计数里尤其重要：有些查询应该只读历史状态，不能悄悄插入一个 0。

源码阅读也能帮助理解算法复杂度的真实成本。字符串 \code{substr} 会创建新字符串，放在双层循环里可能把 \complexity{O(n^2)} 变成接近 \complexity{O(n^3)} 的实际成本；\code{vector<vector<int>>} 的每一行是独立分配，矩阵遍历仍然可用，但内存不一定完全连续；\code{priority_queue} 没有删除任意元素和 decrease-key，所以 Dijkstra 常用“压入新状态，弹出时跳过过期状态”的写法。理解容器能力边界，才能解释为什么代码这样写。

\code{std::deque} 的名字容易误导初学者。它不是双向链表，而是分段连续数组。典型实现会维护一个“map”，其中每个槽指向一段固定大小的缓冲区；首尾插入删除大多只移动段内指针，必要时分配新段。它支持随机访问，但随机访问需要先定位段再定位段内偏移，局部性通常不如 \code{vector}。这解释了为什么滑动窗口最大值适合用 \code{deque} 存下标，而普通顺序 DP 数组不应该为了“可能两端操作”随手换成 \code{deque}。

\begin{lstlisting}[language=C++,caption={deque 在单调队列中的正确用法}]
std::vector<int> max_sliding_window(const std::vector<int>& nums, int k)
{
    std::deque<int> deque;
    std::vector<int> answer;
    answer.reserve(nums.size());

    for (int i = 0; i < static_cast<int>(nums.size()); ++i) {
        while (!deque.empty() && deque.front() <= i - k) {
            deque.pop_front();
        }
        while (!deque.empty() && nums[deque.back()] <= nums[i]) {
            deque.pop_back();
        }
        deque.push_back(i);
        if (i + 1 >= k) {
            answer.push_back(nums[deque.front()]);
        }
    }

    return answer;
}
\end{lstlisting}

这里 \code{deque} 保存下标而不是值，原因和单调队列 DP 一样：需要判断过期位置。下标还能处理重复值，若只存值，窗口移出时很难知道队首值是否对应离开的元素。源码层面看，\code{pop_front} 和 \code{push_back} 都是 \code{deque} 的强项；若用 \code{vector} 在头部删除，就会移动大量元素；若用 \code{list}，虽然删除快，但无法高效访问队尾并且缓存局部性差。容器选择和算法不变量应该匹配。

\code{std::string} 在现代实现中通常具有小字符串优化。短字符串可以直接存放在对象内部，不必单独堆分配；长字符串才指向堆内存。这个事实能解释为什么偶尔创建短字符串不一定灾难，但不能因此在热循环里无限 \code{substr}。标准并不要求具体的小字符串容量，不同标准库实现也不同，因此工程代码不能依赖“小于多少字符一定不分配”。源码阅读要区分“标准保证的语义”和“某个实现的优化”。

\begin{lstlisting}[language=C++,caption={字符串扫描：用下标边界避免大量 substr}]
bool has_word_at(const std::string& text,
                 int begin,
                 const std::string& word)
{
    if (begin + static_cast<int>(word.size()) >
        static_cast<int>(text.size())) {
        return false;
    }
    for (int i = 0; i < static_cast<int>(word.size()); ++i) {
        if (text[begin + i] != word[i]) {
            return false;
        }
    }
    return true;
}
\end{lstlisting}

这段代码展示的是思想：当你只需要比较一段文本时，不一定要构造新字符串。C++20 可以使用 \code{std::string_view} 表示非拥有切片，但要注意生命周期；视图不能比原字符串活得更久。刷题中 \code{substr} 更短、更不易错，数据量允许时可以使用；工程中如果它处在双层或三层循环内，就应该估算分配和拷贝成本。

\code{std::priority_queue} 是容器适配器，默认底层容器通常是 \code{vector}，并用堆算法维护最大堆。它只暴露 \code{top}、\code{push}、\code{pop}，不支持遍历、不支持删除任意元素、不支持 decrease-key。这些限制不是缺陷，而是接口选择：标准库给的是简单、稳定、低成本的堆抽象。Dijkstra 中距离变小后重新入堆，是顺应这个接口；双堆中位数中延迟删除，也是因为堆不能直接删中间元素。

\begin{lstlisting}[language=C++,caption={LeetCode 480 滑动窗口中位数：双堆延迟删除}]
class DualHeap {
public:
    explicit DualHeap(int window_size) : window_size_(window_size) {}

    void insert(int value)
    {
        if (this->small_.empty() || value <= this->small_.top()) {
            this->small_.push(value);
            ++this->small_size_;
        } else {
            this->large_.push(value);
            ++this->large_size_;
        }
        this->rebalance();
    }

    void erase(int value)
    {
        ++this->delayed_[value];
        if (value <= this->small_.top()) {
            --this->small_size_;
            if (value == this->small_.top()) {
                this->prune_small();
            }
        } else {
            --this->large_size_;
            if (!this->large_.empty() && value == this->large_.top()) {
                this->prune_large();
            }
        }
        this->rebalance();
    }

    double median() const
    {
        if (this->window_size_ % 2 == 1) {
            return this->small_.top();
        }
        return (static_cast<double>(this->small_.top()) +
                static_cast<double>(this->large_.top())) / 2.0;
    }

private:
    void rebalance()
    {
        if (this->small_size_ > this->large_size_ + 1) {
            this->large_.push(this->small_.top());
            this->small_.pop();
            --this->small_size_;
            ++this->large_size_;
            this->prune_small();
        } else if (this->small_size_ < this->large_size_) {
            this->small_.push(this->large_.top());
            this->large_.pop();
            ++this->small_size_;
            --this->large_size_;
            this->prune_large();
        }
    }

    void prune_small()
    {
        while (!this->small_.empty()) {
            const int value = this->small_.top();
            auto found = this->delayed_.find(value);
            if (found == this->delayed_.end()) {
                return;
            }
            if (--found->second == 0) {
                this->delayed_.erase(found);
            }
            this->small_.pop();
        }
    }

    void prune_large()
    {
        while (!this->large_.empty()) {
            const int value = this->large_.top();
            auto found = this->delayed_.find(value);
            if (found == this->delayed_.end()) {
                return;
            }
            if (--found->second == 0) {
                this->delayed_.erase(found);
            }
            this->large_.pop();
        }
    }

    int window_size_;
    std::priority_queue<int> small_;
    std::priority_queue<int, std::vector<int>, std::greater<int>> large_;
    std::unordered_map<int, int> delayed_;
    int small_size_ = 0;
    int large_size_ = 0;
};
\end{lstlisting}

这段结构比普通数据流中位数复杂，因为窗口滑动时需要删除离开的元素。既然堆不能删除任意位置，就把待删除值记在 \code{delayed_} 中；只有当它出现在堆顶时才真正弹出。 \code{small_size_} 和 \code{large_size_} 记录的是有效元素数量，不是堆内部实际元素数量。成员访问使用 \code{this->} 是为了明确哪些变量属于对象状态。复杂度均摊 \complexity{O(\log k)}。如果用 \code{multiset}，删除任意元素更直接，但常数和迭代器维护方式不同；两种方案都要能解释取舍。

\code{std::map} 和 \code{std::set} 通常由红黑树实现，保证有序遍历和 \complexity{O(\log n)} 插入、删除、查找。红黑树是一种近似平衡二叉搜索树，插入删除后通过旋转和染色维持高度对数级。面试刷题里，使用 \code{map} 的信号通常是需要有序键、前驱后继、区间边界或按顺序输出。若只需要存在性或计数，\code{unordered_map} 平均更快；若哈希可能被攻击、需要稳定顺序或需要 \code{lower_bound}，就选有序树。

\begin{lstlisting}[language=C++,caption={有序映射：用 upper\_bound 找右侧区间}]
int count_covered_points(const std::map<int, int>& intervals, int point)
{
    const auto right = intervals.upper_bound(point);
    if (right == intervals.begin()) {
        return 0;
    }
    const auto candidate = std::prev(right);
    return candidate->second >= point ? 1 : 0;
}
\end{lstlisting}

\code{upper_bound(point)} 返回第一个起点大于 \code{point} 的区间，前一个区间才可能覆盖 \code{point}。这类写法是有序容器的典型优势。注意 \code{std::prev(right)} 之前必须判断 \code{right != begin()}，否则会越界。源码阅读中看到树迭代器时，要记住它不是连续内存指针，自增自减会沿树结构寻找后继前驱，复杂度通常是均摊常数或对数级细节由实现保证，但缓存局部性不如数组。

\code{unordered_map} 的平均常数复杂度不是无条件的。它依赖哈希函数把键分散到桶里，也依赖负载因子保持在合理范围。若大量键落入同一桶，查找会退化。C++ 标准没有要求具体冲突处理方式必须是链表还是开放寻址，常见实现多使用桶加节点。节点式实现的好处是 rehash 时可以重挂节点，坏处是每个元素单独分配，局部性不如连续表。刷题时使用 \code{reserve}，不只是为了速度，也是为了减少 rehash 导致的迭代器失效风险。

\begin{lstlisting}[language=C++,caption={LeetCode 560 和为 K 的子数组：查询和更新顺序不能反}]
int subarray_sum(const std::vector<int>& nums, int target)
{
    std::unordered_map<int, int> count;
    count.reserve(nums.size() * 2);
    count[0] = 1;

    int prefix = 0;
    int answer = 0;
    for (const int x : nums) {
        prefix += x;
        const auto found = count.find(prefix - target);
        if (found != count.end()) {
            answer += found->second;
        }
        ++count[prefix];
    }
    return answer;
}
\end{lstlisting}

这里必须先查询历史前缀，再加入当前前缀。若先更新，\code{target == 0} 时会把空子数组错误计入。这个例子说明容器 API 和算法不变量是绑在一起的：\code{count} 的语义是“当前下标之前出现过的前缀和次数”，不是所有前缀和次数。 \code{find} 用于查询，\code{operator[]} 用于确认要插入或累加的位置；二者角色不同。

源码阅读可以按固定路线进行。第一步读标准语义：这个容器承诺什么复杂度，哪些操作会让迭代器失效。第二步看实现的核心成员：\code{vector} 的三个指针，\code{string} 的大小容量和小字符串缓冲，\code{deque} 的段表，\code{unordered_map} 的桶数组和节点，\code{map} 的树根和节点颜色。第三步跟一条热路径：\code{push_back} 如何判断容量，\code{rehash} 如何搬迁，\code{erase} 如何维护链接，\code{lower_bound} 如何沿树下降。第四步回到题目，解释为什么某个 API 成本可接受或不可接受。

\begin{deepdive}
阅读 libstdc++ 或 libc++ 源码时，不要被模板细节吓住。可以先从头文件里的 public 函数名定位到内部辅助函数，再沿着分配、构造、移动、销毁四类操作看。比如 \code{vector::push_back} 最核心的问题只有两个：容量够不够，元素如何构造；容量不够时，才进入重新分配和移动旧元素的路径。再比如 \code{unordered_map::operator[]} 的关键语义是“没有 key 就插入默认值”，源码里会表现为查找失败后构造节点。把这些行为和题目中的不变量联系起来，源码阅读才会变成能力，而不是机械翻文件。
\end{deepdive}

\begin{labbox}
实验建议：第一，自己写一个小程序，创建 \code{vector<int>}，连续 \code{push_back} 并打印 \code{size}、\code{capacity} 和 \code{data()} 地址，观察扩容时机和地址变化。第二，写前缀和计数题，把“先查询后更新”故意改反，用 \code{target = 0} 的用例验证错误。第三，比较 \code{substr} 和下标比较在长字符串 DP 中的运行时间。第四，阅读本机标准库中 \code{vector}、\code{deque}、\code{unordered_map} 的头文件，只记录核心数据成员和迭代器失效规则，不要求理解每个模板分支。
\end{labbox}

\begin{principlebox}{DP 和容器源码阅读复盘}
DP 先问状态是否稳定，转移是否只依赖已计算状态，遍历顺序是否正确；容器先问数据是否连续、插入删除位置在哪里、查询是否需要顺序、迭代器何时失效。源码阅读不追求背实现细节，而是用实现模型解释复杂度、失效规则和 API 选择。
\end{principlebox}

\topic{代表性 C++20 代码骨架}

下面这些代码骨架不是为了替代题解，而是给读者提供可复用的实现基准。每段代码都尽量把状态命名清楚，把边界放在显式位置，把容器成本控制在题目需要的范围内。真正刷题时不要机械复制，而要先说清楚题目的不变量，再决定是否使用这些骨架。

\begin{lstlisting}[language=C++,caption={并查集：迭代 find、路径压缩和按大小合并}]
class DisjointSetUnion {
public:
    explicit DisjointSetUnion(int n)
        : parent_(n), size_(n, 1)
    {
        std::iota(this->parent_.begin(), this->parent_.end(), 0);
    }

    int find(int x)
    {
        int root = x;
        while (this->parent_[root] != root) {
            root = this->parent_[root];
        }
        while (this->parent_[x] != x) {
            const int next = this->parent_[x];
            this->parent_[x] = root;
            x = next;
        }
        return root;
    }

    bool unite(int lhs, int rhs)
    {
        int root_lhs = this->find(lhs);
        int root_rhs = this->find(rhs);
        if (root_lhs == root_rhs) {
            return false;
        }
        if (this->size_[root_lhs] < this->size_[root_rhs]) {
            std::swap(root_lhs, root_rhs);
        }
        this->parent_[root_rhs] = root_lhs;
        this->size_[root_lhs] += this->size_[root_rhs];
        return true;
    }

private:
    std::vector<int> parent_;
    std::vector<int> size_;
};
\end{lstlisting}

\begin{lstlisting}[language=C++,caption={Dijkstra：跳过 priority queue 中的过期状态}]
std::vector<long long> shortest_paths(
    int n,
    const std::vector<std::vector<std::pair<int, int>>>& graph,
    int source)
{
    constexpr long long INF = 4'000'000'000'000'000'000LL;
    using State = std::pair<long long, int>;

    std::vector<long long> distance(n, INF);
    std::priority_queue<State, std::vector<State>, std::greater<State>> heap;
    distance[source] = 0;
    heap.emplace(0, source);

    while (!heap.empty()) {
        const auto [current_distance, node] = heap.top();
        heap.pop();
        if (current_distance != distance[node]) {
            continue;
        }

        for (const auto& [next, weight] : graph[node]) {
            const long long candidate = current_distance + weight;
            if (candidate >= distance[next]) {
                continue;
            }
            distance[next] = candidate;
            heap.emplace(candidate, next);
        }
    }

    return distance;
}


\end{lstlisting}

\begin{lstlisting}[language=C++,caption={滑动窗口最大值：单调队列保存下标}]
std::vector<int> max_sliding_window(const std::vector<int>& nums, int k)
{
    std::deque<int> candidates;
    std::vector<int> answer;
    answer.reserve(nums.size());

    for (int right = 0; right < static_cast<int>(nums.size()); ++right) {
        while (!candidates.empty() && nums[candidates.back()] <= nums[right]) {
            candidates.pop_back();
        }
        candidates.push_back(right);

        const int expired = right - k;
        if (!candidates.empty() && candidates.front() <= expired) {
            candidates.pop_front();
        }
        if (right + 1 >= k) {
            answer.push_back(nums[candidates.front()]);
        }
    }

    return answer;
}
\end{lstlisting}

\begin{lstlisting}[language=C++,caption={前缀和计数：历史状态先查询后更新}]
int subarray_sum_equals_k(const std::vector<int>& nums, int target)
{
    std::unordered_map<int, int> prefix_count;
    prefix_count.reserve(nums.size() * 2 + 1);
    prefix_count[0] = 1;

    int prefix = 0;
    int answer = 0;
    for (const int x : nums) {
        prefix += x;
        const auto found = prefix_count.find(prefix - target);
        if (found != prefix_count.end()) {
            answer += found->second;
        }
        ++prefix_count[prefix];
    }
    return answer;
}
\end{lstlisting}

\begin{lstlisting}[language=C++,caption={迭代二叉树中序遍历：验证 BST 的基础骨架}]
bool is_valid_bst(TreeNode* root)
{
    std::vector<TreeNode*> stack;
    TreeNode* current = root;
    std::optional<long long> previous;

    while (current != nullptr || !stack.empty()) {
        while (current != nullptr) {
            stack.push_back(current);
            current = current->left;
        }

        current = stack.back();
        stack.pop_back();
        if (previous.has_value() && current->value <= previous.value()) {
            return false;
        }
        previous = current->value;
        current = current->right;
    }

    return true;
}
\end{lstlisting}

\section{C++ 面试代码质量和容器工程细节}

算法面试中的 C++ 代码不只是能 AC，还要让状态、所有权和复杂度可读。函数参数只读大对象应使用 const 引用，返回结果时依赖移动优化即可；循环中遍历复杂对象使用 const auto\&，不要无意复制字符串、向量或节点记录。局部常量要有名字，尤其是方向数组、无穷大、字符状态、取模常数和哨兵值。

整数类型选择要主动。数组下标和 size 返回 size\_t，但很多题需要与 -1 哨兵或差值比较，直接混用无符号可能产生隐蔽 bug。常见做法是在函数开始把 size 转成 int，并保证输入规模在 int 范围内；涉及总和、乘积、面积、距离时使用 \code{std::int64_t}。二分平方、堆距离、图最短路都要优先防溢出。

容器 API 语义要讲清。unordered\_map 的 operator[] 在 key 不存在时会插入默认值，适合计数累加，不适合只检查存在。find 返回迭代器，适合查询后读取值。contains 只判断存在，如果随后还要取值会导致二次查找。vector 的 reserve 只改容量，resize 才改大小。priority\_queue 的 top 只保证极值，不能遍历出有序序列。

迭代器失效是工程题常见追问。vector 扩容后所有指针和迭代器失效；erase 中间元素会让后续位置移动；unordered\_map rehash 会让迭代器失效；list 的 splice 能移动节点且保持其他迭代器稳定。LRU 缓存为什么用 list 加哈希表，核心就是哈希定位和链表节点稳定的组合。

代码组织要让不变量有位置。二分答案把检查函数单独写出来，图最短路把松弛逻辑集中，回溯把 push、搜索、pop 成对出现，DP 把初始化和主转移分开。这样做能让读者看到状态生命周期。把所有逻辑塞进一个巨大循环，短期可能省行数，长期会让边界和证明都变模糊。

面试时可以主动说明替代方案。第 K 大可以排序、堆、快速选择；岛屿数量可以 DFS、BFS、并查集；最小体力路径可以 Dijkstra、二分加 BFS、并查集排序边；单词接龙可以普通 BFS、通配桶、双向 BFS。比较方案不是炫耀，而是说明你知道当前选择依赖什么约束。

Linux 源码阅读放到后面的独立专题中集中学习。面试代码部分只需要先记住一个工程判断：容器选择不只看复杂度，还要看对象生命周期、内存分配、并发保护和数据是否需要被多个索引同时组织。

实验建议：写一个前缀和计数题，分别用 contains 加 operator[]、find、直接 operator[] 三种方式实现，观察语义差异。写一个 vector 扩容实验，保存元素地址后继续 push\_back，验证地址失效。写一个 LRU，先用 vector 保存顺序，再用 list 加 unordered\_map，比较移动到头部的成本。

\begin{principlebox}{本节复盘}
阅读本节后，请回到相邻章节的例题，把暴力瓶颈、优化依据、状态不变量、C++ 容器成本和边界测试重新写一遍。能把同一思想迁移到变体题，才算真正掌握。
\end{principlebox}

\topic{进阶综合题卡}

\noindent\textbf{LeetCode 4 寻找两个正序数组的中位数。}

\textbf{题目说明。} LeetCode 4 寻找两个正序数组的中位数 的题面入口要先还原成具体任务：两个有序数组的中位数可以看成左右两半切分后的边界值。

\textbf{样例入口。} 拿 [1, 3] 和 [2] 手算，合并后中位数是 2；但不用真合并，只要把总长度切成左半两个数、右半一个数。再拿 [1, 2] 和 [3, 4]，如果短数组切 1 个、长数组切 1 个，左半最大 3 大于右半最小 2，切分错了，应把短数组切分往右调。

\textbf{暴力基线。} 慢版本按顺序扫描下标、逐个试答案值，或直接检查候选直到遇到目标边界。它先保证候选空间覆盖完整，但每次只排除一个候选，浪费了题目给出的顺序、比较或可行性边界。放到本题就是：先按“两个有序数组的中位数可以看成左右两半切分后的边界值”完整试慢版本；样例里已经能看到“规律是二分短数组切分点，使左半元素个数正确且左右边界有序”，所以优化前必须先能说清慢版本枚举了哪些候选。

\textbf{暴力过程。} 拿 [1, 3] 和 [2] 手算，合并后中位数是 2；但不用真合并，只要把总长度切成左半两个数、右半一个数。再拿 [1, 2] 和 [3, 4]，如果短数组切 1 个、长数组切 1 个，左半最大 3 大于右半最小 2，切分错了，应把短数组切分往右调。

\textbf{重复工作。} 题面模型：两个有序数组的中位数可以看成左右两半切分后的边界值。暴力版的慢点要落到本题：慢版本会按顺序试“两个有序数组的中位数可以看成左右两半切分后的边界值”里的候选；而样例规律已经指出“规律是二分短数组切分点，使左半元素个数正确且左右边界有序”。一次比较或一次可行性检查能丢掉一整段候选，继续线性试就是浪费。后面的优化只消掉这类重复工作，不能改变答案集合。

\textbf{优化条件。} 本题的关键观察是：二分较短数组的切分位置，让左半元素数量正确且左半最大值不大于右半最小值。这句话必须能翻译成可维护状态；如果题目条件改变后观察不再成立，就不能继续套原来的写法。

\textbf{相邻状态。} 规律是二分短数组切分点，使左半元素个数正确且左右边界有序。把这个特例继续拆开看：每次比较 mid 之后，必须能指出哪一侧整体不可能包含目标，或哪一侧整体已经满足可行性。二分的核心不是 mid 公式，而是被丢弃区间确实不含答案。

\textbf{状态复用。} 把逐个试候选改成维护答案所在区间；边界谓词、可行性检查或局部比较给出分支方向，每个分支都必须能排除一侧候选。本题落点是：规律是二分短数组切分点，使左半元素个数正确且左右边界有序。手算时只改被新输入影响的那一小部分，而不是回头重算完整候选。

\textbf{指针和字段。} 状态字段要能说清：left、right、mid、mid 处比较值或检查返回值、答案区间含义、返回值语义；每一轮更新都要保留目标位置或第一个可行值。手算时按这个协议跟踪：拿 [1, 3] 和 [2] 手算，合并后中位数是 2；但不用真合并，只要把总长度切成左半两个数、右半一个数。再拿 [1, 2] 和 [3, 4]，如果短数组切 1 个、长数组切 1 个，左半最大 3 大于右半最小 2，切分错了，应把短数组切分往右调。然后按协议复查：手算时列出 left、mid、right，以及 mid 处比较结果、可行性检查结果或局部趋势。每一步都要写出被丢弃的一侧为什么不含答案。

\textbf{代码落点。} 检查函数或比较分支单独写清，mid 不溢出，循环退出条件和返回值配套；每个分支都要解释丢弃区间。关键代码行必须能逐句翻译成“旧状态怎样变成新状态”，否则就只是模板。

\textbf{正确性。} 证明从排除一侧开始：答案空间题证明谓词单调，局部比较题证明保留下来的区间仍含目标；不能只靠经验猜测左右边界。

\textbf{边界实验。} 先用某个数组为空、总长度奇偶、切分在边界、哨兵值不能参与真实比较做反例检查。实验时覆盖某个数组为空、总长度奇偶、切分在边界、哨兵值不能参与真实比较，逐轮打印 left、mid、right、比较值或检查结果。只要有一次被排除区间仍含答案，二分不变量就错了。

\textbf{复杂度变化。} 暴力逐个试候选是线性或和值域成正比；二分每次排除一半候选，复杂度变成对数级，答案二分还要乘上单次检查成本。

\textbf{迁移追问。} 如果题目改成“若只要求第 k 小元素，可以用递归丢弃较小前缀或二分切分的同类思想”，先判断原状态是否仍然覆盖新答案集合，再决定是微调边界、改 value 语义，还是换算法。

\noindent\textbf{LeetCode 10 正则表达式匹配。}

\textbf{题目说明。} LeetCode 10 正则表达式匹配 的题面入口要先还原成具体任务：模式中的点号匹配单字符，星号修饰前一个字符并允许出现零次或多次。

\textbf{样例入口。} 拿 s=a、p=a* 手算，星号可以让 a 出现一次；拿 s=空、p=a*，星号又可以让 a 出现零次。再拿 s=aa、p=a*，消耗一个 a 后模式仍停在 a*，还能继续匹配。

\textbf{暴力基线。} 慢版本先写递归搜索树：节点表示处理位置、剩余资源和当前选择。只要不同路径反复到达同一个节点，就说明需要把这个节点命名成 DP 状态。放到本题就是：先按“模式中的点号匹配单字符，星号修饰前一个字符并允许出现零次或多次”完整试慢版本；样例里已经能看到“规律是星号有两条分叉：当零次用时看 dp[i][j-2]，当消耗当前字符时看 dp[i-1][j]”，所以优化前必须先能说清慢版本枚举了哪些候选。

\textbf{暴力过程。} 拿 s=a、p=a* 手算，星号可以让 a 出现一次；拿 s=空、p=a*，星号又可以让 a 出现零次。再拿 s=aa、p=a*，消耗一个 a 后模式仍停在 a*，还能继续匹配。

\textbf{重复工作。} 题面模型：模式中的点号匹配单字符，星号修饰前一个字符并允许出现零次或多次。暴力版的慢点要落到本题：慢版本会在“模式中的点号匹配单字符，星号修饰前一个字符并允许出现零次或多次”的选择树里反复到达相同参数。样例规律是“规律是星号有两条分叉：当零次用时看 dp[i][j-2]，当消耗当前字符时看 dp[i-1][j]”，说明这个参数组合可以命名成状态，算过之后后面直接复用。后面的优化只消掉这类重复工作，不能改变答案集合。

\textbf{优化条件。} 本题的关键观察是：二维 DP 表示两个前缀是否匹配，星号转移同时覆盖零次和消耗一个字符后继续匹配。这句话必须能翻译成可维护状态；如果题目条件改变后观察不再成立，就不能继续套原来的写法。

\textbf{相邻状态。} 规律是星号有两条分叉：当零次用时看 dp[i][j-2]，当消耗当前字符时看 dp[i-1][j]。把这个特例继续拆开看：先问暴力搜索里哪些不同路径到达同一个后续问题，再给这个后续问题命名为状态。状态必须包含未来决策需要的全部信息，转移只从更小且已经算清的状态来。

\textbf{状态复用。} 把重复递归节点命名为状态；遍历顺序只是在保证依赖状态已经算完。本题落点是：规律是星号有两条分叉：当零次用时看 dp[i][j-2]，当消耗当前字符时看 dp[i-1][j]。手算时只改被新输入影响的那一小部分，而不是回头重算完整候选。

\textbf{指针和字段。} 状态字段要能说清：状态下标、状态值语义、初始化、转移来源、遍历顺序；空间压缩时还要指出读到的是旧值还是新值。手算时按这个协议跟踪：拿 s=a、p=a* 手算，星号可以让 a 出现一次；拿 s=空、p=a*，星号又可以让 a 出现零次。再拿 s=aa、p=a*，消耗一个 a 后模式仍停在 a*，还能继续匹配。然后按协议复查：手算时先写一个很小的 DP 表或记忆化搜索记录。每个格子旁边写它来自哪些更小状态。

\textbf{代码落点。} 先写未压缩版本校验转移；压缩前后用小表对拍；不可达哨兵参与加法前要检查溢出。关键代码行必须能逐句翻译成“旧状态怎样变成新状态”，否则就只是模板。

\textbf{正确性。} 证明从最后一步开始：任意最优解的最后一步都在转移枚举中，转移来源已经由更小状态正确计算。

\textbf{边界实验。} 先用空串、能匹配空串的星号链、星号前必须有字符、点号和普通字符的匹配条件做反例检查。实验时覆盖空串、能匹配空串的星号链、星号前必须有字符、点号和普通字符的匹配条件，先打印完整 DP 表，再比较空间压缩版本。若压缩版不同，优先检查遍历顺序和旧值保存。

\textbf{复杂度变化。} 暴力递归会重复计算相同子问题，常见指数级；DP 把每个状态算一次，时间是状态数乘单个转移成本，空间是状态表大小或压缩后的大小。

\textbf{迁移追问。} 如果题目改成“通配符匹配中的星号语义是任意长度字符串，不能直接搬用本题转移”，先判断原状态是否仍然覆盖新答案集合，再决定是微调边界、改 value 语义，还是换算法。

\noindent\textbf{LeetCode 23 合并 K 个升序链表。}

\textbf{题目说明。} LeetCode 23 合并 K 个升序链表 的题面入口要先还原成具体任务：每条链表当前头节点都是可能的下一个最小元素。

\textbf{样例入口。} 拿三条链表头 1、1、2 手算，每次答案只需要当前最小头节点；弹出某条链的头后，只有它的后继会成为新候选。暴力每轮扫 k 个头，堆把“找当前最小头”变成对数操作。

\textbf{暴力基线。} 暴力每轮扫描 K 个链表头，找到当前最小头节点接到结果后推进该链。

\textbf{暴力过程。} 若三条链当前头是 1、1、2，第一步只需从这三个头里选最小；弹出某条链的 1 后，只有它的后继会成为新候选。

\textbf{重复工作。} 题面模型：每条链表当前头节点都是可能的下一个最小元素。暴力版的慢点要落到本题：浪费在每输出一个节点都线性扫 K 个头。候选集合动态变化，但每次只需要最小头节点，适合小顶堆。后面的优化只消掉这类重复工作，不能改变答案集合。

\textbf{优化条件。} 本题的关键观察是：小顶堆保存每条非空链的当前头，弹出最小节点后把它的后继入堆。这句话必须能翻译成可维护状态；如果题目条件改变后观察不再成立，就不能继续套原来的写法。

\textbf{相邻状态。} 规律是堆里始终最多保存每条非空链的一个当前头。把这个特例继续拆开看：堆顶必须正好代表下一步最需要处理的候选；若堆里可能有过期对象，读取堆顶前要先清理。堆只解决动态极值，不自动证明被弹出或被丢弃的元素无用。

\textbf{状态复用。} 小顶堆保存每条非空链的当前头节点。弹出最小节点接入结果；若该节点有 next，把 next 入堆。手算时只改被新输入影响的那一小部分，而不是回头重算完整候选。

\textbf{指针和字段。} 状态字段要能说清：heap 中元素是 ListNode*，比较依据是 node->val；tail 是结果尾；node->next 是该链的新候选。手算时按这个协议跟踪：头节点 1、1、2 入堆；弹出第一个 1 后把它的后继入堆，堆仍保存每条链当前最小暴露节点。

\textbf{代码落点。} 关键是堆里保存节点指针而不是复制值，这样可以复用原节点。比较器要写成小顶堆语义，C++ priority\_queue 默认是大顶堆。关键代码行必须能逐句翻译成“旧状态怎样变成新状态”，否则就只是模板。

\textbf{正确性。} 每条有序链未输出部分的最小值就是当前头；所有链当前头中的最小者就是全局下一个输出节点。

\textbf{边界实验。} 先用空链表数组、某条链为空、重复值、堆里保存节点指针还是值做反例检查。实验时覆盖空链表数组、某条链为空、重复值、堆里保存节点指针还是值，记录堆顶是否过期、比较器是否让正确候选在堆顶、K 或窗口大小为边界值时是否稳定。

\textbf{复杂度变化。} 每轮扫 K 个头是 O(NK)。堆中最多 K 个节点，每个节点入堆出堆一次，总时间 O(N log K)，空间 O(K)。

\textbf{迁移追问。} 如果题目改成“也可以分治两两合并，堆更适合 K 较大且链长不均的场景”，先判断原状态是否仍然覆盖新答案集合，再决定是微调边界、改 value 语义，还是换算法。

\noindent\textbf{LeetCode 25 K 个一组翻转链表。}

\textbf{题目说明。} LeetCode 25 K 个一组翻转链表 的题面入口要先还原成具体任务：链表被切成长度为 K 的连续组，只有完整组才反转。

\textbf{样例入口。} 拿 1->2->3->4、k=2 画图。第一组边界是 1 和 2，下一组入口是 3；反转前如果不保存 3，链会断。反转后上一组尾接 2，旧头 1 接 3。

\textbf{暴力基线。} 暴力可以把每 K 个节点的值放进数组反转，再写回链表；这避开了指针重连，但没有真正按节点反转。

\textbf{暴力过程。} 以 1->2->3->4->5、k=2 为例，第一组 1、2 反转成 2->1，第二组 3、4 反转成 4->3，剩余 5 不动。

\textbf{重复工作。} 题面模型：链表被切成长度为 K 的连续组，只有完整组才反转。暴力版的慢点要落到本题：浪费在数组辅助和节点值交换。题目要求链表节点重排时，应直接反转每一组的 next 指针。后面的优化只消掉这类重复工作，不能改变答案集合。

\textbf{优化条件。} 本题的关键观察是：每轮先探测 K 个节点确认完整，再保存下一组头并做局部反转接回。这句话必须能翻译成可维护状态；如果题目条件改变后观察不再成立，就不能继续套原来的写法。

\textbf{相邻状态。} 规律是每组先找 kth 和 group\_next，再局部反转并接回前后两段。把这个特例继续拆开看：每个指针都要有角色，上一段尾、当前段头、当前节点和后继入口不能混。只要一次改 next 之前没有保存后继，就可能丢掉整段未处理链。

\textbf{状态复用。} dummy 指向头。group\_prev 是上一组尾，先从 group\_prev 出发找第 k 个节点 group\_end；不足 k 个就停止。保存 next\_group 后反转 [group\_start, group\_end]，再接回。手算时只改被新输入影响的那一小部分，而不是回头重算完整候选。

\textbf{指针和字段。} 状态字段要能说清：group\_prev 是已处理部分尾，group\_start 是本组头，group\_end 是本组尾，next\_group 是下一组入口。手算时按这个协议跟踪：第一组 start=1、end=2、next\_group=3。反转后得到 2->1，再让 group\_prev->next=2，1->next=3。

\textbf{代码落点。} 关键顺序是先保存 next\_group，再局部反转；反转结束后原 group\_start 变成本组尾，下一轮 group\_prev=group\_start。关键代码行必须能逐句翻译成“旧状态怎样变成新状态”，否则就只是模板。

\textbf{正确性。} 每次只改变完整 K 组内部指针，并把本组头尾接回已处理链和未处理链；不足 K 个节点不进入反转，符合题意。

\textbf{边界实验。} 先用不足 K 个不反转、K 为一、刚好整除、反转后头尾接错做反例检查。实验时覆盖不足 K 个不反转、K 为一、刚好整除、反转后头尾接错，每步画出前驱、当前、后继和下一段头。链表题不要只看输出值，还要检查节点是否丢失或成环。

\textbf{复杂度变化。} 数组辅助 O(n) 空间；原地分组反转每节点访问常数次，O(n) 时间、O(1) 空间。

\textbf{迁移追问。} 如果题目改成“递归写法短但可能有栈风险，迭代写法更接近工程实现”，先判断原状态是否仍然覆盖新答案集合，再决定是微调边界、改 value 语义，还是换算法。

\noindent\textbf{LeetCode 32 最长有效括号。}

\textbf{题目说明。} LeetCode 32 最长有效括号 的题面入口要先还原成具体任务：有效括号子串需要知道每段非法边界之后的最早可连接位置。

\textbf{样例入口。} 拿字符串 )()()) 手算。第一个 ) 不是可匹配对象，而是非法边界；后面每次匹配成功，当前有效长度由当前位置减去栈顶边界得出。若栈空，要把当前右括号作为新非法边界压入。

\textbf{暴力基线。} 慢版本让每个位置向左或向右线性寻找边界，或者让每个结构反复等待匹配。它能算出答案，但很多尚未完成的候选会被未来同一个事件一起解决。放到本题就是：先按“有效括号子串需要知道每段非法边界之后的最早可连接位置”完整试慢版本；样例里已经能看到“规律是栈里既保存未匹配左括号下标，也保存最近非法边界”，所以优化前必须先能说清慢版本枚举了哪些候选。

\textbf{暴力过程。} 拿字符串 )()()) 手算。第一个 ) 不是可匹配对象，而是非法边界；后面每次匹配成功，当前有效长度由当前位置减去栈顶边界得出。若栈空，要把当前右括号作为新非法边界压入。

\textbf{重复工作。} 题面模型：有效括号子串需要知道每段非法边界之后的最早可连接位置。暴力版的慢点要落到本题：慢版本会围绕“有效括号子串需要知道每段非法边界之后的最早可连接位置”反复寻找最近边界或最近未完成对象。样例规律是“规律是栈里既保存未匹配左括号下标，也保存最近非法边界”，说明旧候选一旦遇到当前输入就能被批量结算；继续为每个位置单独向左或向右扫描就是浪费。后面的优化只消掉这类重复工作，不能改变答案集合。

\textbf{优化条件。} 本题的关键观察是：栈保存下标，先放入哨兵负一；匹配后用当前下标减栈顶得到长度。这句话必须能翻译成可维护状态；如果题目条件改变后观察不再成立，就不能继续套原来的写法。

\textbf{相邻状态。} 规律是栈里既保存未匹配左括号下标，也保存最近非法边界。把这个特例继续拆开看：栈里的对象都是答案尚未确定的候选；一旦当前元素触发弹栈，被弹出的候选必须已经同时拿到了左边界和右边界，未来元素不再有资格改变它的答案。

\textbf{状态复用。} 把反复寻找最近边界改成维护未完成候选；新元素只和栈顶比较，连续弹出一批已经被当前元素解决的对象。本题落点是：规律是栈里既保存未匹配左括号下标，也保存最近非法边界。手算时只改被新输入影响的那一小部分，而不是回头重算完整候选。

\textbf{指针和字段。} 状态字段要能说清：栈内对象的业务含义、当前输入、弹出时结算的对象、左边界和右边界；栈中存下标还是值要明确。手算时按这个协议跟踪：拿字符串 )()()) 手算。第一个 ) 不是可匹配对象，而是非法边界；后面每次匹配成功，当前有效长度由当前位置减去栈顶边界得出。若栈空，要把当前右括号作为新非法边界压入。然后按协议复查：手算时记录栈内元素的业务含义，不只记录数值。入栈表示答案未定，弹栈表示当前事件已经让它的答案定下来。

\textbf{代码落点。} 弹栈前确认栈非空，弹出后再读取新的左边界；相等元素和哨兵高度要有统一策略。关键代码行必须能逐句翻译成“旧状态怎样变成新状态”，否则就只是模板。

\textbf{正确性。} 证明从弹出时机开始：被弹出的对象在当前时刻答案已经确定，未来元素无法再改变它。

\textbf{边界实验。} 先用空串、全左括号、全右括号、多个断开的合法段做反例检查。实验时准备空串、全左括号、全右括号、多个断开的合法段，打印每次入栈、弹栈和结算对象。重点观察相等元素、空栈、哨兵和最后一次结算。

\textbf{复杂度变化。} 暴力为每个位置寻找边界常见为平方级；单调栈让每个元素最多入栈一次、出栈一次，因此主体扫描为线性时间。

\textbf{迁移追问。} 如果题目改成“DP 写法也可做，但栈写法更直接表达最近非法边界”，先判断原状态是否仍然覆盖新答案集合，再决定是微调边界、改 value 语义，还是换算法。

\noindent\textbf{LeetCode 41 缺失的第一个正数。}

\textbf{题目说明。} LeetCode 41 缺失的第一个正数 的题面入口要先还原成具体任务：答案只可能落在一到 n 加一之间，数组本身可以充当值到位置的映射。

\textbf{样例入口。} 拿 [3, 4, -1, 1] 手算。答案只可能在 1 到 n+1；值 3 应放到下标 2，值 4 放到下标 3，-1 不管，1 放到下标 0。放完后第一个位置不等于 i+1 的地方就是缺失值。再试 [1, 1]，若不判断目标位置已有相同值会无限交换。

\textbf{暴力基线。} 暴力可以从 1 开始逐个查找是否出现在数组中，或者用哈希集合记录所有正数后再查。

\textbf{暴力过程。} 以 [3, 4, -1, 1] 为例，答案只可能在 1..5。把 3 放到下标 2，把 4 放到下标 3，把 1 放到下标 0，最后下标 1 缺 2。

\textbf{重复工作。} 题面模型：答案只可能落在一到 n 加一之间，数组本身可以充当值到位置的映射。暴力版的慢点要落到本题：浪费在额外集合或重复查找。长度为 n 的数组只关心 1..n 这些值，数组位置本身可以当哈希桶。后面的优化只消掉这类重复工作，不能改变答案集合。

\textbf{优化条件。} 本题的关键观察是：把合法正数 x 尽量交换到下标 x 减一，最后第一个位置不匹配就是答案。这句话必须能翻译成可维护状态；如果题目条件改变后观察不再成立，就不能继续套原来的写法。

\textbf{相邻状态。} 规律是数组本身当哈希表，只处理合法正数且避免重复交换。把这个特例继续拆开看：先标出已经确认的前缀或后缀，再标出未知区；能被丢掉的候选，必须是以后再也不会改变已确认区域语义的候选。这样才算从例子推出数组不变量，而不是只记一次操作。

\textbf{状态复用。} 扫描每个位置，把合法值 x 交换到 nums[x-1]。若目标位置已经是 x，停止交换，避免重复值死循环。手算时只改被新输入影响的那一小部分，而不是回头重算完整候选。

\textbf{指针和字段。} 状态字段要能说清：i 是当前整理位置，x=nums[i] 是试图归位的值，x-1 是它的目标下标。手算时按这个协议跟踪：[3, 4, -1, 1] 中 i=0，3 应去下标 2，换来 -1；i=1，4 应去下标 3，换来 1；1 再换到下标 0。

\textbf{代码落点。} 关键 while 条件是 1<=nums[i]<=n \&\& nums[nums[i]-1]!=nums[i]。第二遍找第一个 nums[i]!=i+1。关键代码行必须能逐句翻译成“旧状态怎样变成新状态”，否则就只是模板。

\textbf{正确性。} 每次有效交换都会至少让一个合法正数回到正确位置；最终若位置 i 不等于 i+1，则 i+1 没有被任何位置放回来，就是最小缺失正数。

\textbf{边界实验。} 先用非正数、大于 n 的数、重复值导致无限交换、数组已经完整包含一到 n做反例检查。实验时把非正数、大于 n 的数、重复值导致无限交换、数组已经完整包含一到 n列成表，再打印关键下标和有效区间。若某个元素被写入后又被未来步骤破坏，说明区间不变量没有成立。

\textbf{复杂度变化。} 哈希法 O(n) 空间；原地归位每个数交换有限次，总时间 O(n)，额外空间 O(1)。

\textbf{迁移追问。} 如果题目改成“若不能修改输入，只能使用哈希集合或额外布尔数组换取更简单边界”，先判断原状态是否仍然覆盖新答案集合，再决定是微调边界、改 value 语义，还是换算法。

\noindent\textbf{LeetCode 44 通配符匹配。}

\textbf{题目说明。} LeetCode 44 通配符匹配 的题面入口要先还原成具体任务：问号匹配一个字符，星号匹配任意长度字符串，状态是两个前缀能否匹配。

\textbf{样例入口。} 拿 s=ab、p=a* 手算，星号可以吞掉 b；拿 s=空、p=*，星号也可以代表空。DP 中星号的两条路是匹配空串或继续吞当前字符。贪心中失配时回到最近星号，让它多吞一个字符。

\textbf{暴力基线。} 慢版本先写递归搜索树：节点表示处理位置、剩余资源和当前选择。只要不同路径反复到达同一个节点，就说明需要把这个节点命名成 DP 状态。放到本题就是：先按“问号匹配一个字符，星号匹配任意长度字符串，状态是两个前缀能否匹配”完整试慢版本；样例里已经能看到“规律是本题星号独立代表任意串，和正则里修饰前一字符的星号不同”，所以优化前必须先能说清慢版本枚举了哪些候选。

\textbf{暴力过程。} 拿 s=ab、p=a* 手算，星号可以吞掉 b；拿 s=空、p=*，星号也可以代表空。DP 中星号的两条路是匹配空串或继续吞当前字符。贪心中失配时回到最近星号，让它多吞一个字符。

\textbf{重复工作。} 题面模型：问号匹配一个字符，星号匹配任意长度字符串，状态是两个前缀能否匹配。暴力版的慢点要落到本题：慢版本会在“问号匹配一个字符，星号匹配任意长度字符串，状态是两个前缀能否匹配”的选择树里反复到达相同参数。样例规律是“规律是本题星号独立代表任意串，和正则里修饰前一字符的星号不同”，说明这个参数组合可以命名成状态，算过之后后面直接复用。后面的优化只消掉这类重复工作，不能改变答案集合。

\textbf{优化条件。} 本题的关键观察是：DP 中星号可匹配空串或继续吞掉当前字符；贪心写法记录最近星号并在失配时回退扩展星号覆盖范围。这句话必须能翻译成可维护状态；如果题目条件改变后观察不再成立，就不能继续套原来的写法。

\textbf{相邻状态。} 规律是本题星号独立代表任意串，和正则里修饰前一字符的星号不同。把这个特例继续拆开看：先问暴力搜索里哪些不同路径到达同一个后续问题，再给这个后续问题命名为状态。状态必须包含未来决策需要的全部信息，转移只从更小且已经算清的状态来。

\textbf{状态复用。} 把重复递归节点命名为状态；遍历顺序只是在保证依赖状态已经算完。本题落点是：规律是本题星号独立代表任意串，和正则里修饰前一字符的星号不同。手算时只改被新输入影响的那一小部分，而不是回头重算完整候选。

\textbf{指针和字段。} 状态字段要能说清：状态下标、状态值语义、初始化、转移来源、遍历顺序；空间压缩时还要指出读到的是旧值还是新值。手算时按这个协议跟踪：拿 s=ab、p=a* 手算，星号可以吞掉 b；拿 s=空、p=*，星号也可以代表空。DP 中星号的两条路是匹配空串或继续吞当前字符。贪心中失配时回到最近星号，让它多吞一个字符。然后按协议复查：手算时先写一个很小的 DP 表或记忆化搜索记录。每个格子旁边写它来自哪些更小状态。

\textbf{代码落点。} 先写未压缩版本校验转移；压缩前后用小表对拍；不可达哨兵参与加法前要检查溢出。关键代码行必须能逐句翻译成“旧状态怎样变成新状态”，否则就只是模板。

\textbf{正确性。} 证明从最后一步开始：任意最优解的最后一步都在转移枚举中，转移来源已经由更小状态正确计算。

\textbf{边界实验。} 先用连续星号、空串、模式只含星号、贪心回退时字符串位置不能倒退错做反例检查。实验时覆盖连续星号、空串、模式只含星号、贪心回退时字符串位置不能倒退错，先打印完整 DP 表，再比较空间压缩版本。若压缩版不同，优先检查遍历顺序和旧值保存。

\textbf{复杂度变化。} 暴力递归会重复计算相同子问题，常见指数级；DP 把每个状态算一次，时间是状态数乘单个转移成本，空间是状态表大小或压缩后的大小。

\textbf{迁移追问。} 如果题目改成“正则表达式匹配的星号修饰前一个字符，和本题星号独立匹配任意串不同”，先判断原状态是否仍然覆盖新答案集合，再决定是微调边界、改 value 语义，还是换算法。

\noindent\textbf{LeetCode 51 N 皇后。}

\textbf{题目说明。} LeetCode 51 N 皇后 的题面入口要先还原成具体任务：每一行放一个皇后，列和两条对角线不能冲突。

\textbf{样例入口。} 拿 n=4 手算，第一行放某列后，下一行不能放同列、主对角线和副对角线冲突的位置。

\textbf{暴力基线。} 慢版本就是完整搜索树：每层列出选择列表，路径到达终止条件时检查答案。它先保证覆盖所有可能，再把明显无效或重复的分支剪掉。放到本题就是：先按“每一行放一个皇后，列和两条对角线不能冲突”完整试慢版本；样例里已经能看到“规律是按行递归，每行放一个皇后，三个集合分别记录列、row-col、row+col 是否被占用”，所以优化前必须先能说清慢版本枚举了哪些候选。

\textbf{暴力过程。} 拿 n=4 手算，第一行放某列后，下一行不能放同列、主对角线和副对角线冲突的位置。

\textbf{重复工作。} 题面模型：每一行放一个皇后，列和两条对角线不能冲突。暴力版的慢点要落到本题：慢版本会围绕“每一行放一个皇后，列和两条对角线不能冲突”展开完整搜索树。样例规律是“规律是按行递归，每行放一个皇后，三个集合分别记录列、row-col、row+col 是否被占用”，说明一层递归必须对应一个明确决策，剪枝只能删除整棵不可能成功或重复的子树。后面的优化只消掉这类重复工作，不能改变答案集合。

\textbf{优化条件。} 本题的关键观察是：逐行搜索，用列集合、主对角线集合和副对角线集合快速判断合法性。这句话必须能翻译成可维护状态；如果题目条件改变后观察不再成立，就不能继续套原来的写法。

\textbf{相邻状态。} 规律是按行递归，每行放一个皇后，三个集合分别记录列、row-col、row+col 是否被占用。把这个特例继续拆开看：一层递归对应一个明确决策，路径保存已经做出的选择，选择列表保存仍可能扩展的分支。剪枝前必须能说明被剪掉的整棵子树没有合法答案，而不是只因为它看起来不优。

\textbf{状态复用。} 把完整搜索树加上合法性检查、剩余资源剪枝和同层去重；路径状态进入和退出要成对恢复。本题落点是：规律是按行递归，每行放一个皇后，三个集合分别记录列、row-col、row+col 是否被占用。手算时只改被新输入影响的那一小部分，而不是回头重算完整候选。

\textbf{指针和字段。} 状态字段要能说清：路径、起点或使用标记、剩余目标、答案集合、剪枝条件；进入递归和退出递归必须成对恢复。手算时按这个协议跟踪：拿 n=4 手算，第一行放某列后，下一行不能放同列、主对角线和副对角线冲突的位置。然后按协议复查：手算时给每层标出选择列表和路径。剪枝发生前要指出被剪分支为什么已经不可能得到合法答案。

\textbf{代码落点。} 剪枝条件写在扩展前；同层去重不要误写成全局去重；保存答案时复制当前路径。关键代码行必须能逐句翻译成“旧状态怎样变成新状态”，否则就只是模板。

\textbf{正确性。} 证明从搜索树覆盖开始：每个合法答案有且只有一条路径，剪枝只删掉不可能成功的分支。

\textbf{边界实验。} 先用n 为一、n 为二或三无解、对角线编号、回退时状态恢复做反例检查。实验时覆盖n 为一、n 为二或三无解、对角线编号、回退时状态恢复，统计搜索节点数量和剪枝次数。重复元素题要打印同一层跳过哪些值，确认没有误删合法路径。

\textbf{复杂度变化。} 暴力搜索树的规模由输出或选择空间决定；剪枝不改变最坏指数级本质，但能删除不可能成功或重复的分支，复杂度要说明输出规模。

\textbf{迁移追问。} 如果题目改成“可用位掩码压缩三个约束集合，提高常数性能”，先判断原状态是否仍然覆盖新答案集合，再决定是微调边界、改 value 语义，还是换算法。

\noindent\textbf{LeetCode 72 编辑距离。}

\textbf{题目说明。} LeetCode 72 编辑距离 的题面入口要先还原成具体任务：二维状态表示两个前缀之间的最少编辑操作。

\textbf{样例入口。} 拿 word1=ab、word2=ac 手算最后一个字符。若 b 改成 c，是 dp[1][1]+1；若删掉 b，是 dp[1][2]+1；若插入 c，是 dp[2][1]+1。字符相等时才直接继承左上。

\textbf{暴力基线。} 暴力递归比较两个字符串后缀：相等就同时跳过，不等就尝试插入、删除、替换三种操作。

\textbf{暴力过程。} 以 word1=ab、word2=ac 为例，末尾 b 和 c 不同。替换 b 成 c 来自 dp[1][1]+1；删除 b 来自 dp[1][2]+1；插入 c 来自 dp[2][1]+1。

\textbf{重复工作。} 题面模型：二维状态表示两个前缀之间的最少编辑操作。暴力版的慢点要落到本题：浪费在不同编辑序列会反复到达同一对前缀长度。只要 i 和 j 相同，后续最少编辑次数就相同。后面的优化只消掉这类重复工作，不能改变答案集合。

\textbf{优化条件。} 本题的关键观察是：末尾字符相同继承左上；不同则枚举插入、删除、替换。这句话必须能翻译成可维护状态；如果题目条件改变后观察不再成立，就不能继续套原来的写法。

\textbf{相邻状态。} 规律是每个格子枚举最后一次编辑操作，而不是凭感觉填三邻居最小。把这个特例继续拆开看：先问暴力搜索里哪些不同路径到达同一个后续问题，再给这个后续问题命名为状态。状态必须包含未来决策需要的全部信息，转移只从更小且已经算清的状态来。

\textbf{状态复用。} dp[i][j] 表示 word1 前 i 个字符变成 word2 前 j 个字符的最少操作数。按前缀长度从小到大填表。手算时只改被新输入影响的那一小部分，而不是回头重算完整候选。

\textbf{指针和字段。} 状态字段要能说清：i/j 是前缀长度；dp[i-1][j] 对应删除，dp[i][j-1] 对应插入，dp[i-1][j-1] 对应替换或字符相等继承。手算时按这个协议跟踪：空串到 ac 需要插入 2 次；ab 到空串需要删除 2 次。内部格子按末尾字符是否相等决定转移。

\textbf{代码落点。} 关键是第一行 dp[0][j]=j、第一列 dp[i][0]=i。字符相等时不要额外加一。关键代码行必须能逐句翻译成“旧状态怎样变成新状态”，否则就只是模板。

\textbf{正确性。} 任意最优编辑序列的最后一步只有插入、删除、替换或不操作四类，转移枚举了全部最后一步。

\textbf{边界实验。} 先用空串、完全相同、完全不同、操作代价不同做反例检查。实验时覆盖空串、完全相同、完全不同、操作代价不同，先打印完整 DP 表，再比较空间压缩版本。若压缩版不同，优先检查遍历顺序和旧值保存。

\textbf{复杂度变化。} 暴力递归指数级；二维 DP O(mn) 时间和空间，可滚动到 O(min(m, n)) 空间。

\textbf{迁移追问。} 如果题目改成“若只允许插入删除，问题退化到和最长公共子序列相关”，先判断原状态是否仍然覆盖新答案集合，再决定是微调边界、改 value 语义，还是换算法。

\noindent\textbf{LeetCode 85 最大矩形。}

\textbf{题目说明。} LeetCode 85 最大矩形 的题面入口要先还原成具体任务：每一行都可以把连续的一当成柱状图高度。

\textbf{样例入口。} 拿矩阵每一行向上累计高度手算，某行高度数组若为 [2, 1, 2]，就退化成柱状图最大矩形。

\textbf{暴力基线。} 慢版本让每个位置向左或向右线性寻找边界，或者让每个结构反复等待匹配。它能算出答案，但很多尚未完成的候选会被未来同一个事件一起解决。放到本题就是：先按“每一行都可以把连续的一当成柱状图高度”完整试慢版本；样例里已经能看到“规律是逐行更新每列连续 1 的高度，再对每行高度数组跑单调栈；全零行会把高度清零”，所以优化前必须先能说清慢版本枚举了哪些候选。

\textbf{暴力过程。} 拿矩阵每一行向上累计高度手算，某行高度数组若为 [2, 1, 2]，就退化成柱状图最大矩形。

\textbf{重复工作。} 题面模型：每一行都可以把连续的一当成柱状图高度。暴力版的慢点要落到本题：慢版本会围绕“每一行都可以把连续的一当成柱状图高度”反复寻找最近边界或最近未完成对象。样例规律是“规律是逐行更新每列连续 1 的高度，再对每行高度数组跑单调栈；全零行会把高度清零”，说明旧候选一旦遇到当前输入就能被批量结算；继续为每个位置单独向左或向右扫描就是浪费。后面的优化只消掉这类重复工作，不能改变答案集合。

\textbf{优化条件。} 本题的关键观察是：逐行更新每列高度，然后把当前行转成柱状图最大矩形问题。这句话必须能翻译成可维护状态；如果题目条件改变后观察不再成立，就不能继续套原来的写法。

\textbf{相邻状态。} 规律是逐行更新每列连续 1 的高度，再对每行高度数组跑单调栈；全零行会把高度清零。把这个特例继续拆开看：栈里的对象都是答案尚未确定的候选；一旦当前元素触发弹栈，被弹出的候选必须已经同时拿到了左边界和右边界，未来元素不再有资格改变它的答案。

\textbf{状态复用。} 把反复寻找最近边界改成维护未完成候选；新元素只和栈顶比较，连续弹出一批已经被当前元素解决的对象。本题落点是：规律是逐行更新每列连续 1 的高度，再对每行高度数组跑单调栈；全零行会把高度清零。手算时只改被新输入影响的那一小部分，而不是回头重算完整候选。

\textbf{指针和字段。} 状态字段要能说清：栈内对象的业务含义、当前输入、弹出时结算的对象、左边界和右边界；栈中存下标还是值要明确。手算时按这个协议跟踪：拿矩阵每一行向上累计高度手算，某行高度数组若为 [2, 1, 2]，就退化成柱状图最大矩形。然后按协议复查：手算时记录栈内元素的业务含义，不只记录数值。入栈表示答案未定，弹栈表示当前事件已经让它的答案定下来。

\textbf{代码落点。} 弹栈前确认栈非空，弹出后再读取新的左边界；相等元素和哨兵高度要有统一策略。关键代码行必须能逐句翻译成“旧状态怎样变成新状态”，否则就只是模板。

\textbf{正确性。} 证明从弹出时机开始：被弹出的对象在当前时刻答案已经确定，未来元素无法再改变它。

\textbf{边界实验。} 先用全零、全一、单行、单列、宽度和高度结算做反例检查。实验时准备全零、全一、单行、单列、宽度和高度结算，打印每次入栈、弹栈和结算对象。重点观察相等元素、空栈、哨兵和最后一次结算。

\textbf{复杂度变化。} 暴力为每个位置寻找边界常见为平方级；单调栈让每个元素最多入栈一次、出栈一次，因此主体扫描为线性时间。

\textbf{迁移追问。} 如果题目改成“这是二维问题压成一维单调栈的典型做法”，先判断原状态是否仍然覆盖新答案集合，再决定是微调边界、改 value 语义，还是换算法。

\noindent\textbf{LeetCode 97 交错字符串。}

\textbf{题目说明。} LeetCode 97 交错字符串 的题面入口要先还原成具体任务：目标串前缀由两个来源字符串的前缀交错组成。

\textbf{样例入口。} 拿 s1=a、s2=b、s3=ab 手算，最后一个字符 b 可以来自 s2；拿 s1=a、s2=a、s3=aa，两个来源都能匹配，不能贪心固定选一个。

\textbf{暴力基线。} 慢版本先写递归搜索树：节点表示处理位置、剩余资源和当前选择。只要不同路径反复到达同一个节点，就说明需要把这个节点命名成 DP 状态。放到本题就是：先按“目标串前缀由两个来源字符串的前缀交错组成”完整试慢版本；样例里已经能看到“规律是 dp[i][j] 表示两个前缀能否组成目标前 i+j 个字符，转移同时尝试来自 s1 和 s2 的最后一步”，所以优化前必须先能说清慢版本枚举了哪些候选。

\textbf{暴力过程。} 拿 s1=a、s2=b、s3=ab 手算，最后一个字符 b 可以来自 s2；拿 s1=a、s2=a、s3=aa，两个来源都能匹配，不能贪心固定选一个。

\textbf{重复工作。} 题面模型：目标串前缀由两个来源字符串的前缀交错组成。暴力版的慢点要落到本题：慢版本会在“目标串前缀由两个来源字符串的前缀交错组成”的选择树里反复到达相同参数。样例规律是“规律是 dp[i][j] 表示两个前缀能否组成目标前 i+j 个字符，转移同时尝试来自 s1 和 s2 的最后一步”，说明这个参数组合可以命名成状态，算过之后后面直接复用。后面的优化只消掉这类重复工作，不能改变答案集合。

\textbf{优化条件。} 本题的关键观察是：状态 dp[i][j] 表示 s1 前 i 个字符和 s2 前 j 个字符能否组成 s3 前 i 加 j 个字符。这句话必须能翻译成可维护状态；如果题目条件改变后观察不再成立，就不能继续套原来的写法。

\textbf{相邻状态。} 规律是 dp[i][j] 表示两个前缀能否组成目标前 i+j 个字符，转移同时尝试来自 s1 和 s2 的最后一步。把这个特例继续拆开看：先问暴力搜索里哪些不同路径到达同一个后续问题，再给这个后续问题命名为状态。状态必须包含未来决策需要的全部信息，转移只从更小且已经算清的状态来。

\textbf{状态复用。} 把重复递归节点命名为状态；遍历顺序只是在保证依赖状态已经算完。本题落点是：规律是 dp[i][j] 表示两个前缀能否组成目标前 i+j 个字符，转移同时尝试来自 s1 和 s2 的最后一步。手算时只改被新输入影响的那一小部分，而不是回头重算完整候选。

\textbf{指针和字段。} 状态字段要能说清：状态下标、状态值语义、初始化、转移来源、遍历顺序；空间压缩时还要指出读到的是旧值还是新值。手算时按这个协议跟踪：拿 s1=a、s2=b、s3=ab 手算，最后一个字符 b 可以来自 s2；拿 s1=a、s2=a、s3=aa，两个来源都能匹配，不能贪心固定选一个。然后按协议复查：手算时先写一个很小的 DP 表或记忆化搜索记录。每个格子旁边写它来自哪些更小状态。

\textbf{代码落点。} 先写未压缩版本校验转移；压缩前后用小表对拍；不可达哨兵参与加法前要检查溢出。关键代码行必须能逐句翻译成“旧状态怎样变成新状态”，否则就只是模板。

\textbf{正确性。} 证明从最后一步开始：任意最优解的最后一步都在转移枚举中，转移来源已经由更小状态正确计算。

\textbf{边界实验。} 先用总长度不等、第一行第一列初始化、两个来源字符都能匹配时不能贪心只选一个做反例检查。实验时覆盖总长度不等、第一行第一列初始化、两个来源字符都能匹配时不能贪心只选一个，先打印完整 DP 表，再比较空间压缩版本。若压缩版不同，优先检查遍历顺序和旧值保存。

\textbf{复杂度变化。} 暴力递归会重复计算相同子问题，常见指数级；DP 把每个状态算一次，时间是状态数乘单个转移成本，空间是状态表大小或压缩后的大小。

\textbf{迁移追问。} 如果题目改成“若要统计交错方案数，布尔状态要改成计数状态并处理大数或取模”，先判断原状态是否仍然覆盖新答案集合，再决定是微调边界、改 value 语义，还是换算法。

\noindent\textbf{LeetCode 115 不同的子序列。}

\textbf{题目说明。} LeetCode 115 不同的子序列 的题面入口要先还原成具体任务：从源串中删除若干字符后形成目标串，本质是两个前缀之间的计数问题。

\textbf{样例入口。} 拿 s=bab、t=ba 手算。扫描到最后一个 b 时，它不能匹配 t 的 a，只能继承不用它的方案；扫描到 a 时，有两类方案：用这个 a 匹配 t 的末尾，或不用它。

\textbf{暴力基线。} 慢版本先写递归搜索树：节点表示处理位置、剩余资源和当前选择。只要不同路径反复到达同一个节点，就说明需要把这个节点命名成 DP 状态。放到本题就是：先按“从源串中删除若干字符后形成目标串，本质是两个前缀之间的计数问题”完整试慢版本；样例里已经能看到“规律是字符相等时计数相加，不等时只继承不用源字符的方案”，所以优化前必须先能说清慢版本枚举了哪些候选。

\textbf{暴力过程。} 拿 s=bab、t=ba 手算。扫描到最后一个 b 时，它不能匹配 t 的 a，只能继承不用它的方案；扫描到 a 时，有两类方案：用这个 a 匹配 t 的末尾，或不用它。

\textbf{重复工作。} 题面模型：从源串中删除若干字符后形成目标串，本质是两个前缀之间的计数问题。暴力版的慢点要落到本题：慢版本会在“从源串中删除若干字符后形成目标串，本质是两个前缀之间的计数问题”的选择树里反复到达相同参数。样例规律是“规律是字符相等时计数相加，不等时只继承不用源字符的方案”，说明这个参数组合可以命名成状态，算过之后后面直接复用。后面的优化只消掉这类重复工作，不能改变答案集合。

\textbf{优化条件。} 本题的关键观察是：状态 dp[i][j] 表示源串前 i 个字符中有多少个子序列等于目标前 j 个字符。这句话必须能翻译成可维护状态；如果题目条件改变后观察不再成立，就不能继续套原来的写法。

\textbf{相邻状态。} 规律是字符相等时计数相加，不等时只继承不用源字符的方案。把这个特例继续拆开看：先问暴力搜索里哪些不同路径到达同一个后续问题，再给这个后续问题命名为状态。状态必须包含未来决策需要的全部信息，转移只从更小且已经算清的状态来。

\textbf{状态复用。} 把重复递归节点命名为状态；遍历顺序只是在保证依赖状态已经算完。本题落点是：规律是字符相等时计数相加，不等时只继承不用源字符的方案。手算时只改被新输入影响的那一小部分，而不是回头重算完整候选。

\textbf{指针和字段。} 状态字段要能说清：状态下标、状态值语义、初始化、转移来源、遍历顺序；空间压缩时还要指出读到的是旧值还是新值。手算时按这个协议跟踪：拿 s=bab、t=ba 手算。扫描到最后一个 b 时，它不能匹配 t 的 a，只能继承不用它的方案；扫描到 a 时，有两类方案：用这个 a 匹配 t 的末尾，或不用它。然后按协议复查：手算时先写一个很小的 DP 表或记忆化搜索记录。每个格子旁边写它来自哪些更小状态。

\textbf{代码落点。} 先写未压缩版本校验转移；压缩前后用小表对拍；不可达哨兵参与加法前要检查溢出。关键代码行必须能逐句翻译成“旧状态怎样变成新状态”，否则就只是模板。

\textbf{正确性。} 证明从最后一步开始：任意最优解的最后一步都在转移枚举中，转移来源已经由更小状态正确计算。

\textbf{边界实验。} 先用空目标串计数为一、源串为空、重复字符导致计数增长、计数可能超过 int做反例检查。实验时覆盖空目标串计数为一、源串为空、重复字符导致计数增长、计数可能超过 int，先打印完整 DP 表，再比较空间压缩版本。若压缩版不同，优先检查遍历顺序和旧值保存。

\textbf{复杂度变化。} 暴力递归会重复计算相同子问题，常见指数级；DP 把每个状态算一次，时间是状态数乘单个转移成本，空间是状态表大小或压缩后的大小。

\textbf{迁移追问。} 如果题目改成“若只问是否存在目标子序列，可以用双指针，不需要二维计数 DP”，先判断原状态是否仍然覆盖新答案集合，再决定是微调边界、改 value 语义，还是换算法。

\noindent\textbf{LeetCode 126 单词接龙 II。}

\textbf{题目说明。} LeetCode 126 单词接龙 II 的题面入口要先还原成具体任务：不仅要最短转换长度，还要输出所有最短转换路径。

\textbf{样例入口。} 拿 hit 到 cog 的小词表手算，同一层 hot 可以通向 dot 和 lot，后续 dog/log 都可能到 cog。若某个单词在本层第一次见到就立刻删掉，会丢掉另一条同长度父路径。

\textbf{暴力基线。} 慢版本先按题意画出完整状态图，用普通搜索跑通可达性。这个过程用于确认不仅要最短转换长度，还要输出所有最短转换路径的节点和边，之后再讨论访问标记、层次或代价顺序。放到本题就是：先按“不仅要最短转换长度，还要输出所有最短转换路径”完整试慢版本；样例里已经能看到“规律是 BFS 建最短层级，前驱列表保留同层多父节点，最后再回溯所有路径”，所以优化前必须先能说清慢版本枚举了哪些候选。

\textbf{暴力过程。} 拿 hit 到 cog 的小词表手算，同一层 hot 可以通向 dot 和 lot，后续 dog/log 都可能到 cog。若某个单词在本层第一次见到就立刻删掉，会丢掉另一条同长度父路径。

\textbf{重复工作。} 题面模型：不仅要最短转换长度，还要输出所有最短转换路径。暴力版的慢点要落到本题：慢版本会围绕“不仅要最短转换长度，还要输出所有最短转换路径”反复展开同一批状态。样例规律是“规律是 BFS 建最短层级，前驱列表保留同层多父节点，最后再回溯所有路径”，说明必须把节点、边、访问标记和资源维度一次建清；同一状态不应重复入队，不同资源状态也不能误合并。后面的优化只消掉这类重复工作，不能改变答案集合。

\textbf{优化条件。} 本题的关键观察是：BFS 只建立最短层级和前驱列表，同一层的多个前驱都保留，最后再从终点回溯路径。这句话必须能翻译成可维护状态；如果题目条件改变后观察不再成立，就不能继续套原来的写法。

\textbf{相邻状态。} 规律是 BFS 建最短层级，前驱列表保留同层多父节点，最后再回溯所有路径。把这个特例继续拆开看：状态节点不一定等于原始位置，还可能包含钥匙、剩余资源、时间或访问集合。visited 只能合并未来能力完全相同的状态，否则会把合法路径剪掉。

\textbf{状态复用。} 把重复搜索压成访问标记、距离数组或拓扑状态；邻居生成也要从全量扫描变成结构化枚举。本题落点是：规律是 BFS 建最短层级，前驱列表保留同层多父节点，最后再回溯所有路径。手算时只改被新输入影响的那一小部分，而不是回头重算完整候选。

\textbf{指针和字段。} 状态字段要能说清：状态编号、邻接生成方式、访问标记、距离或层数、队列内容；若状态含资源，visited 不能只按位置记录。手算时按这个协议跟踪：拿 hit 到 cog 的小词表手算，同一层 hot 可以通向 dot 和 lot，后续 dog/log 都可能到 cog。若某个单词在本层第一次见到就立刻删掉，会丢掉另一条同长度父路径。然后按协议复查：手算时画队列或栈，状态必须包含题目需要的所有资源。入队时标记还是出队时标记，要和去重语义一致。

\textbf{代码落点。} 入队时标记还是出队时标记必须一致；方向数组集中定义；多源 BFS 要一次性把所有源点放入初始队列。关键代码行必须能逐句翻译成“旧状态怎样变成新状态”，否则就只是模板。

\textbf{正确性。} 证明从可达性开始：搜索覆盖所有合法边能到达的状态，访问标记只删除重复状态而不删除不同未来能力。

\textbf{边界实验。} 先用终点不在词表、同层多父节点、本层访问不能提前删除等长前驱、输出规模可能很大做反例检查。实验时覆盖终点不在词表、同层多父节点、本层访问不能提前删除等长前驱、输出规模可能很大，并打印入队时的完整状态。若同一位置但资源不同的状态被合并，很多图题会出现隐蔽错误。

\textbf{复杂度变化。} 暴力重复从多个起点搜索会反复遍历图；正确建模和访问标记后，BFS 或 DFS 通常按节点数加边数计算，状态图题则按完整状态数计算。

\textbf{迁移追问。} 如果题目改成“若只问最短长度，普通 BFS 或双向 BFS 更轻；若要路径，就必须保存前驱关系”，先判断原状态是否仍然覆盖新答案集合，再决定是微调边界、改 value 语义，还是换算法。

\noindent\textbf{LeetCode 135 分发糖果。}

\textbf{题目说明。} LeetCode 135 分发糖果 的题面入口要先还原成具体任务：相邻评分约束需要同时满足左邻和右邻两个方向。

\textbf{样例入口。} 拿 ratings=[1, 0, 2] 手算，中间孩子评分最低，左右都要比它多糖，结果 [2, 1, 2]。只从左到右会漏掉右侧约束。

\textbf{暴力基线。} 慢版本枚举选择顺序或用 DP 保留多个可能方案，先确认最优值存在。随后才检查局部选择是否能通过交换或领先性固定下来。放到本题就是：先按“相邻评分约束需要同时满足左邻和右邻两个方向”完整试慢版本；样例里已经能看到“规律是左到右满足左邻递增，右到左满足右邻递增，最终取两侧要求最大值”，所以优化前必须先能说清慢版本枚举了哪些候选。

\textbf{暴力过程。} 拿 ratings=[1, 0, 2] 手算，中间孩子评分最低，左右都要比它多糖，结果 [2, 1, 2]。只从左到右会漏掉右侧约束。

\textbf{重复工作。} 题面模型：相邻评分约束需要同时满足左邻和右邻两个方向。暴力版的慢点要落到本题：慢版本会围绕“相邻评分约束需要同时满足左邻和右邻两个方向”枚举选择顺序。样例规律是“规律是左到右满足左邻递增，右到左满足右邻递增，最终取两侧要求最大值”，说明存在一个局部选择或边界状态可以安全保留；不能证明这个选择不变差时，就不能把它叫贪心。后面的优化只消掉这类重复工作，不能改变答案集合。

\textbf{优化条件。} 本题的关键观察是：左到右满足递增关系，右到左满足反向递增关系，最后取两侧要求最大值。这句话必须能翻译成可维护状态；如果题目条件改变后观察不再成立，就不能继续套原来的写法。

\textbf{相邻状态。} 规律是左到右满足左邻递增，右到左满足右邻递增，最终取两侧要求最大值。把这个特例继续拆开看：先构造一个非贪心选择，再比较两者留下的资源、右边界或剩余时间。只有能把非贪心第一步交换成贪心第一步且不变差，局部选择才是规律。

\textbf{状态复用。} 把指数级选择压成一个可证明安全的局部选择；状态通常是当前最远边界、最早结束时间、最小代价或剩余资源。本题落点是：规律是左到右满足左邻递增，右到左满足右邻递增，最终取两侧要求最大值。手算时只改被新输入影响的那一小部分，而不是回头重算完整候选。

\textbf{指针和字段。} 状态字段要能说清：处理顺序、当前已选方案、剩余资源或最远边界、答案计数；若使用排序，要说明排序键；每次局部选择后要能说明当前状态不落后。手算时按这个协议跟踪：拿 ratings=[1, 0, 2] 手算，中间孩子评分最低，左右都要比它多糖，结果 [2, 1, 2]。只从左到右会漏掉右侧约束。然后按协议复查：手算时同时写一个非贪心选择，与贪心选择比较剩余资源。若无法说明贪心后剩余资源不差，就不能直接下结论。

\textbf{代码落点。} 把局部规则写成业务含义；若需要排序，就处理相同关键字；循环中只维护证明需要的状态，不把局部选择和答案更新混在一起。关键代码行必须能逐句翻译成“旧状态怎样变成新状态”，否则就只是模板。

\textbf{正确性。} 证明从交换或领先性开始：任意最优解都能替换成当前贪心选择而不变差，或贪心状态始终不落后。

\textbf{边界实验。} 先用评分相等、严格递增、严格递减、山峰和谷底做反例检查。实验时除了评分相等、严格递增、严格递减、山峰和谷底，还要故意替换成一个看似合理但错误的局部规则，构造反例。能解释错误贪心为何失败，才算理解正确贪心。

\textbf{复杂度变化。} 暴力枚举选择顺序可能指数级；贪心通常先排序，再线性扫描或用堆维护少量候选。

\textbf{迁移追问。} 如果题目改成“这是局部约束两遍扫描，不能只用一个方向贪心”，先判断原状态是否仍然覆盖新答案集合，再决定是微调边界、改 value 语义，还是换算法。

\noindent\textbf{LeetCode 140 单词拆分 II。}

\textbf{题目说明。} LeetCode 140 单词拆分 II 的题面入口要先还原成具体任务：需要输出所有由字典单词拼出的句子，输出规模本身可能指数级。

\textbf{样例入口。} 拿 catsanddog 手算，前缀 cat 和 cats 都能走，但某些后缀继续切会失败。若直接 DFS，每次都会重复探索同一个后缀。

\textbf{暴力基线。} 慢版本就是完整搜索树：每层列出选择列表，路径到达终止条件时检查答案。它先保证覆盖所有可能，再把明显无效或重复的分支剪掉。放到本题就是：先按“需要输出所有由字典单词拼出的句子，输出规模本身可能指数级”完整试慢版本；样例里已经能看到“规律是先用记忆化记录某个起点之后能生成哪些句子，或先用 DP 剪掉不可拆后缀，再枚举路径”，所以优化前必须先能说清慢版本枚举了哪些候选。

\textbf{暴力过程。} 拿 catsanddog 手算，前缀 cat 和 cats 都能走，但某些后缀继续切会失败。若直接 DFS，每次都会重复探索同一个后缀。

\textbf{重复工作。} 题面模型：需要输出所有由字典单词拼出的句子，输出规模本身可能指数级。暴力版的慢点要落到本题：慢版本会围绕“需要输出所有由字典单词拼出的句子，输出规模本身可能指数级”展开完整搜索树。样例规律是“规律是先用记忆化记录某个起点之后能生成哪些句子，或先用 DP 剪掉不可拆后缀，再枚举路径”，说明一层递归必须对应一个明确决策，剪枝只能删除整棵不可能成功或重复的子树。后面的优化只消掉这类重复工作，不能改变答案集合。

\textbf{优化条件。} 本题的关键观察是：先用 DP 或记忆化判断后缀是否可拆，再 DFS 枚举可行前缀并在末尾拼接句子。这句话必须能翻译成可维护状态；如果题目条件改变后观察不再成立，就不能继续套原来的写法。

\textbf{相邻状态。} 规律是先用记忆化记录某个起点之后能生成哪些句子，或先用 DP 剪掉不可拆后缀，再枚举路径。把这个特例继续拆开看：一层递归对应一个明确决策，路径保存已经做出的选择，选择列表保存仍可能扩展的分支。剪枝前必须能说明被剪掉的整棵子树没有合法答案，而不是只因为它看起来不优。

\textbf{状态复用。} 把完整搜索树加上合法性检查、剩余资源剪枝和同层去重；路径状态进入和退出要成对恢复。本题落点是：规律是先用记忆化记录某个起点之后能生成哪些句子，或先用 DP 剪掉不可拆后缀，再枚举路径。手算时只改被新输入影响的那一小部分，而不是回头重算完整候选。

\textbf{指针和字段。} 状态字段要能说清：路径、起点或使用标记、剩余目标、答案集合、剪枝条件；进入递归和退出递归必须成对恢复。手算时按这个协议跟踪：拿 catsanddog 手算，前缀 cat 和 cats 都能走，但某些后缀继续切会失败。若直接 DFS，每次都会重复探索同一个后缀。然后按协议复查：手算时给每层标出选择列表和路径。剪枝发生前要指出被剪分支为什么已经不可能得到合法答案。

\textbf{代码落点。} 剪枝条件写在扩展前；同层去重不要误写成全局去重；保存答案时复制当前路径。关键代码行必须能逐句翻译成“旧状态怎样变成新状态”，否则就只是模板。

\textbf{正确性。} 证明从搜索树覆盖开始：每个合法答案有且只有一条路径，剪枝只删掉不可能成功的分支。

\textbf{边界实验。} 先用无解后缀、重复单词、路径拼接成本、答案数量主导复杂度做反例检查。实验时覆盖无解后缀、重复单词、路径拼接成本、答案数量主导复杂度，统计搜索节点数量和剪枝次数。重复元素题要打印同一层跳过哪些值，确认没有误删合法路径。

\textbf{复杂度变化。} 暴力搜索树的规模由输出或选择空间决定；剪枝不改变最坏指数级本质，但能删除不可能成功或重复的分支，复杂度要说明输出规模。

\textbf{迁移追问。} 如果题目改成“只问能否拆分时用布尔 DP 即可，不要承担枚举所有句子的成本”，先判断原状态是否仍然覆盖新答案集合，再决定是微调边界、改 value 语义，还是换算法。

\noindent\textbf{LeetCode 188 买卖股票 IV。}

\textbf{题目说明。} LeetCode 188 买卖股票 IV 的题面入口要先还原成具体任务：最多 k 次交易让状态必须同时包含交易次数和是否持有股票。

\textbf{样例入口。} 拿 k=2、价格 [3, 2, 6] 手算。某天结束时只说最大利润不够，因为未来能否卖出取决于是否持股，也取决于已经完成几次交易。

\textbf{暴力基线。} 慢版本先写递归搜索树：节点表示处理位置、剩余资源和当前选择。只要不同路径反复到达同一个节点，就说明需要把这个节点命名成 DP 状态。放到本题就是：先按“最多 k 次交易让状态必须同时包含交易次数和是否持有股票”完整试慢版本；样例里已经能看到“规律是状态必须包含交易次数和持有状态；k 很大时才退化为无限次交易”，所以优化前必须先能说清慢版本枚举了哪些候选。

\textbf{暴力过程。} 拿 k=2、价格 [3, 2, 6] 手算。某天结束时只说最大利润不够，因为未来能否卖出取决于是否持股，也取决于已经完成几次交易。

\textbf{重复工作。} 题面模型：最多 k 次交易让状态必须同时包含交易次数和是否持有股票。暴力版的慢点要落到本题：慢版本会在“最多 k 次交易让状态必须同时包含交易次数和是否持有股票”的选择树里反复到达相同参数。样例规律是“规律是状态必须包含交易次数和持有状态；k 很大时才退化为无限次交易”，说明这个参数组合可以命名成状态，算过之后后面直接复用。后面的优化只消掉这类重复工作，不能改变答案集合。

\textbf{优化条件。} 本题的关键观察是：维护每个交易次数下的 buy 和 sell 状态；当 k 足够大时可退化为无限次交易。这句话必须能翻译成可维护状态；如果题目条件改变后观察不再成立，就不能继续套原来的写法。

\textbf{相邻状态。} 规律是状态必须包含交易次数和持有状态；k 很大时才退化为无限次交易。把这个特例继续拆开看：先问暴力搜索里哪些不同路径到达同一个后续问题，再给这个后续问题命名为状态。状态必须包含未来决策需要的全部信息，转移只从更小且已经算清的状态来。

\textbf{状态复用。} 把重复递归节点命名为状态；遍历顺序只是在保证依赖状态已经算完。本题落点是：规律是状态必须包含交易次数和持有状态；k 很大时才退化为无限次交易。手算时只改被新输入影响的那一小部分，而不是回头重算完整候选。

\textbf{指针和字段。} 状态字段要能说清：状态下标、状态值语义、初始化、转移来源、遍历顺序；空间压缩时还要指出读到的是旧值还是新值。手算时按这个协议跟踪：拿 k=2、价格 [3, 2, 6] 手算。某天结束时只说最大利润不够，因为未来能否卖出取决于是否持股，也取决于已经完成几次交易。然后按协议复查：手算时先写一个很小的 DP 表或记忆化搜索记录。每个格子旁边写它来自哪些更小状态。

\textbf{代码落点。} 先写未压缩版本校验转移；压缩前后用小表对拍；不可达哨兵参与加法前要检查溢出。关键代码行必须能逐句翻译成“旧状态怎样变成新状态”，否则就只是模板。

\textbf{正确性。} 证明从最后一步开始：任意最优解的最后一步都在转移枚举中，转移来源已经由更小状态正确计算。

\textbf{边界实验。} 先用k 为零、价格天数很少、k 大于等于 n/2、更新顺序造成同一天状态污染做反例检查。实验时覆盖k 为零、价格天数很少、k 大于等于 n/2、更新顺序造成同一天状态污染，先打印完整 DP 表，再比较空间压缩版本。若压缩版不同，优先检查遍历顺序和旧值保存。

\textbf{复杂度变化。} 暴力递归会重复计算相同子问题，常见指数级；DP 把每个状态算一次，时间是状态数乘单个转移成本，空间是状态表大小或压缩后的大小。

\textbf{迁移追问。} 如果题目改成“手续费、冷冻期和最多两次交易都是同一状态机框架的不同约束”，先判断原状态是否仍然覆盖新答案集合，再决定是微调边界、改 value 语义，还是换算法。

\noindent\textbf{LeetCode 212 单词搜索 II。}

\textbf{题目说明。} LeetCode 212 单词搜索 II 的题面入口要先还原成具体任务：多个单词在同一棋盘上搜索，公共前缀可以共享。

\textbf{样例入口。} 拿 board 上有 oath、pea、eat 手算，单词共享前缀时不应对每个词重复全图搜索；从 o 出发沿 Trie 前缀 o-a-t-h 找到 oath。

\textbf{暴力基线。} 慢版本就是完整搜索树：每层列出选择列表，路径到达终止条件时检查答案。它先保证覆盖所有可能，再把明显无效或重复的分支剪掉。放到本题就是：先按“多个单词在同一棋盘上搜索，公共前缀可以共享”完整试慢版本；样例里已经能看到“规律是 Trie 合并所有单词前缀，DFS 走棋盘同时走 Trie，命中单词后记录并可去重”，所以优化前必须先能说清慢版本枚举了哪些候选。

\textbf{暴力过程。} 拿 board 上有 oath、pea、eat 手算，单词共享前缀时不应对每个词重复全图搜索；从 o 出发沿 Trie 前缀 o-a-t-h 找到 oath。

\textbf{重复工作。} 题面模型：多个单词在同一棋盘上搜索，公共前缀可以共享。暴力版的慢点要落到本题：慢版本会围绕“多个单词在同一棋盘上搜索，公共前缀可以共享”展开完整搜索树。样例规律是“规律是 Trie 合并所有单词前缀，DFS 走棋盘同时走 Trie，命中单词后记录并可去重”，说明一层递归必须对应一个明确决策，剪枝只能删除整棵不可能成功或重复的子树。后面的优化只消掉这类重复工作，不能改变答案集合。

\textbf{优化条件。} 本题的关键观察是：先建 Trie，再从每个格子回溯；路径到达单词结束节点就收集答案。这句话必须能翻译成可维护状态；如果题目条件改变后观察不再成立，就不能继续套原来的写法。

\textbf{相邻状态。} 规律是 Trie 合并所有单词前缀，DFS 走棋盘同时走 Trie，命中单词后记录并可去重。把这个特例继续拆开看：一层递归对应一个明确决策，路径保存已经做出的选择，选择列表保存仍可能扩展的分支。剪枝前必须能说明被剪掉的整棵子树没有合法答案，而不是只因为它看起来不优。

\textbf{状态复用。} 把完整搜索树加上合法性检查、剩余资源剪枝和同层去重；路径状态进入和退出要成对恢复。本题落点是：规律是 Trie 合并所有单词前缀，DFS 走棋盘同时走 Trie，命中单词后记录并可去重。手算时只改被新输入影响的那一小部分，而不是回头重算完整候选。

\textbf{指针和字段。} 状态字段要能说清：路径、起点或使用标记、剩余目标、答案集合、剪枝条件；进入递归和退出递归必须成对恢复。手算时按这个协议跟踪：拿 board 上有 oath、pea、eat 手算，单词共享前缀时不应对每个词重复全图搜索；从 o 出发沿 Trie 前缀 o-a-t-h 找到 oath。然后按协议复查：手算时给每层标出选择列表和路径。剪枝发生前要指出被剪分支为什么已经不可能得到合法答案。

\textbf{代码落点。} 剪枝条件写在扩展前；同层去重不要误写成全局去重；保存答案时复制当前路径。关键代码行必须能逐句翻译成“旧状态怎样变成新状态”，否则就只是模板。

\textbf{正确性。} 证明从搜索树覆盖开始：每个合法答案有且只有一条路径，剪枝只删掉不可能成功的分支。

\textbf{边界实验。} 先用重复单词、同一格不能复用、前缀是完整单词、剪枝删除空子树做反例检查。实验时覆盖重复单词、同一格不能复用、前缀是完整单词、剪枝删除空子树，统计搜索节点数量和剪枝次数。重复元素题要打印同一层跳过哪些值，确认没有误删合法路径。

\textbf{复杂度变化。} 暴力搜索树的规模由输出或选择空间决定；剪枝不改变最坏指数级本质，但能删除不可能成功或重复的分支，复杂度要说明输出规模。

\textbf{迁移追问。} 如果题目改成“这是单词搜索从单目标扩展到多目标后的结构升级”，先判断原状态是否仍然覆盖新答案集合，再决定是微调边界、改 value 语义，还是换算法。

\noindent\textbf{LeetCode 224 基本计算器。}

\textbf{题目说明。} LeetCode 224 基本计算器 的题面入口要先还原成具体任务：括号会开启新的表达式上下文，括号结束时要回到上一层。

\textbf{样例入口。} 拿 1-(2+3) 手算，进入括号前当前结果是 1，符号是 -；括号内算出 5 后要乘以上层符号并合并成 -4。

\textbf{暴力基线。} 慢版本让每个位置向左或向右线性寻找边界，或者让每个结构反复等待匹配。它能算出答案，但很多尚未完成的候选会被未来同一个事件一起解决。放到本题就是：先按“括号会开启新的表达式上下文，括号结束时要回到上一层”完整试慢版本；样例里已经能看到“规律是栈保存进入括号前的结果和符号，一元负号和空格要按字符流处理”，所以优化前必须先能说清慢版本枚举了哪些候选。

\textbf{暴力过程。} 拿 1-(2+3) 手算，进入括号前当前结果是 1，符号是 -；括号内算出 5 后要乘以上层符号并合并成 -4。

\textbf{重复工作。} 题面模型：括号会开启新的表达式上下文，括号结束时要回到上一层。暴力版的慢点要落到本题：慢版本会围绕“括号会开启新的表达式上下文，括号结束时要回到上一层”反复寻找最近边界或最近未完成对象。样例规律是“规律是栈保存进入括号前的结果和符号，一元负号和空格要按字符流处理”，说明旧候选一旦遇到当前输入就能被批量结算；继续为每个位置单独向左或向右扫描就是浪费。后面的优化只消掉这类重复工作，不能改变答案集合。

\textbf{优化条件。} 本题的关键观察是：栈保存进入括号前的结果和符号，当前层算完后乘以上层符号并合并。这句话必须能翻译成可维护状态；如果题目条件改变后观察不再成立，就不能继续套原来的写法。

\textbf{相邻状态。} 规律是栈保存进入括号前的结果和符号，一元负号和空格要按字符流处理。把这个特例继续拆开看：栈里的对象都是答案尚未确定的候选；一旦当前元素触发弹栈，被弹出的候选必须已经同时拿到了左边界和右边界，未来元素不再有资格改变它的答案。

\textbf{状态复用。} 把反复寻找最近边界改成维护未完成候选；新元素只和栈顶比较，连续弹出一批已经被当前元素解决的对象。本题落点是：规律是栈保存进入括号前的结果和符号，一元负号和空格要按字符流处理。手算时只改被新输入影响的那一小部分，而不是回头重算完整候选。

\textbf{指针和字段。} 状态字段要能说清：栈内对象的业务含义、当前输入、弹出时结算的对象、左边界和右边界；栈中存下标还是值要明确。手算时按这个协议跟踪：拿 1-(2+3) 手算，进入括号前当前结果是 1，符号是 -；括号内算出 5 后要乘以上层符号并合并成 -4。然后按协议复查：手算时记录栈内元素的业务含义，不只记录数值。入栈表示答案未定，弹栈表示当前事件已经让它的答案定下来。

\textbf{代码落点。} 弹栈前确认栈非空，弹出后再读取新的左边界；相等元素和哨兵高度要有统一策略。关键代码行必须能逐句翻译成“旧状态怎样变成新状态”，否则就只是模板。

\textbf{正确性。} 证明从弹出时机开始：被弹出的对象在当前时刻答案已经确定，未来元素无法再改变它。

\textbf{边界实验。} 先用多位数、前导空格、一元负号、嵌套括号做反例检查。实验时准备多位数、前导空格、一元负号、嵌套括号，打印每次入栈、弹栈和结算对象。重点观察相等元素、空栈、哨兵和最后一次结算。

\textbf{复杂度变化。} 暴力为每个位置寻找边界常见为平方级；单调栈让每个元素最多入栈一次、出栈一次，因此主体扫描为线性时间。

\textbf{迁移追问。} 如果题目改成“基本计算器 II 没有括号但有乘除优先级，需要不同的栈语义”，先判断原状态是否仍然覆盖新答案集合，再决定是微调边界、改 value 语义，还是换算法。

\noindent\textbf{LeetCode 227 基本计算器 II。}

\textbf{题目说明。} LeetCode 227 基本计算器 II 的题面入口要先还原成具体任务：乘除优先级高于加减，可以在读到下一个运算符时结算上一项。

\textbf{样例入口。} 拿 3+2*2 手算，加号前的 3 可以先入栈，乘号遇到 2 时要立刻把栈顶 2 乘以下一个数 2，而不能等到最后按顺序相加。

\textbf{暴力基线。} 慢版本让每个位置向左或向右线性寻找边界，或者让每个结构反复等待匹配。它能算出答案，但很多尚未完成的候选会被未来同一个事件一起解决。放到本题就是：先按“乘除优先级高于加减，可以在读到下一个运算符时结算上一项”完整试慢版本；样例里已经能看到“规律是栈保存已确定符号的项，乘除立即修改栈顶，加减压入正负项”，所以优化前必须先能说清慢版本枚举了哪些候选。

\textbf{暴力过程。} 拿 3+2*2 手算，加号前的 3 可以先入栈，乘号遇到 2 时要立刻把栈顶 2 乘以下一个数 2，而不能等到最后按顺序相加。

\textbf{重复工作。} 题面模型：乘除优先级高于加减，可以在读到下一个运算符时结算上一项。暴力版的慢点要落到本题：慢版本会围绕“乘除优先级高于加减，可以在读到下一个运算符时结算上一项”反复寻找最近边界或最近未完成对象。样例规律是“规律是栈保存已确定符号的项，乘除立即修改栈顶，加减压入正负项”，说明旧候选一旦遇到当前输入就能被批量结算；继续为每个位置单独向左或向右扫描就是浪费。后面的优化只消掉这类重复工作，不能改变答案集合。

\textbf{优化条件。} 本题的关键观察是：栈保存已经确定符号的项，乘除直接修改栈顶，加减压入正负项。这句话必须能翻译成可维护状态；如果题目条件改变后观察不再成立，就不能继续套原来的写法。

\textbf{相邻状态。} 规律是栈保存已确定符号的项，乘除立即修改栈顶，加减压入正负项。把这个特例继续拆开看：栈里的对象都是答案尚未确定的候选；一旦当前元素触发弹栈，被弹出的候选必须已经同时拿到了左边界和右边界，未来元素不再有资格改变它的答案。

\textbf{状态复用。} 把反复寻找最近边界改成维护未完成候选；新元素只和栈顶比较，连续弹出一批已经被当前元素解决的对象。本题落点是：规律是栈保存已确定符号的项，乘除立即修改栈顶，加减压入正负项。手算时只改被新输入影响的那一小部分，而不是回头重算完整候选。

\textbf{指针和字段。} 状态字段要能说清：栈内对象的业务含义、当前输入、弹出时结算的对象、左边界和右边界；栈中存下标还是值要明确。手算时按这个协议跟踪：拿 3+2*2 手算，加号前的 3 可以先入栈，乘号遇到 2 时要立刻把栈顶 2 乘以下一个数 2，而不能等到最后按顺序相加。然后按协议复查：手算时记录栈内元素的业务含义，不只记录数值。入栈表示答案未定，弹栈表示当前事件已经让它的答案定下来。

\textbf{代码落点。} 弹栈前确认栈非空，弹出后再读取新的左边界；相等元素和哨兵高度要有统一策略。关键代码行必须能逐句翻译成“旧状态怎样变成新状态”，否则就只是模板。

\textbf{正确性。} 证明从弹出时机开始：被弹出的对象在当前时刻答案已经确定，未来元素无法再改变它。

\textbf{边界实验。} 先用多位数、空格、除法向零截断、最后一个数字需要结算做反例检查。实验时准备多位数、空格、除法向零截断、最后一个数字需要结算，打印每次入栈、弹栈和结算对象。重点观察相等元素、空栈、哨兵和最后一次结算。

\textbf{复杂度变化。} 暴力为每个位置寻找边界常见为平方级；单调栈让每个元素最多入栈一次、出栈一次，因此主体扫描为线性时间。

\textbf{迁移追问。} 如果题目改成“若加入括号，需要再引入上下文栈或递归下降解析”，先判断原状态是否仍然覆盖新答案集合，再决定是微调边界、改 value 语义，还是换算法。

\noindent\textbf{LeetCode 295 数据流的中位数。}

\textbf{题目说明。} LeetCode 295 数据流的中位数 的题面入口要先还原成具体任务：数据流需要动态维护较小一半和较大一半。

\textbf{样例入口。} 拿数据流 1、2、3 手算，插入 1 后中位数 1；插入 2 后两半是 [1] 和 [2]，中位数 1.5；插入 3 后右半多一个，要平衡成左半 [2, 1]、右半 [3]。

\textbf{暴力基线。} 暴力每插入一个数就重新排序所有已到达元素，然后按长度奇偶取中位数。

\textbf{暴力过程。} 数据流 1、2、3 中，插入 1 中位数是 1；插入 2 后中位数是 (1+2)/2；插入 3 后中位数是 2。

\textbf{重复工作。} 题面模型：数据流需要动态维护较小一半和较大一半。暴力版的慢点要落到本题：浪费在每次都维护完整排序。中位数只需要较小一半的最大值和较大一半的最小值。后面的优化只消掉这类重复工作，不能改变答案集合。

\textbf{优化条件。} 本题的关键观察是：大顶堆保存较小一半，小顶堆保存较大一半，大小差不超过一。这句话必须能翻译成可维护状态；如果题目条件改变后观察不再成立，就不能继续套原来的写法。

\textbf{相邻状态。} 规律是大顶堆保存较小一半，小顶堆保存较大一半，大小差最多 1。把这个特例继续拆开看：堆顶必须正好代表下一步最需要处理的候选；若堆里可能有过期对象，读取堆顶前要先清理。堆只解决动态极值，不自动证明被弹出或被丢弃的元素无用。

\textbf{状态复用。} 大顶堆 small 保存较小一半，小顶堆 large 保存较大一半。保持 small.size() 与 large.size() 差不超过 1，且 small.top()<=large.top()。手算时只改被新输入影响的那一小部分，而不是回头重算完整候选。

\textbf{指针和字段。} 状态字段要能说清：small.top 是左半最大值，large.top 是右半最小值；奇数时较大堆顶是中位数，偶数时两个堆顶平均。手算时按这个协议跟踪：插入 1 放 small；插入 2 后移动平衡成 small=[1]、large=[2]；插入 3 后 large 多，再移动 2 到 small，中位数 2。

\textbf{代码落点。} 关键是先按顺序不变量放入堆，再重平衡大小。求平均时用 double，并防止两个 int 相加溢出。关键代码行必须能逐句翻译成“旧状态怎样变成新状态”，否则就只是模板。

\textbf{正确性。} 两个堆按大小分割数据流，中位数只可能位于左半最大或右半最小；大小不变量保证取堆顶即可。

\textbf{边界实验。} 先用偶数个元素、负数、重复值、平均值溢出做反例检查。实验时覆盖偶数个元素、负数、重复值、平均值溢出，记录堆顶是否过期、比较器是否让正确候选在堆顶、K 或窗口大小为边界值时是否稳定。

\textbf{复杂度变化。} 每次重新排序 O(n log n)；双堆 addNum O(log n)，findMedian O(1)，空间 O(n)。

\textbf{迁移追问。} 如果题目改成“若还要删除任意元素，需要延迟删除哈希表”，先判断原状态是否仍然覆盖新答案集合，再决定是微调边界、改 value 语义，还是换算法。

\noindent\textbf{LeetCode 312 戳气球。}

\textbf{题目说明。} LeetCode 312 戳气球 的题面入口要先还原成具体任务：先戳哪个会让邻居持续变化，改成最后戳某个气球后区间边界才稳定。

\textbf{样例入口。} 拿 [3, 1, 5] 手算，若先戳 1，左右邻居会变；若改问某个区间里最后戳谁，最后一刻左右边界固定，收益可拆成左右两个子区间。

\textbf{暴力基线。} 慢版本先写递归搜索树：节点表示处理位置、剩余资源和当前选择。只要不同路径反复到达同一个节点，就说明需要把这个节点命名成 DP 状态。放到本题就是：先按“先戳哪个会让邻居持续变化，改成最后戳某个气球后区间边界才稳定”完整试慢版本；样例里已经能看到“规律是区间 DP 枚举最后戳的 mid，而不是第一个戳的气球”，所以优化前必须先能说清慢版本枚举了哪些候选。

\textbf{暴力过程。} 拿 [3, 1, 5] 手算，若先戳 1，左右邻居会变；若改问某个区间里最后戳谁，最后一刻左右边界固定，收益可拆成左右两个子区间。

\textbf{重复工作。} 题面模型：先戳哪个会让邻居持续变化，改成最后戳某个气球后区间边界才稳定。暴力版的慢点要落到本题：慢版本会在“先戳哪个会让邻居持续变化，改成最后戳某个气球后区间边界才稳定”的选择树里反复到达相同参数。样例规律是“规律是区间 DP 枚举最后戳的 mid，而不是第一个戳的气球”，说明这个参数组合可以命名成状态，算过之后后面直接复用。后面的优化只消掉这类重复工作，不能改变答案集合。

\textbf{优化条件。} 本题的关键观察是：区间 DP 定义开区间内全部戳完的最大收益，枚举最后一个被戳的气球作为切分点。这句话必须能翻译成可维护状态；如果题目条件改变后观察不再成立，就不能继续套原来的写法。

\textbf{相邻状态。} 规律是区间 DP 枚举最后戳的 mid，而不是第一个戳的气球。把这个特例继续拆开看：先问暴力搜索里哪些不同路径到达同一个后续问题，再给这个后续问题命名为状态。状态必须包含未来决策需要的全部信息，转移只从更小且已经算清的状态来。

\textbf{状态复用。} 把重复递归节点命名为状态；遍历顺序只是在保证依赖状态已经算完。本题落点是：规律是区间 DP 枚举最后戳的 mid，而不是第一个戳的气球。手算时只改被新输入影响的那一小部分，而不是回头重算完整候选。

\textbf{指针和字段。} 状态字段要能说清：状态下标、状态值语义、初始化、转移来源、遍历顺序；空间压缩时还要指出读到的是旧值还是新值。手算时按这个协议跟踪：拿 [3, 1, 5] 手算，若先戳 1，左右邻居会变；若改问某个区间里最后戳谁，最后一刻左右边界固定，收益可拆成左右两个子区间。然后按协议复查：手算时先写一个很小的 DP 表或记忆化搜索记录。每个格子旁边写它来自哪些更小状态。

\textbf{代码落点。} 先写未压缩版本校验转移；压缩前后用小表对拍；不可达哨兵参与加法前要检查溢出。关键代码行必须能逐句翻译成“旧状态怎样变成新状态”，否则就只是模板。

\textbf{正确性。} 证明从最后一步开始：任意最优解的最后一步都在转移枚举中，转移来源已经由更小状态正确计算。

\textbf{边界实验。} 先用补一的虚拟边界、区间长度遍历顺序、空区间收益为零、乘积可能较大做反例检查。实验时覆盖补一的虚拟边界、区间长度遍历顺序、空区间收益为零、乘积可能较大，先打印完整 DP 表，再比较空间压缩版本。若压缩版不同，优先检查遍历顺序和旧值保存。

\textbf{复杂度变化。} 暴力递归会重复计算相同子问题，常见指数级；DP 把每个状态算一次，时间是状态数乘单个转移成本，空间是状态表大小或压缩后的大小。

\textbf{迁移追问。} 如果题目改成“很多区间 DP 都要换成最后一步思考，避免中间操作改变边界”，先判断原状态是否仍然覆盖新答案集合，再决定是微调边界、改 value 语义，还是换算法。

\noindent\textbf{LeetCode 329 矩阵中的最长递增路径。}

\textbf{题目说明。} LeetCode 329 矩阵中的最长递增路径 的题面入口要先还原成具体任务：每个格子的答案是从它出发能走出的最长严格递增路径。

\textbf{样例入口。} 拿矩阵中 1->2->3 的小路径手算，从 1 出发会用到从 2 出发的答案；如果另一个格子也走到 2，后缀路径被重复计算。严格递增保证不会成环。

\textbf{暴力基线。} 慢版本先写递归搜索树：节点表示处理位置、剩余资源和当前选择。只要不同路径反复到达同一个节点，就说明需要把这个节点命名成 DP 状态。放到本题就是：先按“每个格子的答案是从它出发能走出的最长严格递增路径”完整试慢版本；样例里已经能看到“规律是每个格子记忆化保存从这里出发的最长路径”，所以优化前必须先能说清慢版本枚举了哪些候选。

\textbf{暴力过程。} 拿矩阵中 1->2->3 的小路径手算，从 1 出发会用到从 2 出发的答案；如果另一个格子也走到 2，后缀路径被重复计算。严格递增保证不会成环。

\textbf{重复工作。} 题面模型：每个格子的答案是从它出发能走出的最长严格递增路径。暴力版的慢点要落到本题：慢版本会在“每个格子的答案是从它出发能走出的最长严格递增路径”的选择树里反复到达相同参数。样例规律是“规律是每个格子记忆化保存从这里出发的最长路径”，说明这个参数组合可以命名成状态，算过之后后面直接复用。后面的优化只消掉这类重复工作，不能改变答案集合。

\textbf{优化条件。} 本题的关键观察是：把严格递增关系看成无环状态图，用记忆化 DFS 或按值排序的 DAG DP 复用后续格子结果。这句话必须能翻译成可维护状态；如果题目条件改变后观察不再成立，就不能继续套原来的写法。

\textbf{相邻状态。} 规律是每个格子记忆化保存从这里出发的最长路径。把这个特例继续拆开看：先问暴力搜索里哪些不同路径到达同一个后续问题，再给这个后续问题命名为状态。状态必须包含未来决策需要的全部信息，转移只从更小且已经算清的状态来。

\textbf{状态复用。} 把重复递归节点命名为状态；遍历顺序只是在保证依赖状态已经算完。本题落点是：规律是每个格子记忆化保存从这里出发的最长路径。手算时只改被新输入影响的那一小部分，而不是回头重算完整候选。

\textbf{指针和字段。} 状态字段要能说清：状态下标、状态值语义、初始化、转移来源、遍历顺序；空间压缩时还要指出读到的是旧值还是新值。手算时按这个协议跟踪：拿矩阵中 1->2->3 的小路径手算，从 1 出发会用到从 2 出发的答案；如果另一个格子也走到 2，后缀路径被重复计算。严格递增保证不会成环。然后按协议复查：手算时先写一个很小的 DP 表或记忆化搜索记录。每个格子旁边写它来自哪些更小状态。

\textbf{代码落点。} 先写未压缩版本校验转移；压缩前后用小表对拍；不可达哨兵参与加法前要检查溢出。关键代码行必须能逐句翻译成“旧状态怎样变成新状态”，否则就只是模板。

\textbf{正确性。} 证明从最后一步开始：任意最优解的最后一步都在转移枚举中，转移来源已经由更小状态正确计算。

\textbf{边界实验。} 先用平台相等不能走、四方向越界、递归深度、允许不下降时会产生环做反例检查。实验时覆盖平台相等不能走、四方向越界、递归深度、允许不下降时会产生环，先打印完整 DP 表，再比较空间压缩版本。若压缩版不同，优先检查遍历顺序和旧值保存。

\textbf{复杂度变化。} 暴力递归会重复计算相同子问题，常见指数级；DP 把每个状态算一次，时间是状态数乘单个转移成本，空间是状态表大小或压缩后的大小。

\textbf{迁移追问。} 如果题目改成“若改成求最短路径，严格递增的 DP 模型不再直接适用”，先判断原状态是否仍然覆盖新答案集合，再决定是微调边界、改 value 语义，还是换算法。

\noindent\textbf{LeetCode 394 字符串解码。}

\textbf{题目说明。} LeetCode 394 字符串解码 的题面入口要先还原成具体任务：嵌套结构需要保存上一层字符串和重复次数。

\textbf{样例入口。} 拿 3[a2[c]] 手算，遇到左括号前保存当前字符串和重复次数；遇到右括号时弹出上一层，把当前 cc 重复 2 次拼成 acc，再被外层重复 3 次。

\textbf{暴力基线。} 慢版本让每个位置向左或向右线性寻找边界，或者让每个结构反复等待匹配。它能算出答案，但很多尚未完成的候选会被未来同一个事件一起解决。放到本题就是：先按“嵌套结构需要保存上一层字符串和重复次数”完整试慢版本；样例里已经能看到“规律是栈保存进入括号前的上下文，右括号触发一次完整展开”，所以优化前必须先能说清慢版本枚举了哪些候选。

\textbf{暴力过程。} 拿 3[a2[c]] 手算，遇到左括号前保存当前字符串和重复次数；遇到右括号时弹出上一层，把当前 cc 重复 2 次拼成 acc，再被外层重复 3 次。

\textbf{重复工作。} 题面模型：嵌套结构需要保存上一层字符串和重复次数。暴力版的慢点要落到本题：慢版本会围绕“嵌套结构需要保存上一层字符串和重复次数”反复寻找最近边界或最近未完成对象。样例规律是“规律是栈保存进入括号前的上下文，右括号触发一次完整展开”，说明旧候选一旦遇到当前输入就能被批量结算；继续为每个位置单独向左或向右扫描就是浪费。后面的优化只消掉这类重复工作，不能改变答案集合。

\textbf{优化条件。} 本题的关键观察是：遇到左括号入栈当前状态，右括号弹出并拼接展开。这句话必须能翻译成可维护状态；如果题目条件改变后观察不再成立，就不能继续套原来的写法。

\textbf{相邻状态。} 规律是栈保存进入括号前的上下文，右括号触发一次完整展开。把这个特例继续拆开看：栈里的对象都是答案尚未确定的候选；一旦当前元素触发弹栈，被弹出的候选必须已经同时拿到了左边界和右边界，未来元素不再有资格改变它的答案。

\textbf{状态复用。} 把反复寻找最近边界改成维护未完成候选；新元素只和栈顶比较，连续弹出一批已经被当前元素解决的对象。本题落点是：规律是栈保存进入括号前的上下文，右括号触发一次完整展开。手算时只改被新输入影响的那一小部分，而不是回头重算完整候选。

\textbf{指针和字段。} 状态字段要能说清：栈内对象的业务含义、当前输入、弹出时结算的对象、左边界和右边界；栈中存下标还是值要明确。手算时按这个协议跟踪：拿 3[a2[c]] 手算，遇到左括号前保存当前字符串和重复次数；遇到右括号时弹出上一层，把当前 cc 重复 2 次拼成 acc，再被外层重复 3 次。然后按协议复查：手算时记录栈内元素的业务含义，不只记录数值。入栈表示答案未定，弹栈表示当前事件已经让它的答案定下来。

\textbf{代码落点。} 弹栈前确认栈非空，弹出后再读取新的左边界；相等元素和哨兵高度要有统一策略。关键代码行必须能逐句翻译成“旧状态怎样变成新状态”，否则就只是模板。

\textbf{正确性。} 证明从弹出时机开始：被弹出的对象在当前时刻答案已经确定，未来元素无法再改变它。

\textbf{边界实验。} 先用多位数字、嵌套、空片段、数字后必须跟括号做反例检查。实验时准备多位数字、嵌套、空片段、数字后必须跟括号，打印每次入栈、弹栈和结算对象。重点观察相等元素、空栈、哨兵和最后一次结算。

\textbf{复杂度变化。} 暴力为每个位置寻找边界常见为平方级；单调栈让每个元素最多入栈一次、出栈一次，因此主体扫描为线性时间。

\textbf{迁移追问。} 如果题目改成“递归下降解析更通用，但迭代栈更符合工程栈控制”，先判断原状态是否仍然覆盖新答案集合，再决定是微调边界、改 value 语义，还是换算法。

\noindent\textbf{LeetCode 460 LFU 缓存。}

\textbf{题目说明。} LeetCode 460 LFU 缓存 的题面入口要先还原成具体任务：缓存淘汰同时按访问频次和同频最近性排序，需要维护两个维度的索引。

\textbf{样例入口。} 拿容量 2，put(1)、put(2)、get(1)、put(3) 手算。此时 1 的频次高于 2，应淘汰 2；若两个 key 频次相同，还要淘汰同频里最旧的。

\textbf{暴力基线。} 慢版本可以先把节点顺序抄到数组里，按数组推导目标顺序。回到链表后，难点不再是值的顺序，而是节点身份和 next 指针不能丢。放到本题就是：先按“缓存淘汰同时按访问频次和同频最近性排序，需要维护两个维度的索引”完整试慢版本；样例里已经能看到“规律是 key 哈希定位节点，频次哈希定位同频链表，minFreq 指向当前最低频次”，所以优化前必须先能说清慢版本枚举了哪些候选。

\textbf{暴力过程。} 拿容量 2，put(1)、put(2)、get(1)、put(3) 手算。此时 1 的频次高于 2，应淘汰 2；若两个 key 频次相同，还要淘汰同频里最旧的。

\textbf{重复工作。} 题面模型：缓存淘汰同时按访问频次和同频最近性排序，需要维护两个维度的索引。暴力版的慢点要落到本题：慢版本会围绕“缓存淘汰同时按访问频次和同频最近性排序，需要维护两个维度的索引”借助数组或多次扫描确认目标顺序。样例规律是“规律是 key 哈希定位节点，频次哈希定位同频链表，minFreq 指向当前最低频次”，说明优化点不在节点值，而在少量指针角色能否保持已处理链、未处理链和接回位置都不丢。后面的优化只消掉这类重复工作，不能改变答案集合。

\textbf{优化条件。} 本题的关键观察是：哈希表按 key 定位节点，频次到双向链表保存同频节点顺序，并维护当前最小频次。这句话必须能翻译成可维护状态；如果题目条件改变后观察不再成立，就不能继续套原来的写法。

\textbf{相邻状态。} 规律是 key 哈希定位节点，频次哈希定位同频链表，minFreq 指向当前最低频次。把这个特例继续拆开看：每个指针都要有角色，上一段尾、当前段头、当前节点和后继入口不能混。只要一次改 next 之前没有保存后继，就可能丢掉整段未处理链。

\textbf{状态复用。} 把数组辅助结构换成几个稳定指针角色；重连顺序必须保证已处理链合法、未处理链仍可达。本题落点是：规律是 key 哈希定位节点，频次哈希定位同频链表，minFreq 指向当前最低频次。手算时只改被新输入影响的那一小部分，而不是回头重算完整候选。

\textbf{指针和字段。} 状态字段要能说清：虚拟头、上一段尾、当前段头、当前节点、后继节点；任何改 next 的操作之前都要保存后续入口。手算时按这个协议跟踪：拿容量 2，put(1)、put(2)、get(1)、put(3) 手算。此时 1 的频次高于 2，应淘汰 2；若两个 key 频次相同，还要淘汰同频里最旧的。然后按协议复查：手算时画出 prev、cur、next 和下一段头节点。每次改 next 前先保存后继，这是链表推导的安全线。

\textbf{代码落点。} 所有重连步骤按保存后继、断开或反转、接回前缀、接回后缀的顺序写；最后检查是否形成环或丢节点。关键代码行必须能逐句翻译成“旧状态怎样变成新状态”，否则就只是模板。

\textbf{正确性。} 证明从链结构开始：已处理链连通、未处理链可达、二者连接点明确，节点没有丢失也没有成环。

\textbf{边界实验。} 先用容量为零、更新已有 key、频次链表变空时最小频次更新、同频淘汰最旧节点做反例检查。实验时覆盖容量为零、更新已有 key、频次链表变空时最小频次更新、同频淘汰最旧节点，每步画出前驱、当前、后继和下一段头。链表题不要只看输出值，还要检查节点是否丢失或成环。

\textbf{复杂度变化。} 暴力可能多次扫描链表或拷贝到数组；优化后每个节点通常只被访问、重连常数次，目标是线性时间和常数额外空间。

\textbf{迁移追问。} 如果题目改成“LRU 只有时间顺序一个维度，LFU 多了频次维度，结构维护明显更复杂”，先判断原状态是否仍然覆盖新答案集合，再决定是微调边界、改 value 语义，还是换算法。

\noindent\textbf{LeetCode 560 和为 K 的子数组。}

\textbf{题目说明。} LeetCode 560 和为 K 的子数组给定整数数组 nums 和整数 k，统计和等于 k 的连续子数组个数；数组中可以有负数。

\textbf{样例入口。} 样例 nums=[1, 1, 1]，k=2，输出 2；两个合法连续子数组分别是前两个 1 和后两个 1。

\textbf{暴力基线。} 暴力固定 left，向右累计 sum；每当 sum==k 就计数。

\textbf{暴力过程。} 以 nums=[1, 1, 1]、k=2 为例，从 left=0 可找到 [1, 1]，从 left=1 又找到后两个 1。

\textbf{重复工作。} 题面模型：当前前缀和减去目标值，就是需要匹配的历史前缀。暴力版的慢点要落到本题：浪费在每个右端都回头枚举所有 left。区间和等于 k 可以改写成 prefix[right]-prefix[left]=k。后面的优化只消掉这类重复工作，不能改变答案集合。

\textbf{优化条件。} 本题的关键观察是：哈希表保存前缀和出现次数，先查询再更新。这句话必须能翻译成可维护状态；如果题目条件改变后观察不再成立，就不能继续套原来的写法。

\textbf{相邻状态。} 规律是哈希表保存前缀和频次，且初始要放 0:1。把这个特例继续拆开看：每个合法答案都要还原成当前状态和某个历史状态的配对；哈希表的 key 负责定位历史状态，value 负责表达次数、最早位置或存在性。先查还是先写，必须由是否允许当前前缀和自己配对来决定。

\textbf{状态复用。} 扫描前缀和 prefix。哈希表 count[old\_prefix] 保存历史前缀出现次数；当前位置需要查 prefix-k 出现过几次。手算时只改被新输入影响的那一小部分，而不是回头重算完整候选。

\textbf{指针和字段。} 状态字段要能说清：prefix 是当前前缀和，need=prefix-k 是所需历史前缀，count 保存历史前缀频次，answer 累加匹配次数。手算时按这个协议跟踪：[1, 1, 1] 中扫描到第二个元素时 prefix=2，需要历史 0，count[0]=1，所以找到第一个长度 2 子数组。

\textbf{代码落点。} 关键是先 answer += count[prefix-k]，再 count[prefix]++。初始化 count[0]=1，表示从下标 0 开始的区间。关键代码行必须能逐句翻译成“旧状态怎样变成新状态”，否则就只是模板。

\textbf{正确性。} 任意和为 k 的子数组都对应一对前缀差 k；扫描到右端时，所有合法左端前缀都已在表里。

\textbf{边界实验。} 先用负数、目标为零、从下标零开始的区间、重复前缀做反例检查。实验时把负数、目标为零、从下标零开始的区间、重复前缀加入对拍集合，用平方暴力和优化版比较。失败时先看初始历史状态、查询更新顺序和哈希表 value 类型。

\textbf{复杂度变化。} 暴力 O(\ensuremath{n^{2}})，前缀哈希 O(n) 均摊时间、O(n) 空间。

\textbf{迁移追问。} 如果题目改成“若要求最长区间，哈希表应保存最早位置而不是次数”，先判断原状态是否仍然覆盖新答案集合，再决定是微调边界、改 value 语义，还是换算法。

\noindent\textbf{LeetCode 632 最小区间覆盖 K 个列表。}

\textbf{题目说明。} LeetCode 632 最小区间覆盖 K 个列表 的题面入口要先还原成具体任务：每个列表至少取一个元素时，当前区间由最小当前值和最大当前值决定。

\textbf{样例入口。} 拿三组 [4, 10]、[0, 9]、[5, 18] 手算，当前头是 4、0、5，区间 [0, 5] 覆盖三组；弹出最小头 0 后推进该组。

\textbf{暴力基线。} 慢版本每轮扫描全部候选来找最大或最小，再更新候选集合。它的答案选择过程清楚，但动态集合每变化一次就重新找极值，代价太高。放到本题就是：先按“每个列表至少取一个元素时，当前区间由最小当前值和最大当前值决定”完整试慢版本；样例里已经能看到“规律是堆保存每组当前元素，同时维护当前最大值，最小头决定下一步推进哪一组”，所以优化前必须先能说清慢版本枚举了哪些候选。

\textbf{暴力过程。} 拿三组 [4, 10]、[0, 9]、[5, 18] 手算，当前头是 4、0、5，区间 [0, 5] 覆盖三组；弹出最小头 0 后推进该组。

\textbf{重复工作。} 题面模型：每个列表至少取一个元素时，当前区间由最小当前值和最大当前值决定。暴力版的慢点要落到本题：慢版本会围绕“每个列表至少取一个元素时，当前区间由最小当前值和最大当前值决定”反复扫描动态候选集找极值。样例规律是“规律是堆保存每组当前元素，同时维护当前最大值，最小头决定下一步推进哪一组”，说明每轮真正需要的是堆顶边界，而不是把候选全集重新排序。后面的优化只消掉这类重复工作，不能改变答案集合。

\textbf{优化条件。} 本题的关键观察是：小顶堆保存各列表当前元素，同时维护当前最大值；弹出最小所在列表并推进。这句话必须能翻译成可维护状态；如果题目条件改变后观察不再成立，就不能继续套原来的写法。

\textbf{相邻状态。} 规律是堆保存每组当前元素，同时维护当前最大值，最小头决定下一步推进哪一组。把这个特例继续拆开看：堆顶必须正好代表下一步最需要处理的候选；若堆里可能有过期对象，读取堆顶前要先清理。堆只解决动态极值，不自动证明被弹出或被丢弃的元素无用。

\textbf{状态复用。} 把每轮全量扫描改成堆顶查询；插入、弹出和过期清理共同维护动态候选集。本题落点是：规律是堆保存每组当前元素，同时维护当前最大值，最小头决定下一步推进哪一组。手算时只改被新输入影响的那一小部分，而不是回头重算完整候选。

\textbf{指针和字段。} 状态字段要能说清：堆元素字段、堆顶比较规则、候选是否过期、答案读取时机；如果需要懒删除，还要保存删除计数。手算时按这个协议跟踪：拿三组 [4, 10]、[0, 9]、[5, 18] 手算，当前头是 4、0、5，区间 [0, 5] 覆盖三组；弹出最小头 0 后推进该组。然后按协议复查：手算时记录堆顶、候选集合和是否有过期元素。每次使用堆顶前都要确认它仍然属于当前问题。

\textbf{代码落点。} 比较器方向先用小样例验证；每次使用堆顶前清理过期项；堆中保存下标时要能回查原数据。关键代码行必须能逐句翻译成“旧状态怎样变成新状态”，否则就只是模板。

\textbf{正确性。} 证明从候选覆盖开始：所有可能成为下一步答案的对象都在堆中，堆顶过期时不会被使用。

\textbf{边界实验。} 先用某个列表为空、重复值、区间长度相同的 tie-breaker、列表耗尽做反例检查。实验时覆盖某个列表为空、重复值、区间长度相同的 tie-breaker、列表耗尽，记录堆顶是否过期、比较器是否让正确候选在堆顶、K 或窗口大小为边界值时是否稳定。

\textbf{复杂度变化。} 暴力每轮扫描或排序动态候选，会把线性扫描或排序反复发生；堆把取极值、插入和删除边界控制在对数级。

\textbf{迁移追问。} 如果题目改成“这是多路归并和覆盖窗口的结合”，先判断原状态是否仍然覆盖新答案集合，再决定是微调边界、改 value 语义，还是换算法。

\noindent\textbf{LeetCode 678 有效的括号字符串。}

\textbf{题目说明。} LeetCode 678 有效的括号字符串 的题面入口要先还原成具体任务：星号可以作为左括号、右括号或空，使未匹配左括号数量变成一个可能区间。

\textbf{样例入口。} 拿 (*) 手算，星号可以当空或左括号；拿 (*))，星号必须当左括号才能匹配后面的右括号。与其枚举星号三种选择，不如维护可能未匹配左括号数的区间 [low, high]。

\textbf{暴力基线。} 慢版本枚举选择顺序或用 DP 保留多个可能方案，先确认最优值存在。随后才检查局部选择是否能通过交换或领先性固定下来。放到本题就是：先按“星号可以作为左括号、右括号或空，使未匹配左括号数量变成一个可能区间”完整试慢版本；样例里已经能看到“规律是左括号让区间加一，右括号让区间减一，星号让下界减一上界加一，且下界不能低于零”，所以优化前必须先能说清慢版本枚举了哪些候选。

\textbf{暴力过程。} 拿 (*) 手算，星号可以当空或左括号；拿 (*))，星号必须当左括号才能匹配后面的右括号。与其枚举星号三种选择，不如维护可能未匹配左括号数的区间 [low, high]。

\textbf{重复工作。} 题面模型：星号可以作为左括号、右括号或空，使未匹配左括号数量变成一个可能区间。暴力版的慢点要落到本题：慢版本会围绕“星号可以作为左括号、右括号或空，使未匹配左括号数量变成一个可能区间”枚举选择顺序。样例规律是“规律是左括号让区间加一，右括号让区间减一，星号让下界减一上界加一，且下界不能低于零”，说明存在一个局部选择或边界状态可以安全保留；不能证明这个选择不变差时，就不能把它叫贪心。后面的优化只消掉这类重复工作，不能改变答案集合。

\textbf{优化条件。} 本题的关键观察是：贪心维护可能未匹配左括号数的下界和上界，遇到星号扩大区间并把下界截到零。这句话必须能翻译成可维护状态；如果题目条件改变后观察不再成立，就不能继续套原来的写法。

\textbf{相邻状态。} 规律是左括号让区间加一，右括号让区间减一，星号让下界减一上界加一，且下界不能低于零。把这个特例继续拆开看：先构造一个非贪心选择，再比较两者留下的资源、右边界或剩余时间。只有能把非贪心第一步交换成贪心第一步且不变差，局部选择才是规律。

\textbf{状态复用。} 把指数级选择压成一个可证明安全的局部选择；状态通常是当前最远边界、最早结束时间、最小代价或剩余资源。本题落点是：规律是左括号让区间加一，右括号让区间减一，星号让下界减一上界加一，且下界不能低于零。手算时只改被新输入影响的那一小部分，而不是回头重算完整候选。

\textbf{指针和字段。} 状态字段要能说清：处理顺序、当前已选方案、剩余资源或最远边界、答案计数；若使用排序，要说明排序键；每次局部选择后要能说明当前状态不落后。手算时按这个协议跟踪：拿 (*) 手算，星号可以当空或左括号；拿 (*))，星号必须当左括号才能匹配后面的右括号。与其枚举星号三种选择，不如维护可能未匹配左括号数的区间 [low, high]。然后按协议复查：手算时同时写一个非贪心选择，与贪心选择比较剩余资源。若无法说明贪心后剩余资源不差，就不能直接下结论。

\textbf{代码落点。} 把局部规则写成业务含义；若需要排序，就处理相同关键字；循环中只维护证明需要的状态，不把局部选择和答案更新混在一起。关键代码行必须能逐句翻译成“旧状态怎样变成新状态”，否则就只是模板。

\textbf{正确性。} 证明从交换或领先性开始：任意最优解都能替换成当前贪心选择而不变差，或贪心状态始终不落后。

\textbf{边界实验。} 先用上界变负、最终下界不为零、连续星号、前缀右括号过多做反例检查。实验时除了上界变负、最终下界不为零、连续星号、前缀右括号过多，还要故意替换成一个看似合理但错误的局部规则，构造反例。能解释错误贪心为何失败，才算理解正确贪心。

\textbf{复杂度变化。} 暴力枚举选择顺序可能指数级；贪心通常先排序，再线性扫描或用堆维护少量候选。

\textbf{迁移追问。} 如果题目改成“也可用双栈保存左括号和星号位置，但区间贪心更能说明状态压缩”，先判断原状态是否仍然覆盖新答案集合，再决定是微调边界、改 value 语义，还是换算法。

\noindent\textbf{LeetCode 721 账户合并。}

\textbf{题目说明。} LeetCode 721 账户合并 的题面入口要先还原成具体任务：共享邮箱说明账户属于同一人，邮箱之间形成连通分量。

\textbf{样例入口。} 拿 John 的账号 [a, b] 和另一个 [b, c] 手算，邮箱 b 重叠，所以 a、b、c 属于同一人；名字只是输出标签，合并依据是邮箱。

\textbf{暴力基线。} 慢版本每次合并后用 DFS 或 BFS 重新判断连通性。它正确但重复遍历已有连接；当关系只会合并不会拆开时，可以压成集合代表。放到本题就是：先按“共享邮箱说明账户属于同一人，邮箱之间形成连通分量”完整试慢版本；样例里已经能看到“规律是邮箱映射到编号后并查集合并，最后按代表收集并排序邮箱”，所以优化前必须先能说清慢版本枚举了哪些候选。

\textbf{暴力过程。} 拿 John 的账号 [a, b] 和另一个 [b, c] 手算，邮箱 b 重叠，所以 a、b、c 属于同一人；名字只是输出标签，合并依据是邮箱。

\textbf{重复工作。} 题面模型：共享邮箱说明账户属于同一人，邮箱之间形成连通分量。暴力版的慢点要落到本题：慢版本会围绕“共享邮箱说明账户属于同一人，邮箱之间形成连通分量”反复搜索连通性。样例规律是“规律是邮箱映射到编号后并查集合并，最后按代表收集并排序邮箱”，说明关系只需要被压成集合代表；每次 union 失败或成功都要能翻译回题意。后面的优化只消掉这类重复工作，不能改变答案集合。

\textbf{优化条件。} 本题的关键观察是：给邮箱编号并查集合并，同根邮箱收集后排序输出。这句话必须能翻译成可维护状态；如果题目条件改变后观察不再成立，就不能继续套原来的写法。

\textbf{相邻状态。} 规律是邮箱映射到编号后并查集合并，最后按代表收集并排序邮箱。把这个特例继续拆开看：union 只是在维护等价类，不是在保存具体路径。两个节点代表元相同只能说明连通或关系一致；如果题目还要距离、比例或奇偶性，必须把附加信息挂到根关系上。

\textbf{状态复用。} 把连通分量压成代表元；查询只比较代表，合并只发生在两个根之间。本题落点是：规律是邮箱映射到编号后并查集合并，最后按代表收集并排序邮箱。手算时只改被新输入影响的那一小部分，而不是回头重算完整候选。

\textbf{指针和字段。} 状态字段要能说清：parent、size 或 rank、find 返回的代表元、合并时的附加信息；带权或奇偶并查集还要维护到根的关系。手算时按这个协议跟踪：拿 John 的账号 [a, b] 和另一个 [b, c] 手算，邮箱 b 重叠，所以 a、b、c 属于同一人；名字只是输出标签，合并依据是邮箱。然后按协议复查：手算时写 parent 表和集合代表。每次 union 只改变根之间的关系，find 后代表相同才说明连通。

\textbf{代码落点。} find 路径压缩不能破坏附加权值；union 只能连接两个根；合并失败和重复边要保留题意解释。关键代码行必须能逐句翻译成“旧状态怎样变成新状态”，否则就只是模板。

\textbf{正确性。} 证明从等价关系开始：合并维护连通性的传递性，find 返回的代表元就是当前集合身份。

\textbf{边界实验。} 先用同名不同人、一个账户多个邮箱、重复邮箱、输出排序做反例检查。实验时覆盖同名不同人、一个账户多个邮箱、重复邮箱、输出排序，打印每个元素的代表元和集合大小。重复合并、已经连通的边和字符串编号最容易暴露实现问题。

\textbf{复杂度变化。} 暴力每次查询都搜索连通性代价高；并查集把合并和查询摊还到近似常数，整体接近线性。

\textbf{迁移追问。} 如果题目改成“姓名不能作为合并依据，只能作为输出字段”，先判断原状态是否仍然覆盖新答案集合，再决定是微调边界、改 value 语义，还是换算法。

\noindent\textbf{LeetCode 815 公交路线。}

\textbf{题目说明。} LeetCode 815 公交路线 的题面入口要先还原成具体任务：公交路线和站点构成二分关系，换乘次数比站点步数更重要。

\textbf{样例入口。} 拿起点站可坐路线 A，A 和 B 在某站相交，B 到终点手算。按站点一步步 BFS 会反复展开同一条路线的所有站；按路线 BFS，一次乘车层数就是换乘次数。

\textbf{暴力基线。} 慢版本先按题意画出完整状态图，用普通搜索跑通可达性。这个过程用于确认公交路线和站点构成二分关系，换乘次数比站点步数更重要的节点和边，之后再讨论访问标记、层次或代价顺序。放到本题就是：先按“公交路线和站点构成二分关系，换乘次数比站点步数更重要”完整试慢版本；样例里已经能看到“规律是访问过的路线不再展开，站点到路线表只用于找到下一批路线”，所以优化前必须先能说清慢版本枚举了哪些候选。

\textbf{暴力过程。} 拿起点站可坐路线 A，A 和 B 在某站相交，B 到终点手算。按站点一步步 BFS 会反复展开同一条路线的所有站；按路线 BFS，一次乘车层数就是换乘次数。

\textbf{重复工作。} 题面模型：公交路线和站点构成二分关系，换乘次数比站点步数更重要。暴力版的慢点要落到本题：慢版本会围绕“公交路线和站点构成二分关系，换乘次数比站点步数更重要”反复展开同一批状态。样例规律是“规律是访问过的路线不再展开，站点到路线表只用于找到下一批路线”，说明必须把节点、边、访问标记和资源维度一次建清；同一状态不应重复入队，不同资源状态也不能误合并。后面的优化只消掉这类重复工作，不能改变答案集合。

\textbf{优化条件。} 本题的关键观察是：用站点到路线列表的哈希表找可乘路线，以路线或乘车层数做 BFS，访问过的路线不再重复展开。这句话必须能翻译成可维护状态；如果题目条件改变后观察不再成立，就不能继续套原来的写法。

\textbf{相邻状态。} 规律是访问过的路线不再展开，站点到路线表只用于找到下一批路线。把这个特例继续拆开看：状态节点不一定等于原始位置，还可能包含钥匙、剩余资源、时间或访问集合。visited 只能合并未来能力完全相同的状态，否则会把合法路径剪掉。

\textbf{状态复用。} 把重复搜索压成访问标记、距离数组或拓扑状态；邻居生成也要从全量扫描变成结构化枚举。本题落点是：规律是访问过的路线不再展开，站点到路线表只用于找到下一批路线。手算时只改被新输入影响的那一小部分，而不是回头重算完整候选。

\textbf{指针和字段。} 状态字段要能说清：状态编号、邻接生成方式、访问标记、距离或层数、队列内容；若状态含资源，visited 不能只按位置记录。手算时按这个协议跟踪：拿起点站可坐路线 A，A 和 B 在某站相交，B 到终点手算。按站点一步步 BFS 会反复展开同一条路线的所有站；按路线 BFS，一次乘车层数就是换乘次数。然后按协议复查：手算时画队列或栈，状态必须包含题目需要的所有资源。入队时标记还是出队时标记，要和去重语义一致。

\textbf{代码落点。} 入队时标记还是出队时标记必须一致；方向数组集中定义；多源 BFS 要一次性把所有源点放入初始队列。关键代码行必须能逐句翻译成“旧状态怎样变成新状态”，否则就只是模板。

\textbf{正确性。} 证明从可达性开始：搜索覆盖所有合法边能到达的状态，访问标记只删除重复状态而不删除不同未来能力。

\textbf{边界实验。} 先用起点等于终点、某站经过路线很多、路线重复访问、目标站不可达做反例检查。实验时覆盖起点等于终点、某站经过路线很多、路线重复访问、目标站不可达，并打印入队时的完整状态。若同一位置但资源不同的状态被合并，很多图题会出现隐蔽错误。

\textbf{复杂度变化。} 暴力重复从多个起点搜索会反复遍历图；正确建模和访问标记后，BFS 或 DFS 通常按节点数加边数计算，状态图题则按完整状态数计算。

\textbf{迁移追问。} 如果题目改成“若把每个站点之间完全建边会爆炸，应利用路线这一层压缩邻居生成”，先判断原状态是否仍然覆盖新答案集合，再决定是微调边界、改 value 语义，还是换算法。

\noindent\textbf{LeetCode 862 和至少为 K 的最短子数组。}

\textbf{题目说明。} LeetCode 862 和至少为 K 的最短子数组 的题面入口要先还原成具体任务：数组允许负数，普通滑动窗口的和不再随左端移动单调变化。

\textbf{样例入口。} 拿 [2, -1, 2]、K=3 手算，普通正数窗口会被 -1 破坏单调性；前缀和为 0, 2, 1, 3，只有 3-0 达到 3。

\textbf{暴力基线。} 暴力枚举所有 left/right，用前缀和 O(1) 求区间和，检查是否至少为 K 并更新最短长度。

\textbf{暴力过程。} 以 [2, -1, 2]、K=3 为例，前缀和是 [0, 2, 1, 3]。普通正数窗口会被 -1 破坏：右端增加后和可能变小。

\textbf{重复工作。} 题面模型：数组允许负数，普通滑动窗口的和不再随左端移动单调变化。暴力版的慢点要落到本题：浪费在每个右端都回头扫描所有历史左端。需要的是某些前缀和尽量小且下标尽量靠后的候选左端。后面的优化只消掉这类重复工作，不能改变答案集合。

\textbf{优化条件。} 本题的关键观察是：在前缀和数组上用单调队列维护可能成为最优左端的下标。这句话必须能翻译成可维护状态；如果题目条件改变后观察不再成立，就不能继续套原来的写法。

\textbf{相邻状态。} 规律是在前缀和上用单调队列维护可能的最优左端，队首能满足时结算长度。把这个特例继续拆开看：右端进入只带来一个新元素，左端离开只删掉一个旧元素；只有当旧左端被移走后不可能再产生更优答案时，窗口才可以单向推进。若这个单调淘汰条件不存在，就要换成前缀和、队列或其他状态。

\textbf{状态复用。} 构造前缀和 prefix。用单调递增队列保存可能作为左端的前缀下标；当前前缀减队首前缀 >=K 时更新答案并弹出队首；队尾前缀不小于当前前缀时被当前支配，弹出。手算时只改被新输入影响的那一小部分，而不是回头重算完整候选。

\textbf{指针和字段。} 状态字段要能说清：deque 中保存前缀下标，且对应 prefix 值递增；i 是当前右端后一位，i-deque.front 是候选长度。手算时按这个协议跟踪：prefix 到 3 时，3-0>=3，长度 3；若出现更小的新前缀，它会支配队尾较大前缀，因为未来用它能得到更大区间和且更短或不更差。

\textbf{代码落点。} 关键有两个 while：先从队首结算满足 K 的最短长度，再从队尾删除 prefix 值不小于当前值的被支配候选。关键代码行必须能逐句翻译成“旧状态怎样变成新状态”，否则就只是模板。

\textbf{正确性。} 队首是当前最早且前缀较小的候选；一旦满足 K，继续保留它只会让未来长度更长，所以可以弹出。队尾被更小且更新的前缀支配，永远不会更优。

\textbf{边界实验。} 先用负数、答案不存在、K 为负或零、队尾支配关系做反例检查。实验时重点跑负数、答案不存在、K 为负或零、队尾支配关系，并打印左右端、窗口计数和答案。若左端发生回退，或者窗口失去合法性后仍更新答案，就能很快定位错误。

\textbf{复杂度变化。} 暴力 O(\ensuremath{n^{2}})，单调队列中每个前缀入队出队各一次，时间 O(n)，空间 O(n)。

\textbf{迁移追问。} 如果题目改成“这是从窗口题升级到前缀和加单调队列的代表”，先判断原状态是否仍然覆盖新答案集合，再决定是微调边界、改 value 语义，还是换算法。

\noindent\textbf{LeetCode 895 最大频率栈。}

\textbf{题目说明。} LeetCode 895 最大频率栈 的题面入口要先还原成具体任务：弹出时既要频率最高，又要在同频中最近入栈。

\textbf{样例入口。} 拿 push 5, 7, 5, 7, 4, 5 后 pop 手算，频率最高是 5，先弹 5；下一次 5 和 7 同频，弹最近达到该频率的 7。

\textbf{暴力基线。} 慢版本让每个位置向左或向右线性寻找边界，或者让每个结构反复等待匹配。它能算出答案，但很多尚未完成的候选会被未来同一个事件一起解决。放到本题就是：先按“弹出时既要频率最高，又要在同频中最近入栈”完整试慢版本；样例里已经能看到“规律是值到频率、频率到栈、当前最大频率三者共同维护”，所以优化前必须先能说清慢版本枚举了哪些候选。

\textbf{暴力过程。} 拿 push 5, 7, 5, 7, 4, 5 后 pop 手算，频率最高是 5，先弹 5；下一次 5 和 7 同频，弹最近达到该频率的 7。

\textbf{重复工作。} 题面模型：弹出时既要频率最高，又要在同频中最近入栈。暴力版的慢点要落到本题：慢版本会围绕“弹出时既要频率最高，又要在同频中最近入栈”反复寻找最近边界或最近未完成对象。样例规律是“规律是值到频率、频率到栈、当前最大频率三者共同维护”，说明旧候选一旦遇到当前输入就能被批量结算；继续为每个位置单独向左或向右扫描就是浪费。后面的优化只消掉这类重复工作，不能改变答案集合。

\textbf{优化条件。} 本题的关键观察是：哈希表记录值频率，频率到栈记录达到该频率的入栈顺序，维护当前最大频率。这句话必须能翻译成可维护状态；如果题目条件改变后观察不再成立，就不能继续套原来的写法。

\textbf{相邻状态。} 规律是值到频率、频率到栈、当前最大频率三者共同维护。把这个特例继续拆开看：栈里的对象都是答案尚未确定的候选；一旦当前元素触发弹栈，被弹出的候选必须已经同时拿到了左边界和右边界，未来元素不再有资格改变它的答案。

\textbf{状态复用。} 把反复寻找最近边界改成维护未完成候选；新元素只和栈顶比较，连续弹出一批已经被当前元素解决的对象。本题落点是：规律是值到频率、频率到栈、当前最大频率三者共同维护。手算时只改被新输入影响的那一小部分，而不是回头重算完整候选。

\textbf{指针和字段。} 状态字段要能说清：栈内对象的业务含义、当前输入、弹出时结算的对象、左边界和右边界；栈中存下标还是值要明确。手算时按这个协议跟踪：拿 push 5, 7, 5, 7, 4, 5 后 pop 手算，频率最高是 5，先弹 5；下一次 5 和 7 同频，弹最近达到该频率的 7。然后按协议复查：手算时记录栈内元素的业务含义，不只记录数值。入栈表示答案未定，弹栈表示当前事件已经让它的答案定下来。

\textbf{代码落点。} 弹栈前确认栈非空，弹出后再读取新的左边界；相等元素和哨兵高度要有统一策略。关键代码行必须能逐句翻译成“旧状态怎样变成新状态”，否则就只是模板。

\textbf{正确性。} 证明从弹出时机开始：被弹出的对象在当前时刻答案已经确定，未来元素无法再改变它。

\textbf{边界实验。} 先用重复 push、pop 后最大频率下降、同频最近性、空栈不在题意内做反例检查。实验时准备重复 push、pop 后最大频率下降、同频最近性、空栈不在题意内，打印每次入栈、弹栈和结算对象。重点观察相等元素、空栈、哨兵和最后一次结算。

\textbf{复杂度变化。} 暴力为每个位置寻找边界常见为平方级；单调栈让每个元素最多入栈一次、出栈一次，因此主体扫描为线性时间。

\textbf{迁移追问。} 如果题目改成“它是栈和哈希表按频率维度组合的设计题”，先判断原状态是否仍然覆盖新答案集合，再决定是微调边界、改 value 语义，还是换算法。

\noindent\textbf{LeetCode 981 基于时间的键值存储。}

\textbf{题目说明。} LeetCode 981 基于时间的键值存储 的题面入口要先还原成具体任务：同一个 key 下的记录按时间戳递增，查询是不超过目标时间的最新值。

\textbf{样例入口。} 拿 key=foo 的记录 (1, bar)、(4, baz)，查询 t=3 手算，答案应是时间不超过 3 的最后一条，也就是 bar。查询 t=4 才是 baz。

\textbf{暴力基线。} 慢版本按顺序扫描下标、逐个试答案值，或直接检查候选直到遇到目标边界。它先保证候选空间覆盖完整，但每次只排除一个候选，浪费了题目给出的顺序、比较或可行性边界。放到本题就是：先按“同一个 key 下的记录按时间戳递增，查询是不超过目标时间的最新值”完整试慢版本；样例里已经能看到“规律是每个 key 内部时间递增，二分第一个大于查询时间的位置，再回退一格”，所以优化前必须先能说清慢版本枚举了哪些候选。

\textbf{暴力过程。} 拿 key=foo 的记录 (1, bar)、(4, baz)，查询 t=3 手算，答案应是时间不超过 3 的最后一条，也就是 bar。查询 t=4 才是 baz。

\textbf{重复工作。} 题面模型：同一个 key 下的记录按时间戳递增，查询是不超过目标时间的最新值。暴力版的慢点要落到本题：慢版本会按顺序试“同一个 key 下的记录按时间戳递增，查询是不超过目标时间的最新值”里的候选；而样例规律已经指出“规律是每个 key 内部时间递增，二分第一个大于查询时间的位置，再回退一格”。一次比较或一次可行性检查能丢掉一整段候选，继续线性试就是浪费。后面的优化只消掉这类重复工作，不能改变答案集合。

\textbf{优化条件。} 本题的关键观察是：哈希表定位 key 后，在该 key 的时间列表中二分第一个大于目标时间的位置并回退一格。这句话必须能翻译成可维护状态；如果题目条件改变后观察不再成立，就不能继续套原来的写法。

\textbf{相邻状态。} 规律是每个 key 内部时间递增，二分第一个大于查询时间的位置，再回退一格。把这个特例继续拆开看：每次比较 mid 之后，必须能指出哪一侧整体不可能包含目标，或哪一侧整体已经满足可行性。二分的核心不是 mid 公式，而是被丢弃区间确实不含答案。

\textbf{状态复用。} 把逐个试候选改成维护答案所在区间；边界谓词、可行性检查或局部比较给出分支方向，每个分支都必须能排除一侧候选。本题落点是：规律是每个 key 内部时间递增，二分第一个大于查询时间的位置，再回退一格。手算时只改被新输入影响的那一小部分，而不是回头重算完整候选。

\textbf{指针和字段。} 状态字段要能说清：left、right、mid、mid 处比较值或检查返回值、答案区间含义、返回值语义；每一轮更新都要保留目标位置或第一个可行值。手算时按这个协议跟踪：拿 key=foo 的记录 (1, bar)、(4, baz)，查询 t=3 手算，答案应是时间不超过 3 的最后一条，也就是 bar。查询 t=4 才是 baz。然后按协议复查：手算时列出 left、mid、right，以及 mid 处比较结果、可行性检查结果或局部趋势。每一步都要写出被丢弃的一侧为什么不含答案。

\textbf{代码落点。} 检查函数或比较分支单独写清，mid 不溢出，循环退出条件和返回值配套；每个分支都要解释丢弃区间。关键代码行必须能逐句翻译成“旧状态怎样变成新状态”，否则就只是模板。

\textbf{正确性。} 证明从排除一侧开始：答案空间题证明谓词单调，局部比较题证明保留下来的区间仍含目标；不能只靠经验猜测左右边界。

\textbf{边界实验。} 先用key 不存在、查询早于第一条记录、相同 key 多次 set、时间戳递增假设做反例检查。实验时覆盖key 不存在、查询早于第一条记录、相同 key 多次 set、时间戳递增假设，逐轮打印 left、mid、right、比较值或检查结果。只要有一次被排除区间仍含答案，二分不变量就错了。

\textbf{复杂度变化。} 暴力逐个试候选是线性或和值域成正比；二分每次排除一半候选，复杂度变成对数级，答案二分还要乘上单次检查成本。

\textbf{迁移追问。} 如果题目改成“如果时间戳无序到达，必须先排序或改成有序表维护每个 key 的记录”，先判断原状态是否仍然覆盖新答案集合，再决定是微调边界、改 value 语义，还是换算法。

\noindent\textbf{LeetCode 1091 二进制矩阵中的最短路径。}

\textbf{题目说明。} LeetCode 1091 二进制矩阵中的最短路径 的题面入口要先还原成具体任务：八方向网格无权边，答案是从左上到右下的最少格子数。

\textbf{样例入口。} 拿 [[0, 1], [1, 0]] 手算，从左上可以斜着一步到右下，路径长度是 2；若起点或终点为 1 直接不可达。

\textbf{暴力基线。} 慢版本先按题意画出完整状态图，用普通搜索跑通可达性。这个过程用于确认八方向网格无权边，答案是从左上到右下的最少格子数的节点和边，之后再讨论访问标记、层次或代价顺序。放到本题就是：先按“八方向网格无权边，答案是从左上到右下的最少格子数”完整试慢版本；样例里已经能看到“规律是八方向 BFS，第一次到达终点就是最短路径长度”，所以优化前必须先能说清慢版本枚举了哪些候选。

\textbf{暴力过程。} 拿 [[0, 1], [1, 0]] 手算，从左上可以斜着一步到右下，路径长度是 2；若起点或终点为 1 直接不可达。

\textbf{重复工作。} 题面模型：八方向网格无权边，答案是从左上到右下的最少格子数。暴力版的慢点要落到本题：慢版本会围绕“八方向网格无权边，答案是从左上到右下的最少格子数”反复展开同一批状态。样例规律是“规律是八方向 BFS，第一次到达终点就是最短路径长度”，说明必须把节点、边、访问标记和资源维度一次建清；同一状态不应重复入队，不同资源状态也不能误合并。后面的优化只消掉这类重复工作，不能改变答案集合。

\textbf{优化条件。} 本题的关键观察是：BFS 入队时标记访问，层数或距离随状态保存。这句话必须能翻译成可维护状态；如果题目条件改变后观察不再成立，就不能继续套原来的写法。

\textbf{相邻状态。} 规律是八方向 BFS，第一次到达终点就是最短路径长度。把这个特例继续拆开看：状态节点不一定等于原始位置，还可能包含钥匙、剩余资源、时间或访问集合。visited 只能合并未来能力完全相同的状态，否则会把合法路径剪掉。

\textbf{状态复用。} 把重复搜索压成访问标记、距离数组或拓扑状态；邻居生成也要从全量扫描变成结构化枚举。本题落点是：规律是八方向 BFS，第一次到达终点就是最短路径长度。手算时只改被新输入影响的那一小部分，而不是回头重算完整候选。

\textbf{指针和字段。} 状态字段要能说清：状态编号、邻接生成方式、访问标记、距离或层数、队列内容；若状态含资源，visited 不能只按位置记录。手算时按这个协议跟踪：拿 [[0, 1], [1, 0]] 手算，从左上可以斜着一步到右下，路径长度是 2；若起点或终点为 1 直接不可达。然后按协议复查：手算时画队列或栈，状态必须包含题目需要的所有资源。入队时标记还是出队时标记，要和去重语义一致。

\textbf{代码落点。} 入队时标记还是出队时标记必须一致；方向数组集中定义；多源 BFS 要一次性把所有源点放入初始队列。关键代码行必须能逐句翻译成“旧状态怎样变成新状态”，否则就只是模板。

\textbf{正确性。} 证明从可达性开始：搜索覆盖所有合法边能到达的状态，访问标记只删除重复状态而不删除不同未来能力。

\textbf{边界实验。} 先用起点或终点阻塞、单格、八方向越界、无路可达做反例检查。实验时覆盖起点或终点阻塞、单格、八方向越界、无路可达，并打印入队时的完整状态。若同一位置但资源不同的状态被合并，很多图题会出现隐蔽错误。

\textbf{复杂度变化。} 暴力重复从多个起点搜索会反复遍历图；正确建模和访问标记后，BFS 或 DFS 通常按节点数加边数计算，状态图题则按完整状态数计算。

\textbf{迁移追问。} 如果题目改成“边权相同才能用 BFS，若每格有代价就要最短路”，先判断原状态是否仍然覆盖新答案集合，再决定是微调边界、改 value 语义，还是换算法。

\noindent\textbf{LeetCode 1235 规划兼职工作。}

\textbf{题目说明。} LeetCode 1235 规划兼职工作 的题面入口要先还原成具体任务：每份工作有开始、结束和收益，选择不冲突工作使收益最大。

\textbf{样例入口。} 拿 jobs (1, 3, 50), (2, 4, 10), (3, 5, 40), (3, 6, 70) 手算，选 (1, 3, 50) 后还能接 (3, 6, 70)，总 120。

\textbf{暴力基线。} 慢版本先写递归搜索树：节点表示处理位置、剩余资源和当前选择。只要不同路径反复到达同一个节点，就说明需要把这个节点命名成 DP 状态。放到本题就是：先按“每份工作有开始、结束和收益，选择不冲突工作使收益最大”完整试慢版本；样例里已经能看到“规律是按结束时间排序，dp[i] 在不选当前和选当前加上前一个不冲突工作之间取最大”，所以优化前必须先能说清慢版本枚举了哪些候选。

\textbf{暴力过程。} 拿 jobs (1, 3, 50), (2, 4, 10), (3, 5, 40), (3, 6, 70) 手算，选 (1, 3, 50) 后还能接 (3, 6, 70)，总 120。

\textbf{重复工作。} 题面模型：每份工作有开始、结束和收益，选择不冲突工作使收益最大。暴力版的慢点要落到本题：慢版本会在“每份工作有开始、结束和收益，选择不冲突工作使收益最大”的选择树里反复到达相同参数。样例规律是“规律是按结束时间排序，dp[i] 在不选当前和选当前加上前一个不冲突工作之间取最大”，说明这个参数组合可以命名成状态，算过之后后面直接复用。后面的优化只消掉这类重复工作，不能改变答案集合。

\textbf{优化条件。} 本题的关键观察是：按结束时间排序，dp[i] 在不选当前和选当前加上兼容前驱之间取最大，前驱用二分找。这句话必须能翻译成可维护状态；如果题目条件改变后观察不再成立，就不能继续套原来的写法。

\textbf{相邻状态。} 规律是按结束时间排序，dp[i] 在不选当前和选当前加上前一个不冲突工作之间取最大。把这个特例继续拆开看：先问暴力搜索里哪些不同路径到达同一个后续问题，再给这个后续问题命名为状态。状态必须包含未来决策需要的全部信息，转移只从更小且已经算清的状态来。

\textbf{状态复用。} 把重复递归节点命名为状态；遍历顺序只是在保证依赖状态已经算完。本题落点是：规律是按结束时间排序，dp[i] 在不选当前和选当前加上前一个不冲突工作之间取最大。手算时只改被新输入影响的那一小部分，而不是回头重算完整候选。

\textbf{指针和字段。} 状态字段要能说清：状态下标、状态值语义、初始化、转移来源、遍历顺序；空间压缩时还要指出读到的是旧值还是新值。手算时按这个协议跟踪：拿 jobs (1, 3, 50), (2, 4, 10), (3, 5, 40), (3, 6, 70) 手算，选 (1, 3, 50) 后还能接 (3, 6, 70)，总 120。然后按协议复查：手算时先写一个很小的 DP 表或记忆化搜索记录。每个格子旁边写它来自哪些更小状态。

\textbf{代码落点。} 先写未压缩版本校验转移；压缩前后用小表对拍；不可达哨兵参与加法前要检查溢出。关键代码行必须能逐句翻译成“旧状态怎样变成新状态”，否则就只是模板。

\textbf{正确性。} 证明从最后一步开始：任意最优解的最后一步都在转移枚举中，转移来源已经由更小状态正确计算。

\textbf{边界实验。} 先用结束等于开始是否兼容、收益相同、工作排序、空列表做反例检查。实验时覆盖结束等于开始是否兼容、收益相同、工作排序、空列表，先打印完整 DP 表，再比较空间压缩版本。若压缩版不同，优先检查遍历顺序和旧值保存。

\textbf{复杂度变化。} 暴力递归会重复计算相同子问题，常见指数级；DP 把每个状态算一次，时间是状态数乘单个转移成本，空间是状态表大小或压缩后的大小。

\textbf{迁移追问。} 如果题目改成“这是排序、二分和 DP 的组合，不是按收益贪心”，先判断原状态是否仍然覆盖新答案集合，再决定是微调边界、改 value 语义，还是换算法。

\noindent\textbf{LeetCode 1248 统计优美子数组。}

\textbf{题目说明。} LeetCode 1248 统计优美子数组 的题面入口要先还原成具体任务：恰好 K 个奇数可转成前缀奇数计数差为 K。

\textbf{样例入口。} 拿 [1, 1, 2, 1, 1]、k=3 手算，奇数个数前缀从 0 到 3 时，需要历史奇数个数 0；每个匹配历史前缀对应一个优美子数组。

\textbf{暴力基线。} 慢版本按题意两两比较、枚举区间、扫描历史或持续迭代并查重。把检查条件写出来后，通常能看到当前状态只缺一个历史状态、规范键或已经访问过的中间状态；这正是从平方枚举或重复迭代走向哈希表的入口。放到本题就是：先按“恰好 K 个奇数可转成前缀奇数计数差为 K”完整试慢版本；样例里已经能看到“规律是把奇数看成 1、偶数看成 0，转成前缀和计数”，所以优化前必须先能说清慢版本枚举了哪些候选。

\textbf{暴力过程。} 拿 [1, 1, 2, 1, 1]、k=3 手算，奇数个数前缀从 0 到 3 时，需要历史奇数个数 0；每个匹配历史前缀对应一个优美子数组。

\textbf{重复工作。} 题面模型：恰好 K 个奇数可转成前缀奇数计数差为 K。暴力版的慢点要落到本题：慢版本会为“恰好 K 个奇数可转成前缀奇数计数差为 K”反复寻找历史状态；而样例规律已经指出“规律是把奇数看成 1、偶数看成 0，转成前缀和计数”。这些补数、前缀状态、规范键、半边和或访问过的中间状态可以用一个 key 归类，不必每次重新扫描或重复比较。后面的优化只消掉这类重复工作，不能改变答案集合。

\textbf{优化条件。} 本题的关键观察是：哈希表保存某个奇数计数出现次数，当前计数查询 count - K。这句话必须能翻译成可维护状态；如果题目条件改变后观察不再成立，就不能继续套原来的写法。

\textbf{相邻状态。} 规律是把奇数看成 1、偶数看成 0，转成前缀和计数。把这个特例继续拆开看：每个合法答案都要还原成当前状态和某个历史状态的配对；哈希表的 key 负责定位历史状态，value 负责表达次数、最早位置或存在性。先查还是先写，必须由是否允许当前前缀和自己配对来决定。

\textbf{状态复用。} 把回头扫描、两两比较或重复状态检测改成查表、分桶或访问记录；key 是当前题目需要的补数、前缀状态、规范表示、半边和或中间状态，value 根据目标选择次数、下标、列表或存在性。本题落点是：规律是把奇数看成 1、偶数看成 0，转成前缀和计数。手算时只改被新输入影响的那一小部分，而不是回头重算完整候选。

\textbf{指针和字段。} 状态字段要能说清：当前状态或规范键、需要查询的历史 key、哈希表 value 语义、答案累计值或分桶结果；空前缀、空字符串、重复状态或初始状态要单独说明。手算时按这个协议跟踪：拿 [1, 1, 2, 1, 1]、k=3 手算，奇数个数前缀从 0 到 3 时，需要历史奇数个数 0；每个匹配历史前缀对应一个优美子数组。然后按协议复查：手算时列出当前输入、当前状态或规范键、需要查询的历史 key、查到的 value、更新后的哈希表。这样能看出先查后插、先分桶后输出，还是先检测重复状态。

\textbf{代码落点。} 先查询还是先插入要由题意决定；key 构造必须无歧义；哈希表 value 的默认插入副作用要可控。关键代码行必须能逐句翻译成“旧状态怎样变成新状态”，否则就只是模板。

\textbf{正确性。} 证明从状态等价开始：任意合法答案都对应当前状态和某个历史 key、规范键或访问状态，查到的 key 也能反推出合法答案或合法分组。

\textbf{边界实验。} 先用K 为零、全偶数、全奇数、从开头开始的区间做反例检查。实验时把K 为零、全偶数、全奇数、从开头开始的区间加入对拍集合，用平方暴力和优化版比较。失败时先看初始历史状态、查询更新顺序和哈希表 value 类型。

\textbf{复杂度变化。} 暴力两两比较、枚举区间、扫描历史或重复迭代常见为平方级或更多；把历史状态、规范键或访问状态放进哈希表后，每步只做常数次查询和更新，通常是线性时间、线性空间。

\textbf{迁移追问。} 如果题目改成“也可用至多 K 减至多 K-1 的滑动窗口”，先判断原状态是否仍然覆盖新答案集合，再决定是微调边界、改 value 语义，还是换算法。

\noindent\textbf{LeetCode 1438 绝对差不超过限制的最长连续子数组。}

\textbf{题目说明。} LeetCode 1438 绝对差不超过限制的最长连续子数组 的题面入口要先还原成具体任务：窗口合法性取决于窗口内最大值和最小值之差。

\textbf{样例入口。} 拿 [8, 2, 4, 7]、limit=4 手算，窗口 [8, 2] 最大最小差 6，不合法，左端必须右移；窗口 [2, 4] 合法。

\textbf{暴力基线。} 暴力枚举每个窗口，重新扫描窗口求最大值和最小值，若 max-min<=limit 就更新最长长度。

\textbf{暴力过程。} 以 [8, 2, 4, 7]、limit=4 为例，窗口 [8, 2] 差为 6 不合法，窗口 [2, 4] 合法。暴力每次都重新找最大最小。

\textbf{重复工作。} 题面模型：窗口合法性取决于窗口内最大值和最小值之差。暴力版的慢点要落到本题：浪费在相邻窗口重复求极值。右端新增一个元素后，只有被它支配的旧极值候选需要删除；左端移动时只需清理过期下标。后面的优化只消掉这类重复工作，不能改变答案集合。

\textbf{优化条件。} 本题的关键观察是：两个单调队列分别维护最大和最小，超过限制时移动左端并清理过期下标。这句话必须能翻译成可维护状态；如果题目条件改变后观察不再成立，就不能继续套原来的写法。

\textbf{相邻状态。} 规律是两个单调队列分别维护最大值和最小值，差值超限时移动左端并清理过期下标。把这个特例继续拆开看：右端进入只带来一个新元素，左端离开只删掉一个旧元素；只有当旧左端被移走后不可能再产生更优答案时，窗口才可以单向推进。若这个单调淘汰条件不存在，就要换成前缀和、队列或其他状态。

\textbf{状态复用。} 用 maxDeque 维护递减下标队列，队首是窗口最大值；minDeque 维护递增下标队列，队首是窗口最小值。差值超限时移动 left 并清理过期。手算时只改被新输入影响的那一小部分，而不是回头重算完整候选。

\textbf{指针和字段。} 状态字段要能说清：left/right 是窗口边界，maxDeque.front/minDeque.front 是当前极值下标，answer 是最长合法窗口长度。手算时按这个协议跟踪：加入 2 后，最大队首为 8、最小队首为 2，差 6 超限；left 从 0 移到 1，过期的 8 被清理，窗口 [2] 合法。

\textbf{代码落点。} 关键是入队前从队尾弹出被新值支配的下标；收缩时如果队首下标小于 left 就弹出。关键代码行必须能逐句翻译成“旧状态怎样变成新状态”，否则就只是模板。

\textbf{正确性。} 单调队列保留的都是可能成为当前或未来极值的下标，被新元素支配的旧下标既更旧又不更优，永远不会再成为极值。

\textbf{边界实验。} 先用limit 为零、重复元素、最大最小同时过期、窗口收缩多次做反例检查。实验时重点跑limit 为零、重复元素、最大最小同时过期、窗口收缩多次，并打印左右端、窗口计数和答案。若左端发生回退，或者窗口失去合法性后仍更新答案，就能很快定位错误。

\textbf{复杂度变化。} 暴力 O(\ensuremath{n^{2}})，两个单调队列每个下标入队出队各一次，时间 O(n)，空间 O(n)。

\textbf{迁移追问。} 如果题目改成“它说明窗口状态可以是动态极值，不只是频次或总和”，先判断原状态是否仍然覆盖新答案集合，再决定是微调边界、改 value 语义，还是换算法。

\noindent\textbf{LeetCode 1462 课程表 IV。}

\textbf{题目说明。} LeetCode 1462 课程表 IV 的题面入口要先还原成具体任务：大量先修关系查询需要预处理课程之间的可达性。

\textbf{样例入口。} 拿 0->1->2 和查询 0 是否先修 2 手算，直接边只能回答 0->1、1->2，不能回答传递关系。若查询很多，每次 DFS 会重复走相同路径。

\textbf{暴力基线。} 慢版本先按题意画出完整状态图，用普通搜索跑通可达性。这个过程用于确认大量先修关系查询需要预处理课程之间的可达性的节点和边，之后再讨论访问标记、层次或代价顺序。放到本题就是：先按“大量先修关系查询需要预处理课程之间的可达性”完整试慢版本；样例里已经能看到“规律是先做可达性闭包或拓扑合并祖先集合，再 O(1) 回答查询”，所以优化前必须先能说清慢版本枚举了哪些候选。

\textbf{暴力过程。} 拿 0->1->2 和查询 0 是否先修 2 手算，直接边只能回答 0->1、1->2，不能回答传递关系。若查询很多，每次 DFS 会重复走相同路径。

\textbf{重复工作。} 题面模型：大量先修关系查询需要预处理课程之间的可达性。暴力版的慢点要落到本题：慢版本会围绕“大量先修关系查询需要预处理课程之间的可达性”反复展开同一批状态。样例规律是“规律是先做可达性闭包或拓扑合并祖先集合，再 O(1) 回答查询”，说明必须把节点、边、访问标记和资源维度一次建清；同一状态不应重复入队，不同资源状态也不能误合并。后面的优化只消掉这类重复工作，不能改变答案集合。

\textbf{优化条件。} 本题的关键观察是：可以用 Floyd 传递闭包、每个点 BFS，或拓扑 DP 合并祖先集合，查询时直接查可达表。这句话必须能翻译成可维护状态；如果题目条件改变后观察不再成立，就不能继续套原来的写法。

\textbf{相邻状态。} 规律是先做可达性闭包或拓扑合并祖先集合，再 O(1) 回答查询。把这个特例继续拆开看：状态节点不一定等于原始位置，还可能包含钥匙、剩余资源、时间或访问集合。visited 只能合并未来能力完全相同的状态，否则会把合法路径剪掉。

\textbf{状态复用。} 把重复搜索压成访问标记、距离数组或拓扑状态；邻居生成也要从全量扫描变成结构化枚举。本题落点是：规律是先做可达性闭包或拓扑合并祖先集合，再 O(1) 回答查询。手算时只改被新输入影响的那一小部分，而不是回头重算完整候选。

\textbf{指针和字段。} 状态字段要能说清：状态编号、邻接生成方式、访问标记、距离或层数、队列内容；若状态含资源，visited 不能只按位置记录。手算时按这个协议跟踪：拿 0->1->2 和查询 0 是否先修 2 手算，直接边只能回答 0->1、1->2，不能回答传递关系。若查询很多，每次 DFS 会重复走相同路径。然后按协议复查：手算时画队列或栈，状态必须包含题目需要的所有资源。入队时标记还是出队时标记，要和去重语义一致。

\textbf{代码落点。} 入队时标记还是出队时标记必须一致；方向数组集中定义；多源 BFS 要一次性把所有源点放入初始队列。关键代码行必须能逐句翻译成“旧状态怎样变成新状态”，否则就只是模板。

\textbf{正确性。} 证明从可达性开始：搜索覆盖所有合法边能到达的状态，访问标记只删除重复状态而不删除不同未来能力。

\textbf{边界实验。} 先用课程之间没有关系、自身是否算先修、查询很多、图中按题意应为 DAG做反例检查。实验时覆盖课程之间没有关系、自身是否算先修、查询很多、图中按题意应为 DAG，并打印入队时的完整状态。若同一位置但资源不同的状态被合并，很多图题会出现隐蔽错误。

\textbf{复杂度变化。} 暴力重复从多个起点搜索会反复遍历图；正确建模和访问标记后，BFS 或 DFS 通常按节点数加边数计算，状态图题则按完整状态数计算。

\textbf{迁移追问。} 如果题目改成“如果查询很少，逐次搜索可能更省预处理成本；多查询时闭包更稳”，先判断原状态是否仍然覆盖新答案集合，再决定是微调边界、改 value 语义，还是换算法。

\noindent\textbf{LeetCode 1514 概率最大的路径。}

\textbf{题目说明。} LeetCode 1514 概率最大的路径 的题面入口要先还原成具体任务：路径概率是边概率乘积，目标是最大乘积路径。

\textbf{样例入口。} 拿 0->1 概率 0.5、1->2 概率 0.5、0->2 概率 0.2 手算，经 1 到 2 的概率 0.25 更大。

\textbf{暴力基线。} 慢版本可以枚举路径，或先用普通队列试跑一个小图。它会暴露步数最少和代价最小是否一致；一旦不一致，就必须围绕代价组织状态。放到本题就是：先按“路径概率是边概率乘积，目标是最大乘积路径”完整试慢版本；样例里已经能看到“规律是把距离改成最大概率，优先队列每次弹出当前概率最高的节点并做乘法松弛”，所以优化前必须先能说清慢版本枚举了哪些候选。

\textbf{暴力过程。} 拿 0->1 概率 0.5、1->2 概率 0.5、0->2 概率 0.2 手算，经 1 到 2 的概率 0.25 更大。

\textbf{重复工作。} 题面模型：路径概率是边概率乘积，目标是最大乘积路径。暴力版的慢点要落到本题：慢版本会围绕“路径概率是边概率乘积，目标是最大乘积路径”枚举路径或按错误顺序扩展状态。样例规律是“规律是把距离改成最大概率，优先队列每次弹出当前概率最高的节点并做乘法松弛”，说明队列或堆的弹出顺序必须和代价定义一致，旧状态为什么能跳过也要讲清。后面的优化只消掉这类重复工作，不能改变答案集合。

\textbf{优化条件。} 本题的关键观察是：用大顶堆每次扩展当前概率最高节点，乘以边概率得到候选概率。这句话必须能翻译成可维护状态；如果题目条件改变后观察不再成立，就不能继续套原来的写法。

\textbf{相邻状态。} 规律是把距离改成最大概率，优先队列每次弹出当前概率最高的节点并做乘法松弛。把这个特例继续拆开看：队列或堆弹出的顺序必须和代价规则一致；旧状态被跳过，是因为已经有更小代价到达同一状态。只要边权或代价合成方式改变，就要重新证明扩展顺序。

\textbf{状态复用。} 把路径枚举压成距离数组和优先队列；每次只扩展当前最有希望的未过期状态。本题落点是：规律是把距离改成最大概率，优先队列每次弹出当前概率最高的节点并做乘法松弛。手算时只改被新输入影响的那一小部分，而不是回头重算完整候选。

\textbf{指针和字段。} 状态字段要能说清：距离数组、优先队列状态、松弛候选、旧状态判定、不可达哨兵；资源约束题还要把资源维度放进状态。手算时按这个协议跟踪：拿 0->1 概率 0.5、1->2 概率 0.5、0->2 概率 0.2 手算，经 1 到 2 的概率 0.25 更大。然后按协议复查：手算时记录每次弹出的代价、当前距离数组和松弛结果。旧状态被弹出时为什么跳过，也要写在表里。

\textbf{代码落点。} 松弛时只在候选更优时更新；priority\_queue 弹出旧状态要跳过；距离加法用更宽整数。关键代码行必须能逐句翻译成“旧状态怎样变成新状态”，否则就只是模板。

\textbf{正确性。} 证明从扩展顺序开始：当前弹出的状态在代价规则下已经不会被未来路径改得更优。

\textbf{边界实验。} 先用不可达、概率为零、起点等于终点、浮点比较做反例检查。实验时覆盖不可达、概率为零、起点等于终点、浮点比较，打印每次弹出的代价和是否过期。若普通 BFS 与 Dijkstra 结果不同，要能解释边权条件造成的差异。

\textbf{复杂度变化。} 暴力枚举路径可能指数级；Dijkstra 在邻接表和堆实现下按边数、点数和堆操作计算，0-1 BFS 可按节点数加边数计算。

\textbf{迁移追问。} 如果题目改成“取负对数后可转成最短路，但直接最大堆更直观”，先判断原状态是否仍然覆盖新答案集合，再决定是微调边界、改 value 语义，还是换算法。

\noindent\textbf{LeetCode 1675 数组的最小偏移量。}

\textbf{题目说明。} LeetCode 1675 数组的最小偏移量 的题面入口要先还原成具体任务：奇数可以先乘二，偶数可以不断除二，目标是缩小当前最大最小差。

\textbf{样例入口。} 拿 [1, 2, 3, 4] 手算，奇数可先翻倍成偶数，之后只能不断把当前最大偶数减半；每次记录最大值和当前最小值差。

\textbf{暴力基线。} 慢版本每轮扫描全部候选来找最大或最小，再更新候选集合。它的答案选择过程清楚，但动态集合每变化一次就重新找极值，代价太高。放到本题就是：先按“奇数可以先乘二，偶数可以不断除二，目标是缩小当前最大最小差”完整试慢版本；样例里已经能看到“规律是把所有数变到可下降的最高状态，再用大顶堆逐步降低最大值”，所以优化前必须先能说清慢版本枚举了哪些候选。

\textbf{暴力过程。} 拿 [1, 2, 3, 4] 手算，奇数可先翻倍成偶数，之后只能不断把当前最大偶数减半；每次记录最大值和当前最小值差。

\textbf{重复工作。} 题面模型：奇数可以先乘二，偶数可以不断除二，目标是缩小当前最大最小差。暴力版的慢点要落到本题：慢版本会围绕“奇数可以先乘二，偶数可以不断除二，目标是缩小当前最大最小差”反复扫描动态候选集找极值。样例规律是“规律是把所有数变到可下降的最高状态，再用大顶堆逐步降低最大值”，说明每轮真正需要的是堆顶边界，而不是把候选全集重新排序。后面的优化只消掉这类重复工作，不能改变答案集合。

\textbf{优化条件。} 本题的关键观察是：把所有数规范到可降低的偶数形态，用大顶堆维护当前最大值，同时记录当前最小值。这句话必须能翻译成可维护状态；如果题目条件改变后观察不再成立，就不能继续套原来的写法。

\textbf{相邻状态。} 规律是把所有数变到可下降的最高状态，再用大顶堆逐步降低最大值。把这个特例继续拆开看：堆顶必须正好代表下一步最需要处理的候选；若堆里可能有过期对象，读取堆顶前要先清理。堆只解决动态极值，不自动证明被弹出或被丢弃的元素无用。

\textbf{状态复用。} 把每轮全量扫描改成堆顶查询；插入、弹出和过期清理共同维护动态候选集。本题落点是：规律是把所有数变到可下降的最高状态，再用大顶堆逐步降低最大值。手算时只改被新输入影响的那一小部分，而不是回头重算完整候选。

\textbf{指针和字段。} 状态字段要能说清：堆元素字段、堆顶比较规则、候选是否过期、答案读取时机；如果需要懒删除，还要保存删除计数。手算时按这个协议跟踪：拿 [1, 2, 3, 4] 手算，奇数可先翻倍成偶数，之后只能不断把当前最大偶数减半；每次记录最大值和当前最小值差。然后按协议复查：手算时记录堆顶、候选集合和是否有过期元素。每次使用堆顶前都要确认它仍然属于当前问题。

\textbf{代码落点。} 比较器方向先用小样例验证；每次使用堆顶前清理过期项；堆中保存下标时要能回查原数据。关键代码行必须能逐句翻译成“旧状态怎样变成新状态”，否则就只是模板。

\textbf{正确性。} 证明从候选覆盖开始：所有可能成为下一步答案的对象都在堆中，堆顶过期时不会被使用。

\textbf{边界实验。} 先用全奇数、全偶数、最大值变奇数停止、数值范围做反例检查。实验时覆盖全奇数、全偶数、最大值变奇数停止、数值范围，记录堆顶是否过期、比较器是否让正确候选在堆顶、K 或窗口大小为边界值时是否稳定。

\textbf{复杂度变化。} 暴力每轮扫描或排序动态候选，会把线性扫描或排序反复发生；堆把取极值、插入和删除边界控制在对数级。

\textbf{迁移追问。} 如果题目改成“这是堆维护动态极值和状态空间规范化的结合”，先判断原状态是否仍然覆盖新答案集合，再决定是微调边界、改 value 语义，还是换算法。

\begin{principlebox}{本节复盘}
阅读本节后，请回到相邻章节的 LeetCode 案例，例如 LeetCode 875 爱吃香蕉、LeetCode 72 编辑距离、LeetCode 743 网络延迟时间、LeetCode 215 数组第 K 大，把暴力瓶颈、优化依据、状态不变量、C++ 容器成本和边界测试重新写一遍。能从这些具体题迁移到变体题，才算真正掌握。
\end{principlebox}

\section{二刷复盘总览}

这一组材料用于第二遍和第三遍复习。结构上不再把 265 道题直接堆成题号列表，而是按题型分成若干大节；每个题型大节内部只保留复盘目标、核心题卡和自测标准几个小节。

\section{数组与原地维护：二刷复盘卡}
这一大节只围绕一个题型复盘，不把题号直接铺成目录。阅读时先看复盘目标，再读核心题卡，最后用自测标准检查自己是否能独立重讲。

\subsection{复盘目标}

追问通常会落在原地修改、稳定顺序、输出规模和整数溢出上。请准备说明为什么保存下标比保存指针更稳，为什么某个位置一旦写入就不会再被未来元素破坏。

\subsection{核心题卡}

\noindent\textbf{LeetCode 5 最长回文子串。} 模型：回文围绕中心对称，中心可以是字符也可以是两个字符之间。推导重点：中心扩展避免枚举子串后再逐个检查；每次扩展只比较新增左右边界。

\textbf{二刷读题。} 读题时先问答案落在哪些下标上，是保留元素、重排元素、计算片段，还是从两侧合成贡献。本题可以压成“回文围绕中心对称，中心可以是字符也可以是两个字符之间”，所以后续所有指针都必须有区间语义。先遮住答案，只保留题面，复述候选对象、合法性条件和输出目标。

\textbf{暴力回放。} 慢版本从下标枚举开始：把每个可能位置、片段、中心或边界都按题意检查一遍。它能直接验证回文围绕中心对称，中心可以是字符也可以是两个字符之间，缺点是同一段信息会被重复读取、重复搬移或重复比较。二刷时先写慢版本或口述慢版本，用它确认不会漏答案。

\textbf{特例回放。} 拿 aba 和 abba 手算，回文不是任意子串都要重查，而是围绕中心向两边同步比较。aba 的中心是 b，abba 的中心在两个 b 中间；每扩一层只新增左右两个字符。规律是枚举 2n-1 个中心，每个中心只在新增边界相等时继续扩展。把这个特例继续拆开看：先标出已经确认的前缀或后缀，再标出未知区；能被丢掉的候选，必须是以后再也不会改变已确认区域语义的候选。这样才算从例子推出数组不变量，而不是只记一次操作。复盘时把这个特例压成状态字段：当前下标、已处理前缀边界、未知区边界、答案变量；若有前后缀贡献，还要说明左侧累计值和右侧累计值分别何时更新；再用边界 偶数长度回文、单字符、全相同字符、多个同长答案 反查初始化、更新顺序和退出条件。

\textbf{状态回放。} 手算路线：手算时把下标语义写在纸上：当前候选来自哪个位置、哪段区间、哪个中心或哪条边界；每轮状态变化后，都要能指出哪些候选已经确定、哪些候选还没有看。状态字段：当前下标、已处理前缀边界、未知区边界、答案变量；若有前后缀贡献，还要说明左侧累计值和右侧累计值分别何时更新。状态升级：把重复工作收进边界变量、前后缀值、中心扩展结果或慢指针中；状态只保存已经被证明稳定的信息，不让相同区间或相同位置被反复完整检查。

\textbf{证明和边界。} 证明从区间不变量开始：已处理区域已经满足题意，未知区域还没有承诺，每次推进不会破坏已处理区域。边界：偶数长度回文、单字符、全相同字符、多个同长答案。变体：若要求回文子序列，应转成区间 DP，不能用中心扩展。

\noindent\textbf{LeetCode 26 删除有序数组重复项。} 模型：有序数组让重复元素相邻，慢指针表示下一个唯一值写入位置。推导重点：快指针扫描，只有发现新值才写入慢指针并推进。

\textbf{二刷读题。} 读题时先问答案落在哪些下标上，是保留元素、重排元素、计算片段，还是从两侧合成贡献。本题可以压成“有序数组让重复元素相邻，慢指针表示下一个唯一值写入位置”，所以后续所有指针都必须有区间语义。先遮住答案，只保留题面，复述候选对象、合法性条件和输出目标。

\textbf{暴力回放。} 慢版本从下标枚举开始：把每个可能位置、片段、中心或边界都按题意检查一遍。它能直接验证有序数组让重复元素相邻，慢指针表示下一个唯一值写入位置，缺点是同一段信息会被重复读取、重复搬移或重复比较。二刷时先写慢版本或口述慢版本，用它确认不会漏答案。

\textbf{特例回放。} 拿 [1,1,2] 手算，write 先停在第一个 1 后面；读到第二个 1 时，它和已保留尾值相同，不写入；读到 2 时写到 write 位置，前缀变成 [1,2]。规律是慢指针左侧始终是不重复前缀，快指针只在遇到新值时推进慢指针；空数组和全重复数组只改变初始化，不改变不变量。把这个特例继续拆开看：先标出已经确认的前缀或后缀，再标出未知区；能被丢掉的候选，必须是以后再也不会改变已确认区域语义的候选。这样才算从例子推出数组不变量，而不是只记一次操作。复盘时把这个特例压成状态字段：当前下标、已处理前缀边界、未知区边界、答案变量；若有前后缀贡献，还要说明左侧累计值和右侧累计值分别何时更新；再用边界 空数组、全重复、无重复、返回长度不等于物理容量 反查初始化、更新顺序和退出条件。

\textbf{状态回放。} 手算路线：手算时把下标语义写在纸上：当前候选来自哪个位置、哪段区间、哪个中心或哪条边界；每轮状态变化后，都要能指出哪些候选已经确定、哪些候选还没有看。状态字段：当前下标、已处理前缀边界、未知区边界、答案变量；若有前后缀贡献，还要说明左侧累计值和右侧累计值分别何时更新。状态升级：把重复工作收进边界变量、前后缀值、中心扩展结果或慢指针中；状态只保存已经被证明稳定的信息，不让相同区间或相同位置被反复完整检查。

\textbf{证明和边界。} 证明从区间不变量开始：已处理区域已经满足题意，未知区域还没有承诺，每次推进不会破坏已处理区域。边界：空数组、全重复、无重复、返回长度不等于物理容量。变体：若允许每个元素最多保留两次，只需改变写入条件为和前两个比较。

\noindent\textbf{LeetCode 27 移除元素。} 模型：原地过滤，慢指针左侧始终是不等于目标值的稳定前缀。推导重点：保持顺序用快慢指针；不保持顺序可用左右交换减少写入。

\textbf{二刷读题。} 读题时先问答案落在哪些下标上，是保留元素、重排元素、计算片段，还是从两侧合成贡献。本题可以压成“原地过滤，慢指针左侧始终是不等于目标值的稳定前缀”，所以后续所有指针都必须有区间语义。先遮住答案，只保留题面，复述候选对象、合法性条件和输出目标。

\textbf{暴力回放。} 慢版本从下标枚举开始：把每个可能位置、片段、中心或边界都按题意检查一遍。它能直接验证原地过滤，慢指针左侧始终是不等于目标值的稳定前缀，缺点是同一段信息会被重复读取、重复搬移或重复比较。二刷时先写慢版本或口述慢版本，用它确认不会漏答案。

\textbf{特例回放。} 拿 [3,2,2,3]、val=3 手算，稳定写法读到 3 不写，读到两个 2 依次写入前缀，最后有效区间是 [2,2]。若题目不要求顺序，也可以把左侧 3 和右侧非 3 交换来减少写入。规律是先确认是否稳定，再决定快慢指针还是左右交换；所有元素都被移除时返回长度 0。把这个特例继续拆开看：先标出已经确认的前缀或后缀，再标出未知区；能被丢掉的候选，必须是以后再也不会改变已确认区域语义的候选。这样才算从例子推出数组不变量，而不是只记一次操作。复盘时把这个特例压成状态字段：当前下标、已处理前缀边界、未知区边界、答案变量；若有前后缀贡献，还要说明左侧累计值和右侧累计值分别何时更新；再用边界 所有元素都被移除、没有元素被移除、目标值在尾部密集 反查初始化、更新顺序和退出条件。

\textbf{状态回放。} 手算路线：手算时把下标语义写在纸上：当前候选来自哪个位置、哪段区间、哪个中心或哪条边界；每轮状态变化后，都要能指出哪些候选已经确定、哪些候选还没有看。状态字段：当前下标、已处理前缀边界、未知区边界、答案变量；若有前后缀贡献，还要说明左侧累计值和右侧累计值分别何时更新。状态升级：把重复工作收进边界变量、前后缀值、中心扩展结果或慢指针中；状态只保存已经被证明稳定的信息，不让相同区间或相同位置被反复完整检查。

\textbf{证明和边界。} 证明从区间不变量开始：已处理区域已经满足题意，未知区域还没有承诺，每次推进不会破坏已处理区域。边界：所有元素都被移除、没有元素被移除、目标值在尾部密集。变体：工程里要先确认是否要求稳定顺序，不同要求对应不同写法。

\noindent\textbf{LeetCode 31 下一个排列。} 模型：字典序结构由最长非递增后缀决定，枢轴是后缀前一个位置。推导重点：从右找第一个上升点，再从右找刚好更大的元素交换，最后反转后缀。

\textbf{二刷读题。} 读题时先问答案落在哪些下标上，是保留元素、重排元素、计算片段，还是从两侧合成贡献。本题可以压成“字典序结构由最长非递增后缀决定，枢轴是后缀前一个位置”，所以后续所有指针都必须有区间语义。先遮住答案，只保留题面，复述候选对象、合法性条件和输出目标。

\textbf{暴力回放。} 慢版本从下标枚举开始：把每个可能位置、片段、中心或边界都按题意检查一遍。它能直接验证字典序结构由最长非递增后缀决定，枢轴是后缀前一个位置，缺点是同一段信息会被重复读取、重复搬移或重复比较。二刷时先写慢版本或口述慢版本，用它确认不会漏答案。

\textbf{特例回放。} 拿 [1,3,2] 手算，后缀 [3,2] 已经非递增，枢轴是 1；从右找刚好比 1 大的 2 交换，得到 [2,3,1]，再反转后缀成 [2,1,3]。规律是最长非递增后缀已经是该前缀下的最大排列，必须提升枢轴并把后缀变成最小结构；完全降序时整体反转成最小排列。把这个特例继续拆开看：先标出已经确认的前缀或后缀，再标出未知区；能被丢掉的候选，必须是以后再也不会改变已确认区域语义的候选。这样才算从例子推出数组不变量，而不是只记一次操作。复盘时把这个特例压成状态字段：当前下标、已处理前缀边界、未知区边界、答案变量；若有前后缀贡献，还要说明左侧累计值和右侧累计值分别何时更新；再用边界 完全降序、重复数字、只有一个元素、交换后后缀必须最小化 反查初始化、更新顺序和退出条件。

\textbf{状态回放。} 手算路线：手算时把下标语义写在纸上：当前候选来自哪个位置、哪段区间、哪个中心或哪条边界；每轮状态变化后，都要能指出哪些候选已经确定、哪些候选还没有看。状态字段：当前下标、已处理前缀边界、未知区边界、答案变量；若有前后缀贡献，还要说明左侧累计值和右侧累计值分别何时更新。状态升级：把重复工作收进边界变量、前后缀值、中心扩展结果或慢指针中；状态只保存已经被证明稳定的信息，不让相同区间或相同位置被反复完整检查。

\textbf{证明和边界。} 证明从区间不变量开始：已处理区域已经满足题意，未知区域还没有承诺，每次推进不会破坏已处理区域。边界：完全降序、重复数字、只有一个元素、交换后后缀必须最小化。变体：若要求上一个排列，比较方向和后缀处理对称变化。

\noindent\textbf{LeetCode 41 缺失的第一个正数。} 模型：答案只可能落在一到 n 加一之间，数组本身可以充当值到位置的映射。推导重点：把合法正数 x 尽量交换到下标 x 减一，最后第一个位置不匹配就是答案。

\textbf{二刷读题。} 读题时先问答案落在哪些下标上，是保留元素、重排元素、计算片段，还是从两侧合成贡献。本题可以压成“答案只可能落在一到 n 加一之间，数组本身可以充当值到位置的映射”，所以后续所有指针都必须有区间语义。先遮住答案，只保留题面，复述候选对象、合法性条件和输出目标。

\textbf{暴力回放。} 慢版本从下标枚举开始：把每个可能位置、片段、中心或边界都按题意检查一遍。它能直接验证答案只可能落在一到 n 加一之间，数组本身可以充当值到位置的映射，缺点是同一段信息会被重复读取、重复搬移或重复比较。二刷时先写慢版本或口述慢版本，用它确认不会漏答案。

\textbf{特例回放。} 拿 [3,4,-1,1] 手算。答案只可能在 1 到 n+1；值 3 应放到下标 2，值 4 放到下标 3，-1 不管，1 放到下标 0。放完后第一个位置不等于 i+1 的地方就是缺失值。再试 [1,1]，若不判断目标位置已有相同值会无限交换。规律是数组本身当哈希表，只处理合法正数且避免重复交换。把这个特例继续拆开看：先标出已经确认的前缀或后缀，再标出未知区；能被丢掉的候选，必须是以后再也不会改变已确认区域语义的候选。这样才算从例子推出数组不变量，而不是只记一次操作。复盘时把这个特例压成状态字段：当前下标、已处理前缀边界、未知区边界、答案变量；若有前后缀贡献，还要说明左侧累计值和右侧累计值分别何时更新；再用边界 非正数、大于 n 的数、重复值导致无限交换、数组已经完整包含一到 n 反查初始化、更新顺序和退出条件。

\textbf{状态回放。} 手算路线：手算时把下标语义写在纸上：当前候选来自哪个位置、哪段区间、哪个中心或哪条边界；每轮状态变化后，都要能指出哪些候选已经确定、哪些候选还没有看。状态字段：当前下标、已处理前缀边界、未知区边界、答案变量；若有前后缀贡献，还要说明左侧累计值和右侧累计值分别何时更新。状态升级：把重复工作收进边界变量、前后缀值、中心扩展结果或慢指针中；状态只保存已经被证明稳定的信息，不让相同区间或相同位置被反复完整检查。

\textbf{证明和边界。} 证明从区间不变量开始：已处理区域已经满足题意，未知区域还没有承诺，每次推进不会破坏已处理区域。边界：非正数、大于 n 的数、重复值导致无限交换、数组已经完整包含一到 n。变体：若不能修改输入，只能使用哈希集合或额外布尔数组换取更简单边界。

\noindent\textbf{LeetCode 54 螺旋矩阵。} 模型：四个边界收缩模拟一圈一圈访问。推导重点：每走完一条边就收缩对应边界，并在走反方向前检查是否仍合法。

\textbf{二刷读题。} 读题时先问答案落在哪些下标上，是保留元素、重排元素、计算片段，还是从两侧合成贡献。本题可以压成“四个边界收缩模拟一圈一圈访问”，所以后续所有指针都必须有区间语义。先遮住答案，只保留题面，复述候选对象、合法性条件和输出目标。

\textbf{暴力回放。} 慢版本从下标枚举开始：把每个可能位置、片段、中心或边界都按题意检查一遍。它能直接验证四个边界收缩模拟一圈一圈访问，缺点是同一段信息会被重复读取、重复搬移或重复比较。二刷时先写慢版本或口述慢版本，用它确认不会漏答案。

\textbf{特例回放。} 拿 3x3 矩阵手算，先走上边一整行，再走右边一整列，然后收缩 top 和 right；反向走下边和左边前要检查 top<=bottom、left<=right。规律是四个边界每走完一条边就收缩一次，最后剩一行或一列时不能重复访问。把这个特例继续拆开看：先标出已经确认的前缀或后缀，再标出未知区；能被丢掉的候选，必须是以后再也不会改变已确认区域语义的候选。这样才算从例子推出数组不变量，而不是只记一次操作。复盘时把这个特例压成状态字段：当前下标、已处理前缀边界、未知区边界、答案变量；若有前后缀贡献，还要说明左侧累计值和右侧累计值分别何时更新；再用边界 单行、单列、空矩阵、最后一圈只剩一个元素 反查初始化、更新顺序和退出条件。

\textbf{状态回放。} 手算路线：手算时把下标语义写在纸上：当前候选来自哪个位置、哪段区间、哪个中心或哪条边界；每轮状态变化后，都要能指出哪些候选已经确定、哪些候选还没有看。状态字段：当前下标、已处理前缀边界、未知区边界、答案变量；若有前后缀贡献，还要说明左侧累计值和右侧累计值分别何时更新。状态升级：把重复工作收进边界变量、前后缀值、中心扩展结果或慢指针中；状态只保存已经被证明稳定的信息，不让相同区间或相同位置被反复完整检查。

\textbf{证明和边界。} 证明从区间不变量开始：已处理区域已经满足题意，未知区域还没有承诺，每次推进不会破坏已处理区域。边界：单行、单列、空矩阵、最后一圈只剩一个元素。变体：若改成生成矩阵，边界逻辑相同，只是读变成写。

\noindent\textbf{LeetCode 75 颜色分类。} 模型：三指针维护零区、一号区、未知区和二区。推导重点：遇到二时只移动右边界，不移动扫描指针，因为换回来的元素未知。

\textbf{二刷读题。} 读题时先问答案落在哪些下标上，是保留元素、重排元素、计算片段，还是从两侧合成贡献。本题可以压成“三指针维护零区、一号区、未知区和二区”，所以后续所有指针都必须有区间语义。先遮住答案，只保留题面，复述候选对象、合法性条件和输出目标。

\textbf{暴力回放。} 慢版本从下标枚举开始：把每个可能位置、片段、中心或边界都按题意检查一遍。它能直接验证三指针维护零区、一号区、未知区和二区，缺点是同一段信息会被重复读取、重复搬移或重复比较。二刷时先写慢版本或口述慢版本，用它确认不会漏答案。

\textbf{特例回放。} 拿 [2,0,1] 手算，mid 指向 2 时和右边界交换，换回来的元素未知，所以 mid 不能前进；再看到 1 才前进。规律是 [0,low) 是 0， [low,mid) 是 1， (high,n) 是 2，未知区只在 mid 到 high 之间收缩。把这个特例继续拆开看：先标出已经确认的前缀或后缀，再标出未知区；能被丢掉的候选，必须是以后再也不会改变已确认区域语义的候选。这样才算从例子推出数组不变量，而不是只记一次操作。复盘时把这个特例压成状态字段：当前下标、已处理前缀边界、未知区边界、答案变量；若有前后缀贡献，还要说明左侧累计值和右侧累计值分别何时更新；再用边界 全零、全二、已经有序、逆序、交换后漏检查 反查初始化、更新顺序和退出条件。

\textbf{状态回放。} 手算路线：手算时把下标语义写在纸上：当前候选来自哪个位置、哪段区间、哪个中心或哪条边界；每轮状态变化后，都要能指出哪些候选已经确定、哪些候选还没有看。状态字段：当前下标、已处理前缀边界、未知区边界、答案变量；若有前后缀贡献，还要说明左侧累计值和右侧累计值分别何时更新。状态升级：把重复工作收进边界变量、前后缀值、中心扩展结果或慢指针中；状态只保存已经被证明稳定的信息，不让相同区间或相同位置被反复完整检查。

\textbf{证明和边界。} 证明从区间不变量开始：已处理区域已经满足题意，未知区域还没有承诺，每次推进不会破坏已处理区域。边界：全零、全二、已经有序、逆序、交换后漏检查。变体：这是荷兰国旗问题，能推广到三类分区的原地维护。

\noindent\textbf{LeetCode 88 合并两个有序数组。} 模型：第一个数组尾部有空位，适合从后往前写入最终最大值。推导重点：逆向填充避免覆盖尚未读取的 nums1 元素。

\textbf{二刷读题。} 读题时先问答案落在哪些下标上，是保留元素、重排元素、计算片段，还是从两侧合成贡献。本题可以压成“第一个数组尾部有空位，适合从后往前写入最终最大值”，所以后续所有指针都必须有区间语义。先遮住答案，只保留题面，复述候选对象、合法性条件和输出目标。

\textbf{暴力回放。} 慢版本从下标枚举开始：把每个可能位置、片段、中心或边界都按题意检查一遍。它能直接验证第一个数组尾部有空位，适合从后往前写入最终最大值，缺点是同一段信息会被重复读取、重复搬移或重复比较。二刷时先写慢版本或口述慢版本，用它确认不会漏答案。

\textbf{特例回放。} 拿 nums1=[1,3,0]、nums2=[2] 手算，如果从前往后写，2 会覆盖 nums1 中尚未比较的 3；从尾部写，先比较 3 和 2，把 3 放到最后，再把 2 放到中间。规律是利用尾部空位逆向填充，避免覆盖未读元素。把这个特例继续拆开看：先标出已经确认的前缀或后缀，再标出未知区；能被丢掉的候选，必须是以后再也不会改变已确认区域语义的候选。这样才算从例子推出数组不变量，而不是只记一次操作。复盘时把这个特例压成状态字段：当前下标、已处理前缀边界、未知区边界、答案变量；若有前后缀贡献，还要说明左侧累计值和右侧累计值分别何时更新；再用边界 第二个数组为空、第一个有效元素为空、重复值、尾部剩余 反查初始化、更新顺序和退出条件。

\textbf{状态回放。} 手算路线：手算时把下标语义写在纸上：当前候选来自哪个位置、哪段区间、哪个中心或哪条边界；每轮状态变化后，都要能指出哪些候选已经确定、哪些候选还没有看。状态字段：当前下标、已处理前缀边界、未知区边界、答案变量；若有前后缀贡献，还要说明左侧累计值和右侧累计值分别何时更新。状态升级：把重复工作收进边界变量、前后缀值、中心扩展结果或慢指针中；状态只保存已经被证明稳定的信息，不让相同区间或相同位置被反复完整检查。

\textbf{证明和边界。} 证明从区间不变量开始：已处理区域已经满足题意，未知区域还没有承诺，每次推进不会破坏已处理区域。边界：第二个数组为空、第一个有效元素为空、重复值、尾部剩余。变体：若没有额外空位，只能新开数组或使用更复杂原地归并。

\noindent\textbf{LeetCode 189 轮转数组。} 模型：右移 k 位等价于把数组切成两段后换序。推导重点：三次反转能原地完成：整体、前 k、后 n 减 k。

\textbf{二刷读题。} 读题时先问答案落在哪些下标上，是保留元素、重排元素、计算片段，还是从两侧合成贡献。本题可以压成“右移 k 位等价于把数组切成两段后换序”，所以后续所有指针都必须有区间语义。先遮住答案，只保留题面，复述候选对象、合法性条件和输出目标。

\textbf{暴力回放。} 慢版本从下标枚举开始：把每个可能位置、片段、中心或边界都按题意检查一遍。它能直接验证右移 k 位等价于把数组切成两段后换序，缺点是同一段信息会被重复读取、重复搬移或重复比较。二刷时先写慢版本或口述慢版本，用它确认不会漏答案。

\textbf{特例回放。} 拿 [1,2,3,4]、k=2 手算，目标是 [3,4,1,2]；整体反转得 [4,3,2,1]，再反转前 2 个得 [3,4,2,1]，反转后半得 [3,4,1,2]。规律是右移 k 等价于两段换序，k 要先对 n 取模。把这个特例继续拆开看：先标出已经确认的前缀或后缀，再标出未知区；能被丢掉的候选，必须是以后再也不会改变已确认区域语义的候选。这样才算从例子推出数组不变量，而不是只记一次操作。复盘时把这个特例压成状态字段：当前下标、已处理前缀边界、未知区边界、答案变量；若有前后缀贡献，还要说明左侧累计值和右侧累计值分别何时更新；再用边界 k 为零、k 大于 n、空数组、单元素 反查初始化、更新顺序和退出条件。

\textbf{状态回放。} 手算路线：手算时把下标语义写在纸上：当前候选来自哪个位置、哪段区间、哪个中心或哪条边界；每轮状态变化后，都要能指出哪些候选已经确定、哪些候选还没有看。状态字段：当前下标、已处理前缀边界、未知区边界、答案变量；若有前后缀贡献，还要说明左侧累计值和右侧累计值分别何时更新。状态升级：把重复工作收进边界变量、前后缀值、中心扩展结果或慢指针中；状态只保存已经被证明稳定的信息，不让相同区间或相同位置被反复完整检查。

\textbf{证明和边界。} 证明从区间不变量开始：已处理区域已经满足题意，未知区域还没有承诺，每次推进不会破坏已处理区域。边界：k 为零、k 大于 n、空数组、单元素。变体：环状替换也可做，但要处理最大公约数个循环。

\noindent\textbf{LeetCode 238 除自身以外数组的乘积。} 模型：不能用除法时，答案由左侧乘积和右侧乘积组成。推导重点：先写入左侧前缀积，再从右向左乘右侧后缀积。

\textbf{二刷读题。} 读题时先问答案落在哪些下标上，是保留元素、重排元素、计算片段，还是从两侧合成贡献。本题可以压成“不能用除法时，答案由左侧乘积和右侧乘积组成”，所以后续所有指针都必须有区间语义。先遮住答案，只保留题面，复述候选对象、合法性条件和输出目标。

\textbf{暴力回放。} 慢版本从下标枚举开始：把每个可能位置、片段、中心或边界都按题意检查一遍。它能直接验证不能用除法时，答案由左侧乘积和右侧乘积组成，缺点是同一段信息会被重复读取、重复搬移或重复比较。二刷时先写慢版本或口述慢版本，用它确认不会漏答案。

\textbf{特例回放。} 拿 [1,2,3,4] 手算，第一遍把每个位置左侧乘积写进答案，得到 [1,1,2,6]；第二遍从右侧累乘 4、12、24 并乘回去。规律是答案等于左侧乘积乘右侧乘积，不能使用除法时就用两次扫描合并贡献。把这个特例继续拆开看：先标出已经确认的前缀或后缀，再标出未知区；能被丢掉的候选，必须是以后再也不会改变已确认区域语义的候选。这样才算从例子推出数组不变量，而不是只记一次操作。复盘时把这个特例压成状态字段：当前下标、已处理前缀边界、未知区边界、答案变量；若有前后缀贡献，还要说明左侧累计值和右侧累计值分别何时更新；再用边界 一个零、多个零、负数、乘积溢出 反查初始化、更新顺序和退出条件。

\textbf{状态回放。} 手算路线：手算时把下标语义写在纸上：当前候选来自哪个位置、哪段区间、哪个中心或哪条边界；每轮状态变化后，都要能指出哪些候选已经确定、哪些候选还没有看。状态字段：当前下标、已处理前缀边界、未知区边界、答案变量；若有前后缀贡献，还要说明左侧累计值和右侧累计值分别何时更新。状态升级：把重复工作收进边界变量、前后缀值、中心扩展结果或慢指针中；状态只保存已经被证明稳定的信息，不让相同区间或相同位置被反复完整检查。

\textbf{证明和边界。} 证明从区间不变量开始：已处理区域已经满足题意，未知区域还没有承诺，每次推进不会破坏已处理区域。边界：一个零、多个零、负数、乘积溢出。变体：这个题展示输出数组可作为状态存储的一部分。

\noindent\textbf{LeetCode 283 移动零。} 模型：稳定保留非零元素顺序，把零集中到尾部。推导重点：慢指针写入下一个非零位置，最后补零或用交换。

\textbf{二刷读题。} 读题时先问答案落在哪些下标上，是保留元素、重排元素、计算片段，还是从两侧合成贡献。本题可以压成“稳定保留非零元素顺序，把零集中到尾部”，所以后续所有指针都必须有区间语义。先遮住答案，只保留题面，复述候选对象、合法性条件和输出目标。

\textbf{暴力回放。} 慢版本从下标枚举开始：把每个可能位置、片段、中心或边界都按题意检查一遍。它能直接验证稳定保留非零元素顺序，把零集中到尾部，缺点是同一段信息会被重复读取、重复搬移或重复比较。二刷时先写慢版本或口述慢版本，用它确认不会漏答案。

\textbf{特例回放。} 拿 [0,1,0,3] 手算，write 指向下一个非零写入位置，读到 1 写到 0 号位，读到 3 写到 1 号位，最后把后面补 0。规律是先稳定压缩非零前缀，再填零；不能在扫描时盲目交换导致顺序变化。把这个特例继续拆开看：先标出已经确认的前缀或后缀，再标出未知区；能被丢掉的候选，必须是以后再也不会改变已确认区域语义的候选。这样才算从例子推出数组不变量，而不是只记一次操作。复盘时把这个特例压成状态字段：当前下标、已处理前缀边界、未知区边界、答案变量；若有前后缀贡献，还要说明左侧累计值和右侧累计值分别何时更新；再用边界 全零、无零、零在开头结尾、稳定顺序要求 反查初始化、更新顺序和退出条件。

\textbf{状态回放。} 手算路线：手算时把下标语义写在纸上：当前候选来自哪个位置、哪段区间、哪个中心或哪条边界；每轮状态变化后，都要能指出哪些候选已经确定、哪些候选还没有看。状态字段：当前下标、已处理前缀边界、未知区边界、答案变量；若有前后缀贡献，还要说明左侧累计值和右侧累计值分别何时更新。状态升级：把重复工作收进边界变量、前后缀值、中心扩展结果或慢指针中；状态只保存已经被证明稳定的信息，不让相同区间或相同位置被反复完整检查。

\textbf{证明和边界。} 证明从区间不变量开始：已处理区域已经满足题意，未知区域还没有承诺，每次推进不会破坏已处理区域。边界：全零、无零、零在开头结尾、稳定顺序要求。变体：若目标是移动指定值，模型与移除元素相同。

\noindent\textbf{LeetCode 581 最短无序连续子数组。} 模型：要找破坏全局有序性的最左和最右位置。推导重点：从左维护历史最大值确定右边界，从右维护历史最小值确定左边界。

\textbf{二刷读题。} 读题时先问答案落在哪些下标上，是保留元素、重排元素、计算片段，还是从两侧合成贡献。本题可以压成“要找破坏全局有序性的最左和最右位置”，所以后续所有指针都必须有区间语义。先遮住答案，只保留题面，复述候选对象、合法性条件和输出目标。

\textbf{暴力回放。} 慢版本从下标枚举开始：把每个可能位置、片段、中心或边界都按题意检查一遍。它能直接验证要找破坏全局有序性的最左和最右位置，缺点是同一段信息会被重复读取、重复搬移或重复比较。二刷时先写慢版本或口述慢版本，用它确认不会漏答案。

\textbf{特例回放。} 拿 [2,6,4,8,10,9,15] 手算，从左扫描时 4 小于此前最大 6，说明右边界至少到 4；9 小于此前最大 10，右边界更新到 9。规律是左右扫描分别找破坏有序的位置，最后得到最短需要排序的连续段。把这个特例继续拆开看：先标出已经确认的前缀或后缀，再标出未知区；能被丢掉的候选，必须是以后再也不会改变已确认区域语义的候选。这样才算从例子推出数组不变量，而不是只记一次操作。复盘时把这个特例压成状态字段：当前下标、已处理前缀边界、未知区边界、答案变量；若有前后缀贡献，还要说明左侧累计值和右侧累计值分别何时更新；再用边界 已经有序、逆序、重复值、局部乱序 反查初始化、更新顺序和退出条件。

\textbf{状态回放。} 手算路线：手算时把下标语义写在纸上：当前候选来自哪个位置、哪段区间、哪个中心或哪条边界；每轮状态变化后，都要能指出哪些候选已经确定、哪些候选还没有看。状态字段：当前下标、已处理前缀边界、未知区边界、答案变量；若有前后缀贡献，还要说明左侧累计值和右侧累计值分别何时更新。状态升级：把重复工作收进边界变量、前后缀值、中心扩展结果或慢指针中；状态只保存已经被证明稳定的信息，不让相同区间或相同位置被反复完整检查。

\textbf{证明和边界。} 证明从区间不变量开始：已处理区域已经满足题意，未知区域还没有承诺，每次推进不会破坏已处理区域。边界：已经有序、逆序、重复值、局部乱序。变体：排序比较法更直观但复杂度更高。

\noindent\textbf{LeetCode 647 回文子串。} 模型：每个回文都有唯一中心或中心间隙。推导重点：对每个中心向两边扩展，扩展成功一次就计数一次。

\textbf{二刷读题。} 读题时先问答案落在哪些下标上，是保留元素、重排元素、计算片段，还是从两侧合成贡献。本题可以压成“每个回文都有唯一中心或中心间隙”，所以后续所有指针都必须有区间语义。先遮住答案，只保留题面，复述候选对象、合法性条件和输出目标。

\textbf{暴力回放。} 慢版本从下标枚举开始：把每个可能位置、片段、中心或边界都按题意检查一遍。它能直接验证每个回文都有唯一中心或中心间隙，缺点是同一段信息会被重复读取、重复搬移或重复比较。二刷时先写慢版本或口述慢版本，用它确认不会漏答案。

\textbf{特例回放。} 拿 aaa 手算，以每个字符为中心有 3 个单字符回文，以两个字符中间为中心有 2 个 aa，以中间 a 为中心还有 aaa。规律是枚举字符中心和间隙中心，每次扩展成功就新增一个回文子串。把这个特例继续拆开看：先标出已经确认的前缀或后缀，再标出未知区；能被丢掉的候选，必须是以后再也不会改变已确认区域语义的候选。这样才算从例子推出数组不变量，而不是只记一次操作。复盘时把这个特例压成状态字段：当前下标、已处理前缀边界、未知区边界、答案变量；若有前后缀贡献，还要说明左侧累计值和右侧累计值分别何时更新；再用边界 空串、全相同字符、偶数回文、奇数回文 反查初始化、更新顺序和退出条件。

\textbf{状态回放。} 手算路线：手算时把下标语义写在纸上：当前候选来自哪个位置、哪段区间、哪个中心或哪条边界；每轮状态变化后，都要能指出哪些候选已经确定、哪些候选还没有看。状态字段：当前下标、已处理前缀边界、未知区边界、答案变量；若有前后缀贡献，还要说明左侧累计值和右侧累计值分别何时更新。状态升级：把重复工作收进边界变量、前后缀值、中心扩展结果或慢指针中；状态只保存已经被证明稳定的信息，不让相同区间或相同位置被反复完整检查。

\textbf{证明和边界。} 证明从区间不变量开始：已处理区域已经满足题意，未知区域还没有承诺，每次推进不会破坏已处理区域。边界：空串、全相同字符、偶数回文、奇数回文。变体：Manacher 算法可线性解决，但中心扩展更适合面试基础。

\subsection{自测标准}

二刷时不要只确认自己见过这道题。请遮住答案，先讲题面模型，再讲暴力基线和瓶颈，然后说出优化结构保存了什么、不变量如何排除错误候选、边界实验如何打破错误实现。

\section{滑动窗口与双指针：二刷复盘卡}
这一大节只围绕一个题型复盘，不把题号直接铺成目录。阅读时先看复盘目标，再读核心题卡，最后用自测标准检查自己是否能独立重讲。

\subsection{复盘目标}

追问通常会问窗口为什么不会漏掉左端、为什么有负数会失效、固定窗口和可变窗口有什么区别。请准备一个反例证明错误窗口条件会漏解。

\subsection{核心题卡}

\noindent\textbf{LeetCode 3 无重复字符的最长子串。} 模型：窗口保存当前无重复子串，右端每次加入一个字符。推导重点：用上次出现位置直接跳过无效左端，左边界只能向右不能回退。

\textbf{二刷读题。} 读题时先确认左右端的语义：它们可能表示连续窗口，也可能表示排序后的双指针候选。本题可以压成“窗口保存当前无重复子串，右端每次加入一个字符”，因此左右端不是模板变量，而是候选排除和状态变化的触发点。先遮住答案，只保留题面，复述候选对象、合法性条件和输出目标。

\textbf{暴力回放。} 慢版本枚举所有左端和右端，或在排序后枚举固定点与另一侧配对。它先完整覆盖题目要求的窗口、片段或有序候选；真正浪费的是相邻候选之间有大量状态可以复用，却被重新统计或重新搜索。二刷时先写慢版本或口述慢版本，用它确认不会漏答案。

\textbf{特例回放。} 拿字符串 abca 手算。窗口 abc 合法，右端加入第二个 a 后，真正冲突的是上一次 a 的位置；左端从 0 跳到 1 后，bca 重新合法。再试 abba：处理第二个 b 时左端跳到 2，后面遇到 a 的上次位置在窗口左侧之外，不能让 left 回退。规律是左边界只向右，状态保存每个字符最近出现位置。把这个特例继续拆开看：右端进入只带来一个新元素，左端离开只删掉一个旧元素；只有当旧左端被移走后不可能再产生更优答案时，窗口才可以单向推进。若这个单调淘汰条件不存在，就要换成前缀和、队列或其他状态。复盘时把这个特例压成状态字段：左端、右端、窗口计数或当前双指针和、合法性指标、当前答案；每次移动边界后，状态必须和候选区间或窗口内容一致；再用边界 空串、全相同字符、重复字符出现在窗口左侧之外、扩展字符集 反查初始化、更新顺序和退出条件。

\textbf{状态回放。} 手算路线：手算时画一张表，列出 left、right、窗口内容、窗口状态和答案。每次右端进入只加一个元素，每次左端离开只删一个元素。状态字段：左端、右端、窗口计数或当前双指针和、合法性指标、当前答案；每次移动边界后，状态必须和候选区间或窗口内容一致。状态升级：把完整窗口检查改成增量维护；右端负责引入新信息，左端负责恢复合法性或压缩答案，窗口状态必须始终和边界一致。

\textbf{证明和边界。} 证明从左端淘汰开始：被移走的左端对应窗口已经检查完，或继续保留只会让状态更差。边界：空串、全相同字符、重复字符出现在窗口左侧之外、扩展字符集。变体：若要求恰好包含 k 种字符，窗口条件从无重复变成种类计数。

\noindent\textbf{LeetCode 11 盛最多水的容器。} 模型：左右指针代表当前选取的两条边，面积由短板和宽度共同决定。推导重点：每次移动短板，因为移动长板只会缩小宽度且不会提高短板。

\textbf{二刷读题。} 读题时先确认左右端的语义：它们可能表示连续窗口，也可能表示排序后的双指针候选。本题可以压成“左右指针代表当前选取的两条边，面积由短板和宽度共同决定”，因此左右端不是模板变量，而是候选排除和状态变化的触发点。先遮住答案，只保留题面，复述候选对象、合法性条件和输出目标。

\textbf{暴力回放。} 慢版本枚举所有左端和右端，或在排序后枚举固定点与另一侧配对。它先完整覆盖题目要求的窗口、片段或有序候选；真正浪费的是相邻候选之间有大量状态可以复用，却被重新统计或重新搜索。二刷时先写慢版本或口述慢版本，用它确认不会漏答案。

\textbf{特例回放。} 拿高度 [1,8,6,2] 的两端手算。若固定宽度从最外侧开始，面积受短板限制；移动高板只会让宽度变小，短板仍不变或更差，不可能得到更大面积。只有移动短板，才有机会提高限制高度。规律是每轮淘汰短板那一侧，因为所有保留它且宽度更小的候选都被当前面积上界支配。把这个特例继续拆开看：右端进入只带来一个新元素，左端离开只删掉一个旧元素；只有当旧左端被移走后不可能再产生更优答案时，窗口才可以单向推进。若这个单调淘汰条件不存在，就要换成前缀和、队列或其他状态。复盘时把这个特例压成状态字段：左端、右端、窗口计数或当前双指针和、合法性指标、当前答案；每次移动边界后，状态必须和候选区间或窗口内容一致；再用边界 两端高度相等、数组长度为二、面积乘法可能溢出 反查初始化、更新顺序和退出条件。

\textbf{状态回放。} 手算路线：手算时画一张表，列出 left、right、窗口内容、窗口状态和答案。每次右端进入只加一个元素，每次左端离开只删一个元素。状态字段：左端、右端、窗口计数或当前双指针和、合法性指标、当前答案；每次移动边界后，状态必须和候选区间或窗口内容一致。状态升级：把完整窗口检查改成增量维护；右端负责引入新信息，左端负责恢复合法性或压缩答案，窗口状态必须始终和边界一致。

\textbf{证明和边界。} 证明从左端淘汰开始：被移走的左端对应窗口已经检查完，或继续保留只会让状态更差。边界：两端高度相等、数组长度为二、面积乘法可能溢出。变体：接雨水计算每个位置贡献，不能把本题双指针证明直接搬过去。

\noindent\textbf{LeetCode 15 三数之和。} 模型：排序后固定第一个数，剩余两数转成有序双指针。推导重点：和太小移动左指针，和太大移动右指针；固定值和左右值都要跳过重复。

\textbf{二刷读题。} 读题时先确认左右端的语义：它们可能表示连续窗口，也可能表示排序后的双指针候选。本题可以压成“排序后固定第一个数，剩余两数转成有序双指针”，因此左右端不是模板变量，而是候选排除和状态变化的触发点。先遮住答案，只保留题面，复述候选对象、合法性条件和输出目标。

\textbf{暴力回放。} 慢版本枚举所有左端和右端，或在排序后枚举固定点与另一侧配对。它先完整覆盖题目要求的窗口、片段或有序候选；真正浪费的是相邻候选之间有大量状态可以复用，却被重新统计或重新搜索。二刷时先写慢版本或口述慢版本，用它确认不会漏答案。

\textbf{特例回放。} 拿 [-1,0,1,2,-1,-4] 排序后手算。固定第一个 -1，剩下问题变成有序两数之和；和太小只能左指针右移，和太大只能右指针左移。再遇到第二个 -1 作为固定点时，会生成重复三元组，所以同层固定点要跳过重复。规律是排序把三数问题拆成枚举固定点加双指针，并在每层处理去重。把这个特例继续拆开看：右端进入只带来一个新元素，左端离开只删掉一个旧元素；只有当旧左端被移走后不可能再产生更优答案时，窗口才可以单向推进。若这个单调淘汰条件不存在，就要换成前缀和、队列或其他状态。复盘时把这个特例压成状态字段：左端、右端、窗口计数或当前双指针和、合法性指标、当前答案；每次移动边界后，状态必须和候选区间或窗口内容一致；再用边界 全零、重复元素很多、没有答案、结果顺序不要求但组合不能重复 反查初始化、更新顺序和退出条件。

\textbf{状态回放。} 手算路线：手算时画一张表，列出 left、right、窗口内容、窗口状态和答案。每次右端进入只加一个元素，每次左端离开只删一个元素。状态字段：左端、右端、窗口计数或当前双指针和、合法性指标、当前答案；每次移动边界后，状态必须和候选区间或窗口内容一致。状态升级：把完整窗口检查改成增量维护；右端负责引入新信息，左端负责恢复合法性或压缩答案，窗口状态必须始终和边界一致。

\textbf{证明和边界。} 证明从左端淘汰开始：被移走的左端对应窗口已经检查完，或继续保留只会让状态更差。边界：全零、重复元素很多、没有答案、结果顺序不要求但组合不能重复。变体：四数之和再外层固定一个数，并用 long long 防止目标和溢出。

\noindent\textbf{LeetCode 16 最接近的三数之和。} 模型：排序后固定一数，双指针寻找最接近目标的两数和。推导重点：每次根据当前和与目标的大小移动指针，同时更新最小绝对差。

\textbf{二刷读题。} 读题时先确认左右端的语义：它们可能表示连续窗口，也可能表示排序后的双指针候选。本题可以压成“排序后固定一数，双指针寻找最接近目标的两数和”，因此左右端不是模板变量，而是候选排除和状态变化的触发点。先遮住答案，只保留题面，复述候选对象、合法性条件和输出目标。

\textbf{暴力回放。} 慢版本枚举所有左端和右端，或在排序后枚举固定点与另一侧配对。它先完整覆盖题目要求的窗口、片段或有序候选；真正浪费的是相邻候选之间有大量状态可以复用，却被重新统计或重新搜索。二刷时先写慢版本或口述慢版本，用它确认不会漏答案。

\textbf{特例回放。} 拿 [-1,2,1,-4]、target=1 手算，排序后固定 -1，左右指针在 1 和 2 上得到和 2，距离目标只差 1；若和偏大，只能让右指针左移，因为右侧数再大只会更远。规律是每次先更新最小差，再按当前和与目标的大小排除一侧候选；差值相等和三数和溢出要单独检查。把这个特例继续拆开看：右端进入只带来一个新元素，左端离开只删掉一个旧元素；只有当旧左端被移走后不可能再产生更优答案时，窗口才可以单向推进。若这个单调淘汰条件不存在，就要换成前缀和、队列或其他状态。复盘时把这个特例压成状态字段：左端、右端、窗口计数或当前双指针和、合法性指标、当前答案；每次移动边界后，状态必须和候选区间或窗口内容一致；再用边界 差值相等、负数、目标很大、三数和溢出 反查初始化、更新顺序和退出条件。

\textbf{状态回放。} 手算路线：手算时画一张表，列出 left、right、窗口内容、窗口状态和答案。每次右端进入只加一个元素，每次左端离开只删一个元素。状态字段：左端、右端、窗口计数或当前双指针和、合法性指标、当前答案；每次移动边界后，状态必须和候选区间或窗口内容一致。状态升级：把完整窗口检查改成增量维护；右端负责引入新信息，左端负责恢复合法性或压缩答案，窗口状态必须始终和边界一致。

\textbf{证明和边界。} 证明从左端淘汰开始：被移走的左端对应窗口已经检查完，或继续保留只会让状态更差。边界：差值相等、负数、目标很大、三数和溢出。变体：若要求返回所有最接近组合，需要记录同差值并去重。

\noindent\textbf{LeetCode 18 四数之和。} 模型：两层固定加一层双指针，关键是剪枝和去重。推导重点：排序提供单调性，外层可用最小可能和和最大可能和提前跳过。

\textbf{二刷读题。} 读题时先确认左右端的语义：它们可能表示连续窗口，也可能表示排序后的双指针候选。本题可以压成“两层固定加一层双指针，关键是剪枝和去重”，因此左右端不是模板变量，而是候选排除和状态变化的触发点。先遮住答案，只保留题面，复述候选对象、合法性条件和输出目标。

\textbf{暴力回放。} 慢版本枚举所有左端和右端，或在排序后枚举固定点与另一侧配对。它先完整覆盖题目要求的窗口、片段或有序候选；真正浪费的是相邻候选之间有大量状态可以复用，却被重新统计或重新搜索。二刷时先写慢版本或口述慢版本，用它确认不会漏答案。

\textbf{特例回放。} 拿 [1,0,-1,0,-2,2]、target=0 手算，排序后先固定 -2，再固定 -1，剩余两数用双指针找 3。找到 [-2,-1,1,2] 后，左右指针相同值必须跳过，否则会重复输出。规律是两层固定把四数降成有序两数，外层用最小可能和与最大可能和剪枝，内外层都要去重。把这个特例继续拆开看：右端进入只带来一个新元素，左端离开只删掉一个旧元素；只有当旧左端被移走后不可能再产生更优答案时，窗口才可以单向推进。若这个单调淘汰条件不存在，就要换成前缀和、队列或其他状态。复盘时把这个特例压成状态字段：左端、右端、窗口计数或当前双指针和、合法性指标、当前答案；每次移动边界后，状态必须和候选区间或窗口内容一致；再用边界 重复元素、目标超出 int、答案很多导致输出空间主导复杂度 反查初始化、更新顺序和退出条件。

\textbf{状态回放。} 手算路线：手算时画一张表，列出 left、right、窗口内容、窗口状态和答案。每次右端进入只加一个元素，每次左端离开只删一个元素。状态字段：左端、右端、窗口计数或当前双指针和、合法性指标、当前答案；每次移动边界后，状态必须和候选区间或窗口内容一致。状态升级：把完整窗口检查改成增量维护；右端负责引入新信息，左端负责恢复合法性或压缩答案，窗口状态必须始终和边界一致。

\textbf{证明和边界。} 证明从左端淘汰开始：被移走的左端对应窗口已经检查完，或继续保留只会让状态更差。边界：重复元素、目标超出 int、答案很多导致输出空间主导复杂度。变体：k 数之和可以递归降维，但工程实现要控制重复和边界。

\noindent\textbf{LeetCode 76 最小覆盖子串。} 模型：窗口内字符计数必须覆盖目标计数，满足后尽量收缩左端。推导重点：用 formed 记录满足需求的字符种类，避免每次全量比较。

\textbf{二刷读题。} 读题时先确认左右端的语义：它们可能表示连续窗口，也可能表示排序后的双指针候选。本题可以压成“窗口内字符计数必须覆盖目标计数，满足后尽量收缩左端”，因此左右端不是模板变量，而是候选排除和状态变化的触发点。先遮住答案，只保留题面，复述候选对象、合法性条件和输出目标。

\textbf{暴力回放。} 慢版本枚举所有左端和右端，或在排序后枚举固定点与另一侧配对。它先完整覆盖题目要求的窗口、片段或有序候选；真正浪费的是相邻候选之间有大量状态可以复用，却被重新统计或重新搜索。二刷时先写慢版本或口述慢版本，用它确认不会漏答案。

\textbf{特例回放。} 拿 s=ADOBEC、t=ABC 手算，右端扩展直到第一次覆盖 A/B/C，此时窗口可行；再移动左端，去掉 A 后窗口不再覆盖，就必须停止收缩。规律是右端负责获得可行，左端负责在可行时尽量缩短；计数表要区分已有字符和需求字符。把这个特例继续拆开看：右端进入只带来一个新元素，左端离开只删掉一个旧元素；只有当旧左端被移走后不可能再产生更优答案时，窗口才可以单向推进。若这个单调淘汰条件不存在，就要换成前缀和、队列或其他状态。复盘时把这个特例压成状态字段：左端、右端、窗口计数或当前双指针和、合法性指标、当前答案；每次移动边界后，状态必须和候选区间或窗口内容一致；再用边界 目标有重复字符、源串缺字符、最小窗口在开头或结尾、大小写 反查初始化、更新顺序和退出条件。

\textbf{状态回放。} 手算路线：手算时画一张表，列出 left、right、窗口内容、窗口状态和答案。每次右端进入只加一个元素，每次左端离开只删一个元素。状态字段：左端、右端、窗口计数或当前双指针和、合法性指标、当前答案；每次移动边界后，状态必须和候选区间或窗口内容一致。状态升级：把完整窗口检查改成增量维护；右端负责引入新信息，左端负责恢复合法性或压缩答案，窗口状态必须始终和边界一致。

\textbf{证明和边界。} 证明从左端淘汰开始：被移走的左端对应窗口已经检查完，或继续保留只会让状态更差。边界：目标有重复字符、源串缺字符、最小窗口在开头或结尾、大小写。变体：若字符集固定，数组计数更快；若是 Unicode，哈希表更通用。

\noindent\textbf{LeetCode 167 两数之和 II 输入有序数组。} 模型：有序数组中左右指针的和随指针移动单调变化。推导重点：当前和太小左指针右移，太大右指针左移，等于目标返回。

\textbf{二刷读题。} 读题时先确认左右端的语义：它们可能表示连续窗口，也可能表示排序后的双指针候选。本题可以压成“有序数组中左右指针的和随指针移动单调变化”，因此左右端不是模板变量，而是候选排除和状态变化的触发点。先遮住答案，只保留题面，复述候选对象、合法性条件和输出目标。

\textbf{暴力回放。} 慢版本枚举所有左端和右端，或在排序后枚举固定点与另一侧配对。它先完整覆盖题目要求的窗口、片段或有序候选；真正浪费的是相邻候选之间有大量状态可以复用，却被重新统计或重新搜索。二刷时先写慢版本或口述慢版本，用它确认不会漏答案。

\textbf{特例回放。} 拿 [2,7,11,15]、target=9 手算，左右和为 17 太大，右指针左移；2+7 正好等于 9。规律是有序数组让和随左指针右移变大、随右指针左移变小，因此每次能排除一侧候选。把这个特例继续拆开看：右端进入只带来一个新元素，左端离开只删掉一个旧元素；只有当旧左端被移走后不可能再产生更优答案时，窗口才可以单向推进。若这个单调淘汰条件不存在，就要换成前缀和、队列或其他状态。复盘时把这个特例压成状态字段：左端、右端、窗口计数或当前双指针和、合法性指标、当前答案；每次移动边界后，状态必须和候选区间或窗口内容一致；再用边界 两个数在边界、重复值、题目下标从一开始、保证有解 反查初始化、更新顺序和退出条件。

\textbf{状态回放。} 手算路线：手算时画一张表，列出 left、right、窗口内容、窗口状态和答案。每次右端进入只加一个元素，每次左端离开只删一个元素。状态字段：左端、右端、窗口计数或当前双指针和、合法性指标、当前答案；每次移动边界后，状态必须和候选区间或窗口内容一致。状态升级：把完整窗口检查改成增量维护；右端负责引入新信息，左端负责恢复合法性或压缩答案，窗口状态必须始终和边界一致。

\textbf{证明和边界。} 证明从左端淘汰开始：被移走的左端对应窗口已经检查完，或继续保留只会让状态更差。边界：两个数在边界、重复值、题目下标从一开始、保证有解。变体：无序版本通常用哈希表，不能直接套双指针。

\noindent\textbf{LeetCode 209 长度最小的子数组。} 模型：正数数组让窗口和随右端增加、随左端减少。推导重点：右端扩展到满足目标后，不断左移缩短并更新答案。

\textbf{二刷读题。} 读题时先确认左右端的语义：它们可能表示连续窗口，也可能表示排序后的双指针候选。本题可以压成“正数数组让窗口和随右端增加、随左端减少”，因此左右端不是模板变量，而是候选排除和状态变化的触发点。先遮住答案，只保留题面，复述候选对象、合法性条件和输出目标。

\textbf{暴力回放。} 慢版本枚举所有左端和右端，或在排序后枚举固定点与另一侧配对。它先完整覆盖题目要求的窗口、片段或有序候选；真正浪费的是相邻候选之间有大量状态可以复用，却被重新统计或重新搜索。二刷时先写慢版本或口述慢版本，用它确认不会漏答案。

\textbf{特例回放。} 拿 [2,3,1,2,4,3]、target=7 手算，右端扩到 4 时窗口和达到 12，可以不断左缩；最后 [4,3] 长度 2 仍满足。规律是全正数保证右扩和变大、左缩和变小，所以每个左端只需被移动一次。把这个特例继续拆开看：右端进入只带来一个新元素，左端离开只删掉一个旧元素；只有当旧左端被移走后不可能再产生更优答案时，窗口才可以单向推进。若这个单调淘汰条件不存在，就要换成前缀和、队列或其他状态。复盘时把这个特例压成状态字段：左端、右端、窗口计数或当前双指针和、合法性指标、当前答案；每次移动边界后，状态必须和候选区间或窗口内容一致；再用边界 不存在答案、单元素达到目标、全正是必要条件 反查初始化、更新顺序和退出条件。

\textbf{状态回放。} 手算路线：手算时画一张表，列出 left、right、窗口内容、窗口状态和答案。每次右端进入只加一个元素，每次左端离开只删一个元素。状态字段：左端、右端、窗口计数或当前双指针和、合法性指标、当前答案；每次移动边界后，状态必须和候选区间或窗口内容一致。状态升级：把完整窗口检查改成增量维护；右端负责引入新信息，左端负责恢复合法性或压缩答案，窗口状态必须始终和边界一致。

\textbf{证明和边界。} 证明从左端淘汰开始：被移走的左端对应窗口已经检查完，或继续保留只会让状态更差。边界：不存在答案、单元素达到目标、全正是必要条件。变体：若含负数，需要前缀和加单调队列。

\noindent\textbf{LeetCode 438 找到字符串中所有字母异位词。} 模型：固定长度窗口内字符频次等于目标频次即为答案。推导重点：右进左出维护计数，窗口长度达到目标长度后检查差异。

\textbf{二刷读题。} 读题时先确认左右端的语义：它们可能表示连续窗口，也可能表示排序后的双指针候选。本题可以压成“固定长度窗口内字符频次等于目标频次即为答案”，因此左右端不是模板变量，而是候选排除和状态变化的触发点。先遮住答案，只保留题面，复述候选对象、合法性条件和输出目标。

\textbf{暴力回放。} 慢版本枚举所有左端和右端，或在排序后枚举固定点与另一侧配对。它先完整覆盖题目要求的窗口、片段或有序候选；真正浪费的是相邻候选之间有大量状态可以复用，却被重新统计或重新搜索。二刷时先写慢版本或口述慢版本，用它确认不会漏答案。

\textbf{特例回放。} 拿 s=cbaebabacd、p=abc 手算，长度为 3 的窗口 cba 频次和 abc 相同，起点 0 是答案；窗口右移时只删除 c、加入 e。规律是固定长度窗口维护字符频次差，重叠答案不能漏。把这个特例继续拆开看：右端进入只带来一个新元素，左端离开只删掉一个旧元素；只有当旧左端被移走后不可能再产生更优答案时，窗口才可以单向推进。若这个单调淘汰条件不存在，就要换成前缀和、队列或其他状态。复盘时把这个特例压成状态字段：左端、右端、窗口计数或当前双指针和、合法性指标、当前答案；每次移动边界后，状态必须和候选区间或窗口内容一致；再用边界 目标比源串长、重复字符、字符集固定、答案重叠 反查初始化、更新顺序和退出条件。

\textbf{状态回放。} 手算路线：手算时画一张表，列出 left、right、窗口内容、窗口状态和答案。每次右端进入只加一个元素，每次左端离开只删一个元素。状态字段：左端、右端、窗口计数或当前双指针和、合法性指标、当前答案；每次移动边界后，状态必须和候选区间或窗口内容一致。状态升级：把完整窗口检查改成增量维护；右端负责引入新信息，左端负责恢复合法性或压缩答案，窗口状态必须始终和边界一致。

\textbf{证明和边界。} 证明从左端淘汰开始：被移走的左端对应窗口已经检查完，或继续保留只会让状态更差。边界：目标比源串长、重复字符、字符集固定、答案重叠。变体：维护差异计数可以避免每次比较整个数组。

\noindent\textbf{LeetCode 567 字符串的排列。} 模型：目标串任意排列出现在源串中，等价于固定长度窗口频次相同。推导重点：窗口长度固定为目标长度，右进左出维护字符计数差异。

\textbf{二刷读题。} 读题时先确认左右端的语义：它们可能表示连续窗口，也可能表示排序后的双指针候选。本题可以压成“目标串任意排列出现在源串中，等价于固定长度窗口频次相同”，因此左右端不是模板变量，而是候选排除和状态变化的触发点。先遮住答案，只保留题面，复述候选对象、合法性条件和输出目标。

\textbf{暴力回放。} 慢版本枚举所有左端和右端，或在排序后枚举固定点与另一侧配对。它先完整覆盖题目要求的窗口、片段或有序候选；真正浪费的是相邻候选之间有大量状态可以复用，却被重新统计或重新搜索。二刷时先写慢版本或口述慢版本，用它确认不会漏答案。

\textbf{特例回放。} 拿 s=eidbaooo、p=ab 手算，窗口 ba 的频次和 ab 相同，返回真；窗口长度固定为 2，每步只进一个字符出一个字符。规律是排列出现等价于固定长度频次相同，不需要枚举所有排列。把这个特例继续拆开看：右端进入只带来一个新元素，左端离开只删掉一个旧元素；只有当旧左端被移走后不可能再产生更优答案时，窗口才可以单向推进。若这个单调淘汰条件不存在，就要换成前缀和、队列或其他状态。复盘时把这个特例压成状态字段：左端、右端、窗口计数或当前双指针和、合法性指标、当前答案；每次移动边界后，状态必须和候选区间或窗口内容一致；再用边界 目标比源串长、重复字符、字符集固定、窗口重叠 反查初始化、更新顺序和退出条件。

\textbf{状态回放。} 手算路线：手算时画一张表，列出 left、right、窗口内容、窗口状态和答案。每次右端进入只加一个元素，每次左端离开只删一个元素。状态字段：左端、右端、窗口计数或当前双指针和、合法性指标、当前答案；每次移动边界后，状态必须和候选区间或窗口内容一致。状态升级：把完整窗口检查改成增量维护；右端负责引入新信息，左端负责恢复合法性或压缩答案，窗口状态必须始终和边界一致。

\textbf{证明和边界。} 证明从左端淘汰开始：被移走的左端对应窗口已经检查完，或继续保留只会让状态更差。边界：目标比源串长、重复字符、字符集固定、窗口重叠。变体：与找到所有异位词同型，只是返回布尔值。

\noindent\textbf{LeetCode 862 和至少为 K 的最短子数组。} 模型：数组允许负数，普通滑动窗口的和不再随左端移动单调变化。推导重点：在前缀和数组上用单调队列维护可能成为最优左端的下标。

\textbf{二刷读题。} 读题时先确认左右端的语义：它们可能表示连续窗口，也可能表示排序后的双指针候选。本题可以压成“数组允许负数，普通滑动窗口的和不再随左端移动单调变化”，因此左右端不是模板变量，而是候选排除和状态变化的触发点。先遮住答案，只保留题面，复述候选对象、合法性条件和输出目标。

\textbf{暴力回放。} 慢版本枚举所有左端和右端，或在排序后枚举固定点与另一侧配对。它先完整覆盖题目要求的窗口、片段或有序候选；真正浪费的是相邻候选之间有大量状态可以复用，却被重新统计或重新搜索。二刷时先写慢版本或口述慢版本，用它确认不会漏答案。

\textbf{特例回放。} 拿 [2,-1,2]、K=3 手算，普通正数窗口会被 -1 破坏单调性；前缀和为 0,2,1,3，只有 3-0 达到 3。规律是在前缀和上用单调队列维护可能的最优左端，队首能满足时结算长度。把这个特例继续拆开看：右端进入只带来一个新元素，左端离开只删掉一个旧元素；只有当旧左端被移走后不可能再产生更优答案时，窗口才可以单向推进。若这个单调淘汰条件不存在，就要换成前缀和、队列或其他状态。复盘时把这个特例压成状态字段：左端、右端、窗口计数或当前双指针和、合法性指标、当前答案；每次移动边界后，状态必须和候选区间或窗口内容一致；再用边界 负数、答案不存在、K 为负或零、队尾支配关系 反查初始化、更新顺序和退出条件。

\textbf{状态回放。} 手算路线：手算时画一张表，列出 left、right、窗口内容、窗口状态和答案。每次右端进入只加一个元素，每次左端离开只删一个元素。状态字段：左端、右端、窗口计数或当前双指针和、合法性指标、当前答案；每次移动边界后，状态必须和候选区间或窗口内容一致。状态升级：把完整窗口检查改成增量维护；右端负责引入新信息，左端负责恢复合法性或压缩答案，窗口状态必须始终和边界一致。

\textbf{证明和边界。} 证明从左端淘汰开始：被移走的左端对应窗口已经检查完，或继续保留只会让状态更差。边界：负数、答案不存在、K 为负或零、队尾支配关系。变体：这是从窗口题升级到前缀和加单调队列的代表。

\noindent\textbf{LeetCode 1438 绝对差不超过限制的最长连续子数组。} 模型：窗口合法性取决于窗口内最大值和最小值之差。推导重点：两个单调队列分别维护最大和最小，超过限制时移动左端并清理过期下标。

\textbf{二刷读题。} 读题时先确认左右端的语义：它们可能表示连续窗口，也可能表示排序后的双指针候选。本题可以压成“窗口合法性取决于窗口内最大值和最小值之差”，因此左右端不是模板变量，而是候选排除和状态变化的触发点。先遮住答案，只保留题面，复述候选对象、合法性条件和输出目标。

\textbf{暴力回放。} 慢版本枚举所有左端和右端，或在排序后枚举固定点与另一侧配对。它先完整覆盖题目要求的窗口、片段或有序候选；真正浪费的是相邻候选之间有大量状态可以复用，却被重新统计或重新搜索。二刷时先写慢版本或口述慢版本，用它确认不会漏答案。

\textbf{特例回放。} 拿 [8,2,4,7]、limit=4 手算，窗口 [8,2] 最大最小差 6，不合法，左端必须右移；窗口 [2,4] 合法。规律是两个单调队列分别维护最大值和最小值，差值超限时移动左端并清理过期下标。把这个特例继续拆开看：右端进入只带来一个新元素，左端离开只删掉一个旧元素；只有当旧左端被移走后不可能再产生更优答案时，窗口才可以单向推进。若这个单调淘汰条件不存在，就要换成前缀和、队列或其他状态。复盘时把这个特例压成状态字段：左端、右端、窗口计数或当前双指针和、合法性指标、当前答案；每次移动边界后，状态必须和候选区间或窗口内容一致；再用边界 limit 为零、重复元素、最大最小同时过期、窗口收缩多次 反查初始化、更新顺序和退出条件。

\textbf{状态回放。} 手算路线：手算时画一张表，列出 left、right、窗口内容、窗口状态和答案。每次右端进入只加一个元素，每次左端离开只删一个元素。状态字段：左端、右端、窗口计数或当前双指针和、合法性指标、当前答案；每次移动边界后，状态必须和候选区间或窗口内容一致。状态升级：把完整窗口检查改成增量维护；右端负责引入新信息，左端负责恢复合法性或压缩答案，窗口状态必须始终和边界一致。

\textbf{证明和边界。} 证明从左端淘汰开始：被移走的左端对应窗口已经检查完，或继续保留只会让状态更差。边界：limit 为零、重复元素、最大最小同时过期、窗口收缩多次。变体：它说明窗口状态可以是动态极值，不只是频次或总和。

\subsection{自测标准}

二刷时不要只确认自己见过这道题。请遮住答案，先讲题面模型，再讲暴力基线和瓶颈，然后说出优化结构保存了什么、不变量如何排除错误候选、边界实验如何打破错误实现。

\section{前缀和与哈希表：二刷复盘卡}
这一大节只围绕一个题型复盘，不把题号直接铺成目录。阅读时先看复盘目标，再读核心题卡，最后用自测标准检查自己是否能独立重讲。

\subsection{复盘目标}

追问通常会问先查询还是先更新、哈希表保存次数还是最早下标、为什么负数不影响前缀和。请准备说明从任意合法区间如何反推到两个前缀状态。

\subsection{核心题卡}

\noindent\textbf{LeetCode 1 两数之和。} 模型：当前元素只需要寻找唯一补数，历史值到下标的映射就是最小状态。推导重点：先查再插入，避免同一个下标被用两次；若数组有序才考虑双指针。

\textbf{二刷读题。} 读题时先把条件改写成当前状态与历史状态的关系。本题可以压成“当前元素只需要寻找唯一补数，历史值到下标的映射就是最小状态”，所以重点是历史表的 key 和 value 分别代表什么。先遮住答案，只保留题面，复述候选对象、合法性条件和输出目标。

\textbf{暴力回放。} 慢版本枚举所有区间或候选对，再逐个检查条件。把检查式写出来后，可以看到当前状态只缺一个历史状态；这正是从平方枚举走向哈希表的入口。二刷时先写慢版本或口述慢版本，用它确认不会漏答案。

\textbf{特例回放。} 拿数组 [2, 7]、target=9 手算：站在 2 时历史为空，不能配对；站在 7 时只需要问历史里有没有 2。再试 [3, 3]、target=6：如果先把当前 3 插入再查，就会把同一个下标用两次。于是规律不是“用哈希表”四个字，而是每个当前位置只和它左侧历史配对，代码必须先查补数再插入当前值。把这个特例继续拆开看：每个合法答案都要还原成当前状态和某个历史状态的配对；哈希表的 key 负责定位历史状态，value 负责表达次数、最早位置或存在性。先查还是先写，必须由是否允许当前前缀和自己配对来决定。复盘时把这个特例压成状态字段：当前前缀状态、需要查询的历史 key、哈希表 value 语义、答案累计值；初始化的空前缀必须单独说明；再用边界 重复数字、目标为偶数且补数等于当前值、没有答案时返回空结果 反查初始化、更新顺序和退出条件。

\textbf{状态回放。} 手算路线：手算时按下标列出当前状态、需要的历史 key、查到的 value、更新后的哈希表。这样能看出先查后插还是先插后查。状态字段：当前前缀状态、需要查询的历史 key、哈希表 value 语义、答案累计值；初始化的空前缀必须单独说明。状态升级：把回头扫描历史改成查表；key 是当前状态需要的历史状态，value 根据目标选择次数、最早位置、最近位置或存在性。

\textbf{证明和边界。} 证明从代数等价开始：任意合法答案都对应当前状态和某个历史 key，查到的历史 key 也能反推出合法答案。边界：重复数字、目标为偶数且补数等于当前值、没有答案时返回空结果。变体：若改成返回所有不重复数对，要先排序或用频次表处理重复。

\noindent\textbf{LeetCode 49 字母异位词分组。} 模型：异位词共享同一个规范表示，可以是排序字符串或字符计数。推导重点：哈希表按规范键分桶，避免两两比较字符串。

\textbf{二刷读题。} 读题时先把条件改写成当前状态与历史状态的关系。本题可以压成“异位词共享同一个规范表示，可以是排序字符串或字符计数”，所以重点是历史表的 key 和 value 分别代表什么。先遮住答案，只保留题面，复述候选对象、合法性条件和输出目标。

\textbf{暴力回放。} 慢版本枚举所有区间或候选对，再逐个检查条件。把检查式写出来后，可以看到当前状态只缺一个历史状态；这正是从平方枚举走向哈希表的入口。二刷时先写慢版本或口述慢版本，用它确认不会漏答案。

\textbf{特例回放。} 拿 eat、tea、tan 手算，eat 和 tea 排序后都是 aet，或计数向量都是 a1e1t1；tan 的键不同。两两比较会重复比较很多字符串，哈希分桶只比较规范键。规律是异位词共享规范表示，字符集固定时计数键更直接；空字符串和重复单词也必须落到同一个键语义里。把这个特例继续拆开看：每个合法答案都要还原成当前状态和某个历史状态的配对；哈希表的 key 负责定位历史状态，value 负责表达次数、最早位置或存在性。先查还是先写，必须由是否允许当前前缀和自己配对来决定。复盘时把这个特例压成状态字段：当前前缀状态、需要查询的历史 key、哈希表 value 语义、答案累计值；初始化的空前缀必须单独说明；再用边界 空字符串、重复单词、字符集不止小写字母、计数键必须无歧义 反查初始化、更新顺序和退出条件。

\textbf{状态回放。} 手算路线：手算时按下标列出当前状态、需要的历史 key、查到的 value、更新后的哈希表。这样能看出先查后插还是先插后查。状态字段：当前前缀状态、需要查询的历史 key、哈希表 value 语义、答案累计值；初始化的空前缀必须单独说明。状态升级：把回头扫描历史改成查表；key 是当前状态需要的历史状态，value 根据目标选择次数、最早位置、最近位置或存在性。

\textbf{证明和边界。} 证明从代数等价开始：任意合法答案都对应当前状态和某个历史 key，查到的历史 key 也能反推出合法答案。边界：空字符串、重复单词、字符集不止小写字母、计数键必须无歧义。变体：若字符串很长且字符集固定，计数键通常比排序键更稳定。

\noindent\textbf{LeetCode 128 最长连续序列。} 模型：连续段只需要从没有前驱的数开始扩展。推导重点：哈希集合判断 x 减一 是否存在，跳过非起点避免重复扫描。

\textbf{二刷读题。} 读题时先把条件改写成当前状态与历史状态的关系。本题可以压成“连续段只需要从没有前驱的数开始扩展”，所以重点是历史表的 key 和 value 分别代表什么。先遮住答案，只保留题面，复述候选对象、合法性条件和输出目标。

\textbf{暴力回放。} 慢版本枚举所有区间或候选对，再逐个检查条件。把检查式写出来后，可以看到当前状态只缺一个历史状态；这正是从平方枚举走向哈希表的入口。二刷时先写慢版本或口述慢版本，用它确认不会漏答案。

\textbf{特例回放。} 拿 [100,4,200,1,3,2] 手算，4 不是起点，因为 3 存在；1 没有前驱，从 1 扩到 4 得到长度 4。规律是只从没有 x-1 的数开始扩展，避免每个连续段被重复扫描；重复元素先由集合去掉。把这个特例继续拆开看：每个合法答案都要还原成当前状态和某个历史状态的配对；哈希表的 key 负责定位历史状态，value 负责表达次数、最早位置或存在性。先查还是先写，必须由是否允许当前前缀和自己配对来决定。复盘时把这个特例压成状态字段：当前前缀状态、需要查询的历史 key、哈希表 value 语义、答案累计值；初始化的空前缀必须单独说明；再用边界 重复元素、负数、空数组、整数边界 反查初始化、更新顺序和退出条件。

\textbf{状态回放。} 手算路线：手算时按下标列出当前状态、需要的历史 key、查到的 value、更新后的哈希表。这样能看出先查后插还是先插后查。状态字段：当前前缀状态、需要查询的历史 key、哈希表 value 语义、答案累计值；初始化的空前缀必须单独说明。状态升级：把回头扫描历史改成查表；key 是当前状态需要的历史状态，value 根据目标选择次数、最早位置、最近位置或存在性。

\textbf{证明和边界。} 证明从代数等价开始：任意合法答案都对应当前状态和某个历史 key，查到的历史 key 也能反推出合法答案。边界：重复元素、负数、空数组、整数边界。变体：若数据流动态插入，可以用并查集或区间合并维护长度。

\noindent\textbf{LeetCode 202 快乐数。} 模型：反复变换数字会进入一或循环，状态空间最终有限。推导重点：用集合检测重复状态，或用快慢指针看成函数图。

\textbf{二刷读题。} 读题时先把条件改写成当前状态与历史状态的关系。本题可以压成“反复变换数字会进入一或循环，状态空间最终有限”，所以重点是历史表的 key 和 value 分别代表什么。先遮住答案，只保留题面，复述候选对象、合法性条件和输出目标。

\textbf{暴力回放。} 慢版本枚举所有区间或候选对，再逐个检查条件。把检查式写出来后，可以看到当前状态只缺一个历史状态；这正是从平方枚举走向哈希表的入口。二刷时先写慢版本或口述慢版本，用它确认不会漏答案。

\textbf{特例回放。} 拿 19 手算，1 的平方加 9 的平方得到 82，82 再变成 68、100、1，所以快乐；拿 2 会进入循环而不到 1。规律是用集合记录出现过的中间值，或用快慢指针检测平方和序列是否成环。把这个特例继续拆开看：每个合法答案都要还原成当前状态和某个历史状态的配对；哈希表的 key 负责定位历史状态，value 负责表达次数、最早位置或存在性。先查还是先写，必须由是否允许当前前缀和自己配对来决定。复盘时把这个特例压成状态字段：当前前缀状态、需要查询的历史 key、哈希表 value 语义、答案累计值；初始化的空前缀必须单独说明；再用边界 输入为一、进入循环、位平方和函数、负数不在题意内 反查初始化、更新顺序和退出条件。

\textbf{状态回放。} 手算路线：手算时按下标列出当前状态、需要的历史 key、查到的 value、更新后的哈希表。这样能看出先查后插还是先插后查。状态字段：当前前缀状态、需要查询的历史 key、哈希表 value 语义、答案累计值；初始化的空前缀必须单独说明。状态升级：把回头扫描历史改成查表；key 是当前状态需要的历史状态，value 根据目标选择次数、最早位置、最近位置或存在性。

\textbf{证明和边界。} 证明从代数等价开始：任意合法答案都对应当前状态和某个历史 key，查到的历史 key 也能反推出合法答案。边界：输入为一、进入循环、位平方和函数、负数不在题意内。变体：快慢指针版本能训练把数值迭代看成链表。

\noindent\textbf{LeetCode 217 存在重复元素。} 模型：判断是否出现过是哈希集合最直接的职责。推导重点：扫描时若插入前已经存在则立即返回真。

\textbf{二刷读题。} 读题时先把条件改写成当前状态与历史状态的关系。本题可以压成“判断是否出现过是哈希集合最直接的职责”，所以重点是历史表的 key 和 value 分别代表什么。先遮住答案，只保留题面，复述候选对象、合法性条件和输出目标。

\textbf{暴力回放。} 慢版本枚举所有区间或候选对，再逐个检查条件。把检查式写出来后，可以看到当前状态只缺一个历史状态；这正是从平方枚举走向哈希表的入口。二刷时先写慢版本或口述慢版本，用它确认不会漏答案。

\textbf{特例回放。} 拿 [1,2,3,1] 手算，扫描到最后一个 1 时，哈希集合里已经有 1，立即返回重复。规律是集合保存已经出现过的历史值，插入前查询可以在第一次重复处停止；空数组和无重复数组返回假。把这个特例继续拆开看：每个合法答案都要还原成当前状态和某个历史状态的配对；哈希表的 key 负责定位历史状态，value 负责表达次数、最早位置或存在性。先查还是先写，必须由是否允许当前前缀和自己配对来决定。复盘时把这个特例压成状态字段：当前前缀状态、需要查询的历史 key、哈希表 value 语义、答案累计值；初始化的空前缀必须单独说明；再用边界 空数组、全唯一、重复在末尾、值域很小 反查初始化、更新顺序和退出条件。

\textbf{状态回放。} 手算路线：手算时按下标列出当前状态、需要的历史 key、查到的 value、更新后的哈希表。这样能看出先查后插还是先插后查。状态字段：当前前缀状态、需要查询的历史 key、哈希表 value 语义、答案累计值；初始化的空前缀必须单独说明。状态升级：把回头扫描历史改成查表；key 是当前状态需要的历史状态，value 根据目标选择次数、最早位置、最近位置或存在性。

\textbf{证明和边界。} 证明从代数等价开始：任意合法答案都对应当前状态和某个历史 key，查到的历史 key 也能反推出合法答案。边界：空数组、全唯一、重复在末尾、值域很小。变体：若允许修改输入，排序后看相邻可以降低额外空间。

\noindent\textbf{LeetCode 242 有效的字母异位词。} 模型：两个字符串是异位词等价于每个字符出现次数相同。推导重点：固定小写字母集时用定长计数数组，字符集未知时用哈希表。

\textbf{二刷读题。} 读题时先把条件改写成当前状态与历史状态的关系。本题可以压成“两个字符串是异位词等价于每个字符出现次数相同”，所以重点是历史表的 key 和 value 分别代表什么。先遮住答案，只保留题面，复述候选对象、合法性条件和输出目标。

\textbf{暴力回放。} 慢版本枚举所有区间或候选对，再逐个检查条件。把检查式写出来后，可以看到当前状态只缺一个历史状态；这正是从平方枚举走向哈希表的入口。二刷时先写慢版本或口述慢版本，用它确认不会漏答案。

\textbf{特例回放。} 拿 anagram 和 nagaram 手算，两个字符串每个字符计数都相同；拿 rat 和 car，r 和 c 的计数不同。规律是异位词比较的是字符频次，不是字符顺序；字符集固定可用数组，泛化字符集用哈希表。把这个特例继续拆开看：每个合法答案都要还原成当前状态和某个历史状态的配对；哈希表的 key 负责定位历史状态，value 负责表达次数、最早位置或存在性。先查还是先写，必须由是否允许当前前缀和自己配对来决定。复盘时把这个特例压成状态字段：当前前缀状态、需要查询的历史 key、哈希表 value 语义、答案累计值；初始化的空前缀必须单独说明；再用边界 长度不同、空串、Unicode 字符、计数减到负数 反查初始化、更新顺序和退出条件。

\textbf{状态回放。} 手算路线：手算时按下标列出当前状态、需要的历史 key、查到的 value、更新后的哈希表。这样能看出先查后插还是先插后查。状态字段：当前前缀状态、需要查询的历史 key、哈希表 value 语义、答案累计值；初始化的空前缀必须单独说明。状态升级：把回头扫描历史改成查表；key 是当前状态需要的历史状态，value 根据目标选择次数、最早位置、最近位置或存在性。

\textbf{证明和边界。} 证明从代数等价开始：任意合法答案都对应当前状态和某个历史 key，查到的历史 key 也能反推出合法答案。边界：长度不同、空串、Unicode 字符、计数减到负数。变体：异位词分组把这个二元比较扩展成规范 key 分桶。

\noindent\textbf{LeetCode 325 和等于 k 的最长子数组长度。} 模型：最长区间要求在同一前缀差关系下尽量拉大左右端距离。推导重点：哈希表保存每个前缀和最早出现下标，当前前缀查询 prefix 减 k。

\textbf{二刷读题。} 读题时先把条件改写成当前状态与历史状态的关系。本题可以压成“最长区间要求在同一前缀差关系下尽量拉大左右端距离”，所以重点是历史表的 key 和 value 分别代表什么。先遮住答案，只保留题面，复述候选对象、合法性条件和输出目标。

\textbf{暴力回放。} 慢版本枚举所有区间或候选对，再逐个检查条件。把检查式写出来后，可以看到当前状态只缺一个历史状态；这正是从平方枚举走向哈希表的入口。二刷时先写慢版本或口述慢版本，用它确认不会漏答案。

\textbf{特例回放。} 拿 [1,-1,5,-2,3]、k=3 手算，到前缀和 3 时需要历史前缀 0；到前缀和 4 时需要历史前缀 1，最早的 1 能给出更长区间。规律是最长区间要保存每个前缀和的最早下标，而不是最新下标。把这个特例继续拆开看：每个合法答案都要还原成当前状态和某个历史状态的配对；哈希表的 key 负责定位历史状态，value 负责表达次数、最早位置或存在性。先查还是先写，必须由是否允许当前前缀和自己配对来决定。复盘时把这个特例压成状态字段：当前前缀状态、需要查询的历史 key、哈希表 value 语义、答案累计值；初始化的空前缀必须单独说明；再用边界 从下标零开始、负数、前缀重复时不能覆盖最早位置、答案不存在 反查初始化、更新顺序和退出条件。

\textbf{状态回放。} 手算路线：手算时按下标列出当前状态、需要的历史 key、查到的 value、更新后的哈希表。这样能看出先查后插还是先插后查。状态字段：当前前缀状态、需要查询的历史 key、哈希表 value 语义、答案累计值；初始化的空前缀必须单独说明。状态升级：把回头扫描历史改成查表；key 是当前状态需要的历史状态，value 根据目标选择次数、最早位置、最近位置或存在性。

\textbf{证明和边界。} 证明从代数等价开始：任意合法答案都对应当前状态和某个历史 key，查到的历史 key 也能反推出合法答案。边界：从下标零开始、负数、前缀重复时不能覆盖最早位置、答案不存在。变体：计数版本保存频次，最长版本保存最早位置，这是核心差别。

\noindent\textbf{LeetCode 454 四数相加 II。} 模型：四个数组的和为零可以拆成两个二数组和互为相反数。推导重点：先统计 A 和 B 的所有两数和频次，再枚举 C 和 D 查询相反数。

\textbf{二刷读题。} 读题时先把条件改写成当前状态与历史状态的关系。本题可以压成“四个数组的和为零可以拆成两个二数组和互为相反数”，所以重点是历史表的 key 和 value 分别代表什么。先遮住答案，只保留题面，复述候选对象、合法性条件和输出目标。

\textbf{暴力回放。} 慢版本枚举所有区间或候选对，再逐个检查条件。把检查式写出来后，可以看到当前状态只缺一个历史状态；这正是从平方枚举走向哈希表的入口。二刷时先写慢版本或口述慢版本，用它确认不会漏答案。

\textbf{特例回放。} 拿 A=[1,2]、B=[-2,-1]、C=[-1,2]、D=[0,2] 手算，先统计 A+B 的和，后面枚举 C+D 时只查相反数。规律是四重枚举拆成两个二重枚举，用哈希表保存一半的和频次。把这个特例继续拆开看：每个合法答案都要还原成当前状态和某个历史状态的配对；哈希表的 key 负责定位历史状态，value 负责表达次数、最早位置或存在性。先查还是先写，必须由是否允许当前前缀和自己配对来决定。复盘时把这个特例压成状态字段：当前前缀状态、需要查询的历史 key、哈希表 value 语义、答案累计值；初始化的空前缀必须单独说明；再用边界 重复值、负数、答案数量很大、两数和范围 反查初始化、更新顺序和退出条件。

\textbf{状态回放。} 手算路线：手算时按下标列出当前状态、需要的历史 key、查到的 value、更新后的哈希表。这样能看出先查后插还是先插后查。状态字段：当前前缀状态、需要查询的历史 key、哈希表 value 语义、答案累计值；初始化的空前缀必须单独说明。状态升级：把回头扫描历史改成查表；key 是当前状态需要的历史状态，value 根据目标选择次数、最早位置、最近位置或存在性。

\textbf{证明和边界。} 证明从代数等价开始：任意合法答案都对应当前状态和某个历史 key，查到的历史 key 也能反推出合法答案。边界：重复值、负数、答案数量很大、两数和范围。变体：这是 meet-in-the-middle 思想，比四层枚举少两个维度。

\noindent\textbf{LeetCode 523 连续的子数组和。} 模型：长度至少二的子数组和为 k 的倍数，等价于两个前缀和同余。推导重点：哈希表保存某个余数最早出现下标，当前余数若重复且距离至少二则成立。

\textbf{二刷读题。} 读题时先把条件改写成当前状态与历史状态的关系。本题可以压成“长度至少二的子数组和为 k 的倍数，等价于两个前缀和同余”，所以重点是历史表的 key 和 value 分别代表什么。先遮住答案，只保留题面，复述候选对象、合法性条件和输出目标。

\textbf{暴力回放。} 慢版本枚举所有区间或候选对，再逐个检查条件。把检查式写出来后，可以看到当前状态只缺一个历史状态；这正是从平方枚举走向哈希表的入口。二刷时先写慢版本或口述慢版本，用它确认不会漏答案。

\textbf{特例回放。} 拿 [23,2,4,6,7]、k=6 手算，前缀余数 5 在下标 0 和 2 重复，中间 [2,4] 的和能被 6 整除。规律是两个前缀余数相同表示区间和为 k 的倍数，且区间长度至少为 2。把这个特例继续拆开看：每个合法答案都要还原成当前状态和某个历史状态的配对；哈希表的 key 负责定位历史状态，value 负责表达次数、最早位置或存在性。先查还是先写，必须由是否允许当前前缀和自己配对来决定。复盘时把这个特例压成状态字段：当前前缀状态、需要查询的历史 key、哈希表 value 语义、答案累计值；初始化的空前缀必须单独说明；再用边界 k 为零、负数余数、长度限制、从开头开始的区间 反查初始化、更新顺序和退出条件。

\textbf{状态回放。} 手算路线：手算时按下标列出当前状态、需要的历史 key、查到的 value、更新后的哈希表。这样能看出先查后插还是先插后查。状态字段：当前前缀状态、需要查询的历史 key、哈希表 value 语义、答案累计值；初始化的空前缀必须单独说明。状态升级：把回头扫描历史改成查表；key 是当前状态需要的历史状态，value 根据目标选择次数、最早位置、最近位置或存在性。

\textbf{证明和边界。} 证明从代数等价开始：任意合法答案都对应当前状态和某个历史 key，查到的历史 key 也能反推出合法答案。边界：k 为零、负数余数、长度限制、从开头开始的区间。变体：保存最早下标是为了满足长度条件并保留最长距离。

\noindent\textbf{LeetCode 525 连续数组。} 模型：把零看成负一后，零和一数量相等等价于前缀和相同。推导重点：哈希表保存每个前缀和最早下标，重复出现时更新最长长度。

\textbf{二刷读题。} 读题时先把条件改写成当前状态与历史状态的关系。本题可以压成“把零看成负一后，零和一数量相等等价于前缀和相同”，所以重点是历史表的 key 和 value 分别代表什么。先遮住答案，只保留题面，复述候选对象、合法性条件和输出目标。

\textbf{暴力回放。} 慢版本枚举所有区间或候选对，再逐个检查条件。把检查式写出来后，可以看到当前状态只缺一个历史状态；这正是从平方枚举走向哈希表的入口。二刷时先写慢版本或口述慢版本，用它确认不会漏答案。

\textbf{特例回放。} 拿 [0,1,0] 手算，把 0 当 -1，前缀差值序列 -1、0、-1；差值再次出现说明中间 0 和 1 数量相同。规律是保存差值第一次出现的位置，最长区间用最早位置。把这个特例继续拆开看：每个合法答案都要还原成当前状态和某个历史状态的配对；哈希表的 key 负责定位历史状态，value 负责表达次数、最早位置或存在性。先查还是先写，必须由是否允许当前前缀和自己配对来决定。复盘时把这个特例压成状态字段：当前前缀状态、需要查询的历史 key、哈希表 value 语义、答案累计值；初始化的空前缀必须单独说明；再用边界 全零、全一、从下标零开始、不要覆盖最早位置 反查初始化、更新顺序和退出条件。

\textbf{状态回放。} 手算路线：手算时按下标列出当前状态、需要的历史 key、查到的 value、更新后的哈希表。这样能看出先查后插还是先插后查。状态字段：当前前缀状态、需要查询的历史 key、哈希表 value 语义、答案累计值；初始化的空前缀必须单独说明。状态升级：把回头扫描历史改成查表；key 是当前状态需要的历史状态，value 根据目标选择次数、最早位置、最近位置或存在性。

\textbf{证明和边界。} 证明从代数等价开始：任意合法答案都对应当前状态和某个历史 key，查到的历史 key 也能反推出合法答案。边界：全零、全一、从下标零开始、不要覆盖最早位置。变体：这是前缀和最早位置语义，不是频次语义。

\noindent\textbf{LeetCode 560 和为 K 的子数组。} 模型：当前前缀和减去目标值，就是需要匹配的历史前缀。推导重点：哈希表保存前缀和出现次数，先查询再更新。

\textbf{二刷读题。} 读题时先把条件改写成当前状态与历史状态的关系。本题可以压成“当前前缀和减去目标值，就是需要匹配的历史前缀”，所以重点是历史表的 key 和 value 分别代表什么。先遮住答案，只保留题面，复述候选对象、合法性条件和输出目标。

\textbf{暴力回放。} 慢版本枚举所有区间或候选对，再逐个检查条件。把检查式写出来后，可以看到当前状态只缺一个历史状态；这正是从平方枚举走向哈希表的入口。二刷时先写慢版本或口述慢版本，用它确认不会漏答案。

\textbf{特例回放。} 拿 [1,2,3]、k=3 手算，到前缀和 3 时，需要历史前缀 0；到前缀和 6 时，需要历史前缀 3。每个历史前缀出现几次，就新增几个区间。规律是哈希表保存前缀和频次，且初始要放 0:1。把这个特例继续拆开看：每个合法答案都要还原成当前状态和某个历史状态的配对；哈希表的 key 负责定位历史状态，value 负责表达次数、最早位置或存在性。先查还是先写，必须由是否允许当前前缀和自己配对来决定。复盘时把这个特例压成状态字段：当前前缀状态、需要查询的历史 key、哈希表 value 语义、答案累计值；初始化的空前缀必须单独说明；再用边界 负数、目标为零、从下标零开始的区间、重复前缀 反查初始化、更新顺序和退出条件。

\textbf{状态回放。} 手算路线：手算时按下标列出当前状态、需要的历史 key、查到的 value、更新后的哈希表。这样能看出先查后插还是先插后查。状态字段：当前前缀状态、需要查询的历史 key、哈希表 value 语义、答案累计值；初始化的空前缀必须单独说明。状态升级：把回头扫描历史改成查表；key 是当前状态需要的历史状态，value 根据目标选择次数、最早位置、最近位置或存在性。

\textbf{证明和边界。} 证明从代数等价开始：任意合法答案都对应当前状态和某个历史 key，查到的历史 key 也能反推出合法答案。边界：负数、目标为零、从下标零开始的区间、重复前缀。变体：若要求最长区间，哈希表应保存最早位置而不是次数。

\noindent\textbf{LeetCode 974 和可被 K 整除的子数组。} 模型：区间和能被 K 整除等价于两个前缀和余数相同。推导重点：统计每个余数出现次数，当前余数能和所有历史同余前缀组成答案。

\textbf{二刷读题。} 读题时先把条件改写成当前状态与历史状态的关系。本题可以压成“区间和能被 K 整除等价于两个前缀和余数相同”，所以重点是历史表的 key 和 value 分别代表什么。先遮住答案，只保留题面，复述候选对象、合法性条件和输出目标。

\textbf{暴力回放。} 慢版本枚举所有区间或候选对，再逐个检查条件。把检查式写出来后，可以看到当前状态只缺一个历史状态；这正是从平方枚举走向哈希表的入口。二刷时先写慢版本或口述慢版本，用它确认不会漏答案。

\textbf{特例回放。} 拿 [4,5,0,-2,-3,1]、k=5 手算，前缀余数重复时，中间区间和能被 5 整除；负数前缀要先把余数归一化。规律是哈希表保存余数频次，每到一个余数就增加已有次数。把这个特例继续拆开看：每个合法答案都要还原成当前状态和某个历史状态的配对；哈希表的 key 负责定位历史状态，value 负责表达次数、最早位置或存在性。先查还是先写，必须由是否允许当前前缀和自己配对来决定。复盘时把这个特例压成状态字段：当前前缀状态、需要查询的历史 key、哈希表 value 语义、答案累计值；初始化的空前缀必须单独说明；再用边界 负数余数归一化、K 为一、从开头开始、重复余数很多 反查初始化、更新顺序和退出条件。

\textbf{状态回放。} 手算路线：手算时按下标列出当前状态、需要的历史 key、查到的 value、更新后的哈希表。这样能看出先查后插还是先插后查。状态字段：当前前缀状态、需要查询的历史 key、哈希表 value 语义、答案累计值；初始化的空前缀必须单独说明。状态升级：把回头扫描历史改成查表；key 是当前状态需要的历史状态，value 根据目标选择次数、最早位置、最近位置或存在性。

\textbf{证明和边界。} 证明从代数等价开始：任意合法答案都对应当前状态和某个历史 key，查到的历史 key 也能反推出合法答案。边界：负数余数归一化、K 为一、从开头开始、重复余数很多。变体：如果问最长长度，保存最早位置而不是次数。

\noindent\textbf{LeetCode 1074 元素和为目标值的子矩阵数量。} 模型：二维矩阵子区域和可以固定上下边界后转成一维子数组和。推导重点：枚举行边界，压缩每列区间和，再用前缀和哈希统计目标和子数组。

\textbf{二刷读题。} 读题时先把条件改写成当前状态与历史状态的关系。本题可以压成“二维矩阵子区域和可以固定上下边界后转成一维子数组和”，所以重点是历史表的 key 和 value 分别代表什么。先遮住答案，只保留题面，复述候选对象、合法性条件和输出目标。

\textbf{暴力回放。} 慢版本枚举所有区间或候选对，再逐个检查条件。把检查式写出来后，可以看到当前状态只缺一个历史状态；这正是从平方枚举走向哈希表的入口。二刷时先写慢版本或口述慢版本，用它确认不会漏答案。

\textbf{特例回放。} 拿 3x3 矩阵和 target=0 手算，先固定上下两行，把每列压成列和，问题变成一维子数组和为 0 的计数。规律是枚举行边界，列方向用前缀和加哈希计数，二维问题降成多次一维问题。把这个特例继续拆开看：每个合法答案都要还原成当前状态和某个历史状态的配对；哈希表的 key 负责定位历史状态，value 负责表达次数、最早位置或存在性。先查还是先写，必须由是否允许当前前缀和自己配对来决定。复盘时把这个特例压成状态字段：当前前缀状态、需要查询的历史 key、哈希表 value 语义、答案累计值；初始化的空前缀必须单独说明；再用边界 单行、单列、负数、目标为零、矩阵较大 反查初始化、更新顺序和退出条件。

\textbf{状态回放。} 手算路线：手算时按下标列出当前状态、需要的历史 key、查到的 value、更新后的哈希表。这样能看出先查后插还是先插后查。状态字段：当前前缀状态、需要查询的历史 key、哈希表 value 语义、答案累计值；初始化的空前缀必须单独说明。状态升级：把回头扫描历史改成查表；key 是当前状态需要的历史状态，value 根据目标选择次数、最早位置、最近位置或存在性。

\textbf{证明和边界。} 证明从代数等价开始：任意合法答案都对应当前状态和某个历史 key，查到的历史 key 也能反推出合法答案。边界：单行、单列、负数、目标为零、矩阵较大。变体：这是二维前缀思想和一维哈希计数的组合。

\noindent\textbf{LeetCode 1248 统计优美子数组。} 模型：恰好 K 个奇数可转成前缀奇数计数差为 K。推导重点：哈希表保存某个奇数计数出现次数，当前计数查询 count - K。

\textbf{二刷读题。} 读题时先把条件改写成当前状态与历史状态的关系。本题可以压成“恰好 K 个奇数可转成前缀奇数计数差为 K”，所以重点是历史表的 key 和 value 分别代表什么。先遮住答案，只保留题面，复述候选对象、合法性条件和输出目标。

\textbf{暴力回放。} 慢版本枚举所有区间或候选对，再逐个检查条件。把检查式写出来后，可以看到当前状态只缺一个历史状态；这正是从平方枚举走向哈希表的入口。二刷时先写慢版本或口述慢版本，用它确认不会漏答案。

\textbf{特例回放。} 拿 [1,1,2,1,1]、k=3 手算，奇数个数前缀从 0 到 3 时，需要历史奇数个数 0；每个匹配历史前缀对应一个优美子数组。规律是把奇数看成 1、偶数看成 0，转成前缀和计数。把这个特例继续拆开看：每个合法答案都要还原成当前状态和某个历史状态的配对；哈希表的 key 负责定位历史状态，value 负责表达次数、最早位置或存在性。先查还是先写，必须由是否允许当前前缀和自己配对来决定。复盘时把这个特例压成状态字段：当前前缀状态、需要查询的历史 key、哈希表 value 语义、答案累计值；初始化的空前缀必须单独说明；再用边界 K 为零、全偶数、全奇数、从开头开始的区间 反查初始化、更新顺序和退出条件。

\textbf{状态回放。} 手算路线：手算时按下标列出当前状态、需要的历史 key、查到的 value、更新后的哈希表。这样能看出先查后插还是先插后查。状态字段：当前前缀状态、需要查询的历史 key、哈希表 value 语义、答案累计值；初始化的空前缀必须单独说明。状态升级：把回头扫描历史改成查表；key 是当前状态需要的历史状态，value 根据目标选择次数、最早位置、最近位置或存在性。

\textbf{证明和边界。} 证明从代数等价开始：任意合法答案都对应当前状态和某个历史 key，查到的历史 key 也能反推出合法答案。边界：K 为零、全偶数、全奇数、从开头开始的区间。变体：也可用至多 K 减至多 K-1 的滑动窗口。

\subsection{自测标准}

二刷时不要只确认自己见过这道题。请遮住答案，先讲题面模型，再讲暴力基线和瓶颈，然后说出优化结构保存了什么、不变量如何排除错误候选、边界实验如何打破错误实现。

\section{二分查找与答案空间：二刷复盘卡}
这一大节只围绕一个题型复盘，不把题号直接铺成目录。阅读时先看复盘目标，再读核心题卡，最后用自测标准检查自己是否能独立重讲。

\subsection{复盘目标}

追问通常会要求你讲清答案边界、可行性谓词或局部排除规则。请准备说明 left、right 的语义，为什么 mid 被排除或保留，以及返回值是否还需要二次验证。

\subsection{核心题卡}

\noindent\textbf{LeetCode 4 寻找两个正序数组的中位数。} 模型：两个有序数组的中位数可以看成左右两半切分后的边界值。推导重点：二分较短数组的切分位置，让左半元素数量正确且左半最大值不大于右半最小值。

\textbf{二刷读题。} 读题时先找候选空间和目标边界。本题可以压成“两个有序数组的中位数可以看成左右两半切分后的边界值”。然后把关键观察写成可执行规则：二分较短数组的切分位置，让左半元素数量正确且左半最大值不大于右半最小值。这条规则必须能安全排除一侧候选，或收缩到可检查的答案边界。先遮住答案，只保留题面，复述候选对象、合法性条件和输出目标。

\textbf{暴力回放。} 慢版本按顺序扫描下标、逐个试答案值，或直接检查候选直到遇到目标边界。它先保证候选空间覆盖完整，但每次只排除一个候选，浪费了题目给出的顺序、比较或可行性边界。二刷时先写慢版本或口述慢版本，用它确认不会漏答案。

\textbf{特例回放。} 拿 [1, 3] 和 [2] 手算，合并后中位数是 2；但不用真合并，只要把总长度切成左半两个数、右半一个数。再拿 [1, 2] 和 [3, 4]，如果短数组切 1 个、长数组切 1 个，左半最大 3 大于右半最小 2，切分错了，应把短数组切分往右调。规律是二分短数组切分点，使左半元素个数正确且左右边界有序。把这个特例继续拆开看：每次比较 mid 之后，必须能指出哪一侧整体不可能包含目标，或哪一侧整体已经满足可行性。二分的核心不是 mid 公式，而是被丢弃区间确实不含答案。复盘时把这个特例压成状态字段：left、right、mid、mid 处比较值或检查返回值、答案区间含义、返回值语义；每一轮更新都要保留目标位置或第一个可行值；再用边界 某个数组为空、总长度奇偶、切分在边界、哨兵值不能参与真实比较 反查初始化、更新顺序和退出条件。

\textbf{状态回放。} 手算路线：手算时列出 left、mid、right，以及 mid 处比较结果、可行性检查结果或局部趋势。每一步都要写出被丢弃的一侧为什么不含答案。状态字段：left、right、mid、mid 处比较值或检查返回值、答案区间含义、返回值语义；每一轮更新都要保留目标位置或第一个可行值。状态升级：把逐个试候选改成维护答案所在区间；边界谓词、可行性检查或局部比较给出分支方向，每个分支都必须能排除一侧候选。

\textbf{证明和边界。} 证明从排除一侧开始：答案空间题证明谓词单调，局部比较题证明保留下来的区间仍含目标；不能只靠经验猜测左右边界。边界：某个数组为空、总长度奇偶、切分在边界、哨兵值不能参与真实比较。变体：若只要求第 k 小元素，可以用递归丢弃较小前缀或二分切分的同类思想。

\noindent\textbf{LeetCode 33 搜索旋转排序数组。} 模型：旋转数组每次二分至少有一半保持有序。推导重点：判断哪一半有序，再看目标是否落在有序半边范围内。

\textbf{二刷读题。} 读题时先找候选空间和目标边界。本题可以压成“旋转数组每次二分至少有一半保持有序”。然后把关键观察写成可执行规则：判断哪一半有序，再看目标是否落在有序半边范围内。这条规则必须能安全排除一侧候选，或收缩到可检查的答案边界。先遮住答案，只保留题面，复述候选对象、合法性条件和输出目标。

\textbf{暴力回放。} 慢版本按顺序扫描下标、逐个试答案值，或直接检查候选直到遇到目标边界。它先保证候选空间覆盖完整，但每次只排除一个候选，浪费了题目给出的顺序、比较或可行性边界。二刷时先写慢版本或口述慢版本，用它确认不会漏答案。

\textbf{特例回放。} 拿 [4,5,6,7,0,1,2] 找 0 手算，mid=7 时左半 [4,5,6,7] 有序，但目标 0 不在这个范围，所以整段左半都能丢掉。下一轮目标落在右半。规律是每次先判断哪一半有序，再用目标是否落入有序半边来排除一整段；有重复元素时有序半边可能判断不出来。把这个特例继续拆开看：每次比较 mid 之后，必须能指出哪一侧整体不可能包含目标，或哪一侧整体已经满足可行性。二分的核心不是 mid 公式，而是被丢弃区间确实不含答案。复盘时把这个特例压成状态字段：left、right、mid、mid 处比较值或检查返回值、答案区间含义、返回值语义；每一轮更新都要保留目标位置或第一个可行值；再用边界 未旋转、长度一、目标不存在、边界等号、存在重复元素时退化 反查初始化、更新顺序和退出条件。

\textbf{状态回放。} 手算路线：手算时列出 left、mid、right，以及 mid 处比较结果、可行性检查结果或局部趋势。每一步都要写出被丢弃的一侧为什么不含答案。状态字段：left、right、mid、mid 处比较值或检查返回值、答案区间含义、返回值语义；每一轮更新都要保留目标位置或第一个可行值。状态升级：把逐个试候选改成维护答案所在区间；边界谓词、可行性检查或局部比较给出分支方向，每个分支都必须能排除一侧候选。

\textbf{证明和边界。} 证明从排除一侧开始：答案空间题证明谓词单调，局部比较题证明保留下来的区间仍含目标；不能只靠经验猜测左右边界。边界：未旋转、长度一、目标不存在、边界等号、存在重复元素时退化。变体：若重复元素很多，需要在无法判断有序半边时收缩边界。

\noindent\textbf{LeetCode 34 排序数组中查找首尾位置。} 模型：首尾位置可以拆成两个边界：第一个大于等于和第一个大于。推导重点：不要找到一个目标后向两边线性扩展，否则重复元素很多时退化。

\textbf{二刷读题。} 读题时先找候选空间和目标边界。本题可以压成“首尾位置可以拆成两个边界：第一个大于等于和第一个大于”。然后把关键观察写成可执行规则：不要找到一个目标后向两边线性扩展，否则重复元素很多时退化。这条规则必须能安全排除一侧候选，或收缩到可检查的答案边界。先遮住答案，只保留题面，复述候选对象、合法性条件和输出目标。

\textbf{暴力回放。} 慢版本按顺序扫描下标、逐个试答案值，或直接检查候选直到遇到目标边界。它先保证候选空间覆盖完整，但每次只排除一个候选，浪费了题目给出的顺序、比较或可行性边界。二刷时先写慢版本或口述慢版本，用它确认不会漏答案。

\textbf{特例回放。} 拿 [5,7,7,8,8,10] 找 8 手算，找到一个 8 后向两边扩展能得到答案，但全数组都是 8 时会退化成线性。改成两次边界二分：第一个 >=8 的位置是 3，第一个 >8 的位置是 5，区间就是 [3,4]。规律是首尾位置不是一次查找，而是两个单调边界。把这个特例继续拆开看：每次比较 mid 之后，必须能指出哪一侧整体不可能包含目标，或哪一侧整体已经满足可行性。二分的核心不是 mid 公式，而是被丢弃区间确实不含答案。复盘时把这个特例压成状态字段：left、right、mid、mid 处比较值或检查返回值、答案区间含义、返回值语义；每一轮更新都要保留目标位置或第一个可行值；再用边界 目标不存在、目标在开头或结尾、数组为空、全是目标 反查初始化、更新顺序和退出条件。

\textbf{状态回放。} 手算路线：手算时列出 left、mid、right，以及 mid 处比较结果、可行性检查结果或局部趋势。每一步都要写出被丢弃的一侧为什么不含答案。状态字段：left、right、mid、mid 处比较值或检查返回值、答案区间含义、返回值语义；每一轮更新都要保留目标位置或第一个可行值。状态升级：把逐个试候选改成维护答案所在区间；边界谓词、可行性检查或局部比较给出分支方向，每个分支都必须能排除一侧候选。

\textbf{证明和边界。} 证明从排除一侧开始：答案空间题证明谓词单调，局部比较题证明保留下来的区间仍含目标；不能只靠经验猜测左右边界。边界：目标不存在、目标在开头或结尾、数组为空、全是目标。变体：这个模型能迁移到计数某值出现次数和插入区间边界。

\noindent\textbf{LeetCode 35 搜索插入位置。} 模型：题目等价于寻找第一个大于等于目标的位置。推导重点：统一用 lower bound 语义，不在相等时提前返回也能保持边界一致。

\textbf{二刷读题。} 读题时先找候选空间和目标边界。本题可以压成“题目等价于寻找第一个大于等于目标的位置”。然后把关键观察写成可执行规则：统一用 lower bound 语义，不在相等时提前返回也能保持边界一致。这条规则必须能安全排除一侧候选，或收缩到可检查的答案边界。先遮住答案，只保留题面，复述候选对象、合法性条件和输出目标。

\textbf{暴力回放。} 慢版本按顺序扫描下标、逐个试答案值，或直接检查候选直到遇到目标边界。它先保证候选空间覆盖完整，但每次只排除一个候选，浪费了题目给出的顺序、比较或可行性边界。二刷时先写慢版本或口述慢版本，用它确认不会漏答案。

\textbf{特例回放。} 拿 [1,3,5,6] 插入 2 手算，线性看 1 还小、3 已经不小，所以答案是下标 1。二分时保持 left 是第一个可能不小于目标的位置，mid 值小于目标才丢左侧。规律是统一成 lower\_bound，不需要在相等时提前返回；目标大于所有元素时 left 会走到 n。把这个特例继续拆开看：每次比较 mid 之后，必须能指出哪一侧整体不可能包含目标，或哪一侧整体已经满足可行性。二分的核心不是 mid 公式，而是被丢弃区间确实不含答案。复盘时把这个特例压成状态字段：left、right、mid、mid 处比较值或检查返回值、答案区间含义、返回值语义；每一轮更新都要保留目标位置或第一个可行值；再用边界 目标小于所有元素、目标大于所有元素、重复值、空数组 反查初始化、更新顺序和退出条件。

\textbf{状态回放。} 手算路线：手算时列出 left、mid、right，以及 mid 处比较结果、可行性检查结果或局部趋势。每一步都要写出被丢弃的一侧为什么不含答案。状态字段：left、right、mid、mid 处比较值或检查返回值、答案区间含义、返回值语义；每一轮更新都要保留目标位置或第一个可行值。状态升级：把逐个试候选改成维护答案所在区间；边界谓词、可行性检查或局部比较给出分支方向，每个分支都必须能排除一侧候选。

\textbf{证明和边界。} 证明从排除一侧开始：答案空间题证明谓词单调，局部比较题证明保留下来的区间仍含目标；不能只靠经验猜测左右边界。边界：目标小于所有元素、目标大于所有元素、重复值、空数组。变体：手写二分时用左闭右开可直接返回 left。

\noindent\textbf{LeetCode 69 x 的平方根。} 模型：答案是最大的整数 r，使得 r 的平方不超过 x。推导重点：在整数范围上二分右边界，比较时用除法或 long long 避免平方溢出。

\textbf{二刷读题。} 读题时先找候选空间和目标边界。本题可以压成“答案是最大的整数 r，使得 r 的平方不超过 x”。然后把关键观察写成可执行规则：在整数范围上二分右边界，比较时用除法或 long long 避免平方溢出。这条规则必须能安全排除一侧候选，或收缩到可检查的答案边界。先遮住答案，只保留题面，复述候选对象、合法性条件和输出目标。

\textbf{暴力回放。} 慢版本按顺序扫描下标、逐个试答案值，或直接检查候选直到遇到目标边界。它先保证候选空间覆盖完整，但每次只排除一个候选，浪费了题目给出的顺序、比较或可行性边界。二刷时先写慢版本或口述慢版本，用它确认不会漏答案。

\textbf{特例回放。} 拿 x=8 手算，二的平方不超过 8，而三的平方已经超过 8，所以答案是 2。规律是找最后一个平方不超过 x 的整数，比较平方时用 long long 或除法避免溢出。把这个特例继续拆开看：每次比较 mid 之后，必须能指出哪一侧整体不可能包含目标，或哪一侧整体已经满足可行性。二分的核心不是 mid 公式，而是被丢弃区间确实不含答案。复盘时把这个特例压成状态字段：left、right、mid、mid 处比较值或检查返回值、答案区间含义、返回值语义；每一轮更新都要保留目标位置或第一个可行值；再用边界 x 为零、一、非完全平方、接近 int 最大值 反查初始化、更新顺序和退出条件。

\textbf{状态回放。} 手算路线：手算时列出 left、mid、right，以及 mid 处比较结果、可行性检查结果或局部趋势。每一步都要写出被丢弃的一侧为什么不含答案。状态字段：left、right、mid、mid 处比较值或检查返回值、答案区间含义、返回值语义；每一轮更新都要保留目标位置或第一个可行值。状态升级：把逐个试候选改成维护答案所在区间；边界谓词、可行性检查或局部比较给出分支方向，每个分支都必须能排除一侧候选。

\textbf{证明和边界。} 证明从排除一侧开始：答案空间题证明谓词单调，局部比较题证明保留下来的区间仍含目标；不能只靠经验猜测左右边界。边界：x 为零、一、非完全平方、接近 int 最大值。变体：若要求浮点精度，可以二分实数区间或用牛顿迭代。

\noindent\textbf{LeetCode 74 搜索二维矩阵。} 模型：矩阵每行有序且下一行首元素大于上一行末元素，可视为一维有序数组。推导重点：把下标 mid 映射到 row 和 col，再做标准二分。

\textbf{二刷读题。} 读题时先找候选空间和目标边界。本题可以压成“矩阵每行有序且下一行首元素大于上一行末元素，可视为一维有序数组”。然后把关键观察写成可执行规则：把下标 mid 映射到 row 和 col，再做标准二分。这条规则必须能安全排除一侧候选，或收缩到可检查的答案边界。先遮住答案，只保留题面，复述候选对象、合法性条件和输出目标。

\textbf{暴力回放。} 慢版本按顺序扫描下标、逐个试答案值，或直接检查候选直到遇到目标边界。它先保证候选空间覆盖完整，但每次只排除一个候选，浪费了题目给出的顺序、比较或可行性边界。二刷时先写慢版本或口述慢版本，用它确认不会漏答案。

\textbf{特例回放。} 拿矩阵 [[1,3,5],[7,9,11]] 找 9 手算，把它看成一维有序数组 [1,3,5,7,9,11]，下标 4 映射到 row=1、col=1。规律是行尾小于下一行行首时可以整体展平二分，空矩阵和单行要先处理。把这个特例继续拆开看：每次比较 mid 之后，必须能指出哪一侧整体不可能包含目标，或哪一侧整体已经满足可行性。二分的核心不是 mid 公式，而是被丢弃区间确实不含答案。复盘时把这个特例压成状态字段：left、right、mid、mid 处比较值或检查返回值、答案区间含义、返回值语义；每一轮更新都要保留目标位置或第一个可行值；再用边界 空矩阵、单行、单列、目标小于全部或大于全部 反查初始化、更新顺序和退出条件。

\textbf{状态回放。} 手算路线：手算时列出 left、mid、right，以及 mid 处比较结果、可行性检查结果或局部趋势。每一步都要写出被丢弃的一侧为什么不含答案。状态字段：left、right、mid、mid 处比较值或检查返回值、答案区间含义、返回值语义；每一轮更新都要保留目标位置或第一个可行值。状态升级：把逐个试候选改成维护答案所在区间；边界谓词、可行性检查或局部比较给出分支方向，每个分支都必须能排除一侧候选。

\textbf{证明和边界。} 证明从排除一侧开始：答案空间题证明谓词单调，局部比较题证明保留下来的区间仍含目标；不能只靠经验猜测左右边界。边界：空矩阵、单行、单列、目标小于全部或大于全部。变体：若只是行列分别有序但行之间不连续，应改用右上角排除法。

\noindent\textbf{LeetCode 81 搜索旋转排序数组 II。} 模型：旋转数组存在重复值时，二分可能无法判断哪一半有序。推导重点：当左右和中点值相同导致信息不足时，收缩边界；否则仍按有序半边排除。

\textbf{二刷读题。} 读题时先找候选空间和目标边界。本题可以压成“旋转数组存在重复值时，二分可能无法判断哪一半有序”。然后把关键观察写成可执行规则：当左右和中点值相同导致信息不足时，收缩边界；否则仍按有序半边排除。这条规则必须能安全排除一侧候选，或收缩到可检查的答案边界。先遮住答案，只保留题面，复述候选对象、合法性条件和输出目标。

\textbf{暴力回放。} 慢版本按顺序扫描下标、逐个试答案值，或直接检查候选直到遇到目标边界。它先保证候选空间覆盖完整，但每次只排除一个候选，浪费了题目给出的顺序、比较或可行性边界。二刷时先写慢版本或口述慢版本，用它确认不会漏答案。

\textbf{特例回放。} 拿 [2,5,6,0,0,1,2] 找 0 手算，可以判断一侧有序并排除；拿 [1,1,1,1] 时 mid、left、right 相同，无法判断，只能收缩边界。规律是重复元素让旋转二分可能退化，但每次丢掉一个相同边界不会漏目标。把这个特例继续拆开看：每次比较 mid 之后，必须能指出哪一侧整体不可能包含目标，或哪一侧整体已经满足可行性。二分的核心不是 mid 公式，而是被丢弃区间确实不含答案。复盘时把这个特例压成状态字段：left、right、mid、mid 处比较值或检查返回值、答案区间含义、返回值语义；每一轮更新都要保留目标位置或第一个可行值；再用边界 全相同元素、目标不存在、重复值跨过旋转点、退化到线性 反查初始化、更新顺序和退出条件。

\textbf{状态回放。} 手算路线：手算时列出 left、mid、right，以及 mid 处比较结果、可行性检查结果或局部趋势。每一步都要写出被丢弃的一侧为什么不含答案。状态字段：left、right、mid、mid 处比较值或检查返回值、答案区间含义、返回值语义；每一轮更新都要保留目标位置或第一个可行值。状态升级：把逐个试候选改成维护答案所在区间；边界谓词、可行性检查或局部比较给出分支方向，每个分支都必须能排除一侧候选。

\textbf{证明和边界。} 证明从排除一侧开始：答案空间题证明谓词单调，局部比较题证明保留下来的区间仍含目标；不能只靠经验猜测左右边界。边界：全相同元素、目标不存在、重复值跨过旋转点、退化到线性。变体：这题说明二分依赖信息量，重复值会削弱每次排除一半的能力。

\noindent\textbf{LeetCode 153 寻找旋转排序数组中的最小值。} 模型：最小值是旋转断点，右端点可帮助判断中点在哪一段。推导重点：若 mid 大于 right，最小值在右侧；否则在左侧含 mid。

\textbf{二刷读题。} 读题时先找候选空间和目标边界。本题可以压成“最小值是旋转断点，右端点可帮助判断中点在哪一段”。然后把关键观察写成可执行规则：若 mid 大于 right，最小值在右侧；否则在左侧含 mid。这条规则必须能安全排除一侧候选，或收缩到可检查的答案边界。先遮住答案，只保留题面，复述候选对象、合法性条件和输出目标。

\textbf{暴力回放。} 慢版本按顺序扫描下标、逐个试答案值，或直接检查候选直到遇到目标边界。它先保证候选空间覆盖完整，但每次只排除一个候选，浪费了题目给出的顺序、比较或可行性边界。二刷时先写慢版本或口述慢版本，用它确认不会漏答案。

\textbf{特例回放。} 拿 [3,4,5,1,2] 手算，mid=5 大于右端 2，说明最小值一定在 mid 右侧；若 mid 小于右端，最小值在左半含 mid。规律是和右端比较能判断旋转点所在侧，未旋转数组会一路收缩到第一个元素。把这个特例继续拆开看：每次比较 mid 之后，必须能指出哪一侧整体不可能包含目标，或哪一侧整体已经满足可行性。二分的核心不是 mid 公式，而是被丢弃区间确实不含答案。复盘时把这个特例压成状态字段：left、right、mid、mid 处比较值或检查返回值、答案区间含义、返回值语义；每一轮更新都要保留目标位置或第一个可行值；再用边界 未旋转、长度一、旋转一位、无重复条件 反查初始化、更新顺序和退出条件。

\textbf{状态回放。} 手算路线：手算时列出 left、mid、right，以及 mid 处比较结果、可行性检查结果或局部趋势。每一步都要写出被丢弃的一侧为什么不含答案。状态字段：left、right、mid、mid 处比较值或检查返回值、答案区间含义、返回值语义；每一轮更新都要保留目标位置或第一个可行值。状态升级：把逐个试候选改成维护答案所在区间；边界谓词、可行性检查或局部比较给出分支方向，每个分支都必须能排除一侧候选。

\textbf{证明和边界。} 证明从排除一侧开始：答案空间题证明谓词单调，局部比较题证明保留下来的区间仍含目标；不能只靠经验猜测左右边界。边界：未旋转、长度一、旋转一位、无重复条件。变体：有重复时遇到相等只能收缩 right，最坏退化。

\noindent\textbf{LeetCode 154 寻找旋转排序数组中的最小值 II。} 模型：重复元素会让 mid 和 right 相等时无法判断最小值方向。推导重点：相等时只能安全地丢掉一个 right，因为它不是唯一必要候选。

\textbf{二刷读题。} 读题时先找候选空间和目标边界。本题可以压成“重复元素会让 mid 和 right 相等时无法判断最小值方向”。然后把关键观察写成可执行规则：相等时只能安全地丢掉一个 right，因为它不是唯一必要候选。这条规则必须能安全排除一侧候选，或收缩到可检查的答案边界。先遮住答案，只保留题面，复述候选对象、合法性条件和输出目标。

\textbf{暴力回放。} 慢版本按顺序扫描下标、逐个试答案值，或直接检查候选直到遇到目标边界。它先保证候选空间覆盖完整，但每次只排除一个候选，浪费了题目给出的顺序、比较或可行性边界。二刷时先写慢版本或口述慢版本，用它确认不会漏答案。

\textbf{特例回放。} 拿 [2,2,2,0,1] 手算，mid 和 right 都是 2 时无法判断最小值在哪侧，只能 right-- 丢掉一个重复值。规律是重复元素会削弱二分信息，最坏会退化成线性，但不会丢失最小值。把这个特例继续拆开看：每次比较 mid 之后，必须能指出哪一侧整体不可能包含目标，或哪一侧整体已经满足可行性。二分的核心不是 mid 公式，而是被丢弃区间确实不含答案。复盘时把这个特例压成状态字段：left、right、mid、mid 处比较值或检查返回值、答案区间含义、返回值语义；每一轮更新都要保留目标位置或第一个可行值；再用边界 全相等、重复值跨断点、最小值出现多次 反查初始化、更新顺序和退出条件。

\textbf{状态回放。} 手算路线：手算时列出 left、mid、right，以及 mid 处比较结果、可行性检查结果或局部趋势。每一步都要写出被丢弃的一侧为什么不含答案。状态字段：left、right、mid、mid 处比较值或检查返回值、答案区间含义、返回值语义；每一轮更新都要保留目标位置或第一个可行值。状态升级：把逐个试候选改成维护答案所在区间；边界谓词、可行性检查或局部比较给出分支方向，每个分支都必须能排除一侧候选。

\textbf{证明和边界。} 证明从排除一侧开始：答案空间题证明谓词单调，局部比较题证明保留下来的区间仍含目标；不能只靠经验猜测左右边界。边界：全相等、重复值跨断点、最小值出现多次。变体：这题说明二分依赖的信息不足时会退化，不是所有旋转题都严格对数。

\noindent\textbf{LeetCode 162 寻找峰值。} 模型：比较 mid 和 mid 加一 的斜率能判断某侧一定存在峰值。推导重点：上升则右侧有峰，下降则左侧含 mid 有峰。

\textbf{二刷读题。} 读题时先找候选空间和目标边界。本题可以压成“比较 mid 和 mid 加一 的斜率能判断某侧一定存在峰值”。然后把关键观察写成可执行规则：上升则右侧有峰，下降则左侧含 mid 有峰。这条规则必须能安全排除一侧候选，或收缩到可检查的答案边界。先遮住答案，只保留题面，复述候选对象、合法性条件和输出目标。

\textbf{暴力回放。} 慢版本按顺序扫描下标、逐个试答案值，或直接检查候选直到遇到目标边界。它先保证候选空间覆盖完整，但每次只排除一个候选，浪费了题目给出的顺序、比较或可行性边界。二刷时先写慢版本或口述慢版本，用它确认不会漏答案。

\textbf{特例回放。} 拿 [1,2,3,1] 手算，mid=2 时 nums[mid] 小于右邻 3，说明右侧上坡方向必有峰值；若大于右邻，左侧含 mid 必有峰值。规律是比较 mid 和 mid+1 就能保留一侧峰值区间，边界外可视为负无穷。把这个特例继续拆开看：每次比较 mid 之后，必须能指出哪一侧整体不可能包含目标，或哪一侧整体已经满足可行性。二分的核心不是 mid 公式，而是被丢弃区间确实不含答案。复盘时把这个特例压成状态字段：left、right、mid、mid 处比较值或检查返回值、答案区间含义、返回值语义；每一轮更新都要保留目标位置或第一个可行值；再用边界 长度一、峰在边界、多个峰、相邻不相等假设 反查初始化、更新顺序和退出条件。

\textbf{状态回放。} 手算路线：手算时列出 left、mid、right，以及 mid 处比较结果、可行性检查结果或局部趋势。每一步都要写出被丢弃的一侧为什么不含答案。状态字段：left、right、mid、mid 处比较值或检查返回值、答案区间含义、返回值语义；每一轮更新都要保留目标位置或第一个可行值。状态升级：把逐个试候选改成维护答案所在区间；边界谓词、可行性检查或局部比较给出分支方向，每个分支都必须能排除一侧候选。

\textbf{证明和边界。} 证明从排除一侧开始：答案空间题证明谓词单调，局部比较题证明保留下来的区间仍含目标；不能只靠经验猜测左右边界。边界：长度一、峰在边界、多个峰、相邻不相等假设。变体：二维峰值需要按行列最大值进行更复杂二分。

\noindent\textbf{LeetCode 240 搜索二维矩阵 II。} 模型：右上角同时具有向左变小、向下变大的排除能力。推导重点：当前值大于目标就左移，小于目标就下移。

\textbf{二刷读题。} 读题时先找候选空间和目标边界。本题可以压成“右上角同时具有向左变小、向下变大的排除能力”。然后把关键观察写成可执行规则：当前值大于目标就左移，小于目标就下移。这条规则必须能安全排除一侧候选，或收缩到可检查的答案边界。先遮住答案，只保留题面，复述候选对象、合法性条件和输出目标。

\textbf{暴力回放。} 慢版本按顺序扫描下标、逐个试答案值，或直接检查候选直到遇到目标边界。它先保证候选空间覆盖完整，但每次只排除一个候选，浪费了题目给出的顺序、比较或可行性边界。二刷时先写慢版本或口述慢版本，用它确认不会漏答案。

\textbf{特例回放。} 拿矩阵 [[1,4,7],[2,5,8],[3,6,9]] 找 6 手算，从右上角 7 开始，7 太大就左移到 4，4 太小就下移到 5，再下移到 6。规律是右上角一列向下增大、一行向左减小，每步能排除一行或一列。把这个特例继续拆开看：每次比较 mid 之后，必须能指出哪一侧整体不可能包含目标，或哪一侧整体已经满足可行性。二分的核心不是 mid 公式，而是被丢弃区间确实不含答案。复盘时把这个特例压成状态字段：left、right、mid、mid 处比较值或检查返回值、答案区间含义、返回值语义；每一轮更新都要保留目标位置或第一个可行值；再用边界 空矩阵、单行、单列、目标在角落、重复值 反查初始化、更新顺序和退出条件。

\textbf{状态回放。} 手算路线：手算时列出 left、mid、right，以及 mid 处比较结果、可行性检查结果或局部趋势。每一步都要写出被丢弃的一侧为什么不含答案。状态字段：left、right、mid、mid 处比较值或检查返回值、答案区间含义、返回值语义；每一轮更新都要保留目标位置或第一个可行值。状态升级：把逐个试候选改成维护答案所在区间；边界谓词、可行性检查或局部比较给出分支方向，每个分支都必须能排除一侧候选。

\textbf{证明和边界。} 证明从排除一侧开始：答案空间题证明谓词单调，局部比较题证明保留下来的区间仍含目标；不能只靠经验猜测左右边界。边界：空矩阵、单行、单列、目标在角落、重复值。变体：从左上角无法一次排除一行或一列。

\noindent\textbf{LeetCode 410 分割数组的最大值。} 模型：连续分成 m 段后最大段和越大，越容易完成分割。推导重点：二分最大段和上限，检查函数贪心累计，超过上限就新开一段。

\textbf{二刷读题。} 读题时先找候选空间和目标边界。本题可以压成“连续分成 m 段后最大段和越大，越容易完成分割”。然后把关键观察写成可执行规则：二分最大段和上限，检查函数贪心累计，超过上限就新开一段。这条规则必须能安全排除一侧候选，或收缩到可检查的答案边界。先遮住答案，只保留题面，复述候选对象、合法性条件和输出目标。

\textbf{暴力回放。} 慢版本按顺序扫描下标、逐个试答案值，或直接检查候选直到遇到目标边界。它先保证候选空间覆盖完整，但每次只排除一个候选，浪费了题目给出的顺序、比较或可行性边界。二刷时先写慢版本或口述慢版本，用它确认不会漏答案。

\textbf{特例回放。} 拿 [7,2,5,10,8] 分成 2 段手算，若最大段和限制为 18，可以分成 [7,2,5] 和 [10,8]；限制为 17 就需要三段。规律是答案空间是最大段和，可行性随容量增大单调变真，所以二分最小可行容量。把这个特例继续拆开看：每次比较 mid 之后，必须能指出哪一侧整体不可能包含目标，或哪一侧整体已经满足可行性。二分的核心不是 mid 公式，而是被丢弃区间确实不含答案。复盘时把这个特例压成状态字段：left、right、mid、mid 处比较值或检查返回值、答案区间含义、返回值语义；每一轮更新都要保留目标位置或第一个可行值；再用边界 m 为一、m 等于数组长度、单个元素超过猜测上限、总和溢出 反查初始化、更新顺序和退出条件。

\textbf{状态回放。} 手算路线：手算时列出 left、mid、right，以及 mid 处比较结果、可行性检查结果或局部趋势。每一步都要写出被丢弃的一侧为什么不含答案。状态字段：left、right、mid、mid 处比较值或检查返回值、答案区间含义、返回值语义；每一轮更新都要保留目标位置或第一个可行值。状态升级：把逐个试候选改成维护答案所在区间；边界谓词、可行性检查或局部比较给出分支方向，每个分支都必须能排除一侧候选。

\textbf{证明和边界。} 证明从排除一侧开始：答案空间题证明谓词单调，局部比较题证明保留下来的区间仍含目标；不能只靠经验猜测左右边界。边界：m 为一、m 等于数组长度、单个元素超过猜测上限、总和溢出。变体：也可用 DP 求精确最优，但答案二分更适合大输入。

\noindent\textbf{LeetCode 540 有序数组中的单一元素。} 模型：除一个元素外其余元素都成对出现，单一元素会改变配对下标的奇偶规律。推导重点：把 mid 调整到偶数位，若 nums[mid] 等于 nums[mid+1]，单一元素在右侧，否则在左侧。

\textbf{二刷读题。} 读题时先找候选空间和目标边界。本题可以压成“除一个元素外其余元素都成对出现，单一元素会改变配对下标的奇偶规律”。然后把关键观察写成可执行规则：把 mid 调整到偶数位，若 nums[mid] 等于 nums[mid+1]，单一元素在右侧，否则在左侧。这条规则必须能安全排除一侧候选，或收缩到可检查的答案边界。先遮住答案，只保留题面，复述候选对象、合法性条件和输出目标。

\textbf{暴力回放。} 慢版本按顺序扫描下标、逐个试答案值，或直接检查候选直到遇到目标边界。它先保证候选空间覆盖完整，但每次只排除一个候选，浪费了题目给出的顺序、比较或可行性边界。二刷时先写慢版本或口述慢版本，用它确认不会漏答案。

\textbf{特例回放。} 拿 [1,1,2,3,3] 手算，单一元素 2 之前成对元素从偶数下标开始，之后配对位置会错位。规律是用 mid 的偶偶配对关系判断单一元素在哪边，保持答案区间长度为奇数。把这个特例继续拆开看：每次比较 mid 之后，必须能指出哪一侧整体不可能包含目标，或哪一侧整体已经满足可行性。二分的核心不是 mid 公式，而是被丢弃区间确实不含答案。复盘时把这个特例压成状态字段：left、right、mid、mid 处比较值或检查返回值、答案区间含义、返回值语义；每一轮更新都要保留目标位置或第一个可行值；再用边界 单一元素在开头或结尾、数组长度一、mid 加一越界 反查初始化、更新顺序和退出条件。

\textbf{状态回放。} 手算路线：手算时列出 left、mid、right，以及 mid 处比较结果、可行性检查结果或局部趋势。每一步都要写出被丢弃的一侧为什么不含答案。状态字段：left、right、mid、mid 处比较值或检查返回值、答案区间含义、返回值语义；每一轮更新都要保留目标位置或第一个可行值。状态升级：把逐个试候选改成维护答案所在区间；边界谓词、可行性检查或局部比较给出分支方向，每个分支都必须能排除一侧候选。

\textbf{证明和边界。} 证明从排除一侧开始：答案空间题证明谓词单调，局部比较题证明保留下来的区间仍含目标；不能只靠经验猜测左右边界。边界：单一元素在开头或结尾、数组长度一、mid 加一越界。变体：异或也能求值，但二分利用有序性达到对数时间。

\noindent\textbf{LeetCode 704 二分查找。} 模型：基础题的重点是固定区间语义。推导重点：左闭右开写法能统一 lower bound 和存在性检查。

\textbf{二刷读题。} 读题时先找候选空间和目标边界。本题可以压成“基础题的重点是固定区间语义”。然后把关键观察写成可执行规则：左闭右开写法能统一 lower bound 和存在性检查。这条规则必须能安全排除一侧候选，或收缩到可检查的答案边界。先遮住答案，只保留题面，复述候选对象、合法性条件和输出目标。

\textbf{暴力回放。} 慢版本按顺序扫描下标、逐个试答案值，或直接检查候选直到遇到目标边界。它先保证候选空间覆盖完整，但每次只排除一个候选，浪费了题目给出的顺序、比较或可行性边界。二刷时先写慢版本或口述慢版本，用它确认不会漏答案。

\textbf{特例回放。} 拿 [1,3,5] 找 3 手算，mid 正好命中返回；找 4 时，3 太小丢左侧，5 太大丢右侧，区间为空返回 -1。规律是标准二分每次排除一半，边界写法必须和退出条件匹配。把这个特例继续拆开看：每次比较 mid 之后，必须能指出哪一侧整体不可能包含目标，或哪一侧整体已经满足可行性。二分的核心不是 mid 公式，而是被丢弃区间确实不含答案。复盘时把这个特例压成状态字段：left、right、mid、mid 处比较值或检查返回值、答案区间含义、返回值语义；每一轮更新都要保留目标位置或第一个可行值；再用边界 空数组、长度一、目标在边界、目标不存在 反查初始化、更新顺序和退出条件。

\textbf{状态回放。} 手算路线：手算时列出 left、mid、right，以及 mid 处比较结果、可行性检查结果或局部趋势。每一步都要写出被丢弃的一侧为什么不含答案。状态字段：left、right、mid、mid 处比较值或检查返回值、答案区间含义、返回值语义；每一轮更新都要保留目标位置或第一个可行值。状态升级：把逐个试候选改成维护答案所在区间；边界谓词、可行性检查或局部比较给出分支方向，每个分支都必须能排除一侧候选。

\textbf{证明和边界。} 证明从排除一侧开始：答案空间题证明谓词单调，局部比较题证明保留下来的区间仍含目标；不能只靠经验猜测左右边界。边界：空数组、长度一、目标在边界、目标不存在。变体：所有复杂二分都应该能退回到这道题的边界不变量。

\noindent\textbf{LeetCode 875 爱吃香蕉的珂珂。} 模型：速度越大，总耗时越小，答案具有单调可行性。推导重点：二分速度，检查函数用向上取整累加小时数。

\textbf{二刷读题。} 读题时先找候选空间和目标边界。本题可以压成“速度越大，总耗时越小，答案具有单调可行性”。然后把关键观察写成可执行规则：二分速度，检查函数用向上取整累加小时数。这条规则必须能安全排除一侧候选，或收缩到可检查的答案边界。先遮住答案，只保留题面，复述候选对象、合法性条件和输出目标。

\textbf{暴力回放。} 慢版本按顺序扫描下标、逐个试答案值，或直接检查候选直到遇到目标边界。它先保证候选空间覆盖完整，但每次只排除一个候选，浪费了题目给出的顺序、比较或可行性边界。二刷时先写慢版本或口述慢版本，用它确认不会漏答案。

\textbf{特例回放。} 拿 piles=[3,6,7,11]、h=8 手算，速度 4 时四堆分别需要 1、2、2、3 小时，总计 8 小时，刚好可行；速度 3 会超过 8 小时。规律是速度越大耗时越少，可行性单调，二分最小可行速度。把这个特例继续拆开看：每次比较 mid 之后，必须能指出哪一侧整体不可能包含目标，或哪一侧整体已经满足可行性。二分的核心不是 mid 公式，而是被丢弃区间确实不含答案。复盘时把这个特例压成状态字段：left、right、mid、mid 处比较值或检查返回值、答案区间含义、返回值语义；每一轮更新都要保留目标位置或第一个可行值；再用边界 速度下界、堆很大、h 等于堆数、乘法加法溢出 反查初始化、更新顺序和退出条件。

\textbf{状态回放。} 手算路线：手算时列出 left、mid、right，以及 mid 处比较结果、可行性检查结果或局部趋势。每一步都要写出被丢弃的一侧为什么不含答案。状态字段：left、right、mid、mid 处比较值或检查返回值、答案区间含义、返回值语义；每一轮更新都要保留目标位置或第一个可行值。状态升级：把逐个试候选改成维护答案所在区间；边界谓词、可行性检查或局部比较给出分支方向，每个分支都必须能排除一侧候选。

\textbf{证明和边界。} 证明从排除一侧开始：答案空间题证明谓词单调，局部比较题证明保留下来的区间仍含目标；不能只靠经验猜测左右边界。边界：速度下界、堆很大、h 等于堆数、乘法加法溢出。变体：船运包裹和分割数组都是同类最小可行上限。

\noindent\textbf{LeetCode 981 基于时间的键值存储。} 模型：同一个 key 下的记录按时间戳递增，查询是不超过目标时间的最新值。推导重点：哈希表定位 key 后，在该 key 的时间列表中二分第一个大于目标时间的位置并回退一格。

\textbf{二刷读题。} 读题时先找候选空间和目标边界。本题可以压成“同一个 key 下的记录按时间戳递增，查询是不超过目标时间的最新值”。然后把关键观察写成可执行规则：哈希表定位 key 后，在该 key 的时间列表中二分第一个大于目标时间的位置并回退一格。这条规则必须能安全排除一侧候选，或收缩到可检查的答案边界。先遮住答案，只保留题面，复述候选对象、合法性条件和输出目标。

\textbf{暴力回放。} 慢版本按顺序扫描下标、逐个试答案值，或直接检查候选直到遇到目标边界。它先保证候选空间覆盖完整，但每次只排除一个候选，浪费了题目给出的顺序、比较或可行性边界。二刷时先写慢版本或口述慢版本，用它确认不会漏答案。

\textbf{特例回放。} 拿 key=foo 的记录 (1,bar)、(4,baz)，查询 t=3 手算，答案应是时间不超过 3 的最后一条，也就是 bar。查询 t=4 才是 baz。规律是每个 key 内部时间递增，二分第一个大于查询时间的位置，再回退一格。把这个特例继续拆开看：每次比较 mid 之后，必须能指出哪一侧整体不可能包含目标，或哪一侧整体已经满足可行性。二分的核心不是 mid 公式，而是被丢弃区间确实不含答案。复盘时把这个特例压成状态字段：left、right、mid、mid 处比较值或检查返回值、答案区间含义、返回值语义；每一轮更新都要保留目标位置或第一个可行值；再用边界 key 不存在、查询早于第一条记录、相同 key 多次 set、时间戳递增假设 反查初始化、更新顺序和退出条件。

\textbf{状态回放。} 手算路线：手算时列出 left、mid、right，以及 mid 处比较结果、可行性检查结果或局部趋势。每一步都要写出被丢弃的一侧为什么不含答案。状态字段：left、right、mid、mid 处比较值或检查返回值、答案区间含义、返回值语义；每一轮更新都要保留目标位置或第一个可行值。状态升级：把逐个试候选改成维护答案所在区间；边界谓词、可行性检查或局部比较给出分支方向，每个分支都必须能排除一侧候选。

\textbf{证明和边界。} 证明从排除一侧开始：答案空间题证明谓词单调，局部比较题证明保留下来的区间仍含目标；不能只靠经验猜测左右边界。边界：key 不存在、查询早于第一条记录、相同 key 多次 set、时间戳递增假设。变体：如果时间戳无序到达，必须先排序或改成有序表维护每个 key 的记录。

\noindent\textbf{LeetCode 1011 在 D 天内送达包裹。} 模型：容量越大，需要天数越少，存在最小可行容量。推导重点：下界是最大包裹，上界是总重量，检查函数按顺序贪心装船。

\textbf{二刷读题。} 读题时先找候选空间和目标边界。本题可以压成“容量越大，需要天数越少，存在最小可行容量”。然后把关键观察写成可执行规则：下界是最大包裹，上界是总重量，检查函数按顺序贪心装船。这条规则必须能安全排除一侧候选，或收缩到可检查的答案边界。先遮住答案，只保留题面，复述候选对象、合法性条件和输出目标。

\textbf{暴力回放。} 慢版本按顺序扫描下标、逐个试答案值，或直接检查候选直到遇到目标边界。它先保证候选空间覆盖完整，但每次只排除一个候选，浪费了题目给出的顺序、比较或可行性边界。二刷时先写慢版本或口述慢版本，用它确认不会漏答案。

\textbf{特例回放。} 拿 weights=[1,2,3,4,5]、D=3 手算，容量 6 可以分成 [1,2,3]、[4]、[5] 三天；容量 5 需要更多天。规律是容量越大所需天数越少，可行性单调，二分最小容量。把这个特例继续拆开看：每次比较 mid 之后，必须能指出哪一侧整体不可能包含目标，或哪一侧整体已经满足可行性。二分的核心不是 mid 公式，而是被丢弃区间确实不含答案。复盘时把这个特例压成状态字段：left、right、mid、mid 处比较值或检查返回值、答案区间含义、返回值语义；每一轮更新都要保留目标位置或第一个可行值；再用边界 D 为一、D 等于包裹数、单个包裹超过猜测容量 反查初始化、更新顺序和退出条件。

\textbf{状态回放。} 手算路线：手算时列出 left、mid、right，以及 mid 处比较结果、可行性检查结果或局部趋势。每一步都要写出被丢弃的一侧为什么不含答案。状态字段：left、right、mid、mid 处比较值或检查返回值、答案区间含义、返回值语义；每一轮更新都要保留目标位置或第一个可行值。状态升级：把逐个试候选改成维护答案所在区间；边界谓词、可行性检查或局部比较给出分支方向，每个分支都必须能排除一侧候选。

\textbf{证明和边界。} 证明从排除一侧开始：答案空间题证明谓词单调，局部比较题证明保留下来的区间仍含目标；不能只靠经验猜测左右边界。边界：D 为一、D 等于包裹数、单个包裹超过猜测容量。变体：它和分割数组最大值都是连续分段最大和模型。

\subsection{自测标准}

二刷时不要只确认自己见过这道题。请遮住答案，先讲题面模型，再讲暴力基线和瓶颈，然后说出优化结构保存了什么、不变量如何排除错误候选、边界实验如何打破错误实现。

\section{栈、单调栈与表达式：二刷复盘卡}
这一大节只围绕一个题型复盘，不把题号直接铺成目录。阅读时先看复盘目标，再读核心题卡，最后用自测标准检查自己是否能独立重讲。

\subsection{复盘目标}

追问通常会问栈里为什么存下标、相等元素如何处理、弹出时到底结算了谁。请准备用一个严格递增和一个严格递减样例解释入栈出栈次数。

\subsection{核心题卡}

\noindent\textbf{LeetCode 20 有效的括号。} 模型：括号匹配要求最近打开的左括号先被关闭。推导重点：栈保存尚未匹配的左括号，遇到右括号必须和栈顶类型一致。

\textbf{二刷读题。} 读题时先找未完成候选：哪些对象必须等到右侧、未来字符或后续运算符出现。本题可以压成“括号匹配要求最近打开的左括号先被关闭”，栈只保存这些尚未结算的对象。先遮住答案，只保留题面，复述候选对象、合法性条件和输出目标。

\textbf{暴力回放。} 慢版本让每个位置向左或向右线性寻找边界，或者让每个结构反复等待匹配。它能算出答案，但很多尚未完成的候选会被未来同一个事件一起解决。二刷时先写慢版本或口述慢版本，用它确认不会漏答案。

\textbf{特例回放。} 拿 ([)] 手算：只计数左右括号会误判，因为闭合顺序错了。栈顶必须保存最近的未匹配左括号；遇到 ] 时栈顶若不是 [，立即失败。规律是括号匹配不是数量问题，而是最近未闭合结构的后进先出问题。把这个特例继续拆开看：栈里的对象都是答案尚未确定的候选；一旦当前元素触发弹栈，被弹出的候选必须已经同时拿到了左边界和右边界，未来元素不再有资格改变它的答案。复盘时把这个特例压成状态字段：栈内对象的业务含义、当前输入、弹出时结算的对象、左边界和右边界；栈中存下标还是值要明确；再用边界 空串、以右括号开头、交叉匹配、结束后栈非空 反查初始化、更新顺序和退出条件。

\textbf{状态回放。} 手算路线：手算时记录栈内元素的业务含义，不只记录数值。入栈表示答案未定，弹栈表示当前事件已经让它的答案定下来。状态字段：栈内对象的业务含义、当前输入、弹出时结算的对象、左边界和右边界；栈中存下标还是值要明确。状态升级：把反复寻找最近边界改成维护未完成候选；新元素只和栈顶比较，连续弹出一批已经被当前元素解决的对象。

\textbf{证明和边界。} 证明从弹出时机开始：被弹出的对象在当前时刻答案已经确定，未来元素无法再改变它。边界：空串、以右括号开头、交叉匹配、结束后栈非空。变体：若加入通配符星号，需要额外维护可行区间或双栈。

\noindent\textbf{LeetCode 32 最长有效括号。} 模型：有效括号子串需要知道每段非法边界之后的最早可连接位置。推导重点：栈保存下标，先放入哨兵负一；匹配后用当前下标减栈顶得到长度。

\textbf{二刷读题。} 读题时先找未完成候选：哪些对象必须等到右侧、未来字符或后续运算符出现。本题可以压成“有效括号子串需要知道每段非法边界之后的最早可连接位置”，栈只保存这些尚未结算的对象。先遮住答案，只保留题面，复述候选对象、合法性条件和输出目标。

\textbf{暴力回放。} 慢版本让每个位置向左或向右线性寻找边界，或者让每个结构反复等待匹配。它能算出答案，但很多尚未完成的候选会被未来同一个事件一起解决。二刷时先写慢版本或口述慢版本，用它确认不会漏答案。

\textbf{特例回放。} 拿字符串 )()()) 手算。第一个 ) 不是可匹配对象，而是非法边界；后面每次匹配成功，当前有效长度由当前位置减去栈顶边界得出。若栈空，要把当前右括号作为新非法边界压入。规律是栈里既保存未匹配左括号下标，也保存最近非法边界。把这个特例继续拆开看：栈里的对象都是答案尚未确定的候选；一旦当前元素触发弹栈，被弹出的候选必须已经同时拿到了左边界和右边界，未来元素不再有资格改变它的答案。复盘时把这个特例压成状态字段：栈内对象的业务含义、当前输入、弹出时结算的对象、左边界和右边界；栈中存下标还是值要明确；再用边界 空串、全左括号、全右括号、多个断开的合法段 反查初始化、更新顺序和退出条件。

\textbf{状态回放。} 手算路线：手算时记录栈内元素的业务含义，不只记录数值。入栈表示答案未定，弹栈表示当前事件已经让它的答案定下来。状态字段：栈内对象的业务含义、当前输入、弹出时结算的对象、左边界和右边界；栈中存下标还是值要明确。状态升级：把反复寻找最近边界改成维护未完成候选；新元素只和栈顶比较，连续弹出一批已经被当前元素解决的对象。

\textbf{证明和边界。} 证明从弹出时机开始：被弹出的对象在当前时刻答案已经确定，未来元素无法再改变它。边界：空串、全左括号、全右括号、多个断开的合法段。变体：DP 写法也可做，但栈写法更直接表达最近非法边界。

\noindent\textbf{LeetCode 42 接雨水。} 模型：每个位置的水由左侧最高和右侧最高的较小者决定。推导重点：前后缀、双指针、单调栈都是不同状态维护方式；要能说明各自结算对象。

\textbf{二刷读题。} 读题时先找未完成候选：哪些对象必须等到右侧、未来字符或后续运算符出现。本题可以压成“每个位置的水由左侧最高和右侧最高的较小者决定”，栈只保存这些尚未结算的对象。先遮住答案，只保留题面，复述候选对象、合法性条件和输出目标。

\textbf{暴力回放。} 慢版本让每个位置向左或向右线性寻找边界，或者让每个结构反复等待匹配。它能算出答案，但很多尚未完成的候选会被未来同一个事件一起解决。二刷时先写慢版本或口述慢版本，用它确认不会漏答案。

\textbf{特例回放。} 拿高度 [0,2,0,3] 手算，中间的 0 左侧最高是 2，右侧最高是 3，所以能接 2 格水。若用单调栈，遇到 3 时弹出 0，左边界是 2，右边界是 3，宽度为 1，高度差为 2。规律是每个凹槽必须同时等到左右边界，弹栈那一刻才真正结算水量。把这个特例继续拆开看：栈里的对象都是答案尚未确定的候选；一旦当前元素触发弹栈，被弹出的候选必须已经同时拿到了左边界和右边界，未来元素不再有资格改变它的答案。复盘时把这个特例压成状态字段：栈内对象的业务含义、当前输入、弹出时结算的对象、左边界和右边界；栈中存下标还是值要明确；再用边界 单调递增、单调递减、平台高度、多个凹槽、数组长度不足三 反查初始化、更新顺序和退出条件。

\textbf{状态回放。} 手算路线：手算时记录栈内元素的业务含义，不只记录数值。入栈表示答案未定，弹栈表示当前事件已经让它的答案定下来。状态字段：栈内对象的业务含义、当前输入、弹出时结算的对象、左边界和右边界；栈中存下标还是值要明确。状态升级：把反复寻找最近边界改成维护未完成候选；新元素只和栈顶比较，连续弹出一批已经被当前元素解决的对象。

\textbf{证明和边界。} 证明从弹出时机开始：被弹出的对象在当前时刻答案已经确定，未来元素无法再改变它。边界：单调递增、单调递减、平台高度、多个凹槽、数组长度不足三。变体：若柱子变成二维网格，需要最小堆从边界向内扩展。

\noindent\textbf{LeetCode 84 柱状图中最大的矩形。} 模型：每根柱子作为最矮高度时，需要知道左右第一个更矮位置。推导重点：单调递增栈在遇到更矮柱时结算被弹出柱子的最大宽度。

\textbf{二刷读题。} 读题时先找未完成候选：哪些对象必须等到右侧、未来字符或后续运算符出现。本题可以压成“每根柱子作为最矮高度时，需要知道左右第一个更矮位置”，栈只保存这些尚未结算的对象。先遮住答案，只保留题面，复述候选对象、合法性条件和输出目标。

\textbf{暴力回放。} 慢版本让每个位置向左或向右线性寻找边界，或者让每个结构反复等待匹配。它能算出答案，但很多尚未完成的候选会被未来同一个事件一起解决。二刷时先写慢版本或口述慢版本，用它确认不会漏答案。

\textbf{特例回放。} 拿高度 [2,1,2] 手算。第一个 2 遇到更矮的 1 时，右边界已经确定在当前位置，左边界是弹出后新栈顶的后一位，所以面积可结算。若一直递增，最后需要哨兵触发结算。规律是栈保存递增高度下标，弹出时才知道该柱子的最大扩展宽度。把这个特例继续拆开看：栈里的对象都是答案尚未确定的候选；一旦当前元素触发弹栈，被弹出的候选必须已经同时拿到了左边界和右边界，未来元素不再有资格改变它的答案。复盘时把这个特例压成状态字段：栈内对象的业务含义、当前输入、弹出时结算的对象、左边界和右边界；栈中存下标还是值要明确；再用边界 递增、递减、相等高度、哨兵零、宽度计算 反查初始化、更新顺序和退出条件。

\textbf{状态回放。} 手算路线：手算时记录栈内元素的业务含义，不只记录数值。入栈表示答案未定，弹栈表示当前事件已经让它的答案定下来。状态字段：栈内对象的业务含义、当前输入、弹出时结算的对象、左边界和右边界；栈中存下标还是值要明确。状态升级：把反复寻找最近边界改成维护未完成候选；新元素只和栈顶比较，连续弹出一批已经被当前元素解决的对象。

\textbf{证明和边界。} 证明从弹出时机开始：被弹出的对象在当前时刻答案已经确定，未来元素无法再改变它。边界：递增、递减、相等高度、哨兵零、宽度计算。变体：最大矩形可以作为二维矩阵最大矩形的行压缩子问题。

\noindent\textbf{LeetCode 85 最大矩形。} 模型：每一行都可以把连续的一当成柱状图高度。推导重点：逐行更新每列高度，然后把当前行转成柱状图最大矩形问题。

\textbf{二刷读题。} 读题时先找未完成候选：哪些对象必须等到右侧、未来字符或后续运算符出现。本题可以压成“每一行都可以把连续的一当成柱状图高度”，栈只保存这些尚未结算的对象。先遮住答案，只保留题面，复述候选对象、合法性条件和输出目标。

\textbf{暴力回放。} 慢版本让每个位置向左或向右线性寻找边界，或者让每个结构反复等待匹配。它能算出答案，但很多尚未完成的候选会被未来同一个事件一起解决。二刷时先写慢版本或口述慢版本，用它确认不会漏答案。

\textbf{特例回放。} 拿矩阵每一行向上累计高度手算，某行高度数组若为 [2,1,2]，就退化成柱状图最大矩形。规律是逐行更新每列连续 1 的高度，再对每行高度数组跑单调栈；全零行会把高度清零。把这个特例继续拆开看：栈里的对象都是答案尚未确定的候选；一旦当前元素触发弹栈，被弹出的候选必须已经同时拿到了左边界和右边界，未来元素不再有资格改变它的答案。复盘时把这个特例压成状态字段：栈内对象的业务含义、当前输入、弹出时结算的对象、左边界和右边界；栈中存下标还是值要明确；再用边界 全零、全一、单行、单列、宽度和高度结算 反查初始化、更新顺序和退出条件。

\textbf{状态回放。} 手算路线：手算时记录栈内元素的业务含义，不只记录数值。入栈表示答案未定，弹栈表示当前事件已经让它的答案定下来。状态字段：栈内对象的业务含义、当前输入、弹出时结算的对象、左边界和右边界；栈中存下标还是值要明确。状态升级：把反复寻找最近边界改成维护未完成候选；新元素只和栈顶比较，连续弹出一批已经被当前元素解决的对象。

\textbf{证明和边界。} 证明从弹出时机开始：被弹出的对象在当前时刻答案已经确定，未来元素无法再改变它。边界：全零、全一、单行、单列、宽度和高度结算。变体：这是二维问题压成一维单调栈的典型做法。

\noindent\textbf{LeetCode 150 逆波兰表达式求值。} 模型：后缀表达式把优先级和括号都编码进 token 顺序。推导重点：遇到数字入栈，遇到运算符弹出右操作数和左操作数计算。

\textbf{二刷读题。} 读题时先找未完成候选：哪些对象必须等到右侧、未来字符或后续运算符出现。本题可以压成“后缀表达式把优先级和括号都编码进 token 顺序”，栈只保存这些尚未结算的对象。先遮住答案，只保留题面，复述候选对象、合法性条件和输出目标。

\textbf{暴力回放。} 慢版本让每个位置向左或向右线性寻找边界，或者让每个结构反复等待匹配。它能算出答案，但很多尚未完成的候选会被未来同一个事件一起解决。二刷时先写慢版本或口述慢版本，用它确认不会漏答案。

\textbf{特例回放。} 拿 tokens=[2,1,+,3,*] 手算，2 和 1 入栈，遇到 + 弹出右操作数 1 和左操作数 2 得 3，再与 3 相乘得 9。规律是后缀表达式遇到运算符时，栈顶两个数已经是它的完整左右操作数，减法和除法顺序不能反。把这个特例继续拆开看：栈里的对象都是答案尚未确定的候选；一旦当前元素触发弹栈，被弹出的候选必须已经同时拿到了左边界和右边界，未来元素不再有资格改变它的答案。复盘时把这个特例压成状态字段：栈内对象的业务含义、当前输入、弹出时结算的对象、左边界和右边界；栈中存下标还是值要明确；再用边界 负数 token、多位数、除法向零截断、减法顺序 反查初始化、更新顺序和退出条件。

\textbf{状态回放。} 手算路线：手算时记录栈内元素的业务含义，不只记录数值。入栈表示答案未定，弹栈表示当前事件已经让它的答案定下来。状态字段：栈内对象的业务含义、当前输入、弹出时结算的对象、左边界和右边界；栈中存下标还是值要明确。状态升级：把反复寻找最近边界改成维护未完成候选；新元素只和栈顶比较，连续弹出一批已经被当前元素解决的对象。

\textbf{证明和边界。} 证明从弹出时机开始：被弹出的对象在当前时刻答案已经确定，未来元素无法再改变它。边界：负数 token、多位数、除法向零截断、减法顺序。变体：中缀表达式求值需要另一个运算符栈处理优先级。

\noindent\textbf{LeetCode 155 最小栈。} 模型：普通栈只能看到栈顶，不能常数时间知道历史最小值。推导重点：辅助栈同步保存每个深度对应的最小值。

\textbf{二刷读题。} 读题时先找未完成候选：哪些对象必须等到右侧、未来字符或后续运算符出现。本题可以压成“普通栈只能看到栈顶，不能常数时间知道历史最小值”，栈只保存这些尚未结算的对象。先遮住答案，只保留题面，复述候选对象、合法性条件和输出目标。

\textbf{暴力回放。} 慢版本让每个位置向左或向右线性寻找边界，或者让每个结构反复等待匹配。它能算出答案，但很多尚未完成的候选会被未来同一个事件一起解决。二刷时先写慢版本或口述慢版本，用它确认不会漏答案。

\textbf{特例回放。} 拿 push 2、push 0、push 3、push 0 手算，当前最小值序列是 2、0、0、0；弹出最后一个 0 后，最小值仍是 0。规律是辅助栈要在每个深度同步保存当前最小值，重复最小值不能只保存一次。把这个特例继续拆开看：栈里的对象都是答案尚未确定的候选；一旦当前元素触发弹栈，被弹出的候选必须已经同时拿到了左边界和右边界，未来元素不再有资格改变它的答案。复盘时把这个特例压成状态字段：栈内对象的业务含义、当前输入、弹出时结算的对象、左边界和右边界；栈中存下标还是值要明确；再用边界 重复最小值、弹出最小值、空栈操作约定 反查初始化、更新顺序和退出条件。

\textbf{状态回放。} 手算路线：手算时记录栈内元素的业务含义，不只记录数值。入栈表示答案未定，弹栈表示当前事件已经让它的答案定下来。状态字段：栈内对象的业务含义、当前输入、弹出时结算的对象、左边界和右边界；栈中存下标还是值要明确。状态升级：把反复寻找最近边界改成维护未完成候选；新元素只和栈顶比较，连续弹出一批已经被当前元素解决的对象。

\textbf{证明和边界。} 证明从弹出时机开始：被弹出的对象在当前时刻答案已经确定，未来元素无法再改变它。边界：重复最小值、弹出最小值、空栈操作约定。变体：也可保存差值压缩空间，但可读性和溢出处理更复杂。

\noindent\textbf{LeetCode 224 基本计算器。} 模型：括号会开启新的表达式上下文，括号结束时要回到上一层。推导重点：栈保存进入括号前的结果和符号，当前层算完后乘以上层符号并合并。

\textbf{二刷读题。} 读题时先找未完成候选：哪些对象必须等到右侧、未来字符或后续运算符出现。本题可以压成“括号会开启新的表达式上下文，括号结束时要回到上一层”，栈只保存这些尚未结算的对象。先遮住答案，只保留题面，复述候选对象、合法性条件和输出目标。

\textbf{暴力回放。} 慢版本让每个位置向左或向右线性寻找边界，或者让每个结构反复等待匹配。它能算出答案，但很多尚未完成的候选会被未来同一个事件一起解决。二刷时先写慢版本或口述慢版本，用它确认不会漏答案。

\textbf{特例回放。} 拿 1-(2+3) 手算，进入括号前当前结果是 1，符号是 -；括号内算出 5 后要乘以上层符号并合并成 -4。规律是栈保存进入括号前的结果和符号，一元负号和空格要按字符流处理。把这个特例继续拆开看：栈里的对象都是答案尚未确定的候选；一旦当前元素触发弹栈，被弹出的候选必须已经同时拿到了左边界和右边界，未来元素不再有资格改变它的答案。复盘时把这个特例压成状态字段：栈内对象的业务含义、当前输入、弹出时结算的对象、左边界和右边界；栈中存下标还是值要明确；再用边界 多位数、前导空格、一元负号、嵌套括号 反查初始化、更新顺序和退出条件。

\textbf{状态回放。} 手算路线：手算时记录栈内元素的业务含义，不只记录数值。入栈表示答案未定，弹栈表示当前事件已经让它的答案定下来。状态字段：栈内对象的业务含义、当前输入、弹出时结算的对象、左边界和右边界；栈中存下标还是值要明确。状态升级：把反复寻找最近边界改成维护未完成候选；新元素只和栈顶比较，连续弹出一批已经被当前元素解决的对象。

\textbf{证明和边界。} 证明从弹出时机开始：被弹出的对象在当前时刻答案已经确定，未来元素无法再改变它。边界：多位数、前导空格、一元负号、嵌套括号。变体：基本计算器 II 没有括号但有乘除优先级，需要不同的栈语义。

\noindent\textbf{LeetCode 227 基本计算器 II。} 模型：乘除优先级高于加减，可以在读到下一个运算符时结算上一项。推导重点：栈保存已经确定符号的项，乘除直接修改栈顶，加减压入正负项。

\textbf{二刷读题。} 读题时先找未完成候选：哪些对象必须等到右侧、未来字符或后续运算符出现。本题可以压成“乘除优先级高于加减，可以在读到下一个运算符时结算上一项”，栈只保存这些尚未结算的对象。先遮住答案，只保留题面，复述候选对象、合法性条件和输出目标。

\textbf{暴力回放。} 慢版本让每个位置向左或向右线性寻找边界，或者让每个结构反复等待匹配。它能算出答案，但很多尚未完成的候选会被未来同一个事件一起解决。二刷时先写慢版本或口述慢版本，用它确认不会漏答案。

\textbf{特例回放。} 拿 3+2*2 手算，加号前的 3 可以先入栈，乘号遇到 2 时要立刻把栈顶 2 乘以下一个数 2，而不能等到最后按顺序相加。规律是栈保存已确定符号的项，乘除立即修改栈顶，加减压入正负项。把这个特例继续拆开看：栈里的对象都是答案尚未确定的候选；一旦当前元素触发弹栈，被弹出的候选必须已经同时拿到了左边界和右边界，未来元素不再有资格改变它的答案。复盘时把这个特例压成状态字段：栈内对象的业务含义、当前输入、弹出时结算的对象、左边界和右边界；栈中存下标还是值要明确；再用边界 多位数、空格、除法向零截断、最后一个数字需要结算 反查初始化、更新顺序和退出条件。

\textbf{状态回放。} 手算路线：手算时记录栈内元素的业务含义，不只记录数值。入栈表示答案未定，弹栈表示当前事件已经让它的答案定下来。状态字段：栈内对象的业务含义、当前输入、弹出时结算的对象、左边界和右边界；栈中存下标还是值要明确。状态升级：把反复寻找最近边界改成维护未完成候选；新元素只和栈顶比较，连续弹出一批已经被当前元素解决的对象。

\textbf{证明和边界。} 证明从弹出时机开始：被弹出的对象在当前时刻答案已经确定，未来元素无法再改变它。边界：多位数、空格、除法向零截断、最后一个数字需要结算。变体：若加入括号，需要再引入上下文栈或递归下降解析。

\noindent\textbf{LeetCode 232 用栈实现队列。} 模型：两个栈分别负责输入顺序和输出顺序。推导重点：只有输出栈为空时，才把输入栈整体倒过去，保证均摊常数。

\textbf{二刷读题。} 读题时先找未完成候选：哪些对象必须等到右侧、未来字符或后续运算符出现。本题可以压成“两个栈分别负责输入顺序和输出顺序”，栈只保存这些尚未结算的对象。先遮住答案，只保留题面，复述候选对象、合法性条件和输出目标。

\textbf{暴力回放。} 慢版本让每个位置向左或向右线性寻找边界，或者让每个结构反复等待匹配。它能算出答案，但很多尚未完成的候选会被未来同一个事件一起解决。二刷时先写慢版本或口述慢版本，用它确认不会漏答案。

\textbf{特例回放。} 拿 push 1、push 2、pop 手算，如果输出栈为空，把输入栈整体倒过去得到栈顶 1，正好符合队列先进先出。规律是只有输出栈为空才搬运，单个元素最多进出两个栈各一次，所以均摊常数。把这个特例继续拆开看：栈里的对象都是答案尚未确定的候选；一旦当前元素触发弹栈，被弹出的候选必须已经同时拿到了左边界和右边界，未来元素不再有资格改变它的答案。复盘时把这个特例压成状态字段：栈内对象的业务含义、当前输入、弹出时结算的对象、左边界和右边界；栈中存下标还是值要明确；再用边界 连续 peek、空队列、交替 push pop 反查初始化、更新顺序和退出条件。

\textbf{状态回放。} 手算路线：手算时记录栈内元素的业务含义，不只记录数值。入栈表示答案未定，弹栈表示当前事件已经让它的答案定下来。状态字段：栈内对象的业务含义、当前输入、弹出时结算的对象、左边界和右边界；栈中存下标还是值要明确。状态升级：把反复寻找最近边界改成维护未完成候选；新元素只和栈顶比较，连续弹出一批已经被当前元素解决的对象。

\textbf{证明和边界。} 证明从弹出时机开始：被弹出的对象在当前时刻答案已经确定，未来元素无法再改变它。边界：连续 peek、空队列、交替 push pop。变体：用队列实现栈则要通过轮转保持最近元素在队首。

\noindent\textbf{LeetCode 239 滑动窗口最大值。} 模型：每个窗口最大值来自仍未过期且没有被新元素支配的候选。推导重点：单调队列保存下标，队尾小于等于新值的候选被弹出。

\textbf{二刷读题。} 读题时先找未完成候选：哪些对象必须等到右侧、未来字符或后续运算符出现。本题可以压成“每个窗口最大值来自仍未过期且没有被新元素支配的候选”，栈只保存这些尚未结算的对象。先遮住答案，只保留题面，复述候选对象、合法性条件和输出目标。

\textbf{暴力回放。} 慢版本让每个位置向左或向右线性寻找边界，或者让每个结构反复等待匹配。它能算出答案，但很多尚未完成的候选会被未来同一个事件一起解决。二刷时先写慢版本或口述慢版本，用它确认不会漏答案。

\textbf{特例回放。} 拿 [1,3,-1]、窗口 3 手算，队尾的 1 在 3 进入后永远不可能成为最大值，因为 3 更新且更靠右。于是 1 可以弹出。规律是单调队列保存可能成为当前或未来窗口最大值的下标，过期下标从队首清理。把这个特例继续拆开看：栈里的对象都是答案尚未确定的候选；一旦当前元素触发弹栈，被弹出的候选必须已经同时拿到了左边界和右边界，未来元素不再有资格改变它的答案。复盘时把这个特例压成状态字段：栈内对象的业务含义、当前输入、弹出时结算的对象、左边界和右边界；栈中存下标还是值要明确；再用边界 k 为一、k 等于数组长度、重复最大值、队首过期 反查初始化、更新顺序和退出条件。

\textbf{状态回放。} 手算路线：手算时记录栈内元素的业务含义，不只记录数值。入栈表示答案未定，弹栈表示当前事件已经让它的答案定下来。状态字段：栈内对象的业务含义、当前输入、弹出时结算的对象、左边界和右边界；栈中存下标还是值要明确。状态升级：把反复寻找最近边界改成维护未完成候选；新元素只和栈顶比较，连续弹出一批已经被当前元素解决的对象。

\textbf{证明和边界。} 证明从弹出时机开始：被弹出的对象在当前时刻答案已经确定，未来元素无法再改变它。边界：k 为一、k 等于数组长度、重复最大值、队首过期。变体：受限子序列和使用同样队列维护最近 k 个 DP 最大值。

\noindent\textbf{LeetCode 394 字符串解码。} 模型：嵌套结构需要保存上一层字符串和重复次数。推导重点：遇到左括号入栈当前状态，右括号弹出并拼接展开。

\textbf{二刷读题。} 读题时先找未完成候选：哪些对象必须等到右侧、未来字符或后续运算符出现。本题可以压成“嵌套结构需要保存上一层字符串和重复次数”，栈只保存这些尚未结算的对象。先遮住答案，只保留题面，复述候选对象、合法性条件和输出目标。

\textbf{暴力回放。} 慢版本让每个位置向左或向右线性寻找边界，或者让每个结构反复等待匹配。它能算出答案，但很多尚未完成的候选会被未来同一个事件一起解决。二刷时先写慢版本或口述慢版本，用它确认不会漏答案。

\textbf{特例回放。} 拿 3[a2[c]] 手算，遇到左括号前保存当前字符串和重复次数；遇到右括号时弹出上一层，把当前 cc 重复 2 次拼成 acc，再被外层重复 3 次。规律是栈保存进入括号前的上下文，右括号触发一次完整展开。把这个特例继续拆开看：栈里的对象都是答案尚未确定的候选；一旦当前元素触发弹栈，被弹出的候选必须已经同时拿到了左边界和右边界，未来元素不再有资格改变它的答案。复盘时把这个特例压成状态字段：栈内对象的业务含义、当前输入、弹出时结算的对象、左边界和右边界；栈中存下标还是值要明确；再用边界 多位数字、嵌套、空片段、数字后必须跟括号 反查初始化、更新顺序和退出条件。

\textbf{状态回放。} 手算路线：手算时记录栈内元素的业务含义，不只记录数值。入栈表示答案未定，弹栈表示当前事件已经让它的答案定下来。状态字段：栈内对象的业务含义、当前输入、弹出时结算的对象、左边界和右边界；栈中存下标还是值要明确。状态升级：把反复寻找最近边界改成维护未完成候选；新元素只和栈顶比较，连续弹出一批已经被当前元素解决的对象。

\textbf{证明和边界。} 证明从弹出时机开始：被弹出的对象在当前时刻答案已经确定，未来元素无法再改变它。边界：多位数字、嵌套、空片段、数字后必须跟括号。变体：递归下降解析更通用，但迭代栈更符合工程栈控制。

\noindent\textbf{LeetCode 402 移掉 K 位数字。} 模型：要删除 k 位使剩余数字最小，本质是让高位尽量小。推导重点：单调递增栈中遇到更小数字时，弹出前面较大的数字并消耗删除次数。

\textbf{二刷读题。} 读题时先找未完成候选：哪些对象必须等到右侧、未来字符或后续运算符出现。本题可以压成“要删除 k 位使剩余数字最小，本质是让高位尽量小”，栈只保存这些尚未结算的对象。先遮住答案，只保留题面，复述候选对象、合法性条件和输出目标。

\textbf{暴力回放。} 慢版本让每个位置向左或向右线性寻找边界，或者让每个结构反复等待匹配。它能算出答案，但很多尚未完成的候选会被未来同一个事件一起解决。二刷时先写慢版本或口述慢版本，用它确认不会漏答案。

\textbf{特例回放。} 拿 1432219、k=3 手算，遇到 3 后面的 2 时，前面的 4 和 3 都比 2 大且可删，弹出能让高位更小；最后去掉前导零。规律是单调递增栈让更小数字尽量靠前，k 未用完时从尾部继续删。把这个特例继续拆开看：栈里的对象都是答案尚未确定的候选；一旦当前元素触发弹栈，被弹出的候选必须已经同时拿到了左边界和右边界，未来元素不再有资格改变它的答案。复盘时把这个特例压成状态字段：栈内对象的业务含义、当前输入、弹出时结算的对象、左边界和右边界；栈中存下标还是值要明确；再用边界 前导零、k 未用完、k 等于长度、数字已经递增 反查初始化、更新顺序和退出条件。

\textbf{状态回放。} 手算路线：手算时记录栈内元素的业务含义，不只记录数值。入栈表示答案未定，弹栈表示当前事件已经让它的答案定下来。状态字段：栈内对象的业务含义、当前输入、弹出时结算的对象、左边界和右边界；栈中存下标还是值要明确。状态升级：把反复寻找最近边界改成维护未完成候选；新元素只和栈顶比较，连续弹出一批已经被当前元素解决的对象。

\textbf{证明和边界。} 证明从弹出时机开始：被弹出的对象在当前时刻答案已经确定，未来元素无法再改变它。边界：前导零、k 未用完、k 等于长度、数字已经递增。变体：去除重复字母也用单调栈，但还要保证每个字符保留一次。

\noindent\textbf{LeetCode 496 下一个更大元素 I。} 模型：第二个数组提供每个值右侧第一个更大值的全局关系。推导重点：单调栈预处理值到下一个更大值映射，再回答第一个数组查询。

\textbf{二刷读题。} 读题时先找未完成候选：哪些对象必须等到右侧、未来字符或后续运算符出现。本题可以压成“第二个数组提供每个值右侧第一个更大值的全局关系”，栈只保存这些尚未结算的对象。先遮住答案，只保留题面，复述候选对象、合法性条件和输出目标。

\textbf{暴力回放。} 慢版本让每个位置向左或向右线性寻找边界，或者让每个结构反复等待匹配。它能算出答案，但很多尚未完成的候选会被未来同一个事件一起解决。二刷时先写慢版本或口述慢版本，用它确认不会漏答案。

\textbf{特例回放。} 拿 nums2=[1,3,4,2] 手算，1 的下一个更大是 3，3 的下一个更大是 4，4 没有。规律是单调栈在扫描 nums2 时建立值到下一个更大值的映射，再回答 nums1 查询。把这个特例继续拆开看：栈里的对象都是答案尚未确定的候选；一旦当前元素触发弹栈，被弹出的候选必须已经同时拿到了左边界和右边界，未来元素不再有资格改变它的答案。复盘时把这个特例压成状态字段：栈内对象的业务含义、当前输入、弹出时结算的对象、左边界和右边界；栈中存下标还是值要明确；再用边界 不存在更大值、递减数组、值唯一假设 反查初始化、更新顺序和退出条件。

\textbf{状态回放。} 手算路线：手算时记录栈内元素的业务含义，不只记录数值。入栈表示答案未定，弹栈表示当前事件已经让它的答案定下来。状态字段：栈内对象的业务含义、当前输入、弹出时结算的对象、左边界和右边界；栈中存下标还是值要明确。状态升级：把反复寻找最近边界改成维护未完成候选；新元素只和栈顶比较，连续弹出一批已经被当前元素解决的对象。

\textbf{证明和边界。} 证明从弹出时机开始：被弹出的对象在当前时刻答案已经确定，未来元素无法再改变它。边界：不存在更大值、递减数组、值唯一假设。变体：若值不唯一，应按下标而不是值建映射。

\noindent\textbf{LeetCode 503 下一个更大元素 II。} 模型：环形数组可以通过遍历两倍下标模拟绕回。推导重点：第一遍入栈，第二遍只帮助未解决下标找到答案。

\textbf{二刷读题。} 读题时先找未完成候选：哪些对象必须等到右侧、未来字符或后续运算符出现。本题可以压成“环形数组可以通过遍历两倍下标模拟绕回”，栈只保存这些尚未结算的对象。先遮住答案，只保留题面，复述候选对象、合法性条件和输出目标。

\textbf{暴力回放。} 慢版本让每个位置向左或向右线性寻找边界，或者让每个结构反复等待匹配。它能算出答案，但很多尚未完成的候选会被未来同一个事件一起解决。二刷时先写慢版本或口述慢版本，用它确认不会漏答案。

\textbf{特例回放。} 拿 [1,2,1] 手算，第一个 1 的下一个更大是 2，最后一个 1 环形回到前面也能找到 2。规律是扫描两遍模拟环形，第一遍入栈，第二遍只帮助未解决下标结算。把这个特例继续拆开看：栈里的对象都是答案尚未确定的候选；一旦当前元素触发弹栈，被弹出的候选必须已经同时拿到了左边界和右边界，未来元素不再有资格改变它的答案。复盘时把这个特例压成状态字段：栈内对象的业务含义、当前输入、弹出时结算的对象、左边界和右边界；栈中存下标还是值要明确；再用边界 全递减、全相等、长度一、取模下标 反查初始化、更新顺序和退出条件。

\textbf{状态回放。} 手算路线：手算时记录栈内元素的业务含义，不只记录数值。入栈表示答案未定，弹栈表示当前事件已经让它的答案定下来。状态字段：栈内对象的业务含义、当前输入、弹出时结算的对象、左边界和右边界；栈中存下标还是值要明确。状态升级：把反复寻找最近边界改成维护未完成候选；新元素只和栈顶比较，连续弹出一批已经被当前元素解决的对象。

\textbf{证明和边界。} 证明从弹出时机开始：被弹出的对象在当前时刻答案已经确定，未来元素无法再改变它。边界：全递减、全相等、长度一、取模下标。变体：不要在第二遍重复入栈，否则可能死循环或重复计算。

\noindent\textbf{LeetCode 739 每日温度。} 模型：每一天等待右侧第一个更高温度。推导重点：单调递减栈保存尚未找到答案的下标。

\textbf{二刷读题。} 读题时先找未完成候选：哪些对象必须等到右侧、未来字符或后续运算符出现。本题可以压成“每一天等待右侧第一个更高温度”，栈只保存这些尚未结算的对象。先遮住答案，只保留题面，复述候选对象、合法性条件和输出目标。

\textbf{暴力回放。} 慢版本让每个位置向左或向右线性寻找边界，或者让每个结构反复等待匹配。它能算出答案，但很多尚未完成的候选会被未来同一个事件一起解决。二刷时先写慢版本或口述慢版本，用它确认不会漏答案。

\textbf{特例回放。} 拿 [73,74,75,71,69,72,76,73] 手算，73 遇到 74 时等待 1 天；71 要等到 72 才解决。规律是单调递减栈保存尚未等到更高温的下标，遇到更高温时连续弹栈结算天数。把这个特例继续拆开看：栈里的对象都是答案尚未确定的候选；一旦当前元素触发弹栈，被弹出的候选必须已经同时拿到了左边界和右边界，未来元素不再有资格改变它的答案。复盘时把这个特例压成状态字段：栈内对象的业务含义、当前输入、弹出时结算的对象、左边界和右边界；栈中存下标还是值要明确；再用边界 没有更高温、相等温度、递增、递减 反查初始化、更新顺序和退出条件。

\textbf{状态回放。} 手算路线：手算时记录栈内元素的业务含义，不只记录数值。入栈表示答案未定，弹栈表示当前事件已经让它的答案定下来。状态字段：栈内对象的业务含义、当前输入、弹出时结算的对象、左边界和右边界；栈中存下标还是值要明确。状态升级：把反复寻找最近边界改成维护未完成候选；新元素只和栈顶比较，连续弹出一批已经被当前元素解决的对象。

\textbf{证明和边界。} 证明从弹出时机开始：被弹出的对象在当前时刻答案已经确定，未来元素无法再改变它。边界：没有更高温、相等温度、递增、递减。变体：下一个更大元素与它同型，只是输出值或距离不同。

\noindent\textbf{LeetCode 895 最大频率栈。} 模型：弹出时既要频率最高，又要在同频中最近入栈。推导重点：哈希表记录值频率，频率到栈记录达到该频率的入栈顺序，维护当前最大频率。

\textbf{二刷读题。} 读题时先找未完成候选：哪些对象必须等到右侧、未来字符或后续运算符出现。本题可以压成“弹出时既要频率最高，又要在同频中最近入栈”，栈只保存这些尚未结算的对象。先遮住答案，只保留题面，复述候选对象、合法性条件和输出目标。

\textbf{暴力回放。} 慢版本让每个位置向左或向右线性寻找边界，或者让每个结构反复等待匹配。它能算出答案，但很多尚未完成的候选会被未来同一个事件一起解决。二刷时先写慢版本或口述慢版本，用它确认不会漏答案。

\textbf{特例回放。} 拿 push 5,7,5,7,4,5 后 pop 手算，频率最高是 5，先弹 5；下一次 5 和 7 同频，弹最近达到该频率的 7。规律是值到频率、频率到栈、当前最大频率三者共同维护。把这个特例继续拆开看：栈里的对象都是答案尚未确定的候选；一旦当前元素触发弹栈，被弹出的候选必须已经同时拿到了左边界和右边界，未来元素不再有资格改变它的答案。复盘时把这个特例压成状态字段：栈内对象的业务含义、当前输入、弹出时结算的对象、左边界和右边界；栈中存下标还是值要明确；再用边界 重复 push、pop 后最大频率下降、同频最近性、空栈不在题意内 反查初始化、更新顺序和退出条件。

\textbf{状态回放。} 手算路线：手算时记录栈内元素的业务含义，不只记录数值。入栈表示答案未定，弹栈表示当前事件已经让它的答案定下来。状态字段：栈内对象的业务含义、当前输入、弹出时结算的对象、左边界和右边界；栈中存下标还是值要明确。状态升级：把反复寻找最近边界改成维护未完成候选；新元素只和栈顶比较，连续弹出一批已经被当前元素解决的对象。

\textbf{证明和边界。} 证明从弹出时机开始：被弹出的对象在当前时刻答案已经确定，未来元素无法再改变它。边界：重复 push、pop 后最大频率下降、同频最近性、空栈不在题意内。变体：它是栈和哈希表按频率维度组合的设计题。

\subsection{自测标准}

二刷时不要只确认自己见过这道题。请遮住答案，先讲题面模型，再讲暴力基线和瓶颈，然后说出优化结构保存了什么、不变量如何排除错误候选、边界实验如何打破错误实现。

\section{堆与动态极值：二刷复盘卡}
这一大节只围绕一个题型复盘，不把题号直接铺成目录。阅读时先看复盘目标，再读核心题卡，最后用自测标准检查自己是否能独立重讲。

\subsection{复盘目标}

追问通常会问为什么不完整排序、堆顶是否可能过期、比较器方向是否写反。请准备说明堆里元素的业务含义，而不是只说 priority\_queue。

\subsection{核心题卡}

\noindent\textbf{LeetCode 23 合并 K 个升序链表。} 模型：每条链表当前头节点都是可能的下一个最小元素。推导重点：小顶堆保存每条非空链的当前头，弹出最小节点后把它的后继入堆。

\textbf{二刷读题。} 读题时先找动态候选集合以及每一步需要的极值。本题可以压成“每条链表当前头节点都是可能的下一个最小元素”，堆顶必须正好回答下一步选择。先遮住答案，只保留题面，复述候选对象、合法性条件和输出目标。

\textbf{暴力回放。} 慢版本每轮扫描全部候选来找最大或最小，再更新候选集合。它的答案选择过程清楚，但动态集合每变化一次就重新找极值，代价太高。二刷时先写慢版本或口述慢版本，用它确认不会漏答案。

\textbf{特例回放。} 拿三条链表头 1、1、2 手算，每次答案只需要当前最小头节点；弹出某条链的头后，只有它的后继会成为新候选。暴力每轮扫 k 个头，堆把“找当前最小头”变成对数操作。规律是堆里始终最多保存每条非空链的一个当前头。把这个特例继续拆开看：堆顶必须正好代表下一步最需要处理的候选；若堆里可能有过期对象，读取堆顶前要先清理。堆只解决动态极值，不自动证明被弹出或被丢弃的元素无用。复盘时把这个特例压成状态字段：堆元素字段、堆顶比较规则、候选是否过期、答案读取时机；如果需要懒删除，还要保存删除计数；再用边界 空链表数组、某条链为空、重复值、堆里保存节点指针还是值 反查初始化、更新顺序和退出条件。

\textbf{状态回放。} 手算路线：手算时记录堆顶、候选集合和是否有过期元素。每次使用堆顶前都要确认它仍然属于当前问题。状态字段：堆元素字段、堆顶比较规则、候选是否过期、答案读取时机；如果需要懒删除，还要保存删除计数。状态升级：把每轮全量扫描改成堆顶查询；插入、弹出和过期清理共同维护动态候选集。

\textbf{证明和边界。} 证明从候选覆盖开始：所有可能成为下一步答案的对象都在堆中，堆顶过期时不会被使用。边界：空链表数组、某条链为空、重复值、堆里保存节点指针还是值。变体：也可以分治两两合并，堆更适合 K 较大且链长不均的场景。

\noindent\textbf{LeetCode 215 数组中的第 K 个最大元素。} 模型：只关心第 K 大，不需要完整排序。推导重点：大小为 K 的小顶堆保存当前前 K 大边界；快速选择可做到平均线性。

\textbf{二刷读题。} 读题时先找动态候选集合以及每一步需要的极值。本题可以压成“只关心第 K 大，不需要完整排序”，堆顶必须正好回答下一步选择。先遮住答案，只保留题面，复述候选对象、合法性条件和输出目标。

\textbf{暴力回放。} 慢版本每轮扫描全部候选来找最大或最小，再更新候选集合。它的答案选择过程清楚，但动态集合每变化一次就重新找极值，代价太高。二刷时先写慢版本或口述慢版本，用它确认不会漏答案。

\textbf{特例回放。} 拿 [3,2,1,5,6,4]、k=2 手算，固定大小为 2 的小顶堆最终保存最大的两个数，堆顶就是第 2 大。规律是堆里只保留当前前 k 大候选，比堆顶小的数不可能进入答案。把这个特例继续拆开看：堆顶必须正好代表下一步最需要处理的候选；若堆里可能有过期对象，读取堆顶前要先清理。堆只解决动态极值，不自动证明被弹出或被丢弃的元素无用。复盘时把这个特例压成状态字段：堆元素字段、堆顶比较规则、候选是否过期、答案读取时机；如果需要懒删除，还要保存删除计数；再用边界 K 为一、K 为 n、重复值、随机 pivot 反查初始化、更新顺序和退出条件。

\textbf{状态回放。} 手算路线：手算时记录堆顶、候选集合和是否有过期元素。每次使用堆顶前都要确认它仍然属于当前问题。状态字段：堆元素字段、堆顶比较规则、候选是否过期、答案读取时机；如果需要懒删除，还要保存删除计数。状态升级：把每轮全量扫描改成堆顶查询；插入、弹出和过期清理共同维护动态候选集。

\textbf{证明和边界。} 证明从候选覆盖开始：所有可能成为下一步答案的对象都在堆中，堆顶过期时不会被使用。边界：K 为一、K 为 n、重复值、随机 pivot。变体：若数据流持续到来，堆方案比一次性快速选择更自然。

\noindent\textbf{LeetCode 295 数据流的中位数。} 模型：数据流需要动态维护较小一半和较大一半。推导重点：大顶堆保存较小一半，小顶堆保存较大一半，大小差不超过一。

\textbf{二刷读题。} 读题时先找动态候选集合以及每一步需要的极值。本题可以压成“数据流需要动态维护较小一半和较大一半”，堆顶必须正好回答下一步选择。先遮住答案，只保留题面，复述候选对象、合法性条件和输出目标。

\textbf{暴力回放。} 慢版本每轮扫描全部候选来找最大或最小，再更新候选集合。它的答案选择过程清楚，但动态集合每变化一次就重新找极值，代价太高。二刷时先写慢版本或口述慢版本，用它确认不会漏答案。

\textbf{特例回放。} 拿数据流 1、2、3 手算，插入 1 后中位数 1；插入 2 后两半是 [1] 和 [2]，中位数 1.5；插入 3 后右半多一个，要平衡成左半 [2,1]、右半 [3]。规律是大顶堆保存较小一半，小顶堆保存较大一半，大小差最多 1。把这个特例继续拆开看：堆顶必须正好代表下一步最需要处理的候选；若堆里可能有过期对象，读取堆顶前要先清理。堆只解决动态极值，不自动证明被弹出或被丢弃的元素无用。复盘时把这个特例压成状态字段：堆元素字段、堆顶比较规则、候选是否过期、答案读取时机；如果需要懒删除，还要保存删除计数；再用边界 偶数个元素、负数、重复值、平均值溢出 反查初始化、更新顺序和退出条件。

\textbf{状态回放。} 手算路线：手算时记录堆顶、候选集合和是否有过期元素。每次使用堆顶前都要确认它仍然属于当前问题。状态字段：堆元素字段、堆顶比较规则、候选是否过期、答案读取时机；如果需要懒删除，还要保存删除计数。状态升级：把每轮全量扫描改成堆顶查询；插入、弹出和过期清理共同维护动态候选集。

\textbf{证明和边界。} 证明从候选覆盖开始：所有可能成为下一步答案的对象都在堆中，堆顶过期时不会被使用。边界：偶数个元素、负数、重复值、平均值溢出。变体：若还要删除任意元素，需要延迟删除哈希表。

\noindent\textbf{LeetCode 347 前 K 个高频元素。} 模型：先统计频次，再只维护前 K 个候选。推导重点：小顶堆大小超过 K 就弹出最低频；桶排序利用频次上界。

\textbf{二刷读题。} 读题时先找动态候选集合以及每一步需要的极值。本题可以压成“先统计频次，再只维护前 K 个候选”，堆顶必须正好回答下一步选择。先遮住答案，只保留题面，复述候选对象、合法性条件和输出目标。

\textbf{暴力回放。} 慢版本每轮扫描全部候选来找最大或最小，再更新候选集合。它的答案选择过程清楚，但动态集合每变化一次就重新找极值，代价太高。二刷时先写慢版本或口述慢版本，用它确认不会漏答案。

\textbf{特例回放。} 拿 [1,1,1,2,2,3]、k=2 手算，频次是 1->3、2->2、3->1，前两个高频是 1 和 2。规律是先用哈希计数，再用大小为 k 的小顶堆或桶按频次取 Top K；频次相同不影响题目集合答案。把这个特例继续拆开看：堆顶必须正好代表下一步最需要处理的候选；若堆里可能有过期对象，读取堆顶前要先清理。堆只解决动态极值，不自动证明被弹出或被丢弃的元素无用。复盘时把这个特例压成状态字段：堆元素字段、堆顶比较规则、候选是否过期、答案读取时机；如果需要懒删除，还要保存删除计数；再用边界 K 等于不同元素数、频次相同、结果顺序不要求 反查初始化、更新顺序和退出条件。

\textbf{状态回放。} 手算路线：手算时记录堆顶、候选集合和是否有过期元素。每次使用堆顶前都要确认它仍然属于当前问题。状态字段：堆元素字段、堆顶比较规则、候选是否过期、答案读取时机；如果需要懒删除，还要保存删除计数。状态升级：把每轮全量扫描改成堆顶查询；插入、弹出和过期清理共同维护动态候选集。

\textbf{证明和边界。} 证明从候选覆盖开始：所有可能成为下一步答案的对象都在堆中，堆顶过期时不会被使用。边界：K 等于不同元素数、频次相同、结果顺序不要求。变体：若要按频次和字典序排序，比较器要同时处理两个键。

\noindent\textbf{LeetCode 480 滑动窗口中位数。} 模型：窗口移动时既要加入新元素，也要删除滑出的旧元素并查询中位数。推导重点：双堆维护左右两半，延迟删除表记录已离窗但尚未到堆顶的元素。

\textbf{二刷读题。} 读题时先找动态候选集合以及每一步需要的极值。本题可以压成“窗口移动时既要加入新元素，也要删除滑出的旧元素并查询中位数”，堆顶必须正好回答下一步选择。先遮住答案，只保留题面，复述候选对象、合法性条件和输出目标。

\textbf{暴力回放。} 慢版本每轮扫描全部候选来找最大或最小，再更新候选集合。它的答案选择过程清楚，但动态集合每变化一次就重新找极值，代价太高。二刷时先写慢版本或口述慢版本，用它确认不会漏答案。

\textbf{特例回放。} 拿 [1,3,-1]、k=3 手算，中位数是 1；窗口右移时要删除 1、加入 -3，普通堆不能任意删除。规律是双堆维护两半，再用延迟删除清理过期元素，堆顶使用前必须有效。把这个特例继续拆开看：堆顶必须正好代表下一步最需要处理的候选；若堆里可能有过期对象，读取堆顶前要先清理。堆只解决动态极值，不自动证明被弹出或被丢弃的元素无用。复盘时把这个特例压成状态字段：堆元素字段、堆顶比较规则、候选是否过期、答案读取时机；如果需要懒删除，还要保存删除计数；再用边界 重复值、偶数窗口、平均值溢出、堆顶过期 反查初始化、更新顺序和退出条件。

\textbf{状态回放。} 手算路线：手算时记录堆顶、候选集合和是否有过期元素。每次使用堆顶前都要确认它仍然属于当前问题。状态字段：堆元素字段、堆顶比较规则、候选是否过期、答案读取时机；如果需要懒删除，还要保存删除计数。状态升级：把每轮全量扫描改成堆顶查询；插入、弹出和过期清理共同维护动态候选集。

\textbf{证明和边界。} 证明从候选覆盖开始：所有可能成为下一步答案的对象都在堆中，堆顶过期时不会被使用。边界：重复值、偶数窗口、平均值溢出、堆顶过期。变体：普通 priority\_queue 不能任意删除，所以必须解释懒删除。

\noindent\textbf{LeetCode 621 任务调度器。} 模型：相同任务之间有冷却间隔，核心受最高频任务骨架限制。推导重点：可用大顶堆模拟，也可用最高频公式计算最短时间。

\textbf{二刷读题。} 读题时先找动态候选集合以及每一步需要的极值。本题可以压成“相同任务之间有冷却间隔，核心受最高频任务骨架限制”，堆顶必须正好回答下一步选择。先遮住答案，只保留题面，复述候选对象、合法性条件和输出目标。

\textbf{暴力回放。} 慢版本每轮扫描全部候选来找最大或最小，再更新候选集合。它的答案选择过程清楚，但动态集合每变化一次就重新找极值，代价太高。二刷时先写慢版本或口述慢版本，用它确认不会漏答案。

\textbf{特例回放。} 拿 A,A,A,B,B,B、n=2 手算，若只按顺序执行 A 会违反冷却，必须在相同任务之间间隔两个时间单位。规律是最高频任务决定骨架长度，也可用堆按剩余次数模拟，每轮安排最多 n+1 个不同任务。把这个特例继续拆开看：堆顶必须正好代表下一步最需要处理的候选；若堆里可能有过期对象，读取堆顶前要先清理。堆只解决动态极值，不自动证明被弹出或被丢弃的元素无用。复盘时把这个特例压成状态字段：堆元素字段、堆顶比较规则、候选是否过期、答案读取时机；如果需要懒删除，还要保存删除计数；再用边界 冷却为零、多个最高频、空闲被其他任务填满 反查初始化、更新顺序和退出条件。

\textbf{状态回放。} 手算路线：手算时记录堆顶、候选集合和是否有过期元素。每次使用堆顶前都要确认它仍然属于当前问题。状态字段：堆元素字段、堆顶比较规则、候选是否过期、答案读取时机；如果需要懒删除，还要保存删除计数。状态升级：把每轮全量扫描改成堆顶查询；插入、弹出和过期清理共同维护动态候选集。

\textbf{证明和边界。} 证明从候选覆盖开始：所有可能成为下一步答案的对象都在堆中，堆顶过期时不会被使用。边界：冷却为零、多个最高频、空闲被其他任务填满。变体：若任务有不同释放时间或执行时长，就变成更一般的调度问题。

\noindent\textbf{LeetCode 632 最小区间覆盖 K 个列表。} 模型：每个列表至少取一个元素时，当前区间由最小当前值和最大当前值决定。推导重点：小顶堆保存各列表当前元素，同时维护当前最大值；弹出最小所在列表并推进。

\textbf{二刷读题。} 读题时先找动态候选集合以及每一步需要的极值。本题可以压成“每个列表至少取一个元素时，当前区间由最小当前值和最大当前值决定”，堆顶必须正好回答下一步选择。先遮住答案，只保留题面，复述候选对象、合法性条件和输出目标。

\textbf{暴力回放。} 慢版本每轮扫描全部候选来找最大或最小，再更新候选集合。它的答案选择过程清楚，但动态集合每变化一次就重新找极值，代价太高。二刷时先写慢版本或口述慢版本，用它确认不会漏答案。

\textbf{特例回放。} 拿三组 [4,10]、[0,9]、[5,18] 手算，当前头是 4、0、5，区间 [0,5] 覆盖三组；弹出最小头 0 后推进该组。规律是堆保存每组当前元素，同时维护当前最大值，最小头决定下一步推进哪一组。把这个特例继续拆开看：堆顶必须正好代表下一步最需要处理的候选；若堆里可能有过期对象，读取堆顶前要先清理。堆只解决动态极值，不自动证明被弹出或被丢弃的元素无用。复盘时把这个特例压成状态字段：堆元素字段、堆顶比较规则、候选是否过期、答案读取时机；如果需要懒删除，还要保存删除计数；再用边界 某个列表为空、重复值、区间长度相同的 tie-breaker、列表耗尽 反查初始化、更新顺序和退出条件。

\textbf{状态回放。} 手算路线：手算时记录堆顶、候选集合和是否有过期元素。每次使用堆顶前都要确认它仍然属于当前问题。状态字段：堆元素字段、堆顶比较规则、候选是否过期、答案读取时机；如果需要懒删除，还要保存删除计数。状态升级：把每轮全量扫描改成堆顶查询；插入、弹出和过期清理共同维护动态候选集。

\textbf{证明和边界。} 证明从候选覆盖开始：所有可能成为下一步答案的对象都在堆中，堆顶过期时不会被使用。边界：某个列表为空、重复值、区间长度相同的 tie-breaker、列表耗尽。变体：这是多路归并和覆盖窗口的结合。

\noindent\textbf{LeetCode 703 数据流中的第 K 大元素。} 模型：数据不断到达，只关心当前第 K 大而非完整排序。推导重点：维护大小为 K 的小顶堆，堆顶就是当前第 K 大元素。

\textbf{二刷读题。} 读题时先找动态候选集合以及每一步需要的极值。本题可以压成“数据不断到达，只关心当前第 K 大而非完整排序”，堆顶必须正好回答下一步选择。先遮住答案，只保留题面，复述候选对象、合法性条件和输出目标。

\textbf{暴力回放。} 慢版本每轮扫描全部候选来找最大或最小，再更新候选集合。它的答案选择过程清楚，但动态集合每变化一次就重新找极值，代价太高。二刷时先写慢版本或口述慢版本，用它确认不会漏答案。

\textbf{特例回放。} 拿 k=3，初始 4,5,8,2 手算，小顶堆保存 4,5,8，堆顶 4 是第 3 大；add 3 不进入堆，add 5 会替换 4。规律是数据流只保留前 k 大，堆顶就是当前第 k 大。把这个特例继续拆开看：堆顶必须正好代表下一步最需要处理的候选；若堆里可能有过期对象，读取堆顶前要先清理。堆只解决动态极值，不自动证明被弹出或被丢弃的元素无用。复盘时把这个特例压成状态字段：堆元素字段、堆顶比较规则、候选是否过期、答案读取时机；如果需要懒删除，还要保存删除计数；再用边界 初始元素少于 K、重复值、K 为一、连续 add 反查初始化、更新顺序和退出条件。

\textbf{状态回放。} 手算路线：手算时记录堆顶、候选集合和是否有过期元素。每次使用堆顶前都要确认它仍然属于当前问题。状态字段：堆元素字段、堆顶比较规则、候选是否过期、答案读取时机；如果需要懒删除，还要保存删除计数。状态升级：把每轮全量扫描改成堆顶查询；插入、弹出和过期清理共同维护动态候选集。

\textbf{证明和边界。} 证明从候选覆盖开始：所有可能成为下一步答案的对象都在堆中，堆顶过期时不会被使用。边界：初始元素少于 K、重复值、K 为一、连续 add。变体：这题是 Top K 的在线版本，和一次性排序的适用场景不同。

\noindent\textbf{LeetCode 857 雇佣 K 名工人的最低成本。} 模型：工资比例由当前队伍中最高工资质量比决定，总成本等于比例乘质量和。推导重点：按比例升序扫描，用大顶堆维护当前 K 个工人的最小质量和。

\textbf{二刷读题。} 读题时先找动态候选集合以及每一步需要的极值。本题可以压成“工资比例由当前队伍中最高工资质量比决定，总成本等于比例乘质量和”，堆顶必须正好回答下一步选择。先遮住答案，只保留题面，复述候选对象、合法性条件和输出目标。

\textbf{暴力回放。} 慢版本每轮扫描全部候选来找最大或最小，再更新候选集合。它的答案选择过程清楚，但动态集合每变化一次就重新找极值，代价太高。二刷时先写慢版本或口述慢版本，用它确认不会漏答案。

\textbf{特例回放。} 拿 quality=[10,20,5]、wage=[70,50,30] 手算，按工资质量比选组时，当前最大比率决定组内所有人的单位工资；再用堆保留质量和最小的 k 人。规律是排序枚举最高比率，堆优化质量总和。把这个特例继续拆开看：堆顶必须正好代表下一步最需要处理的候选；若堆里可能有过期对象，读取堆顶前要先清理。堆只解决动态极值，不自动证明被弹出或被丢弃的元素无用。复盘时把这个特例压成状态字段：堆元素字段、堆顶比较规则、候选是否过期、答案读取时机；如果需要懒删除，还要保存删除计数；再用边界 K 为一、比例相同、浮点精度、质量和更新 反查初始化、更新顺序和退出条件。

\textbf{状态回放。} 手算路线：手算时记录堆顶、候选集合和是否有过期元素。每次使用堆顶前都要确认它仍然属于当前问题。状态字段：堆元素字段、堆顶比较规则、候选是否过期、答案读取时机；如果需要懒删除，还要保存删除计数。状态升级：把每轮全量扫描改成堆顶查询；插入、弹出和过期清理共同维护动态候选集。

\textbf{证明和边界。} 证明从候选覆盖开始：所有可能成为下一步答案的对象都在堆中，堆顶过期时不会被使用。边界：K 为一、比例相同、浮点精度、质量和更新。变体：这是排序确定价格比例，堆维护可替换候选的组合题。

\noindent\textbf{LeetCode 871 最低加油次数。} 模型：行驶过程中，经过的加油站形成可反悔的燃料候选。推导重点：按位置扫描，油量不足到达下一点时，从已路过站点中取最大燃料补充。

\textbf{二刷读题。} 读题时先找动态候选集合以及每一步需要的极值。本题可以压成“行驶过程中，经过的加油站形成可反悔的燃料候选”，堆顶必须正好回答下一步选择。先遮住答案，只保留题面，复述候选对象、合法性条件和输出目标。

\textbf{暴力回放。} 慢版本每轮扫描全部候选来找最大或最小，再更新候选集合。它的答案选择过程清楚，但动态集合每变化一次就重新找极值，代价太高。二刷时先写慢版本或口述慢版本，用它确认不会漏答案。

\textbf{特例回放。} 拿 target=100、startFuel=10、站点 (10,60),(20,30),(30,30),(60,40) 手算，到达 10 后把 60 放入堆，燃料不够去下一站时先加此前可达站里最大油量。规律是能到的站都作为候选，加油时选最大油量，保证次数最少。把这个特例继续拆开看：堆顶必须正好代表下一步最需要处理的候选；若堆里可能有过期对象，读取堆顶前要先清理。堆只解决动态极值，不自动证明被弹出或被丢弃的元素无用。复盘时把这个特例压成状态字段：堆元素字段、堆顶比较规则、候选是否过期、答案读取时机；如果需要懒删除，还要保存删除计数；再用边界 起点油足够、无法到达第一个站、多个站同位置、目标作为终点事件 反查初始化、更新顺序和退出条件。

\textbf{状态回放。} 手算路线：手算时记录堆顶、候选集合和是否有过期元素。每次使用堆顶前都要确认它仍然属于当前问题。状态字段：堆元素字段、堆顶比较规则、候选是否过期、答案读取时机；如果需要懒删除，还要保存删除计数。状态升级：把每轮全量扫描改成堆顶查询；插入、弹出和过期清理共同维护动态候选集。

\textbf{证明和边界。} 证明从候选覆盖开始：所有可能成为下一步答案的对象都在堆中，堆顶过期时不会被使用。边界：起点油足够、无法到达第一个站、多个站同位置、目标作为终点事件。变体：这是反悔贪心，堆里保存的是已经经过但尚未使用的资源。

\noindent\textbf{LeetCode 973 最接近原点的 K 个点。} 模型：距离只需比较平方，避免开方误差和成本。推导重点：大小为 K 的大顶堆保存当前最近 K 个点中最远者。

\textbf{二刷读题。} 读题时先找动态候选集合以及每一步需要的极值。本题可以压成“距离只需比较平方，避免开方误差和成本”，堆顶必须正好回答下一步选择。先遮住答案，只保留题面，复述候选对象、合法性条件和输出目标。

\textbf{暴力回放。} 慢版本每轮扫描全部候选来找最大或最小，再更新候选集合。它的答案选择过程清楚，但动态集合每变化一次就重新找极值，代价太高。二刷时先写慢版本或口述慢版本，用它确认不会漏答案。

\textbf{特例回放。} 拿点 (1,3) 和 (-2,2) 手算，距离平方分别是 10 和 8，所以 k=1 时选 (-2,2)。规律是比较距离平方避免开方，用大小为 k 的大顶堆保留当前最近点。把这个特例继续拆开看：堆顶必须正好代表下一步最需要处理的候选；若堆里可能有过期对象，读取堆顶前要先清理。堆只解决动态极值，不自动证明被弹出或被丢弃的元素无用。复盘时把这个特例压成状态字段：堆元素字段、堆顶比较规则、候选是否过期、答案读取时机；如果需要懒删除，还要保存删除计数；再用边界 K 为一、K 为全部、距离相同、坐标平方溢出 反查初始化、更新顺序和退出条件。

\textbf{状态回放。} 手算路线：手算时记录堆顶、候选集合和是否有过期元素。每次使用堆顶前都要确认它仍然属于当前问题。状态字段：堆元素字段、堆顶比较规则、候选是否过期、答案读取时机；如果需要懒删除，还要保存删除计数。状态升级：把每轮全量扫描改成堆顶查询；插入、弹出和过期清理共同维护动态候选集。

\textbf{证明和边界。} 证明从候选覆盖开始：所有可能成为下一步答案的对象都在堆中，堆顶过期时不会被使用。边界：K 为一、K 为全部、距离相同、坐标平方溢出。变体：快速选择可平均线性，但堆更适合数据流。

\noindent\textbf{LeetCode 1046 最后一块石头的重量。} 模型：每次都要取当前最重的两块石头碰撞。推导重点：大顶堆保存所有重量，弹出两个最大值，若不同则把差值放回。

\textbf{二刷读题。} 读题时先找动态候选集合以及每一步需要的极值。本题可以压成“每次都要取当前最重的两块石头碰撞”，堆顶必须正好回答下一步选择。先遮住答案，只保留题面，复述候选对象、合法性条件和输出目标。

\textbf{暴力回放。} 慢版本每轮扫描全部候选来找最大或最小，再更新候选集合。它的答案选择过程清楚，但动态集合每变化一次就重新找极值，代价太高。二刷时先写慢版本或口述慢版本，用它确认不会漏答案。

\textbf{特例回放。} 拿 [2,7,4,1,8,1] 手算，每次取两块最大石头 8 和 7，撞剩 1 再放回。规律是动态最大值查询，用大顶堆保存当前石头重量，直到剩 0 或 1 块。把这个特例继续拆开看：堆顶必须正好代表下一步最需要处理的候选；若堆里可能有过期对象，读取堆顶前要先清理。堆只解决动态极值，不自动证明被弹出或被丢弃的元素无用。复盘时把这个特例压成状态字段：堆元素字段、堆顶比较规则、候选是否过期、答案读取时机；如果需要懒删除，还要保存删除计数；再用边界 没有石头、一块石头、两块相等、多次产生相同重量 反查初始化、更新顺序和退出条件。

\textbf{状态回放。} 手算路线：手算时记录堆顶、候选集合和是否有过期元素。每次使用堆顶前都要确认它仍然属于当前问题。状态字段：堆元素字段、堆顶比较规则、候选是否过期、答案读取时机；如果需要懒删除，还要保存删除计数。状态升级：把每轮全量扫描改成堆顶查询；插入、弹出和过期清理共同维护动态候选集。

\textbf{证明和边界。} 证明从候选覆盖开始：所有可能成为下一步答案的对象都在堆中，堆顶过期时不会被使用。边界：没有石头、一块石头、两块相等、多次产生相同重量。变体：最后一块石头 II 则是分组差值最小化，模型转为 0-1 背包。

\noindent\textbf{LeetCode 1353 最多可以参加的会议数目。} 模型：每天只能参加一个正在进行的会议，应优先选择最早结束的会议。推导重点：按开始日排序扫描日期，把当天可参加会议的结束日入小顶堆，弹出最早结束者。

\textbf{二刷读题。} 读题时先找动态候选集合以及每一步需要的极值。本题可以压成“每天只能参加一个正在进行的会议，应优先选择最早结束的会议”，堆顶必须正好回答下一步选择。先遮住答案，只保留题面，复述候选对象、合法性条件和输出目标。

\textbf{暴力回放。} 慢版本每轮扫描全部候选来找最大或最小，再更新候选集合。它的答案选择过程清楚，但动态集合每变化一次就重新找极值，代价太高。二刷时先写慢版本或口述慢版本，用它确认不会漏答案。

\textbf{特例回放。} 拿 events=[[1,2],[2,3],[3,4]] 手算，第 1 天把开始的会议入堆并参加最早结束的；第 2 天再清理过期并参加一个。规律是按天扫描，用小顶堆保存当前可参加会议的结束日，每天选最早结束者。把这个特例继续拆开看：堆顶必须正好代表下一步最需要处理的候选；若堆里可能有过期对象，读取堆顶前要先清理。堆只解决动态极值，不自动证明被弹出或被丢弃的元素无用。复盘时把这个特例压成状态字段：堆元素字段、堆顶比较规则、候选是否过期、答案读取时机；如果需要懒删除，还要保存删除计数；再用边界 同一天多个会议、过期会议清理、空闲日期跳跃、结束日相同 反查初始化、更新顺序和退出条件。

\textbf{状态回放。} 手算路线：手算时记录堆顶、候选集合和是否有过期元素。每次使用堆顶前都要确认它仍然属于当前问题。状态字段：堆元素字段、堆顶比较规则、候选是否过期、答案读取时机；如果需要懒删除，还要保存删除计数。状态升级：把每轮全量扫描改成堆顶查询；插入、弹出和过期清理共同维护动态候选集。

\textbf{证明和边界。} 证明从候选覆盖开始：所有可能成为下一步答案的对象都在堆中，堆顶过期时不会被使用。边界：同一天多个会议、过期会议清理、空闲日期跳跃、结束日相同。变体：这是扫描线和堆结合的调度题。

\noindent\textbf{LeetCode 1675 数组的最小偏移量。} 模型：奇数可以先乘二，偶数可以不断除二，目标是缩小当前最大最小差。推导重点：把所有数规范到可降低的偶数形态，用大顶堆维护当前最大值，同时记录当前最小值。

\textbf{二刷读题。} 读题时先找动态候选集合以及每一步需要的极值。本题可以压成“奇数可以先乘二，偶数可以不断除二，目标是缩小当前最大最小差”，堆顶必须正好回答下一步选择。先遮住答案，只保留题面，复述候选对象、合法性条件和输出目标。

\textbf{暴力回放。} 慢版本每轮扫描全部候选来找最大或最小，再更新候选集合。它的答案选择过程清楚，但动态集合每变化一次就重新找极值，代价太高。二刷时先写慢版本或口述慢版本，用它确认不会漏答案。

\textbf{特例回放。} 拿 [1,2,3,4] 手算，奇数可先翻倍成偶数，之后只能不断把当前最大偶数减半；每次记录最大值和当前最小值差。规律是把所有数变到可下降的最高状态，再用大顶堆逐步降低最大值。把这个特例继续拆开看：堆顶必须正好代表下一步最需要处理的候选；若堆里可能有过期对象，读取堆顶前要先清理。堆只解决动态极值，不自动证明被弹出或被丢弃的元素无用。复盘时把这个特例压成状态字段：堆元素字段、堆顶比较规则、候选是否过期、答案读取时机；如果需要懒删除，还要保存删除计数；再用边界 全奇数、全偶数、最大值变奇数停止、数值范围 反查初始化、更新顺序和退出条件。

\textbf{状态回放。} 手算路线：手算时记录堆顶、候选集合和是否有过期元素。每次使用堆顶前都要确认它仍然属于当前问题。状态字段：堆元素字段、堆顶比较规则、候选是否过期、答案读取时机；如果需要懒删除，还要保存删除计数。状态升级：把每轮全量扫描改成堆顶查询；插入、弹出和过期清理共同维护动态候选集。

\textbf{证明和边界。} 证明从候选覆盖开始：所有可能成为下一步答案的对象都在堆中，堆顶过期时不会被使用。边界：全奇数、全偶数、最大值变奇数停止、数值范围。变体：这是堆维护动态极值和状态空间规范化的结合。

\noindent\textbf{LeetCode 1851 包含每个查询的最小区间。} 模型：每个查询点需要找覆盖它的最短区间，区间和查询都可离线排序。推导重点：按查询值升序扫描，把左端不大于查询的区间入堆，堆顶过期区间弹出。

\textbf{二刷读题。} 读题时先找动态候选集合以及每一步需要的极值。本题可以压成“每个查询点需要找覆盖它的最短区间，区间和查询都可离线排序”，堆顶必须正好回答下一步选择。先遮住答案，只保留题面，复述候选对象、合法性条件和输出目标。

\textbf{暴力回放。} 慢版本每轮扫描全部候选来找最大或最小，再更新候选集合。它的答案选择过程清楚，但动态集合每变化一次就重新找极值，代价太高。二刷时先写慢版本或口述慢版本，用它确认不会漏答案。

\textbf{特例回放。} 拿 intervals=[[1,4],[2,4],[3,6]]、queries=[2,3,5] 手算，查询 2 时可选 [1,4] 和 [2,4]，较短的是 [2,4]。规律是离线按查询升序加入左端不大于查询的区间，堆顶清理右端已过期区间。把这个特例继续拆开看：堆顶必须正好代表下一步最需要处理的候选；若堆里可能有过期对象，读取堆顶前要先清理。堆只解决动态极值，不自动证明被弹出或被丢弃的元素无用。复盘时把这个特例压成状态字段：堆元素字段、堆顶比较规则、候选是否过期、答案读取时机；如果需要懒删除，还要保存删除计数；再用边界 没有覆盖区间、区间长度相同、查询重复、右端过期 反查初始化、更新顺序和退出条件。

\textbf{状态回放。} 手算路线：手算时记录堆顶、候选集合和是否有过期元素。每次使用堆顶前都要确认它仍然属于当前问题。状态字段：堆元素字段、堆顶比较规则、候选是否过期、答案读取时机；如果需要懒删除，还要保存删除计数。状态升级：把每轮全量扫描改成堆顶查询；插入、弹出和过期清理共同维护动态候选集。

\textbf{证明和边界。} 证明从候选覆盖开始：所有可能成为下一步答案的对象都在堆中，堆顶过期时不会被使用。边界：没有覆盖区间、区间长度相同、查询重复、右端过期。变体：这是离线扫描线、堆和区间覆盖的组合题。

\subsection{自测标准}

二刷时不要只确认自己见过这道题。请遮住答案，先讲题面模型，再讲暴力基线和瓶颈，然后说出优化结构保存了什么、不变量如何排除错误候选、边界实验如何打破错误实现。

\section{链表与指针重连：二刷复盘卡}
这一大节只围绕一个题型复盘，不把题号直接铺成目录。阅读时先看复盘目标，再读核心题卡，最后用自测标准检查自己是否能独立重讲。

\subsection{复盘目标}

追问通常会让你画指针。请准备说明上一段尾、当前段头、当前节点、后继节点分别是谁，以及哪一步保存 next 防止断链。

\subsection{核心题卡}

\noindent\textbf{LeetCode 2 两数相加。} 模型：链表按低位到高位保存数字，本质是竖式加法而不是整数转换。推导重点：维护两个游标和进位，任一链未结束或进位非零都继续生成节点。

\textbf{二刷读题。} 读题时先区分节点值、节点身份和连接关系。本题可以压成“链表按低位到高位保存数字，本质是竖式加法而不是整数转换”，后续推导要围绕哪些指针会断开、哪些节点仍要可达。先遮住答案，只保留题面，复述候选对象、合法性条件和输出目标。

\textbf{暴力回放。} 慢版本可以先把节点顺序抄到数组里，按数组推导目标顺序。回到链表后，难点不再是值的顺序，而是节点身份和 next 指针不能丢。二刷时先写慢版本或口述慢版本，用它确认不会漏答案。

\textbf{特例回放。} 拿 2->4->3 加 5->6->4 手算，个位 2+5 得到 7，十位 4+6 得到 10，所以当前位写 0 并把进位 1 带到下一位。再走百位 3+4+1 得到 8。规律是链表从低位开始，状态只需要两个游标、当前进位和结果尾节点；两条链长度不同或最后还有进位时，循环仍不能停。把这个特例继续拆开看：每个指针都要有角色，上一段尾、当前段头、当前节点和后继入口不能混。只要一次改 next 之前没有保存后继，就可能丢掉整段未处理链。复盘时把这个特例压成状态字段：虚拟头、上一段尾、当前段头、当前节点、后继节点；任何改 next 的操作之前都要保存后续入口；再用边界 两链长度不同、最后进位、输入为零、不能用内置整数承载超长数字 反查初始化、更新顺序和退出条件。

\textbf{状态回放。} 手算路线：手算时画出 prev、cur、next 和下一段头节点。每次改 next 前先保存后继，这是链表推导的安全线。状态字段：虚拟头、上一段尾、当前段头、当前节点、后继节点；任何改 next 的操作之前都要保存后续入口。状态升级：把数组辅助结构换成几个稳定指针角色；重连顺序必须保证已处理链合法、未处理链仍可达。

\textbf{证明和边界。} 证明从链结构开始：已处理链连通、未处理链可达、二者连接点明确，节点没有丢失也没有成环。边界：两链长度不同、最后进位、输入为零、不能用内置整数承载超长数字。变体：若链表高位在前，可以用栈反向处理或先反转链表。

\noindent\textbf{LeetCode 19 删除链表的倒数第 N 个结点。} 模型：倒数位置可以转成两个指针之间固定 n 步的距离。推导重点：快指针先走 n 步，再让快慢指针同步移动，慢指针停在待删节点前驱。

\textbf{二刷读题。} 读题时先区分节点值、节点身份和连接关系。本题可以压成“倒数位置可以转成两个指针之间固定 n 步的距离”，后续推导要围绕哪些指针会断开、哪些节点仍要可达。先遮住答案，只保留题面，复述候选对象、合法性条件和输出目标。

\textbf{暴力回放。} 慢版本可以先把节点顺序抄到数组里，按数组推导目标顺序。回到链表后，难点不再是值的顺序，而是节点身份和 next 指针不能丢。二刷时先写慢版本或口述慢版本，用它确认不会漏答案。

\textbf{特例回放。} 拿长度 5 删除倒数第 2 个手算。快指针先走 2 步，慢指针和快指针保持固定距离；当快指针到尾后，慢指针正好停在待删节点前驱。再试删除头节点，若没有虚拟头就缺前驱。规律是虚拟头统一头部删除，快慢指针固定距离定位前驱。把这个特例继续拆开看：每个指针都要有角色，上一段尾、当前段头、当前节点和后继入口不能混。只要一次改 next 之前没有保存后继，就可能丢掉整段未处理链。复盘时把这个特例压成状态字段：虚拟头、上一段尾、当前段头、当前节点、后继节点；任何改 next 的操作之前都要保存后续入口；再用边界 删除头节点、n 等于链表长度、只有一个节点、题目保证 n 合法 反查初始化、更新顺序和退出条件。

\textbf{状态回放。} 手算路线：手算时画出 prev、cur、next 和下一段头节点。每次改 next 前先保存后继，这是链表推导的安全线。状态字段：虚拟头、上一段尾、当前段头、当前节点、后继节点；任何改 next 的操作之前都要保存后续入口。状态升级：把数组辅助结构换成几个稳定指针角色；重连顺序必须保证已处理链合法、未处理链仍可达。

\textbf{证明和边界。} 证明从链结构开始：已处理链连通、未处理链可达、二者连接点明确，节点没有丢失也没有成环。边界：删除头节点、n 等于链表长度、只有一个节点、题目保证 n 合法。变体：若 n 不保证合法，需要在快指针预走阶段检测长度不足。

\noindent\textbf{LeetCode 21 合并两个有序链表。} 模型：两个链表已经各自有序，下一节点只能来自两个当前头中较小者。推导重点：虚拟头统一结果链表，尾指针不断接上较小节点并推进来源链。

\textbf{二刷读题。} 读题时先区分节点值、节点身份和连接关系。本题可以压成“两个链表已经各自有序，下一节点只能来自两个当前头中较小者”，后续推导要围绕哪些指针会断开、哪些节点仍要可达。先遮住答案，只保留题面，复述候选对象、合法性条件和输出目标。

\textbf{暴力回放。} 慢版本可以先把节点顺序抄到数组里，按数组推导目标顺序。回到链表后，难点不再是值的顺序，而是节点身份和 next 指针不能丢。二刷时先写慢版本或口述慢版本，用它确认不会漏答案。

\textbf{特例回放。} 拿 1->3 和 2->4 手算，先接 1，再比较 3 和 2 接 2，之后接 3 和 4。规律是结果尾指针每次接较小头节点，某条链耗尽后直接接另一条剩余部分；虚拟头统一处理首节点。把这个特例继续拆开看：每个指针都要有角色，上一段尾、当前段头、当前节点和后继入口不能混。只要一次改 next 之前没有保存后继，就可能丢掉整段未处理链。复盘时把这个特例压成状态字段：虚拟头、上一段尾、当前段头、当前节点、后继节点；任何改 next 的操作之前都要保存后续入口；再用边界 某一条链为空、重复值、剩余整段接回、是否复用原节点 反查初始化、更新顺序和退出条件。

\textbf{状态回放。} 手算路线：手算时画出 prev、cur、next 和下一段头节点。每次改 next 前先保存后继，这是链表推导的安全线。状态字段：虚拟头、上一段尾、当前段头、当前节点、后继节点；任何改 next 的操作之前都要保存后续入口。状态升级：把数组辅助结构换成几个稳定指针角色；重连顺序必须保证已处理链合法、未处理链仍可达。

\textbf{证明和边界。} 证明从链结构开始：已处理链连通、未处理链可达、二者连接点明确，节点没有丢失也没有成环。边界：某一条链为空、重复值、剩余整段接回、是否复用原节点。变体：合并 K 条链表可用分治或堆，核心仍是维护每条链的当前头。

\noindent\textbf{LeetCode 24 两两交换链表中的节点。} 模型：链表按固定长度为二的小组做局部重连。推导重点：使用组前驱、第一节点、第二节点和下一组头，按顺序改指针并推进前驱。

\textbf{二刷读题。} 读题时先区分节点值、节点身份和连接关系。本题可以压成“链表按固定长度为二的小组做局部重连”，后续推导要围绕哪些指针会断开、哪些节点仍要可达。先遮住答案，只保留题面，复述候选对象、合法性条件和输出目标。

\textbf{暴力回放。} 慢版本可以先把节点顺序抄到数组里，按数组推导目标顺序。回到链表后，难点不再是值的顺序，而是节点身份和 next 指针不能丢。二刷时先写慢版本或口述慢版本，用它确认不会漏答案。

\textbf{特例回放。} 拿 1->2->3->4 手算，第一组交换后应是 2->1，并且 1 要接到下一组交换后的头。规律是每次固定 pre、first、second、next 四个指针，先保存 next，再重连 second->first，first->next。把这个特例继续拆开看：每个指针都要有角色，上一段尾、当前段头、当前节点和后继入口不能混。只要一次改 next 之前没有保存后继，就可能丢掉整段未处理链。复盘时把这个特例压成状态字段：虚拟头、上一段尾、当前段头、当前节点、后继节点；任何改 next 的操作之前都要保存后续入口；再用边界 空链、单节点、奇数长度、交换后前驱更新到组尾 反查初始化、更新顺序和退出条件。

\textbf{状态回放。} 手算路线：手算时画出 prev、cur、next 和下一段头节点。每次改 next 前先保存后继，这是链表推导的安全线。状态字段：虚拟头、上一段尾、当前段头、当前节点、后继节点；任何改 next 的操作之前都要保存后续入口。状态升级：把数组辅助结构换成几个稳定指针角色；重连顺序必须保证已处理链合法、未处理链仍可达。

\textbf{证明和边界。} 证明从链结构开始：已处理链连通、未处理链可达、二者连接点明确，节点没有丢失也没有成环。边界：空链、单节点、奇数长度、交换后前驱更新到组尾。变体：K 个一组翻转是同一思想的推广，只是组内反转更长。

\noindent\textbf{LeetCode 25 K 个一组翻转链表。} 模型：链表被切成长度为 K 的连续组，只有完整组才反转。推导重点：每轮先探测 K 个节点确认完整，再保存下一组头并做局部反转接回。

\textbf{二刷读题。} 读题时先区分节点值、节点身份和连接关系。本题可以压成“链表被切成长度为 K 的连续组，只有完整组才反转”，后续推导要围绕哪些指针会断开、哪些节点仍要可达。先遮住答案，只保留题面，复述候选对象、合法性条件和输出目标。

\textbf{暴力回放。} 慢版本可以先把节点顺序抄到数组里，按数组推导目标顺序。回到链表后，难点不再是值的顺序，而是节点身份和 next 指针不能丢。二刷时先写慢版本或口述慢版本，用它确认不会漏答案。

\textbf{特例回放。} 拿 1->2->3->4、k=2 画图。第一组边界是 1 和 2，下一组入口是 3；反转前如果不保存 3，链会断。反转后上一组尾接 2，旧头 1 接 3。规律是每组先找 kth 和 group\_next，再局部反转并接回前后两段。把这个特例继续拆开看：每个指针都要有角色，上一段尾、当前段头、当前节点和后继入口不能混。只要一次改 next 之前没有保存后继，就可能丢掉整段未处理链。复盘时把这个特例压成状态字段：虚拟头、上一段尾、当前段头、当前节点、后继节点；任何改 next 的操作之前都要保存后续入口；再用边界 不足 K 个不反转、K 为一、刚好整除、反转后头尾接错 反查初始化、更新顺序和退出条件。

\textbf{状态回放。} 手算路线：手算时画出 prev、cur、next 和下一段头节点。每次改 next 前先保存后继，这是链表推导的安全线。状态字段：虚拟头、上一段尾、当前段头、当前节点、后继节点；任何改 next 的操作之前都要保存后续入口。状态升级：把数组辅助结构换成几个稳定指针角色；重连顺序必须保证已处理链合法、未处理链仍可达。

\textbf{证明和边界。} 证明从链结构开始：已处理链连通、未处理链可达、二者连接点明确，节点没有丢失也没有成环。边界：不足 K 个不反转、K 为一、刚好整除、反转后头尾接错。变体：递归写法短但可能有栈风险，迭代写法更接近工程实现。

\noindent\textbf{LeetCode 61 旋转链表。} 模型：右旋 k 位等价于把链表尾部一段切到头部。推导重点：先求长度并让 k 对长度取模，再找到新尾节点并重连成新头。

\textbf{二刷读题。} 读题时先区分节点值、节点身份和连接关系。本题可以压成“右旋 k 位等价于把链表尾部一段切到头部”，后续推导要围绕哪些指针会断开、哪些节点仍要可达。先遮住答案，只保留题面，复述候选对象、合法性条件和输出目标。

\textbf{暴力回放。} 慢版本可以先把节点顺序抄到数组里，按数组推导目标顺序。回到链表后，难点不再是值的顺序，而是节点身份和 next 指针不能丢。二刷时先写慢版本或口述慢版本，用它确认不会漏答案。

\textbf{特例回放。} 拿 1->2->3->4->5、k=2 手算，右移后应为 4->5->1->2->3；先连成环，再从新尾 3 处断开。规律是 k 先对长度取模，新头是第 n-k 个节点的后继，断环前必须保存新头。把这个特例继续拆开看：每个指针都要有角色，上一段尾、当前段头、当前节点和后继入口不能混。只要一次改 next 之前没有保存后继，就可能丢掉整段未处理链。复盘时把这个特例压成状态字段：虚拟头、上一段尾、当前段头、当前节点、后继节点；任何改 next 的操作之前都要保存后续入口；再用边界 空链、单节点、k 为零、k 大于长度、刚好整圈 反查初始化、更新顺序和退出条件。

\textbf{状态回放。} 手算路线：手算时画出 prev、cur、next 和下一段头节点。每次改 next 前先保存后继，这是链表推导的安全线。状态字段：虚拟头、上一段尾、当前段头、当前节点、后继节点；任何改 next 的操作之前都要保存后续入口。状态升级：把数组辅助结构换成几个稳定指针角色；重连顺序必须保证已处理链合法、未处理链仍可达。

\textbf{证明和边界。} 证明从链结构开始：已处理链连通、未处理链可达、二者连接点明确，节点没有丢失也没有成环。边界：空链、单节点、k 为零、k 大于长度、刚好整圈。变体：若本来把链表接成环，最后必须断开新尾的 next。

\noindent\textbf{LeetCode 82 删除排序链表中的重复元素 II。} 模型：有序链表中重复值连续出现，题目要求把所有重复值整组删除。推导重点：虚拟头保存结果前驱，遇到一段重复值时跳过整段，否则前驱前进。

\textbf{二刷读题。} 读题时先区分节点值、节点身份和连接关系。本题可以压成“有序链表中重复值连续出现，题目要求把所有重复值整组删除”，后续推导要围绕哪些指针会断开、哪些节点仍要可达。先遮住答案，只保留题面，复述候选对象、合法性条件和输出目标。

\textbf{暴力回放。} 慢版本可以先把节点顺序抄到数组里，按数组推导目标顺序。回到链表后，难点不再是值的顺序，而是节点身份和 next 指针不能丢。二刷时先写慢版本或口述慢版本，用它确认不会漏答案。

\textbf{特例回放。} 拿 1->2->3->3->4 手算，两个 3 都要删除，不是保留一个；pre 停在 2，cur 扫完整段重复 3 后接到 4。规律是虚拟头加 pre 指向已确认无重复前缀尾，遇到重复段时整段跳过。把这个特例继续拆开看：每个指针都要有角色，上一段尾、当前段头、当前节点和后继入口不能混。只要一次改 next 之前没有保存后继，就可能丢掉整段未处理链。复盘时把这个特例压成状态字段：虚拟头、上一段尾、当前段头、当前节点、后继节点；任何改 next 的操作之前都要保存后续入口；再用边界 头部重复、尾部重复、全部重复、没有重复 反查初始化、更新顺序和退出条件。

\textbf{状态回放。} 手算路线：手算时画出 prev、cur、next 和下一段头节点。每次改 next 前先保存后继，这是链表推导的安全线。状态字段：虚拟头、上一段尾、当前段头、当前节点、后继节点；任何改 next 的操作之前都要保存后续入口。状态升级：把数组辅助结构换成几个稳定指针角色；重连顺序必须保证已处理链合法、未处理链仍可达。

\textbf{证明和边界。} 证明从链结构开始：已处理链连通、未处理链可达、二者连接点明确，节点没有丢失也没有成环。边界：头部重复、尾部重复、全部重复、没有重复。变体：若只保留一个重复值，就是删除重复元素的较简单版本。

\noindent\textbf{LeetCode 83 删除排序链表中的重复元素。} 模型：有序链表让重复节点相邻，只需保留每段相同值的第一个节点。推导重点：当前节点和后继值相同就跳过后继，否则当前节点前进。

\textbf{二刷读题。} 读题时先区分节点值、节点身份和连接关系。本题可以压成“有序链表让重复节点相邻，只需保留每段相同值的第一个节点”，后续推导要围绕哪些指针会断开、哪些节点仍要可达。先遮住答案，只保留题面，复述候选对象、合法性条件和输出目标。

\textbf{暴力回放。} 慢版本可以先把节点顺序抄到数组里，按数组推导目标顺序。回到链表后，难点不再是值的顺序，而是节点身份和 next 指针不能丢。二刷时先写慢版本或口述慢版本，用它确认不会漏答案。

\textbf{特例回放。} 拿 1->1->2 手算，当前节点 1 的 next 值相同，就让当前节点跳过 next；遇到 2 时再前进。规律是有序链表让重复值相邻，只需比较当前节点和下一个节点，保留每组第一个。把这个特例继续拆开看：每个指针都要有角色，上一段尾、当前段头、当前节点和后继入口不能混。只要一次改 next 之前没有保存后继，就可能丢掉整段未处理链。复盘时把这个特例压成状态字段：虚拟头、上一段尾、当前段头、当前节点、后继节点；任何改 next 的操作之前都要保存后续入口；再用边界 空链、单节点、全重复、重复段在尾部 反查初始化、更新顺序和退出条件。

\textbf{状态回放。} 手算路线：手算时画出 prev、cur、next 和下一段头节点。每次改 next 前先保存后继，这是链表推导的安全线。状态字段：虚拟头、上一段尾、当前段头、当前节点、后继节点；任何改 next 的操作之前都要保存后续入口。状态升级：把数组辅助结构换成几个稳定指针角色；重连顺序必须保证已处理链合法、未处理链仍可达。

\textbf{证明和边界。} 证明从链结构开始：已处理链连通、未处理链可达、二者连接点明确，节点没有丢失也没有成环。边界：空链、单节点、全重复、重复段在尾部。变体：和删除重复元素 II 的区别是是否整段删除。

\noindent\textbf{LeetCode 92 反转链表 II。} 模型：只反转链表中一段连续区间，区间外节点必须保持原连接。推导重点：用虚拟头找到区间前驱，再用头插法把区间内部节点逐个移到前面。

\textbf{二刷读题。} 读题时先区分节点值、节点身份和连接关系。本题可以压成“只反转链表中一段连续区间，区间外节点必须保持原连接”，后续推导要围绕哪些指针会断开、哪些节点仍要可达。先遮住答案，只保留题面，复述候选对象、合法性条件和输出目标。

\textbf{暴力回放。} 慢版本可以先把节点顺序抄到数组里，按数组推导目标顺序。回到链表后，难点不再是值的顺序，而是节点身份和 next 指针不能丢。二刷时先写慢版本或口述慢版本，用它确认不会漏答案。

\textbf{特例回放。} 拿 1->2->3->4->5、left=2、right=4 手算，反转中段后是 1->4->3->2->5；反转前要保存 left 前驱和 right 后继。规律是虚拟头找到子段前驱，局部头插或三指针反转，最后接回前后两段。把这个特例继续拆开看：每个指针都要有角色，上一段尾、当前段头、当前节点和后继入口不能混。只要一次改 next 之前没有保存后继，就可能丢掉整段未处理链。复盘时把这个特例压成状态字段：虚拟头、上一段尾、当前段头、当前节点、后继节点；任何改 next 的操作之前都要保存后续入口；再用边界 从头开始反转、只反转一个节点、反转到尾部、区间边界 反查初始化、更新顺序和退出条件。

\textbf{状态回放。} 手算路线：手算时画出 prev、cur、next 和下一段头节点。每次改 next 前先保存后继，这是链表推导的安全线。状态字段：虚拟头、上一段尾、当前段头、当前节点、后继节点；任何改 next 的操作之前都要保存后续入口。状态升级：把数组辅助结构换成几个稳定指针角色；重连顺序必须保证已处理链合法、未处理链仍可达。

\textbf{证明和边界。} 证明从链结构开始：已处理链连通、未处理链可达、二者连接点明确，节点没有丢失也没有成环。边界：从头开始反转、只反转一个节点、反转到尾部、区间边界。变体：K 个一组反转是在多个固定长度区间上重复这一类局部重连。

\noindent\textbf{LeetCode 141 环形链表。} 模型：快慢指针在有环时必然相遇，无环时快指针先到空。推导重点：不用哈希表也能常数空间检测环。

\textbf{二刷读题。} 读题时先区分节点值、节点身份和连接关系。本题可以压成“快慢指针在有环时必然相遇，无环时快指针先到空”，后续推导要围绕哪些指针会断开、哪些节点仍要可达。先遮住答案，只保留题面，复述候选对象、合法性条件和输出目标。

\textbf{暴力回放。} 慢版本可以先把节点顺序抄到数组里，按数组推导目标顺序。回到链表后，难点不再是值的顺序，而是节点身份和 next 指针不能丢。二刷时先写慢版本或口述慢版本，用它确认不会漏答案。

\textbf{特例回放。} 拿 1->2->3 且 3 指回 2 手算，快指针每次走两步，慢指针每次走一步，进入环后二者距离会被不断缩小直到相遇。规律是快慢指针的相遇说明有环；无环时快指针会先到空。把这个特例继续拆开看：每个指针都要有角色，上一段尾、当前段头、当前节点和后继入口不能混。只要一次改 next 之前没有保存后继，就可能丢掉整段未处理链。复盘时把这个特例压成状态字段：虚拟头、上一段尾、当前段头、当前节点、后继节点；任何改 next 的操作之前都要保存后续入口；再用边界 空链、单节点自环、两节点环、尾部无环 反查初始化、更新顺序和退出条件。

\textbf{状态回放。} 手算路线：手算时画出 prev、cur、next 和下一段头节点。每次改 next 前先保存后继，这是链表推导的安全线。状态字段：虚拟头、上一段尾、当前段头、当前节点、后继节点；任何改 next 的操作之前都要保存后续入口。状态升级：把数组辅助结构换成几个稳定指针角色；重连顺序必须保证已处理链合法、未处理链仍可达。

\textbf{证明和边界。} 证明从链结构开始：已处理链连通、未处理链可达、二者连接点明确，节点没有丢失也没有成环。边界：空链、单节点自环、两节点环、尾部无环。变体：若要求环入口，用相遇后双指针同步走。

\noindent\textbf{LeetCode 142 环形链表 II。} 模型：相遇点到入口的距离关系来自快慢指针速度差。推导重点：相遇后一个指针回头，两个指针每次走一步，再次相遇就是入口。

\textbf{二刷读题。} 读题时先区分节点值、节点身份和连接关系。本题可以压成“相遇点到入口的距离关系来自快慢指针速度差”，后续推导要围绕哪些指针会断开、哪些节点仍要可达。先遮住答案，只保留题面，复述候选对象、合法性条件和输出目标。

\textbf{暴力回放。} 慢版本可以先把节点顺序抄到数组里，按数组推导目标顺序。回到链表后，难点不再是值的顺序，而是节点身份和 next 指针不能丢。二刷时先写慢版本或口述慢版本，用它确认不会漏答案。

\textbf{特例回放。} 拿 1->2->3->4 且 4 指回 2 手算，快慢指针在环内相遇后，把一个指针放回头节点，两个指针同速走，会在入环点 2 相遇。规律是相遇点到入环点的距离和头节点到入环点的距离同余。把这个特例继续拆开看：每个指针都要有角色，上一段尾、当前段头、当前节点和后继入口不能混。只要一次改 next 之前没有保存后继，就可能丢掉整段未处理链。复盘时把这个特例压成状态字段：虚拟头、上一段尾、当前段头、当前节点、后继节点；任何改 next 的操作之前都要保存后续入口；再用边界 入口在头节点、长尾短环、无环、单节点环 反查初始化、更新顺序和退出条件。

\textbf{状态回放。} 手算路线：手算时画出 prev、cur、next 和下一段头节点。每次改 next 前先保存后继，这是链表推导的安全线。状态字段：虚拟头、上一段尾、当前段头、当前节点、后继节点；任何改 next 的操作之前都要保存后续入口。状态升级：把数组辅助结构换成几个稳定指针角色；重连顺序必须保证已处理链合法、未处理链仍可达。

\textbf{证明和边界。} 证明从链结构开始：已处理链连通、未处理链可达、二者连接点明确，节点没有丢失也没有成环。边界：入口在头节点、长尾短环、无环、单节点环。变体：也可用哈希集合，但空间更高。

\noindent\textbf{LeetCode 146 LRU 缓存。} 模型：哈希表负责按 key 定位，双向链表负责维护新旧顺序。推导重点：访问和更新都把节点移动到头部，容量满时删除尾部最旧节点。

\textbf{二刷读题。} 读题时先区分节点值、节点身份和连接关系。本题可以压成“哈希表负责按 key 定位，双向链表负责维护新旧顺序”，后续推导要围绕哪些指针会断开、哪些节点仍要可达。先遮住答案，只保留题面，复述候选对象、合法性条件和输出目标。

\textbf{暴力回放。} 慢版本可以先把节点顺序抄到数组里，按数组推导目标顺序。回到链表后，难点不再是值的顺序，而是节点身份和 next 指针不能丢。二刷时先写慢版本或口述慢版本，用它确认不会漏答案。

\textbf{特例回放。} 拿容量 2 依次 put(1)、put(2)、get(1)、put(3) 手算。get(1) 后 1 变成最新，淘汰时应删 2。数组维护顺序每次移动成本高；链表移动节点到头 O(1)，哈希表按 key 定位节点 O(1)。规律是哈希负责找节点，双向链表负责维护新旧顺序。把这个特例继续拆开看：每个指针都要有角色，上一段尾、当前段头、当前节点和后继入口不能混。只要一次改 next 之前没有保存后继，就可能丢掉整段未处理链。复盘时把这个特例压成状态字段：虚拟头、上一段尾、当前段头、当前节点、后继节点；任何改 next 的操作之前都要保存后续入口；再用边界 容量为一、重复 put、get 不存在、更新已有值 反查初始化、更新顺序和退出条件。

\textbf{状态回放。} 手算路线：手算时画出 prev、cur、next 和下一段头节点。每次改 next 前先保存后继，这是链表推导的安全线。状态字段：虚拟头、上一段尾、当前段头、当前节点、后继节点；任何改 next 的操作之前都要保存后续入口。状态升级：把数组辅助结构换成几个稳定指针角色；重连顺序必须保证已处理链合法、未处理链仍可达。

\textbf{证明和边界。} 证明从链结构开始：已处理链连通、未处理链可达、二者连接点明确，节点没有丢失也没有成环。边界：容量为一、重复 put、get 不存在、更新已有值。变体：C++ 可用 \code{std::list} 加迭代器，也可手写双向链表理解所有权。

\noindent\textbf{LeetCode 148 排序链表。} 模型：链表不能随机访问，适合归并排序而不是快速排序数组 partition。推导重点：自底向上迭代归并能避免递归栈，同时保持 O(n log n)。

\textbf{二刷读题。} 读题时先区分节点值、节点身份和连接关系。本题可以压成“链表不能随机访问，适合归并排序而不是快速排序数组 partition”，后续推导要围绕哪些指针会断开、哪些节点仍要可达。先遮住答案，只保留题面，复述候选对象、合法性条件和输出目标。

\textbf{暴力回放。} 慢版本可以先把节点顺序抄到数组里，按数组推导目标顺序。回到链表后，难点不再是值的顺序，而是节点身份和 next 指针不能丢。二刷时先写慢版本或口述慢版本，用它确认不会漏答案。

\textbf{特例回放。} 拿 4->2->1->3 手算，快慢指针把链表切成 4->2 和 1->3，分别排好后再归并。规律是链表排序不能随机访问，归并排序只需要切断中点和有序合并，复杂度稳定在 O(n log n)。把这个特例继续拆开看：每个指针都要有角色，上一段尾、当前段头、当前节点和后继入口不能混。只要一次改 next 之前没有保存后继，就可能丢掉整段未处理链。复盘时把这个特例压成状态字段：虚拟头、上一段尾、当前段头、当前节点、后继节点；任何改 next 的操作之前都要保存后续入口；再用边界 空链、单节点、重复值、奇数长度、切断子链 反查初始化、更新顺序和退出条件。

\textbf{状态回放。} 手算路线：手算时画出 prev、cur、next 和下一段头节点。每次改 next 前先保存后继，这是链表推导的安全线。状态字段：虚拟头、上一段尾、当前段头、当前节点、后继节点；任何改 next 的操作之前都要保存后续入口。状态升级：把数组辅助结构换成几个稳定指针角色；重连顺序必须保证已处理链合法、未处理链仍可达。

\textbf{证明和边界。} 证明从链结构开始：已处理链连通、未处理链可达、二者连接点明确，节点没有丢失也没有成环。边界：空链、单节点、重复值、奇数长度、切断子链。变体：若把值复制到数组排序，会改变节点身份语义。

\noindent\textbf{LeetCode 160 相交链表。} 模型：相交是节点地址相同，不是节点值相同。推导重点：两个指针走完各自链表后切换到另一条链，走过总长度相同后相遇。

\textbf{二刷读题。} 读题时先区分节点值、节点身份和连接关系。本题可以压成“相交是节点地址相同，不是节点值相同”，后续推导要围绕哪些指针会断开、哪些节点仍要可达。先遮住答案，只保留题面，复述候选对象、合法性条件和输出目标。

\textbf{暴力回放。} 慢版本可以先把节点顺序抄到数组里，按数组推导目标顺序。回到链表后，难点不再是值的顺序，而是节点身份和 next 指针不能丢。二刷时先写慢版本或口述慢版本，用它确认不会漏答案。

\textbf{特例回放。} 拿 A: a1->a2->c1->c2 和 B: b1->c1->c2 手算，交点是节点地址 c1，不是值相等。两个指针走到尾后切换到对方头部，会在总路程相同后于 c1 相遇。规律是比较节点身份，长度差由切换链表自动抵消。把这个特例继续拆开看：每个指针都要有角色，上一段尾、当前段头、当前节点和后继入口不能混。只要一次改 next 之前没有保存后继，就可能丢掉整段未处理链。复盘时把这个特例压成状态字段：虚拟头、上一段尾、当前段头、当前节点、后继节点；任何改 next 的操作之前都要保存后续入口；再用边界 不相交、头节点相交、尾节点相交、长度差很大 反查初始化、更新顺序和退出条件。

\textbf{状态回放。} 手算路线：手算时画出 prev、cur、next 和下一段头节点。每次改 next 前先保存后继，这是链表推导的安全线。状态字段：虚拟头、上一段尾、当前段头、当前节点、后继节点；任何改 next 的操作之前都要保存后续入口。状态升级：把数组辅助结构换成几个稳定指针角色；重连顺序必须保证已处理链合法、未处理链仍可达。

\textbf{证明和边界。} 证明从链结构开始：已处理链连通、未处理链可达、二者连接点明确，节点没有丢失也没有成环。边界：不相交、头节点相交、尾节点相交、长度差很大。变体：哈希集合更直观但空间更高。

\noindent\textbf{LeetCode 206 反转链表。} 模型：每次把当前节点的 next 指向已经反转好的前缀头。推导重点：先保存后继，再改指针，再推进 previous 和 current。

\textbf{二刷读题。} 读题时先区分节点值、节点身份和连接关系。本题可以压成“每次把当前节点的 next 指向已经反转好的前缀头”，后续推导要围绕哪些指针会断开、哪些节点仍要可达。先遮住答案，只保留题面，复述候选对象、合法性条件和输出目标。

\textbf{暴力回放。} 慢版本可以先把节点顺序抄到数组里，按数组推导目标顺序。回到链表后，难点不再是值的顺序，而是节点身份和 next 指针不能丢。二刷时先写慢版本或口述慢版本，用它确认不会漏答案。

\textbf{特例回放。} 拿 1->2->3 手算，处理 1 时先保存 next=2，再让 1->null；处理 2 时保存 3，再让 2->1。规律是 prev 保存已反转前缀头，cur 保存待处理节点，每次改 next 前必须保存后继。把这个特例继续拆开看：每个指针都要有角色，上一段尾、当前段头、当前节点和后继入口不能混。只要一次改 next 之前没有保存后继，就可能丢掉整段未处理链。复盘时把这个特例压成状态字段：虚拟头、上一段尾、当前段头、当前节点、后继节点；任何改 next 的操作之前都要保存后续入口；再用边界 空链、单节点、两个节点、不要丢失剩余链 反查初始化、更新顺序和退出条件。

\textbf{状态回放。} 手算路线：手算时画出 prev、cur、next 和下一段头节点。每次改 next 前先保存后继，这是链表推导的安全线。状态字段：虚拟头、上一段尾、当前段头、当前节点、后继节点；任何改 next 的操作之前都要保存后续入口。状态升级：把数组辅助结构换成几个稳定指针角色；重连顺序必须保证已处理链合法、未处理链仍可达。

\textbf{证明和边界。} 证明从链结构开始：已处理链连通、未处理链可达、二者连接点明确，节点没有丢失也没有成环。边界：空链、单节点、两个节点、不要丢失剩余链。变体：K 组反转是在这段局部反转外增加分组边界。

\noindent\textbf{LeetCode 234 回文链表。} 模型：链表回文要求前半段和后半段逆序后逐一相等。推导重点：快慢指针找中点，反转后半段，再和前半段同步比较。

\textbf{二刷读题。} 读题时先区分节点值、节点身份和连接关系。本题可以压成“链表回文要求前半段和后半段逆序后逐一相等”，后续推导要围绕哪些指针会断开、哪些节点仍要可达。先遮住答案，只保留题面，复述候选对象、合法性条件和输出目标。

\textbf{暴力回放。} 慢版本可以先把节点顺序抄到数组里，按数组推导目标顺序。回到链表后，难点不再是值的顺序，而是节点身份和 next 指针不能丢。二刷时先写慢版本或口述慢版本，用它确认不会漏答案。

\textbf{特例回放。} 拿 1->2->2->1 手算，快慢指针找到后半段起点，再反转后半段得到 1->2，与前半段逐个比较。规律是链表不能随机访问，先找中点、反转后半、比较，再按需要恢复结构。把这个特例继续拆开看：每个指针都要有角色，上一段尾、当前段头、当前节点和后继入口不能混。只要一次改 next 之前没有保存后继，就可能丢掉整段未处理链。复盘时把这个特例压成状态字段：虚拟头、上一段尾、当前段头、当前节点、后继节点；任何改 next 的操作之前都要保存后续入口；再用边界 空链、单节点、奇数长度、偶数长度、是否恢复链表 反查初始化、更新顺序和退出条件。

\textbf{状态回放。} 手算路线：手算时画出 prev、cur、next 和下一段头节点。每次改 next 前先保存后继，这是链表推导的安全线。状态字段：虚拟头、上一段尾、当前段头、当前节点、后继节点；任何改 next 的操作之前都要保存后续入口。状态升级：把数组辅助结构换成几个稳定指针角色；重连顺序必须保证已处理链合法、未处理链仍可达。

\textbf{证明和边界。} 证明从链结构开始：已处理链连通、未处理链可达、二者连接点明确，节点没有丢失也没有成环。边界：空链、单节点、奇数长度、偶数长度、是否恢复链表。变体：如果不允许修改链表，可用栈保存前半段但空间变为线性。

\noindent\textbf{LeetCode 430 扁平化多级双向链表。} 模型：多级链表按深度优先顺序展开成一条双向链表。推导重点：显式栈保存后续 next，遇到 child 先接 child，再在子链结束后恢复原 next。

\textbf{二刷读题。} 读题时先区分节点值、节点身份和连接关系。本题可以压成“多级链表按深度优先顺序展开成一条双向链表”，后续推导要围绕哪些指针会断开、哪些节点仍要可达。先遮住答案，只保留题面，复述候选对象、合法性条件和输出目标。

\textbf{暴力回放。} 慢版本可以先把节点顺序抄到数组里，按数组推导目标顺序。回到链表后，难点不再是值的顺序，而是节点身份和 next 指针不能丢。二刷时先写慢版本或口述慢版本，用它确认不会漏答案。

\textbf{特例回放。} 拿 1-2-3 且 2 有子链 4-5 手算，扁平化顺序应为 1-2-4-5-3；若不保存 2 原来的 next=3，子链接完后会丢失 3。规律是遇到 child 时先保存 next，再把 child 接入，并把 child 尾部接回原 next。把这个特例继续拆开看：每个指针都要有角色，上一段尾、当前段头、当前节点和后继入口不能混。只要一次改 next 之前没有保存后继，就可能丢掉整段未处理链。复盘时把这个特例压成状态字段：虚拟头、上一段尾、当前段头、当前节点、后继节点；任何改 next 的操作之前都要保存后续入口；再用边界 child 为空、next 和 child 同时存在、多层嵌套、prev 指针恢复 反查初始化、更新顺序和退出条件。

\textbf{状态回放。} 手算路线：手算时画出 prev、cur、next 和下一段头节点。每次改 next 前先保存后继，这是链表推导的安全线。状态字段：虚拟头、上一段尾、当前段头、当前节点、后继节点；任何改 next 的操作之前都要保存后续入口。状态升级：把数组辅助结构换成几个稳定指针角色；重连顺序必须保证已处理链合法、未处理链仍可达。

\textbf{证明和边界。} 证明从链结构开始：已处理链连通、未处理链可达、二者连接点明确，节点没有丢失也没有成环。边界：child 为空、next 和 child 同时存在、多层嵌套、prev 指针恢复。变体：它和二叉树前序展开类似，核心是保存未处理分支。

\noindent\textbf{LeetCode 460 LFU 缓存。} 模型：缓存淘汰同时按访问频次和同频最近性排序，需要维护两个维度的索引。推导重点：哈希表按 key 定位节点，频次到双向链表保存同频节点顺序，并维护当前最小频次。

\textbf{二刷读题。} 读题时先区分节点值、节点身份和连接关系。本题可以压成“缓存淘汰同时按访问频次和同频最近性排序，需要维护两个维度的索引”，后续推导要围绕哪些指针会断开、哪些节点仍要可达。先遮住答案，只保留题面，复述候选对象、合法性条件和输出目标。

\textbf{暴力回放。} 慢版本可以先把节点顺序抄到数组里，按数组推导目标顺序。回到链表后，难点不再是值的顺序，而是节点身份和 next 指针不能丢。二刷时先写慢版本或口述慢版本，用它确认不会漏答案。

\textbf{特例回放。} 拿容量 2，put(1)、put(2)、get(1)、put(3) 手算。此时 1 的频次高于 2，应淘汰 2；若两个 key 频次相同，还要淘汰同频里最旧的。规律是 key 哈希定位节点，频次哈希定位同频链表，minFreq 指向当前最低频次。把这个特例继续拆开看：每个指针都要有角色，上一段尾、当前段头、当前节点和后继入口不能混。只要一次改 next 之前没有保存后继，就可能丢掉整段未处理链。复盘时把这个特例压成状态字段：虚拟头、上一段尾、当前段头、当前节点、后继节点；任何改 next 的操作之前都要保存后续入口；再用边界 容量为零、更新已有 key、频次链表变空时最小频次更新、同频淘汰最旧节点 反查初始化、更新顺序和退出条件。

\textbf{状态回放。} 手算路线：手算时画出 prev、cur、next 和下一段头节点。每次改 next 前先保存后继，这是链表推导的安全线。状态字段：虚拟头、上一段尾、当前段头、当前节点、后继节点；任何改 next 的操作之前都要保存后续入口。状态升级：把数组辅助结构换成几个稳定指针角色；重连顺序必须保证已处理链合法、未处理链仍可达。

\textbf{证明和边界。} 证明从链结构开始：已处理链连通、未处理链可达、二者连接点明确，节点没有丢失也没有成环。边界：容量为零、更新已有 key、频次链表变空时最小频次更新、同频淘汰最旧节点。变体：LRU 只有时间顺序一个维度，LFU 多了频次维度，结构维护明显更复杂。

\subsection{自测标准}

二刷时不要只确认自己见过这道题。请遮住答案，先讲题面模型，再讲暴力基线和瓶颈，然后说出优化结构保存了什么、不变量如何排除错误候选、边界实验如何打破错误实现。

\section{二叉树与树形状态：二刷复盘卡}
这一大节只围绕一个题型复盘，不把题号直接铺成目录。阅读时先看复盘目标，再读核心题卡，最后用自测标准检查自己是否能独立重讲。

\subsection{复盘目标}

追问通常会问返回值语义、空节点初值、全负数和极端深树。请准备把递归思想改写成显式栈或队列的口头方案。

\subsection{核心题卡}

\noindent\textbf{LeetCode 94 二叉树中序遍历。} 模型：中序访问顺序是左子树、根、右子树，BST 中会得到有序序列。推导重点：用显式栈模拟一路向左，再回退访问根并转向右子树。

\textbf{二刷读题。} 读题时先判断状态方向：父到子、子到父，还是按层传播。本题可以压成“中序访问顺序是左子树、根、右子树，BST 中会得到有序序列”，每个节点的输入和输出状态必须写清。先遮住答案，只保留题面，复述候选对象、合法性条件和输出目标。

\textbf{暴力回放。} 慢版本在每个节点重新扫描子树或路径，先求出局部答案。重复扫描暴露了真正状态：每个节点只需要从孩子或父亲那里拿到少量信息。二刷时先写慢版本或口述慢版本，用它确认不会漏答案。

\textbf{特例回放。} 拿根 2、左 1、右 3 手算，中序结果必须是 1,2,3。显式栈做法先一路压到最左，弹出 1 后访问，再回到 2。规律是栈保存尚未访问但左侧已经处理完的祖先，BST 题里这个顺序会产生递增序列。把这个特例继续拆开看：节点向父亲返回的量和全局答案通常不是同一个量。先写叶子或空节点的返回值，再写父节点如何合并左右孩子，递归式才有边界基础。复盘时把这个特例压成状态字段：节点指针、传入约束或返回值、当前答案、层号或路径状态；空节点的返回值必须和状态定义匹配；再用边界 空树、单节点、链式左树、链式右树 反查初始化、更新顺序和退出条件。

\textbf{状态回放。} 手算路线：手算时从叶子开始写返回值，或者从根开始写传入约束。不要只写访问顺序，要写每条边上传递的信息。状态字段：节点指针、传入约束或返回值、当前答案、层号或路径状态；空节点的返回值必须和状态定义匹配。状态升级：把对子树的重复扫描压成返回值或队列状态；每个节点只合并孩子给出的稳定信息。

\textbf{证明和边界。} 证明从子树归纳开始：孩子状态正确时，父节点按题目规则合并后也正确。边界：空树、单节点、链式左树、链式右树。变体：Morris 遍历可做到常数额外空间，但会临时改树结构。

\noindent\textbf{LeetCode 98 验证二叉搜索树。} 模型：BST 约束是全局上下界，不是只比较父子节点。推导重点：中序严格递增或显式传递上下界都可以；中序版本要保存前一个访问值。

\textbf{二刷读题。} 读题时先判断状态方向：父到子、子到父，还是按层传播。本题可以压成“BST 约束是全局上下界，不是只比较父子节点”，每个节点的输入和输出状态必须写清。先遮住答案，只保留题面，复述候选对象、合法性条件和输出目标。

\textbf{暴力回放。} 慢版本在每个节点重新扫描子树或路径，先求出局部答案。重复扫描暴露了真正状态：每个节点只需要从孩子或父亲那里拿到少量信息。二刷时先写慢版本或口述慢版本，用它确认不会漏答案。

\textbf{特例回放。} 拿 [5,1,4,null,null,3,6] 手算，节点 3 小于根 5 却在右子树里，只比较父子 4 和 3 会误判。规律是 BST 约束来自整条祖先上下界，或中序序列必须严格递增；重复值和 int 边界要按题意单独处理。把这个特例继续拆开看：节点向父亲返回的量和全局答案通常不是同一个量。先写叶子或空节点的返回值，再写父节点如何合并左右孩子，递归式才有边界基础。复盘时把这个特例压成状态字段：节点指针、传入约束或返回值、当前答案、层号或路径状态；空节点的返回值必须和状态定义匹配；再用边界 重复值、int 最小最大值、局部合法但全局非法 反查初始化、更新顺序和退出条件。

\textbf{状态回放。} 手算路线：手算时从叶子开始写返回值，或者从根开始写传入约束。不要只写访问顺序，要写每条边上传递的信息。状态字段：节点指针、传入约束或返回值、当前答案、层号或路径状态；空节点的返回值必须和状态定义匹配。状态升级：把对子树的重复扫描压成返回值或队列状态；每个节点只合并孩子给出的稳定信息。

\textbf{证明和边界。} 证明从子树归纳开始：孩子状态正确时，父节点按题目规则合并后也正确。边界：重复值、int 最小最大值、局部合法但全局非法。变体：若允许重复值，需要明确重复放左还是放右。

\noindent\textbf{LeetCode 102 二叉树层序遍历。} 模型：队列在每轮开始时保存当前层所有节点。推导重点：记录 level size 固定层边界，避免下一层节点混入当前层。

\textbf{二刷读题。} 读题时先判断状态方向：父到子、子到父，还是按层传播。本题可以压成“队列在每轮开始时保存当前层所有节点”，每个节点的输入和输出状态必须写清。先遮住答案，只保留题面，复述候选对象、合法性条件和输出目标。

\textbf{暴力回放。} 慢版本在每个节点重新扫描子树或路径，先求出局部答案。重复扫描暴露了真正状态：每个节点只需要从孩子或父亲那里拿到少量信息。二刷时先写慢版本或口述慢版本，用它确认不会漏答案。

\textbf{特例回放。} 拿根 3、左 9、右 20 手算，队列第一轮只有 3，第二轮才是 9 和 20。若不固定本层 size，新入队的下一层节点会混进当前层。规律是每轮开始记录队列长度，这个长度就是当前层边界。把这个特例继续拆开看：节点向父亲返回的量和全局答案通常不是同一个量。先写叶子或空节点的返回值，再写父节点如何合并左右孩子，递归式才有边界基础。复盘时把这个特例压成状态字段：节点指针、传入约束或返回值、当前答案、层号或路径状态；空节点的返回值必须和状态定义匹配；再用边界 空树、只有根、最后一层不满、左右顺序 反查初始化、更新顺序和退出条件。

\textbf{状态回放。} 手算路线：手算时从叶子开始写返回值，或者从根开始写传入约束。不要只写访问顺序，要写每条边上传递的信息。状态字段：节点指针、传入约束或返回值、当前答案、层号或路径状态；空节点的返回值必须和状态定义匹配。状态升级：把对子树的重复扫描压成返回值或队列状态；每个节点只合并孩子给出的稳定信息。

\textbf{证明和边界。} 证明从子树归纳开始：孩子状态正确时，父节点按题目规则合并后也正确。边界：空树、只有根、最后一层不满、左右顺序。变体：锯齿形层序只是在每层输出顺序上增加方向控制。

\noindent\textbf{LeetCode 103 二叉树锯齿形层序遍历。} 模型：BFS 层边界不变，变化的是当前层写入结果的方向。推导重点：可以层内反转，也可以预分配数组按方向写入目标下标。

\textbf{二刷读题。} 读题时先判断状态方向：父到子、子到父，还是按层传播。本题可以压成“BFS 层边界不变，变化的是当前层写入结果的方向”，每个节点的输入和输出状态必须写清。先遮住答案，只保留题面，复述候选对象、合法性条件和输出目标。

\textbf{暴力回放。} 慢版本在每个节点重新扫描子树或路径，先求出局部答案。重复扫描暴露了真正状态：每个节点只需要从孩子或父亲那里拿到少量信息。二刷时先写慢版本或口述慢版本，用它确认不会漏答案。

\textbf{特例回放。} 拿三层树手算，第一层从左到右，第二层从右到左，第三层再从左到右；BFS 层边界不变，只是写入结果的位置反向。规律是层序遍历负责分层，方向标志只影响当前层数组的填充顺序。把这个特例继续拆开看：节点向父亲返回的量和全局答案通常不是同一个量。先写叶子或空节点的返回值，再写父节点如何合并左右孩子，递归式才有边界基础。复盘时把这个特例压成状态字段：节点指针、传入约束或返回值、当前答案、层号或路径状态；空节点的返回值必须和状态定义匹配；再用边界 单层、偶数层、奇数层、空树 反查初始化、更新顺序和退出条件。

\textbf{状态回放。} 手算路线：手算时从叶子开始写返回值，或者从根开始写传入约束。不要只写访问顺序，要写每条边上传递的信息。状态字段：节点指针、传入约束或返回值、当前答案、层号或路径状态；空节点的返回值必须和状态定义匹配。状态升级：把对子树的重复扫描压成返回值或队列状态；每个节点只合并孩子给出的稳定信息。

\textbf{证明和边界。} 证明从子树归纳开始：孩子状态正确时，父节点按题目规则合并后也正确。边界：单层、偶数层、奇数层、空树。变体：若要求从底向上输出，收集后反转层列表即可。

\noindent\textbf{LeetCode 104 二叉树最大深度。} 模型：深度可以看成层数，也可以作为栈中携带的状态。推导重点：显式栈保存节点和深度，避免极端链式树递归栈过深。

\textbf{二刷读题。} 读题时先判断状态方向：父到子、子到父，还是按层传播。本题可以压成“深度可以看成层数，也可以作为栈中携带的状态”，每个节点的输入和输出状态必须写清。先遮住答案，只保留题面，复述候选对象、合法性条件和输出目标。

\textbf{暴力回放。} 慢版本在每个节点重新扫描子树或路径，先求出局部答案。重复扫描暴露了真正状态：每个节点只需要从孩子或父亲那里拿到少量信息。二刷时先写慢版本或口述慢版本，用它确认不会漏答案。

\textbf{特例回放。} 拿链式树 1->2->3 手算，最大深度是 3；空树是 0。递归返回左右子树深度的较大值加一，显式栈则把节点和当前深度一起压栈。规律是状态必须携带深度，不能只记录访问过哪些节点。把这个特例继续拆开看：节点向父亲返回的量和全局答案通常不是同一个量。先写叶子或空节点的返回值，再写父节点如何合并左右孩子，递归式才有边界基础。复盘时把这个特例压成状态字段：节点指针、传入约束或返回值、当前答案、层号或路径状态；空节点的返回值必须和状态定义匹配；再用边界 空树返回零、根深度是一、链式树 反查初始化、更新顺序和退出条件。

\textbf{状态回放。} 手算路线：手算时从叶子开始写返回值，或者从根开始写传入约束。不要只写访问顺序，要写每条边上传递的信息。状态字段：节点指针、传入约束或返回值、当前答案、层号或路径状态；空节点的返回值必须和状态定义匹配。状态升级：把对子树的重复扫描压成返回值或队列状态；每个节点只合并孩子给出的稳定信息。

\textbf{证明和边界。} 证明从子树归纳开始：孩子状态正确时，父节点按题目规则合并后也正确。边界：空树返回零、根深度是一、链式树。变体：最小深度不能简单取左右最小，因为空子树不构成叶子路径。

\noindent\textbf{LeetCode 105 前序中序构造二叉树。} 模型：前序给根，中序把左右子树切开。推导重点：用哈希表记录中序下标，避免每次线性寻找根位置。

\textbf{二刷读题。} 读题时先判断状态方向：父到子、子到父，还是按层传播。本题可以压成“前序给根，中序把左右子树切开”，每个节点的输入和输出状态必须写清。先遮住答案，只保留题面，复述候选对象、合法性条件和输出目标。

\textbf{暴力回放。} 慢版本在每个节点重新扫描子树或路径，先求出局部答案。重复扫描暴露了真正状态：每个节点只需要从孩子或父亲那里拿到少量信息。二刷时先写慢版本或口述慢版本，用它确认不会漏答案。

\textbf{特例回放。} 拿 preorder=[3,9,20,15,7]、inorder=[9,3,15,20,7] 手算，前序第一个 3 是根，中序中 3 左边是左子树，右边是右子树。规律是用哈希表快速找根在中序的位置，再用左右子树长度切分前序区间。把这个特例继续拆开看：节点向父亲返回的量和全局答案通常不是同一个量。先写叶子或空节点的返回值，再写父节点如何合并左右孩子，递归式才有边界基础。复盘时把这个特例压成状态字段：节点指针、传入约束或返回值、当前答案、层号或路径状态；空节点的返回值必须和状态定义匹配；再用边界 节点值唯一、空子树、左右子树长度计算、输入不合法 反查初始化、更新顺序和退出条件。

\textbf{状态回放。} 手算路线：手算时从叶子开始写返回值，或者从根开始写传入约束。不要只写访问顺序，要写每条边上传递的信息。状态字段：节点指针、传入约束或返回值、当前答案、层号或路径状态；空节点的返回值必须和状态定义匹配。状态升级：把对子树的重复扫描压成返回值或队列状态；每个节点只合并孩子给出的稳定信息。

\textbf{证明和边界。} 证明从子树归纳开始：孩子状态正确时，父节点按题目规则合并后也正确。边界：节点值唯一、空子树、左右子树长度计算、输入不合法。变体：后序中序构造类似，只是根来自后序末尾。

\noindent\textbf{LeetCode 106 中序后序构造二叉树。} 模型：后序末尾是根，中序位置确定左右子树规模。推导重点：构造时要小心后序左右区间边界，避免切分错位。

\textbf{二刷读题。} 读题时先判断状态方向：父到子、子到父，还是按层传播。本题可以压成“后序末尾是根，中序位置确定左右子树规模”，每个节点的输入和输出状态必须写清。先遮住答案，只保留题面，复述候选对象、合法性条件和输出目标。

\textbf{暴力回放。} 慢版本在每个节点重新扫描子树或路径，先求出局部答案。重复扫描暴露了真正状态：每个节点只需要从孩子或父亲那里拿到少量信息。二刷时先写慢版本或口述慢版本，用它确认不会漏答案。

\textbf{特例回放。} 拿 inorder=[9,3,15,20,7]、postorder=[9,15,7,20,3] 手算，后序最后一个 3 是根，中序把左右子树切开。规律是根来自后序末尾，左右区间长度必须同步更新；全左链和全右链最容易暴露边界错位。把这个特例继续拆开看：节点向父亲返回的量和全局答案通常不是同一个量。先写叶子或空节点的返回值，再写父节点如何合并左右孩子，递归式才有边界基础。复盘时把这个特例压成状态字段：节点指针、传入约束或返回值、当前答案、层号或路径状态；空节点的返回值必须和状态定义匹配；再用边界 空树、单节点、全左链、全右链 反查初始化、更新顺序和退出条件。

\textbf{状态回放。} 手算路线：手算时从叶子开始写返回值，或者从根开始写传入约束。不要只写访问顺序，要写每条边上传递的信息。状态字段：节点指针、传入约束或返回值、当前答案、层号或路径状态；空节点的返回值必须和状态定义匹配。状态升级：把对子树的重复扫描压成返回值或队列状态；每个节点只合并孩子给出的稳定信息。

\textbf{证明和边界。} 证明从子树归纳开始：孩子状态正确时，父节点按题目规则合并后也正确。边界：空树、单节点、全左链、全右链。变体：若给前序和后序，二叉树通常不能唯一确定，除非额外限制。

\noindent\textbf{LeetCode 108 有序数组转二叉搜索树。} 模型：有序数组中点作为根能让左右规模接近。推导重点：每次选择区间中点，左右子区间构造左右子树。

\textbf{二刷读题。} 读题时先判断状态方向：父到子、子到父，还是按层传播。本题可以压成“有序数组中点作为根能让左右规模接近”，每个节点的输入和输出状态必须写清。先遮住答案，只保留题面，复述候选对象、合法性条件和输出目标。

\textbf{暴力回放。} 慢版本在每个节点重新扫描子树或路径，先求出局部答案。重复扫描暴露了真正状态：每个节点只需要从孩子或父亲那里拿到少量信息。二刷时先写慢版本或口述慢版本，用它确认不会漏答案。

\textbf{特例回放。} 拿 [-10,-3,0,5,9] 手算，选 0 做根，左边两个数构成左子树，右边两个数构成右子树，高度接近平衡。规律是有序数组的中点天然把值域分成 BST 左右两半，空区间返回空节点。把这个特例继续拆开看：节点向父亲返回的量和全局答案通常不是同一个量。先写叶子或空节点的返回值，再写父节点如何合并左右孩子，递归式才有边界基础。复盘时把这个特例压成状态字段：节点指针、传入约束或返回值、当前答案、层号或路径状态；空节点的返回值必须和状态定义匹配；再用边界 空区间、偶数长度选左中还是右中、高度平衡定义 反查初始化、更新顺序和退出条件。

\textbf{状态回放。} 手算路线：手算时从叶子开始写返回值，或者从根开始写传入约束。不要只写访问顺序，要写每条边上传递的信息。状态字段：节点指针、传入约束或返回值、当前答案、层号或路径状态；空节点的返回值必须和状态定义匹配。状态升级：把对子树的重复扫描压成返回值或队列状态；每个节点只合并孩子给出的稳定信息。

\textbf{证明和边界。} 证明从子树归纳开始：孩子状态正确时，父节点按题目规则合并后也正确。边界：空区间、偶数长度选左中还是右中、高度平衡定义。变体：工程上递归可读，若严格避免递归可用队列保存区间和父指针。

\noindent\textbf{LeetCode 110 平衡二叉树。} 模型：每个节点需要知道子树高度以及是否已经失衡。推导重点：后序汇总时一旦子树失衡就向上返回失败标记，避免重复算高度。

\textbf{二刷读题。} 读题时先判断状态方向：父到子、子到父，还是按层传播。本题可以压成“每个节点需要知道子树高度以及是否已经失衡”，每个节点的输入和输出状态必须写清。先遮住答案，只保留题面，复述候选对象、合法性条件和输出目标。

\textbf{暴力回放。} 慢版本在每个节点重新扫描子树或路径，先求出局部答案。重复扫描暴露了真正状态：每个节点只需要从孩子或父亲那里拿到少量信息。二刷时先写慢版本或口述慢版本，用它确认不会漏答案。

\textbf{特例回放。} 拿链式树 1->2->3 手算，根的左右高度差为 2，已经不平衡；若子树已经失衡，父节点没必要继续精算高度。规律是后序返回高度或失败标记，空树高度为 0，左右差正好 1 仍合法。把这个特例继续拆开看：节点向父亲返回的量和全局答案通常不是同一个量。先写叶子或空节点的返回值，再写父节点如何合并左右孩子，递归式才有边界基础。复盘时把这个特例压成状态字段：节点指针、传入约束或返回值、当前答案、层号或路径状态；空节点的返回值必须和状态定义匹配；再用边界 空树、单节点、左右高度差正好一、链式树 反查初始化、更新顺序和退出条件。

\textbf{状态回放。} 手算路线：手算时从叶子开始写返回值，或者从根开始写传入约束。不要只写访问顺序，要写每条边上传递的信息。状态字段：节点指针、传入约束或返回值、当前答案、层号或路径状态；空节点的返回值必须和状态定义匹配。状态升级：把对子树的重复扫描压成返回值或队列状态；每个节点只合并孩子给出的稳定信息。

\textbf{证明和边界。} 证明从子树归纳开始：孩子状态正确时，父节点按题目规则合并后也正确。边界：空树、单节点、左右高度差正好一、链式树。变体：可以把返回值设计成 optional 高度，空表示不平衡。

\noindent\textbf{LeetCode 111 二叉树最小深度。} 模型：最小深度是到最近叶子的节点数，不是左右深度简单取最小。推导重点：BFS 第一次遇到叶子就是最小深度，因为按层扩展。

\textbf{二刷读题。} 读题时先判断状态方向：父到子、子到父，还是按层传播。本题可以压成“最小深度是到最近叶子的节点数，不是左右深度简单取最小”，每个节点的输入和输出状态必须写清。先遮住答案，只保留题面，复述候选对象、合法性条件和输出目标。

\textbf{暴力回放。} 慢版本在每个节点重新扫描子树或路径，先求出局部答案。重复扫描暴露了真正状态：每个节点只需要从孩子或父亲那里拿到少量信息。二刷时先写慢版本或口述慢版本，用它确认不会漏答案。

\textbf{特例回放。} 拿根 1 只有左孩子 2 手算，最小深度是 2，不是 min(左深度1,右深度0)+1。规律是空子树不能当作叶子路径；BFS 第一次遇到叶子最直接，因为层次递增。把这个特例继续拆开看：节点向父亲返回的量和全局答案通常不是同一个量。先写叶子或空节点的返回值，再写父节点如何合并左右孩子，递归式才有边界基础。复盘时把这个特例压成状态字段：节点指针、传入约束或返回值、当前答案、层号或路径状态；空节点的返回值必须和状态定义匹配；再用边界 只有左子树、只有右子树、根为空、根是叶子 反查初始化、更新顺序和退出条件。

\textbf{状态回放。} 手算路线：手算时从叶子开始写返回值，或者从根开始写传入约束。不要只写访问顺序，要写每条边上传递的信息。状态字段：节点指针、传入约束或返回值、当前答案、层号或路径状态；空节点的返回值必须和状态定义匹配。状态升级：把对子树的重复扫描压成返回值或队列状态；每个节点只合并孩子给出的稳定信息。

\textbf{证明和边界。} 证明从子树归纳开始：孩子状态正确时，父节点按题目规则合并后也正确。边界：只有左子树、只有右子树、根为空、根是叶子。变体：DFS 写法必须特殊处理空子树不能作为最小路径。

\noindent\textbf{LeetCode 112 路径总和。} 模型：路径必须从根到叶子，状态是剩余目标值。推导重点：沿路径减去当前节点值，到叶子时检查剩余是否等于叶子值。

\textbf{二刷读题。} 读题时先判断状态方向：父到子、子到父，还是按层传播。本题可以压成“路径必须从根到叶子，状态是剩余目标值”，每个节点的输入和输出状态必须写清。先遮住答案，只保留题面，复述候选对象、合法性条件和输出目标。

\textbf{暴力回放。} 慢版本在每个节点重新扫描子树或路径，先求出局部答案。重复扫描暴露了真正状态：每个节点只需要从孩子或父亲那里拿到少量信息。二刷时先写慢版本或口述慢版本，用它确认不会漏答案。

\textbf{特例回放。} 拿路径 5->4->11、target=20 手算，沿途把目标减成 15、11、0，到叶子时正好为 0 才成功；非叶子提前等于目标不算。规律是状态是剩余目标值，终止条件必须同时满足叶子节点。把这个特例继续拆开看：节点向父亲返回的量和全局答案通常不是同一个量。先写叶子或空节点的返回值，再写父节点如何合并左右孩子，递归式才有边界基础。复盘时把这个特例压成状态字段：节点指针、传入约束或返回值、当前答案、层号或路径状态；空节点的返回值必须和状态定义匹配；再用边界 负数节点、目标为零、非叶子提前达到目标、空树 反查初始化、更新顺序和退出条件。

\textbf{状态回放。} 手算路线：手算时从叶子开始写返回值，或者从根开始写传入约束。不要只写访问顺序，要写每条边上传递的信息。状态字段：节点指针、传入约束或返回值、当前答案、层号或路径状态；空节点的返回值必须和状态定义匹配。状态升级：把对子树的重复扫描压成返回值或队列状态；每个节点只合并孩子给出的稳定信息。

\textbf{证明和边界。} 证明从子树归纳开始：孩子状态正确时，父节点按题目规则合并后也正确。边界：负数节点、目标为零、非叶子提前达到目标、空树。变体：若要求列出所有路径，需要保存路径数组并在回退时弹出。

\noindent\textbf{LeetCode 113 路径总和 II。} 模型：相比判定题，状态还要保存当前路径。推导重点：到叶子且和满足时复制路径；回溯时恢复路径。

\textbf{二刷读题。} 读题时先判断状态方向：父到子、子到父，还是按层传播。本题可以压成“相比判定题，状态还要保存当前路径”，每个节点的输入和输出状态必须写清。先遮住答案，只保留题面，复述候选对象、合法性条件和输出目标。

\textbf{暴力回放。} 慢版本在每个节点重新扫描子树或路径，先求出局部答案。重复扫描暴露了真正状态：每个节点只需要从孩子或父亲那里拿到少量信息。二刷时先写慢版本或口述慢版本，用它确认不会漏答案。

\textbf{特例回放。} 拿 target=22 和路径 5->4->11->2 手算，到叶子时剩余为 0，需要复制当前路径；回退后要弹出 2，继续找别的分支。规律是判定题多了路径数组，保存答案必须复制，不能保存同一个可变容器引用。把这个特例继续拆开看：节点向父亲返回的量和全局答案通常不是同一个量。先写叶子或空节点的返回值，再写父节点如何合并左右孩子，递归式才有边界基础。复盘时把这个特例压成状态字段：节点指针、传入约束或返回值、当前答案、层号或路径状态；空节点的返回值必须和状态定义匹配；再用边界 多个答案、负数、空树、路径数组复制时机 反查初始化、更新顺序和退出条件。

\textbf{状态回放。} 手算路线：手算时从叶子开始写返回值，或者从根开始写传入约束。不要只写访问顺序，要写每条边上传递的信息。状态字段：节点指针、传入约束或返回值、当前答案、层号或路径状态；空节点的返回值必须和状态定义匹配。状态升级：把对子树的重复扫描压成返回值或队列状态；每个节点只合并孩子给出的稳定信息。

\textbf{证明和边界。} 证明从子树归纳开始：孩子状态正确时，父节点按题目规则合并后也正确。边界：多个答案、负数、空树、路径数组复制时机。变体：若路径不要求从根到叶，模型会变成前缀和加哈希。

\noindent\textbf{LeetCode 114 二叉树展开为链表。} 模型：目标是按前序顺序把树原地改成右指针链。推导重点：显式栈按右后左先入顺序处理，逐个接到前驱右侧。

\textbf{二刷读题。} 读题时先判断状态方向：父到子、子到父，还是按层传播。本题可以压成“目标是按前序顺序把树原地改成右指针链”，每个节点的输入和输出状态必须写清。先遮住答案，只保留题面，复述候选对象、合法性条件和输出目标。

\textbf{暴力回放。} 慢版本在每个节点重新扫描子树或路径，先求出局部答案。重复扫描暴露了真正状态：每个节点只需要从孩子或父亲那里拿到少量信息。二刷时先写慢版本或口述慢版本，用它确认不会漏答案。

\textbf{特例回放。} 拿 1 的左子树 2 和右子树 5 手算，前序链应该是 1->2->...->5；若直接把左子树接到右边而不保存原右子树，5 会丢失。规律是按前序顺序重连，左指针置空，原右子树必须接到左子树展开后的尾部。把这个特例继续拆开看：节点向父亲返回的量和全局答案通常不是同一个量。先写叶子或空节点的返回值，再写父节点如何合并左右孩子，递归式才有边界基础。复盘时把这个特例压成状态字段：节点指针、传入约束或返回值、当前答案、层号或路径状态；空节点的返回值必须和状态定义匹配；再用边界 空树、单节点、原有右子树不能丢、左指针置空 反查初始化、更新顺序和退出条件。

\textbf{状态回放。} 手算路线：手算时从叶子开始写返回值，或者从根开始写传入约束。不要只写访问顺序，要写每条边上传递的信息。状态字段：节点指针、传入约束或返回值、当前答案、层号或路径状态；空节点的返回值必须和状态定义匹配。状态升级：把对子树的重复扫描压成返回值或队列状态；每个节点只合并孩子给出的稳定信息。

\textbf{证明和边界。} 证明从子树归纳开始：孩子状态正确时，父节点按题目规则合并后也正确。边界：空树、单节点、原有右子树不能丢、左指针置空。变体：也可用寻找左子树最右节点的原地方法。

\noindent\textbf{LeetCode 124 二叉树最大路径和。} 模型：路径可以从任意节点到任意节点，但不能分叉向父节点贡献两边。推导重点：每个节点向父亲贡献单边最大增益，全局答案可用左右两边加当前节点更新。

\textbf{二刷读题。} 读题时先判断状态方向：父到子、子到父，还是按层传播。本题可以压成“路径可以从任意节点到任意节点，但不能分叉向父节点贡献两边”，每个节点的输入和输出状态必须写清。先遮住答案，只保留题面，复述候选对象、合法性条件和输出目标。

\textbf{暴力回放。} 慢版本在每个节点重新扫描子树或路径，先求出局部答案。重复扫描暴露了真正状态：每个节点只需要从孩子或父亲那里拿到少量信息。二刷时先写慢版本或口述慢版本，用它确认不会漏答案。

\textbf{特例回放。} 拿根 -10、左 9、右 20 手算，经过 20 的路径可以同时吃左增益 15 和右增益 7 更新全局答案，但向父亲只能贡献一条单边路径。规律是返回给父亲的是单边最大增益，全局答案才允许左右两边加当前节点。把这个特例继续拆开看：节点向父亲返回的量和全局答案通常不是同一个量。先写叶子或空节点的返回值，再写父节点如何合并左右孩子，递归式才有边界基础。复盘时把这个特例压成状态字段：节点指针、传入约束或返回值、当前答案、层号或路径状态；空节点的返回值必须和状态定义匹配；再用边界 全负数、单节点、左右增益为负时取零、答案初始值 反查初始化、更新顺序和退出条件。

\textbf{状态回放。} 手算路线：手算时从叶子开始写返回值，或者从根开始写传入约束。不要只写访问顺序，要写每条边上传递的信息。状态字段：节点指针、传入约束或返回值、当前答案、层号或路径状态；空节点的返回值必须和状态定义匹配。状态升级：把对子树的重复扫描压成返回值或队列状态；每个节点只合并孩子给出的稳定信息。

\textbf{证明和边界。} 证明从子树归纳开始：孩子状态正确时，父节点按题目规则合并后也正确。边界：全负数、单节点、左右增益为负时取零、答案初始值。变体：若路径必须从根到叶，状态和答案位置完全不同。

\noindent\textbf{LeetCode 199 二叉树右视图。} 模型：每层从右侧能看到的就是该层最后访问或最先访问的右侧节点。推导重点：BFS 固定层边界，取每层最后一个；或 DFS 按右优先首次到达深度。

\textbf{二刷读题。} 读题时先判断状态方向：父到子、子到父，还是按层传播。本题可以压成“每层从右侧能看到的就是该层最后访问或最先访问的右侧节点”，每个节点的输入和输出状态必须写清。先遮住答案，只保留题面，复述候选对象、合法性条件和输出目标。

\textbf{暴力回放。} 慢版本在每个节点重新扫描子树或路径，先求出局部答案。重复扫描暴露了真正状态：每个节点只需要从孩子或父亲那里拿到少量信息。二刷时先写慢版本或口述慢版本，用它确认不会漏答案。

\textbf{特例回放。} 拿根 1、左 2、右 3 手算，从右侧看第一层是 1，第二层是 3；若右边缺失，才会看到左侧节点。规律是 BFS 每层最后一个节点，或 DFS 优先右子树并第一次到达某层时记录。把这个特例继续拆开看：节点向父亲返回的量和全局答案通常不是同一个量。先写叶子或空节点的返回值，再写父节点如何合并左右孩子，递归式才有边界基础。复盘时把这个特例压成状态字段：节点指针、传入约束或返回值、当前答案、层号或路径状态；空节点的返回值必须和状态定义匹配；再用边界 空树、左链、右链、左右交错 反查初始化、更新顺序和退出条件。

\textbf{状态回放。} 手算路线：手算时从叶子开始写返回值，或者从根开始写传入约束。不要只写访问顺序，要写每条边上传递的信息。状态字段：节点指针、传入约束或返回值、当前答案、层号或路径状态；空节点的返回值必须和状态定义匹配。状态升级：把对子树的重复扫描压成返回值或队列状态；每个节点只合并孩子给出的稳定信息。

\textbf{证明和边界。} 证明从子树归纳开始：孩子状态正确时，父节点按题目规则合并后也正确。边界：空树、左链、右链、左右交错。变体：若求左视图，只改变层内取值方向。

\noindent\textbf{LeetCode 208 实现 Trie。} 模型：前缀树把字符串公共前缀共享成路径。推导重点：插入和查询都逐字符沿边移动，不存在的边按需要创建。

\textbf{二刷读题。} 读题时先判断状态方向：父到子、子到父，还是按层传播。本题可以压成“前缀树把字符串公共前缀共享成路径”，每个节点的输入和输出状态必须写清。先遮住答案，只保留题面，复述候选对象、合法性条件和输出目标。

\textbf{暴力回放。} 慢版本在每个节点重新扫描子树或路径，先求出局部答案。重复扫描暴露了真正状态：每个节点只需要从孩子或父亲那里拿到少量信息。二刷时先写慢版本或口述慢版本，用它确认不会漏答案。

\textbf{特例回放。} 拿插入 apple 再查 app 手算，app 的路径存在，但如果 p 节点没有终止标记，search(app) 应为假，startsWith(app) 才为真。规律是 Trie 节点既要有孩子边，也要有单词终止标记。把这个特例继续拆开看：节点向父亲返回的量和全局答案通常不是同一个量。先写叶子或空节点的返回值，再写父节点如何合并左右孩子，递归式才有边界基础。复盘时把这个特例压成状态字段：节点指针、传入约束或返回值、当前答案、层号或路径状态；空节点的返回值必须和状态定义匹配；再用边界 空字符串约定、单词结束标记、前缀存在不等于单词存在 反查初始化、更新顺序和退出条件。

\textbf{状态回放。} 手算路线：手算时从叶子开始写返回值，或者从根开始写传入约束。不要只写访问顺序，要写每条边上传递的信息。状态字段：节点指针、传入约束或返回值、当前答案、层号或路径状态；空节点的返回值必须和状态定义匹配。状态升级：把对子树的重复扫描压成返回值或队列状态；每个节点只合并孩子给出的稳定信息。

\textbf{证明和边界。} 证明从子树归纳开始：孩子状态正确时，父节点按题目规则合并后也正确。边界：空字符串约定、单词结束标记、前缀存在不等于单词存在。变体：若字符集很大，用哈希子节点；小写字母用数组更快。

\noindent\textbf{LeetCode 211 添加与搜索单词。} 模型：点号通配符让查询可能分叉，但前缀树仍能剪掉不匹配前缀。推导重点：普通字符沿唯一边走，点号枚举当前节点所有孩子。

\textbf{二刷读题。} 读题时先判断状态方向：父到子、子到父，还是按层传播。本题可以压成“点号通配符让查询可能分叉，但前缀树仍能剪掉不匹配前缀”，每个节点的输入和输出状态必须写清。先遮住答案，只保留题面，复述候选对象、合法性条件和输出目标。

\textbf{暴力回放。} 慢版本在每个节点重新扫描子树或路径，先求出局部答案。重复扫描暴露了真正状态：每个节点只需要从孩子或父亲那里拿到少量信息。二刷时先写慢版本或口述慢版本，用它确认不会漏答案。

\textbf{特例回放。} 拿 add bad、dad、mad 后查 .ad 手算，点号可以走 b、d、m 三条边；查 b.. 时后两个点继续在子树里展开。规律是普通字符沿唯一边走，点号触发 DFS 枚举孩子，终点仍要检查单词结束标记。把这个特例继续拆开看：节点向父亲返回的量和全局答案通常不是同一个量。先写叶子或空节点的返回值，再写父节点如何合并左右孩子，递归式才有边界基础。复盘时把这个特例压成状态字段：节点指针、传入约束或返回值、当前答案、层号或路径状态；空节点的返回值必须和状态定义匹配；再用边界 连续通配符、空结果、单词结束标记、字符集大小 反查初始化、更新顺序和退出条件。

\textbf{状态回放。} 手算路线：手算时从叶子开始写返回值，或者从根开始写传入约束。不要只写访问顺序，要写每条边上传递的信息。状态字段：节点指针、传入约束或返回值、当前答案、层号或路径状态；空节点的返回值必须和状态定义匹配。状态升级：把对子树的重复扫描压成返回值或队列状态；每个节点只合并孩子给出的稳定信息。

\textbf{证明和边界。} 证明从子树归纳开始：孩子状态正确时，父节点按题目规则合并后也正确。边界：连续通配符、空结果、单词结束标记、字符集大小。变体：大量通配符会接近遍历整棵 Trie，需要规模约束。

\noindent\textbf{LeetCode 226 翻转二叉树。} 模型：每个节点交换左右孩子即可，遍历顺序不影响最终结构。推导重点：显式队列或栈逐个节点交换，避免递归深度。

\textbf{二刷读题。} 读题时先判断状态方向：父到子、子到父，还是按层传播。本题可以压成“每个节点交换左右孩子即可，遍历顺序不影响最终结构”，每个节点的输入和输出状态必须写清。先遮住答案，只保留题面，复述候选对象、合法性条件和输出目标。

\textbf{暴力回放。} 慢版本在每个节点重新扫描子树或路径，先求出局部答案。重复扫描暴露了真正状态：每个节点只需要从孩子或父亲那里拿到少量信息。二刷时先写慢版本或口述慢版本，用它确认不会漏答案。

\textbf{特例回放。} 拿根 4、左 2、右 7 手算，翻转后左右孩子交换，2 和 7 的子树也要递归交换。规律是每个节点只负责交换自己的左右指针，子树正确性由递归返回；空节点直接返回。把这个特例继续拆开看：节点向父亲返回的量和全局答案通常不是同一个量。先写叶子或空节点的返回值，再写父节点如何合并左右孩子，递归式才有边界基础。复盘时把这个特例压成状态字段：节点指针、传入约束或返回值、当前答案、层号或路径状态；空节点的返回值必须和状态定义匹配；再用边界 空树、单节点、只有一侧子树、重复值 反查初始化、更新顺序和退出条件。

\textbf{状态回放。} 手算路线：手算时从叶子开始写返回值，或者从根开始写传入约束。不要只写访问顺序，要写每条边上传递的信息。状态字段：节点指针、传入约束或返回值、当前答案、层号或路径状态；空节点的返回值必须和状态定义匹配。状态升级：把对子树的重复扫描压成返回值或队列状态；每个节点只合并孩子给出的稳定信息。

\textbf{证明和边界。} 证明从子树归纳开始：孩子状态正确时，父节点按题目规则合并后也正确。边界：空树、单节点、只有一侧子树、重复值。变体：这题简单但适合训练树节点指针修改。

\noindent\textbf{LeetCode 230 二叉搜索树中第 K 小。} 模型：BST 中序遍历按升序访问节点。推导重点：显式栈中序，访问第 k 个节点时返回。

\textbf{二刷读题。} 读题时先判断状态方向：父到子、子到父，还是按层传播。本题可以压成“BST 中序遍历按升序访问节点”，每个节点的输入和输出状态必须写清。先遮住答案，只保留题面，复述候选对象、合法性条件和输出目标。

\textbf{暴力回放。} 慢版本在每个节点重新扫描子树或路径，先求出局部答案。重复扫描暴露了真正状态：每个节点只需要从孩子或父亲那里拿到少量信息。二刷时先写慢版本或口述慢版本，用它确认不会漏答案。

\textbf{特例回放。} 拿 BST [3,1,4,null,2]、k=1 手算，中序访问顺序是 1,2,3,4，第一个就是答案。规律是 BST 的中序严格递增，计数到第 k 个即可停止；重复值是否允许要按题意定义。把这个特例继续拆开看：节点向父亲返回的量和全局答案通常不是同一个量。先写叶子或空节点的返回值，再写父节点如何合并左右孩子，递归式才有边界基础。复盘时把这个特例压成状态字段：节点指针、传入约束或返回值、当前答案、层号或路径状态；空节点的返回值必须和状态定义匹配；再用边界 k 为一、k 为节点数、树不平衡、是否有重复值 反查初始化、更新顺序和退出条件。

\textbf{状态回放。} 手算路线：手算时从叶子开始写返回值，或者从根开始写传入约束。不要只写访问顺序，要写每条边上传递的信息。状态字段：节点指针、传入约束或返回值、当前答案、层号或路径状态；空节点的返回值必须和状态定义匹配。状态升级：把对子树的重复扫描压成返回值或队列状态；每个节点只合并孩子给出的稳定信息。

\textbf{证明和边界。} 证明从子树归纳开始：孩子状态正确时，父节点按题目规则合并后也正确。边界：k 为一、k 为节点数、树不平衡、是否有重复值。变体：若频繁查询，可在节点维护子树大小成为顺序统计树。

\noindent\textbf{LeetCode 235 二叉搜索树最近公共祖先。} 模型：BST 的值域关系能决定两个目标在当前节点的哪一侧。推导重点：两个值都小去左，都大去右，否则当前节点就是分叉点。

\textbf{二刷读题。} 读题时先判断状态方向：父到子、子到父，还是按层传播。本题可以压成“BST 的值域关系能决定两个目标在当前节点的哪一侧”，每个节点的输入和输出状态必须写清。先遮住答案，只保留题面，复述候选对象、合法性条件和输出目标。

\textbf{暴力回放。} 慢版本在每个节点重新扫描子树或路径，先求出局部答案。重复扫描暴露了真正状态：每个节点只需要从孩子或父亲那里拿到少量信息。二刷时先写慢版本或口述慢版本，用它确认不会漏答案。

\textbf{特例回放。} 拿 BST 根 6，p=2，q=8 手算，2 在左侧、8 在右侧，根 6 就是最近公共祖先；若 p 和 q 都小于根，就继续去左子树。规律是 BST 的值域能直接决定搜索方向。把这个特例继续拆开看：节点向父亲返回的量和全局答案通常不是同一个量。先写叶子或空节点的返回值，再写父节点如何合并左右孩子，递归式才有边界基础。复盘时把这个特例压成状态字段：节点指针、传入约束或返回值、当前答案、层号或路径状态；空节点的返回值必须和状态定义匹配；再用边界 一个节点是另一个祖先、目标顺序、重复值约定 反查初始化、更新顺序和退出条件。

\textbf{状态回放。} 手算路线：手算时从叶子开始写返回值，或者从根开始写传入约束。不要只写访问顺序，要写每条边上传递的信息。状态字段：节点指针、传入约束或返回值、当前答案、层号或路径状态；空节点的返回值必须和状态定义匹配。状态升级：把对子树的重复扫描压成返回值或队列状态；每个节点只合并孩子给出的稳定信息。

\textbf{证明和边界。} 证明从子树归纳开始：孩子状态正确时，父节点按题目规则合并后也正确。边界：一个节点是另一个祖先、目标顺序、重复值约定。变体：普通二叉树不能用值域剪枝，需要父指针或后序状态。

\noindent\textbf{LeetCode 236 二叉树最近公共祖先。} 模型：普通树没有有序性，只能依靠祖先关系。推导重点：父指针集合或后序返回目标命中状态都可行。

\textbf{二刷读题。} 读题时先判断状态方向：父到子、子到父，还是按层传播。本题可以压成“普通树没有有序性，只能依靠祖先关系”，每个节点的输入和输出状态必须写清。先遮住答案，只保留题面，复述候选对象、合法性条件和输出目标。

\textbf{暴力回放。} 慢版本在每个节点重新扫描子树或路径，先求出局部答案。重复扫描暴露了真正状态：每个节点只需要从孩子或父亲那里拿到少量信息。二刷时先写慢版本或口述慢版本，用它确认不会漏答案。

\textbf{特例回放。} 拿普通二叉树根 3，p=5、q=1 手算，左右子树分别找到一个目标，根 3 就是最近公共祖先；若两个目标都在左子树，答案应由左子树返回。规律是后序返回是否找到目标或祖先，不能用 BST 的大小关系。把这个特例继续拆开看：节点向父亲返回的量和全局答案通常不是同一个量。先写叶子或空节点的返回值，再写父节点如何合并左右孩子，递归式才有边界基础。复盘时把这个特例压成状态字段：节点指针、传入约束或返回值、当前答案、层号或路径状态；空节点的返回值必须和状态定义匹配；再用边界 目标不存在、p 等于 q、一个是另一个祖先、比较节点地址 反查初始化、更新顺序和退出条件。

\textbf{状态回放。} 手算路线：手算时从叶子开始写返回值，或者从根开始写传入约束。不要只写访问顺序，要写每条边上传递的信息。状态字段：节点指针、传入约束或返回值、当前答案、层号或路径状态；空节点的返回值必须和状态定义匹配。状态升级：把对子树的重复扫描压成返回值或队列状态；每个节点只合并孩子给出的稳定信息。

\textbf{证明和边界。} 证明从子树归纳开始：孩子状态正确时，父节点按题目规则合并后也正确。边界：目标不存在、p 等于 q、一个是另一个祖先、比较节点地址。变体：多次 LCA 查询可预处理倍增祖先表。

\noindent\textbf{LeetCode 337 打家劫舍 III。} 模型：树上相邻节点不能同时选，每个节点需要偷和不偷两个状态。推导重点：后序汇总：偷当前则孩子不能偷，不偷当前则孩子可偷可不偷。

\textbf{二刷读题。} 读题时先判断状态方向：父到子、子到父，还是按层传播。本题可以压成“树上相邻节点不能同时选，每个节点需要偷和不偷两个状态”，每个节点的输入和输出状态必须写清。先遮住答案，只保留题面，复述候选对象、合法性条件和输出目标。

\textbf{暴力回放。} 慢版本在每个节点重新扫描子树或路径，先求出局部答案。重复扫描暴露了真正状态：每个节点只需要从孩子或父亲那里拿到少量信息。二刷时先写慢版本或口述慢版本，用它确认不会漏答案。

\textbf{特例回放。} 拿树 [3,2,3,null,3,null,1] 手算，偷根 3 时不能偷两个孩子，只能接孙子；不偷根时可以分别取左右孩子的最优。规律是每个节点返回两个状态：偷当前和不偷当前，父节点再组合。把这个特例继续拆开看：节点向父亲返回的量和全局答案通常不是同一个量。先写叶子或空节点的返回值，再写父节点如何合并左右孩子，递归式才有边界基础。复盘时把这个特例压成状态字段：节点指针、传入约束或返回值、当前答案、层号或路径状态；空节点的返回值必须和状态定义匹配；再用边界 空树、负值是否可能、链式树、递归深度 反查初始化、更新顺序和退出条件。

\textbf{状态回放。} 手算路线：手算时从叶子开始写返回值，或者从根开始写传入约束。不要只写访问顺序，要写每条边上传递的信息。状态字段：节点指针、传入约束或返回值、当前答案、层号或路径状态；空节点的返回值必须和状态定义匹配。状态升级：把对子树的重复扫描压成返回值或队列状态；每个节点只合并孩子给出的稳定信息。

\textbf{证明和边界。} 证明从子树归纳开始：孩子状态正确时，父节点按题目规则合并后也正确。边界：空树、负值是否可能、链式树、递归深度。变体：可用显式栈后序保存节点状态，避免调用栈。

\noindent\textbf{LeetCode 450 删除二叉搜索树中的节点。} 模型：删除节点后仍要保持 BST 中序有序。推导重点：若有两个孩子，用后继或前驱替换当前值，再删除那个后继节点。

\textbf{二刷读题。} 读题时先判断状态方向：父到子、子到父，还是按层传播。本题可以压成“删除节点后仍要保持 BST 中序有序”，每个节点的输入和输出状态必须写清。先遮住答案，只保留题面，复述候选对象、合法性条件和输出目标。

\textbf{暴力回放。} 慢版本在每个节点重新扫描子树或路径，先求出局部答案。重复扫描暴露了真正状态：每个节点只需要从孩子或父亲那里拿到少量信息。二刷时先写慢版本或口述慢版本，用它确认不会漏答案。

\textbf{特例回放。} 拿 BST 删除 3 且 3 有左右孩子手算，不能直接丢掉整棵子树；可用右子树最小节点接替 3，再在右子树删除这个后继。规律是删除分三类：叶子、单孩子、双孩子，双孩子要用前驱或后继保持 BST 顺序。把这个特例继续拆开看：节点向父亲返回的量和全局答案通常不是同一个量。先写叶子或空节点的返回值，再写父节点如何合并左右孩子，递归式才有边界基础。复盘时把这个特例压成状态字段：节点指针、传入约束或返回值、当前答案、层号或路径状态；空节点的返回值必须和状态定义匹配；再用边界 删除叶子、一个孩子、两个孩子、删除根节点 反查初始化、更新顺序和退出条件。

\textbf{状态回放。} 手算路线：手算时从叶子开始写返回值，或者从根开始写传入约束。不要只写访问顺序，要写每条边上传递的信息。状态字段：节点指针、传入约束或返回值、当前答案、层号或路径状态；空节点的返回值必须和状态定义匹配。状态升级：把对子树的重复扫描压成返回值或队列状态；每个节点只合并孩子给出的稳定信息。

\textbf{证明和边界。} 证明从子树归纳开始：孩子状态正确时，父节点按题目规则合并后也正确。边界：删除叶子、一个孩子、两个孩子、删除根节点。变体：工程实现要清楚是移动值还是重连节点。

\noindent\textbf{LeetCode 543 二叉树直径。} 模型：直径经过某节点时等于左高加右高，但向父亲只能贡献单边高度。推导重点：后序遍历同时更新全局直径并返回高度。

\textbf{二刷读题。} 读题时先判断状态方向：父到子、子到父，还是按层传播。本题可以压成“直径经过某节点时等于左高加右高，但向父亲只能贡献单边高度”，每个节点的输入和输出状态必须写清。先遮住答案，只保留题面，复述候选对象、合法性条件和输出目标。

\textbf{暴力回放。} 慢版本在每个节点重新扫描子树或路径，先求出局部答案。重复扫描暴露了真正状态：每个节点只需要从孩子或父亲那里拿到少量信息。二刷时先写慢版本或口述慢版本，用它确认不会漏答案。

\textbf{特例回放。} 拿树 1 的左右子树高度分别为 2 和 1 手算，经过根的直径是 3 条边；但向父节点只能返回单侧高度。规律是后序返回高度，全局答案用 left\_height+right\_height 更新。把这个特例继续拆开看：节点向父亲返回的量和全局答案通常不是同一个量。先写叶子或空节点的返回值，再写父节点如何合并左右孩子，递归式才有边界基础。复盘时把这个特例压成状态字段：节点指针、传入约束或返回值、当前答案、层号或路径状态；空节点的返回值必须和状态定义匹配；再用边界 空树、单节点、链式树、边数还是节点数定义 反查初始化、更新顺序和退出条件。

\textbf{状态回放。} 手算路线：手算时从叶子开始写返回值，或者从根开始写传入约束。不要只写访问顺序，要写每条边上传递的信息。状态字段：节点指针、传入约束或返回值、当前答案、层号或路径状态；空节点的返回值必须和状态定义匹配。状态升级：把对子树的重复扫描压成返回值或队列状态；每个节点只合并孩子给出的稳定信息。

\textbf{证明和边界。} 证明从子树归纳开始：孩子状态正确时，父节点按题目规则合并后也正确。边界：空树、单节点、链式树、边数还是节点数定义。变体：最大路径和是同型题，但贡献值和全局更新有负数处理。

\subsection{自测标准}

二刷时不要只确认自己见过这道题。请遮住答案，先讲题面模型，再讲暴力基线和瓶颈，然后说出优化结构保存了什么、不变量如何排除错误候选、边界实验如何打破错误实现。

\section{图建模、BFS 与 DFS：二刷复盘卡}
这一大节只围绕一个题型复盘，不把题号直接铺成目录。阅读时先看复盘目标，再读核心题卡，最后用自测标准检查自己是否能独立重讲。

\subsection{复盘目标}

追问通常会问节点和边的建模、访问标记时机、BFS 为什么最短。请准备说明状态是否还包含资源，为什么不能只按位置去重。

\subsection{核心题卡}

\noindent\textbf{LeetCode 126 单词接龙 II。} 模型：不仅要最短转换长度，还要输出所有最短转换路径。推导重点：BFS 只建立最短层级和前驱列表，同一层的多个前驱都保留，最后再从终点回溯路径。

\textbf{二刷读题。} 读题时先定义状态图。本题可以压成“不仅要最短转换长度，还要输出所有最短转换路径”，状态节点和题目原始位置不一定相同。先遮住答案，只保留题面，复述候选对象、合法性条件和输出目标。

\textbf{暴力回放。} 慢版本先按题意画出完整状态图，用普通搜索跑通可达性。这个过程用于确认不仅要最短转换长度，还要输出所有最短转换路径的节点和边，之后再讨论访问标记、层次或代价顺序。二刷时先写慢版本或口述慢版本，用它确认不会漏答案。

\textbf{特例回放。} 拿 hit 到 cog 的小词表手算，同一层 hot 可以通向 dot 和 lot，后续 dog/log 都可能到 cog。若某个单词在本层第一次见到就立刻删掉，会丢掉另一条同长度父路径。规律是 BFS 建最短层级，前驱列表保留同层多父节点，最后再回溯所有路径。把这个特例继续拆开看：状态节点不一定等于原始位置，还可能包含钥匙、剩余资源、时间或访问集合。visited 只能合并未来能力完全相同的状态，否则会把合法路径剪掉。复盘时把这个特例压成状态字段：状态编号、邻接生成方式、访问标记、距离或层数、队列内容；若状态含资源，visited 不能只按位置记录；再用边界 终点不在词表、同层多父节点、本层访问不能提前删除等长前驱、输出规模可能很大 反查初始化、更新顺序和退出条件。

\textbf{状态回放。} 手算路线：手算时画队列或栈，状态必须包含题目需要的所有资源。入队时标记还是出队时标记，要和去重语义一致。状态字段：状态编号、邻接生成方式、访问标记、距离或层数、队列内容；若状态含资源，visited 不能只按位置记录。状态升级：把重复搜索压成访问标记、距离数组或拓扑状态；邻居生成也要从全量扫描变成结构化枚举。

\textbf{证明和边界。} 证明从可达性开始：搜索覆盖所有合法边能到达的状态，访问标记只删除重复状态而不删除不同未来能力。边界：终点不在词表、同层多父节点、本层访问不能提前删除等长前驱、输出规模可能很大。变体：若只问最短长度，普通 BFS 或双向 BFS 更轻；若要路径，就必须保存前驱关系。

\noindent\textbf{LeetCode 127 单词接龙。} 模型：每个单词是节点，一次修改一个字符形成边。推导重点：通配模式桶加速邻居查询，BFS 第一次到达终点就是最短转换长度。

\textbf{二刷读题。} 读题时先定义状态图。本题可以压成“每个单词是节点，一次修改一个字符形成边”，状态节点和题目原始位置不一定相同。先遮住答案，只保留题面，复述候选对象、合法性条件和输出目标。

\textbf{暴力回放。} 慢版本先按题意画出完整状态图，用普通搜索跑通可达性。这个过程用于确认每个单词是节点，一次修改一个字符形成边的节点和边，之后再讨论访问标记、层次或代价顺序。二刷时先写慢版本或口述慢版本，用它确认不会漏答案。

\textbf{特例回放。} 拿 hit -> hot -> dot -> dog -> cog 手算，每次只能改一个字符，所以 BFS 第一层是 hot，第二层才到 dot 或 lot。规律是单词作为节点，通配桶如 h*t 把邻居快速找出；第一次到达终点就是最短转换长度。把这个特例继续拆开看：状态节点不一定等于原始位置，还可能包含钥匙、剩余资源、时间或访问集合。visited 只能合并未来能力完全相同的状态，否则会把合法路径剪掉。复盘时把这个特例压成状态字段：状态编号、邻接生成方式、访问标记、距离或层数、队列内容；若状态含资源，visited 不能只按位置记录；再用边界 终点不在词表、起点和终点相同、桶重复扫描、单词长度一致 反查初始化、更新顺序和退出条件。

\textbf{状态回放。} 手算路线：手算时画队列或栈，状态必须包含题目需要的所有资源。入队时标记还是出队时标记，要和去重语义一致。状态字段：状态编号、邻接生成方式、访问标记、距离或层数、队列内容；若状态含资源，visited 不能只按位置记录。状态升级：把重复搜索压成访问标记、距离数组或拓扑状态；邻居生成也要从全量扫描变成结构化枚举。

\textbf{证明和边界。} 证明从可达性开始：搜索覆盖所有合法边能到达的状态，访问标记只删除重复状态而不删除不同未来能力。边界：终点不在词表、起点和终点相同、桶重复扫描、单词长度一致。变体：双向 BFS 可以进一步降低搜索空间。

\noindent\textbf{LeetCode 130 被围绕的区域。} 模型：只有不连通到边界的 O 才会被翻转。推导重点：从边界 O 出发标记安全区域，再翻转剩余 O。

\textbf{二刷读题。} 读题时先定义状态图。本题可以压成“只有不连通到边界的 O 才会被翻转”，状态节点和题目原始位置不一定相同。先遮住答案，只保留题面，复述候选对象、合法性条件和输出目标。

\textbf{暴力回放。} 慢版本先按题意画出完整状态图，用普通搜索跑通可达性。这个过程用于确认只有不连通到边界的 O 才会被翻转的节点和边，之后再讨论访问标记、层次或代价顺序。二刷时先写慢版本或口述慢版本，用它确认不会漏答案。

\textbf{特例回放。} 拿边界上有 O、内部也有 O 的棋盘手算，边界 O 以及和它连通的 O 不能翻转，真正要翻的是不接触边界的内部块。规律是反向从边界 O 出发标记安全区，再翻转剩下的 O；单行单列全部接触边界。把这个特例继续拆开看：状态节点不一定等于原始位置，还可能包含钥匙、剩余资源、时间或访问集合。visited 只能合并未来能力完全相同的状态，否则会把合法路径剪掉。复盘时把这个特例压成状态字段：状态编号、邻接生成方式、访问标记、距离或层数、队列内容；若状态含资源，visited 不能只按位置记录；再用边界 全边界、无 O、单行单列、内部岛屿 反查初始化、更新顺序和退出条件。

\textbf{状态回放。} 手算路线：手算时画队列或栈，状态必须包含题目需要的所有资源。入队时标记还是出队时标记，要和去重语义一致。状态字段：状态编号、邻接生成方式、访问标记、距离或层数、队列内容；若状态含资源，visited 不能只按位置记录。状态升级：把重复搜索压成访问标记、距离数组或拓扑状态；邻居生成也要从全量扫描变成结构化枚举。

\textbf{证明和边界。} 证明从可达性开始：搜索覆盖所有合法边能到达的状态，访问标记只删除重复状态而不删除不同未来能力。边界：全边界、无 O、单行单列、内部岛屿。变体：这是反向思维：找不能被翻转的区域比直接判断每块区域更简单。

\noindent\textbf{LeetCode 133 克隆图。} 模型：每个原节点需要唯一对应一个新节点，边按原图连接。推导重点：哈希表保存原节点到克隆节点的映射，BFS 或 DFS 遍历图。

\textbf{二刷读题。} 读题时先定义状态图。本题可以压成“每个原节点需要唯一对应一个新节点，边按原图连接”，状态节点和题目原始位置不一定相同。先遮住答案，只保留题面，复述候选对象、合法性条件和输出目标。

\textbf{暴力回放。} 慢版本先按题意画出完整状态图，用普通搜索跑通可达性。这个过程用于确认每个原节点需要唯一对应一个新节点，边按原图连接的节点和边，之后再讨论访问标记、层次或代价顺序。二刷时先写慢版本或口述慢版本，用它确认不会漏答案。

\textbf{特例回放。} 拿三角形 1-2-3-1 手算，克隆 1 时会遇到 2 和 3；克隆 2 时又会遇到 1，如果没有哈希表会无限复制。规律是哈希表保存原节点到新节点的映射，先建节点再连邻居，环和自环都依赖这个映射。把这个特例继续拆开看：状态节点不一定等于原始位置，还可能包含钥匙、剩余资源、时间或访问集合。visited 只能合并未来能力完全相同的状态，否则会把合法路径剪掉。复盘时把这个特例压成状态字段：状态编号、邻接生成方式、访问标记、距离或层数、队列内容；若状态含资源，visited 不能只按位置记录；再用边界 空图、自环、重边、环结构、不能按节点值唯一假设 反查初始化、更新顺序和退出条件。

\textbf{状态回放。} 手算路线：手算时画队列或栈，状态必须包含题目需要的所有资源。入队时标记还是出队时标记，要和去重语义一致。状态字段：状态编号、邻接生成方式、访问标记、距离或层数、队列内容；若状态含资源，visited 不能只按位置记录。状态升级：把重复搜索压成访问标记、距离数组或拓扑状态；邻居生成也要从全量扫描变成结构化枚举。

\textbf{证明和边界。} 证明从可达性开始：搜索覆盖所有合法边能到达的状态，访问标记只删除重复状态而不删除不同未来能力。边界：空图、自环、重边、环结构、不能按节点值唯一假设。变体：若图节点值不唯一，必须用地址作为键。

\noindent\textbf{LeetCode 200 岛屿数量。} 模型：陆地格子是节点，四方向相邻是边，岛屿是连通块。推导重点：扫描遇到未访问陆地就启动搜索并标记整块。

\textbf{二刷读题。} 读题时先定义状态图。本题可以压成“陆地格子是节点，四方向相邻是边，岛屿是连通块”，状态节点和题目原始位置不一定相同。先遮住答案，只保留题面，复述候选对象、合法性条件和输出目标。

\textbf{暴力回放。} 慢版本先按题意画出完整状态图，用普通搜索跑通可达性。这个过程用于确认陆地格子是节点，四方向相邻是边，岛屿是连通块的节点和边，之后再讨论访问标记、层次或代价顺序。二刷时先写慢版本或口述慢版本，用它确认不会漏答案。

\textbf{特例回放。} 拿一个 2x2 小岛手算，从一个陆地出发能沿上下左右标记同一块岛。若不标记访问，会重复计数；若对角线也连通，会误合并。规律是每发现一个未访问陆地，答案加一，并用 DFS/BFS 标记与它四连通的全部陆地。把这个特例继续拆开看：状态节点不一定等于原始位置，还可能包含钥匙、剩余资源、时间或访问集合。visited 只能合并未来能力完全相同的状态，否则会把合法路径剪掉。复盘时把这个特例压成状态字段：状态编号、邻接生成方式、访问标记、距离或层数、队列内容；若状态含资源，visited 不能只按位置记录；再用边界 空网格、全水、全陆、斜对角不连通、是否允许修改输入 反查初始化、更新顺序和退出条件。

\textbf{状态回放。} 手算路线：手算时画队列或栈，状态必须包含题目需要的所有资源。入队时标记还是出队时标记，要和去重语义一致。状态字段：状态编号、邻接生成方式、访问标记、距离或层数、队列内容；若状态含资源，visited 不能只按位置记录。状态升级：把重复搜索压成访问标记、距离数组或拓扑状态；邻居生成也要从全量扫描变成结构化枚举。

\textbf{证明和边界。} 证明从可达性开始：搜索覆盖所有合法边能到达的状态，访问标记只删除重复状态而不删除不同未来能力。边界：空网格、全水、全陆、斜对角不连通、是否允许修改输入。变体：也可用并查集合并相邻陆地，适合动态岛屿扩展。

\noindent\textbf{LeetCode 207 课程表。} 模型：课程是节点，先修关系是有向边，问题是判断是否有环。推导重点：入度拓扑排序不断学习当前无依赖课程。

\textbf{二刷读题。} 读题时先定义状态图。本题可以压成“课程是节点，先修关系是有向边，问题是判断是否有环”，状态节点和题目原始位置不一定相同。先遮住答案，只保留题面，复述候选对象、合法性条件和输出目标。

\textbf{暴力回放。} 慢版本先按题意画出完整状态图，用普通搜索跑通可达性。这个过程用于确认课程是节点，先修关系是有向边，问题是判断是否有环的节点和边，之后再讨论访问标记、层次或代价顺序。二刷时先写慢版本或口述慢版本，用它确认不会漏答案。

\textbf{特例回放。} 拿课程 0->1 手算，入度为 0 的课程 0 先入队，学完后 1 的入度降为 0；若有 0->1 和 1->0，队列一开始为空。规律是拓扑排序用入度为 0 的节点推进，处理数量小于课程数说明有环。把这个特例继续拆开看：状态节点不一定等于原始位置，还可能包含钥匙、剩余资源、时间或访问集合。visited 只能合并未来能力完全相同的状态，否则会把合法路径剪掉。复盘时把这个特例压成状态字段：状态编号、邻接生成方式、访问标记、距离或层数、队列内容；若状态含资源，visited 不能只按位置记录；再用边界 边方向、孤立课程、环、自依赖、重复边 反查初始化、更新顺序和退出条件。

\textbf{状态回放。} 手算路线：手算时画队列或栈，状态必须包含题目需要的所有资源。入队时标记还是出队时标记，要和去重语义一致。状态字段：状态编号、邻接生成方式、访问标记、距离或层数、队列内容；若状态含资源，visited 不能只按位置记录。状态升级：把重复搜索压成访问标记、距离数组或拓扑状态；邻居生成也要从全量扫描变成结构化枚举。

\textbf{证明和边界。} 证明从可达性开始：搜索覆盖所有合法边能到达的状态，访问标记只删除重复状态而不删除不同未来能力。边界：边方向、孤立课程、环、自依赖、重复边。变体：DFS 三色标记也能判环，但要明确访问状态。

\noindent\textbf{LeetCode 210 课程表 II。} 模型：在判环基础上还要输出一个可行拓扑序。推导重点：每弹出一个入度为零课程，就把它加入顺序并删除出边。

\textbf{二刷读题。} 读题时先定义状态图。本题可以压成“在判环基础上还要输出一个可行拓扑序”，状态节点和题目原始位置不一定相同。先遮住答案，只保留题面，复述候选对象、合法性条件和输出目标。

\textbf{暴力回放。} 慢版本先按题意画出完整状态图，用普通搜索跑通可达性。这个过程用于确认在判环基础上还要输出一个可行拓扑序的节点和边，之后再讨论访问标记、层次或代价顺序。二刷时先写慢版本或口述慢版本，用它确认不会漏答案。

\textbf{特例回放。} 拿 0->1、0->2 手算，课程 0 入度为 0，先输出 0，再把 1 和 2 的入度减到 0。若存在环，输出数量会少于课程数。规律是课程表 II 比判定题多了输出顺序，队列弹出的顺序就是一种可行拓扑序。把这个特例继续拆开看：状态节点不一定等于原始位置，还可能包含钥匙、剩余资源、时间或访问集合。visited 只能合并未来能力完全相同的状态，否则会把合法路径剪掉。复盘时把这个特例压成状态字段：状态编号、邻接生成方式、访问标记、距离或层数、队列内容；若状态含资源，visited 不能只按位置记录；再用边界 有环返回空数组、多个可行顺序、边方向 反查初始化、更新顺序和退出条件。

\textbf{状态回放。} 手算路线：手算时画队列或栈，状态必须包含题目需要的所有资源。入队时标记还是出队时标记，要和去重语义一致。状态字段：状态编号、邻接生成方式、访问标记、距离或层数、队列内容；若状态含资源，visited 不能只按位置记录。状态升级：把重复搜索压成访问标记、距离数组或拓扑状态；邻居生成也要从全量扫描变成结构化枚举。

\textbf{证明和边界。} 证明从可达性开始：搜索覆盖所有合法边能到达的状态，访问标记只删除重复状态而不删除不同未来能力。边界：有环返回空数组、多个可行顺序、边方向。变体：若需要字典序最小拓扑序，把队列换成小顶堆。

\noindent\textbf{LeetCode 417 太平洋大西洋水流问题。} 模型：从每个格子向海搜索很慢，反向从海边沿非递减高度搜索。推导重点：分别标记能到太平洋和大西洋的格子，交集就是答案。

\textbf{二刷读题。} 读题时先定义状态图。本题可以压成“从每个格子向海搜索很慢，反向从海边沿非递减高度搜索”，状态节点和题目原始位置不一定相同。先遮住答案，只保留题面，复述候选对象、合法性条件和输出目标。

\textbf{暴力回放。} 慢版本先按题意画出完整状态图，用普通搜索跑通可达性。这个过程用于确认从每个格子向海搜索很慢，反向从海边沿非递减高度搜索的节点和边，之后再讨论访问标记、层次或代价顺序。二刷时先写慢版本或口述慢版本，用它确认不会漏答案。

\textbf{特例回放。} 拿高度矩阵 [[1,2],[4,3]] 手算，水从高处流向低处；反向从太平洋边界向内，只能走到更高或相等的格子。规律是分别从两个海洋边界反向 BFS/DFS，两个访问集合的交集就是答案。把这个特例继续拆开看：状态节点不一定等于原始位置，还可能包含钥匙、剩余资源、时间或访问集合。visited 只能合并未来能力完全相同的状态，否则会把合法路径剪掉。复盘时把这个特例压成状态字段：状态编号、邻接生成方式、访问标记、距离或层数、队列内容；若状态含资源，visited 不能只按位置记录；再用边界 单行单列、平台高度、边界重复入队、方向条件反向 反查初始化、更新顺序和退出条件。

\textbf{状态回放。} 手算路线：手算时画队列或栈，状态必须包含题目需要的所有资源。入队时标记还是出队时标记，要和去重语义一致。状态字段：状态编号、邻接生成方式、访问标记、距离或层数、队列内容；若状态含资源，visited 不能只按位置记录。状态升级：把重复搜索压成访问标记、距离数组或拓扑状态；邻居生成也要从全量扫描变成结构化枚举。

\textbf{证明和边界。} 证明从可达性开始：搜索覆盖所有合法边能到达的状态，访问标记只删除重复状态而不删除不同未来能力。边界：单行单列、平台高度、边界重复入队、方向条件反向。变体：反向搜索是很多可达性题的重要技巧。

\noindent\textbf{LeetCode 695 岛屿的最大面积。} 模型：面积就是一个连通块中陆地格子的数量。推导重点：每发现新陆地就搜索完整连通块并计数，取最大值。

\textbf{二刷读题。} 读题时先定义状态图。本题可以压成“面积就是一个连通块中陆地格子的数量”，状态节点和题目原始位置不一定相同。先遮住答案，只保留题面，复述候选对象、合法性条件和输出目标。

\textbf{暴力回放。} 慢版本先按题意画出完整状态图，用普通搜索跑通可达性。这个过程用于确认面积就是一个连通块中陆地格子的数量的节点和边，之后再讨论访问标记、层次或代价顺序。二刷时先写慢版本或口述慢版本，用它确认不会漏答案。

\textbf{特例回放。} 拿网格 [[1,1,0],[0,1,0]] 手算，从左上陆地出发能访问 3 个四连通格子；对角线不算连通。规律是每发现未访问陆地就 DFS/BFS 统计面积并标记，最大值随每个连通块更新。把这个特例继续拆开看：状态节点不一定等于原始位置，还可能包含钥匙、剩余资源、时间或访问集合。visited 只能合并未来能力完全相同的状态，否则会把合法路径剪掉。复盘时把这个特例压成状态字段：状态编号、邻接生成方式、访问标记、距离或层数、队列内容；若状态含资源，visited 不能只按位置记录；再用边界 全水、全陆、细长岛、是否允许修改输入 反查初始化、更新顺序和退出条件。

\textbf{状态回放。} 手算路线：手算时画队列或栈，状态必须包含题目需要的所有资源。入队时标记还是出队时标记，要和去重语义一致。状态字段：状态编号、邻接生成方式、访问标记、距离或层数、队列内容；若状态含资源，visited 不能只按位置记录。状态升级：把重复搜索压成访问标记、距离数组或拓扑状态；邻居生成也要从全量扫描变成结构化枚举。

\textbf{证明和边界。} 证明从可达性开始：搜索覆盖所有合法边能到达的状态，访问标记只删除重复状态而不删除不同未来能力。边界：全水、全陆、细长岛、是否允许修改输入。变体：与岛屿数量相比，只是连通块聚合值不同。

\noindent\textbf{LeetCode 752 打开转盘锁。} 模型：四位密码是状态，每次转动一位形成边。推导重点：BFS 从起点扩展，死亡状态不能入队，访问状态避免重复。

\textbf{二刷读题。} 读题时先定义状态图。本题可以压成“四位密码是状态，每次转动一位形成边”，状态节点和题目原始位置不一定相同。先遮住答案，只保留题面，复述候选对象、合法性条件和输出目标。

\textbf{暴力回放。} 慢版本先按题意画出完整状态图，用普通搜索跑通可达性。这个过程用于确认四位密码是状态，每次转动一位形成边的节点和边，之后再讨论访问标记、层次或代价顺序。二刷时先写慢版本或口述慢版本，用它确认不会漏答案。

\textbf{特例回放。} 拿锁 0000 到 0001 手算，每一步只能拨动一位加一或减一，0000 的邻居有 1000、9000、0100 等。规律是状态是四位字符串，BFS 按步数扩展，deadend 入队前就要排除。把这个特例继续拆开看：状态节点不一定等于原始位置，还可能包含钥匙、剩余资源、时间或访问集合。visited 只能合并未来能力完全相同的状态，否则会把合法路径剪掉。复盘时把这个特例压成状态字段：状态编号、邻接生成方式、访问标记、距离或层数、队列内容；若状态含资源，visited 不能只按位置记录；再用边界 起点死亡、目标就是起点、数字环绕、死亡集合很大 反查初始化、更新顺序和退出条件。

\textbf{状态回放。} 手算路线：手算时画队列或栈，状态必须包含题目需要的所有资源。入队时标记还是出队时标记，要和去重语义一致。状态字段：状态编号、邻接生成方式、访问标记、距离或层数、队列内容；若状态含资源，visited 不能只按位置记录。状态升级：把重复搜索压成访问标记、距离数组或拓扑状态；邻居生成也要从全量扫描变成结构化枚举。

\textbf{证明和边界。} 证明从可达性开始：搜索覆盖所有合法边能到达的状态，访问标记只删除重复状态而不删除不同未来能力。边界：起点死亡、目标就是起点、数字环绕、死亡集合很大。变体：双向 BFS 能明显减少状态展开。

\noindent\textbf{LeetCode 815 公交路线。} 模型：公交路线和站点构成二分关系，换乘次数比站点步数更重要。推导重点：用站点到路线列表的哈希表找可乘路线，以路线或乘车层数做 BFS，访问过的路线不再重复展开。

\textbf{二刷读题。} 读题时先定义状态图。本题可以压成“公交路线和站点构成二分关系，换乘次数比站点步数更重要”，状态节点和题目原始位置不一定相同。先遮住答案，只保留题面，复述候选对象、合法性条件和输出目标。

\textbf{暴力回放。} 慢版本先按题意画出完整状态图，用普通搜索跑通可达性。这个过程用于确认公交路线和站点构成二分关系，换乘次数比站点步数更重要的节点和边，之后再讨论访问标记、层次或代价顺序。二刷时先写慢版本或口述慢版本，用它确认不会漏答案。

\textbf{特例回放。} 拿起点站可坐路线 A，A 和 B 在某站相交，B 到终点手算。按站点一步步 BFS 会反复展开同一条路线的所有站；按路线 BFS，一次乘车层数就是换乘次数。规律是访问过的路线不再展开，站点到路线表只用于找到下一批路线。把这个特例继续拆开看：状态节点不一定等于原始位置，还可能包含钥匙、剩余资源、时间或访问集合。visited 只能合并未来能力完全相同的状态，否则会把合法路径剪掉。复盘时把这个特例压成状态字段：状态编号、邻接生成方式、访问标记、距离或层数、队列内容；若状态含资源，visited 不能只按位置记录；再用边界 起点等于终点、某站经过路线很多、路线重复访问、目标站不可达 反查初始化、更新顺序和退出条件。

\textbf{状态回放。} 手算路线：手算时画队列或栈，状态必须包含题目需要的所有资源。入队时标记还是出队时标记，要和去重语义一致。状态字段：状态编号、邻接生成方式、访问标记、距离或层数、队列内容；若状态含资源，visited 不能只按位置记录。状态升级：把重复搜索压成访问标记、距离数组或拓扑状态；邻居生成也要从全量扫描变成结构化枚举。

\textbf{证明和边界。} 证明从可达性开始：搜索覆盖所有合法边能到达的状态，访问标记只删除重复状态而不删除不同未来能力。边界：起点等于终点、某站经过路线很多、路线重复访问、目标站不可达。变体：若把每个站点之间完全建边会爆炸，应利用路线这一层压缩邻居生成。

\noindent\textbf{LeetCode 864 获取所有钥匙的最短路径。} 模型：位置相同但钥匙集合不同，未来能开的门不同。推导重点：BFS 状态包含行、列和钥匙掩码，访问标记也必须包含掩码。

\textbf{二刷读题。} 读题时先定义状态图。本题可以压成“位置相同但钥匙集合不同，未来能开的门不同”，状态节点和题目原始位置不一定相同。先遮住答案，只保留题面，复述候选对象、合法性条件和输出目标。

\textbf{暴力回放。} 慢版本先按题意画出完整状态图，用普通搜索跑通可达性。这个过程用于确认位置相同但钥匙集合不同，未来能开的门不同的节点和边，之后再讨论访问标记、层次或代价顺序。二刷时先写慢版本或口述慢版本，用它确认不会漏答案。

\textbf{特例回放。} 拿网格 @.a / \#.\# / b.A 手算，走到 a 后状态不是位置变了而是钥匙集合多了一把；同一位置带不同钥匙未来能力不同。规律是 BFS 的 visited 必须按 (row,col,keymask) 记录，拿齐所有钥匙第一次出队就是最短步数。把这个特例继续拆开看：状态节点不一定等于原始位置，还可能包含钥匙、剩余资源、时间或访问集合。visited 只能合并未来能力完全相同的状态，否则会把合法路径剪掉。复盘时把这个特例压成状态字段：状态编号、邻接生成方式、访问标记、距离或层数、队列内容；若状态含资源，visited 不能只按位置记录；再用边界 起点旁边有门、钥匙顺序、不可达、钥匙数量决定掩码大小 反查初始化、更新顺序和退出条件。

\textbf{状态回放。} 手算路线：手算时画队列或栈，状态必须包含题目需要的所有资源。入队时标记还是出队时标记，要和去重语义一致。状态字段：状态编号、邻接生成方式、访问标记、距离或层数、队列内容；若状态含资源，visited 不能只按位置记录。状态升级：把重复搜索压成访问标记、距离数组或拓扑状态；邻居生成也要从全量扫描变成结构化枚举。

\textbf{证明和边界。} 证明从可达性开始：搜索覆盖所有合法边能到达的状态，访问标记只删除重复状态而不删除不同未来能力。边界：起点旁边有门、钥匙顺序、不可达、钥匙数量决定掩码大小。变体：这是状态图维度不能只按位置去重的典型题。

\noindent\textbf{LeetCode 994 腐烂的橘子。} 模型：所有初始腐烂橘子同时作为 BFS 源点。推导重点：按层扩散，分钟数等于最后一个新鲜橘子被首次访问的层数。

\textbf{二刷读题。} 读题时先定义状态图。本题可以压成“所有初始腐烂橘子同时作为 BFS 源点”，状态节点和题目原始位置不一定相同。先遮住答案，只保留题面，复述候选对象、合法性条件和输出目标。

\textbf{暴力回放。} 慢版本先按题意画出完整状态图，用普通搜索跑通可达性。这个过程用于确认所有初始腐烂橘子同时作为 BFS 源点的节点和边，之后再讨论访问标记、层次或代价顺序。二刷时先写慢版本或口述慢版本，用它确认不会漏答案。

\textbf{特例回放。} 拿 [[2,1,1],[1,1,0],[0,1,1]] 手算，所有腐烂橘子同时向四周扩散，第一分钟会感染相邻新鲜橘子。规律是多源 BFS，初始所有腐烂橘子同时入队，层数就是分钟数。把这个特例继续拆开看：状态节点不一定等于原始位置，还可能包含钥匙、剩余资源、时间或访问集合。visited 只能合并未来能力完全相同的状态，否则会把合法路径剪掉。复盘时把这个特例压成状态字段：状态编号、邻接生成方式、访问标记、距离或层数、队列内容；若状态含资源，visited 不能只按位置记录；再用边界 没有新鲜橘子、没有腐烂源、隔离新鲜橘子、空格阻断 反查初始化、更新顺序和退出条件。

\textbf{状态回放。} 手算路线：手算时画队列或栈，状态必须包含题目需要的所有资源。入队时标记还是出队时标记，要和去重语义一致。状态字段：状态编号、邻接生成方式、访问标记、距离或层数、队列内容；若状态含资源，visited 不能只按位置记录。状态升级：把重复搜索压成访问标记、距离数组或拓扑状态；邻居生成也要从全量扫描变成结构化枚举。

\textbf{证明和边界。} 证明从可达性开始：搜索覆盖所有合法边能到达的状态，访问标记只删除重复状态而不删除不同未来能力。边界：没有新鲜橘子、没有腐烂源、隔离新鲜橘子、空格阻断。变体：多源 BFS 也适用于最近零距离、地图扩散等题。

\noindent\textbf{LeetCode 1091 二进制矩阵中的最短路径。} 模型：八方向网格无权边，答案是从左上到右下的最少格子数。推导重点：BFS 入队时标记访问，层数或距离随状态保存。

\textbf{二刷读题。} 读题时先定义状态图。本题可以压成“八方向网格无权边，答案是从左上到右下的最少格子数”，状态节点和题目原始位置不一定相同。先遮住答案，只保留题面，复述候选对象、合法性条件和输出目标。

\textbf{暴力回放。} 慢版本先按题意画出完整状态图，用普通搜索跑通可达性。这个过程用于确认八方向网格无权边，答案是从左上到右下的最少格子数的节点和边，之后再讨论访问标记、层次或代价顺序。二刷时先写慢版本或口述慢版本，用它确认不会漏答案。

\textbf{特例回放。} 拿 [[0,1],[1,0]] 手算，从左上可以斜着一步到右下，路径长度是 2；若起点或终点为 1 直接不可达。规律是八方向 BFS，第一次到达终点就是最短路径长度。把这个特例继续拆开看：状态节点不一定等于原始位置，还可能包含钥匙、剩余资源、时间或访问集合。visited 只能合并未来能力完全相同的状态，否则会把合法路径剪掉。复盘时把这个特例压成状态字段：状态编号、邻接生成方式、访问标记、距离或层数、队列内容；若状态含资源，visited 不能只按位置记录；再用边界 起点或终点阻塞、单格、八方向越界、无路可达 反查初始化、更新顺序和退出条件。

\textbf{状态回放。} 手算路线：手算时画队列或栈，状态必须包含题目需要的所有资源。入队时标记还是出队时标记，要和去重语义一致。状态字段：状态编号、邻接生成方式、访问标记、距离或层数、队列内容；若状态含资源，visited 不能只按位置记录。状态升级：把重复搜索压成访问标记、距离数组或拓扑状态；邻居生成也要从全量扫描变成结构化枚举。

\textbf{证明和边界。} 证明从可达性开始：搜索覆盖所有合法边能到达的状态，访问标记只删除重复状态而不删除不同未来能力。边界：起点或终点阻塞、单格、八方向越界、无路可达。变体：边权相同才能用 BFS，若每格有代价就要最短路。

\noindent\textbf{LeetCode 1293 网格中的最短路径。} 模型：走到同一格时，剩余可消除障碍数量不同会影响未来可达性。推导重点：BFS 状态包含位置和剩余消除次数，visited 记录该维度或记录到达该格的最大剩余次数。

\textbf{二刷读题。} 读题时先定义状态图。本题可以压成“走到同一格时，剩余可消除障碍数量不同会影响未来可达性”，状态节点和题目原始位置不一定相同。先遮住答案，只保留题面，复述候选对象、合法性条件和输出目标。

\textbf{暴力回放。} 慢版本先按题意画出完整状态图，用普通搜索跑通可达性。这个过程用于确认走到同一格时，剩余可消除障碍数量不同会影响未来可达性的节点和边，之后再讨论访问标记、层次或代价顺序。二刷时先写慢版本或口述慢版本，用它确认不会漏答案。

\textbf{特例回放。} 拿网格三列中间有障碍且 k=1 手算，走到同一格时，如果剩余消障次数不同，未来可走路径也不同。规律是 BFS 的状态必须包含剩余 k，visited 不能只按位置记录。把这个特例继续拆开看：状态节点不一定等于原始位置，还可能包含钥匙、剩余资源、时间或访问集合。visited 只能合并未来能力完全相同的状态，否则会把合法路径剪掉。复盘时把这个特例压成状态字段：状态编号、邻接生成方式、访问标记、距离或层数、队列内容；若状态含资源，visited 不能只按位置记录；再用边界 起点终点、k 很大、障碍密集、二维 visited 错误剪枝 反查初始化、更新顺序和退出条件。

\textbf{状态回放。} 手算路线：手算时画队列或栈，状态必须包含题目需要的所有资源。入队时标记还是出队时标记，要和去重语义一致。状态字段：状态编号、邻接生成方式、访问标记、距离或层数、队列内容；若状态含资源，visited 不能只按位置记录。状态升级：把重复搜索压成访问标记、距离数组或拓扑状态；邻居生成也要从全量扫描变成结构化枚举。

\textbf{证明和边界。} 证明从可达性开始：搜索覆盖所有合法边能到达的状态，访问标记只删除重复状态而不删除不同未来能力。边界：起点终点、k 很大、障碍密集、二维 visited 错误剪枝。变体：它和钥匙题一样说明状态维度不能丢。

\noindent\textbf{LeetCode 1462 课程表 IV。} 模型：大量先修关系查询需要预处理课程之间的可达性。推导重点：可以用 Floyd 传递闭包、每个点 BFS，或拓扑 DP 合并祖先集合，查询时直接查可达表。

\textbf{二刷读题。} 读题时先定义状态图。本题可以压成“大量先修关系查询需要预处理课程之间的可达性”，状态节点和题目原始位置不一定相同。先遮住答案，只保留题面，复述候选对象、合法性条件和输出目标。

\textbf{暴力回放。} 慢版本先按题意画出完整状态图，用普通搜索跑通可达性。这个过程用于确认大量先修关系查询需要预处理课程之间的可达性的节点和边，之后再讨论访问标记、层次或代价顺序。二刷时先写慢版本或口述慢版本，用它确认不会漏答案。

\textbf{特例回放。} 拿 0->1->2 和查询 0 是否先修 2 手算，直接边只能回答 0->1、1->2，不能回答传递关系。若查询很多，每次 DFS 会重复走相同路径。规律是先做可达性闭包或拓扑合并祖先集合，再 O(1) 回答查询。把这个特例继续拆开看：状态节点不一定等于原始位置，还可能包含钥匙、剩余资源、时间或访问集合。visited 只能合并未来能力完全相同的状态，否则会把合法路径剪掉。复盘时把这个特例压成状态字段：状态编号、邻接生成方式、访问标记、距离或层数、队列内容；若状态含资源，visited 不能只按位置记录；再用边界 课程之间没有关系、自身是否算先修、查询很多、图中按题意应为 DAG 反查初始化、更新顺序和退出条件。

\textbf{状态回放。} 手算路线：手算时画队列或栈，状态必须包含题目需要的所有资源。入队时标记还是出队时标记，要和去重语义一致。状态字段：状态编号、邻接生成方式、访问标记、距离或层数、队列内容；若状态含资源，visited 不能只按位置记录。状态升级：把重复搜索压成访问标记、距离数组或拓扑状态；邻居生成也要从全量扫描变成结构化枚举。

\textbf{证明和边界。} 证明从可达性开始：搜索覆盖所有合法边能到达的状态，访问标记只删除重复状态而不删除不同未来能力。边界：课程之间没有关系、自身是否算先修、查询很多、图中按题意应为 DAG。变体：如果查询很少，逐次搜索可能更省预处理成本；多查询时闭包更稳。

\subsection{自测标准}

二刷时不要只确认自己见过这道题。请遮住答案，先讲题面模型，再讲暴力基线和瓶颈，然后说出优化结构保存了什么、不变量如何排除错误候选、边界实验如何打破错误实现。

\section{最短路与带权图：二刷复盘卡}
这一大节只围绕一个题型复盘，不把题号直接铺成目录。阅读时先看复盘目标，再读核心题卡，最后用自测标准检查自己是否能独立重讲。

\subsection{复盘目标}

追问通常会问为什么普通 BFS 不行、Dijkstra 的非负边条件、堆中旧状态如何处理。请准备一个不同边权反例。

\subsection{核心题卡}

\noindent\textbf{LeetCode 505 迷宫 II。} 模型：小球一旦选择方向就会滚到墙前停止，边权是滚动经过的格子数。推导重点：邻居生成不是走一步，而是沿四个方向模拟滚动到停点，再用 Dijkstra 按累计距离松弛。

\textbf{二刷读题。} 读题时先区分可达、最少边数和最小代价。本题可以压成“小球一旦选择方向就会滚到墙前停止，边权是滚动经过的格子数”，邻居生成不是走一步，而是沿四个方向模拟滚动到停点，再用 Dijkstra 按累计距离松弛决定扩展顺序。先遮住答案，只保留题面，复述候选对象、合法性条件和输出目标。

\textbf{暴力回放。} 慢版本可以枚举路径，或先用普通队列试跑一个小图。它会暴露步数最少和代价最小是否一致；一旦不一致，就必须围绕代价组织状态。二刷时先写慢版本或口述慢版本，用它确认不会漏答案。

\textbf{特例回放。} 拿迷宫中球从 (0,4) 滚向左手算，它不会每格停下，只会直到撞墙才停。普通 BFS 按格走会错，边权是滚动距离。规律是状态是停球位置，邻居是四个方向滚到的终点，距离要用 Dijkstra 更新。把这个特例继续拆开看：队列或堆弹出的顺序必须和代价规则一致；旧状态被跳过，是因为已经有更小代价到达同一状态。只要边权或代价合成方式改变，就要重新证明扩展顺序。复盘时把这个特例压成状态字段：距离数组、优先队列状态、松弛候选、旧状态判定、不可达哨兵；资源约束题还要把资源维度放进状态；再用边界 起点就是终点、目标无法停下、绕远路更短、同一停点多次被松弛 反查初始化、更新顺序和退出条件。

\textbf{状态回放。} 手算路线：手算时记录每次弹出的代价、当前距离数组和松弛结果。旧状态被弹出时为什么跳过，也要写在表里。状态字段：距离数组、优先队列状态、松弛候选、旧状态判定、不可达哨兵；资源约束题还要把资源维度放进状态。状态升级：把路径枚举压成距离数组和优先队列；每次只扩展当前最有希望的未过期状态。

\textbf{证明和边界。} 证明从扩展顺序开始：当前弹出的状态在代价规则下已经不会被未来路径改得更优。边界：起点就是终点、目标无法停下、绕远路更短、同一停点多次被松弛。变体：若只问能否到达，可用 BFS 或 DFS；若问最少距离，就必须处理带权边。

\noindent\textbf{LeetCode 743 网络延迟时间。} 模型：有向正权图从起点到所有节点的最短路最大值。推导重点：Dijkstra 求每个节点最短距离，最后取最大，若有无穷则不可达。

\textbf{二刷读题。} 读题时先区分可达、最少边数和最小代价。本题可以压成“有向正权图从起点到所有节点的最短路最大值”，Dijkstra 求每个节点最短距离，最后取最大，若有无穷则不可达决定扩展顺序。先遮住答案，只保留题面，复述候选对象、合法性条件和输出目标。

\textbf{暴力回放。} 慢版本可以枚举路径，或先用普通队列试跑一个小图。它会暴露步数最少和代价最小是否一致；一旦不一致，就必须围绕代价组织状态。二刷时先写慢版本或口述慢版本，用它确认不会漏答案。

\textbf{特例回放。} 拿 times=[[2,1,1],[2,3,1],[3,4,1]]、起点 2 手算，2 到 1 和 3 距离都是 1，到 4 距离是 2。规律是非负边最短路用 Dijkstra，最后取所有最短距离的最大值；不可达节点导致答案 -1。把这个特例继续拆开看：队列或堆弹出的顺序必须和代价规则一致；旧状态被跳过，是因为已经有更小代价到达同一状态。只要边权或代价合成方式改变，就要重新证明扩展顺序。复盘时把这个特例压成状态字段：距离数组、优先队列状态、松弛候选、旧状态判定、不可达哨兵；资源约束题还要把资源维度放进状态；再用边界 节点编号从一开始、不可达、重边、过期堆元素 反查初始化、更新顺序和退出条件。

\textbf{状态回放。} 手算路线：手算时记录每次弹出的代价、当前距离数组和松弛结果。旧状态被弹出时为什么跳过，也要写在表里。状态字段：距离数组、优先队列状态、松弛候选、旧状态判定、不可达哨兵；资源约束题还要把资源维度放进状态。状态升级：把路径枚举压成距离数组和优先队列；每次只扩展当前最有希望的未过期状态。

\textbf{证明和边界。} 证明从扩展顺序开始：当前弹出的状态在代价规则下已经不会被未来路径改得更优。边界：节点编号从一开始、不可达、重边、过期堆元素。变体：边权相同才可用普通 BFS。

\noindent\textbf{LeetCode 778 水位上升的泳池中游泳。} 模型：到达某格子的代价是路径上最大高度，而不是高度和。推导重点：用 Dijkstra 按当前路径最大高度扩展，松弛时取 max 当前代价和邻格高度。

\textbf{二刷读题。} 读题时先区分可达、最少边数和最小代价。本题可以压成“到达某格子的代价是路径上最大高度，而不是高度和”，用 Dijkstra 按当前路径最大高度扩展，松弛时取 max 当前代价和邻格高度决定扩展顺序。先遮住答案，只保留题面，复述候选对象、合法性条件和输出目标。

\textbf{暴力回放。} 慢版本可以枚举路径，或先用普通队列试跑一个小图。它会暴露步数最少和代价最小是否一致；一旦不一致，就必须围绕代价组织状态。二刷时先写慢版本或口述慢版本，用它确认不会漏答案。

\textbf{特例回放。} 拿 [[0,2],[1,3]] 手算，时间必须至少到 3 才能进入右下角；路径代价不是边权和，而是沿路最大高度。规律是 Dijkstra 的距离定义为到达某格所需的最小最大高度，或二分时间做可达性。把这个特例继续拆开看：队列或堆弹出的顺序必须和代价规则一致；旧状态被跳过，是因为已经有更小代价到达同一状态。只要边权或代价合成方式改变，就要重新证明扩展顺序。复盘时把这个特例压成状态字段：距离数组、优先队列状态、松弛候选、旧状态判定、不可达哨兵；资源约束题还要把资源维度放进状态；再用边界 单格、起点高度很高、需要绕路、多个相同高度 反查初始化、更新顺序和退出条件。

\textbf{状态回放。} 手算路线：手算时记录每次弹出的代价、当前距离数组和松弛结果。旧状态被弹出时为什么跳过，也要写在表里。状态字段：距离数组、优先队列状态、松弛候选、旧状态判定、不可达哨兵；资源约束题还要把资源维度放进状态。状态升级：把路径枚举压成距离数组和优先队列；每次只扩展当前最有希望的未过期状态。

\textbf{证明和边界。} 证明从扩展顺序开始：当前弹出的状态在代价规则下已经不会被未来路径改得更优。边界：单格、起点高度很高、需要绕路、多个相同高度。变体：也可以按水位二分可达，或按高度排序并查集合并。

\noindent\textbf{LeetCode 787 K 站中转内最便宜的航班。} 模型：路径代价受中转次数限制，不能只按费用最小弹出节点后永久确定。推导重点：用按边数分层的 Bellman-Ford 或带层数状态的优先队列，状态包含城市和已用边数。

\textbf{二刷读题。} 读题时先区分可达、最少边数和最小代价。本题可以压成“路径代价受中转次数限制，不能只按费用最小弹出节点后永久确定”，用按边数分层的 Bellman-Ford 或带层数状态的优先队列，状态包含城市和已用边数决定扩展顺序。先遮住答案，只保留题面，复述候选对象、合法性条件和输出目标。

\textbf{暴力回放。} 慢版本可以枚举路径，或先用普通队列试跑一个小图。它会暴露步数最少和代价最小是否一致；一旦不一致，就必须围绕代价组织状态。二刷时先写慢版本或口述慢版本，用它确认不会漏答案。

\textbf{特例回放。} 拿航班 0->1 价格 100、1->2 价格 100、0->2 价格 500，K=1 手算，允许一次中转时 0->1->2 更便宜。规律是状态必须包含已用中转次数或边数，不能只保存每个城市一个最低价格。把这个特例继续拆开看：队列或堆弹出的顺序必须和代价规则一致；旧状态被跳过，是因为已经有更小代价到达同一状态。只要边权或代价合成方式改变，就要重新证明扩展顺序。复盘时把这个特例压成状态字段：距离数组、优先队列状态、松弛候选、旧状态判定、不可达哨兵；资源约束题还要把资源维度放进状态；再用边界 起点终点相同、不可达、K 为零、便宜路径中转过多 反查初始化、更新顺序和退出条件。

\textbf{状态回放。} 手算路线：手算时记录每次弹出的代价、当前距离数组和松弛结果。旧状态被弹出时为什么跳过，也要写在表里。状态字段：距离数组、优先队列状态、松弛候选、旧状态判定、不可达哨兵；资源约束题还要把资源维度放进状态。状态升级：把路径枚举压成距离数组和优先队列；每次只扩展当前最有希望的未过期状态。

\textbf{证明和边界。} 证明从扩展顺序开始：当前弹出的状态在代价规则下已经不会被未来路径改得更优。边界：起点终点相同、不可达、K 为零、便宜路径中转过多。变体：若去掉中转限制，普通 Dijkstra 才能直接按费用组织扩展。

\noindent\textbf{LeetCode 1102 得分最高的路径。} 模型：路径得分是沿途最小格子值，要最大化这个最小值。推导重点：用最大堆按当前路径得分扩展，进入邻格时取当前得分和邻格值的较小者。

\textbf{二刷读题。} 读题时先区分可达、最少边数和最小代价。本题可以压成“路径得分是沿途最小格子值，要最大化这个最小值”，用最大堆按当前路径得分扩展，进入邻格时取当前得分和邻格值的较小者决定扩展顺序。先遮住答案，只保留题面，复述候选对象、合法性条件和输出目标。

\textbf{暴力回放。} 慢版本可以枚举路径，或先用普通队列试跑一个小图。它会暴露步数最少和代价最小是否一致；一旦不一致，就必须围绕代价组织状态。二刷时先写慢版本或口述慢版本，用它确认不会漏答案。

\textbf{特例回放。} 拿 [[5,4,5],[1,2,6],[7,4,6]] 手算，路径分数是沿路最小值，要最大化这个最小值。规律是优先扩展当前路径分数最高的格子，或二分阈值判断能否只走不低于阈值的格子。把这个特例继续拆开看：队列或堆弹出的顺序必须和代价规则一致；旧状态被跳过，是因为已经有更小代价到达同一状态。只要边权或代价合成方式改变，就要重新证明扩展顺序。复盘时把这个特例压成状态字段：距离数组、优先队列状态、松弛候选、旧状态判定、不可达哨兵；资源约束题还要把资源维度放进状态；再用边界 单格、起点或终点很小、必须绕路、多个相同得分路径 反查初始化、更新顺序和退出条件。

\textbf{状态回放。} 手算路线：手算时记录每次弹出的代价、当前距离数组和松弛结果。旧状态被弹出时为什么跳过，也要写在表里。状态字段：距离数组、优先队列状态、松弛候选、旧状态判定、不可达哨兵；资源约束题还要把资源维度放进状态。状态升级：把路径枚举压成距离数组和优先队列；每次只扩展当前最有希望的未过期状态。

\textbf{证明和边界。} 证明从扩展顺序开始：当前弹出的状态在代价规则下已经不会被未来路径改得更优。边界：单格、起点或终点很小、必须绕路、多个相同得分路径。变体：也可把阈值二分后检查连通性，或按值从高到低并查集合并。

\noindent\textbf{LeetCode 1334 阈值距离内邻居最少的城市。} 模型：每个城市都需要知道到其他城市的最短距离是否不超过阈值。推导重点：节点数较小时可用 Floyd，边稀疏时也可从每个点跑 Dijkstra。

\textbf{二刷读题。} 读题时先区分可达、最少边数和最小代价。本题可以压成“每个城市都需要知道到其他城市的最短距离是否不超过阈值”，节点数较小时可用 Floyd，边稀疏时也可从每个点跑 Dijkstra决定扩展顺序。先遮住答案，只保留题面，复述候选对象、合法性条件和输出目标。

\textbf{暴力回放。} 慢版本可以枚举路径，或先用普通队列试跑一个小图。它会暴露步数最少和代价最小是否一致；一旦不一致，就必须围绕代价组织状态。二刷时先写慢版本或口述慢版本，用它确认不会漏答案。

\textbf{特例回放。} 拿四个城市和距离阈值 4 手算，每个城市需要统计最短距离不超过 4 的其他城市数；并列时选编号更大的城市。规律是先求全源最短路，再按邻居数量和编号规则选答案。把这个特例继续拆开看：队列或堆弹出的顺序必须和代价规则一致；旧状态被跳过，是因为已经有更小代价到达同一状态。只要边权或代价合成方式改变，就要重新证明扩展顺序。复盘时把这个特例压成状态字段：距离数组、优先队列状态、松弛候选、旧状态判定、不可达哨兵；资源约束题还要把资源维度放进状态；再用边界 距离等于阈值、并列时选编号最大、不可达、无向边 反查初始化、更新顺序和退出条件。

\textbf{状态回放。} 手算路线：手算时记录每次弹出的代价、当前距离数组和松弛结果。旧状态被弹出时为什么跳过，也要写在表里。状态字段：距离数组、优先队列状态、松弛候选、旧状态判定、不可达哨兵；资源约束题还要把资源维度放进状态。状态升级：把路径枚举压成距离数组和优先队列；每次只扩展当前最有希望的未过期状态。

\textbf{证明和边界。} 证明从扩展顺序开始：当前弹出的状态在代价规则下已经不会被未来路径改得更优。边界：距离等于阈值、并列时选编号最大、不可达、无向边。变体：这是多源最短路选择问题，算法取决于 n、边数和查询密度。

\noindent\textbf{LeetCode 1368 使网格图至少有一条有效路径的最小代价。} 模型：沿着格子箭头走代价为零，改方向走代价为一。推导重点：边权只有零和一，用 0-1 BFS：零代价邻居入队首，一代价邻居入队尾。

\textbf{二刷读题。} 读题时先区分可达、最少边数和最小代价。本题可以压成“沿着格子箭头走代价为零，改方向走代价为一”，边权只有零和一，用 0-1 BFS：零代价邻居入队首，一代价邻居入队尾决定扩展顺序。先遮住答案，只保留题面，复述候选对象、合法性条件和输出目标。

\textbf{暴力回放。} 慢版本可以枚举路径，或先用普通队列试跑一个小图。它会暴露步数最少和代价最小是否一致；一旦不一致，就必须围绕代价组织状态。二刷时先写慢版本或口述慢版本，用它确认不会漏答案。

\textbf{特例回放。} 拿网格箭头指向右和下手算，沿箭头走的代价是 0，改方向才付 1。规律是边权只有 0 和 1，可用 0-1 BFS；状态是格子，松弛代价取决于当前箭头是否指向邻居。把这个特例继续拆开看：队列或堆弹出的顺序必须和代价规则一致；旧状态被跳过，是因为已经有更小代价到达同一状态。只要边权或代价合成方式改变，就要重新证明扩展顺序。复盘时把这个特例压成状态字段：距离数组、优先队列状态、松弛候选、旧状态判定、不可达哨兵；资源约束题还要把资源维度放进状态；再用边界 单格、箭头指向越界、多个零代价路径、过期状态 反查初始化、更新顺序和退出条件。

\textbf{状态回放。} 手算路线：手算时记录每次弹出的代价、当前距离数组和松弛结果。旧状态被弹出时为什么跳过，也要写在表里。状态字段：距离数组、优先队列状态、松弛候选、旧状态判定、不可达哨兵；资源约束题还要把资源维度放进状态。状态升级：把路径枚举压成距离数组和优先队列；每次只扩展当前最有希望的未过期状态。

\textbf{证明和边界。} 证明从扩展顺序开始：当前弹出的状态在代价规则下已经不会被未来路径改得更优。边界：单格、箭头指向越界、多个零代价路径、过期状态。变体：普通 BFS 按步数扩展，不等于总代价最小。

\noindent\textbf{LeetCode 1514 概率最大的路径。} 模型：路径概率是边概率乘积，目标是最大乘积路径。推导重点：用大顶堆每次扩展当前概率最高节点，乘以边概率得到候选概率。

\textbf{二刷读题。} 读题时先区分可达、最少边数和最小代价。本题可以压成“路径概率是边概率乘积，目标是最大乘积路径”，用大顶堆每次扩展当前概率最高节点，乘以边概率得到候选概率决定扩展顺序。先遮住答案，只保留题面，复述候选对象、合法性条件和输出目标。

\textbf{暴力回放。} 慢版本可以枚举路径，或先用普通队列试跑一个小图。它会暴露步数最少和代价最小是否一致；一旦不一致，就必须围绕代价组织状态。二刷时先写慢版本或口述慢版本，用它确认不会漏答案。

\textbf{特例回放。} 拿 0->1 概率 0.5、1->2 概率 0.5、0->2 概率 0.2 手算，经 1 到 2 的概率 0.25 更大。规律是把距离改成最大概率，优先队列每次弹出当前概率最高的节点并做乘法松弛。把这个特例继续拆开看：队列或堆弹出的顺序必须和代价规则一致；旧状态被跳过，是因为已经有更小代价到达同一状态。只要边权或代价合成方式改变，就要重新证明扩展顺序。复盘时把这个特例压成状态字段：距离数组、优先队列状态、松弛候选、旧状态判定、不可达哨兵；资源约束题还要把资源维度放进状态；再用边界 不可达、概率为零、起点等于终点、浮点比较 反查初始化、更新顺序和退出条件。

\textbf{状态回放。} 手算路线：手算时记录每次弹出的代价、当前距离数组和松弛结果。旧状态被弹出时为什么跳过，也要写在表里。状态字段：距离数组、优先队列状态、松弛候选、旧状态判定、不可达哨兵；资源约束题还要把资源维度放进状态。状态升级：把路径枚举压成距离数组和优先队列；每次只扩展当前最有希望的未过期状态。

\textbf{证明和边界。} 证明从扩展顺序开始：当前弹出的状态在代价规则下已经不会被未来路径改得更优。边界：不可达、概率为零、起点等于终点、浮点比较。变体：取负对数后可转成最短路，但直接最大堆更直观。

\noindent\textbf{LeetCode 1631 最小体力消耗路径。} 模型：路径代价是沿途最大高度差，不是边权和。推导重点：Dijkstra 的松弛改成取当前代价和新边代价的最大值。

\textbf{二刷读题。} 读题时先区分可达、最少边数和最小代价。本题可以压成“路径代价是沿途最大高度差，不是边权和”，Dijkstra 的松弛改成取当前代价和新边代价的最大值决定扩展顺序。先遮住答案，只保留题面，复述候选对象、合法性条件和输出目标。

\textbf{暴力回放。} 慢版本可以枚举路径，或先用普通队列试跑一个小图。它会暴露步数最少和代价最小是否一致；一旦不一致，就必须围绕代价组织状态。二刷时先写慢版本或口述慢版本，用它确认不会漏答案。

\textbf{特例回放。} 拿 heights=[[1,2,2],[3,8,2],[5,3,5]] 手算，路径代价是相邻高度差的最大值，不是差值之和。规律是 Dijkstra 的距离定义为到达某格的最小最大边权，或二分最大体力做可达性。把这个特例继续拆开看：队列或堆弹出的顺序必须和代价规则一致；旧状态被跳过，是因为已经有更小代价到达同一状态。只要边权或代价合成方式改变，就要重新证明扩展顺序。复盘时把这个特例压成状态字段：距离数组、优先队列状态、松弛候选、旧状态判定、不可达哨兵；资源约束题还要把资源维度放进状态；再用边界 单格、平坦网格、必须绕路、多个相同代价 反查初始化、更新顺序和退出条件。

\textbf{状态回放。} 手算路线：手算时记录每次弹出的代价、当前距离数组和松弛结果。旧状态被弹出时为什么跳过，也要写在表里。状态字段：距离数组、优先队列状态、松弛候选、旧状态判定、不可达哨兵；资源约束题还要把资源维度放进状态。状态升级：把路径枚举压成距离数组和优先队列；每次只扩展当前最有希望的未过期状态。

\textbf{证明和边界。} 证明从扩展顺序开始：当前弹出的状态在代价规则下已经不会被未来路径改得更优。边界：单格、平坦网格、必须绕路、多个相同代价。变体：也可二分体力上限后 BFS 检查可达。

\noindent\textbf{LeetCode 1928 规定时间内到达终点的最小花费。} 模型：路径同时有时间约束和费用目标，单一距离无法描述状态。推导重点：状态包含节点和已用时间，DP 或 Dijkstra 在时间维度内松弛最小费用。

\textbf{二刷读题。} 读题时先区分可达、最少边数和最小代价。本题可以压成“路径同时有时间约束和费用目标，单一距离无法描述状态”，状态包含节点和已用时间，DP 或 Dijkstra 在时间维度内松弛最小费用决定扩展顺序。先遮住答案，只保留题面，复述候选对象、合法性条件和输出目标。

\textbf{暴力回放。} 慢版本可以枚举路径，或先用普通队列试跑一个小图。它会暴露步数最少和代价最小是否一致；一旦不一致，就必须围绕代价组织状态。二刷时先写慢版本或口述慢版本，用它确认不会漏答案。

\textbf{特例回放。} 拿 0->1 用时 10 费用 5，1->2 用时 10 费用 5，0->2 用时 25 费用 100、maxTime=30 手算，费用最小的是经 1 到 2。规律是每个节点不能只保存一个最小费用，状态要包含已用时间。把这个特例继续拆开看：队列或堆弹出的顺序必须和代价规则一致；旧状态被跳过，是因为已经有更小代价到达同一状态。只要边权或代价合成方式改变，就要重新证明扩展顺序。复盘时把这个特例压成状态字段：距离数组、优先队列状态、松弛候选、旧状态判定、不可达哨兵；资源约束题还要把资源维度放进状态；再用边界 时间刚好到达、费用高但时间短、不可达、状态数量由 maxTime 决定 反查初始化、更新顺序和退出条件。

\textbf{状态回放。} 手算路线：手算时记录每次弹出的代价、当前距离数组和松弛结果。旧状态被弹出时为什么跳过，也要写在表里。状态字段：距离数组、优先队列状态、松弛候选、旧状态判定、不可达哨兵；资源约束题还要把资源维度放进状态。状态升级：把路径枚举压成距离数组和优先队列；每次只扩展当前最有希望的未过期状态。

\textbf{证明和边界。} 证明从扩展顺序开始：当前弹出的状态在代价规则下已经不会被未来路径改得更优。边界：时间刚好到达、费用高但时间短、不可达、状态数量由 maxTime 决定。变体：这是资源约束最短路，不能只保存每个节点一个最小费用。

\noindent\textbf{LeetCode 1976 到达目的地的方案数。} 模型：需要在最短距离之外统计有多少条最短路径。推导重点：Dijkstra 松弛时若得到更短距离就重置方案数，若等于最短距离就累加方案数。

\textbf{二刷读题。} 读题时先区分可达、最少边数和最小代价。本题可以压成“需要在最短距离之外统计有多少条最短路径”，Dijkstra 松弛时若得到更短距离就重置方案数，若等于最短距离就累加方案数决定扩展顺序。先遮住答案，只保留题面，复述候选对象、合法性条件和输出目标。

\textbf{暴力回放。} 慢版本可以枚举路径，或先用普通队列试跑一个小图。它会暴露步数最少和代价最小是否一致；一旦不一致，就必须围绕代价组织状态。二刷时先写慢版本或口述慢版本，用它确认不会漏答案。

\textbf{特例回放。} 拿 0->1->3 和 0->2->3 两条长度都为 2 的路手算，最短距离相同，所以到 3 的方案数应累加为 2。规律是 Dijkstra 松弛到更短距离时重置方案数，等长时累加方案数。把这个特例继续拆开看：队列或堆弹出的顺序必须和代价规则一致；旧状态被跳过，是因为已经有更小代价到达同一状态。只要边权或代价合成方式改变，就要重新证明扩展顺序。复盘时把这个特例压成状态字段：距离数组、优先队列状态、松弛候选、旧状态判定、不可达哨兵；资源约束题还要把资源维度放进状态；再用边界 多条等长路径、取模、旧堆状态、不可达 反查初始化、更新顺序和退出条件。

\textbf{状态回放。} 手算路线：手算时记录每次弹出的代价、当前距离数组和松弛结果。旧状态被弹出时为什么跳过，也要写在表里。状态字段：距离数组、优先队列状态、松弛候选、旧状态判定、不可达哨兵；资源约束题还要把资源维度放进状态。状态升级：把路径枚举压成距离数组和优先队列；每次只扩展当前最有希望的未过期状态。

\textbf{证明和边界。} 证明从扩展顺序开始：当前弹出的状态在代价规则下已经不会被未来路径改得更优。边界：多条等长路径、取模、旧堆状态、不可达。变体：这是最短路和计数 DP 的结合，更新顺序必须和距离松弛一致。

\noindent\textbf{LeetCode 2045 到达目的地的第二短时间。} 模型：每个节点需要保留最短和次短到达步数，再考虑红绿灯等待时间。推导重点：BFS 式扩展边数或时间，只有候选时间不同且能进入前两名时才入队。

\textbf{二刷读题。} 读题时先区分可达、最少边数和最小代价。本题可以压成“每个节点需要保留最短和次短到达步数，再考虑红绿灯等待时间”，BFS 式扩展边数或时间，只有候选时间不同且能进入前两名时才入队决定扩展顺序。先遮住答案，只保留题面，复述候选对象、合法性条件和输出目标。

\textbf{暴力回放。} 慢版本可以枚举路径，或先用普通队列试跑一个小图。它会暴露步数最少和代价最小是否一致；一旦不一致，就必须围绕代价组织状态。二刷时先写慢版本或口述慢版本，用它确认不会漏答案。

\textbf{特例回放。} 拿 1--2、time=3、change=2 手算，第一次到 2 是最短；第二短必须从 2 回 1 再到 2，且出发时可能遇到红灯等待。规律是每个节点保存两个不同到达时间，不能把次短等同于最短。把这个特例继续拆开看：队列或堆弹出的顺序必须和代价规则一致；旧状态被跳过，是因为已经有更小代价到达同一状态。只要边权或代价合成方式改变，就要重新证明扩展顺序。复盘时把这个特例压成状态字段：距离数组、优先队列状态、松弛候选、旧状态判定、不可达哨兵；资源约束题还要把资源维度放进状态；再用边界 第二短不能等于最短、红灯等待、无向图往返、到达终点后继续找次短 反查初始化、更新顺序和退出条件。

\textbf{状态回放。} 手算路线：手算时记录每次弹出的代价、当前距离数组和松弛结果。旧状态被弹出时为什么跳过，也要写在表里。状态字段：距离数组、优先队列状态、松弛候选、旧状态判定、不可达哨兵；资源约束题还要把资源维度放进状态。状态升级：把路径枚举压成距离数组和优先队列；每次只扩展当前最有希望的未过期状态。

\textbf{证明和边界。} 证明从扩展顺序开始：当前弹出的状态在代价规则下已经不会被未来路径改得更优。边界：第二短不能等于最短、红灯等待、无向图往返、到达终点后继续找次短。变体：它说明最短路状态有时不是一个最优值，而是前两个不同值。

\subsection{自测标准}

二刷时不要只确认自己见过这道题。请遮住答案，先讲题面模型，再讲暴力基线和瓶颈，然后说出优化结构保存了什么、不变量如何排除错误候选、边界实验如何打破错误实现。

\section{并查集与动态连通性：二刷复盘卡}
这一大节只围绕一个题型复盘，不把题号直接铺成目录。阅读时先看复盘目标，再读核心题卡，最后用自测标准检查自己是否能独立重讲。

\subsection{复盘目标}

追问通常会问并查集能否回答路径长度、路径压缩是否改变集合语义、字符串如何编号。请准备说明合并失败为什么意味着环。

\subsection{核心题卡}

\noindent\textbf{LeetCode 305 岛屿数量 II。} 模型：陆地逐步加入，每次加入只可能影响相邻陆地所在集合。推导重点：新增陆地先让岛屿数加一，再和四邻已存在陆地合并，合并成功则岛屿数减一。

\textbf{二刷读题。} 读题时先确认关系是否只增加不删除，且问题只关心同组关系。本题可以压成“陆地逐步加入，每次加入只可能影响相邻陆地所在集合”，并查集的代表元就是集合身份。先遮住答案，只保留题面，复述候选对象、合法性条件和输出目标。

\textbf{暴力回放。} 慢版本每次合并后用 DFS 或 BFS 重新判断连通性。它正确但重复遍历已有连接；当关系只会合并不会拆开时，可以压成集合代表。二刷时先写慢版本或口述慢版本，用它确认不会漏答案。

\textbf{特例回放。} 拿 3x3 网格依次加 (0,0)、(0,1)、(1,2) 手算，前两个陆地相邻要合并，岛屿数从 2 降回 1；第三个不相邻，岛屿数加 1。规律是每次加陆地先新增一个集合，再和四邻陆地 union，重复加点不能重复计数。把这个特例继续拆开看：union 只是在维护等价类，不是在保存具体路径。两个节点代表元相同只能说明连通或关系一致；如果题目还要距离、比例或奇偶性，必须把附加信息挂到根关系上。复盘时把这个特例压成状态字段：parent、size 或 rank、find 返回的代表元、合并时的附加信息；带权或奇偶并查集还要维护到根的关系；再用边界 重复加入同一格、边界位置、四邻属于同一岛、空操作列表 反查初始化、更新顺序和退出条件。

\textbf{状态回放。} 手算路线：手算时写 parent 表和集合代表。每次 union 只改变根之间的关系，find 后代表相同才说明连通。状态字段：parent、size 或 rank、find 返回的代表元、合并时的附加信息；带权或奇偶并查集还要维护到根的关系。状态升级：把连通分量压成代表元；查询只比较代表，合并只发生在两个根之间。

\textbf{证明和边界。} 证明从等价关系开始：合并维护连通性的传递性，find 返回的代表元就是当前集合身份。边界：重复加入同一格、边界位置、四邻属于同一岛、空操作列表。变体：这是动态连通性，比每次重新 BFS 更适合并查集。

\noindent\textbf{LeetCode 399 除法求值。} 模型：变量之间的除法关系可以看成带权连通分量。推导重点：并查集维护节点到根的乘积权值，合并时根据等式调整根之间权重。

\textbf{二刷读题。} 读题时先确认关系是否只增加不删除，且问题只关心同组关系。本题可以压成“变量之间的除法关系可以看成带权连通分量”，并查集的代表元就是集合身份。先遮住答案，只保留题面，复述候选对象、合法性条件和输出目标。

\textbf{暴力回放。} 慢版本每次合并后用 DFS 或 BFS 重新判断连通性。它正确但重复遍历已有连接；当关系只会合并不会拆开时，可以压成集合代表。二刷时先写慢版本或口述慢版本，用它确认不会漏答案。

\textbf{特例回放。} 拿 a/b=2、b/c=3 手算，a 到 c 的比值是沿边相乘 6；若查询 c/a，则沿反向边乘 1/3 和 1/2。规律是并查集或图搜索都要维护到代表元的倍率关系，合并时同步更新权值。把这个特例继续拆开看：union 只是在维护等价类，不是在保存具体路径。两个节点代表元相同只能说明连通或关系一致；如果题目还要距离、比例或奇偶性，必须把附加信息挂到根关系上。复盘时把这个特例压成状态字段：parent、size 或 rank、find 返回的代表元、合并时的附加信息；带权或奇偶并查集还要维护到根的关系；再用边界 查询未知变量、自除、矛盾数据、路径压缩时权值同步更新 反查初始化、更新顺序和退出条件。

\textbf{状态回放。} 手算路线：手算时写 parent 表和集合代表。每次 union 只改变根之间的关系，find 后代表相同才说明连通。状态字段：parent、size 或 rank、find 返回的代表元、合并时的附加信息；带权或奇偶并查集还要维护到根的关系。状态升级：把连通分量压成代表元；查询只比较代表，合并只发生在两个根之间。

\textbf{证明和边界。} 证明从等价关系开始：合并维护连通性的传递性，find 返回的代表元就是当前集合身份。边界：查询未知变量、自除、矛盾数据、路径压缩时权值同步更新。变体：也可建图 DFS，但多次查询时带权并查集更适合。

\noindent\textbf{LeetCode 547 省份数量。} 模型：城市连接关系形成无向图，省份就是连通分量。推导重点：并查集合并所有直接连接的城市，最后统计根数量。

\textbf{二刷读题。} 读题时先确认关系是否只增加不删除，且问题只关心同组关系。本题可以压成“城市连接关系形成无向图，省份就是连通分量”，并查集的代表元就是集合身份。先遮住答案，只保留题面，复述候选对象、合法性条件和输出目标。

\textbf{暴力回放。} 慢版本每次合并后用 DFS 或 BFS 重新判断连通性。它正确但重复遍历已有连接；当关系只会合并不会拆开时，可以压成集合代表。二刷时先写慢版本或口述慢版本，用它确认不会漏答案。

\textbf{特例回放。} 拿 isConnected 中 0 和 1 相连、2 独立手算，0 和 1 union 后属于同一代表，2 保持自己，省份数为 2。规律是矩阵中的一条连接边对应集合合并，最后代表元数量就是连通分量数。把这个特例继续拆开看：union 只是在维护等价类，不是在保存具体路径。两个节点代表元相同只能说明连通或关系一致；如果题目还要距离、比例或奇偶性，必须把附加信息挂到根关系上。复盘时把这个特例压成状态字段：parent、size 或 rank、find 返回的代表元、合并时的附加信息；带权或奇偶并查集还要维护到根的关系；再用边界 邻接矩阵对称、自连接、孤立城市、只扫描上三角 反查初始化、更新顺序和退出条件。

\textbf{状态回放。} 手算路线：手算时写 parent 表和集合代表。每次 union 只改变根之间的关系，find 后代表相同才说明连通。状态字段：parent、size 或 rank、find 返回的代表元、合并时的附加信息；带权或奇偶并查集还要维护到根的关系。状态升级：把连通分量压成代表元；查询只比较代表，合并只发生在两个根之间。

\textbf{证明和边界。} 证明从等价关系开始：合并维护连通性的传递性，find 返回的代表元就是当前集合身份。边界：邻接矩阵对称、自连接、孤立城市、只扫描上三角。变体：BFS 也可做；并查集更突出动态合并。

\noindent\textbf{LeetCode 684 冗余连接。} 模型：树加一条边会形成唯一环。推导重点：按顺序加边，若一条边两端已连通，则它就是冗余边。

\textbf{二刷读题。} 读题时先确认关系是否只增加不删除，且问题只关心同组关系。本题可以压成“树加一条边会形成唯一环”，并查集的代表元就是集合身份。先遮住答案，只保留题面，复述候选对象、合法性条件和输出目标。

\textbf{暴力回放。} 慢版本每次合并后用 DFS 或 BFS 重新判断连通性。它正确但重复遍历已有连接；当关系只会合并不会拆开时，可以压成集合代表。二刷时先写慢版本或口述慢版本，用它确认不会漏答案。

\textbf{特例回放。} 拿 edges=[[1,2],[1,3],[2,3]] 手算，前两条边把 1、2、3 连成一棵树，第三条边连接的两个点已经同组，所以它就是冗余边。规律是无向图加边时，第一次 union 失败的边形成环。把这个特例继续拆开看：union 只是在维护等价类，不是在保存具体路径。两个节点代表元相同只能说明连通或关系一致；如果题目还要距离、比例或奇偶性，必须把附加信息挂到根关系上。复盘时把这个特例压成状态字段：parent、size 或 rank、find 返回的代表元、合并时的附加信息；带权或奇偶并查集还要维护到根的关系；再用边界 节点编号、最后一条成环边、重复边、并查集初始化大小 反查初始化、更新顺序和退出条件。

\textbf{状态回放。} 手算路线：手算时写 parent 表和集合代表。每次 union 只改变根之间的关系，find 后代表相同才说明连通。状态字段：parent、size 或 rank、find 返回的代表元、合并时的附加信息；带权或奇偶并查集还要维护到根的关系。状态升级：把连通分量压成代表元；查询只比较代表，合并只发生在两个根之间。

\textbf{证明和边界。} 证明从等价关系开始：合并维护连通性的传递性，find 返回的代表元就是当前集合身份。边界：节点编号、最后一条成环边、重复边、并查集初始化大小。变体：有向冗余连接需要同时处理入度为二和环，复杂度更高。

\noindent\textbf{LeetCode 685 冗余连接 II。} 模型：有向图中额外边可能导致某节点入度为二，也可能导致环。推导重点：先记录入度为二的两条候选边，再用并查集排除其中一条测试是否成树。

\textbf{二刷读题。} 读题时先确认关系是否只增加不删除，且问题只关心同组关系。本题可以压成“有向图中额外边可能导致某节点入度为二，也可能导致环”，并查集的代表元就是集合身份。先遮住答案，只保留题面，复述候选对象、合法性条件和输出目标。

\textbf{暴力回放。} 慢版本每次合并后用 DFS 或 BFS 重新判断连通性。它正确但重复遍历已有连接；当关系只会合并不会拆开时，可以压成集合代表。二刷时先写慢版本或口述慢版本，用它确认不会漏答案。

\textbf{特例回放。} 拿有向边 [[1,2],[1,3],[2,3]] 手算，节点 3 有两个父亲 1 和 2；需要判断去掉哪条边后剩余图是有根树。规律是先处理双父节点候选，再用并查集检测环，双父和成环组合时要按题意返回正确候选边。把这个特例继续拆开看：union 只是在维护等价类，不是在保存具体路径。两个节点代表元相同只能说明连通或关系一致；如果题目还要距离、比例或奇偶性，必须把附加信息挂到根关系上。复盘时把这个特例压成状态字段：parent、size 或 rank、find 返回的代表元、合并时的附加信息；带权或奇偶并查集还要维护到根的关系；再用边界 只有环、只有入度二、两者同时存在、返回最后出现的边 反查初始化、更新顺序和退出条件。

\textbf{状态回放。} 手算路线：手算时写 parent 表和集合代表。每次 union 只改变根之间的关系，find 后代表相同才说明连通。状态字段：parent、size 或 rank、find 返回的代表元、合并时的附加信息；带权或奇偶并查集还要维护到根的关系。状态升级：把连通分量压成代表元；查询只比较代表，合并只发生在两个根之间。

\textbf{证明和边界。} 证明从等价关系开始：合并维护连通性的传递性，find 返回的代表元就是当前集合身份。边界：只有环、只有入度二、两者同时存在、返回最后出现的边。变体：它比无向冗余连接多了有向树入度约束。

\noindent\textbf{LeetCode 721 账户合并。} 模型：共享邮箱说明账户属于同一人，邮箱之间形成连通分量。推导重点：给邮箱编号并查集合并，同根邮箱收集后排序输出。

\textbf{二刷读题。} 读题时先确认关系是否只增加不删除，且问题只关心同组关系。本题可以压成“共享邮箱说明账户属于同一人，邮箱之间形成连通分量”，并查集的代表元就是集合身份。先遮住答案，只保留题面，复述候选对象、合法性条件和输出目标。

\textbf{暴力回放。} 慢版本每次合并后用 DFS 或 BFS 重新判断连通性。它正确但重复遍历已有连接；当关系只会合并不会拆开时，可以压成集合代表。二刷时先写慢版本或口述慢版本，用它确认不会漏答案。

\textbf{特例回放。} 拿 John 的账号 [a,b] 和另一个 [b,c] 手算，邮箱 b 重叠，所以 a、b、c 属于同一人；名字只是输出标签，合并依据是邮箱。规律是邮箱映射到编号后并查集合并，最后按代表收集并排序邮箱。把这个特例继续拆开看：union 只是在维护等价类，不是在保存具体路径。两个节点代表元相同只能说明连通或关系一致；如果题目还要距离、比例或奇偶性，必须把附加信息挂到根关系上。复盘时把这个特例压成状态字段：parent、size 或 rank、find 返回的代表元、合并时的附加信息；带权或奇偶并查集还要维护到根的关系；再用边界 同名不同人、一个账户多个邮箱、重复邮箱、输出排序 反查初始化、更新顺序和退出条件。

\textbf{状态回放。} 手算路线：手算时写 parent 表和集合代表。每次 union 只改变根之间的关系，find 后代表相同才说明连通。状态字段：parent、size 或 rank、find 返回的代表元、合并时的附加信息；带权或奇偶并查集还要维护到根的关系。状态升级：把连通分量压成代表元；查询只比较代表，合并只发生在两个根之间。

\textbf{证明和边界。} 证明从等价关系开始：合并维护连通性的传递性，find 返回的代表元就是当前集合身份。边界：同名不同人、一个账户多个邮箱、重复邮箱、输出排序。变体：姓名不能作为合并依据，只能作为输出字段。

\noindent\textbf{LeetCode 947 移除最多的同行或同列石头。} 模型：同一行或同一列的石头属于同一连通分量，每个分量至少留一块。推导重点：把行和列编号后合并，答案是石头数减去连通分量数。

\textbf{二刷读题。} 读题时先确认关系是否只增加不删除，且问题只关心同组关系。本题可以压成“同一行或同一列的石头属于同一连通分量，每个分量至少留一块”，并查集的代表元就是集合身份。先遮住答案，只保留题面，复述候选对象、合法性条件和输出目标。

\textbf{暴力回放。} 慢版本每次合并后用 DFS 或 BFS 重新判断连通性。它正确但重复遍历已有连接；当关系只会合并不会拆开时，可以压成集合代表。二刷时先写慢版本或口述慢版本，用它确认不会漏答案。

\textbf{特例回放。} 拿石头 (0,0),(0,1),(1,0) 手算，它们通过同一行或同一列连成一个分量，三个石头最多移除 2 个。规律是每个连通分量至少留一个，答案是石头数减分量数。把这个特例继续拆开看：union 只是在维护等价类，不是在保存具体路径。两个节点代表元相同只能说明连通或关系一致；如果题目还要距离、比例或奇偶性，必须把附加信息挂到根关系上。复盘时把这个特例压成状态字段：parent、size 或 rank、find 返回的代表元、合并时的附加信息；带权或奇偶并查集还要维护到根的关系；再用边界 单块石头、行列编号冲突、多个分量、重复坐标 反查初始化、更新顺序和退出条件。

\textbf{状态回放。} 手算路线：手算时写 parent 表和集合代表。每次 union 只改变根之间的关系，find 后代表相同才说明连通。状态字段：parent、size 或 rank、find 返回的代表元、合并时的附加信息；带权或奇偶并查集还要维护到根的关系。状态升级：把连通分量压成代表元；查询只比较代表，合并只发生在两个根之间。

\textbf{证明和边界。} 证明从等价关系开始：合并维护连通性的传递性，find 返回的代表元就是当前集合身份。边界：单块石头、行列编号冲突、多个分量、重复坐标。变体：也可把石头当节点建图，但行列并查集更省边。

\noindent\textbf{LeetCode 959 由斜杠划分区域。} 模型：每个网格可拆成四个小三角，斜杠决定内部哪些三角相连。推导重点：给每格四个区域编号，先按字符合并内部，再合并相邻格边界区域。

\textbf{二刷读题。} 读题时先确认关系是否只增加不删除，且问题只关心同组关系。本题可以压成“每个网格可拆成四个小三角，斜杠决定内部哪些三角相连”，并查集的代表元就是集合身份。先遮住答案，只保留题面，复述候选对象、合法性条件和输出目标。

\textbf{暴力回放。} 慢版本每次合并后用 DFS 或 BFS 重新判断连通性。它正确但重复遍历已有连接；当关系只会合并不会拆开时，可以压成集合代表。二刷时先写慢版本或口述慢版本，用它确认不会漏答案。

\textbf{特例回放。} 拿 1x1 格子里一条斜杠手算，它把小方格分成两个区域；把每格拆成 4 个三角形后，斜杠决定哪些三角形不能合并。规律是格内按字符合并，格间相邻边合并，最终代表元数量就是区域数。把这个特例继续拆开看：union 只是在维护等价类，不是在保存具体路径。两个节点代表元相同只能说明连通或关系一致；如果题目还要距离、比例或奇偶性，必须把附加信息挂到根关系上。复盘时把这个特例压成状态字段：parent、size 或 rank、find 返回的代表元、合并时的附加信息；带权或奇偶并查集还要维护到根的关系；再用边界 空格、正斜杠、反斜杠、边界合并、区域计数 反查初始化、更新顺序和退出条件。

\textbf{状态回放。} 手算路线：手算时写 parent 表和集合代表。每次 union 只改变根之间的关系，find 后代表相同才说明连通。状态字段：parent、size 或 rank、find 返回的代表元、合并时的附加信息；带权或奇偶并查集还要维护到根的关系。状态升级：把连通分量压成代表元；查询只比较代表，合并只发生在两个根之间。

\textbf{证明和边界。} 证明从等价关系开始：合并维护连通性的传递性，find 返回的代表元就是当前集合身份。边界：空格、正斜杠、反斜杠、边界合并、区域计数。变体：这是把几何连通性离散成并查集元素的典型题。

\noindent\textbf{LeetCode 990 等式方程的可满足性。} 模型：相等关系形成集合，不等关系要求两端不在同一集合。推导重点：先合并所有等式，再检查每个不等式是否冲突。

\textbf{二刷读题。} 读题时先确认关系是否只增加不删除，且问题只关心同组关系。本题可以压成“相等关系形成集合，不等关系要求两端不在同一集合”，并查集的代表元就是集合身份。先遮住答案，只保留题面，复述候选对象、合法性条件和输出目标。

\textbf{暴力回放。} 慢版本每次合并后用 DFS 或 BFS 重新判断连通性。它正确但重复遍历已有连接；当关系只会合并不会拆开时，可以压成集合代表。二刷时先写慢版本或口述慢版本，用它确认不会漏答案。

\textbf{特例回放。} 拿 a==b、b==c、a!=c 手算，前两条等式把 a,b,c 合并，最后不等式发现二者同组，矛盾。规律是先处理所有等式建立集合，再检查不等式；顺序反过来会漏掉后续合并造成的冲突。把这个特例继续拆开看：union 只是在维护等价类，不是在保存具体路径。两个节点代表元相同只能说明连通或关系一致；如果题目还要距离、比例或奇偶性，必须把附加信息挂到根关系上。复盘时把这个特例压成状态字段：parent、size 或 rank、find 返回的代表元、合并时的附加信息；带权或奇偶并查集还要维护到根的关系；再用边界 自等式、自不等式、传递相等、多组变量 反查初始化、更新顺序和退出条件。

\textbf{状态回放。} 手算路线：手算时写 parent 表和集合代表。每次 union 只改变根之间的关系，find 后代表相同才说明连通。状态字段：parent、size 或 rank、find 返回的代表元、合并时的附加信息；带权或奇偶并查集还要维护到根的关系。状态升级：把连通分量压成代表元；查询只比较代表，合并只发生在两个根之间。

\textbf{证明和边界。} 证明从等价关系开始：合并维护连通性的传递性，find 返回的代表元就是当前集合身份。边界：自等式、自不等式、传递相等、多组变量。变体：先处理不等式会漏掉后续等式带来的传递冲突。

\noindent\textbf{LeetCode 1061 按字典序排列最小的等效字符串。} 模型：等价字符形成集合，每个集合应选择字典序最小代表。推导重点：合并两个字符集合时让较小根成为代表，扫描目标串时替换为根。

\textbf{二刷读题。} 读题时先确认关系是否只增加不删除，且问题只关心同组关系。本题可以压成“等价字符形成集合，每个集合应选择字典序最小代表”，并查集的代表元就是集合身份。先遮住答案，只保留题面，复述候选对象、合法性条件和输出目标。

\textbf{暴力回放。} 慢版本每次合并后用 DFS 或 BFS 重新判断连通性。它正确但重复遍历已有连接；当关系只会合并不会拆开时，可以压成集合代表。二刷时先写慢版本或口述慢版本，用它确认不会漏答案。

\textbf{特例回放。} 拿 s1=parker、s2=morris、base=parser 手算，p 和 m 等价，a 和 o 等价；每个等价类输出字典序最小字符。规律是并查集合并字符时让较小根做代表，转换 base 时查代表。把这个特例继续拆开看：union 只是在维护等价类，不是在保存具体路径。两个节点代表元相同只能说明连通或关系一致；如果题目还要距离、比例或奇偶性，必须把附加信息挂到根关系上。复盘时把这个特例压成状态字段：parent、size 或 rank、find 返回的代表元、合并时的附加信息；带权或奇偶并查集还要维护到根的关系；再用边界 传递等价、重复等价关系、自等价、未出现字符 反查初始化、更新顺序和退出条件。

\textbf{状态回放。} 手算路线：手算时写 parent 表和集合代表。每次 union 只改变根之间的关系，find 后代表相同才说明连通。状态字段：parent、size 或 rank、find 返回的代表元、合并时的附加信息；带权或奇偶并查集还要维护到根的关系。状态升级：把连通分量压成代表元；查询只比较代表，合并只发生在两个根之间。

\textbf{证明和边界。} 证明从等价关系开始：合并维护连通性的传递性，find 返回的代表元就是当前集合身份。边界：传递等价、重复等价关系、自等价、未出现字符。变体：如果集合代表还要保存其他权值，合并规则需要同步维护。

\noindent\textbf{LeetCode 1202 交换字符串中的元素。} 模型：可交换下标形成连通分量，同一分量内字符可以任意重排。推导重点：并查集合并所有可交换下标，对每个分量收集字符并按字典序分配回最小下标。

\textbf{二刷读题。} 读题时先确认关系是否只增加不删除，且问题只关心同组关系。本题可以压成“可交换下标形成连通分量，同一分量内字符可以任意重排”，并查集的代表元就是集合身份。先遮住答案，只保留题面，复述候选对象、合法性条件和输出目标。

\textbf{暴力回放。} 慢版本每次合并后用 DFS 或 BFS 重新判断连通性。它正确但重复遍历已有连接；当关系只会合并不会拆开时，可以压成集合代表。二刷时先写慢版本或口述慢版本，用它确认不会漏答案。

\textbf{特例回放。} 拿 s=dcab、pairs=[[0,3],[1,2]] 手算，0 和 3 可交换成一组，1 和 2 可交换成一组；每组内部字符排序后放回最小位置。规律是并查集找可交换连通分量，每个分量的字符可任意重排。把这个特例继续拆开看：union 只是在维护等价类，不是在保存具体路径。两个节点代表元相同只能说明连通或关系一致；如果题目还要距离、比例或奇偶性，必须把附加信息挂到根关系上。复盘时把这个特例压成状态字段：parent、size 或 rank、find 返回的代表元、合并时的附加信息；带权或奇偶并查集还要维护到根的关系；再用边界 没有交换、多个分量、重复字符、分量下标排序 反查初始化、更新顺序和退出条件。

\textbf{状态回放。} 手算路线：手算时写 parent 表和集合代表。每次 union 只改变根之间的关系，find 后代表相同才说明连通。状态字段：parent、size 或 rank、find 返回的代表元、合并时的附加信息；带权或奇偶并查集还要维护到根的关系。状态升级：把连通分量压成代表元；查询只比较代表，合并只发生在两个根之间。

\textbf{证明和边界。} 证明从等价关系开始：合并维护连通性的传递性，find 返回的代表元就是当前集合身份。边界：没有交换、多个分量、重复字符、分量下标排序。变体：它展示并查集不仅能回答连通性，还能把分量内资源重新组织。

\noindent\textbf{LeetCode 1319 连通网络的操作次数。} 模型：多余网线可以挪用来连接不同连通分量。推导重点：先判断边数至少为 n 减一，再用并查集统计连通分量，答案是分量数减一。

\textbf{二刷读题。} 读题时先确认关系是否只增加不删除，且问题只关心同组关系。本题可以压成“多余网线可以挪用来连接不同连通分量”，并查集的代表元就是集合身份。先遮住答案，只保留题面，复述候选对象、合法性条件和输出目标。

\textbf{暴力回放。} 慢版本每次合并后用 DFS 或 BFS 重新判断连通性。它正确但重复遍历已有连接；当关系只会合并不会拆开时，可以压成集合代表。二刷时先写慢版本或口述慢版本，用它确认不会漏答案。

\textbf{特例回放。} 拿 n=4、连接 [[0,1],[0,2],[1,2]] 手算，前三个点已有一条冗余边，但节点 3 孤立；总边数只有 3，刚好可以拿冗余边连上 3。规律是边数少于 n-1 直接失败，否则答案是连通分量数减一。把这个特例继续拆开看：union 只是在维护等价类，不是在保存具体路径。两个节点代表元相同只能说明连通或关系一致；如果题目还要距离、比例或奇偶性，必须把附加信息挂到根关系上。复盘时把这个特例压成状态字段：parent、size 或 rank、find 返回的代表元、合并时的附加信息；带权或奇偶并查集还要维护到根的关系；再用边界 边数不足、已经连通、重复边、多分量 反查初始化、更新顺序和退出条件。

\textbf{状态回放。} 手算路线：手算时写 parent 表和集合代表。每次 union 只改变根之间的关系，find 后代表相同才说明连通。状态字段：parent、size 或 rank、find 返回的代表元、合并时的附加信息；带权或奇偶并查集还要维护到根的关系。状态升级：把连通分量压成代表元；查询只比较代表，合并只发生在两个根之间。

\textbf{证明和边界。} 证明从等价关系开始：合并维护连通性的传递性，find 返回的代表元就是当前集合身份。边界：边数不足、已经连通、重复边、多分量。变体：合并失败的边代表冗余资源，但只要总边数足够即可。

\noindent\textbf{LeetCode 1579 保证图可完全遍历。} 模型：Alice 和 Bob 各自需要图连通，公共边应优先服务两人。推导重点：先处理类型三公共边同步合并两个并查集，再分别处理各自边，最后检查两图连通。

\textbf{二刷读题。} 读题时先确认关系是否只增加不删除，且问题只关心同组关系。本题可以压成“Alice 和 Bob 各自需要图连通，公共边应优先服务两人”，并查集的代表元就是集合身份。先遮住答案，只保留题面，复述候选对象、合法性条件和输出目标。

\textbf{暴力回放。} 慢版本每次合并后用 DFS 或 BFS 重新判断连通性。它正确但重复遍历已有连接；当关系只会合并不会拆开时，可以压成集合代表。二刷时先写慢版本或口述慢版本，用它确认不会漏答案。

\textbf{特例回放。} 拿公共边能同时连接 Alice 和 Bob 的两个分量手算，先用类型 3 边可以让两个人都受益；若先处理私有边，公共边可能变成冗余。规律是先处理公共边，再分别处理 Alice 和 Bob 的边，最后两套并查集都必须连通。把这个特例继续拆开看：union 只是在维护等价类，不是在保存具体路径。两个节点代表元相同只能说明连通或关系一致；如果题目还要距离、比例或奇偶性，必须把附加信息挂到根关系上。复盘时把这个特例压成状态字段：parent、size 或 rank、find 返回的代表元、合并时的附加信息；带权或奇偶并查集还要维护到根的关系；再用边界 公共边冗余、某一方不连通、边数很多、节点从一开始编号 反查初始化、更新顺序和退出条件。

\textbf{状态回放。} 手算路线：手算时写 parent 表和集合代表。每次 union 只改变根之间的关系，find 后代表相同才说明连通。状态字段：parent、size 或 rank、find 返回的代表元、合并时的附加信息；带权或奇偶并查集还要维护到根的关系。状态升级：把连通分量压成代表元；查询只比较代表，合并只发生在两个根之间。

\textbf{证明和边界。} 证明从等价关系开始：合并维护连通性的传递性，find 返回的代表元就是当前集合身份。边界：公共边冗余、某一方不连通、边数很多、节点从一开始编号。变体：这是双并查集和资源优先级的组合。

\subsection{自测标准}

二刷时不要只确认自己见过这道题。请遮住答案，先讲题面模型，再讲暴力基线和瓶颈，然后说出优化结构保存了什么、不变量如何排除错误候选、边界实验如何打破错误实现。

\section{回溯、枚举与剪枝：二刷复盘卡}
这一大节只围绕一个题型复盘，不把题号直接铺成目录。阅读时先看复盘目标，再读核心题卡，最后用自测标准检查自己是否能独立重讲。

\subsection{复盘目标}

追问通常会问去重发生在同层还是路径、剪枝是否会删掉合法解、输出规模是多少。请准备画出前三层搜索树。

\subsection{核心题卡}

\noindent\textbf{LeetCode 17 电话号码的字母组合。} 模型：每个数字对应若干字母，答案是所有按位选择形成的字符串。推导重点：暴力就是完整笛卡尔积，优化重点是路径长度和下标同步推进，避免在每层重新拼接大量中间字符串。

\textbf{二刷读题。} 读题时先画搜索树：层数代表什么，路径保存什么，终止条件是什么。本题可以压成“每个数字对应若干字母，答案是所有按位选择形成的字符串”，剪枝只能在这个搜索树上说明。先遮住答案，只保留题面，复述候选对象、合法性条件和输出目标。

\textbf{暴力回放。} 慢版本就是完整搜索树：每层列出选择列表，路径到达终止条件时检查答案。它先保证覆盖所有可能，再把明显无效或重复的分支剪掉。二刷时先写慢版本或口述慢版本，用它确认不会漏答案。

\textbf{特例回放。} 拿 digits=23 手算，2 对应 a/b/c，3 对应 d/e/f，第一层选 a 后第二层依次选 d/e/f，得到 ad、ae、af。规律是层数对应输入数字位置，路径保存已经选择的字母；空输入应返回空结果而不是包含空串。把这个特例继续拆开看：一层递归对应一个明确决策，路径保存已经做出的选择，选择列表保存仍可能扩展的分支。剪枝前必须能说明被剪掉的整棵子树没有合法答案，而不是只因为它看起来不优。复盘时把这个特例压成状态字段：路径、起点或使用标记、剩余目标、答案集合、剪枝条件；进入递归和退出递归必须成对恢复；再用边界 空输入、包含 7 或 9、输入中不应出现 0 和 1、输出规模指数增长 反查初始化、更新顺序和退出条件。

\textbf{状态回放。} 手算路线：手算时给每层标出选择列表和路径。剪枝发生前要指出被剪分支为什么已经不可能得到合法答案。状态字段：路径、起点或使用标记、剩余目标、答案集合、剪枝条件；进入递归和退出递归必须成对恢复。状态升级：把完整搜索树加上合法性检查、剩余资源剪枝和同层去重；路径状态进入和退出要成对恢复。

\textbf{证明和边界。} 证明从搜索树覆盖开始：每个合法答案有且只有一条路径，剪枝只删掉不可能成功的分支。边界：空输入、包含 7 或 9、输入中不应出现 0 和 1、输出规模指数增长。变体：若字符映射来自键盘以外的字典，只需要替换候选表，搜索骨架不变。

\noindent\textbf{LeetCode 22 括号生成。} 模型：合法括号串的前缀中右括号数量不能超过左括号数量。推导重点：状态保存已用左括号和右括号数量，只有满足约束的选择才继续递归。

\textbf{二刷读题。} 读题时先画搜索树：层数代表什么，路径保存什么，终止条件是什么。本题可以压成“合法括号串的前缀中右括号数量不能超过左括号数量”，剪枝只能在这个搜索树上说明。先遮住答案，只保留题面，复述候选对象、合法性条件和输出目标。

\textbf{暴力回放。} 慢版本就是完整搜索树：每层列出选择列表，路径到达终止条件时检查答案。它先保证覆盖所有可能，再把明显无效或重复的分支剪掉。二刷时先写慢版本或口述慢版本，用它确认不会漏答案。

\textbf{特例回放。} 拿 n=2 手算，第一步只能放左括号，之后可以放左括号得到 (())，也可以在已有左括号数量大于右括号时放右括号得到 ()()。规律是状态包含已放左括号数和右括号数，右括号不能超过左括号，左括号不能超过 n。把这个特例继续拆开看：一层递归对应一个明确决策，路径保存已经做出的选择，选择列表保存仍可能扩展的分支。剪枝前必须能说明被剪掉的整棵子树没有合法答案，而不是只因为它看起来不优。复盘时把这个特例压成状态字段：路径、起点或使用标记、剩余目标、答案集合、剪枝条件；进入递归和退出递归必须成对恢复；再用边界 n 为零或一、右括号提前超过左括号、输出规模是卡特兰数 反查初始化、更新顺序和退出条件。

\textbf{状态回放。} 手算路线：手算时给每层标出选择列表和路径。剪枝发生前要指出被剪分支为什么已经不可能得到合法答案。状态字段：路径、起点或使用标记、剩余目标、答案集合、剪枝条件；进入递归和退出递归必须成对恢复。状态升级：把完整搜索树加上合法性检查、剩余资源剪枝和同层去重；路径状态进入和退出要成对恢复。

\textbf{证明和边界。} 证明从搜索树覆盖开始：每个合法答案有且只有一条路径，剪枝只删掉不可能成功的分支。边界：n 为零或一、右括号提前超过左括号、输出规模是卡特兰数。变体：若有多种括号类型，状态还需要栈保存类型匹配关系。

\noindent\textbf{LeetCode 39 组合总和。} 模型：候选数字可重复使用，答案不关心排列顺序。推导重点：排序后按起点枚举，递归时仍传当前下标表示可以继续使用同一数字。

\textbf{二刷读题。} 读题时先画搜索树：层数代表什么，路径保存什么，终止条件是什么。本题可以压成“候选数字可重复使用，答案不关心排列顺序”，剪枝只能在这个搜索树上说明。先遮住答案，只保留题面，复述候选对象、合法性条件和输出目标。

\textbf{暴力回放。} 慢版本就是完整搜索树：每层列出选择列表，路径到达终止条件时检查答案。它先保证覆盖所有可能，再把明显无效或重复的分支剪掉。二刷时先写慢版本或口述慢版本，用它确认不会漏答案。

\textbf{特例回放。} 拿 candidates=[2,3,6,7]、target=7 手算，选择 2 后仍可以继续选择 2，因为每个数可重复；路径 2,2,3 成功，7 单独也成功。规律是回溯起点不回退，但递归下一层仍传当前下标以允许重复使用。把这个特例继续拆开看：一层递归对应一个明确决策，路径保存已经做出的选择，选择列表保存仍可能扩展的分支。剪枝前必须能说明被剪掉的整棵子树没有合法答案，而不是只因为它看起来不优。复盘时把这个特例压成状态字段：路径、起点或使用标记、剩余目标、答案集合、剪枝条件；进入递归和退出递归必须成对恢复；再用边界 目标小于最小数、刚好等于、候选有重复时题意如何处理、输出规模 反查初始化、更新顺序和退出条件。

\textbf{状态回放。} 手算路线：手算时给每层标出选择列表和路径。剪枝发生前要指出被剪分支为什么已经不可能得到合法答案。状态字段：路径、起点或使用标记、剩余目标、答案集合、剪枝条件；进入递归和退出递归必须成对恢复。状态升级：把完整搜索树加上合法性检查、剩余资源剪枝和同层去重；路径状态进入和退出要成对恢复。

\textbf{证明和边界。} 证明从搜索树覆盖开始：每个合法答案有且只有一条路径，剪枝只删掉不可能成功的分支。边界：目标小于最小数、刚好等于、候选有重复时题意如何处理、输出规模。变体：组合总和 II 每个数只能用一次，起点推进和去重条件都不同。

\noindent\textbf{LeetCode 40 组合总和 II。} 模型：每个候选只能使用一次，输入中可能存在重复值。推导重点：排序后递归下一层从 i 加一开始，同层遇到相同值要跳过。

\textbf{二刷读题。} 读题时先画搜索树：层数代表什么，路径保存什么，终止条件是什么。本题可以压成“每个候选只能使用一次，输入中可能存在重复值”，剪枝只能在这个搜索树上说明。先遮住答案，只保留题面，复述候选对象、合法性条件和输出目标。

\textbf{暴力回放。} 慢版本就是完整搜索树：每层列出选择列表，路径到达终止条件时检查答案。它先保证覆盖所有可能，再把明显无效或重复的分支剪掉。二刷时先写慢版本或口述慢版本，用它确认不会漏答案。

\textbf{特例回放。} 拿 candidates=[10,1,2,7,6,1,5]、target=8 手算，排序后两个 1 在同一层只能选第一个作为起点，否则会生成重复组合；但选了第一个 1 后，下一层可以选第二个 1。规律是每个数只能用一次，递归下一层传 i+1，并做同层去重。把这个特例继续拆开看：一层递归对应一个明确决策，路径保存已经做出的选择，选择列表保存仍可能扩展的分支。剪枝前必须能说明被剪掉的整棵子树没有合法答案，而不是只因为它看起来不优。复盘时把这个特例压成状态字段：路径、起点或使用标记、剩余目标、答案集合、剪枝条件；进入递归和退出递归必须成对恢复；再用边界 全重复、目标为零、同层去重写成跨层去重、答案顺序不要求 反查初始化、更新顺序和退出条件。

\textbf{状态回放。} 手算路线：手算时给每层标出选择列表和路径。剪枝发生前要指出被剪分支为什么已经不可能得到合法答案。状态字段：路径、起点或使用标记、剩余目标、答案集合、剪枝条件；进入递归和退出递归必须成对恢复。状态升级：把完整搜索树加上合法性检查、剩余资源剪枝和同层去重；路径状态进入和退出要成对恢复。

\textbf{证明和边界。} 证明从搜索树覆盖开始：每个合法答案有且只有一条路径，剪枝只删掉不可能成功的分支。边界：全重复、目标为零、同层去重写成跨层去重、答案顺序不要求。变体：子集 II 使用同样的同层去重，只是每个路径节点都可能成为答案。

\noindent\textbf{LeetCode 46 全排列。} 模型：排列关心元素顺序，每个位置从所有未使用元素中选择一个。推导重点：路径长度达到 n 时产生答案，used 数组保证同一元素不重复使用。

\textbf{二刷读题。} 读题时先画搜索树：层数代表什么，路径保存什么，终止条件是什么。本题可以压成“排列关心元素顺序，每个位置从所有未使用元素中选择一个”，剪枝只能在这个搜索树上说明。先遮住答案，只保留题面，复述候选对象、合法性条件和输出目标。

\textbf{暴力回放。} 慢版本就是完整搜索树：每层列出选择列表，路径到达终止条件时检查答案。它先保证覆盖所有可能，再把明显无效或重复的分支剪掉。二刷时先写慢版本或口述慢版本，用它确认不会漏答案。

\textbf{特例回放。} 拿 [1,2,3] 手算，第一位可选 1、2、3；选 1 后，第二位只能从未使用的 2、3 中选。规律是路径长度等于 n 时得到一个排列，used 数组表示哪些下标已经在当前路径里。把这个特例继续拆开看：一层递归对应一个明确决策，路径保存已经做出的选择，选择列表保存仍可能扩展的分支。剪枝前必须能说明被剪掉的整棵子树没有合法答案，而不是只因为它看起来不优。复盘时把这个特例压成状态字段：路径、起点或使用标记、剩余目标、答案集合、剪枝条件；进入递归和退出递归必须成对恢复；再用边界 空数组、单元素、输出规模为 n 阶乘、元素值唯一假设 反查初始化、更新顺序和退出条件。

\textbf{状态回放。} 手算路线：手算时给每层标出选择列表和路径。剪枝发生前要指出被剪分支为什么已经不可能得到合法答案。状态字段：路径、起点或使用标记、剩余目标、答案集合、剪枝条件；进入递归和退出递归必须成对恢复。状态升级：把完整搜索树加上合法性检查、剩余资源剪枝和同层去重；路径状态进入和退出要成对恢复。

\textbf{证明和边界。} 证明从搜索树覆盖开始：每个合法答案有且只有一条路径，剪枝只删掉不可能成功的分支。边界：空数组、单元素、输出规模为 n 阶乘、元素值唯一假设。变体：若有重复元素，需要排序并在同层跳过等价选择。

\noindent\textbf{LeetCode 47 全排列 II。} 模型：重复元素会让不同下标选择生成相同排列。推导重点：排序后使用 used 数组，同层遇到相同值且前一个相同值未使用时跳过。

\textbf{二刷读题。} 读题时先画搜索树：层数代表什么，路径保存什么，终止条件是什么。本题可以压成“重复元素会让不同下标选择生成相同排列”，剪枝只能在这个搜索树上说明。先遮住答案，只保留题面，复述候选对象、合法性条件和输出目标。

\textbf{暴力回放。} 慢版本就是完整搜索树：每层列出选择列表，路径到达终止条件时检查答案。它先保证覆盖所有可能，再把明显无效或重复的分支剪掉。二刷时先写慢版本或口述慢版本，用它确认不会漏答案。

\textbf{特例回放。} 拿 [1,1,2] 手算，第一位选第一个 1 和第二个 1 会生成同样分支，所以同层必须跳过后一个 1；但第一层选 1 后，第二层仍可选另一个 1。规律是排序后用 used 区分路径内使用，同层重复值只保留第一个未使用代表。把这个特例继续拆开看：一层递归对应一个明确决策，路径保存已经做出的选择，选择列表保存仍可能扩展的分支。剪枝前必须能说明被剪掉的整棵子树没有合法答案，而不是只因为它看起来不优。复盘时把这个特例压成状态字段：路径、起点或使用标记、剩余目标、答案集合、剪枝条件；进入递归和退出递归必须成对恢复；再用边界 全重复、部分重复、同层去重条件写反、输出顺序不要求 反查初始化、更新顺序和退出条件。

\textbf{状态回放。} 手算路线：手算时给每层标出选择列表和路径。剪枝发生前要指出被剪分支为什么已经不可能得到合法答案。状态字段：路径、起点或使用标记、剩余目标、答案集合、剪枝条件；进入递归和退出递归必须成对恢复。状态升级：把完整搜索树加上合法性检查、剩余资源剪枝和同层去重；路径状态进入和退出要成对恢复。

\textbf{证明和边界。} 证明从搜索树覆盖开始：每个合法答案有且只有一条路径，剪枝只删掉不可能成功的分支。边界：全重复、部分重复、同层去重条件写反、输出顺序不要求。变体：子集 II 和组合总和 II 也是同层去重，但起点推进方式不同。

\noindent\textbf{LeetCode 51 N 皇后。} 模型：每一行放一个皇后，列和两条对角线不能冲突。推导重点：逐行搜索，用列集合、主对角线集合和副对角线集合快速判断合法性。

\textbf{二刷读题。} 读题时先画搜索树：层数代表什么，路径保存什么，终止条件是什么。本题可以压成“每一行放一个皇后，列和两条对角线不能冲突”，剪枝只能在这个搜索树上说明。先遮住答案，只保留题面，复述候选对象、合法性条件和输出目标。

\textbf{暴力回放。} 慢版本就是完整搜索树：每层列出选择列表，路径到达终止条件时检查答案。它先保证覆盖所有可能，再把明显无效或重复的分支剪掉。二刷时先写慢版本或口述慢版本，用它确认不会漏答案。

\textbf{特例回放。} 拿 n=4 手算，第一行放某列后，下一行不能放同列、主对角线和副对角线冲突的位置。规律是按行递归，每行放一个皇后，三个集合分别记录列、row-col、row+col 是否被占用。把这个特例继续拆开看：一层递归对应一个明确决策，路径保存已经做出的选择，选择列表保存仍可能扩展的分支。剪枝前必须能说明被剪掉的整棵子树没有合法答案，而不是只因为它看起来不优。复盘时把这个特例压成状态字段：路径、起点或使用标记、剩余目标、答案集合、剪枝条件；进入递归和退出递归必须成对恢复；再用边界 n 为一、n 为二或三无解、对角线编号、回退时状态恢复 反查初始化、更新顺序和退出条件。

\textbf{状态回放。} 手算路线：手算时给每层标出选择列表和路径。剪枝发生前要指出被剪分支为什么已经不可能得到合法答案。状态字段：路径、起点或使用标记、剩余目标、答案集合、剪枝条件；进入递归和退出递归必须成对恢复。状态升级：把完整搜索树加上合法性检查、剩余资源剪枝和同层去重；路径状态进入和退出要成对恢复。

\textbf{证明和边界。} 证明从搜索树覆盖开始：每个合法答案有且只有一条路径，剪枝只删掉不可能成功的分支。边界：n 为一、n 为二或三无解、对角线编号、回退时状态恢复。变体：可用位掩码压缩三个约束集合，提高常数性能。

\noindent\textbf{LeetCode 78 子集。} 模型：每个元素都有选或不选两种选择，答案是所有路径节点。推导重点：回溯按起点向后枚举，避免不同顺序生成同一集合。

\textbf{二刷读题。} 读题时先画搜索树：层数代表什么，路径保存什么，终止条件是什么。本题可以压成“每个元素都有选或不选两种选择，答案是所有路径节点”，剪枝只能在这个搜索树上说明。先遮住答案，只保留题面，复述候选对象、合法性条件和输出目标。

\textbf{暴力回放。} 慢版本就是完整搜索树：每层列出选择列表，路径到达终止条件时检查答案。它先保证覆盖所有可能，再把明显无效或重复的分支剪掉。二刷时先写慢版本或口述慢版本，用它确认不会漏答案。

\textbf{特例回放。} 拿 [1,2] 画搜索树，第一层可以不选 1 或选 1；在选 1 的分支里再决定 2，得到 [1] 和 [1,2]。规律是每个元素对应选或不选，或用起点向后枚举，路径上的每个节点都是一个子集；空数组也要输出空集。把这个特例继续拆开看：一层递归对应一个明确决策，路径保存已经做出的选择，选择列表保存仍可能扩展的分支。剪枝前必须能说明被剪掉的整棵子树没有合法答案，而不是只因为它看起来不优。复盘时把这个特例压成状态字段：路径、起点或使用标记、剩余目标、答案集合、剪枝条件；进入递归和退出递归必须成对恢复；再用边界 空数组、单元素、输出数量为二的 n 次方、顺序不要求 反查初始化、更新顺序和退出条件。

\textbf{状态回放。} 手算路线：手算时给每层标出选择列表和路径。剪枝发生前要指出被剪分支为什么已经不可能得到合法答案。状态字段：路径、起点或使用标记、剩余目标、答案集合、剪枝条件；进入递归和退出递归必须成对恢复。状态升级：把完整搜索树加上合法性检查、剩余资源剪枝和同层去重；路径状态进入和退出要成对恢复。

\textbf{证明和边界。} 证明从搜索树覆盖开始：每个合法答案有且只有一条路径，剪枝只删掉不可能成功的分支。边界：空数组、单元素、输出数量为二的 n 次方、顺序不要求。变体：也可用位掩码枚举所有子集，适合 n 较小的场景。

\noindent\textbf{LeetCode 79 单词搜索。} 模型：网格路径不能重复使用同一格，状态包含位置、匹配下标和访问标记。推导重点：剪枝包括首字符不匹配、剩余长度不足、字符频次不足和越界访问。

\textbf{二刷读题。} 读题时先画搜索树：层数代表什么，路径保存什么，终止条件是什么。本题可以压成“网格路径不能重复使用同一格，状态包含位置、匹配下标和访问标记”，剪枝只能在这个搜索树上说明。先遮住答案，只保留题面，复述候选对象、合法性条件和输出目标。

\textbf{暴力回放。} 慢版本就是完整搜索树：每层列出选择列表，路径到达终止条件时检查答案。它先保证覆盖所有可能，再把明显无效或重复的分支剪掉。二刷时先写慢版本或口述慢版本，用它确认不会漏答案。

\textbf{特例回放。} 拿 board=[A B; C D]、word=AB 手算，从 A 出发只能走到相邻的 B，不能跳到 D；若搜索 ABA，第二个 A 不能复用起点格。规律是状态包含位置、匹配下标和访问标记，递归退出时要恢复访问标记。把这个特例继续拆开看：一层递归对应一个明确决策，路径保存已经做出的选择，选择列表保存仍可能扩展的分支。剪枝前必须能说明被剪掉的整棵子树没有合法答案，而不是只因为它看起来不优。复盘时把这个特例压成状态字段：路径、起点或使用标记、剩余目标、答案集合、剪枝条件；进入递归和退出递归必须成对恢复；再用边界 单字符单词、路径需要回退、重复字母、单元格不能复用 反查初始化、更新顺序和退出条件。

\textbf{状态回放。} 手算路线：手算时给每层标出选择列表和路径。剪枝发生前要指出被剪分支为什么已经不可能得到合法答案。状态字段：路径、起点或使用标记、剩余目标、答案集合、剪枝条件；进入递归和退出递归必须成对恢复。状态升级：把完整搜索树加上合法性检查、剩余资源剪枝和同层去重；路径状态进入和退出要成对恢复。

\textbf{证明和边界。} 证明从搜索树覆盖开始：每个合法答案有且只有一条路径，剪枝只删掉不可能成功的分支。边界：单字符单词、路径需要回退、重复字母、单元格不能复用。变体：若要查很多单词，应使用 Trie 合并搜索前缀。

\noindent\textbf{LeetCode 90 子集 II。} 模型：重复元素会让不同下标路径生成相同集合。推导重点：排序后在同一层跳过相同值，保留不同层使用重复值的可能。

\textbf{二刷读题。} 读题时先画搜索树：层数代表什么，路径保存什么，终止条件是什么。本题可以压成“重复元素会让不同下标路径生成相同集合”，剪枝只能在这个搜索树上说明。先遮住答案，只保留题面，复述候选对象、合法性条件和输出目标。

\textbf{暴力回放。} 慢版本就是完整搜索树：每层列出选择列表，路径到达终止条件时检查答案。它先保证覆盖所有可能，再把明显无效或重复的分支剪掉。二刷时先写慢版本或口述慢版本，用它确认不会漏答案。

\textbf{特例回放。} 拿 [1,2,2] 手算，排序后第二个 2 如果在同一层再次作为起点，会生成和第一个 2 同样的子集；但在已经选择第一个 2 的下一层，第二个 2 又可以被选出 [2,2]。规律是同层跳过重复值，不是全局禁止重复值。把这个特例继续拆开看：一层递归对应一个明确决策，路径保存已经做出的选择，选择列表保存仍可能扩展的分支。剪枝前必须能说明被剪掉的整棵子树没有合法答案，而不是只因为它看起来不优。复盘时把这个特例压成状态字段：路径、起点或使用标记、剩余目标、答案集合、剪枝条件；进入递归和退出递归必须成对恢复；再用边界 全重复、空集、重复值相邻、去重条件写成同层而不是全局 反查初始化、更新顺序和退出条件。

\textbf{状态回放。} 手算路线：手算时给每层标出选择列表和路径。剪枝发生前要指出被剪分支为什么已经不可能得到合法答案。状态字段：路径、起点或使用标记、剩余目标、答案集合、剪枝条件；进入递归和退出递归必须成对恢复。状态升级：把完整搜索树加上合法性检查、剩余资源剪枝和同层去重；路径状态进入和退出要成对恢复。

\textbf{证明和边界。} 证明从搜索树覆盖开始：每个合法答案有且只有一条路径，剪枝只删掉不可能成功的分支。边界：全重复、空集、重复值相邻、去重条件写成同层而不是全局。变体：组合总和 II 使用同样的同层去重思想。

\noindent\textbf{LeetCode 93 复原 IP 地址。} 模型：字符串要切成四段，每段必须是 0 到 255 且无非法前导零。推导重点：按段数和剩余长度剪枝，每次尝试长度一到三的片段。

\textbf{二刷读题。} 读题时先画搜索树：层数代表什么，路径保存什么，终止条件是什么。本题可以压成“字符串要切成四段，每段必须是 0 到 255 且无非法前导零”，剪枝只能在这个搜索树上说明。先遮住答案，只保留题面，复述候选对象、合法性条件和输出目标。

\textbf{暴力回放。} 慢版本就是完整搜索树：每层列出选择列表，路径到达终止条件时检查答案。它先保证覆盖所有可能，再把明显无效或重复的分支剪掉。二刷时先写慢版本或口述慢版本，用它确认不会漏答案。

\textbf{特例回放。} 拿 25525511135 手算，第一段可以是 255，但不能是 2552；段值必须在 0..255，且 01 这种前导零非法。规律是回溯选四段，每段长度 1 到 3，剩余字符数必须能填满剩余段数。把这个特例继续拆开看：一层递归对应一个明确决策，路径保存已经做出的选择，选择列表保存仍可能扩展的分支。剪枝前必须能说明被剪掉的整棵子树没有合法答案，而不是只因为它看起来不优。复盘时把这个特例压成状态字段：路径、起点或使用标记、剩余目标、答案集合、剪枝条件；进入递归和退出递归必须成对恢复；再用边界 前导零、数值超过 255、剩余字符太多或太少、刚好四段 反查初始化、更新顺序和退出条件。

\textbf{状态回放。} 手算路线：手算时给每层标出选择列表和路径。剪枝发生前要指出被剪分支为什么已经不可能得到合法答案。状态字段：路径、起点或使用标记、剩余目标、答案集合、剪枝条件；进入递归和退出递归必须成对恢复。状态升级：把完整搜索树加上合法性检查、剩余资源剪枝和同层去重；路径状态进入和退出要成对恢复。

\textbf{证明和边界。} 证明从搜索树覆盖开始：每个合法答案有且只有一条路径，剪枝只删掉不可能成功的分支。边界：前导零、数值超过 255、剩余字符太多或太少、刚好四段。变体：分割回文串也是前缀分割，但合法性检查不同。

\noindent\textbf{LeetCode 131 分割回文串。} 模型：每次选择一个从当前位置开始的回文前缀。推导重点：预处理回文 DP 表，让回溯合法性检查为常数。

\textbf{二刷读题。} 读题时先画搜索树：层数代表什么，路径保存什么，终止条件是什么。本题可以压成“每次选择一个从当前位置开始的回文前缀”，剪枝只能在这个搜索树上说明。先遮住答案，只保留题面，复述候选对象、合法性条件和输出目标。

\textbf{暴力回放。} 慢版本就是完整搜索树：每层列出选择列表，路径到达终止条件时检查答案。它先保证覆盖所有可能，再把明显无效或重复的分支剪掉。二刷时先写慢版本或口述慢版本，用它确认不会漏答案。

\textbf{特例回放。} 拿 aab 手算，第一刀可以切 a，后面 ab 不全是回文；也可以切 aa，后面 b 是回文，得到 [aa,b]。规律是回溯枚举切分点，但每段必须先通过回文检查；预处理回文表可以避免重复检查子串。把这个特例继续拆开看：一层递归对应一个明确决策，路径保存已经做出的选择，选择列表保存仍可能扩展的分支。剪枝前必须能说明被剪掉的整棵子树没有合法答案，而不是只因为它看起来不优。复盘时把这个特例压成状态字段：路径、起点或使用标记、剩余目标、答案集合、剪枝条件；进入递归和退出递归必须成对恢复；再用边界 空串、单字符、全相同字符、重复分割路径 反查初始化、更新顺序和退出条件。

\textbf{状态回放。} 手算路线：手算时给每层标出选择列表和路径。剪枝发生前要指出被剪分支为什么已经不可能得到合法答案。状态字段：路径、起点或使用标记、剩余目标、答案集合、剪枝条件；进入递归和退出递归必须成对恢复。状态升级：把完整搜索树加上合法性检查、剩余资源剪枝和同层去重；路径状态进入和退出要成对恢复。

\textbf{证明和边界。} 证明从搜索树覆盖开始：每个合法答案有且只有一条路径，剪枝只删掉不可能成功的分支。边界：空串、单字符、全相同字符、重复分割路径。变体：若只求最少切割次数，可改成 DP 最短切分。

\noindent\textbf{LeetCode 140 单词拆分 II。} 模型：需要输出所有由字典单词拼出的句子，输出规模本身可能指数级。推导重点：先用 DP 或记忆化判断后缀是否可拆，再 DFS 枚举可行前缀并在末尾拼接句子。

\textbf{二刷读题。} 读题时先画搜索树：层数代表什么，路径保存什么，终止条件是什么。本题可以压成“需要输出所有由字典单词拼出的句子，输出规模本身可能指数级”，剪枝只能在这个搜索树上说明。先遮住答案，只保留题面，复述候选对象、合法性条件和输出目标。

\textbf{暴力回放。} 慢版本就是完整搜索树：每层列出选择列表，路径到达终止条件时检查答案。它先保证覆盖所有可能，再把明显无效或重复的分支剪掉。二刷时先写慢版本或口述慢版本，用它确认不会漏答案。

\textbf{特例回放。} 拿 catsanddog 手算，前缀 cat 和 cats 都能走，但某些后缀继续切会失败。若直接 DFS，每次都会重复探索同一个后缀。规律是先用记忆化记录某个起点之后能生成哪些句子，或先用 DP 剪掉不可拆后缀，再枚举路径。把这个特例继续拆开看：一层递归对应一个明确决策，路径保存已经做出的选择，选择列表保存仍可能扩展的分支。剪枝前必须能说明被剪掉的整棵子树没有合法答案，而不是只因为它看起来不优。复盘时把这个特例压成状态字段：路径、起点或使用标记、剩余目标、答案集合、剪枝条件；进入递归和退出递归必须成对恢复；再用边界 无解后缀、重复单词、路径拼接成本、答案数量主导复杂度 反查初始化、更新顺序和退出条件。

\textbf{状态回放。} 手算路线：手算时给每层标出选择列表和路径。剪枝发生前要指出被剪分支为什么已经不可能得到合法答案。状态字段：路径、起点或使用标记、剩余目标、答案集合、剪枝条件；进入递归和退出递归必须成对恢复。状态升级：把完整搜索树加上合法性检查、剩余资源剪枝和同层去重；路径状态进入和退出要成对恢复。

\textbf{证明和边界。} 证明从搜索树覆盖开始：每个合法答案有且只有一条路径，剪枝只删掉不可能成功的分支。边界：无解后缀、重复单词、路径拼接成本、答案数量主导复杂度。变体：只问能否拆分时用布尔 DP 即可，不要承担枚举所有句子的成本。

\noindent\textbf{LeetCode 212 单词搜索 II。} 模型：多个单词在同一棋盘上搜索，公共前缀可以共享。推导重点：先建 Trie，再从每个格子回溯；路径到达单词结束节点就收集答案。

\textbf{二刷读题。} 读题时先画搜索树：层数代表什么，路径保存什么，终止条件是什么。本题可以压成“多个单词在同一棋盘上搜索，公共前缀可以共享”，剪枝只能在这个搜索树上说明。先遮住答案，只保留题面，复述候选对象、合法性条件和输出目标。

\textbf{暴力回放。} 慢版本就是完整搜索树：每层列出选择列表，路径到达终止条件时检查答案。它先保证覆盖所有可能，再把明显无效或重复的分支剪掉。二刷时先写慢版本或口述慢版本，用它确认不会漏答案。

\textbf{特例回放。} 拿 board 上有 oath、pea、eat 手算，单词共享前缀时不应对每个词重复全图搜索；从 o 出发沿 Trie 前缀 o-a-t-h 找到 oath。规律是 Trie 合并所有单词前缀，DFS 走棋盘同时走 Trie，命中单词后记录并可去重。把这个特例继续拆开看：一层递归对应一个明确决策，路径保存已经做出的选择，选择列表保存仍可能扩展的分支。剪枝前必须能说明被剪掉的整棵子树没有合法答案，而不是只因为它看起来不优。复盘时把这个特例压成状态字段：路径、起点或使用标记、剩余目标、答案集合、剪枝条件；进入递归和退出递归必须成对恢复；再用边界 重复单词、同一格不能复用、前缀是完整单词、剪枝删除空子树 反查初始化、更新顺序和退出条件。

\textbf{状态回放。} 手算路线：手算时给每层标出选择列表和路径。剪枝发生前要指出被剪分支为什么已经不可能得到合法答案。状态字段：路径、起点或使用标记、剩余目标、答案集合、剪枝条件；进入递归和退出递归必须成对恢复。状态升级：把完整搜索树加上合法性检查、剩余资源剪枝和同层去重；路径状态进入和退出要成对恢复。

\textbf{证明和边界。} 证明从搜索树覆盖开始：每个合法答案有且只有一条路径，剪枝只删掉不可能成功的分支。边界：重复单词、同一格不能复用、前缀是完整单词、剪枝删除空子树。变体：这是单词搜索从单目标扩展到多目标后的结构升级。

\noindent\textbf{LeetCode 216 组合总和 III。} 模型：从一到九中选 k 个数和为 n，每个数最多用一次。推导重点：按递增起点枚举，利用剩余个数和剩余和剪枝。

\textbf{二刷读题。} 读题时先画搜索树：层数代表什么，路径保存什么，终止条件是什么。本题可以压成“从一到九中选 k 个数和为 n，每个数最多用一次”，剪枝只能在这个搜索树上说明。先遮住答案，只保留题面，复述候选对象、合法性条件和输出目标。

\textbf{暴力回放。} 慢版本就是完整搜索树：每层列出选择列表，路径到达终止条件时检查答案。它先保证覆盖所有可能，再把明显无效或重复的分支剪掉。二刷时先写慢版本或口述慢版本，用它确认不会漏答案。

\textbf{特例回放。} 拿 k=3、n=7 手算，选择 1 后还需两个数凑 6，可以选 2 和 4；若当前和已经超过 7，就不用再往下扩展。规律是路径长度、剩余和和起点共同定义状态，数字只能向后选避免重复组合。把这个特例继续拆开看：一层递归对应一个明确决策，路径保存已经做出的选择，选择列表保存仍可能扩展的分支。剪枝前必须能说明被剪掉的整棵子树没有合法答案，而不是只因为它看起来不优。复盘时把这个特例压成状态字段：路径、起点或使用标记、剩余目标、答案集合、剪枝条件；进入递归和退出递归必须成对恢复；再用边界 n 太小、n 太大、k 为零、路径长度已经超过 k 反查初始化、更新顺序和退出条件。

\textbf{状态回放。} 手算路线：手算时给每层标出选择列表和路径。剪枝发生前要指出被剪分支为什么已经不可能得到合法答案。状态字段：路径、起点或使用标记、剩余目标、答案集合、剪枝条件；进入递归和退出递归必须成对恢复。状态升级：把完整搜索树加上合法性检查、剩余资源剪枝和同层去重；路径状态进入和退出要成对恢复。

\textbf{证明和边界。} 证明从搜索树覆盖开始：每个合法答案有且只有一条路径，剪枝只删掉不可能成功的分支。边界：n 太小、n 太大、k 为零、路径长度已经超过 k。变体：可预估剩余最小和最大和进一步剪枝。

\noindent\textbf{LeetCode 698 划分为 k 个相等的子集。} 模型：每个元素要分配到 k 个桶中，所有桶目标和相同。推导重点：先判断总和可整除并排序降序，再用桶剩余容量回溯放置元素。

\textbf{二刷读题。} 读题时先画搜索树：层数代表什么，路径保存什么，终止条件是什么。本题可以压成“每个元素要分配到 k 个桶中，所有桶目标和相同”，剪枝只能在这个搜索树上说明。先遮住答案，只保留题面，复述候选对象、合法性条件和输出目标。

\textbf{暴力回放。} 慢版本就是完整搜索树：每层列出选择列表，路径到达终止条件时检查答案。它先保证覆盖所有可能，再把明显无效或重复的分支剪掉。二刷时先写慢版本或口述慢版本，用它确认不会漏答案。

\textbf{特例回放。} 拿 [4,3,2,3,5,2,1]、k=4 手算，总和 20，每桶目标 5；先放 5 单独成桶，再尝试 4+1、3+2。规律是回溯状态是桶剩余容量和已用元素，排序后先放大数能更早剪枝。把这个特例继续拆开看：一层递归对应一个明确决策，路径保存已经做出的选择，选择列表保存仍可能扩展的分支。剪枝前必须能说明被剪掉的整棵子树没有合法答案，而不是只因为它看起来不优。复盘时把这个特例压成状态字段：路径、起点或使用标记、剩余目标、答案集合、剪枝条件；进入递归和退出递归必须成对恢复；再用边界 存在元素大于目标、重复元素、空桶对称、k 为一 反查初始化、更新顺序和退出条件。

\textbf{状态回放。} 手算路线：手算时给每层标出选择列表和路径。剪枝发生前要指出被剪分支为什么已经不可能得到合法答案。状态字段：路径、起点或使用标记、剩余目标、答案集合、剪枝条件；进入递归和退出递归必须成对恢复。状态升级：把完整搜索树加上合法性检查、剩余资源剪枝和同层去重；路径状态进入和退出要成对恢复。

\textbf{证明和边界。} 证明从搜索树覆盖开始：每个合法答案有且只有一条路径，剪枝只删掉不可能成功的分支。边界：存在元素大于目标、重复元素、空桶对称、k 为一。变体：状态压缩 DP 也能做，适合用掩码表示已选元素集合。

\subsection{自测标准}

二刷时不要只确认自己见过这道题。请遮住答案，先讲题面模型，再讲暴力基线和瓶颈，然后说出优化结构保存了什么、不变量如何排除错误候选、边界实验如何打破错误实现。

\section{动态规划与状态定义：二刷复盘卡}
这一大节只围绕一个题型复盘，不把题号直接铺成目录。阅读时先看复盘目标，再读核心题卡，最后用自测标准检查自己是否能独立重讲。

\subsection{复盘目标}

追问通常会问状态是否足够、初始化为什么这样、空间压缩读到的是旧值还是新值。请准备从暴力搜索树指出重复子问题。

\subsection{核心题卡}

\noindent\textbf{LeetCode 10 正则表达式匹配。} 模型：模式中的点号匹配单字符，星号修饰前一个字符并允许出现零次或多次。推导重点：二维 DP 表示两个前缀是否匹配，星号转移同时覆盖零次和消耗一个字符后继续匹配。

\textbf{二刷读题。} 读题时先找重复子问题，而不是直接背公式。本题可以压成“模式中的点号匹配单字符，星号修饰前一个字符并允许出现零次或多次”，二维 DP 表示两个前缀是否匹配，星号转移同时覆盖零次和消耗一个字符后继续匹配就是状态之间的依赖关系。先遮住答案，只保留题面，复述候选对象、合法性条件和输出目标。

\textbf{暴力回放。} 慢版本先写递归搜索树：节点表示处理位置、剩余资源和当前选择。只要不同路径反复到达同一个节点，就说明需要把这个节点命名成 DP 状态。二刷时先写慢版本或口述慢版本，用它确认不会漏答案。

\textbf{特例回放。} 拿 s=a、p=a* 手算，星号可以让 a 出现一次；拿 s=空、p=a*，星号又可以让 a 出现零次。再拿 s=aa、p=a*，消耗一个 a 后模式仍停在 a*，还能继续匹配。规律是星号有两条分叉：当零次用时看 dp[i][j-2]，当消耗当前字符时看 dp[i-1][j]。把这个特例继续拆开看：先问暴力搜索里哪些不同路径到达同一个后续问题，再给这个后续问题命名为状态。状态必须包含未来决策需要的全部信息，转移只从更小且已经算清的状态来。复盘时把这个特例压成状态字段：状态下标、状态值语义、初始化、转移来源、遍历顺序；空间压缩时还要指出读到的是旧值还是新值；再用边界 空串、能匹配空串的星号链、星号前必须有字符、点号和普通字符的匹配条件 反查初始化、更新顺序和退出条件。

\textbf{状态回放。} 手算路线：手算时先写一个很小的 DP 表或记忆化搜索记录。每个格子旁边写它来自哪些更小状态。状态字段：状态下标、状态值语义、初始化、转移来源、遍历顺序；空间压缩时还要指出读到的是旧值还是新值。状态升级：把重复递归节点命名为状态；遍历顺序只是在保证依赖状态已经算完。

\textbf{证明和边界。} 证明从最后一步开始：任意最优解的最后一步都在转移枚举中，转移来源已经由更小状态正确计算。边界：空串、能匹配空串的星号链、星号前必须有字符、点号和普通字符的匹配条件。变体：通配符匹配中的星号语义是任意长度字符串，不能直接搬用本题转移。

\noindent\textbf{LeetCode 44 通配符匹配。} 模型：问号匹配一个字符，星号匹配任意长度字符串，状态是两个前缀能否匹配。推导重点：DP 中星号可匹配空串或继续吞掉当前字符；贪心写法记录最近星号并在失配时回退扩展星号覆盖范围。

\textbf{二刷读题。} 读题时先找重复子问题，而不是直接背公式。本题可以压成“问号匹配一个字符，星号匹配任意长度字符串，状态是两个前缀能否匹配”，DP 中星号可匹配空串或继续吞掉当前字符；贪心写法记录最近星号并在失配时回退扩展星号覆盖范围就是状态之间的依赖关系。先遮住答案，只保留题面，复述候选对象、合法性条件和输出目标。

\textbf{暴力回放。} 慢版本先写递归搜索树：节点表示处理位置、剩余资源和当前选择。只要不同路径反复到达同一个节点，就说明需要把这个节点命名成 DP 状态。二刷时先写慢版本或口述慢版本，用它确认不会漏答案。

\textbf{特例回放。} 拿 s=ab、p=a* 手算，星号可以吞掉 b；拿 s=空、p=*，星号也可以代表空。DP 中星号的两条路是匹配空串或继续吞当前字符。贪心中失配时回到最近星号，让它多吞一个字符。规律是本题星号独立代表任意串，和正则里修饰前一字符的星号不同。把这个特例继续拆开看：先问暴力搜索里哪些不同路径到达同一个后续问题，再给这个后续问题命名为状态。状态必须包含未来决策需要的全部信息，转移只从更小且已经算清的状态来。复盘时把这个特例压成状态字段：状态下标、状态值语义、初始化、转移来源、遍历顺序；空间压缩时还要指出读到的是旧值还是新值；再用边界 连续星号、空串、模式只含星号、贪心回退时字符串位置不能倒退错 反查初始化、更新顺序和退出条件。

\textbf{状态回放。} 手算路线：手算时先写一个很小的 DP 表或记忆化搜索记录。每个格子旁边写它来自哪些更小状态。状态字段：状态下标、状态值语义、初始化、转移来源、遍历顺序；空间压缩时还要指出读到的是旧值还是新值。状态升级：把重复递归节点命名为状态；遍历顺序只是在保证依赖状态已经算完。

\textbf{证明和边界。} 证明从最后一步开始：任意最优解的最后一步都在转移枚举中，转移来源已经由更小状态正确计算。边界：连续星号、空串、模式只含星号、贪心回退时字符串位置不能倒退错。变体：正则表达式匹配的星号修饰前一个字符，和本题星号独立匹配任意串不同。

\noindent\textbf{LeetCode 53 最大子数组和。} 模型：状态是必须以当前位置结尾的最大连续和。推导重点：用 [-2,3] 和 [2,3] 对比，推出当前位置只在单独开始与接上旧后缀之间取较大值。

\textbf{二刷读题。} 读题时先找重复子问题，而不是直接背公式。本题可以压成“状态是必须以当前位置结尾的最大连续和”，用 [-2,3] 和 [2,3] 对比，推出当前位置只在单独开始与接上旧后缀之间取较大值就是状态之间的依赖关系。先遮住答案，只保留题面，复述候选对象、合法性条件和输出目标。

\textbf{暴力回放。} 慢版本先写递归搜索树：节点表示处理位置、剩余资源和当前选择。只要不同路径反复到达同一个节点，就说明需要把这个节点命名成 DP 状态。二刷时先写慢版本或口述慢版本，用它确认不会漏答案。

\textbf{特例回放。} 拿 [-2, 3] 手算。以 3 结尾的连续段只有两个候选：单独 [3]，或者接上前一个最佳后缀 [-2,3]，后者和为 1，比 3 更差，所以负后缀必须丢。再拿 [2, 3]，候选是 [3] 与 [2,3]，后者为 5，更好，所以正后缀要接。再拿 [-1, 0]，接不接都不增益，可以从当前重新开始而不影响最大和。规律是当前位置只需要比较 nums[i] 与 bestEndingHere + nums[i]，全局答案再和 bestEndingHere 比较；全负数组要求答案初始化为第一个数，不能初始化为 0。把这个特例继续拆开看：先问暴力搜索里哪些不同路径到达同一个后续问题，再给这个后续问题命名为状态。状态必须包含未来决策需要的全部信息，转移只从更小且已经算清的状态来。复盘时把这个特例压成状态字段：状态下标、状态值语义、初始化、转移来源、遍历顺序；空间压缩时还要指出读到的是旧值还是新值；再用边界 全负数组、单元素、和溢出、答案不能初始化为零 反查初始化、更新顺序和退出条件。

\textbf{状态回放。} 手算路线：手算时先写一个很小的 DP 表或记忆化搜索记录。每个格子旁边写它来自哪些更小状态。状态字段：状态下标、状态值语义、初始化、转移来源、遍历顺序；空间压缩时还要指出读到的是旧值还是新值。状态升级：把重复递归节点命名为状态；遍历顺序只是在保证依赖状态已经算完。

\textbf{证明和边界。} 证明从最后一步开始：任意最优解的最后一步都在转移枚举中，转移来源已经由更小状态正确计算。边界：全负数组、单元素、和溢出、答案不能初始化为零。变体：若要求返回区间下标，同步记录重新开始位置和最优左右端。

\noindent\textbf{LeetCode 62 不同路径。} 模型：网格状态由上方和左方两种最后一步转移而来。推导重点：二维表可压缩为一维，因为当前行只依赖上一行同列和当前行左侧。

\textbf{二刷读题。} 读题时先找重复子问题，而不是直接背公式。本题可以压成“网格状态由上方和左方两种最后一步转移而来”，二维表可压缩为一维，因为当前行只依赖上一行同列和当前行左侧就是状态之间的依赖关系。先遮住答案，只保留题面，复述候选对象、合法性条件和输出目标。

\textbf{暴力回放。} 慢版本先写递归搜索树：节点表示处理位置、剩余资源和当前选择。只要不同路径反复到达同一个节点，就说明需要把这个节点命名成 DP 状态。二刷时先写慢版本或口述慢版本，用它确认不会漏答案。

\textbf{特例回放。} 拿 2x3 网格手算，到每个格子的最后一步只能从上方或左方来；第一行和第一列都只有 1 条路，右下角等于左边 2 加上上边 1，得到 3。规律是 dp[row][col]=上方+左方，滚动数组只是在保存上一行同列和当前行左侧。把这个特例继续拆开看：先问暴力搜索里哪些不同路径到达同一个后续问题，再给这个后续问题命名为状态。状态必须包含未来决策需要的全部信息，转移只从更小且已经算清的状态来。复盘时把这个特例压成状态字段：状态下标、状态值语义、初始化、转移来源、遍历顺序；空间压缩时还要指出读到的是旧值还是新值；再用边界 一行、一列、较大网格结果溢出、障碍物版本边界 反查初始化、更新顺序和退出条件。

\textbf{状态回放。} 手算路线：手算时先写一个很小的 DP 表或记忆化搜索记录。每个格子旁边写它来自哪些更小状态。状态字段：状态下标、状态值语义、初始化、转移来源、遍历顺序；空间压缩时还要指出读到的是旧值还是新值。状态升级：把重复递归节点命名为状态；遍历顺序只是在保证依赖状态已经算完。

\textbf{证明和边界。} 证明从最后一步开始：任意最优解的最后一步都在转移枚举中，转移来源已经由更小状态正确计算。边界：一行、一列、较大网格结果溢出、障碍物版本边界。变体：带障碍物时遇到障碍状态清零，第一行第一列要按阻断处理。

\noindent\textbf{LeetCode 63 不同路径 II。} 模型：障碍物让某些格子的路径数变为零。推导重点：滚动数组中遇到障碍直接把当前位置清零，否则累加左侧。

\textbf{二刷读题。} 读题时先找重复子问题，而不是直接背公式。本题可以压成“障碍物让某些格子的路径数变为零”，滚动数组中遇到障碍直接把当前位置清零，否则累加左侧就是状态之间的依赖关系。先遮住答案，只保留题面，复述候选对象、合法性条件和输出目标。

\textbf{暴力回放。} 慢版本先写递归搜索树：节点表示处理位置、剩余资源和当前选择。只要不同路径反复到达同一个节点，就说明需要把这个节点命名成 DP 状态。二刷时先写慢版本或口述慢版本，用它确认不会漏答案。

\textbf{特例回放。} 拿 [[0,0],[1,0]] 手算，左下角是障碍，路径数必须清零；右下角只能从上方到达，所以答案是 1。规律是障碍格不是普通转移，而是把当前状态置零；起点或终点是障碍时也按同一清零语义处理。把这个特例继续拆开看：先问暴力搜索里哪些不同路径到达同一个后续问题，再给这个后续问题命名为状态。状态必须包含未来决策需要的全部信息，转移只从更小且已经算清的状态来。复盘时把这个特例压成状态字段：状态下标、状态值语义、初始化、转移来源、遍历顺序；空间压缩时还要指出读到的是旧值还是新值；再用边界 起点或终点是障碍、整行被截断、单行障碍 反查初始化、更新顺序和退出条件。

\textbf{状态回放。} 手算路线：手算时先写一个很小的 DP 表或记忆化搜索记录。每个格子旁边写它来自哪些更小状态。状态字段：状态下标、状态值语义、初始化、转移来源、遍历顺序；空间压缩时还要指出读到的是旧值还是新值。状态升级：把重复递归节点命名为状态；遍历顺序只是在保证依赖状态已经算完。

\textbf{证明和边界。} 证明从最后一步开始：任意最优解的最后一步都在转移枚举中，转移来源已经由更小状态正确计算。边界：起点或终点是障碍、整行被截断、单行障碍。变体：若允许向上或向左走，图会有环，普通填表不再成立。

\noindent\textbf{LeetCode 64 最小路径和。} 模型：状态是到达当前格子的最小代价，最后一步来自上方或左方。推导重点：先初始化第一行和第一列，再处理内部，主转移保持简单。

\textbf{二刷读题。} 读题时先找重复子问题，而不是直接背公式。本题可以压成“状态是到达当前格子的最小代价，最后一步来自上方或左方”，先初始化第一行和第一列，再处理内部，主转移保持简单就是状态之间的依赖关系。先遮住答案，只保留题面，复述候选对象、合法性条件和输出目标。

\textbf{暴力回放。} 慢版本先写递归搜索树：节点表示处理位置、剩余资源和当前选择。只要不同路径反复到达同一个节点，就说明需要把这个节点命名成 DP 状态。二刷时先写慢版本或口述慢版本，用它确认不会漏答案。

\textbf{特例回放。} 拿 [[1,3,1],[1,5,1],[4,2,1]] 手算，到 5 的最小代价来自左侧 4 或上方 4，再加 5；右下角只需要比较上方和左方的已知最小代价。规律是最后一步来自上或左，先初始化第一行第一列，内部统一取 min。把这个特例继续拆开看：先问暴力搜索里哪些不同路径到达同一个后续问题，再给这个后续问题命名为状态。状态必须包含未来决策需要的全部信息，转移只从更小且已经算清的状态来。复盘时把这个特例压成状态字段：状态下标、状态值语义、初始化、转移来源、遍历顺序；空间压缩时还要指出读到的是旧值还是新值；再用边界 单行、单列、权值为零、路径和溢出 反查初始化、更新顺序和退出条件。

\textbf{状态回放。} 手算路线：手算时先写一个很小的 DP 表或记忆化搜索记录。每个格子旁边写它来自哪些更小状态。状态字段：状态下标、状态值语义、初始化、转移来源、遍历顺序；空间压缩时还要指出读到的是旧值还是新值。状态升级：把重复递归节点命名为状态；遍历顺序只是在保证依赖状态已经算完。

\textbf{证明和边界。} 证明从最后一步开始：任意最优解的最后一步都在转移枚举中，转移来源已经由更小状态正确计算。边界：单行、单列、权值为零、路径和溢出。变体：若允许四方向移动，变成最短路而不是普通网格 DP。

\noindent\textbf{LeetCode 70 爬楼梯。} 模型：最后一步来自前一阶或前两阶，是最小的重叠子问题模型。推导重点：滚动变量保存最近两个状态即可，不需要完整数组。

\textbf{二刷读题。} 读题时先找重复子问题，而不是直接背公式。本题可以压成“最后一步来自前一阶或前两阶，是最小的重叠子问题模型”，滚动变量保存最近两个状态即可，不需要完整数组就是状态之间的依赖关系。先遮住答案，只保留题面，复述候选对象、合法性条件和输出目标。

\textbf{暴力回放。} 慢版本先写递归搜索树：节点表示处理位置、剩余资源和当前选择。只要不同路径反复到达同一个节点，就说明需要把这个节点命名成 DP 状态。二刷时先写慢版本或口述慢版本，用它确认不会漏答案。

\textbf{特例回放。} 拿 n=4 手算，最后一步要么从第 3 阶跨 1 阶，要么从第 2 阶跨 2 阶，这两类互斥且覆盖所有走法。于是 ways[4]=ways[3]+ways[2]。规律不是背斐波那契，而是最后一步只有两个来源，状态只需保存前两项。把这个特例继续拆开看：先问暴力搜索里哪些不同路径到达同一个后续问题，再给这个后续问题命名为状态。状态必须包含未来决策需要的全部信息，转移只从更小且已经算清的状态来。复盘时把这个特例压成状态字段：状态下标、状态值语义、初始化、转移来源、遍历顺序；空间压缩时还要指出读到的是旧值还是新值；再用边界 n 为一、n 为二、题目是否允许零阶、结果范围 反查初始化、更新顺序和退出条件。

\textbf{状态回放。} 手算路线：手算时先写一个很小的 DP 表或记忆化搜索记录。每个格子旁边写它来自哪些更小状态。状态字段：状态下标、状态值语义、初始化、转移来源、遍历顺序；空间压缩时还要指出读到的是旧值还是新值。状态升级：把重复递归节点命名为状态；遍历顺序只是在保证依赖状态已经算完。

\textbf{证明和边界。} 证明从最后一步开始：任意最优解的最后一步都在转移枚举中，转移来源已经由更小状态正确计算。边界：n 为一、n 为二、题目是否允许零阶、结果范围。变体：若每次可走一到 k 阶，转移变成固定宽度窗口求和。

\noindent\textbf{LeetCode 72 编辑距离。} 模型：二维状态表示两个前缀之间的最少编辑操作。推导重点：末尾字符相同继承左上；不同则枚举插入、删除、替换。

\textbf{二刷读题。} 读题时先找重复子问题，而不是直接背公式。本题可以压成“二维状态表示两个前缀之间的最少编辑操作”，末尾字符相同继承左上；不同则枚举插入、删除、替换就是状态之间的依赖关系。先遮住答案，只保留题面，复述候选对象、合法性条件和输出目标。

\textbf{暴力回放。} 慢版本先写递归搜索树：节点表示处理位置、剩余资源和当前选择。只要不同路径反复到达同一个节点，就说明需要把这个节点命名成 DP 状态。二刷时先写慢版本或口述慢版本，用它确认不会漏答案。

\textbf{特例回放。} 拿 word1=ab、word2=ac 手算最后一个字符。若 b 改成 c，是 dp[1][1]+1；若删掉 b，是 dp[1][2]+1；若插入 c，是 dp[2][1]+1。字符相等时才直接继承左上。规律是每个格子枚举最后一次编辑操作，而不是凭感觉填三邻居最小。把这个特例继续拆开看：先问暴力搜索里哪些不同路径到达同一个后续问题，再给这个后续问题命名为状态。状态必须包含未来决策需要的全部信息，转移只从更小且已经算清的状态来。复盘时把这个特例压成状态字段：状态下标、状态值语义、初始化、转移来源、遍历顺序；空间压缩时还要指出读到的是旧值还是新值；再用边界 空串、完全相同、完全不同、操作代价不同 反查初始化、更新顺序和退出条件。

\textbf{状态回放。} 手算路线：手算时先写一个很小的 DP 表或记忆化搜索记录。每个格子旁边写它来自哪些更小状态。状态字段：状态下标、状态值语义、初始化、转移来源、遍历顺序；空间压缩时还要指出读到的是旧值还是新值。状态升级：把重复递归节点命名为状态；遍历顺序只是在保证依赖状态已经算完。

\textbf{证明和边界。} 证明从最后一步开始：任意最优解的最后一步都在转移枚举中，转移来源已经由更小状态正确计算。边界：空串、完全相同、完全不同、操作代价不同。变体：若只允许插入删除，问题退化到和最长公共子序列相关。

\noindent\textbf{LeetCode 97 交错字符串。} 模型：目标串前缀由两个来源字符串的前缀交错组成。推导重点：状态 dp[i][j] 表示 s1 前 i 个字符和 s2 前 j 个字符能否组成 s3 前 i 加 j 个字符。

\textbf{二刷读题。} 读题时先找重复子问题，而不是直接背公式。本题可以压成“目标串前缀由两个来源字符串的前缀交错组成”，状态 dp[i][j] 表示 s1 前 i 个字符和 s2 前 j 个字符能否组成 s3 前 i 加 j 个字符就是状态之间的依赖关系。先遮住答案，只保留题面，复述候选对象、合法性条件和输出目标。

\textbf{暴力回放。} 慢版本先写递归搜索树：节点表示处理位置、剩余资源和当前选择。只要不同路径反复到达同一个节点，就说明需要把这个节点命名成 DP 状态。二刷时先写慢版本或口述慢版本，用它确认不会漏答案。

\textbf{特例回放。} 拿 s1=a、s2=b、s3=ab 手算，最后一个字符 b 可以来自 s2；拿 s1=a、s2=a、s3=aa，两个来源都能匹配，不能贪心固定选一个。规律是 dp[i][j] 表示两个前缀能否组成目标前 i+j 个字符，转移同时尝试来自 s1 和 s2 的最后一步。把这个特例继续拆开看：先问暴力搜索里哪些不同路径到达同一个后续问题，再给这个后续问题命名为状态。状态必须包含未来决策需要的全部信息，转移只从更小且已经算清的状态来。复盘时把这个特例压成状态字段：状态下标、状态值语义、初始化、转移来源、遍历顺序；空间压缩时还要指出读到的是旧值还是新值；再用边界 总长度不等、第一行第一列初始化、两个来源字符都能匹配时不能贪心只选一个 反查初始化、更新顺序和退出条件。

\textbf{状态回放。} 手算路线：手算时先写一个很小的 DP 表或记忆化搜索记录。每个格子旁边写它来自哪些更小状态。状态字段：状态下标、状态值语义、初始化、转移来源、遍历顺序；空间压缩时还要指出读到的是旧值还是新值。状态升级：把重复递归节点命名为状态；遍历顺序只是在保证依赖状态已经算完。

\textbf{证明和边界。} 证明从最后一步开始：任意最优解的最后一步都在转移枚举中，转移来源已经由更小状态正确计算。边界：总长度不等、第一行第一列初始化、两个来源字符都能匹配时不能贪心只选一个。变体：若要统计交错方案数，布尔状态要改成计数状态并处理大数或取模。

\noindent\textbf{LeetCode 115 不同的子序列。} 模型：从源串中删除若干字符后形成目标串，本质是两个前缀之间的计数问题。推导重点：状态 dp[i][j] 表示源串前 i 个字符中有多少个子序列等于目标前 j 个字符。

\textbf{二刷读题。} 读题时先找重复子问题，而不是直接背公式。本题可以压成“从源串中删除若干字符后形成目标串，本质是两个前缀之间的计数问题”，状态 dp[i][j] 表示源串前 i 个字符中有多少个子序列等于目标前 j 个字符就是状态之间的依赖关系。先遮住答案，只保留题面，复述候选对象、合法性条件和输出目标。

\textbf{暴力回放。} 慢版本先写递归搜索树：节点表示处理位置、剩余资源和当前选择。只要不同路径反复到达同一个节点，就说明需要把这个节点命名成 DP 状态。二刷时先写慢版本或口述慢版本，用它确认不会漏答案。

\textbf{特例回放。} 拿 s=bab、t=ba 手算。扫描到最后一个 b 时，它不能匹配 t 的 a，只能继承不用它的方案；扫描到 a 时，有两类方案：用这个 a 匹配 t 的末尾，或不用它。规律是字符相等时计数相加，不等时只继承不用源字符的方案。把这个特例继续拆开看：先问暴力搜索里哪些不同路径到达同一个后续问题，再给这个后续问题命名为状态。状态必须包含未来决策需要的全部信息，转移只从更小且已经算清的状态来。复盘时把这个特例压成状态字段：状态下标、状态值语义、初始化、转移来源、遍历顺序；空间压缩时还要指出读到的是旧值还是新值；再用边界 空目标串计数为一、源串为空、重复字符导致计数增长、计数可能超过 int 反查初始化、更新顺序和退出条件。

\textbf{状态回放。} 手算路线：手算时先写一个很小的 DP 表或记忆化搜索记录。每个格子旁边写它来自哪些更小状态。状态字段：状态下标、状态值语义、初始化、转移来源、遍历顺序；空间压缩时还要指出读到的是旧值还是新值。状态升级：把重复递归节点命名为状态；遍历顺序只是在保证依赖状态已经算完。

\textbf{证明和边界。} 证明从最后一步开始：任意最优解的最后一步都在转移枚举中，转移来源已经由更小状态正确计算。边界：空目标串计数为一、源串为空、重复字符导致计数增长、计数可能超过 int。变体：若只问是否存在目标子序列，可以用双指针，不需要二维计数 DP。

\noindent\textbf{LeetCode 121 买卖股票的最佳时机。} 模型：卖出日固定时，最佳买入日是此前最低价格。推导重点：扫描维护历史最低价和最大利润，一次交易无需复杂状态机。

\textbf{二刷读题。} 读题时先找重复子问题，而不是直接背公式。本题可以压成“卖出日固定时，最佳买入日是此前最低价格”，扫描维护历史最低价和最大利润，一次交易无需复杂状态机就是状态之间的依赖关系。先遮住答案，只保留题面，复述候选对象、合法性条件和输出目标。

\textbf{暴力回放。} 慢版本先写递归搜索树：节点表示处理位置、剩余资源和当前选择。只要不同路径反复到达同一个节点，就说明需要把这个节点命名成 DP 状态。二刷时先写慢版本或口述慢版本，用它确认不会漏答案。

\textbf{特例回放。} 拿价格 [7,1,5] 手算。站在 7 不能卖出获利；站在 1 更新最低买入价；站在 5 时利润是 5-1。规律是固定卖出日后，最佳买入日只需要历史最低价，状态不是所有买卖对。把这个特例继续拆开看：先问暴力搜索里哪些不同路径到达同一个后续问题，再给这个后续问题命名为状态。状态必须包含未来决策需要的全部信息，转移只从更小且已经算清的状态来。复盘时把这个特例压成状态字段：状态下标、状态值语义、初始化、转移来源、遍历顺序；空间压缩时还要指出读到的是旧值还是新值；再用边界 单日价格、一直下跌、一直上涨、价格相同 反查初始化、更新顺序和退出条件。

\textbf{状态回放。} 手算路线：手算时先写一个很小的 DP 表或记忆化搜索记录。每个格子旁边写它来自哪些更小状态。状态字段：状态下标、状态值语义、初始化、转移来源、遍历顺序；空间压缩时还要指出读到的是旧值还是新值。状态升级：把重复递归节点命名为状态；遍历顺序只是在保证依赖状态已经算完。

\textbf{证明和边界。} 证明从最后一步开始：任意最优解的最后一步都在转移枚举中，转移来源已经由更小状态正确计算。边界：单日价格、一直下跌、一直上涨、价格相同。变体：多次交易、冷冻期和手续费版本都要扩展成状态机。

\noindent\textbf{LeetCode 123 买卖股票 III。} 模型：最多两次交易，状态需要包含交易次数和持有状态。推导重点：维护 buy1、sell1、buy2、sell2 四个状态，按天更新。

\textbf{二刷读题。} 读题时先找重复子问题，而不是直接背公式。本题可以压成“最多两次交易，状态需要包含交易次数和持有状态”，维护 buy1、sell1、buy2、sell2 四个状态，按天更新就是状态之间的依赖关系。先遮住答案，只保留题面，复述候选对象、合法性条件和输出目标。

\textbf{暴力回放。} 慢版本先写递归搜索树：节点表示处理位置、剩余资源和当前选择。只要不同路径反复到达同一个节点，就说明需要把这个节点命名成 DP 状态。二刷时先写慢版本或口述慢版本，用它确认不会漏答案。

\textbf{特例回放。} 拿 [3,3,5,0,0,3,1,4] 手算，只保存一个最大利润不够，因为第二次买入依赖第一次卖出后的余额。规律是维护 buy1、sell1、buy2、sell2 四个状态，每天都表示到今天为止对应状态的最优值；两次交易不能重叠。把这个特例继续拆开看：先问暴力搜索里哪些不同路径到达同一个后续问题，再给这个后续问题命名为状态。状态必须包含未来决策需要的全部信息，转移只从更小且已经算清的状态来。复盘时把这个特例压成状态字段：状态下标、状态值语义、初始化、转移来源、遍历顺序；空间压缩时还要指出读到的是旧值还是新值；再用边界 价格为空、只够一次交易、两次交易间不能重叠 反查初始化、更新顺序和退出条件。

\textbf{状态回放。} 手算路线：手算时先写一个很小的 DP 表或记忆化搜索记录。每个格子旁边写它来自哪些更小状态。状态字段：状态下标、状态值语义、初始化、转移来源、遍历顺序；空间压缩时还要指出读到的是旧值还是新值。状态升级：把重复递归节点命名为状态；遍历顺序只是在保证依赖状态已经算完。

\textbf{证明和边界。} 证明从最后一步开始：任意最优解的最后一步都在转移枚举中，转移来源已经由更小状态正确计算。边界：价格为空、只够一次交易、两次交易间不能重叠。变体：最多 k 次交易是它的泛化，k 很大时退化为无限次交易。

\noindent\textbf{LeetCode 139 单词拆分。} 模型：状态表示前缀能否由字典单词拼出。推导重点：枚举最后一个单词的起点，左侧前缀可达且子串在字典中则当前可达。

\textbf{二刷读题。} 读题时先找重复子问题，而不是直接背公式。本题可以压成“状态表示前缀能否由字典单词拼出”，枚举最后一个单词的起点，左侧前缀可达且子串在字典中则当前可达就是状态之间的依赖关系。先遮住答案，只保留题面，复述候选对象、合法性条件和输出目标。

\textbf{暴力回放。} 慢版本先写递归搜索树：节点表示处理位置、剩余资源和当前选择。只要不同路径反复到达同一个节点，就说明需要把这个节点命名成 DP 状态。二刷时先写慢版本或口述慢版本，用它确认不会漏答案。

\textbf{特例回放。} 拿 s=leetcode、字典 [leet,code] 手算，dp[4] 因为 leet 可拆而为真，dp[8] 又因为 dp[4] 且 code 在字典里为真。规律是 dp[i] 表示前 i 个字符能否拆分，转移枚举最后一个单词起点。把这个特例继续拆开看：先问暴力搜索里哪些不同路径到达同一个后续问题，再给这个后续问题命名为状态。状态必须包含未来决策需要的全部信息，转移只从更小且已经算清的状态来。复盘时把这个特例压成状态字段：状态下标、状态值语义、初始化、转移来源、遍历顺序；空间压缩时还要指出读到的是旧值还是新值；再用边界 空串、重复单词、长单词、substr 复制成本 反查初始化、更新顺序和退出条件。

\textbf{状态回放。} 手算路线：手算时先写一个很小的 DP 表或记忆化搜索记录。每个格子旁边写它来自哪些更小状态。状态字段：状态下标、状态值语义、初始化、转移来源、遍历顺序；空间压缩时还要指出读到的是旧值还是新值。状态升级：把重复递归节点命名为状态；遍历顺序只是在保证依赖状态已经算完。

\textbf{证明和边界。} 证明从最后一步开始：任意最优解的最后一步都在转移枚举中，转移来源已经由更小状态正确计算。边界：空串、重复单词、长单词、substr 复制成本。变体：若要输出所有句子，需要 DFS 加记忆化或前驱列表。

\noindent\textbf{LeetCode 152 乘积最大子数组。} 模型：负数会让最大乘积和最小乘积互换。推导重点：同时维护以当前位置结尾的最大和最小乘积。

\textbf{二刷读题。} 读题时先找重复子问题，而不是直接背公式。本题可以压成“负数会让最大乘积和最小乘积互换”，同时维护以当前位置结尾的最大和最小乘积就是状态之间的依赖关系。先遮住答案，只保留题面，复述候选对象、合法性条件和输出目标。

\textbf{暴力回放。} 慢版本先写递归搜索树：节点表示处理位置、剩余资源和当前选择。只要不同路径反复到达同一个节点，就说明需要把这个节点命名成 DP 状态。二刷时先写慢版本或口述慢版本，用它确认不会漏答案。

\textbf{特例回放。} 拿 [-2,3,-4] 手算。到 3 时最大乘积是 3、最小是 -6；再乘 -4 后，原来的最小 -6 反而变成最大 24。规律是负数会交换最大和最小，状态必须同时保存以当前位置结尾的最大乘积和最小乘积。把这个特例继续拆开看：先问暴力搜索里哪些不同路径到达同一个后续问题，再给这个后续问题命名为状态。状态必须包含未来决策需要的全部信息，转移只从更小且已经算清的状态来。复盘时把这个特例压成状态字段：状态下标、状态值语义、初始化、转移来源、遍历顺序；空间压缩时还要指出读到的是旧值还是新值；再用边界 零、负数个数奇偶、单元素、乘积溢出 反查初始化、更新顺序和退出条件。

\textbf{状态回放。} 手算路线：手算时先写一个很小的 DP 表或记忆化搜索记录。每个格子旁边写它来自哪些更小状态。状态字段：状态下标、状态值语义、初始化、转移来源、遍历顺序；空间压缩时还要指出读到的是旧值还是新值。状态升级：把重复递归节点命名为状态；遍历顺序只是在保证依赖状态已经算完。

\textbf{证明和边界。} 证明从最后一步开始：任意最优解的最后一步都在转移枚举中，转移来源已经由更小状态正确计算。边界：零、负数个数奇偶、单元素、乘积溢出。变体：如果要求返回区间，同步记录状态来源。

\noindent\textbf{LeetCode 188 买卖股票 IV。} 模型：最多 k 次交易让状态必须同时包含交易次数和是否持有股票。推导重点：维护每个交易次数下的 buy 和 sell 状态；当 k 足够大时可退化为无限次交易。

\textbf{二刷读题。} 读题时先找重复子问题，而不是直接背公式。本题可以压成“最多 k 次交易让状态必须同时包含交易次数和是否持有股票”，维护每个交易次数下的 buy 和 sell 状态；当 k 足够大时可退化为无限次交易就是状态之间的依赖关系。先遮住答案，只保留题面，复述候选对象、合法性条件和输出目标。

\textbf{暴力回放。} 慢版本先写递归搜索树：节点表示处理位置、剩余资源和当前选择。只要不同路径反复到达同一个节点，就说明需要把这个节点命名成 DP 状态。二刷时先写慢版本或口述慢版本，用它确认不会漏答案。

\textbf{特例回放。} 拿 k=2、价格 [3,2,6] 手算。某天结束时只说最大利润不够，因为未来能否卖出取决于是否持股，也取决于已经完成几次交易。规律是状态必须包含交易次数和持有状态；k 很大时才退化为无限次交易。把这个特例继续拆开看：先问暴力搜索里哪些不同路径到达同一个后续问题，再给这个后续问题命名为状态。状态必须包含未来决策需要的全部信息，转移只从更小且已经算清的状态来。复盘时把这个特例压成状态字段：状态下标、状态值语义、初始化、转移来源、遍历顺序；空间压缩时还要指出读到的是旧值还是新值；再用边界 k 为零、价格天数很少、k 大于等于 n/2、更新顺序造成同一天状态污染 反查初始化、更新顺序和退出条件。

\textbf{状态回放。} 手算路线：手算时先写一个很小的 DP 表或记忆化搜索记录。每个格子旁边写它来自哪些更小状态。状态字段：状态下标、状态值语义、初始化、转移来源、遍历顺序；空间压缩时还要指出读到的是旧值还是新值。状态升级：把重复递归节点命名为状态；遍历顺序只是在保证依赖状态已经算完。

\textbf{证明和边界。} 证明从最后一步开始：任意最优解的最后一步都在转移枚举中，转移来源已经由更小状态正确计算。边界：k 为零、价格天数很少、k 大于等于 n/2、更新顺序造成同一天状态污染。变体：手续费、冷冻期和最多两次交易都是同一状态机框架的不同约束。

\noindent\textbf{LeetCode 198 打家劫舍。} 模型：每间房只有偷和不偷两种结局，偷当前就不能偷前一间。推导重点：滚动保存前一状态和前两状态。

\textbf{二刷读题。} 读题时先找重复子问题，而不是直接背公式。本题可以压成“每间房只有偷和不偷两种结局，偷当前就不能偷前一间”，滚动保存前一状态和前两状态就是状态之间的依赖关系。先遮住答案，只保留题面，复述候选对象、合法性条件和输出目标。

\textbf{暴力回放。} 慢版本先写递归搜索树：节点表示处理位置、剩余资源和当前选择。只要不同路径反复到达同一个节点，就说明需要把这个节点命名成 DP 状态。二刷时先写慢版本或口述慢版本，用它确认不会漏答案。

\textbf{特例回放。} 拿 [2,7,9] 手算，到房屋 9 时只有两类合法结局：不偷它，答案沿用前一间；偷它，就只能接前两间的最优。规律是 dp[i]=max(dp[i-1], dp[i-2]+nums[i])，相邻限制来自题面而不是公式本身。把这个特例继续拆开看：先问暴力搜索里哪些不同路径到达同一个后续问题，再给这个后续问题命名为状态。状态必须包含未来决策需要的全部信息，转移只从更小且已经算清的状态来。复盘时把这个特例压成状态字段：状态下标、状态值语义、初始化、转移来源、遍历顺序；空间压缩时还要指出读到的是旧值还是新值；再用边界 空数组、单间、全零、金额很大 反查初始化、更新顺序和退出条件。

\textbf{状态回放。} 手算路线：手算时先写一个很小的 DP 表或记忆化搜索记录。每个格子旁边写它来自哪些更小状态。状态字段：状态下标、状态值语义、初始化、转移来源、遍历顺序；空间压缩时还要指出读到的是旧值还是新值。状态升级：把重复递归节点命名为状态；遍历顺序只是在保证依赖状态已经算完。

\textbf{证明和边界。} 证明从最后一步开始：任意最优解的最后一步都在转移枚举中，转移来源已经由更小状态正确计算。边界：空数组、单间、全零、金额很大。变体：环形房屋要拆成不含首和不含尾两个线性问题。

\noindent\textbf{LeetCode 221 最大正方形。} 模型：以当前位置为右下角的最大正方形由上、左、左上三个状态限制。推导重点：当前位置为一时取三者最小加一，否则为零。

\textbf{二刷读题。} 读题时先找重复子问题，而不是直接背公式。本题可以压成“以当前位置为右下角的最大正方形由上、左、左上三个状态限制”，当前位置为一时取三者最小加一，否则为零就是状态之间的依赖关系。先遮住答案，只保留题面，复述候选对象、合法性条件和输出目标。

\textbf{暴力回放。} 慢版本先写递归搜索树：节点表示处理位置、剩余资源和当前选择。只要不同路径反复到达同一个节点，就说明需要把这个节点命名成 DP 状态。二刷时先写慢版本或口述慢版本，用它确认不会漏答案。

\textbf{特例回放。} 拿 2x2 全 1 方块手算，右下角能形成边长 2，因为上、左、左上三个格子都至少能形成边长 1。若其中任一方向为 0，只能形成边长 1。规律是 dp[row][col]=min(上,左,左上)+1，状态值是以当前格为右下角的最大正方形边长。把这个特例继续拆开看：先问暴力搜索里哪些不同路径到达同一个后续问题，再给这个后续问题命名为状态。状态必须包含未来决策需要的全部信息，转移只从更小且已经算清的状态来。复盘时把这个特例压成状态字段：状态下标、状态值语义、初始化、转移来源、遍历顺序；空间压缩时还要指出读到的是旧值还是新值；再用边界 全零、全一、单行单列、字符一和数字一不要混淆 反查初始化、更新顺序和退出条件。

\textbf{状态回放。} 手算路线：手算时先写一个很小的 DP 表或记忆化搜索记录。每个格子旁边写它来自哪些更小状态。状态字段：状态下标、状态值语义、初始化、转移来源、遍历顺序；空间压缩时还要指出读到的是旧值还是新值。状态升级：把重复递归节点命名为状态；遍历顺序只是在保证依赖状态已经算完。

\textbf{证明和边界。} 证明从最后一步开始：任意最优解的最后一步都在转移枚举中，转移来源已经由更小状态正确计算。边界：全零、全一、单行单列、字符一和数字一不要混淆。变体：最大矩形则需要按行构造柱状图再用单调栈。

\noindent\textbf{LeetCode 300 最长递增子序列。} 模型：二次 DP 固定结尾，二分优化维护每个长度的最小结尾。推导重点：tails 不是实际答案序列，而是未来扩展潜力表。

\textbf{二刷读题。} 读题时先找重复子问题，而不是直接背公式。本题可以压成“二次 DP 固定结尾，二分优化维护每个长度的最小结尾”，tails 不是实际答案序列，而是未来扩展潜力表就是状态之间的依赖关系。先遮住答案，只保留题面，复述候选对象、合法性条件和输出目标。

\textbf{暴力回放。} 慢版本先写递归搜索树：节点表示处理位置、剩余资源和当前选择。只要不同路径反复到达同一个节点，就说明需要把这个节点命名成 DP 状态。二刷时先写慢版本或口述慢版本，用它确认不会漏答案。

\textbf{特例回放。} 拿 [10,9,2,5,3] 手算。二次 DP 固定每个位置为结尾，5 可接在 2 后；优化时 tails[长度] 保存该长度递增子序列的最小尾值，尾值越小未来越容易接。规律是二分替换 tails，不代表真实序列完整内容，而是维护未来潜力。把这个特例继续拆开看：先问暴力搜索里哪些不同路径到达同一个后续问题，再给这个后续问题命名为状态。状态必须包含未来决策需要的全部信息，转移只从更小且已经算清的状态来。复盘时把这个特例压成状态字段：状态下标、状态值语义、初始化、转移来源、遍历顺序；空间压缩时还要指出读到的是旧值还是新值；再用边界 重复元素、严格递增和非递减区别、空数组 反查初始化、更新顺序和退出条件。

\textbf{状态回放。} 手算路线：手算时先写一个很小的 DP 表或记忆化搜索记录。每个格子旁边写它来自哪些更小状态。状态字段：状态下标、状态值语义、初始化、转移来源、遍历顺序；空间压缩时还要指出读到的是旧值还是新值。状态升级：把重复递归节点命名为状态；遍历顺序只是在保证依赖状态已经算完。

\textbf{证明和边界。} 证明从最后一步开始：任意最优解的最后一步都在转移枚举中，转移来源已经由更小状态正确计算。边界：重复元素、严格递增和非递减区别、空数组。变体：若要还原路径，需要额外前驱数组，不能只看 tails。

\noindent\textbf{LeetCode 309 最佳买卖股票含冷冻期。} 模型：冷冻期让买入来源受昨天卖出状态限制。推导重点：状态机维护持有、刚卖出、休息三种日终状态。

\textbf{二刷读题。} 读题时先找重复子问题，而不是直接背公式。本题可以压成“冷冻期让买入来源受昨天卖出状态限制”，状态机维护持有、刚卖出、休息三种日终状态就是状态之间的依赖关系。先遮住答案，只保留题面，复述候选对象、合法性条件和输出目标。

\textbf{暴力回放。} 慢版本先写递归搜索树：节点表示处理位置、剩余资源和当前选择。只要不同路径反复到达同一个节点，就说明需要把这个节点命名成 DP 状态。二刷时先写慢版本或口述慢版本，用它确认不会漏答案。

\textbf{特例回放。} 拿 [1,2,3,0,2] 手算，第 2 天卖出后第 3 天不能立刻买入，因为有冷冻期；状态必须区分持股、刚卖出、空仓可买。规律是冷冻期把普通买卖股票拆成多个日末状态，转移不能跨过限制。把这个特例继续拆开看：先问暴力搜索里哪些不同路径到达同一个后续问题，再给这个后续问题命名为状态。状态必须包含未来决策需要的全部信息，转移只从更小且已经算清的状态来。复盘时把这个特例压成状态字段：状态下标、状态值语义、初始化、转移来源、遍历顺序；空间压缩时还要指出读到的是旧值还是新值；再用边界 空价格、一天、连续上涨、卖出后不能次日买入 反查初始化、更新顺序和退出条件。

\textbf{状态回放。} 手算路线：手算时先写一个很小的 DP 表或记忆化搜索记录。每个格子旁边写它来自哪些更小状态。状态字段：状态下标、状态值语义、初始化、转移来源、遍历顺序；空间压缩时还要指出读到的是旧值还是新值。状态升级：把重复递归节点命名为状态；遍历顺序只是在保证依赖状态已经算完。

\textbf{证明和边界。} 证明从最后一步开始：任意最优解的最后一步都在转移枚举中，转移来源已经由更小状态正确计算。边界：空价格、一天、连续上涨、卖出后不能次日买入。变体：手续费版本可以在买入或卖出转移中扣费。

\noindent\textbf{LeetCode 312 戳气球。} 模型：先戳哪个会让邻居持续变化，改成最后戳某个气球后区间边界才稳定。推导重点：区间 DP 定义开区间内全部戳完的最大收益，枚举最后一个被戳的气球作为切分点。

\textbf{二刷读题。} 读题时先找重复子问题，而不是直接背公式。本题可以压成“先戳哪个会让邻居持续变化，改成最后戳某个气球后区间边界才稳定”，区间 DP 定义开区间内全部戳完的最大收益，枚举最后一个被戳的气球作为切分点就是状态之间的依赖关系。先遮住答案，只保留题面，复述候选对象、合法性条件和输出目标。

\textbf{暴力回放。} 慢版本先写递归搜索树：节点表示处理位置、剩余资源和当前选择。只要不同路径反复到达同一个节点，就说明需要把这个节点命名成 DP 状态。二刷时先写慢版本或口述慢版本，用它确认不会漏答案。

\textbf{特例回放。} 拿 [3,1,5] 手算，若先戳 1，左右邻居会变；若改问某个区间里最后戳谁，最后一刻左右边界固定，收益可拆成左右两个子区间。规律是区间 DP 枚举最后戳的 mid，而不是第一个戳的气球。把这个特例继续拆开看：先问暴力搜索里哪些不同路径到达同一个后续问题，再给这个后续问题命名为状态。状态必须包含未来决策需要的全部信息，转移只从更小且已经算清的状态来。复盘时把这个特例压成状态字段：状态下标、状态值语义、初始化、转移来源、遍历顺序；空间压缩时还要指出读到的是旧值还是新值；再用边界 补一的虚拟边界、区间长度遍历顺序、空区间收益为零、乘积可能较大 反查初始化、更新顺序和退出条件。

\textbf{状态回放。} 手算路线：手算时先写一个很小的 DP 表或记忆化搜索记录。每个格子旁边写它来自哪些更小状态。状态字段：状态下标、状态值语义、初始化、转移来源、遍历顺序；空间压缩时还要指出读到的是旧值还是新值。状态升级：把重复递归节点命名为状态；遍历顺序只是在保证依赖状态已经算完。

\textbf{证明和边界。} 证明从最后一步开始：任意最优解的最后一步都在转移枚举中，转移来源已经由更小状态正确计算。边界：补一的虚拟边界、区间长度遍历顺序、空区间收益为零、乘积可能较大。变体：很多区间 DP 都要换成最后一步思考，避免中间操作改变边界。

\noindent\textbf{LeetCode 329 矩阵中的最长递增路径。} 模型：每个格子的答案是从它出发能走出的最长严格递增路径。推导重点：把严格递增关系看成无环状态图，用记忆化 DFS 或按值排序的 DAG DP 复用后续格子结果。

\textbf{二刷读题。} 读题时先找重复子问题，而不是直接背公式。本题可以压成“每个格子的答案是从它出发能走出的最长严格递增路径”，把严格递增关系看成无环状态图，用记忆化 DFS 或按值排序的 DAG DP 复用后续格子结果就是状态之间的依赖关系。先遮住答案，只保留题面，复述候选对象、合法性条件和输出目标。

\textbf{暴力回放。} 慢版本先写递归搜索树：节点表示处理位置、剩余资源和当前选择。只要不同路径反复到达同一个节点，就说明需要把这个节点命名成 DP 状态。二刷时先写慢版本或口述慢版本，用它确认不会漏答案。

\textbf{特例回放。} 拿矩阵中 1->2->3 的小路径手算，从 1 出发会用到从 2 出发的答案；如果另一个格子也走到 2，后缀路径被重复计算。严格递增保证不会成环。规律是每个格子记忆化保存从这里出发的最长路径。把这个特例继续拆开看：先问暴力搜索里哪些不同路径到达同一个后续问题，再给这个后续问题命名为状态。状态必须包含未来决策需要的全部信息，转移只从更小且已经算清的状态来。复盘时把这个特例压成状态字段：状态下标、状态值语义、初始化、转移来源、遍历顺序；空间压缩时还要指出读到的是旧值还是新值；再用边界 平台相等不能走、四方向越界、递归深度、允许不下降时会产生环 反查初始化、更新顺序和退出条件。

\textbf{状态回放。} 手算路线：手算时先写一个很小的 DP 表或记忆化搜索记录。每个格子旁边写它来自哪些更小状态。状态字段：状态下标、状态值语义、初始化、转移来源、遍历顺序；空间压缩时还要指出读到的是旧值还是新值。状态升级：把重复递归节点命名为状态；遍历顺序只是在保证依赖状态已经算完。

\textbf{证明和边界。} 证明从最后一步开始：任意最优解的最后一步都在转移枚举中，转移来源已经由更小状态正确计算。边界：平台相等不能走、四方向越界、递归深度、允许不下降时会产生环。变体：若改成求最短路径，严格递增的 DP 模型不再直接适用。

\noindent\textbf{LeetCode 714 买卖股票含手续费。} 模型：手续费改变交易收益，仍可用持有和不持有状态。推导重点：卖出时扣费或买入时扣费都可以，但不要重复扣。

\textbf{二刷读题。} 读题时先找重复子问题，而不是直接背公式。本题可以压成“手续费改变交易收益，仍可用持有和不持有状态”，卖出时扣费或买入时扣费都可以，但不要重复扣就是状态之间的依赖关系。先遮住答案，只保留题面，复述候选对象、合法性条件和输出目标。

\textbf{暴力回放。} 慢版本先写递归搜索树：节点表示处理位置、剩余资源和当前选择。只要不同路径反复到达同一个节点，就说明需要把这个节点命名成 DP 状态。二刷时先写慢版本或口述慢版本，用它确认不会漏答案。

\textbf{特例回放。} 拿 [1,3,2,8,4,9]、fee=2 手算，8 卖出时净赚 5；手续费让频繁小波动不一定值得交易。规律是每天维护 hold 和 cash，卖出时扣手续费，买入时从 cash 转到 hold。把这个特例继续拆开看：先问暴力搜索里哪些不同路径到达同一个后续问题，再给这个后续问题命名为状态。状态必须包含未来决策需要的全部信息，转移只从更小且已经算清的状态来。复盘时把这个特例压成状态字段：状态下标、状态值语义、初始化、转移来源、遍历顺序；空间压缩时还要指出读到的是旧值还是新值；再用边界 手续费为零、价格震荡、小利润不足手续费 反查初始化、更新顺序和退出条件。

\textbf{状态回放。} 手算路线：手算时先写一个很小的 DP 表或记忆化搜索记录。每个格子旁边写它来自哪些更小状态。状态字段：状态下标、状态值语义、初始化、转移来源、遍历顺序；空间压缩时还要指出读到的是旧值还是新值。状态升级：把重复递归节点命名为状态；遍历顺序只是在保证依赖状态已经算完。

\textbf{证明和边界。} 证明从最后一步开始：任意最优解的最后一步都在转移枚举中，转移来源已经由更小状态正确计算。边界：手续费为零、价格震荡、小利润不足手续费。变体：这是状态机比局部贪心更稳的例子。

\noindent\textbf{LeetCode 1143 最长公共子序列。} 模型：状态表示两个前缀的最长公共子序列长度。推导重点：末尾相同取左上加一，不同取上方和左方最大。

\textbf{二刷读题。} 读题时先找重复子问题，而不是直接背公式。本题可以压成“状态表示两个前缀的最长公共子序列长度”，末尾相同取左上加一，不同取上方和左方最大就是状态之间的依赖关系。先遮住答案，只保留题面，复述候选对象、合法性条件和输出目标。

\textbf{暴力回放。} 慢版本先写递归搜索树：节点表示处理位置、剩余资源和当前选择。只要不同路径反复到达同一个节点，就说明需要把这个节点命名成 DP 状态。二刷时先写慢版本或口述慢版本，用它确认不会漏答案。

\textbf{特例回放。} 拿 abcde 和 ace 手算，末尾 e 相等时答案来自两个前缀 abcd 和 ac 加一；若末尾不同，就取删掉一侧末尾后的较大值。规律是 dp[i][j] 表示两个前缀的 LCS 长度，转移只看左、上、左上。把这个特例继续拆开看：先问暴力搜索里哪些不同路径到达同一个后续问题，再给这个后续问题命名为状态。状态必须包含未来决策需要的全部信息，转移只从更小且已经算清的状态来。复盘时把这个特例压成状态字段：状态下标、状态值语义、初始化、转移来源、遍历顺序；空间压缩时还要指出读到的是旧值还是新值；再用边界 空串、完全相同、无公共字符、子序列不是子串 反查初始化、更新顺序和退出条件。

\textbf{状态回放。} 手算路线：手算时先写一个很小的 DP 表或记忆化搜索记录。每个格子旁边写它来自哪些更小状态。状态字段：状态下标、状态值语义、初始化、转移来源、遍历顺序；空间压缩时还要指出读到的是旧值还是新值。状态升级：把重复递归节点命名为状态；遍历顺序只是在保证依赖状态已经算完。

\textbf{证明和边界。} 证明从最后一步开始：任意最优解的最后一步都在转移枚举中，转移来源已经由更小状态正确计算。边界：空串、完全相同、无公共字符、子序列不是子串。变体：若要求还原序列，需要从表右下角反向追踪选择。

\noindent\textbf{LeetCode 1235 规划兼职工作。} 模型：每份工作有开始、结束和收益，选择不冲突工作使收益最大。推导重点：按结束时间排序，dp[i] 在不选当前和选当前加上兼容前驱之间取最大，前驱用二分找。

\textbf{二刷读题。} 读题时先找重复子问题，而不是直接背公式。本题可以压成“每份工作有开始、结束和收益，选择不冲突工作使收益最大”，按结束时间排序，dp[i] 在不选当前和选当前加上兼容前驱之间取最大，前驱用二分找就是状态之间的依赖关系。先遮住答案，只保留题面，复述候选对象、合法性条件和输出目标。

\textbf{暴力回放。} 慢版本先写递归搜索树：节点表示处理位置、剩余资源和当前选择。只要不同路径反复到达同一个节点，就说明需要把这个节点命名成 DP 状态。二刷时先写慢版本或口述慢版本，用它确认不会漏答案。

\textbf{特例回放。} 拿 jobs (1,3,50),(2,4,10),(3,5,40),(3,6,70) 手算，选 (1,3,50) 后还能接 (3,6,70)，总 120。规律是按结束时间排序，dp[i] 在不选当前和选当前加上前一个不冲突工作之间取最大。把这个特例继续拆开看：先问暴力搜索里哪些不同路径到达同一个后续问题，再给这个后续问题命名为状态。状态必须包含未来决策需要的全部信息，转移只从更小且已经算清的状态来。复盘时把这个特例压成状态字段：状态下标、状态值语义、初始化、转移来源、遍历顺序；空间压缩时还要指出读到的是旧值还是新值；再用边界 结束等于开始是否兼容、收益相同、工作排序、空列表 反查初始化、更新顺序和退出条件。

\textbf{状态回放。} 手算路线：手算时先写一个很小的 DP 表或记忆化搜索记录。每个格子旁边写它来自哪些更小状态。状态字段：状态下标、状态值语义、初始化、转移来源、遍历顺序；空间压缩时还要指出读到的是旧值还是新值。状态升级：把重复递归节点命名为状态；遍历顺序只是在保证依赖状态已经算完。

\textbf{证明和边界。} 证明从最后一步开始：任意最优解的最后一步都在转移枚举中，转移来源已经由更小状态正确计算。边界：结束等于开始是否兼容、收益相同、工作排序、空列表。变体：这是排序、二分和 DP 的组合，不是按收益贪心。

\subsection{自测标准}

二刷时不要只确认自己见过这道题。请遮住答案，先讲题面模型，再讲暴力基线和瓶颈，然后说出优化结构保存了什么、不变量如何排除错误候选、边界实验如何打破错误实现。

\section{背包模型：二刷复盘卡}
这一大节只围绕一个题型复盘，不把题号直接铺成目录。阅读时先看复盘目标，再读核心题卡，最后用自测标准检查自己是否能独立重讲。

\subsection{复盘目标}

追问通常会问倒序和正序的区别、组合和排列的区别、为什么复杂度是伪多项式。请准备一个循环方向写反的最小反例。

\subsection{核心题卡}

\noindent\textbf{LeetCode 279 完全平方数。} 模型：完全平方数可以看成无限使用的物品，目标是凑出 n 的最少数量。推导重点：完全背包最少物品数，或把数字看成图节点做 BFS。

\textbf{二刷读题。} 读题时先找阶段、容量、物品使用次数和状态值类型。本题可以压成“完全平方数可以看成无限使用的物品，目标是凑出 n 的最少数量”，循环方向由这些条件决定。先遮住答案，只保留题面，复述候选对象、合法性条件和输出目标。

\textbf{暴力回放。} 慢版本枚举每个物品选不选、选几次，并记录阶段和剩余容量。相同阶段和容量反复出现时，就可以把指数搜索压成表。二刷时先写慢版本或口述慢版本，用它确认不会漏答案。

\textbf{特例回放。} 拿 n=12 手算，12 可以由 4+4+4 组成，不能只贪心选 9 后剩 3。规律是完全背包，dp[x] 表示和为 x 的最少平方数，枚举每个平方数时允许重复使用。把这个特例继续拆开看：阶段是物品，容量是资源，状态值是可达、最值还是计数。一个物品能否在同一轮再次使用，直接决定容量循环方向；组合和排列也由循环顺序表达。复盘时把这个特例压成状态字段：物品阶段、容量、状态值、循环方向、取模或不可达哨兵；组合和排列必须用不同循环顺序表达；再用边界 n 为零或一、完全平方数本身、较大 n 反查初始化、更新顺序和退出条件。

\textbf{状态回放。} 手算路线：手算时用一个容量很小的例子比较正序和倒序。只要循环方向改变后同一物品能否重复使用发生变化，就说明方向不是细节。状态字段：物品阶段、容量、状态值、循环方向、取模或不可达哨兵；组合和排列必须用不同循环顺序表达。状态升级：把选择树压成阶段和容量；物品是否可重复使用决定一维表的遍历方向。

\textbf{证明和边界。} 证明从当前物品使用次数开始：0-1、完全、计数和最值的循环语义必须和题意一致。边界：n 为零或一、完全平方数本身、较大 n。变体：数学四平方定理可给更强解法，但面试 DP 更通用。

\noindent\textbf{LeetCode 322 零钱兑换。} 模型：金额是容量，硬币可无限使用，目标是最少硬币数。推导重点：dp[value] 从 value 减 coin 的已知状态转移。

\textbf{二刷读题。} 读题时先找阶段、容量、物品使用次数和状态值类型。本题可以压成“金额是容量，硬币可无限使用，目标是最少硬币数”，循环方向由这些条件决定。先遮住答案，只保留题面，复述候选对象、合法性条件和输出目标。

\textbf{暴力回放。} 慢版本枚举每个物品选不选、选几次，并记录阶段和剩余容量。相同阶段和容量反复出现时，就可以把指数搜索压成表。二刷时先写慢版本或口述慢版本，用它确认不会漏答案。

\textbf{特例回放。} 拿 coins=[1,2,5]、amount=11 手算，最后一步可以用 1、2 或 5，若用 5，则来源是 dp[6]+1。规律是 dp[x] 表示凑成 x 的最少硬币数，完全背包允许同一硬币重复使用；不可达状态不能参与加一。把这个特例继续拆开看：阶段是物品，容量是资源，状态值是可达、最值还是计数。一个物品能否在同一轮再次使用，直接决定容量循环方向；组合和排列也由循环顺序表达。复盘时把这个特例压成状态字段：物品阶段、容量、状态值、循环方向、取模或不可达哨兵；组合和排列必须用不同循环顺序表达；再用边界 无法凑出、amount 为零、硬币面额大于目标、哨兵溢出 反查初始化、更新顺序和退出条件。

\textbf{状态回放。} 手算路线：手算时用一个容量很小的例子比较正序和倒序。只要循环方向改变后同一物品能否重复使用发生变化，就说明方向不是细节。状态字段：物品阶段、容量、状态值、循环方向、取模或不可达哨兵；组合和排列必须用不同循环顺序表达。状态升级：把选择树压成阶段和容量；物品是否可重复使用决定一维表的遍历方向。

\textbf{证明和边界。} 证明从当前物品使用次数开始：0-1、完全、计数和最值的循环语义必须和题意一致。边界：无法凑出、amount 为零、硬币面额大于目标、哨兵溢出。变体：零钱兑换 II 问组合数，循环语义不同。

\noindent\textbf{LeetCode 377 组合总和 IV。} 模型：题目统计排列数，选择顺序不同算不同方案。推导重点：外层遍历容量，内层遍历数字，让每个容量的方案按最后一个数累加。

\textbf{二刷读题。} 读题时先找阶段、容量、物品使用次数和状态值类型。本题可以压成“题目统计排列数，选择顺序不同算不同方案”，循环方向由这些条件决定。先遮住答案，只保留题面，复述候选对象、合法性条件和输出目标。

\textbf{暴力回放。} 慢版本枚举每个物品选不选、选几次，并记录阶段和剩余容量。相同阶段和容量反复出现时，就可以把指数搜索压成表。二刷时先写慢版本或口述慢版本，用它确认不会漏答案。

\textbf{特例回放。} 拿 nums=[1,2,3]、target=4 手算，1+3 和 3+1 是不同排列，所以 dp[4] 要按最后一步来自 1、2、3 分别累加。规律是组合总和 IV 计数排列，容量在外层、数字在内层，循环顺序不能写成组合题。把这个特例继续拆开看：阶段是物品，容量是资源，状态值是可达、最值还是计数。一个物品能否在同一轮再次使用，直接决定容量循环方向；组合和排列也由循环顺序表达。复盘时把这个特例压成状态字段：物品阶段、容量、状态值、循环方向、取模或不可达哨兵；组合和排列必须用不同循环顺序表达；再用边界 target 为零、数字大于目标、方案数溢出、循环顺序 反查初始化、更新顺序和退出条件。

\textbf{状态回放。} 手算路线：手算时用一个容量很小的例子比较正序和倒序。只要循环方向改变后同一物品能否重复使用发生变化，就说明方向不是细节。状态字段：物品阶段、容量、状态值、循环方向、取模或不可达哨兵；组合和排列必须用不同循环顺序表达。状态升级：把选择树压成阶段和容量；物品是否可重复使用决定一维表的遍历方向。

\textbf{证明和边界。} 证明从当前物品使用次数开始：0-1、完全、计数和最值的循环语义必须和题意一致。边界：target 为零、数字大于目标、方案数溢出、循环顺序。变体：若顺序不同不算新方案，就变成完全背包组合数，循环顺序要交换。

\noindent\textbf{LeetCode 416 分割等和子集。} 模型：总和一半是背包容量，每个数只能使用一次。推导重点：0-1 背包可达性，一维容量必须倒序。

\textbf{二刷读题。} 读题时先找阶段、容量、物品使用次数和状态值类型。本题可以压成“总和一半是背包容量，每个数只能使用一次”，循环方向由这些条件决定。先遮住答案，只保留题面，复述候选对象、合法性条件和输出目标。

\textbf{暴力回放。} 慢版本枚举每个物品选不选、选几次，并记录阶段和剩余容量。相同阶段和容量反复出现时，就可以把指数搜索压成表。二刷时先写慢版本或口述慢版本，用它确认不会漏答案。

\textbf{特例回放。} 拿 [1,5,11,5] 手算，总和 22，目标容量 11；可以用 11 单独凑成，也可以 5+5+1。规律是分割等和转成 0-1 背包是否能凑到 sum/2，容量必须倒序避免重复使用同一元素。把这个特例继续拆开看：阶段是物品，容量是资源，状态值是可达、最值还是计数。一个物品能否在同一轮再次使用，直接决定容量循环方向；组合和排列也由循环顺序表达。复盘时把这个特例压成状态字段：物品阶段、容量、状态值、循环方向、取模或不可达哨兵；组合和排列必须用不同循环顺序表达；再用边界 总和为奇数、目标为零、数值较大、复杂度依赖 target 反查初始化、更新顺序和退出条件。

\textbf{状态回放。} 手算路线：手算时用一个容量很小的例子比较正序和倒序。只要循环方向改变后同一物品能否重复使用发生变化，就说明方向不是细节。状态字段：物品阶段、容量、状态值、循环方向、取模或不可达哨兵；组合和排列必须用不同循环顺序表达。状态升级：把选择树压成阶段和容量；物品是否可重复使用决定一维表的遍历方向。

\textbf{证明和边界。} 证明从当前物品使用次数开始：0-1、完全、计数和最值的循环语义必须和题意一致。边界：总和为奇数、目标为零、数值较大、复杂度依赖 target。变体：负数会破坏普通容量模型，需要重新建模。

\noindent\textbf{LeetCode 474 一和零。} 模型：每个字符串消耗零和一两种容量，每个字符串最多选一次。推导重点：二维 0-1 背包，两个容量都要倒序遍历。

\textbf{二刷读题。} 读题时先找阶段、容量、物品使用次数和状态值类型。本题可以压成“每个字符串消耗零和一两种容量，每个字符串最多选一次”，循环方向由这些条件决定。先遮住答案，只保留题面，复述候选对象、合法性条件和输出目标。

\textbf{暴力回放。} 慢版本枚举每个物品选不选、选几次，并记录阶段和剩余容量。相同阶段和容量反复出现时，就可以把指数搜索压成表。二刷时先写慢版本或口述慢版本，用它确认不会漏答案。

\textbf{特例回放。} 拿 strs=[10,0,1]、m=1、n=1 手算，字符串 10 消耗一个 0 和一个 1，单独就得到长度 1；0 和 1 也能组合成长度 2。规律是二维 0-1 背包，容量是 0 的数量和 1 的数量，两维都要倒序。把这个特例继续拆开看：阶段是物品，容量是资源，状态值是可达、最值还是计数。一个物品能否在同一轮再次使用，直接决定容量循环方向；组合和排列也由循环顺序表达。复盘时把这个特例压成状态字段：物品阶段、容量、状态值、循环方向、取模或不可达哨兵；组合和排列必须用不同循环顺序表达；再用边界 字符串全零或全一、容量为零、倒序方向、计数预处理 反查初始化、更新顺序和退出条件。

\textbf{状态回放。} 手算路线：手算时用一个容量很小的例子比较正序和倒序。只要循环方向改变后同一物品能否重复使用发生变化，就说明方向不是细节。状态字段：物品阶段、容量、状态值、循环方向、取模或不可达哨兵；组合和排列必须用不同循环顺序表达。状态升级：把选择树压成阶段和容量；物品是否可重复使用决定一维表的遍历方向。

\textbf{证明和边界。} 证明从当前物品使用次数开始：0-1、完全、计数和最值的循环语义必须和题意一致。边界：字符串全零或全一、容量为零、倒序方向、计数预处理。变体：这是背包从一维容量扩展到二维容量的代表。

\noindent\textbf{LeetCode 494 目标和。} 模型：正负号选择可代数转换为选若干数凑出特定正和。推导重点：若 S 加 target 为奇数或目标绝对值过大，直接无解。

\textbf{二刷读题。} 读题时先找阶段、容量、物品使用次数和状态值类型。本题可以压成“正负号选择可代数转换为选若干数凑出特定正和”，循环方向由这些条件决定。先遮住答案，只保留题面，复述候选对象、合法性条件和输出目标。

\textbf{暴力回放。} 慢版本枚举每个物品选不选、选几次，并记录阶段和剩余容量。相同阶段和容量反复出现时，就可以把指数搜索压成表。二刷时先写慢版本或口述慢版本，用它确认不会漏答案。

\textbf{特例回放。} 拿 [1,1,1,1,1]、target=3 手算，给其中一个 1 加负号，其余为正可以得到 3，共有 5 种位置选择。也可转成选正数子集和为 (sum+target)/2。规律是符号分配变成 0-1 计数背包，奇偶和范围先判定。把这个特例继续拆开看：阶段是物品，容量是资源，状态值是可达、最值还是计数。一个物品能否在同一轮再次使用，直接决定容量循环方向；组合和排列也由循环顺序表达。复盘时把这个特例压成状态字段：物品阶段、容量、状态值、循环方向、取模或不可达哨兵；组合和排列必须用不同循环顺序表达；再用边界 零元素会让方案数翻倍、目标为负、容量为零 反查初始化、更新顺序和退出条件。

\textbf{状态回放。} 手算路线：手算时用一个容量很小的例子比较正序和倒序。只要循环方向改变后同一物品能否重复使用发生变化，就说明方向不是细节。状态字段：物品阶段、容量、状态值、循环方向、取模或不可达哨兵；组合和排列必须用不同循环顺序表达。状态升级：把选择树压成阶段和容量；物品是否可重复使用决定一维表的遍历方向。

\textbf{证明和边界。} 证明从当前物品使用次数开始：0-1、完全、计数和最值的循环语义必须和题意一致。边界：零元素会让方案数翻倍、目标为负、容量为零。变体：也可用哈希 DP 保存所有可能和，但背包转换更高效。

\noindent\textbf{LeetCode 518 零钱兑换 II。} 模型：硬币无限使用，问组合数而不是最少数量。推导重点：外层遍历硬币、内层金额正序，保证组合按固定硬币顺序统计。

\textbf{二刷读题。} 读题时先找阶段、容量、物品使用次数和状态值类型。本题可以压成“硬币无限使用，问组合数而不是最少数量”，循环方向由这些条件决定。先遮住答案，只保留题面，复述候选对象、合法性条件和输出目标。

\textbf{暴力回放。} 慢版本枚举每个物品选不选、选几次，并记录阶段和剩余容量。相同阶段和容量反复出现时，就可以把指数搜索压成表。二刷时先写慢版本或口述慢版本，用它确认不会漏答案。

\textbf{特例回放。} 拿 amount=5、coins=[1,2,5] 手算，组合 5、2+2+1、2+1+1+1、全 1 一共 4 种；1+2 和 2+1 不能重复算。规律是组合计数时外层枚举硬币，内层容量正序。把这个特例继续拆开看：阶段是物品，容量是资源，状态值是可达、最值还是计数。一个物品能否在同一轮再次使用，直接决定容量循环方向；组合和排列也由循环顺序表达。复盘时把这个特例压成状态字段：物品阶段、容量、状态值、循环方向、取模或不可达哨兵；组合和排列必须用不同循环顺序表达；再用边界 amount 为零、无硬币、组合数较大、循环顺序反例 反查初始化、更新顺序和退出条件。

\textbf{状态回放。} 手算路线：手算时用一个容量很小的例子比较正序和倒序。只要循环方向改变后同一物品能否重复使用发生变化，就说明方向不是细节。状态字段：物品阶段、容量、状态值、循环方向、取模或不可达哨兵；组合和排列必须用不同循环顺序表达。状态升级：把选择树压成阶段和容量；物品是否可重复使用决定一维表的遍历方向。

\textbf{证明和边界。} 证明从当前物品使用次数开始：0-1、完全、计数和最值的循环语义必须和题意一致。边界：amount 为零、无硬币、组合数较大、循环顺序反例。变体：若外层金额内层硬币，会统计排列数。

\noindent\textbf{LeetCode 879 盈利计划。} 模型：每个项目消耗人数并贡献利润，利润达到下限后可截断。推导重点：二维 DP 用人数和已达到利润表示方案数，利润维度超过目标就压到目标。

\textbf{二刷读题。} 读题时先找阶段、容量、物品使用次数和状态值类型。本题可以压成“每个项目消耗人数并贡献利润，利润达到下限后可截断”，循环方向由这些条件决定。先遮住答案，只保留题面，复述候选对象、合法性条件和输出目标。

\textbf{暴力回放。} 慢版本枚举每个物品选不选、选几次，并记录阶段和剩余容量。相同阶段和容量反复出现时，就可以把指数搜索压成表。二刷时先写慢版本或口述慢版本，用它确认不会漏答案。

\textbf{特例回放。} 拿 n=5、minProfit=3、group=[2,2]、profit=[2,3] 手算，第二个项目单独就达标，两个项目一起也达标且人数不超。规律是状态包含已用人数和被截断到 minProfit 的利润，项目只能选一次。把这个特例继续拆开看：阶段是物品，容量是资源，状态值是可达、最值还是计数。一个物品能否在同一轮再次使用，直接决定容量循环方向；组合和排列也由循环顺序表达。复盘时把这个特例压成状态字段：物品阶段、容量、状态值、循环方向、取模或不可达哨兵；组合和排列必须用不同循环顺序表达；再用边界 minProfit 为零、人数不足、方案数取模、项目只能选一次 反查初始化、更新顺序和退出条件。

\textbf{状态回放。} 手算路线：手算时用一个容量很小的例子比较正序和倒序。只要循环方向改变后同一物品能否重复使用发生变化，就说明方向不是细节。状态字段：物品阶段、容量、状态值、循环方向、取模或不可达哨兵；组合和排列必须用不同循环顺序表达。状态升级：把选择树压成阶段和容量；物品是否可重复使用决定一维表的遍历方向。

\textbf{证明和边界。} 证明从当前物品使用次数开始：0-1、完全、计数和最值的循环语义必须和题意一致。边界：minProfit 为零、人数不足、方案数取模、项目只能选一次。变体：这是 0-1 背包的计数扩展，目标维度需要饱和处理。

\noindent\textbf{LeetCode 956 最高的广告牌。} 模型：两组钢筋高度相等时得到广告牌，高度差是关键状态。推导重点：DP 保存某个高度差下较矮一侧的最大高度，新增钢筋可放高侧、低侧或不用。

\textbf{二刷读题。} 读题时先找阶段、容量、物品使用次数和状态值类型。本题可以压成“两组钢筋高度相等时得到广告牌，高度差是关键状态”，循环方向由这些条件决定。先遮住答案，只保留题面，复述候选对象、合法性条件和输出目标。

\textbf{暴力回放。} 慢版本枚举每个物品选不选、选几次，并记录阶段和剩余容量。相同阶段和容量反复出现时，就可以把指数搜索压成表。二刷时先写慢版本或口述慢版本，用它确认不会漏答案。

\textbf{特例回放。} 拿 rods=[1,2,3,6] 手算，左塔用 6，右塔用 1+2+3，也能达到高度 6。规律是状态保存两边高度差对应的较矮塔最大高度，加入一根杆可以放左、放右或不用。把这个特例继续拆开看：阶段是物品，容量是资源，状态值是可达、最值还是计数。一个物品能否在同一轮再次使用，直接决定容量循环方向；组合和排列也由循环顺序表达。复盘时把这个特例压成状态字段：物品阶段、容量、状态值、循环方向、取模或不可达哨兵；组合和排列必须用不同循环顺序表达；再用边界 空数组、无法平衡、重复长度、差值更新顺序 反查初始化、更新顺序和退出条件。

\textbf{状态回放。} 手算路线：手算时用一个容量很小的例子比较正序和倒序。只要循环方向改变后同一物品能否重复使用发生变化，就说明方向不是细节。状态字段：物品阶段、容量、状态值、循环方向、取模或不可达哨兵；组合和排列必须用不同循环顺序表达。状态升级：把选择树压成阶段和容量；物品是否可重复使用决定一维表的遍历方向。

\textbf{证明和边界。} 证明从当前物品使用次数开始：0-1、完全、计数和最值的循环语义必须和题意一致。边界：空数组、无法平衡、重复长度、差值更新顺序。变体：这是背包从容量转成差值状态的代表。

\noindent\textbf{LeetCode 1049 最后一块石头的重量 II。} 模型：把石头分成两堆，最终差值等于两堆重量差。推导重点：0-1 背包找不超过总和一半的最大可达重量。

\textbf{二刷读题。} 读题时先找阶段、容量、物品使用次数和状态值类型。本题可以压成“把石头分成两堆，最终差值等于两堆重量差”，循环方向由这些条件决定。先遮住答案，只保留题面，复述候选对象、合法性条件和输出目标。

\textbf{暴力回放。} 慢版本枚举每个物品选不选、选几次，并记录阶段和剩余容量。相同阶段和容量反复出现时，就可以把指数搜索压成表。二刷时先写慢版本或口述慢版本，用它确认不会漏答案。

\textbf{特例回放。} 拿 [2,7,4,1,8,1] 手算，把石头分成两堆，最后差值越小结果越小；总和 23，尽量凑接近 11 的一堆。规律是 0-1 背包求不超过 sum/2 的最大可达重量。把这个特例继续拆开看：阶段是物品，容量是资源，状态值是可达、最值还是计数。一个物品能否在同一轮再次使用，直接决定容量循环方向；组合和排列也由循环顺序表达。复盘时把这个特例压成状态字段：物品阶段、容量、状态值、循环方向、取模或不可达哨兵；组合和排列必须用不同循环顺序表达；再用边界 单块石头、总和为奇偶、多个相同重量、容量方向 反查初始化、更新顺序和退出条件。

\textbf{状态回放。} 手算路线：手算时用一个容量很小的例子比较正序和倒序。只要循环方向改变后同一物品能否重复使用发生变化，就说明方向不是细节。状态字段：物品阶段、容量、状态值、循环方向、取模或不可达哨兵；组合和排列必须用不同循环顺序表达。状态升级：把选择树压成阶段和容量；物品是否可重复使用决定一维表的遍历方向。

\textbf{证明和边界。} 证明从当前物品使用次数开始：0-1、完全、计数和最值的循环语义必须和题意一致。边界：单块石头、总和为奇偶、多个相同重量、容量方向。变体：它和分割等和子集同源，只是目标从是否可达变成最小差。

\noindent\textbf{LeetCode 1155 掷骰子等于目标和的方法数。} 模型：多个骰子每个只能取一到 faces 的点数，问目标和方案数。推导重点：按骰子数量分阶段，容量为当前点数和，枚举当前骰子的点数转移。

\textbf{二刷读题。} 读题时先找阶段、容量、物品使用次数和状态值类型。本题可以压成“多个骰子每个只能取一到 faces 的点数，问目标和方案数”，循环方向由这些条件决定。先遮住答案，只保留题面，复述候选对象、合法性条件和输出目标。

\textbf{暴力回放。} 慢版本枚举每个物品选不选、选几次，并记录阶段和剩余容量。相同阶段和容量反复出现时，就可以把指数搜索压成表。二刷时先写慢版本或口述慢版本，用它确认不会漏答案。

\textbf{特例回放。} 拿 2 个 6 面骰、target=7 手算，第一颗为 1 时第二颗必须 6，第一颗为 2 时第二颗必须 5。规律是阶段是骰子个数，容量是当前点数和，每个骰子枚举 1..k 的面值。把这个特例继续拆开看：阶段是物品，容量是资源，状态值是可达、最值还是计数。一个物品能否在同一轮再次使用，直接决定容量循环方向；组合和排列也由循环顺序表达。复盘时把这个特例压成状态字段：物品阶段、容量、状态值、循环方向、取模或不可达哨兵；组合和排列必须用不同循环顺序表达；再用边界 target 太小或太大、取模、骰子数为一、滚动数组方向 反查初始化、更新顺序和退出条件。

\textbf{状态回放。} 手算路线：手算时用一个容量很小的例子比较正序和倒序。只要循环方向改变后同一物品能否重复使用发生变化，就说明方向不是细节。状态字段：物品阶段、容量、状态值、循环方向、取模或不可达哨兵；组合和排列必须用不同循环顺序表达。状态升级：把选择树压成阶段和容量；物品是否可重复使用决定一维表的遍历方向。

\textbf{证明和边界。} 证明从当前物品使用次数开始：0-1、完全、计数和最值的循环语义必须和题意一致。边界：target 太小或太大、取模、骰子数为一、滚动数组方向。变体：这是有界选择计数，和完全背包不能混。

\noindent\textbf{LeetCode 1449 数位成本和为目标值的最大数字。} 模型：每个数字可重复选择，容量是成本，目标是组成字典序和长度最大的数字。推导重点：完全背包先最大化位数，再按数字从大到小贪心还原答案。

\textbf{二刷读题。} 读题时先找阶段、容量、物品使用次数和状态值类型。本题可以压成“每个数字可重复选择，容量是成本，目标是组成字典序和长度最大的数字”，循环方向由这些条件决定。先遮住答案，只保留题面，复述候选对象、合法性条件和输出目标。

\textbf{暴力回放。} 慢版本枚举每个物品选不选、选几次，并记录阶段和剩余容量。相同阶段和容量反复出现时，就可以把指数搜索压成表。二刷时先写慢版本或口述慢版本，用它确认不会漏答案。

\textbf{特例回放。} 拿 cost 中数字 7 和 2 成本较低、target=9 手算，先最大化能组成的位数，再从 9 到 1 尝试还原字典序最大的数字。规律是完全背包求最大位数，重建时优先放大数字且保证剩余目标可达。把这个特例继续拆开看：阶段是物品，容量是资源，状态值是可达、最值还是计数。一个物品能否在同一轮再次使用，直接决定容量循环方向；组合和排列也由循环顺序表达。复盘时把这个特例压成状态字段：物品阶段、容量、状态值、循环方向、取模或不可达哨兵；组合和排列必须用不同循环顺序表达；再用边界 无法组成目标、多个数字同成本、target 为零、字符串比较 反查初始化、更新顺序和退出条件。

\textbf{状态回放。} 手算路线：手算时用一个容量很小的例子比较正序和倒序。只要循环方向改变后同一物品能否重复使用发生变化，就说明方向不是细节。状态字段：物品阶段、容量、状态值、循环方向、取模或不可达哨兵；组合和排列必须用不同循环顺序表达。状态升级：把选择树压成阶段和容量；物品是否可重复使用决定一维表的遍历方向。

\textbf{证明和边界。} 证明从当前物品使用次数开始：0-1、完全、计数和最值的循环语义必须和题意一致。边界：无法组成目标、多个数字同成本、target 为零、字符串比较。变体：这题说明背包状态值不一定是数值收益，也可以是长度和字典序。

\subsection{自测标准}

二刷时不要只确认自己见过这道题。请遮住答案，先讲题面模型，再讲暴力基线和瓶颈，然后说出优化结构保存了什么、不变量如何排除错误候选、边界实验如何打破错误实现。

\section{贪心与交换证明：二刷复盘卡}
这一大节只围绕一个题型复盘，不把题号直接铺成目录。阅读时先看复盘目标，再读核心题卡，最后用自测标准检查自己是否能独立重讲。

\subsection{复盘目标}

追问通常会要求证明。请准备交换论证或领先性论证，并给出一个看似合理但错误的贪心规则反例。

\subsection{核心题卡}

\noindent\textbf{LeetCode 45 跳跃游戏 II。} 模型：每一次跳跃可以看成 BFS 的一层，当前层内所有位置共享同一次跳数。推导重点：扫描当前层时维护下一层最远边界，到达当前层末尾才增加跳数。

\textbf{二刷读题。} 读题时先找局部选择消耗什么资源，以及选择之后给未来留下什么边界。本题可以压成“每一次跳跃可以看成 BFS 的一层，当前层内所有位置共享同一次跳数”，扫描当前层时维护下一层最远边界，到达当前层末尾才增加跳数必须能被证明安全。先遮住答案，只保留题面，复述候选对象、合法性条件和输出目标。

\textbf{暴力回放。} 慢版本枚举选择顺序或用 DP 保留多个可能方案，先确认最优值存在。随后才检查局部选择是否能通过交换或领先性固定下来。二刷时先写慢版本或口述慢版本，用它确认不会漏答案。

\textbf{特例回放。} 拿 [2,3,1,1,4] 手算，起点这一层能到下标 1 和 2；扫描这一层时，下一层最远可到 4，走完当前层末尾才把步数加一。规律是把可达区间看成 BFS 层，当前层结束时切到下一层；不能在每个位置都加跳数。把这个特例继续拆开看：先构造一个非贪心选择，再比较两者留下的资源、右边界或剩余时间。只有能把非贪心第一步交换成贪心第一步且不变差，局部选择才是规律。复盘时把这个特例压成状态字段：排序后的候选、当前已选方案、剩余资源或最远边界、答案计数；每次局部选择后要能说明当前状态不落后；再用边界 数组长度一、刚好到达、存在多个可选跳点、不要在每个位置都加跳数 反查初始化、更新顺序和退出条件。

\textbf{状态回放。} 手算路线：手算时同时写一个非贪心选择，与贪心选择比较剩余资源。若无法说明贪心后剩余资源不差，就不能直接下结论。状态字段：排序后的候选、当前已选方案、剩余资源或最远边界、答案计数；每次局部选择后要能说明当前状态不落后。状态升级：把指数级选择压成一个可证明安全的局部选择；状态通常是当前最远边界、最早结束时间、最小代价或剩余资源。

\textbf{证明和边界。} 证明从交换或领先性开始：任意最优解都能替换成当前贪心选择而不变差，或贪心状态始终不落后。边界：数组长度一、刚好到达、存在多个可选跳点、不要在每个位置都加跳数。变体：若问能否到达，只需维护全局最远可达位置。

\noindent\textbf{LeetCode 55 跳跃游戏。} 模型：扫描过程中只关心当前能到达的最远位置。推导重点：如果下标超过最远可达，任何之前路径都不能到达它；否则用它扩展边界。

\textbf{二刷读题。} 读题时先找局部选择消耗什么资源，以及选择之后给未来留下什么边界。本题可以压成“扫描过程中只关心当前能到达的最远位置”，如果下标超过最远可达，任何之前路径都不能到达它；否则用它扩展边界必须能被证明安全。先遮住答案，只保留题面，复述候选对象、合法性条件和输出目标。

\textbf{暴力回放。} 慢版本枚举选择顺序或用 DP 保留多个可能方案，先确认最优值存在。随后才检查局部选择是否能通过交换或领先性固定下来。二刷时先写慢版本或口述慢版本，用它确认不会漏答案。

\textbf{特例回放。} 拿 [2,3,1,1,4] 手算，下标 0 能覆盖到 2；走到下标 1 时最远可达更新到 4，已经覆盖终点。再拿 [0,1]，扫描到下标 1 时已经超过最远可达，失败。规律是只维护当前能到达的最远边界，任何超过边界的位置都不可能由前面路径到达。把这个特例继续拆开看：先构造一个非贪心选择，再比较两者留下的资源、右边界或剩余时间。只有能把非贪心第一步交换成贪心第一步且不变差，局部选择才是规律。复盘时把这个特例压成状态字段：排序后的候选、当前已选方案、剩余资源或最远边界、答案计数；每次局部选择后要能说明当前状态不落后；再用边界 第一个元素为零、数组长度一、恰好到达最后、远超最后 反查初始化、更新顺序和退出条件。

\textbf{状态回放。} 手算路线：手算时同时写一个非贪心选择，与贪心选择比较剩余资源。若无法说明贪心后剩余资源不差，就不能直接下结论。状态字段：排序后的候选、当前已选方案、剩余资源或最远边界、答案计数；每次局部选择后要能说明当前状态不落后。状态升级：把指数级选择压成一个可证明安全的局部选择；状态通常是当前最远边界、最早结束时间、最小代价或剩余资源。

\textbf{证明和边界。} 证明从交换或领先性开始：任意最优解都能替换成当前贪心选择而不变差，或贪心状态始终不落后。边界：第一个元素为零、数组长度一、恰好到达最后、远超最后。变体：求最少跳数时把可达区间看成 BFS 层。

\noindent\textbf{LeetCode 122 买卖股票 II。} 模型：允许多次交易且不能同时持有，所有正收益差都可以收集。推导重点：把长上涨段拆成每天买卖收益相同，因此累加正差值最优。

\textbf{二刷读题。} 读题时先找局部选择消耗什么资源，以及选择之后给未来留下什么边界。本题可以压成“允许多次交易且不能同时持有，所有正收益差都可以收集”，把长上涨段拆成每天买卖收益相同，因此累加正差值最优必须能被证明安全。先遮住答案，只保留题面，复述候选对象、合法性条件和输出目标。

\textbf{暴力回放。} 慢版本枚举选择顺序或用 DP 保留多个可能方案，先确认最优值存在。随后才检查局部选择是否能通过交换或领先性固定下来。二刷时先写慢版本或口述慢版本，用它确认不会漏答案。

\textbf{特例回放。} 拿 [7,1,5,3,6,4] 手算，1 到 5 的利润 4，加上 3 到 6 的利润 3，总利润 7。把 1 到 6 的上涨段拆成每天正差值，收益相同。规律是无限次交易且无手续费时，所有正差值都可收集；下降段和平台不贡献收益。把这个特例继续拆开看：先构造一个非贪心选择，再比较两者留下的资源、右边界或剩余时间。只有能把非贪心第一步交换成贪心第一步且不变差，局部选择才是规律。复盘时把这个特例压成状态字段：排序后的候选、当前已选方案、剩余资源或最远边界、答案计数；每次局部选择后要能说明当前状态不落后；再用边界 下降段、平台价格、只有一天、手续费不存在 反查初始化、更新顺序和退出条件。

\textbf{状态回放。} 手算路线：手算时同时写一个非贪心选择，与贪心选择比较剩余资源。若无法说明贪心后剩余资源不差，就不能直接下结论。状态字段：排序后的候选、当前已选方案、剩余资源或最远边界、答案计数；每次局部选择后要能说明当前状态不落后。状态升级：把指数级选择压成一个可证明安全的局部选择；状态通常是当前最远边界、最早结束时间、最小代价或剩余资源。

\textbf{证明和边界。} 证明从交换或领先性开始：任意最优解都能替换成当前贪心选择而不变差，或贪心状态始终不落后。边界：下降段、平台价格、只有一天、手续费不存在。变体：有手续费时不能简单累加正差，需要状态机或调整贪心阈值。

\noindent\textbf{LeetCode 134 加油站。} 模型：绕行一圈是否可行取决于总油量差和局部剩余油量。推导重点：总和为负必不可行；扫描时若当前剩余为负，起点只能放到下一站。

\textbf{二刷读题。} 读题时先找局部选择消耗什么资源，以及选择之后给未来留下什么边界。本题可以压成“绕行一圈是否可行取决于总油量差和局部剩余油量”，总和为负必不可行；扫描时若当前剩余为负，起点只能放到下一站必须能被证明安全。先遮住答案，只保留题面，复述候选对象、合法性条件和输出目标。

\textbf{暴力回放。} 慢版本枚举选择顺序或用 DP 保留多个可能方案，先确认最优值存在。随后才检查局部选择是否能通过交换或领先性固定下来。二刷时先写慢版本或口述慢版本，用它确认不会漏答案。

\textbf{特例回放。} 拿 gas=[1,2,3,4,5]、cost=[3,4,5,1,2] 手算，从 0 出发中途油量为负，0 到失败点之间任何站都不能作为起点；从 3 出发可以绕完。规律是总油量差为负则无解，局部剩余为负时起点跳到下一站。把这个特例继续拆开看：先构造一个非贪心选择，再比较两者留下的资源、右边界或剩余时间。只有能把非贪心第一步交换成贪心第一步且不变差，局部选择才是规律。复盘时把这个特例压成状态字段：排序后的候选、当前已选方案、剩余资源或最远边界、答案计数；每次局部选择后要能说明当前状态不落后；再用边界 总和刚好为零、多个可行起点、某站油量为零、消耗大于补给 反查初始化、更新顺序和退出条件。

\textbf{状态回放。} 手算路线：手算时同时写一个非贪心选择，与贪心选择比较剩余资源。若无法说明贪心后剩余资源不差，就不能直接下结论。状态字段：排序后的候选、当前已选方案、剩余资源或最远边界、答案计数；每次局部选择后要能说明当前状态不落后。状态升级：把指数级选择压成一个可证明安全的局部选择；状态通常是当前最远边界、最早结束时间、最小代价或剩余资源。

\textbf{证明和边界。} 证明从交换或领先性开始：任意最优解都能替换成当前贪心选择而不变差，或贪心状态始终不落后。边界：总和刚好为零、多个可行起点、某站油量为零、消耗大于补给。变体：它的领先性证明是前一段任意站点都不能作为起点。

\noindent\textbf{LeetCode 135 分发糖果。} 模型：相邻评分约束需要同时满足左邻和右邻两个方向。推导重点：左到右满足递增关系，右到左满足反向递增关系，最后取两侧要求最大值。

\textbf{二刷读题。} 读题时先找局部选择消耗什么资源，以及选择之后给未来留下什么边界。本题可以压成“相邻评分约束需要同时满足左邻和右邻两个方向”，左到右满足递增关系，右到左满足反向递增关系，最后取两侧要求最大值必须能被证明安全。先遮住答案，只保留题面，复述候选对象、合法性条件和输出目标。

\textbf{暴力回放。} 慢版本枚举选择顺序或用 DP 保留多个可能方案，先确认最优值存在。随后才检查局部选择是否能通过交换或领先性固定下来。二刷时先写慢版本或口述慢版本，用它确认不会漏答案。

\textbf{特例回放。} 拿 ratings=[1,0,2] 手算，中间孩子评分最低，左右都要比它多糖，结果 [2,1,2]。只从左到右会漏掉右侧约束。规律是左到右满足左邻递增，右到左满足右邻递增，最终取两侧要求最大值。把这个特例继续拆开看：先构造一个非贪心选择，再比较两者留下的资源、右边界或剩余时间。只有能把非贪心第一步交换成贪心第一步且不变差，局部选择才是规律。复盘时把这个特例压成状态字段：排序后的候选、当前已选方案、剩余资源或最远边界、答案计数；每次局部选择后要能说明当前状态不落后；再用边界 评分相等、严格递增、严格递减、山峰和谷底 反查初始化、更新顺序和退出条件。

\textbf{状态回放。} 手算路线：手算时同时写一个非贪心选择，与贪心选择比较剩余资源。若无法说明贪心后剩余资源不差，就不能直接下结论。状态字段：排序后的候选、当前已选方案、剩余资源或最远边界、答案计数；每次局部选择后要能说明当前状态不落后。状态升级：把指数级选择压成一个可证明安全的局部选择；状态通常是当前最远边界、最早结束时间、最小代价或剩余资源。

\textbf{证明和边界。} 证明从交换或领先性开始：任意最优解都能替换成当前贪心选择而不变差，或贪心状态始终不落后。边界：评分相等、严格递增、严格递减、山峰和谷底。变体：这是局部约束两遍扫描，不能只用一个方向贪心。

\noindent\textbf{LeetCode 169 多数元素。} 模型：超过一半的元素可以和其他元素两两抵消后仍然剩下。推导重点：Boyer-Moore 投票维护候选和计数。

\textbf{二刷读题。} 读题时先找局部选择消耗什么资源，以及选择之后给未来留下什么边界。本题可以压成“超过一半的元素可以和其他元素两两抵消后仍然剩下”，Boyer-Moore 投票维护候选和计数必须能被证明安全。先遮住答案，只保留题面，复述候选对象、合法性条件和输出目标。

\textbf{暴力回放。} 慢版本枚举选择顺序或用 DP 保留多个可能方案，先确认最优值存在。随后才检查局部选择是否能通过交换或领先性固定下来。二刷时先写慢版本或口述慢版本，用它确认不会漏答案。

\textbf{特例回放。} 拿 [2,2,1,1,1,2,2] 手算，候选和不同元素互相抵消，最后剩下的候选是 2。规律是多数元素超过一半，和其他元素两两抵消后仍会剩下；题目不保证存在时需要二次计数验证。把这个特例继续拆开看：先构造一个非贪心选择，再比较两者留下的资源、右边界或剩余时间。只有能把非贪心第一步交换成贪心第一步且不变差，局部选择才是规律。复盘时把这个特例压成状态字段：排序后的候选、当前已选方案、剩余资源或最远边界、答案计数；每次局部选择后要能说明当前状态不落后；再用边界 题目是否保证存在多数、计数归零、全相同 反查初始化、更新顺序和退出条件。

\textbf{状态回放。} 手算路线：手算时同时写一个非贪心选择，与贪心选择比较剩余资源。若无法说明贪心后剩余资源不差，就不能直接下结论。状态字段：排序后的候选、当前已选方案、剩余资源或最远边界、答案计数；每次局部选择后要能说明当前状态不落后。状态升级：把指数级选择压成一个可证明安全的局部选择；状态通常是当前最远边界、最早结束时间、最小代价或剩余资源。

\textbf{证明和边界。} 证明从交换或领先性开始：任意最优解都能替换成当前贪心选择而不变差，或贪心状态始终不落后。边界：题目是否保证存在多数、计数归零、全相同。变体：若找超过三分之一的元素，需要维护两个候选。

\noindent\textbf{LeetCode 316 去除重复字母。} 模型：每个字符必须保留一次，同时结果字典序尽量小。推导重点：单调栈维护当前结果，若栈顶更大且后面还会出现，就可以弹出再等后面放回。

\textbf{二刷读题。} 读题时先找局部选择消耗什么资源，以及选择之后给未来留下什么边界。本题可以压成“每个字符必须保留一次，同时结果字典序尽量小”，单调栈维护当前结果，若栈顶更大且后面还会出现，就可以弹出再等后面放回必须能被证明安全。先遮住答案，只保留题面，复述候选对象、合法性条件和输出目标。

\textbf{暴力回放。} 慢版本枚举选择顺序或用 DP 保留多个可能方案，先确认最优值存在。随后才检查局部选择是否能通过交换或领先性固定下来。二刷时先写慢版本或口述慢版本，用它确认不会漏答案。

\textbf{特例回放。} 拿 cbacdcbc 手算，遇到 b 时，前面的 c 后面还会出现且 c>b，所以可以弹出 c；但遇到只剩最后一次的字符就不能弹。规律是单调栈维护结果字典序，弹出条件必须同时满足栈顶更大且后面还会出现。把这个特例继续拆开看：先构造一个非贪心选择，再比较两者留下的资源、右边界或剩余时间。只有能把非贪心第一步交换成贪心第一步且不变差，局部选择才是规律。复盘时把这个特例压成状态字段：排序后的候选、当前已选方案、剩余资源或最远边界、答案计数；每次局部选择后要能说明当前状态不落后；再用边界 字符只出现一次、重复字符、弹出后 visited 恢复、字典序比较 反查初始化、更新顺序和退出条件。

\textbf{状态回放。} 手算路线：手算时同时写一个非贪心选择，与贪心选择比较剩余资源。若无法说明贪心后剩余资源不差，就不能直接下结论。状态字段：排序后的候选、当前已选方案、剩余资源或最远边界、答案计数；每次局部选择后要能说明当前状态不落后。状态升级：把指数级选择压成一个可证明安全的局部选择；状态通常是当前最远边界、最早结束时间、最小代价或剩余资源。

\textbf{证明和边界。} 证明从交换或领先性开始：任意最优解都能替换成当前贪心选择而不变差，或贪心状态始终不落后。边界：字符只出现一次、重复字符、弹出后 visited 恢复、字典序比较。变体：402 移掉 K 位数字也是单调栈删除，但删除预算和必须保留集合不同。

\noindent\textbf{LeetCode 406 根据身高重建队列。} 模型：高个子的位置先确定后，不会被矮个子影响其前方高个数量。推导重点：按身高降序、k 升序排序，再按 k 插入当前结果位置。

\textbf{二刷读题。} 读题时先找局部选择消耗什么资源，以及选择之后给未来留下什么边界。本题可以压成“高个子的位置先确定后，不会被矮个子影响其前方高个数量”，按身高降序、k 升序排序，再按 k 插入当前结果位置必须能被证明安全。先遮住答案，只保留题面，复述候选对象、合法性条件和输出目标。

\textbf{暴力回放。} 慢版本枚举选择顺序或用 DP 保留多个可能方案，先确认最优值存在。随后才检查局部选择是否能通过交换或领先性固定下来。二刷时先写慢版本或口述慢版本，用它确认不会漏答案。

\textbf{特例回放。} 拿 [7,0]、[7,1]、[5,0] 手算，先放高个子，矮个子插入不会改变高个子前面高个数量；[5,0] 插到最前也不影响 [7,0] 的统计。规律是按身高降序、k 升序排序，再按 k 插入。把这个特例继续拆开看：先构造一个非贪心选择，再比较两者留下的资源、右边界或剩余时间。只有能把非贪心第一步交换成贪心第一步且不变差，局部选择才是规律。复盘时把这个特例压成状态字段：排序后的候选、当前已选方案、剩余资源或最远边界、答案计数；每次局部选择后要能说明当前状态不落后；再用边界 相同身高、k 为零、插入到末尾、结果稳定性 反查初始化、更新顺序和退出条件。

\textbf{状态回放。} 手算路线：手算时同时写一个非贪心选择，与贪心选择比较剩余资源。若无法说明贪心后剩余资源不差，就不能直接下结论。状态字段：排序后的候选、当前已选方案、剩余资源或最远边界、答案计数；每次局部选择后要能说明当前状态不落后。状态升级：把指数级选择压成一个可证明安全的局部选择；状态通常是当前最远边界、最早结束时间、最小代价或剩余资源。

\textbf{证明和边界。} 证明从交换或领先性开始：任意最优解都能替换成当前贪心选择而不变差，或贪心状态始终不落后。边界：相同身高、k 为零、插入到末尾、结果稳定性。变体：这是排序后插入的构造型贪心，证明依赖高个子先固定。

\noindent\textbf{LeetCode 435 无重叠区间。} 模型：保留最多不重叠区间等价于删除最少区间。推导重点：按终点升序选择，结束越早给后续留下空间越多。

\textbf{二刷读题。} 读题时先找局部选择消耗什么资源，以及选择之后给未来留下什么边界。本题可以压成“保留最多不重叠区间等价于删除最少区间”，按终点升序选择，结束越早给后续留下空间越多必须能被证明安全。先遮住答案，只保留题面，复述候选对象、合法性条件和输出目标。

\textbf{暴力回放。} 慢版本枚举选择顺序或用 DP 保留多个可能方案，先确认最优值存在。随后才检查局部选择是否能通过交换或领先性固定下来。二刷时先写慢版本或口述慢版本，用它确认不会漏答案。

\textbf{特例回放。} 拿 [[1,2],[2,3],[3,4],[1,3]] 手算，选 [1,2] 后还能接 [2,3] 和 [3,4]；若先选 [1,3]，会挡住更多区间。规律是按右端点升序选择，结束越早给未来留下的空间越多，删除数等于总数减保留数。把这个特例继续拆开看：先构造一个非贪心选择，再比较两者留下的资源、右边界或剩余时间。只有能把非贪心第一步交换成贪心第一步且不变差，局部选择才是规律。复盘时把这个特例压成状态字段：排序后的候选、当前已选方案、剩余资源或最远边界、答案计数；每次局部选择后要能说明当前状态不落后；再用边界 端点相接是否重叠、完全覆盖、相同终点、空列表 反查初始化、更新顺序和退出条件。

\textbf{状态回放。} 手算路线：手算时同时写一个非贪心选择，与贪心选择比较剩余资源。若无法说明贪心后剩余资源不差，就不能直接下结论。状态字段：排序后的候选、当前已选方案、剩余资源或最远边界、答案计数；每次局部选择后要能说明当前状态不落后。状态升级：把指数级选择压成一个可证明安全的局部选择；状态通常是当前最远边界、最早结束时间、最小代价或剩余资源。

\textbf{证明和边界。} 证明从交换或领先性开始：任意最优解都能替换成当前贪心选择而不变差，或贪心状态始终不落后。边界：端点相接是否重叠、完全覆盖、相同终点、空列表。变体：用交换论证说明最早结束的选择不劣于其他第一选择。

\noindent\textbf{LeetCode 452 用最少数量的箭引爆气球。} 模型：一支箭的位置可以覆盖所有包含该位置的区间。推导重点：按右端点排序，当前箭放在最早结束气球的右端。

\textbf{二刷读题。} 读题时先找局部选择消耗什么资源，以及选择之后给未来留下什么边界。本题可以压成“一支箭的位置可以覆盖所有包含该位置的区间”，按右端点排序，当前箭放在最早结束气球的右端必须能被证明安全。先遮住答案，只保留题面，复述候选对象、合法性条件和输出目标。

\textbf{暴力回放。} 慢版本枚举选择顺序或用 DP 保留多个可能方案，先确认最优值存在。随后才检查局部选择是否能通过交换或领先性固定下来。二刷时先写慢版本或口述慢版本，用它确认不会漏答案。

\textbf{特例回放。} 拿 [[10,16],[2,8],[1,6],[7,12]] 手算，按右端排序后第一箭射在 6，可以覆盖 [1,6] 和 [2,8]；下一箭射在 12。规律是每次把箭放在最早结束气球的右端，能最大化后续空间。把这个特例继续拆开看：先构造一个非贪心选择，再比较两者留下的资源、右边界或剩余时间。只有能把非贪心第一步交换成贪心第一步且不变差，局部选择才是规律。复盘时把这个特例压成状态字段：排序后的候选、当前已选方案、剩余资源或最远边界、答案计数；每次局部选择后要能说明当前状态不落后；再用边界 端点相同、负坐标、区间包含、闭区间边界 反查初始化、更新顺序和退出条件。

\textbf{状态回放。} 手算路线：手算时同时写一个非贪心选择，与贪心选择比较剩余资源。若无法说明贪心后剩余资源不差，就不能直接下结论。状态字段：排序后的候选、当前已选方案、剩余资源或最远边界、答案计数；每次局部选择后要能说明当前状态不落后。状态升级：把指数级选择压成一个可证明安全的局部选择；状态通常是当前最远边界、最早结束时间、最小代价或剩余资源。

\textbf{证明和边界。} 证明从交换或领先性开始：任意最优解都能替换成当前贪心选择而不变差，或贪心状态始终不落后。边界：端点相同、负坐标、区间包含、闭区间边界。变体：它和无重叠区间是同一类最早结束贪心。

\noindent\textbf{LeetCode 630 课程表 III。} 模型：每门课有持续时间和截止日，已选课程总时长不能超过当前截止日。推导重点：按截止日排序，先选当前课；若超时，就用大顶堆丢掉已选中耗时最长的课。

\textbf{二刷读题。} 读题时先找局部选择消耗什么资源，以及选择之后给未来留下什么边界。本题可以压成“每门课有持续时间和截止日，已选课程总时长不能超过当前截止日”，按截止日排序，先选当前课；若超时，就用大顶堆丢掉已选中耗时最长的课必须能被证明安全。先遮住答案，只保留题面，复述候选对象、合法性条件和输出目标。

\textbf{暴力回放。} 慢版本枚举选择顺序或用 DP 保留多个可能方案，先确认最优值存在。随后才检查局部选择是否能通过交换或领先性固定下来。二刷时先写慢版本或口述慢版本，用它确认不会漏答案。

\textbf{特例回放。} 拿课程 [100,200]、[200,1300]、[1000,1250] 手算，按截止日排序后若总时长超出当前截止日，就丢掉已选课程里耗时最长的。规律是大顶堆保存已选课程时长，替换最长课能给未来留下最多时间。把这个特例继续拆开看：先构造一个非贪心选择，再比较两者留下的资源、右边界或剩余时间。只有能把非贪心第一步交换成贪心第一步且不变差，局部选择才是规律。复盘时把这个特例压成状态字段：排序后的候选、当前已选方案、剩余资源或最远边界、答案计数；每次局部选择后要能说明当前状态不落后；再用边界 课程无法单独完成、截止日相同、替换已选课程、空列表 反查初始化、更新顺序和退出条件。

\textbf{状态回放。} 手算路线：手算时同时写一个非贪心选择，与贪心选择比较剩余资源。若无法说明贪心后剩余资源不差，就不能直接下结论。状态字段：排序后的候选、当前已选方案、剩余资源或最远边界、答案计数；每次局部选择后要能说明当前状态不落后。状态升级：把指数级选择压成一个可证明安全的局部选择；状态通常是当前最远边界、最早结束时间、最小代价或剩余资源。

\textbf{证明和边界。} 证明从交换或领先性开始：任意最优解都能替换成当前贪心选择而不变差，或贪心状态始终不落后。边界：课程无法单独完成、截止日相同、替换已选课程、空列表。变体：这是反悔贪心：接受后在约束被破坏时删最差候选。

\noindent\textbf{LeetCode 646 最长数对链。} 模型：每个数对像一个区间，选择链要求前一个右端小于后一个左端。推导重点：按右端升序选择，当前右端越小越容易接后续数对。

\textbf{二刷读题。} 读题时先找局部选择消耗什么资源，以及选择之后给未来留下什么边界。本题可以压成“每个数对像一个区间，选择链要求前一个右端小于后一个左端”，按右端升序选择，当前右端越小越容易接后续数对必须能被证明安全。先遮住答案，只保留题面，复述候选对象、合法性条件和输出目标。

\textbf{暴力回放。} 慢版本枚举选择顺序或用 DP 保留多个可能方案，先确认最优值存在。随后才检查局部选择是否能通过交换或领先性固定下来。二刷时先写慢版本或口述慢版本，用它确认不会漏答案。

\textbf{特例回放。} 拿 [[1,2],[2,3],[3,4]] 手算，[1,2] 不能接 [2,3]，因为要求前一个右端小于后一个左端；可以接 [3,4]。规律是按右端升序贪心选择，右端越小越容易接后续数对。把这个特例继续拆开看：先构造一个非贪心选择，再比较两者留下的资源、右边界或剩余时间。只有能把非贪心第一步交换成贪心第一步且不变差，局部选择才是规律。复盘时把这个特例压成状态字段：排序后的候选、当前已选方案、剩余资源或最远边界、答案计数；每次局部选择后要能说明当前状态不落后；再用边界 端点相等不允许连接、完全覆盖、负数端点、空数组 反查初始化、更新顺序和退出条件。

\textbf{状态回放。} 手算路线：手算时同时写一个非贪心选择，与贪心选择比较剩余资源。若无法说明贪心后剩余资源不差，就不能直接下结论。状态字段：排序后的候选、当前已选方案、剩余资源或最远边界、答案计数；每次局部选择后要能说明当前状态不落后。状态升级：把指数级选择压成一个可证明安全的局部选择；状态通常是当前最远边界、最早结束时间、最小代价或剩余资源。

\textbf{证明和边界。} 证明从交换或领先性开始：任意最优解都能替换成当前贪心选择而不变差，或贪心状态始终不落后。边界：端点相等不允许连接、完全覆盖、负数端点、空数组。变体：也可以 DP，但贪心用最早结束证明更简洁。

\noindent\textbf{LeetCode 678 有效的括号字符串。} 模型：星号可以作为左括号、右括号或空，使未匹配左括号数量变成一个可能区间。推导重点：贪心维护可能未匹配左括号数的下界和上界，遇到星号扩大区间并把下界截到零。

\textbf{二刷读题。} 读题时先找局部选择消耗什么资源，以及选择之后给未来留下什么边界。本题可以压成“星号可以作为左括号、右括号或空，使未匹配左括号数量变成一个可能区间”，贪心维护可能未匹配左括号数的下界和上界，遇到星号扩大区间并把下界截到零必须能被证明安全。先遮住答案，只保留题面，复述候选对象、合法性条件和输出目标。

\textbf{暴力回放。} 慢版本枚举选择顺序或用 DP 保留多个可能方案，先确认最优值存在。随后才检查局部选择是否能通过交换或领先性固定下来。二刷时先写慢版本或口述慢版本，用它确认不会漏答案。

\textbf{特例回放。} 拿 (*) 手算，星号可以当空或左括号；拿 (*))，星号必须当左括号才能匹配后面的右括号。与其枚举星号三种选择，不如维护可能未匹配左括号数的区间 [low, high]。规律是左括号让区间加一，右括号让区间减一，星号让下界减一上界加一，且下界不能低于零。把这个特例继续拆开看：先构造一个非贪心选择，再比较两者留下的资源、右边界或剩余时间。只有能把非贪心第一步交换成贪心第一步且不变差，局部选择才是规律。复盘时把这个特例压成状态字段：排序后的候选、当前已选方案、剩余资源或最远边界、答案计数；每次局部选择后要能说明当前状态不落后；再用边界 上界变负、最终下界不为零、连续星号、前缀右括号过多 反查初始化、更新顺序和退出条件。

\textbf{状态回放。} 手算路线：手算时同时写一个非贪心选择，与贪心选择比较剩余资源。若无法说明贪心后剩余资源不差，就不能直接下结论。状态字段：排序后的候选、当前已选方案、剩余资源或最远边界、答案计数；每次局部选择后要能说明当前状态不落后。状态升级：把指数级选择压成一个可证明安全的局部选择；状态通常是当前最远边界、最早结束时间、最小代价或剩余资源。

\textbf{证明和边界。} 证明从交换或领先性开始：任意最优解都能替换成当前贪心选择而不变差，或贪心状态始终不落后。边界：上界变负、最终下界不为零、连续星号、前缀右括号过多。变体：也可用双栈保存左括号和星号位置，但区间贪心更能说明状态压缩。

\noindent\textbf{LeetCode 763 划分字母区间。} 模型：每个片段必须覆盖其中所有字符的最后出现位置。推导重点：扫描维护当前片段最远右端，到达右端即可切分。

\textbf{二刷读题。} 读题时先找局部选择消耗什么资源，以及选择之后给未来留下什么边界。本题可以压成“每个片段必须覆盖其中所有字符的最后出现位置”，扫描维护当前片段最远右端，到达右端即可切分必须能被证明安全。先遮住答案，只保留题面，复述候选对象、合法性条件和输出目标。

\textbf{暴力回放。} 慢版本枚举选择顺序或用 DP 保留多个可能方案，先确认最优值存在。随后才检查局部选择是否能通过交换或领先性固定下来。二刷时先写慢版本或口述慢版本，用它确认不会漏答案。

\textbf{特例回放。} 拿 ababcbacadefegdehijhklij 手算，字符 a 的最后位置在 8，所以第一个片段至少到 8；扫描中 b、c 的最后位置也不超过 8，到 8 才能切。规律是当前片段右端是片段内所有字符最后出现位置的最大值。把这个特例继续拆开看：先构造一个非贪心选择，再比较两者留下的资源、右边界或剩余时间。只有能把非贪心第一步交换成贪心第一步且不变差，局部选择才是规律。复盘时把这个特例压成状态字段：排序后的候选、当前已选方案、剩余资源或最远边界、答案计数；每次局部选择后要能说明当前状态不落后；再用边界 单字符、所有字符相同、字符多次跨段 反查初始化、更新顺序和退出条件。

\textbf{状态回放。} 手算路线：手算时同时写一个非贪心选择，与贪心选择比较剩余资源。若无法说明贪心后剩余资源不差，就不能直接下结论。状态字段：排序后的候选、当前已选方案、剩余资源或最远边界、答案计数；每次局部选择后要能说明当前状态不落后。状态升级：把指数级选择压成一个可证明安全的局部选择；状态通常是当前最远边界、最早结束时间、最小代价或剩余资源。

\textbf{证明和边界。} 证明从交换或领先性开始：任意最优解都能替换成当前贪心选择而不变差，或贪心状态始终不落后。边界：单字符、所有字符相同、字符多次跨段。变体：这是最早合法切分贪心，能最大化片段数量。

\subsection{自测标准}

二刷时不要只确认自己见过这道题。请遮住答案，先讲题面模型，再讲暴力基线和瓶颈，然后说出优化结构保存了什么、不变量如何排除错误候选、边界实验如何打破错误实现。

\section{区间、排序与合并：二刷复盘卡}
这一大节只围绕一个题型复盘，不把题号直接铺成目录。阅读时先看复盘目标，再读核心题卡，最后用自测标准检查自己是否能独立重讲。

\subsection{复盘目标}

追问通常会问排序键为什么这样、端点相等算不算重叠、比较器是否满足严格弱序。请准备说明排序后哪些候选可以被安全丢弃。

\subsection{核心题卡}

\noindent\textbf{LeetCode 56 合并区间。} 模型：排序后可能与当前区间相交的只会是结果尾部。推导重点：按起点排序，扫描时维护结果最后区间的右端点。

\textbf{二刷读题。} 读题时先定义端点语义和冲突条件。本题可以压成“排序后可能与当前区间相交的只会是结果尾部”，排序或有序表只是在利用端点关系。先遮住答案，只保留题面，复述候选对象、合法性条件和输出目标。

\textbf{暴力回放。} 慢版本对新区间和所有旧区间两两比较，手工判断相交、覆盖或冲突。它能验证端点语义，但排序后很多远离当前边界的区间无需再看。二刷时先写慢版本或口述慢版本，用它确认不会漏答案。

\textbf{特例回放。} 拿 [[1,3],[2,6],[8,10]] 手算，排序后 [2,6] 只可能和结果尾部 [1,3] 相交，合并成 [1,6]；[8,10] 的起点大于尾部右端，另开新区间。规律是排序后当前区间只需要看结果尾部，端点相接是否合并要按闭区间语义决定。把这个特例继续拆开看：排序以后，真正还能影响当前区间的候选必须贴近当前端点、结果尾部或资源堆顶。某个区间一旦被覆盖、合并或弹出，就要说明它为什么不会和后续区间重新产生更优关系。复盘时把这个特例压成状态字段：排序后的区间、当前结果尾、扫描右端或资源堆、端点比较规则；动态版本还要保存前驱后继；再用边界 端点相接、完全覆盖、无重叠、空列表、输入未排序 反查初始化、更新顺序和退出条件。

\textbf{状态回放。} 手算路线：手算时先排序，再逐个区间写当前结果尾部、当前右端或资源堆。每个区间离开视野时都要说明它不会再影响后续。状态字段：排序后的区间、当前结果尾、扫描右端或资源堆、端点比较规则；动态版本还要保存前驱后继。状态升级：把两两比较压成排序后的相邻关系、当前边界或动态有序结构；远离当前边界的区间要被安全归档。

\textbf{证明和边界。} 证明从排序后的邻接关系开始：已经合并、选择或丢弃的区间不会被更后面的区间重新影响。边界：端点相接、完全覆盖、无重叠、空列表、输入未排序。变体：若区间是半开区间，端点相等不一定合并。

\noindent\textbf{LeetCode 57 插入区间。} 模型：新区间只会影响与它相交的一段连续区间。推导重点：先追加左侧完全不相交区间，再合并重叠段，最后追加右侧。

\textbf{二刷读题。} 读题时先定义端点语义和冲突条件。本题可以压成“新区间只会影响与它相交的一段连续区间”，排序或有序表只是在利用端点关系。先遮住答案，只保留题面，复述候选对象、合法性条件和输出目标。

\textbf{暴力回放。} 慢版本对新区间和所有旧区间两两比较，手工判断相交、覆盖或冲突。它能验证端点语义，但排序后很多远离当前边界的区间无需再看。二刷时先写慢版本或口述慢版本，用它确认不会漏答案。

\textbf{特例回放。} 拿 intervals=[[1,3],[6,9]]、new=[2,5] 手算，左侧完全小于新区间的先放入；[1,3] 与 [2,5] 相交，合成 [1,5]；[6,9] 在右侧直接追加。规律是新区间影响的只是一段连续重叠区间，先左、再合并、再右，原列表为空或新区间覆盖全部也服从同一流程。把这个特例继续拆开看：排序以后，真正还能影响当前区间的候选必须贴近当前端点、结果尾部或资源堆顶。某个区间一旦被覆盖、合并或弹出，就要说明它为什么不会和后续区间重新产生更优关系。复盘时把这个特例压成状态字段：排序后的区间、当前结果尾、扫描右端或资源堆、端点比较规则；动态版本还要保存前驱后继；再用边界 新区间在最前、最后、覆盖全部、刚好相邻、原列表为空 反查初始化、更新顺序和退出条件。

\textbf{状态回放。} 手算路线：手算时先排序，再逐个区间写当前结果尾部、当前右端或资源堆。每个区间离开视野时都要说明它不会再影响后续。状态字段：排序后的区间、当前结果尾、扫描右端或资源堆、端点比较规则；动态版本还要保存前驱后继。状态升级：把两两比较压成排序后的相邻关系、当前边界或动态有序结构；远离当前边界的区间要被安全归档。

\textbf{证明和边界。} 证明从排序后的邻接关系开始：已经合并、选择或丢弃的区间不会被更后面的区间重新影响。边界：新区间在最前、最后、覆盖全部、刚好相邻、原列表为空。变体：若输入不是有序且不重叠，必须先排序或换模型。

\noindent\textbf{LeetCode 252 会议室。} 模型：每个会议占用一个时间区间，任意重叠都表示不能全部参加。推导重点：按开始时间排序，只需检查相邻会议是否冲突。

\textbf{二刷读题。} 读题时先定义端点语义和冲突条件。本题可以压成“每个会议占用一个时间区间，任意重叠都表示不能全部参加”，排序或有序表只是在利用端点关系。先遮住答案，只保留题面，复述候选对象、合法性条件和输出目标。

\textbf{暴力回放。} 慢版本对新区间和所有旧区间两两比较，手工判断相交、覆盖或冲突。它能验证端点语义，但排序后很多远离当前边界的区间无需再看。二刷时先写慢版本或口述慢版本，用它确认不会漏答案。

\textbf{特例回放。} 拿 [[0,30],[5,10],[15,20]] 手算，排序后 [5,10] 的开始 5 小于上一场结束 30，会议冲突。规律是按开始时间排序后，只需要检查相邻会议是否重叠；端点相等是否冲突取决于半开区间语义。把这个特例继续拆开看：排序以后，真正还能影响当前区间的候选必须贴近当前端点、结果尾部或资源堆顶。某个区间一旦被覆盖、合并或弹出，就要说明它为什么不会和后续区间重新产生更优关系。复盘时把这个特例压成状态字段：排序后的区间、当前结果尾、扫描右端或资源堆、端点比较规则；动态版本还要保存前驱后继；再用边界 端点相等是否冲突、空列表、单会议、完全重叠 反查初始化、更新顺序和退出条件。

\textbf{状态回放。} 手算路线：手算时先排序，再逐个区间写当前结果尾部、当前右端或资源堆。每个区间离开视野时都要说明它不会再影响后续。状态字段：排序后的区间、当前结果尾、扫描右端或资源堆、端点比较规则；动态版本还要保存前驱后继。状态升级：把两两比较压成排序后的相邻关系、当前边界或动态有序结构；远离当前边界的区间要被安全归档。

\textbf{证明和边界。} 证明从排序后的邻接关系开始：已经合并、选择或丢弃的区间不会被更后面的区间重新影响。边界：端点相等是否冲突、空列表、单会议、完全重叠。变体：若问最少会议室，需要维护同时进行的会议数量。

\noindent\textbf{LeetCode 253 会议室 II。} 模型：同时进行的会议数量决定最少房间数。推导重点：按开始时间扫描，用小顶堆维护当前会议的最早结束时间。

\textbf{二刷读题。} 读题时先定义端点语义和冲突条件。本题可以压成“同时进行的会议数量决定最少房间数”，排序或有序表只是在利用端点关系。先遮住答案，只保留题面，复述候选对象、合法性条件和输出目标。

\textbf{暴力回放。} 慢版本对新区间和所有旧区间两两比较，手工判断相交、覆盖或冲突。它能验证端点语义，但排序后很多远离当前边界的区间无需再看。二刷时先写慢版本或口述慢版本，用它确认不会漏答案。

\textbf{特例回放。} 拿 [[0,30],[5,10],[15,20]] 手算，0-30 占一间，5-10 来时最早结束仍是 30，所以需要第二间；15-20 来时 10 已结束，可复用。规律是小顶堆保存当前占用会议的最早结束时间，堆大小就是房间数。把这个特例继续拆开看：排序以后，真正还能影响当前区间的候选必须贴近当前端点、结果尾部或资源堆顶。某个区间一旦被覆盖、合并或弹出，就要说明它为什么不会和后续区间重新产生更优关系。复盘时把这个特例压成状态字段：排序后的区间、当前结果尾、扫描右端或资源堆、端点比较规则；动态版本还要保存前驱后继；再用边界 会议端点相接、多个同一时间开始、空列表、堆顶释放条件 反查初始化、更新顺序和退出条件。

\textbf{状态回放。} 手算路线：手算时先排序，再逐个区间写当前结果尾部、当前右端或资源堆。每个区间离开视野时都要说明它不会再影响后续。状态字段：排序后的区间、当前结果尾、扫描右端或资源堆、端点比较规则；动态版本还要保存前驱后继。状态升级：把两两比较压成排序后的相邻关系、当前边界或动态有序结构；远离当前边界的区间要被安全归档。

\textbf{证明和边界。} 证明从排序后的邻接关系开始：已经合并、选择或丢弃的区间不会被更后面的区间重新影响。边界：会议端点相接、多个同一时间开始、空列表、堆顶释放条件。变体：也可以把开始和结束分别排序，用双指针统计并发数。

\noindent\textbf{LeetCode 352 将数据流变为多个不相交区间。} 模型：数字逐个到达，需要动态维护已覆盖的有序不相交区间。推导重点：用有序表找到新值附近的前驱和后继，判断是否连接、合并或新建区间。

\textbf{二刷读题。} 读题时先定义端点语义和冲突条件。本题可以压成“数字逐个到达，需要动态维护已覆盖的有序不相交区间”，排序或有序表只是在利用端点关系。先遮住答案，只保留题面，复述候选对象、合法性条件和输出目标。

\textbf{暴力回放。} 慢版本对新区间和所有旧区间两两比较，手工判断相交、覆盖或冲突。它能验证端点语义，但排序后很多远离当前边界的区间无需再看。二刷时先写慢版本或口述慢版本，用它确认不会漏答案。

\textbf{特例回放。} 拿数据流依次加入 1、3、7、2、6 手算，加入 2 时把 [1,1] 和 [3,3] 连成 [1,3]；加入 6 时和 [7,7] 连成 [6,7]。规律是有序表只需要看新值的前驱和后继，重复插入不能改变区间。把这个特例继续拆开看：排序以后，真正还能影响当前区间的候选必须贴近当前端点、结果尾部或资源堆顶。某个区间一旦被覆盖、合并或弹出，就要说明它为什么不会和后续区间重新产生更优关系。复盘时把这个特例压成状态字段：排序后的区间、当前结果尾、扫描右端或资源堆、端点比较规则；动态版本还要保存前驱后继；再用边界 重复插入、连接左右两个区间、插入最小最大值、相邻端点 反查初始化、更新顺序和退出条件。

\textbf{状态回放。} 手算路线：手算时先排序，再逐个区间写当前结果尾部、当前右端或资源堆。每个区间离开视野时都要说明它不会再影响后续。状态字段：排序后的区间、当前结果尾、扫描右端或资源堆、端点比较规则；动态版本还要保存前驱后继。状态升级：把两两比较压成排序后的相邻关系、当前边界或动态有序结构；远离当前边界的区间要被安全归档。

\textbf{证明和边界。} 证明从排序后的邻接关系开始：已经合并、选择或丢弃的区间不会被更后面的区间重新影响。边界：重复插入、连接左右两个区间、插入最小最大值、相邻端点。变体：这类动态区间不能用普通数组扫描长期维护。

\noindent\textbf{LeetCode 436 寻找右区间。} 模型：每个区间需要寻找起点不小于当前终点的最小起点区间。推导重点：把所有起点和原下标排序，对每个终点做 lower bound。

\textbf{二刷读题。} 读题时先定义端点语义和冲突条件。本题可以压成“每个区间需要寻找起点不小于当前终点的最小起点区间”，排序或有序表只是在利用端点关系。先遮住答案，只保留题面，复述候选对象、合法性条件和输出目标。

\textbf{暴力回放。} 慢版本对新区间和所有旧区间两两比较，手工判断相交、覆盖或冲突。它能验证端点语义，但排序后很多远离当前边界的区间无需再看。二刷时先写慢版本或口述慢版本，用它确认不会漏答案。

\textbf{特例回放。} 拿 [[1,2],[2,3],[3,4]] 手算，区间 [1,2] 的右区间是起点最小且不小于 2 的 [2,3]。规律是把所有起点排序，对每个终点做 lower\_bound；不存在时返回 -1。把这个特例继续拆开看：排序以后，真正还能影响当前区间的候选必须贴近当前端点、结果尾部或资源堆顶。某个区间一旦被覆盖、合并或弹出，就要说明它为什么不会和后续区间重新产生更优关系。复盘时把这个特例压成状态字段：排序后的区间、当前结果尾、扫描右端或资源堆、端点比较规则；动态版本还要保存前驱后继；再用边界 不存在右区间、起点相同、负端点、单区间 反查初始化、更新顺序和退出条件。

\textbf{状态回放。} 手算路线：手算时先排序，再逐个区间写当前结果尾部、当前右端或资源堆。每个区间离开视野时都要说明它不会再影响后续。状态字段：排序后的区间、当前结果尾、扫描右端或资源堆、端点比较规则；动态版本还要保存前驱后继。状态升级：把两两比较压成排序后的相邻关系、当前边界或动态有序结构；远离当前边界的区间要被安全归档。

\textbf{证明和边界。} 证明从排序后的邻接关系开始：已经合并、选择或丢弃的区间不会被更后面的区间重新影响。边界：不存在右区间、起点相同、负端点、单区间。变体：如果区间动态插入，则需要有序表而不是一次排序数组。

\noindent\textbf{LeetCode 729 我的日程安排表 I。} 模型：每次插入新区间时，需要判断它是否和已有区间重叠。推导重点：有序表按起点存区间，只检查新区间的前驱和后继。

\textbf{二刷读题。} 读题时先定义端点语义和冲突条件。本题可以压成“每次插入新区间时，需要判断它是否和已有区间重叠”，排序或有序表只是在利用端点关系。先遮住答案，只保留题面，复述候选对象、合法性条件和输出目标。

\textbf{暴力回放。} 慢版本对新区间和所有旧区间两两比较，手工判断相交、覆盖或冲突。它能验证端点语义，但排序后很多远离当前边界的区间无需再看。二刷时先写慢版本或口述慢版本，用它确认不会漏答案。

\textbf{特例回放。} 拿 book(10,20) 成功，再 book(15,25) 手算，它和已有区间重叠应失败；book(20,30) 与 [10,20) 端点相接可成功。规律是半开区间下只需检查前驱和后继是否重叠。把这个特例继续拆开看：排序以后，真正还能影响当前区间的候选必须贴近当前端点、结果尾部或资源堆顶。某个区间一旦被覆盖、合并或弹出，就要说明它为什么不会和后续区间重新产生更优关系。复盘时把这个特例压成状态字段：排序后的区间、当前结果尾、扫描右端或资源堆、端点比较规则；动态版本还要保存前驱后继；再用边界 端点相接、插入最前最后、完全覆盖已有区间、动态多次插入 反查初始化、更新顺序和退出条件。

\textbf{状态回放。} 手算路线：手算时先排序，再逐个区间写当前结果尾部、当前右端或资源堆。每个区间离开视野时都要说明它不会再影响后续。状态字段：排序后的区间、当前结果尾、扫描右端或资源堆、端点比较规则；动态版本还要保存前驱后继。状态升级：把两两比较压成排序后的相邻关系、当前边界或动态有序结构；远离当前边界的区间要被安全归档。

\textbf{证明和边界。} 证明从排序后的邻接关系开始：已经合并、选择或丢弃的区间不会被更后面的区间重新影响。边界：端点相接、插入最前最后、完全覆盖已有区间、动态多次插入。变体：哈希表无法回答前驱后继，动态区间要用有序结构。

\noindent\textbf{LeetCode 731 我的日程安排表 II。} 模型：允许双重预订但不允许三重预订，插入会改变区间重叠层数。推导重点：保存已有预订和已经双重预订的区间；新预订若碰到双重区间则失败，否则记录新产生的重叠。

\textbf{二刷读题。} 读题时先定义端点语义和冲突条件。本题可以压成“允许双重预订但不允许三重预订，插入会改变区间重叠层数”，排序或有序表只是在利用端点关系。先遮住答案，只保留题面，复述候选对象、合法性条件和输出目标。

\textbf{暴力回放。} 慢版本对新区间和所有旧区间两两比较，手工判断相交、覆盖或冲突。它能验证端点语义，但排序后很多远离当前边界的区间无需再看。二刷时先写慢版本或口述慢版本，用它确认不会漏答案。

\textbf{特例回放。} 拿 book(10,20)、book(15,25) 后，重叠段 [15,20) 已经是双重预订；再 book(17,22) 会碰到双重段，必须失败。规律是保存所有单次预订和已形成的双重预订，新区间不能与任何双重区间相交。把这个特例继续拆开看：排序以后，真正还能影响当前区间的候选必须贴近当前端点、结果尾部或资源堆顶。某个区间一旦被覆盖、合并或弹出，就要说明它为什么不会和后续区间重新产生更优关系。复盘时把这个特例压成状态字段：排序后的区间、当前结果尾、扫描右端或资源堆、端点比较规则；动态版本还要保存前驱后继；再用边界 端点相接、多个局部重叠、完全覆盖、回滚失败插入 反查初始化、更新顺序和退出条件。

\textbf{状态回放。} 手算路线：手算时先排序，再逐个区间写当前结果尾部、当前右端或资源堆。每个区间离开视野时都要说明它不会再影响后续。状态字段：排序后的区间、当前结果尾、扫描右端或资源堆、端点比较规则；动态版本还要保存前驱后继。状态升级：把两两比较压成排序后的相邻关系、当前边界或动态有序结构；远离当前边界的区间要被安全归档。

\textbf{证明和边界。} 证明从排序后的邻接关系开始：已经合并、选择或丢弃的区间不会被更后面的区间重新影响。边界：端点相接、多个局部重叠、完全覆盖、回滚失败插入。变体：扫描线差分也能做，但要处理失败时恢复状态。

\noindent\textbf{LeetCode 986 区间列表的交集。} 模型：两个有序且不重叠的区间列表，交集只可能来自当前两个区间。推导重点：每次比较两个当前区间的重叠部分，然后推进右端较小的区间。

\textbf{二刷读题。} 读题时先定义端点语义和冲突条件。本题可以压成“两个有序且不重叠的区间列表，交集只可能来自当前两个区间”，排序或有序表只是在利用端点关系。先遮住答案，只保留题面，复述候选对象、合法性条件和输出目标。

\textbf{暴力回放。} 慢版本对新区间和所有旧区间两两比较，手工判断相交、覆盖或冲突。它能验证端点语义，但排序后很多远离当前边界的区间无需再看。二刷时先写慢版本或口述慢版本，用它确认不会漏答案。

\textbf{特例回放。} 拿 A=[0,2],[5,10] 和 B=[1,5],[8,12] 手算，当前两个区间交集是 [1,2]，然后推进右端较小的 [0,2]。规律是两个有序列表只需要比较当前区间，谁先结束谁不可能再和对方后续产生交集。把这个特例继续拆开看：排序以后，真正还能影响当前区间的候选必须贴近当前端点、结果尾部或资源堆顶。某个区间一旦被覆盖、合并或弹出，就要说明它为什么不会和后续区间重新产生更优关系。复盘时把这个特例压成状态字段：排序后的区间、当前结果尾、扫描右端或资源堆、端点比较规则；动态版本还要保存前驱后继；再用边界 空列表、端点相接、一个区间覆盖多个区间、无交集 反查初始化、更新顺序和退出条件。

\textbf{状态回放。} 手算路线：手算时先排序，再逐个区间写当前结果尾部、当前右端或资源堆。每个区间离开视野时都要说明它不会再影响后续。状态字段：排序后的区间、当前结果尾、扫描右端或资源堆、端点比较规则；动态版本还要保存前驱后继。状态升级：把两两比较压成排序后的相邻关系、当前边界或动态有序结构；远离当前边界的区间要被安全归档。

\textbf{证明和边界。} 证明从排序后的邻接关系开始：已经合并、选择或丢弃的区间不会被更后面的区间重新影响。边界：空列表、端点相接、一个区间覆盖多个区间、无交集。变体：这是区间双指针，不需要把两个列表合并排序。

\noindent\textbf{LeetCode 1094 拼车。} 模型：乘客上下车事件构成沿路线位置变化的载客量。推导重点：用差分数组或扫描线记录每个位置人数变化，累计人数不能超过容量。

\textbf{二刷读题。} 读题时先定义端点语义和冲突条件。本题可以压成“乘客上下车事件构成沿路线位置变化的载客量”，排序或有序表只是在利用端点关系。先遮住答案，只保留题面，复述候选对象、合法性条件和输出目标。

\textbf{暴力回放。} 慢版本对新区间和所有旧区间两两比较，手工判断相交、覆盖或冲突。它能验证端点语义，但排序后很多远离当前边界的区间无需再看。二刷时先写慢版本或口述慢版本，用它确认不会漏答案。

\textbf{特例回放。} 拿 trips=[[2,1,5],[3,3,7]]、capacity=4 手算，位置 3 到 5 之间车上人数是 5，超过容量。规律是上车点加人数、下车点减人数，按位置前缀累计检查容量。把这个特例继续拆开看：排序以后，真正还能影响当前区间的候选必须贴近当前端点、结果尾部或资源堆顶。某个区间一旦被覆盖、合并或弹出，就要说明它为什么不会和后续区间重新产生更优关系。复盘时把这个特例压成状态字段：排序后的区间、当前结果尾、扫描右端或资源堆、端点比较规则；动态版本还要保存前驱后继；再用边界 同站上下车、容量刚好、站点范围、输入未排序 反查初始化、更新顺序和退出条件。

\textbf{状态回放。} 手算路线：手算时先排序，再逐个区间写当前结果尾部、当前右端或资源堆。每个区间离开视野时都要说明它不会再影响后续。状态字段：排序后的区间、当前结果尾、扫描右端或资源堆、端点比较规则；动态版本还要保存前驱后继。状态升级：把两两比较压成排序后的相邻关系、当前边界或动态有序结构；远离当前边界的区间要被安全归档。

\textbf{证明和边界。} 证明从排序后的邻接关系开始：已经合并、选择或丢弃的区间不会被更后面的区间重新影响。边界：同站上下车、容量刚好、站点范围、输入未排序。变体：航班预订统计也是差分思想，只是只需要最终每段累加结果。

\noindent\textbf{LeetCode 1109 航班预订统计。} 模型：每条预订给一个连续航班区间增加座位数。推导重点：差分数组在区间起点加人数，终点后一位减人数，最后前缀还原每航班总数。

\textbf{二刷读题。} 读题时先定义端点语义和冲突条件。本题可以压成“每条预订给一个连续航班区间增加座位数”，排序或有序表只是在利用端点关系。先遮住答案，只保留题面，复述候选对象、合法性条件和输出目标。

\textbf{暴力回放。} 慢版本对新区间和所有旧区间两两比较，手工判断相交、覆盖或冲突。它能验证端点语义，但排序后很多远离当前边界的区间无需再看。二刷时先写慢版本或口述慢版本，用它确认不会漏答案。

\textbf{特例回放。} 拿 bookings=[[1,2,10],[2,3,20]]、n=3 手算，第 1 班加 10，第 2 班加 30，第 3 班加 20。规律是区间加法用差分：起点加，终点后一位减，最后前缀还原。把这个特例继续拆开看：排序以后，真正还能影响当前区间的候选必须贴近当前端点、结果尾部或资源堆顶。某个区间一旦被覆盖、合并或弹出，就要说明它为什么不会和后续区间重新产生更优关系。复盘时把这个特例压成状态字段：排序后的区间、当前结果尾、扫描右端或资源堆、端点比较规则；动态版本还要保存前驱后继；再用边界 区间到最后一个航班、单点区间、多条重叠预订、下标从一开始 反查初始化、更新顺序和退出条件。

\textbf{状态回放。} 手算路线：手算时先排序，再逐个区间写当前结果尾部、当前右端或资源堆。每个区间离开视野时都要说明它不会再影响后续。状态字段：排序后的区间、当前结果尾、扫描右端或资源堆、端点比较规则；动态版本还要保存前驱后继。状态升级：把两两比较压成排序后的相邻关系、当前边界或动态有序结构；远离当前边界的区间要被安全归档。

\textbf{证明和边界。} 证明从排序后的邻接关系开始：已经合并、选择或丢弃的区间不会被更后面的区间重新影响。边界：区间到最后一个航班、单点区间、多条重叠预订、下标从一开始。变体：这是静态区间加法，不需要线段树。

\noindent\textbf{LeetCode 1229 安排会议日程。} 模型：两个人的可用时间段都有序，目标是找最早长度足够的交集。推导重点：双指针比较当前两个区间交集，若长度不足就推进结束更早的一侧。

\textbf{二刷读题。} 读题时先定义端点语义和冲突条件。本题可以压成“两个人的可用时间段都有序，目标是找最早长度足够的交集”，排序或有序表只是在利用端点关系。先遮住答案，只保留题面，复述候选对象、合法性条件和输出目标。

\textbf{暴力回放。} 慢版本对新区间和所有旧区间两两比较，手工判断相交、覆盖或冲突。它能验证端点语义，但排序后很多远离当前边界的区间无需再看。二刷时先写慢版本或口述慢版本，用它确认不会漏答案。

\textbf{特例回放。} 拿 slots1=[10,50],[60,120]，slots2=[0,15],[60,70]，duration=8 手算，第一组交集 [10,15] 不够，第二组交集 [60,70] 足够，返回 [60,68]。规律是双指针比较当前交集，不够时推进结束更早的一侧。把这个特例继续拆开看：排序以后，真正还能影响当前区间的候选必须贴近当前端点、结果尾部或资源堆顶。某个区间一旦被覆盖、合并或弹出，就要说明它为什么不会和后续区间重新产生更优关系。复盘时把这个特例压成状态字段：排序后的区间、当前结果尾、扫描右端或资源堆、端点比较规则；动态版本还要保存前驱后继；再用边界 无可行时间、刚好等于 duration、区间端点相接、输入需要先排序 反查初始化、更新顺序和退出条件。

\textbf{状态回放。} 手算路线：手算时先排序，再逐个区间写当前结果尾部、当前右端或资源堆。每个区间离开视野时都要说明它不会再影响后续。状态字段：排序后的区间、当前结果尾、扫描右端或资源堆、端点比较规则；动态版本还要保存前驱后继。状态升级：把两两比较压成排序后的相邻关系、当前边界或动态有序结构；远离当前边界的区间要被安全归档。

\textbf{证明和边界。} 证明从排序后的邻接关系开始：已经合并、选择或丢弃的区间不会被更后面的区间重新影响。边界：无可行时间、刚好等于 duration、区间端点相接、输入需要先排序。变体：它和区间列表交集相同，但额外带最早可行和长度约束。

\noindent\textbf{LeetCode 1288 删除被覆盖区间。} 模型：被覆盖要求另一区间起点不晚且终点不早。推导重点：按起点升序、终点降序排序，扫描维护最大右端。

\textbf{二刷读题。} 读题时先定义端点语义和冲突条件。本题可以压成“被覆盖要求另一区间起点不晚且终点不早”，排序或有序表只是在利用端点关系。先遮住答案，只保留题面，复述候选对象、合法性条件和输出目标。

\textbf{暴力回放。} 慢版本对新区间和所有旧区间两两比较，手工判断相交、覆盖或冲突。它能验证端点语义，但排序后很多远离当前边界的区间无需再看。二刷时先写慢版本或口述慢版本，用它确认不会漏答案。

\textbf{特例回放。} 拿 [[1,4],[3,6],[2,8]] 手算，[3,6] 被 [2,8] 覆盖，应删除；同起点时要先看更长区间，否则短区间会干扰判断。规律是起点升序、终点降序排序，扫描维护最大右端。把这个特例继续拆开看：排序以后，真正还能影响当前区间的候选必须贴近当前端点、结果尾部或资源堆顶。某个区间一旦被覆盖、合并或弹出，就要说明它为什么不会和后续区间重新产生更优关系。复盘时把这个特例压成状态字段：排序后的区间、当前结果尾、扫描右端或资源堆、端点比较规则；动态版本还要保存前驱后继；再用边界 同起点区间、完全覆盖、无覆盖、端点相等 反查初始化、更新顺序和退出条件。

\textbf{状态回放。} 手算路线：手算时先排序，再逐个区间写当前结果尾部、当前右端或资源堆。每个区间离开视野时都要说明它不会再影响后续。状态字段：排序后的区间、当前结果尾、扫描右端或资源堆、端点比较规则；动态版本还要保存前驱后继。状态升级：把两两比较压成排序后的相邻关系、当前边界或动态有序结构；远离当前边界的区间要被安全归档。

\textbf{证明和边界。} 证明从排序后的邻接关系开始：已经合并、选择或丢弃的区间不会被更后面的区间重新影响。边界：同起点区间、完全覆盖、无覆盖、端点相等。变体：终点降序是关键，否则同起点小区间会先出现造成误判。

\subsection{自测标准}

二刷时不要只确认自己见过这道题。请遮住答案，先讲题面模型，再讲暴力基线和瓶颈，然后说出优化结构保存了什么、不变量如何排除错误候选、边界实验如何打破错误实现。

\section{位运算与状态压缩：二刷复盘卡}
这一大节只围绕一个题型复盘，不把题号直接铺成目录。阅读时先看复盘目标，再读核心题卡，最后用自测标准检查自己是否能独立重讲。

\subsection{复盘目标}

追问通常会问每一位代表什么、负数和移位是否安全、状态相同但资源不同能否合并。请准备说明位运算只是编码，真正关键仍是状态语义。

\subsection{核心题卡}

\noindent\textbf{LeetCode 29 两数相除。} 模型：除法可以看成不断减去除数的倍增值。推导重点：把除数按二进制倍增，优先减去不超过被除数的最大倍数并累加商。

\textbf{二刷读题。} 读题时先用普通状态解释每一项布尔信息。本题可以压成“除法可以看成不断减去除数的倍增值”，位运算只是更紧凑的编码。先遮住答案，只保留题面，复述候选对象、合法性条件和输出目标。

\textbf{暴力回放。} 慢版本先用集合、数组或布尔表保存状态，逐步执行题目动作。确认每个布尔字段的含义后，才能把它压缩成位，而不是直接写位运算公式。二刷时先写慢版本或口述慢版本，用它确认不会漏答案。

\textbf{特例回放。} 拿 43 除以 5 手算，5 的倍增序列是 5、10、20、40，先减 40 商加 8，余 3 停止。规律是用二进制倍增找不超过被除数的最大除数倍数，符号和 int 最小值溢出要单独处理。把这个特例继续拆开看：每一位或每个掩码位都要对应题目里的一个布尔事实。位运算只是压缩表示；如果两个状态在位置相同但掩码不同，未来能力不同，就不能合并。复盘时把这个特例压成状态字段：掩码含义、当前位、目标位集合、状态转移和答案读取；负数或高位参与时要说明整数宽度；再用边界 除数为一、被除数为 int 最小值、结果溢出、符号处理 反查初始化、更新顺序和退出条件。

\textbf{状态回放。} 手算路线：手算时把整数写成二进制，并在每一位旁边标注含义。状态转移只允许改变题目动作允许改变的那些位。状态字段：掩码含义、当前位、目标位集合、状态转移和答案读取；负数或高位参与时要说明整数宽度。状态升级：把集合状态压成掩码；包含、加入、删除、枚举子集都变成位操作，但每一位的题目含义不能丢。

\textbf{证明和边界。} 证明从逐位独立或子集归纳开始：被压缩到位里的信息足以决定未来转移。边界：除数为一、被除数为 int 最小值、结果溢出、符号处理。变体：它和快速幂类似，都利用二进制拆分减少重复加减。

\noindent\textbf{LeetCode 136 只出现一次的数字。} 模型：成对元素可以用异或抵消，剩下单个元素。推导重点：异或满足交换结合，扫描顺序不影响答案。

\textbf{二刷读题。} 读题时先用普通状态解释每一项布尔信息。本题可以压成“成对元素可以用异或抵消，剩下单个元素”，位运算只是更紧凑的编码。先遮住答案，只保留题面，复述候选对象、合法性条件和输出目标。

\textbf{暴力回放。} 慢版本先用集合、数组或布尔表保存状态，逐步执行题目动作。确认每个布尔字段的含义后，才能把它压缩成位，而不是直接写位运算公式。二刷时先写慢版本或口述慢版本，用它确认不会漏答案。

\textbf{特例回放。} 拿 [4,1,2,1,2] 手算，1 和 1 异或抵消，2 和 2 抵消，最后剩 4。规律是异或满足交换结合，且相同数异或后变成零，扫描顺序不影响结果；零和负数也只是按位参与同样运算。把这个特例继续拆开看：每一位或每个掩码位都要对应题目里的一个布尔事实。位运算只是压缩表示；如果两个状态在位置相同但掩码不同，未来能力不同，就不能合并。复盘时把这个特例压成状态字段：掩码含义、当前位、目标位集合、状态转移和答案读取；负数或高位参与时要说明整数宽度；再用边界 零、负数、唯一值在开头或结尾、所有其他值恰好两次 反查初始化、更新顺序和退出条件。

\textbf{状态回放。} 手算路线：手算时把整数写成二进制，并在每一位旁边标注含义。状态转移只允许改变题目动作允许改变的那些位。状态字段：掩码含义、当前位、目标位集合、状态转移和答案读取；负数或高位参与时要说明整数宽度。状态升级：把集合状态压成掩码；包含、加入、删除、枚举子集都变成位操作，但每一位的题目含义不能丢。

\textbf{证明和边界。} 证明从逐位独立或子集归纳开始：被压缩到位里的信息足以决定未来转移。边界：零、负数、唯一值在开头或结尾、所有其他值恰好两次。变体：若其他值出现三次，要改为按位计数取模。

\noindent\textbf{LeetCode 137 只出现一次的数字 II。} 模型：除一个数字外，其余数字都出现三次，单纯异或不能抵消。推导重点：逐位统计一的数量，对三取模后还原目标数字。

\textbf{二刷读题。} 读题时先用普通状态解释每一项布尔信息。本题可以压成“除一个数字外，其余数字都出现三次，单纯异或不能抵消”，位运算只是更紧凑的编码。先遮住答案，只保留题面，复述候选对象、合法性条件和输出目标。

\textbf{暴力回放。} 慢版本先用集合、数组或布尔表保存状态，逐步执行题目动作。确认每个布尔字段的含义后，才能把它压缩成位，而不是直接写位运算公式。二刷时先写慢版本或口述慢版本，用它确认不会漏答案。

\textbf{特例回放。} 拿 [2,2,3,2] 手算，每一位上 2 的贡献出现三次，对 3 取模后消失，只剩 3 的位。规律是逐位计数并对 3 取模，符号位也必须按固定整数宽度处理。把这个特例继续拆开看：每一位或每个掩码位都要对应题目里的一个布尔事实。位运算只是压缩表示；如果两个状态在位置相同但掩码不同，未来能力不同，就不能合并。复盘时把这个特例压成状态字段：掩码含义、当前位、目标位集合、状态转移和答案读取；负数或高位参与时要说明整数宽度；再用边界 负数、符号位、目标为零、所有其他数恰好三次 反查初始化、更新顺序和退出条件。

\textbf{状态回放。} 手算路线：手算时把整数写成二进制，并在每一位旁边标注含义。状态转移只允许改变题目动作允许改变的那些位。状态字段：掩码含义、当前位、目标位集合、状态转移和答案读取；负数或高位参与时要说明整数宽度。状态升级：把集合状态压成掩码；包含、加入、删除、枚举子集都变成位操作，但每一位的题目含义不能丢。

\textbf{证明和边界。} 证明从逐位独立或子集归纳开始：被压缩到位里的信息足以决定未来转移。边界：负数、符号位、目标为零、所有其他数恰好三次。变体：若其他数出现 k 次，思路仍是按位计数取模。

\noindent\textbf{LeetCode 190 颠倒二进制位。} 模型：每一位的新位置由原位置镜像决定。推导重点：循环三十二次，把结果左移并接上当前最低位，再右移输入。

\textbf{二刷读题。} 读题时先用普通状态解释每一项布尔信息。本题可以压成“每一位的新位置由原位置镜像决定”，位运算只是更紧凑的编码。先遮住答案，只保留题面，复述候选对象、合法性条件和输出目标。

\textbf{暴力回放。} 慢版本先用集合、数组或布尔表保存状态，逐步执行题目动作。确认每个布尔字段的含义后，才能把它压缩成位，而不是直接写位运算公式。二刷时先写慢版本或口述慢版本，用它确认不会漏答案。

\textbf{特例回放。} 拿低位是 101 的数手算，颠倒时每次把结果左移一位，再接上当前最低位；原数右移继续取下一位。规律是循环 32 次处理固定宽度，不能因为高位暂时为 0 就提前结束。把这个特例继续拆开看：每一位或每个掩码位都要对应题目里的一个布尔事实。位运算只是压缩表示；如果两个状态在位置相同但掩码不同，未来能力不同，就不能合并。复盘时把这个特例压成状态字段：掩码含义、当前位、目标位集合、状态转移和答案读取；负数或高位参与时要说明整数宽度；再用边界 最高位、最低位、全零、全一、无符号语义 反查初始化、更新顺序和退出条件。

\textbf{状态回放。} 手算路线：手算时把整数写成二进制，并在每一位旁边标注含义。状态转移只允许改变题目动作允许改变的那些位。状态字段：掩码含义、当前位、目标位集合、状态转移和答案读取；负数或高位参与时要说明整数宽度。状态升级：把集合状态压成掩码；包含、加入、删除、枚举子集都变成位操作，但每一位的题目含义不能丢。

\textbf{证明和边界。} 证明从逐位独立或子集归纳开始：被压缩到位里的信息足以决定未来转移。边界：最高位、最低位、全零、全一、无符号语义。变体：若多次调用，可以按字节预处理反转表。

\noindent\textbf{LeetCode 191 位 1 的个数。} 模型：每次清掉最低位的一可以让循环次数等于一的个数。推导重点：使用 n \& (n - 1) 删除最低位一，比逐位检查更直接。

\textbf{二刷读题。} 读题时先用普通状态解释每一项布尔信息。本题可以压成“每次清掉最低位的一可以让循环次数等于一的个数”，位运算只是更紧凑的编码。先遮住答案，只保留题面，复述候选对象、合法性条件和输出目标。

\textbf{暴力回放。} 慢版本先用集合、数组或布尔表保存状态，逐步执行题目动作。确认每个布尔字段的含义后，才能把它压缩成位，而不是直接写位运算公式。二刷时先写慢版本或口述慢版本，用它确认不会漏答案。

\textbf{特例回放。} 拿 11 的二进制 1011 手算，n\&(n-1) 第一次删掉最低位 1 得 1010，第二次得 1000，第三次得 0。规律是每次操作恰好删除一个 1，所以循环次数就是 1 的个数。把这个特例继续拆开看：每一位或每个掩码位都要对应题目里的一个布尔事实。位运算只是压缩表示；如果两个状态在位置相同但掩码不同，未来能力不同，就不能合并。复盘时把这个特例压成状态字段：掩码含义、当前位、目标位集合、状态转移和答案读取；负数或高位参与时要说明整数宽度；再用边界 零、只有最高位、一的个数很多、无符号输入 反查初始化、更新顺序和退出条件。

\textbf{状态回放。} 手算路线：手算时把整数写成二进制，并在每一位旁边标注含义。状态转移只允许改变题目动作允许改变的那些位。状态字段：掩码含义、当前位、目标位集合、状态转移和答案读取；负数或高位参与时要说明整数宽度。状态升级：把集合状态压成掩码；包含、加入、删除、枚举子集都变成位操作，但每一位的题目含义不能丢。

\textbf{证明和边界。} 证明从逐位独立或子集归纳开始：被压缩到位里的信息足以决定未来转移。边界：零、只有最高位、一的个数很多、无符号输入。变体：比特位计数可以把这个过程 DP 化到一整段数。

\noindent\textbf{LeetCode 201 数字范围按位与。} 模型：区间内所有数按位与后，只保留左右端公共二进制前缀。推导重点：不断右移 left 和 right，直到二者相等，再把公共前缀移回原位。

\textbf{二刷读题。} 读题时先用普通状态解释每一项布尔信息。本题可以压成“区间内所有数按位与后，只保留左右端公共二进制前缀”，位运算只是更紧凑的编码。先遮住答案，只保留题面，复述候选对象、合法性条件和输出目标。

\textbf{暴力回放。} 慢版本先用集合、数组或布尔表保存状态，逐步执行题目动作。确认每个布尔字段的含义后，才能把它压缩成位，而不是直接写位运算公式。二刷时先写慢版本或口述慢版本，用它确认不会漏答案。

\textbf{特例回放。} 拿区间 [5,7]，二进制是 101、110、111，只有最高公共前缀 1 能留下，低位在区间内会出现 0 和 1，被按位与清零。规律是不断右移 left 和 right 找公共前缀，再移回原位。把这个特例继续拆开看：每一位或每个掩码位都要对应题目里的一个布尔事实。位运算只是压缩表示；如果两个状态在位置相同但掩码不同，未来能力不同，就不能合并。复盘时把这个特例压成状态字段：掩码含义、当前位、目标位集合、状态转移和答案读取；负数或高位参与时要说明整数宽度；再用边界 left 等于 right、区间跨过二的幂边界、结果为零 反查初始化、更新顺序和退出条件。

\textbf{状态回放。} 手算路线：手算时把整数写成二进制，并在每一位旁边标注含义。状态转移只允许改变题目动作允许改变的那些位。状态字段：掩码含义、当前位、目标位集合、状态转移和答案读取；负数或高位参与时要说明整数宽度。状态升级：把集合状态压成掩码；包含、加入、删除、枚举子集都变成位操作，但每一位的题目含义不能丢。

\textbf{证明和边界。} 证明从逐位独立或子集归纳开始：被压缩到位里的信息足以决定未来转移。边界：left 等于 right、区间跨过二的幂边界、结果为零。变体：也可以不断清除 right 的最低位一，直到 right 不大于 left。

\noindent\textbf{LeetCode 231 2 的幂。} 模型：二的幂在二进制中只有一个一。推导重点：正数 n 满足 n 与 n 减一按位与为零。

\textbf{二刷读题。} 读题时先用普通状态解释每一项布尔信息。本题可以压成“二的幂在二进制中只有一个一”，位运算只是更紧凑的编码。先遮住答案，只保留题面，复述候选对象、合法性条件和输出目标。

\textbf{暴力回放。} 慢版本先用集合、数组或布尔表保存状态，逐步执行题目动作。确认每个布尔字段的含义后，才能把它压缩成位，而不是直接写位运算公式。二刷时先写慢版本或口述慢版本，用它确认不会漏答案。

\textbf{特例回放。} 拿 8 的二进制 1000 手算，8\&(7)=1000\&0111=0；拿 6 的 110，6\&5=100 不为 0。规律是 2 的幂只有一个 1，且 n 必须为正；0 和负数不能只靠按位与判断。把这个特例继续拆开看：每一位或每个掩码位都要对应题目里的一个布尔事实。位运算只是压缩表示；如果两个状态在位置相同但掩码不同，未来能力不同，就不能合并。复盘时把这个特例压成状态字段：掩码含义、当前位、目标位集合、状态转移和答案读取；负数或高位参与时要说明整数宽度；再用边界 n 为零、负数、int 最小值、最高位 反查初始化、更新顺序和退出条件。

\textbf{状态回放。} 手算路线：手算时把整数写成二进制，并在每一位旁边标注含义。状态转移只允许改变题目动作允许改变的那些位。状态字段：掩码含义、当前位、目标位集合、状态转移和答案读取；负数或高位参与时要说明整数宽度。状态升级：把集合状态压成掩码；包含、加入、删除、枚举子集都变成位操作，但每一位的题目含义不能丢。

\textbf{证明和边界。} 证明从逐位独立或子集归纳开始：被压缩到位里的信息足以决定未来转移。边界：n 为零、负数、int 最小值、最高位。变体：判断四的幂还要额外限制一所在的位置为偶数位。

\noindent\textbf{LeetCode 260 只出现一次的数字 III。} 模型：两个只出现一次的数异或后至少有一位不同。推导重点：用最低不同位把数组分成两组，每组内部再异或抵消。

\textbf{二刷读题。} 读题时先用普通状态解释每一项布尔信息。本题可以压成“两个只出现一次的数异或后至少有一位不同”，位运算只是更紧凑的编码。先遮住答案，只保留题面，复述候选对象、合法性条件和输出目标。

\textbf{暴力回放。} 慢版本先用集合、数组或布尔表保存状态，逐步执行题目动作。确认每个布尔字段的含义后，才能把它压缩成位，而不是直接写位运算公式。二刷时先写慢版本或口述慢版本，用它确认不会漏答案。

\textbf{特例回放。} 拿 [1,2,1,3,2,5] 手算，全部异或后只剩 3 和 5 的异或结果，取最低不同位可把 3 和 5 分到两组；每组内部重复数再异或抵消。规律是两个答案至少有一位不同，用这位分组后退化成两个只出现一次的问题。把这个特例继续拆开看：每一位或每个掩码位都要对应题目里的一个布尔事实。位运算只是压缩表示；如果两个状态在位置相同但掩码不同，未来能力不同，就不能合并。复盘时把这个特例压成状态字段：掩码含义、当前位、目标位集合、状态转移和答案读取；负数或高位参与时要说明整数宽度；再用边界 两个答案一正一负、最低位分组、其他数恰好两次 反查初始化、更新顺序和退出条件。

\textbf{状态回放。} 手算路线：手算时把整数写成二进制，并在每一位旁边标注含义。状态转移只允许改变题目动作允许改变的那些位。状态字段：掩码含义、当前位、目标位集合、状态转移和答案读取；负数或高位参与时要说明整数宽度。状态升级：把集合状态压成掩码；包含、加入、删除、枚举子集都变成位操作，但每一位的题目含义不能丢。

\textbf{证明和边界。} 证明从逐位独立或子集归纳开始：被压缩到位里的信息足以决定未来转移。边界：两个答案一正一负、最低位分组、其他数恰好两次。变体：它是在异或抵消基础上增加分组位。

\noindent\textbf{LeetCode 268 丢失的数字。} 模型：下标零到 n 与数组值相比，缺失值会在异或抵消后剩下。推导重点：把所有下标和所有值一起异或，重复出现的数抵消。

\textbf{二刷读题。} 读题时先用普通状态解释每一项布尔信息。本题可以压成“下标零到 n 与数组值相比，缺失值会在异或抵消后剩下”，位运算只是更紧凑的编码。先遮住答案，只保留题面，复述候选对象、合法性条件和输出目标。

\textbf{暴力回放。} 慢版本先用集合、数组或布尔表保存状态，逐步执行题目动作。确认每个布尔字段的含义后，才能把它压缩成位，而不是直接写位运算公式。二刷时先写慢版本或口述慢版本，用它确认不会漏答案。

\textbf{特例回放。} 拿 [3,0,1] 手算，0、1、3 与完整下标 0、1、2、3 一起异或，重复的 0、1、3 抵消，剩 2。规律是缺失数可以由异或抵消得到，也可用求和但要注意溢出。把这个特例继续拆开看：每一位或每个掩码位都要对应题目里的一个布尔事实。位运算只是压缩表示；如果两个状态在位置相同但掩码不同，未来能力不同，就不能合并。复盘时把这个特例压成状态字段：掩码含义、当前位、目标位集合、状态转移和答案读取；负数或高位参与时要说明整数宽度；再用边界 缺失零、缺失 n、数组长度一、也可用求和但要防溢出 反查初始化、更新顺序和退出条件。

\textbf{状态回放。} 手算路线：手算时把整数写成二进制，并在每一位旁边标注含义。状态转移只允许改变题目动作允许改变的那些位。状态字段：掩码含义、当前位、目标位集合、状态转移和答案读取；负数或高位参与时要说明整数宽度。状态升级：把集合状态压成掩码；包含、加入、删除、枚举子集都变成位操作，但每一位的题目含义不能丢。

\textbf{证明和边界。} 证明从逐位独立或子集归纳开始：被压缩到位里的信息足以决定未来转移。边界：缺失零、缺失 n、数组长度一、也可用求和但要防溢出。变体：异或法适合展示值域完整且只有一个缺口的结构。

\noindent\textbf{LeetCode 287 寻找重复数。} 模型：数组值域使下标到值形成函数图，重复数是环入口。推导重点：快慢指针不修改数组且常数空间；值域二分用抽屉原理。

\textbf{二刷读题。} 读题时先用普通状态解释每一项布尔信息。本题可以压成“数组值域使下标到值形成函数图，重复数是环入口”，位运算只是更紧凑的编码。先遮住答案，只保留题面，复述候选对象、合法性条件和输出目标。

\textbf{暴力回放。} 慢版本先用集合、数组或布尔表保存状态，逐步执行题目动作。确认每个布尔字段的含义后，才能把它压缩成位，而不是直接写位运算公式。二刷时先写慢版本或口述慢版本，用它确认不会漏答案。

\textbf{特例回放。} 拿 [1,3,4,2,2] 手算，把值当成 next 指针会形成 1->3->2->4->2 的环，入环点 2 就是重复数。规律是长度 n+1、值域 1..n 强制形成环；不能修改数组时可用快慢指针或值域二分。把这个特例继续拆开看：每一位或每个掩码位都要对应题目里的一个布尔事实。位运算只是压缩表示；如果两个状态在位置相同但掩码不同，未来能力不同，就不能合并。复盘时把这个特例压成状态字段：掩码含义、当前位、目标位集合、状态转移和答案读取；负数或高位参与时要说明整数宽度；再用边界 重复数出现多次、不能排序、不能改数组、长度为 n 加一 反查初始化、更新顺序和退出条件。

\textbf{状态回放。} 手算路线：手算时把整数写成二进制，并在每一位旁边标注含义。状态转移只允许改变题目动作允许改变的那些位。状态字段：掩码含义、当前位、目标位集合、状态转移和答案读取；负数或高位参与时要说明整数宽度。状态升级：把集合状态压成掩码；包含、加入、删除、枚举子集都变成位操作，但每一位的题目含义不能丢。

\textbf{证明和边界。} 证明从逐位独立或子集归纳开始：被压缩到位里的信息足以决定未来转移。边界：重复数出现多次、不能排序、不能改数组、长度为 n 加一。变体：快慢指针要求值都能作为合法下标。

\noindent\textbf{LeetCode 318 最大单词长度乘积。} 模型：两个单词没有公共字母时才可相乘，字母集合可压成位掩码。推导重点：每个单词生成 26 位掩码，两个掩码按位与为零表示无交集。

\textbf{二刷读题。} 读题时先用普通状态解释每一项布尔信息。本题可以压成“两个单词没有公共字母时才可相乘，字母集合可压成位掩码”，位运算只是更紧凑的编码。先遮住答案，只保留题面，复述候选对象、合法性条件和输出目标。

\textbf{暴力回放。} 慢版本先用集合、数组或布尔表保存状态，逐步执行题目动作。确认每个布尔字段的含义后，才能把它压缩成位，而不是直接写位运算公式。二刷时先写慢版本或口述慢版本，用它确认不会漏答案。

\textbf{特例回放。} 拿 abcw 和 xtfn 手算，两个单词的 26 位掩码按位与为 0，说明没有公共字符，乘积 4*4 可作为候选。规律是每个单词压成字符集合掩码，重复字母不增加位，掩码相交才排除。把这个特例继续拆开看：每一位或每个掩码位都要对应题目里的一个布尔事实。位运算只是压缩表示；如果两个状态在位置相同但掩码不同，未来能力不同，就不能合并。复盘时把这个特例压成状态字段：掩码含义、当前位、目标位集合、状态转移和答案读取；负数或高位参与时要说明整数宽度；再用边界 重复字母、空字符串、多个单词同掩码、乘积范围 反查初始化、更新顺序和退出条件。

\textbf{状态回放。} 手算路线：手算时把整数写成二进制，并在每一位旁边标注含义。状态转移只允许改变题目动作允许改变的那些位。状态字段：掩码含义、当前位、目标位集合、状态转移和答案读取；负数或高位参与时要说明整数宽度。状态升级：把集合状态压成掩码；包含、加入、删除、枚举子集都变成位操作，但每一位的题目含义不能丢。

\textbf{证明和边界。} 证明从逐位独立或子集归纳开始：被压缩到位里的信息足以决定未来转移。边界：重复字母、空字符串、多个单词同掩码、乘积范围。变体：若字符集大于整数位数，需要 bitset 或哈希集合。

\noindent\textbf{LeetCode 338 比特位计数。} 模型：每个数的一的个数可由去掉最低位一后的较小数转移。推导重点：dp[x] = dp[x \& (x - 1)] + 1，或 dp[x] = dp[x >> 1] + (x \& 1)。

\textbf{二刷读题。} 读题时先用普通状态解释每一项布尔信息。本题可以压成“每个数的一的个数可由去掉最低位一后的较小数转移”，位运算只是更紧凑的编码。先遮住答案，只保留题面，复述候选对象、合法性条件和输出目标。

\textbf{暴力回放。} 慢版本先用集合、数组或布尔表保存状态，逐步执行题目动作。确认每个布尔字段的含义后，才能把它压缩成位，而不是直接写位运算公式。二刷时先写慢版本或口述慢版本，用它确认不会漏答案。

\textbf{特例回放。} 拿 n=5 手算，5 的二进制 101，比 4 多一个最低位 1，所以 bits[5]=bits[4]+1；也可用 5\&(4)=4 得到 bits[5]=bits[4]+1。规律是每个数的答案来自去掉最低位 1 或右移一位后的更小数。把这个特例继续拆开看：每一位或每个掩码位都要对应题目里的一个布尔事实。位运算只是压缩表示；如果两个状态在位置相同但掩码不同，未来能力不同，就不能合并。复盘时把这个特例压成状态字段：掩码含义、当前位、目标位集合、状态转移和答案读取；负数或高位参与时要说明整数宽度；再用边界 n 为零、二的幂、连续区间、表达式优先级 反查初始化、更新顺序和退出条件。

\textbf{状态回放。} 手算路线：手算时把整数写成二进制，并在每一位旁边标注含义。状态转移只允许改变题目动作允许改变的那些位。状态字段：掩码含义、当前位、目标位集合、状态转移和答案读取；负数或高位参与时要说明整数宽度。状态升级：把集合状态压成掩码；包含、加入、删除、枚举子集都变成位操作，但每一位的题目含义不能丢。

\textbf{证明和边界。} 证明从逐位独立或子集归纳开始：被压缩到位里的信息足以决定未来转移。边界：n 为零、二的幂、连续区间、表达式优先级。变体：它展示位运算和 DP 的结合，而不是逐个数循环数位。

\noindent\textbf{LeetCode 371 两整数之和。} 模型：不用加号时，加法可以拆成无进位和与进位。推导重点：异或得到无进位和，与运算后左移得到进位，迭代直到进位为零。

\textbf{二刷读题。} 读题时先用普通状态解释每一项布尔信息。本题可以压成“不用加号时，加法可以拆成无进位和与进位”，位运算只是更紧凑的编码。先遮住答案，只保留题面，复述候选对象、合法性条件和输出目标。

\textbf{暴力回放。} 慢版本先用集合、数组或布尔表保存状态，逐步执行题目动作。确认每个布尔字段的含义后，才能把它压缩成位，而不是直接写位运算公式。二刷时先写慢版本或口述慢版本，用它确认不会漏答案。

\textbf{特例回放。} 拿 5+7 手算，异或得到无进位和 2，按位与左移得到进位 10；继续迭代直到进位为 0。规律是加法拆成无进位和与进位两部分，负数参与时要用足够明确的整数宽度处理。把这个特例继续拆开看：每一位或每个掩码位都要对应题目里的一个布尔事实。位运算只是压缩表示；如果两个状态在位置相同但掩码不同，未来能力不同，就不能合并。复盘时把这个特例压成状态字段：掩码含义、当前位、目标位集合、状态转移和答案读取；负数或高位参与时要说明整数宽度；再用边界 正负数、进位传播、有符号溢出、需要使用无符号中间值 反查初始化、更新顺序和退出条件。

\textbf{状态回放。} 手算路线：手算时把整数写成二进制，并在每一位旁边标注含义。状态转移只允许改变题目动作允许改变的那些位。状态字段：掩码含义、当前位、目标位集合、状态转移和答案读取；负数或高位参与时要说明整数宽度。状态升级：把集合状态压成掩码；包含、加入、删除、枚举子集都变成位操作，但每一位的题目含义不能丢。

\textbf{证明和边界。} 证明从逐位独立或子集归纳开始：被压缩到位里的信息足以决定未来转移。边界：正负数、进位传播、有符号溢出、需要使用无符号中间值。变体：这是位级算术模拟，和异或抵消类题目的语义不同。

\noindent\textbf{LeetCode 393 UTF-8 编码验证。} 模型：每个字节的高位模式决定一个字符需要的后续字节数量。推导重点：用位掩码判断首字节类型，并维护还需要多少个 continuation 字节。

\textbf{二刷读题。} 读题时先用普通状态解释每一项布尔信息。本题可以压成“每个字节的高位模式决定一个字符需要的后续字节数量”，位运算只是更紧凑的编码。先遮住答案，只保留题面，复述候选对象、合法性条件和输出目标。

\textbf{暴力回放。} 慢版本先用集合、数组或布尔表保存状态，逐步执行题目动作。确认每个布尔字段的含义后，才能把它压缩成位，而不是直接写位运算公式。二刷时先写慢版本或口述慢版本，用它确认不会漏答案。

\textbf{特例回放。} 拿 [197,130,1] 手算，197 的高位模式表示还需要 1 个 continuation 字节，130 满足 10xxxxxx，之后 1 是单字节。规律是状态保存剩余后续字节数量，非法首字节或后续字节不足都失败。把这个特例继续拆开看：每一位或每个掩码位都要对应题目里的一个布尔事实。位运算只是压缩表示；如果两个状态在位置相同但掩码不同，未来能力不同，就不能合并。复盘时把这个特例压成状态字段：掩码含义、当前位、目标位集合、状态转移和答案读取；负数或高位参与时要说明整数宽度；再用边界 非法首字节、后续字节不足、后续字节格式错误、单字节字符 反查初始化、更新顺序和退出条件。

\textbf{状态回放。} 手算路线：手算时把整数写成二进制，并在每一位旁边标注含义。状态转移只允许改变题目动作允许改变的那些位。状态字段：掩码含义、当前位、目标位集合、状态转移和答案读取；负数或高位参与时要说明整数宽度。状态升级：把集合状态压成掩码；包含、加入、删除、枚举子集都变成位操作，但每一位的题目含义不能丢。

\textbf{证明和边界。} 证明从逐位独立或子集归纳开始：被压缩到位里的信息足以决定未来转移。边界：非法首字节、后续字节不足、后续字节格式错误、单字节字符。变体：这是位模式解析题，重点是协议规则和状态机。

\noindent\textbf{LeetCode 421 数组中两个数的最大异或值。} 模型：最大异或值由高位到低位尽量取一决定。推导重点：逐位尝试把当前前缀答案设为一，用哈希集合验证是否存在两个前缀配对。

\textbf{二刷读题。} 读题时先用普通状态解释每一项布尔信息。本题可以压成“最大异或值由高位到低位尽量取一决定”，位运算只是更紧凑的编码。先遮住答案，只保留题面，复述候选对象、合法性条件和输出目标。

\textbf{暴力回放。} 慢版本先用集合、数组或布尔表保存状态，逐步执行题目动作。确认每个布尔字段的含义后，才能把它压缩成位，而不是直接写位运算公式。二刷时先写慢版本或口述慢版本，用它确认不会漏答案。

\textbf{特例回放。} 拿 [3,10,5,25,2,8] 手算，25 和 5 的异或为 28。逐位构造答案时，先假设当前位能为 1，再用哈希集合检查是否存在两个前缀配对。规律是从高位到低位贪心验证前缀，不需要枚举所有数对。把这个特例继续拆开看：每一位或每个掩码位都要对应题目里的一个布尔事实。位运算只是压缩表示；如果两个状态在位置相同但掩码不同，未来能力不同，就不能合并。复盘时把这个特例压成状态字段：掩码含义、当前位、目标位集合、状态转移和答案读取；负数或高位参与时要说明整数宽度；再用边界 零、重复数、最高位、负数语义 反查初始化、更新顺序和退出条件。

\textbf{状态回放。} 手算路线：手算时把整数写成二进制，并在每一位旁边标注含义。状态转移只允许改变题目动作允许改变的那些位。状态字段：掩码含义、当前位、目标位集合、状态转移和答案读取；负数或高位参与时要说明整数宽度。状态升级：把集合状态压成掩码；包含、加入、删除、枚举子集都变成位操作，但每一位的题目含义不能丢。

\textbf{证明和边界。} 证明从逐位独立或子集归纳开始：被压缩到位里的信息足以决定未来转移。边界：零、重复数、最高位、负数语义。变体：也可以建二进制 Trie，每个数尽量走相反位。

\subsection{自测标准}

二刷时不要只确认自己见过这道题。请遮住答案，先讲题面模型，再讲暴力基线和瓶颈，然后说出优化结构保存了什么、不变量如何排除错误候选、边界实验如何打破错误实现。

\begin{principlebox}{二刷标准}
二刷不是重新看一遍答案，而是把题目变成可讲、可证、可调试、可迁移的知识块。每道题至少回答：暴力瓶颈、优化依据、不变量、复杂度、边界、C++ 成本和一个变体。
\end{principlebox}

\section{逐题单点推导手册}

这一节专门解决“每一题思路都要讲清楚”的问题。它不重复完整代码，而是要求每道题都从自己的题面出发，说明暴力基线、优化线索、状态不变量、边界和变体。复习时可以把这一节当作口述题解提纲。

\subsubsection{LeetCode 1 两数之和}

\textbf{题面模型。} 当前元素只需要寻找唯一补数，历史值到下标的映射就是最小状态。

\textbf{暴力瓶颈。} 先查再插入，避免同一个下标被用两次；若数组有序才考虑双指针。

\textbf{边界反例。} 重复数字、目标为偶数且补数等于当前值、没有答案时返回空结果。

\textbf{变体迁移。} 若改成返回所有不重复数对，要先排序或用频次表处理重复。

\subsubsection{LeetCode 2 两数相加}

\textbf{题面模型。} 链表按低位到高位保存数字，本质是竖式加法而不是整数转换。

\textbf{暴力瓶颈。} 维护两个游标和进位，任一链未结束或进位非零都继续生成节点。

\textbf{边界反例。} 两链长度不同、最后进位、输入为零、不能用内置整数承载超长数字。

\textbf{变体迁移。} 若链表高位在前，可以用栈反向处理或先反转链表。

\subsubsection{LeetCode 3 无重复字符的最长子串}

\textbf{题面模型。} 窗口保存当前无重复子串，右端每次加入一个字符。

\textbf{暴力瓶颈。} 用上次出现位置直接跳过无效左端，左边界只能向右不能回退。

\textbf{边界反例。} 空串、全相同字符、重复字符出现在窗口左侧之外、扩展字符集。

\textbf{变体迁移。} 若要求恰好包含 k 种字符，窗口条件从无重复变成种类计数。

\subsubsection{LeetCode 5 最长回文子串}

\textbf{题面模型。} 回文围绕中心对称，中心可以是字符也可以是两个字符之间。

\textbf{暴力瓶颈。} 中心扩展避免枚举子串后再逐个检查；每次扩展只比较新增左右边界。

\textbf{边界反例。} 偶数长度回文、单字符、全相同字符、多个同长答案。

\textbf{变体迁移。} 若要求回文子序列，应转成区间 DP，不能用中心扩展。

\subsubsection{LeetCode 11 盛最多水的容器}

\textbf{题面模型。} 左右指针代表当前选取的两条边，面积由短板和宽度共同决定。

\textbf{暴力瓶颈。} 每次移动短板，因为移动长板只会缩小宽度且不会提高短板。

\textbf{边界反例。} 两端高度相等、数组长度为二、面积乘法可能溢出。

\textbf{变体迁移。} 接雨水计算每个位置贡献，不能把本题双指针证明直接搬过去。

\subsubsection{LeetCode 15 三数之和}

\textbf{题面模型。} 排序后固定第一个数，剩余两数转成有序双指针。

\textbf{暴力瓶颈。} 和太小移动左指针，和太大移动右指针；固定值和左右值都要跳过重复。

\textbf{边界反例。} 全零、重复元素很多、没有答案、结果顺序不要求但组合不能重复。

\textbf{变体迁移。} 四数之和再外层固定一个数，并用 long long 防止目标和溢出。

\subsubsection{LeetCode 16 最接近的三数之和}

\textbf{题面模型。} 排序后固定一数，双指针寻找最接近目标的两数和。

\textbf{暴力瓶颈。} 每次根据当前和与目标的大小移动指针，同时更新最小绝对差。

\textbf{边界反例。} 差值相等、负数、目标很大、三数和溢出。

\textbf{变体迁移。} 若要求返回所有最接近组合，需要记录同差值并去重。

\subsubsection{LeetCode 18 四数之和}

\textbf{题面模型。} 两层固定加一层双指针，关键是剪枝和去重。

\textbf{暴力瓶颈。} 排序提供单调性，外层可用最小可能和和最大可能和提前跳过。

\textbf{边界反例。} 重复元素、目标超出 int、答案很多导致输出空间主导复杂度。

\textbf{变体迁移。} k 数之和可以递归降维，但工程实现要控制重复和边界。

\subsubsection{LeetCode 26 删除有序数组重复项}

\textbf{题面模型。} 有序数组让重复元素相邻，慢指针表示下一个唯一值写入位置。

\textbf{暴力瓶颈。} 快指针扫描，只有发现新值才写入慢指针并推进。

\textbf{边界反例。} 空数组、全重复、无重复、返回长度不等于物理容量。

\textbf{变体迁移。} 若允许每个元素最多保留两次，只需改变写入条件为和前两个比较。

\subsubsection{LeetCode 27 移除元素}

\textbf{题面模型。} 原地过滤，慢指针左侧始终是不等于目标值的稳定前缀。

\textbf{暴力瓶颈。} 保持顺序用快慢指针；不保持顺序可用左右交换减少写入。

\textbf{边界反例。} 所有元素都被移除、没有元素被移除、目标值在尾部密集。

\textbf{变体迁移。} 工程里要先确认是否要求稳定顺序，不同要求对应不同写法。

\subsubsection{LeetCode 31 下一个排列}

\textbf{题面模型。} 字典序结构由最长非递增后缀决定，枢轴是后缀前一个位置。

\textbf{暴力瓶颈。} 从右找第一个上升点，再从右找刚好更大的元素交换，最后反转后缀。

\textbf{边界反例。} 完全降序、重复数字、只有一个元素、交换后后缀必须最小化。

\textbf{变体迁移。} 若要求上一个排列，比较方向和后缀处理对称变化。

\subsubsection{LeetCode 33 搜索旋转排序数组}

\textbf{题面模型。} 旋转数组每次二分至少有一半保持有序。

\textbf{暴力瓶颈。} 判断哪一半有序，再看目标是否落在有序半边范围内。

\textbf{边界反例。} 未旋转、长度一、目标不存在、边界等号、存在重复元素时退化。

\textbf{变体迁移。} 若重复元素很多，需要在无法判断有序半边时收缩边界。

\subsubsection{LeetCode 34 排序数组中查找首尾位置}

\textbf{题面模型。} 首尾位置可以拆成两个边界：第一个大于等于和第一个大于。

\textbf{暴力瓶颈。} 不要找到一个目标后向两边线性扩展，否则重复元素很多时退化。

\textbf{边界反例。} 目标不存在、目标在开头或结尾、数组为空、全是目标。

\textbf{变体迁移。} 这个模型能迁移到计数某值出现次数和插入区间边界。

\subsubsection{LeetCode 35 搜索插入位置}

\textbf{题面模型。} 题目等价于寻找第一个大于等于目标的位置。

\textbf{暴力瓶颈。} 统一用 lower bound 语义，不在相等时提前返回也能保持边界一致。

\textbf{边界反例。} 目标小于所有元素、目标大于所有元素、重复值、空数组。

\textbf{变体迁移。} 手写二分时用左闭右开可直接返回 left。

\subsubsection{LeetCode 42 接雨水}

\textbf{题面模型。} 每个位置的水由左侧最高和右侧最高的较小者决定。

\textbf{暴力瓶颈。} 前后缀、双指针、单调栈都是不同状态维护方式；要能说明各自结算对象。

\textbf{边界反例。} 单调递增、单调递减、平台高度、多个凹槽、数组长度不足三。

\textbf{变体迁移。} 若柱子变成二维网格，需要最小堆从边界向内扩展。

\subsubsection{LeetCode 49 字母异位词分组}

\textbf{题面模型。} 异位词共享同一个规范表示，可以是排序字符串或字符计数。

\textbf{暴力瓶颈。} 哈希表按规范键分桶，避免两两比较字符串。

\textbf{边界反例。} 空字符串、重复单词、字符集不止小写字母、计数键必须无歧义。

\textbf{变体迁移。} 若字符串很长且字符集固定，计数键通常比排序键更稳定。

\subsubsection{LeetCode 53 最大子数组和}

\textbf{题面模型。} 状态是必须以当前位置结尾的最大连续和。

\textbf{暴力瓶颈。} 若前一个后缀为负，接上它只会拖累当前元素，因此可以重新开始。

\textbf{边界反例。} 全负数组、单元素、和溢出、答案不能初始化为零。

\textbf{变体迁移。} 若要求返回区间下标，同步记录重新开始位置和最优左右端。

\subsubsection{LeetCode 54 螺旋矩阵}

\textbf{题面模型。} 四个边界收缩模拟一圈一圈访问。

\textbf{暴力瓶颈。} 每走完一条边就收缩对应边界，并在走反方向前检查是否仍合法。

\textbf{边界反例。} 单行、单列、空矩阵、最后一圈只剩一个元素。

\textbf{变体迁移。} 若改成生成矩阵，边界逻辑相同，只是读变成写。

\subsubsection{LeetCode 55 跳跃游戏}

\textbf{题面模型。} 扫描过程中只关心当前能到达的最远位置。

\textbf{暴力瓶颈。} 如果下标超过最远可达，任何之前路径都不能到达它；否则用它扩展边界。

\textbf{边界反例。} 第一个元素为零、数组长度一、恰好到达最后、远超最后。

\textbf{变体迁移。} 求最少跳数时把可达区间看成 BFS 层。

\subsubsection{LeetCode 56 合并区间}

\textbf{题面模型。} 排序后可能与当前区间相交的只会是结果尾部。

\textbf{暴力瓶颈。} 按起点排序，扫描时维护结果最后区间的右端点。

\textbf{边界反例。} 端点相接、完全覆盖、无重叠、空列表、输入未排序。

\textbf{变体迁移。} 若区间是半开区间，端点相等不一定合并。

\subsubsection{LeetCode 57 插入区间}

\textbf{题面模型。} 新区间只会影响与它相交的一段连续区间。

\textbf{暴力瓶颈。} 先追加左侧完全不相交区间，再合并重叠段，最后追加右侧。

\textbf{边界反例。} 新区间在最前、最后、覆盖全部、刚好相邻、原列表为空。

\textbf{变体迁移。} 若输入不是有序且不重叠，必须先排序或换模型。

\subsubsection{LeetCode 62 不同路径}

\textbf{题面模型。} 网格状态由上方和左方两种最后一步转移而来。

\textbf{暴力瓶颈。} 二维表可压缩为一维，因为当前行只依赖上一行同列和当前行左侧。

\textbf{边界反例。} 一行、一列、较大网格结果溢出、障碍物版本边界。

\textbf{变体迁移。} 带障碍物时遇到障碍状态清零，第一行第一列要按阻断处理。

\subsubsection{LeetCode 63 不同路径 II}

\textbf{题面模型。} 障碍物让某些格子的路径数变为零。

\textbf{暴力瓶颈。} 滚动数组中遇到障碍直接把当前位置清零，否则累加左侧。

\textbf{边界反例。} 起点或终点是障碍、整行被截断、单行障碍。

\textbf{变体迁移。} 若允许向上或向左走，图会有环，普通填表不再成立。

\subsubsection{LeetCode 64 最小路径和}

\textbf{题面模型。} 状态是到达当前格子的最小代价，最后一步来自上方或左方。

\textbf{暴力瓶颈。} 先初始化第一行和第一列，再处理内部，主转移保持简单。

\textbf{边界反例。} 单行、单列、权值为零、路径和溢出。

\textbf{变体迁移。} 若允许四方向移动，变成最短路而不是普通网格 DP。

\subsubsection{LeetCode 70 爬楼梯}

\textbf{题面模型。} 最后一步来自前一阶或前两阶，是最小的重叠子问题模型。

\textbf{暴力瓶颈。} 滚动变量保存最近两个状态即可，不需要完整数组。

\textbf{边界反例。} n 为一、n 为二、题目是否允许零阶、结果范围。

\textbf{变体迁移。} 若每次可走一到 k 阶，转移变成固定宽度窗口求和。

\subsubsection{LeetCode 72 编辑距离}

\textbf{题面模型。} 二维状态表示两个前缀之间的最少编辑操作。

\textbf{暴力瓶颈。} 末尾字符相同继承左上；不同则枚举插入、删除、替换。

\textbf{边界反例。} 空串、完全相同、完全不同、操作代价不同。

\textbf{变体迁移。} 若只允许插入删除，问题退化到和最长公共子序列相关。

\subsubsection{LeetCode 75 颜色分类}

\textbf{题面模型。} 三指针维护零区、一号区、未知区和二区。

\textbf{暴力瓶颈。} 遇到二时只移动右边界，不移动扫描指针，因为换回来的元素未知。

\textbf{边界反例。} 全零、全二、已经有序、逆序、交换后漏检查。

\textbf{变体迁移。} 这是荷兰国旗问题，能推广到三类分区的原地维护。

\subsubsection{LeetCode 76 最小覆盖子串}

\textbf{题面模型。} 窗口内字符计数必须覆盖目标计数，满足后尽量收缩左端。

\textbf{暴力瓶颈。} 用 formed 记录满足需求的字符种类，避免每次全量比较。

\textbf{边界反例。} 目标有重复字符、源串缺字符、最小窗口在开头或结尾、大小写。

\textbf{变体迁移。} 若字符集固定，数组计数更快；若是 Unicode，哈希表更通用。

\subsubsection{LeetCode 78 子集}

\textbf{题面模型。} 每个元素都有选或不选两种选择，答案是所有路径节点。

\textbf{暴力瓶颈。} 回溯按起点向后枚举，避免不同顺序生成同一集合。

\textbf{边界反例。} 空数组、单元素、输出数量为二的 n 次方、顺序不要求。

\textbf{变体迁移。} 也可用位掩码枚举所有子集，适合 n 较小的场景。

\subsubsection{LeetCode 79 单词搜索}

\textbf{题面模型。} 网格路径不能重复使用同一格，状态包含位置、匹配下标和访问标记。

\textbf{暴力瓶颈。} 剪枝包括首字符不匹配、剩余长度不足、字符频次不足和越界访问。

\textbf{边界反例。} 单字符单词、路径需要回退、重复字母、单元格不能复用。

\textbf{变体迁移。} 若要查很多单词，应使用 Trie 合并搜索前缀。

\subsubsection{LeetCode 84 柱状图中最大的矩形}

\textbf{题面模型。} 每根柱子作为最矮高度时，需要知道左右第一个更矮位置。

\textbf{暴力瓶颈。} 单调递增栈在遇到更矮柱时结算被弹出柱子的最大宽度。

\textbf{边界反例。} 递增、递减、相等高度、哨兵零、宽度计算。

\textbf{变体迁移。} 最大矩形可以作为二维矩阵最大矩形的行压缩子问题。

\subsubsection{LeetCode 88 合并两个有序数组}

\textbf{题面模型。} 第一个数组尾部有空位，适合从后往前写入最终最大值。

\textbf{暴力瓶颈。} 逆向填充避免覆盖尚未读取的 nums1 元素。

\textbf{边界反例。} 第二个数组为空、第一个有效元素为空、重复值、尾部剩余。

\textbf{变体迁移。} 若没有额外空位，只能新开数组或使用更复杂原地归并。

\subsubsection{LeetCode 90 子集 II}

\textbf{题面模型。} 重复元素会让不同下标路径生成相同集合。

\textbf{暴力瓶颈。} 排序后在同一层跳过相同值，保留不同层使用重复值的可能。

\textbf{边界反例。} 全重复、空集、重复值相邻、去重条件写成同层而不是全局。

\textbf{变体迁移。} 组合总和 II 使用同样的同层去重思想。

\subsubsection{LeetCode 94 二叉树中序遍历}

\textbf{题面模型。} 中序访问顺序是左子树、根、右子树，BST 中会得到有序序列。

\textbf{暴力瓶颈。} 用显式栈模拟一路向左，再回退访问根并转向右子树。

\textbf{边界反例。} 空树、单节点、链式左树、链式右树。

\textbf{变体迁移。} Morris 遍历可做到常数额外空间，但会临时改树结构。

\subsubsection{LeetCode 98 验证二叉搜索树}

\textbf{题面模型。} BST 约束是全局上下界，不是只比较父子节点。

\textbf{暴力瓶颈。} 中序严格递增或显式传递上下界都可以；中序版本要保存前一个访问值。

\textbf{边界反例。} 重复值、int 最小最大值、局部合法但全局非法。

\textbf{变体迁移。} 若允许重复值，需要明确重复放左还是放右。

\subsubsection{LeetCode 102 二叉树层序遍历}

\textbf{题面模型。} 队列在每轮开始时保存当前层所有节点。

\textbf{暴力瓶颈。} 记录 level size 固定层边界，避免下一层节点混入当前层。

\textbf{边界反例。} 空树、只有根、最后一层不满、左右顺序。

\textbf{变体迁移。} 锯齿形层序只是在每层输出顺序上增加方向控制。

\subsubsection{LeetCode 103 二叉树锯齿形层序遍历}

\textbf{题面模型。} BFS 层边界不变，变化的是当前层写入结果的方向。

\textbf{暴力瓶颈。} 可以层内反转，也可以预分配数组按方向写入目标下标。

\textbf{边界反例。} 单层、偶数层、奇数层、空树。

\textbf{变体迁移。} 若要求从底向上输出，收集后反转层列表即可。

\subsubsection{LeetCode 104 二叉树最大深度}

\textbf{题面模型。} 深度可以看成层数，也可以作为栈中携带的状态。

\textbf{暴力瓶颈。} 显式栈保存节点和深度，避免极端链式树递归栈过深。

\textbf{边界反例。} 空树返回零、根深度是一、链式树。

\textbf{变体迁移。} 最小深度不能简单取左右最小，因为空子树不构成叶子路径。

\subsubsection{LeetCode 105 前序中序构造二叉树}

\textbf{题面模型。} 前序给根，中序把左右子树切开。

\textbf{暴力瓶颈。} 用哈希表记录中序下标，避免每次线性寻找根位置。

\textbf{边界反例。} 节点值唯一、空子树、左右子树长度计算、输入不合法。

\textbf{变体迁移。} 后序中序构造类似，只是根来自后序末尾。

\subsubsection{LeetCode 106 中序后序构造二叉树}

\textbf{题面模型。} 后序末尾是根，中序位置确定左右子树规模。

\textbf{暴力瓶颈。} 构造时要小心后序左右区间边界，避免切分错位。

\textbf{边界反例。} 空树、单节点、全左链、全右链。

\textbf{变体迁移。} 若给前序和后序，二叉树通常不能唯一确定，除非额外限制。

\subsubsection{LeetCode 108 有序数组转二叉搜索树}

\textbf{题面模型。} 有序数组中点作为根能让左右规模接近。

\textbf{暴力瓶颈。} 每次选择区间中点，左右子区间构造左右子树。

\textbf{边界反例。} 空区间、偶数长度选左中还是右中、高度平衡定义。

\textbf{变体迁移。} 工程上递归可读，若严格避免递归可用队列保存区间和父指针。

\subsubsection{LeetCode 110 平衡二叉树}

\textbf{题面模型。} 每个节点需要知道子树高度以及是否已经失衡。

\textbf{暴力瓶颈。} 后序汇总时一旦子树失衡就向上返回失败标记，避免重复算高度。

\textbf{边界反例。} 空树、单节点、左右高度差正好一、链式树。

\textbf{变体迁移。} 可以把返回值设计成 optional 高度，空表示不平衡。

\subsubsection{LeetCode 111 二叉树最小深度}

\textbf{题面模型。} 最小深度是到最近叶子的节点数，不是左右深度简单取最小。

\textbf{暴力瓶颈。} BFS 第一次遇到叶子就是最小深度，因为按层扩展。

\textbf{边界反例。} 只有左子树、只有右子树、根为空、根是叶子。

\textbf{变体迁移。} DFS 写法必须特殊处理空子树不能作为最小路径。

\subsubsection{LeetCode 112 路径总和}

\textbf{题面模型。} 路径必须从根到叶子，状态是剩余目标值。

\textbf{暴力瓶颈。} 沿路径减去当前节点值，到叶子时检查剩余是否等于叶子值。

\textbf{边界反例。} 负数节点、目标为零、非叶子提前达到目标、空树。

\textbf{变体迁移。} 若要求列出所有路径，需要保存路径数组并在回退时弹出。

\subsubsection{LeetCode 113 路径总和 II}

\textbf{题面模型。} 相比判定题，状态还要保存当前路径。

\textbf{暴力瓶颈。} 到叶子且和满足时复制路径；回溯时恢复路径。

\textbf{边界反例。} 多个答案、负数、空树、路径数组复制时机。

\textbf{变体迁移。} 若路径不要求从根到叶，模型会变成前缀和加哈希。

\subsubsection{LeetCode 114 二叉树展开为链表}

\textbf{题面模型。} 目标是按前序顺序把树原地改成右指针链。

\textbf{暴力瓶颈。} 显式栈按右后左先入顺序处理，逐个接到前驱右侧。

\textbf{边界反例。} 空树、单节点、原有右子树不能丢、左指针置空。

\textbf{变体迁移。} 也可用寻找左子树最右节点的原地方法。

\subsubsection{LeetCode 121 买卖股票的最佳时机}

\textbf{题面模型。} 卖出日固定时，最佳买入日是此前最低价格。

\textbf{暴力瓶颈。} 扫描维护历史最低价和最大利润，一次交易无需复杂状态机。

\textbf{边界反例。} 单日价格、一直下跌、一直上涨、价格相同。

\textbf{变体迁移。} 多次交易、冷冻期和手续费版本都要扩展成状态机。

\subsubsection{LeetCode 122 买卖股票 II}

\textbf{题面模型。} 允许多次交易且不能同时持有，所有正收益差都可以收集。

\textbf{暴力瓶颈。} 把长上涨段拆成每天买卖收益相同，因此累加正差值最优。

\textbf{边界反例。} 下降段、平台价格、只有一天、手续费不存在。

\textbf{变体迁移。} 有手续费时不能简单累加正差，需要状态机或调整贪心阈值。

\subsubsection{LeetCode 123 买卖股票 III}

\textbf{题面模型。} 最多两次交易，状态需要包含交易次数和持有状态。

\textbf{暴力瓶颈。} 维护 buy1、sell1、buy2、sell2 四个状态，按天更新。

\textbf{边界反例。} 价格为空、只够一次交易、两次交易间不能重叠。

\textbf{变体迁移。} 最多 k 次交易是它的泛化，k 很大时退化为无限次交易。

\subsubsection{LeetCode 124 二叉树最大路径和}

\textbf{题面模型。} 路径可以从任意节点到任意节点，但不能分叉向父节点贡献两边。

\textbf{暴力瓶颈。} 每个节点向父亲贡献单边最大增益，全局答案可用左右两边加当前节点更新。

\textbf{边界反例。} 全负数、单节点、左右增益为负时取零、答案初始值。

\textbf{变体迁移。} 若路径必须从根到叶，状态和答案位置完全不同。

\subsubsection{LeetCode 127 单词接龙}

\textbf{题面模型。} 每个单词是节点，一次修改一个字符形成边。

\textbf{暴力瓶颈。} 通配模式桶加速邻居查询，BFS 第一次到达终点就是最短转换长度。

\textbf{边界反例。} 终点不在词表、起点和终点相同、桶重复扫描、单词长度一致。

\textbf{变体迁移。} 双向 BFS 可以进一步降低搜索空间。

\subsubsection{LeetCode 128 最长连续序列}

\textbf{题面模型。} 连续段只需要从没有前驱的数开始扩展。

\textbf{暴力瓶颈。} 哈希集合判断 x 减一 是否存在，跳过非起点避免重复扫描。

\textbf{边界反例。} 重复元素、负数、空数组、整数边界。

\textbf{变体迁移。} 若数据流动态插入，可以用并查集或区间合并维护长度。

\subsubsection{LeetCode 130 被围绕的区域}

\textbf{题面模型。} 只有不连通到边界的 O 才会被翻转。

\textbf{暴力瓶颈。} 从边界 O 出发标记安全区域，再翻转剩余 O。

\textbf{边界反例。} 全边界、无 O、单行单列、内部岛屿。

\textbf{变体迁移。} 这是反向思维：找不能被翻转的区域比直接判断每块区域更简单。

\subsubsection{LeetCode 131 分割回文串}

\textbf{题面模型。} 每次选择一个从当前位置开始的回文前缀。

\textbf{暴力瓶颈。} 预处理回文 DP 表，让回溯合法性检查为常数。

\textbf{边界反例。} 空串、单字符、全相同字符、重复分割路径。

\textbf{变体迁移。} 若只求最少切割次数，可改成 DP 最短切分。

\subsubsection{LeetCode 133 克隆图}

\textbf{题面模型。} 每个原节点需要唯一对应一个新节点，边按原图连接。

\textbf{暴力瓶颈。} 哈希表保存原节点到克隆节点的映射，BFS 或 DFS 遍历图。

\textbf{边界反例。} 空图、自环、重边、环结构、不能按节点值唯一假设。

\textbf{变体迁移。} 若图节点值不唯一，必须用地址作为键。

\subsubsection{LeetCode 136 只出现一次的数字}

\textbf{题面模型。} 成对元素可以用异或抵消，剩下单个元素。

\textbf{暴力瓶颈。} 异或满足交换结合，扫描顺序不影响答案。

\textbf{边界反例。} 零、负数、唯一值在开头或结尾、所有其他值恰好两次。

\textbf{变体迁移。} 若其他值出现三次，要改为按位计数取模。

\subsubsection{LeetCode 139 单词拆分}

\textbf{题面模型。} 状态表示前缀能否由字典单词拼出。

\textbf{暴力瓶颈。} 枚举最后一个单词的起点，左侧前缀可达且子串在字典中则当前可达。

\textbf{边界反例。} 空串、重复单词、长单词、substr 复制成本。

\textbf{变体迁移。} 若要输出所有句子，需要 DFS 加记忆化或前驱列表。

\subsubsection{LeetCode 141 环形链表}

\textbf{题面模型。} 快慢指针在有环时必然相遇，无环时快指针先到空。

\textbf{暴力瓶颈。} 不用哈希表也能常数空间检测环。

\textbf{边界反例。} 空链、单节点自环、两节点环、尾部无环。

\textbf{变体迁移。} 若要求环入口，用相遇后双指针同步走。

\subsubsection{LeetCode 142 环形链表 II}

\textbf{题面模型。} 相遇点到入口的距离关系来自快慢指针速度差。

\textbf{暴力瓶颈。} 相遇后一个指针回头，两个指针每次走一步，再次相遇就是入口。

\textbf{边界反例。} 入口在头节点、长尾短环、无环、单节点环。

\textbf{变体迁移。} 也可用哈希集合，但空间更高。

\subsubsection{LeetCode 146 LRU 缓存}

\textbf{题面模型。} 哈希表负责按 key 定位，双向链表负责维护新旧顺序。

\textbf{暴力瓶颈。} 访问和更新都把节点移动到头部，容量满时删除尾部最旧节点。

\textbf{边界反例。} 容量为一、重复 put、get 不存在、更新已有值。

\textbf{变体迁移。} C++ 可用 \code{std::list} 加迭代器，也可手写双向链表理解所有权。

\subsubsection{LeetCode 148 排序链表}

\textbf{题面模型。} 链表不能随机访问，适合归并排序而不是快速排序数组 partition。

\textbf{暴力瓶颈。} 自底向上迭代归并能避免递归栈，同时保持 O(n log n)。

\textbf{边界反例。} 空链、单节点、重复值、奇数长度、切断子链。

\textbf{变体迁移。} 若把值复制到数组排序，会改变节点身份语义。

\subsubsection{LeetCode 150 逆波兰表达式求值}

\textbf{题面模型。} 后缀表达式把优先级和括号都编码进 token 顺序。

\textbf{暴力瓶颈。} 遇到数字入栈，遇到运算符弹出右操作数和左操作数计算。

\textbf{边界反例。} 负数 token、多位数、除法向零截断、减法顺序。

\textbf{变体迁移。} 中缀表达式求值需要另一个运算符栈处理优先级。

\subsubsection{LeetCode 152 乘积最大子数组}

\textbf{题面模型。} 负数会让最大乘积和最小乘积互换。

\textbf{暴力瓶颈。} 同时维护以当前位置结尾的最大和最小乘积。

\textbf{边界反例。} 零、负数个数奇偶、单元素、乘积溢出。

\textbf{变体迁移。} 如果要求返回区间，同步记录状态来源。

\subsubsection{LeetCode 153 寻找旋转排序数组中的最小值}

\textbf{题面模型。} 最小值是旋转断点，右端点可帮助判断中点在哪一段。

\textbf{暴力瓶颈。} 若 mid 大于 right，最小值在右侧；否则在左侧含 mid。

\textbf{边界反例。} 未旋转、长度一、旋转一位、无重复条件。

\textbf{变体迁移。} 有重复时遇到相等只能收缩 right，最坏退化。

\subsubsection{LeetCode 154 寻找旋转排序数组中的最小值 II}

\textbf{题面模型。} 重复元素会让 mid 和 right 相等时无法判断最小值方向。

\textbf{暴力瓶颈。} 相等时只能安全地丢掉一个 right，因为它不是唯一必要候选。

\textbf{边界反例。} 全相等、重复值跨断点、最小值出现多次。

\textbf{变体迁移。} 这题说明二分依赖的信息不足时会退化，不是所有旋转题都严格对数。

\subsubsection{LeetCode 155 最小栈}

\textbf{题面模型。} 普通栈只能看到栈顶，不能常数时间知道历史最小值。

\textbf{暴力瓶颈。} 辅助栈同步保存每个深度对应的最小值。

\textbf{边界反例。} 重复最小值、弹出最小值、空栈操作约定。

\textbf{变体迁移。} 也可保存差值压缩空间，但可读性和溢出处理更复杂。

\subsubsection{LeetCode 160 相交链表}

\textbf{题面模型。} 相交是节点地址相同，不是节点值相同。

\textbf{暴力瓶颈。} 两个指针走完各自链表后切换到另一条链，走过总长度相同后相遇。

\textbf{边界反例。} 不相交、头节点相交、尾节点相交、长度差很大。

\textbf{变体迁移。} 哈希集合更直观但空间更高。

\subsubsection{LeetCode 162 寻找峰值}

\textbf{题面模型。} 比较 mid 和 mid 加一 的斜率能判断某侧一定存在峰值。

\textbf{暴力瓶颈。} 上升则右侧有峰，下降则左侧含 mid 有峰。

\textbf{边界反例。} 长度一、峰在边界、多个峰、相邻不相等假设。

\textbf{变体迁移。} 二维峰值需要按行列最大值进行更复杂二分。

\subsubsection{LeetCode 169 多数元素}

\textbf{题面模型。} 超过一半的元素可以和其他元素两两抵消后仍然剩下。

\textbf{暴力瓶颈。} Boyer-Moore 投票维护候选和计数。

\textbf{边界反例。} 题目是否保证存在多数、计数归零、全相同。

\textbf{变体迁移。} 若找超过三分之一的元素，需要维护两个候选。

\subsubsection{LeetCode 189 轮转数组}

\textbf{题面模型。} 右移 k 位等价于把数组切成两段后换序。

\textbf{暴力瓶颈。} 三次反转能原地完成：整体、前 k、后 n 减 k。

\textbf{边界反例。} k 为零、k 大于 n、空数组、单元素。

\textbf{变体迁移。} 环状替换也可做，但要处理最大公约数个循环。

\subsubsection{LeetCode 198 打家劫舍}

\textbf{题面模型。} 每间房只有偷和不偷两种结局，偷当前就不能偷前一间。

\textbf{暴力瓶颈。} 滚动保存前一状态和前两状态。

\textbf{边界反例。} 空数组、单间、全零、金额很大。

\textbf{变体迁移。} 环形房屋要拆成不含首和不含尾两个线性问题。

\subsubsection{LeetCode 199 二叉树右视图}

\textbf{题面模型。} 每层从右侧能看到的就是该层最后访问或最先访问的右侧节点。

\textbf{暴力瓶颈。} BFS 固定层边界，取每层最后一个；或 DFS 按右优先首次到达深度。

\textbf{边界反例。} 空树、左链、右链、左右交错。

\textbf{变体迁移。} 若求左视图，只改变层内取值方向。

\subsubsection{LeetCode 200 岛屿数量}

\textbf{题面模型。} 陆地格子是节点，四方向相邻是边，岛屿是连通块。

\textbf{暴力瓶颈。} 扫描遇到未访问陆地就启动搜索并标记整块。

\textbf{边界反例。} 空网格、全水、全陆、斜对角不连通、是否允许修改输入。

\textbf{变体迁移。} 也可用并查集合并相邻陆地，适合动态岛屿扩展。

\subsubsection{LeetCode 202 快乐数}

\textbf{题面模型。} 反复变换数字会进入一或循环，状态空间最终有限。

\textbf{暴力瓶颈。} 用集合检测重复状态，或用快慢指针看成函数图。

\textbf{边界反例。} 输入为一、进入循环、位平方和函数、负数不在题意内。

\textbf{变体迁移。} 快慢指针版本能训练把数值迭代看成链表。

\subsubsection{LeetCode 206 反转链表}

\textbf{题面模型。} 每次把当前节点的 next 指向已经反转好的前缀头。

\textbf{暴力瓶颈。} 先保存后继，再改指针，再推进 previous 和 current。

\textbf{边界反例。} 空链、单节点、两个节点、不要丢失剩余链。

\textbf{变体迁移。} K 组反转是在这段局部反转外增加分组边界。

\subsubsection{LeetCode 207 课程表}

\textbf{题面模型。} 课程是节点，先修关系是有向边，问题是判断是否有环。

\textbf{暴力瓶颈。} 入度拓扑排序不断学习当前无依赖课程。

\textbf{边界反例。} 边方向、孤立课程、环、自依赖、重复边。

\textbf{变体迁移。} DFS 三色标记也能判环，但要明确访问状态。

\subsubsection{LeetCode 208 实现 Trie}

\textbf{题面模型。} 前缀树把字符串公共前缀共享成路径。

\textbf{暴力瓶颈。} 插入和查询都逐字符沿边移动，不存在的边按需要创建。

\textbf{边界反例。} 空字符串约定、单词结束标记、前缀存在不等于单词存在。

\textbf{变体迁移。} 若字符集很大，用哈希子节点；小写字母用数组更快。

\subsubsection{LeetCode 209 长度最小的子数组}

\textbf{题面模型。} 正数数组让窗口和随右端增加、随左端减少。

\textbf{暴力瓶颈。} 右端扩展到满足目标后，不断左移缩短并更新答案。

\textbf{边界反例。} 不存在答案、单元素达到目标、全正是必要条件。

\textbf{变体迁移。} 若含负数，需要前缀和加单调队列。

\subsubsection{LeetCode 210 课程表 II}

\textbf{题面模型。} 在判环基础上还要输出一个可行拓扑序。

\textbf{暴力瓶颈。} 每弹出一个入度为零课程，就把它加入顺序并删除出边。

\textbf{边界反例。} 有环返回空数组、多个可行顺序、边方向。

\textbf{变体迁移。} 若需要字典序最小拓扑序，把队列换成小顶堆。

\subsubsection{LeetCode 211 添加与搜索单词}

\textbf{题面模型。} 点号通配符让查询可能分叉，但前缀树仍能剪掉不匹配前缀。

\textbf{暴力瓶颈。} 普通字符沿唯一边走，点号枚举当前节点所有孩子。

\textbf{边界反例。} 连续通配符、空结果、单词结束标记、字符集大小。

\textbf{变体迁移。} 大量通配符会接近遍历整棵 Trie，需要规模约束。

\subsubsection{LeetCode 215 数组中的第 K 个最大元素}

\textbf{题面模型。} 只关心第 K 大，不需要完整排序。

\textbf{暴力瓶颈。} 大小为 K 的小顶堆保存当前前 K 大边界；快速选择可做到平均线性。

\textbf{边界反例。} K 为一、K 为 n、重复值、随机 pivot。

\textbf{变体迁移。} 若数据流持续到来，堆方案比一次性快速选择更自然。

\subsubsection{LeetCode 216 组合总和 III}

\textbf{题面模型。} 从一到九中选 k 个数和为 n，每个数最多用一次。

\textbf{暴力瓶颈。} 按递增起点枚举，利用剩余个数和剩余和剪枝。

\textbf{边界反例。} n 太小、n 太大、k 为零、路径长度已经超过 k。

\textbf{变体迁移。} 可预估剩余最小和最大和进一步剪枝。

\subsubsection{LeetCode 217 存在重复元素}

\textbf{题面模型。} 判断是否出现过是哈希集合最直接的职责。

\textbf{暴力瓶颈。} 扫描时若插入前已经存在则立即返回真。

\textbf{边界反例。} 空数组、全唯一、重复在末尾、值域很小。

\textbf{变体迁移。} 若允许修改输入，排序后看相邻可以降低额外空间。

\subsubsection{LeetCode 221 最大正方形}

\textbf{题面模型。} 以当前位置为右下角的最大正方形由上、左、左上三个状态限制。

\textbf{暴力瓶颈。} 当前位置为一时取三者最小加一，否则为零。

\textbf{边界反例。} 全零、全一、单行单列、字符一和数字一不要混淆。

\textbf{变体迁移。} 最大矩形则需要按行构造柱状图再用单调栈。

\subsubsection{LeetCode 226 翻转二叉树}

\textbf{题面模型。} 每个节点交换左右孩子即可，遍历顺序不影响最终结构。

\textbf{暴力瓶颈。} 显式队列或栈逐个节点交换，避免递归深度。

\textbf{边界反例。} 空树、单节点、只有一侧子树、重复值。

\textbf{变体迁移。} 这题简单但适合训练树节点指针修改。

\subsubsection{LeetCode 230 二叉搜索树中第 K 小}

\textbf{题面模型。} BST 中序遍历按升序访问节点。

\textbf{暴力瓶颈。} 显式栈中序，访问第 k 个节点时返回。

\textbf{边界反例。} k 为一、k 为节点数、树不平衡、是否有重复值。

\textbf{变体迁移。} 若频繁查询，可在节点维护子树大小成为顺序统计树。

\subsubsection{LeetCode 232 用栈实现队列}

\textbf{题面模型。} 两个栈分别负责输入顺序和输出顺序。

\textbf{暴力瓶颈。} 只有输出栈为空时，才把输入栈整体倒过去，保证均摊常数。

\textbf{边界反例。} 连续 peek、空队列、交替 push pop。

\textbf{变体迁移。} 用队列实现栈则要通过轮转保持最近元素在队首。

\subsubsection{LeetCode 235 二叉搜索树最近公共祖先}

\textbf{题面模型。} BST 的值域关系能决定两个目标在当前节点的哪一侧。

\textbf{暴力瓶颈。} 两个值都小去左，都大去右，否则当前节点就是分叉点。

\textbf{边界反例。} 一个节点是另一个祖先、目标顺序、重复值约定。

\textbf{变体迁移。} 普通二叉树不能用值域剪枝，需要父指针或后序状态。

\subsubsection{LeetCode 236 二叉树最近公共祖先}

\textbf{题面模型。} 普通树没有有序性，只能依靠祖先关系。

\textbf{暴力瓶颈。} 父指针集合或后序返回目标命中状态都可行。

\textbf{边界反例。} 目标不存在、p 等于 q、一个是另一个祖先、比较节点地址。

\textbf{变体迁移。} 多次 LCA 查询可预处理倍增祖先表。

\subsubsection{LeetCode 238 除自身以外数组的乘积}

\textbf{题面模型。} 不能用除法时，答案由左侧乘积和右侧乘积组成。

\textbf{暴力瓶颈。} 先写入左侧前缀积，再从右向左乘右侧后缀积。

\textbf{边界反例。} 一个零、多个零、负数、乘积溢出。

\textbf{变体迁移。} 这个题展示输出数组可作为状态存储的一部分。

\subsubsection{LeetCode 239 滑动窗口最大值}

\textbf{题面模型。} 每个窗口最大值来自仍未过期且没有被新元素支配的候选。

\textbf{暴力瓶颈。} 单调队列保存下标，队尾小于等于新值的候选被弹出。

\textbf{边界反例。} k 为一、k 等于数组长度、重复最大值、队首过期。

\textbf{变体迁移。} 受限子序列和使用同样队列维护最近 k 个 DP 最大值。

\subsubsection{LeetCode 240 搜索二维矩阵 II}

\textbf{题面模型。} 右上角同时具有向左变小、向下变大的排除能力。

\textbf{暴力瓶颈。} 当前值大于目标就左移，小于目标就下移。

\textbf{边界反例。} 空矩阵、单行、单列、目标在角落、重复值。

\textbf{变体迁移。} 从左上角无法一次排除一行或一列。

\subsubsection{LeetCode 279 完全平方数}

\textbf{题面模型。} 完全平方数可以看成无限使用的物品，目标是凑出 n 的最少数量。

\textbf{暴力瓶颈。} 完全背包最少物品数，或把数字看成图节点做 BFS。

\textbf{边界反例。} n 为零或一、完全平方数本身、较大 n。

\textbf{变体迁移。} 数学四平方定理可给更强解法，但面试 DP 更通用。

\subsubsection{LeetCode 283 移动零}

\textbf{题面模型。} 稳定保留非零元素顺序，把零集中到尾部。

\textbf{暴力瓶颈。} 慢指针写入下一个非零位置，最后补零或用交换。

\textbf{边界反例。} 全零、无零、零在开头结尾、稳定顺序要求。

\textbf{变体迁移。} 若目标是移动指定值，模型与移除元素相同。

\subsubsection{LeetCode 287 寻找重复数}

\textbf{题面模型。} 数组值域使下标到值形成函数图，重复数是环入口。

\textbf{暴力瓶颈。} 快慢指针不修改数组且常数空间；值域二分用抽屉原理。

\textbf{边界反例。} 重复数出现多次、不能排序、不能改数组、长度为 n 加一。

\textbf{变体迁移。} 快慢指针要求值都能作为合法下标。

\subsubsection{LeetCode 295 数据流的中位数}

\textbf{题面模型。} 数据流需要动态维护较小一半和较大一半。

\textbf{暴力瓶颈。} 大顶堆保存较小一半，小顶堆保存较大一半，大小差不超过一。

\textbf{边界反例。} 偶数个元素、负数、重复值、平均值溢出。

\textbf{变体迁移。} 若还要删除任意元素，需要延迟删除哈希表。

\subsubsection{LeetCode 300 最长递增子序列}

\textbf{题面模型。} 二次 DP 固定结尾，二分优化维护每个长度的最小结尾。

\textbf{暴力瓶颈。} tails 不是实际答案序列，而是未来扩展潜力表。

\textbf{边界反例。} 重复元素、严格递增和非递减区别、空数组。

\textbf{变体迁移。} 若要还原路径，需要额外前驱数组，不能只看 tails。

\subsubsection{LeetCode 309 最佳买卖股票含冷冻期}

\textbf{题面模型。} 冷冻期让买入来源受昨天卖出状态限制。

\textbf{暴力瓶颈。} 状态机维护持有、刚卖出、休息三种日终状态。

\textbf{边界反例。} 空价格、一天、连续上涨、卖出后不能次日买入。

\textbf{变体迁移。} 手续费版本可以在买入或卖出转移中扣费。

\subsubsection{LeetCode 322 零钱兑换}

\textbf{题面模型。} 金额是容量，硬币可无限使用，目标是最少硬币数。

\textbf{暴力瓶颈。} dp[value] 从 value 减 coin 的已知状态转移。

\textbf{边界反例。} 无法凑出、amount 为零、硬币面额大于目标、哨兵溢出。

\textbf{变体迁移。} 零钱兑换 II 问组合数，循环语义不同。

\subsubsection{LeetCode 337 打家劫舍 III}

\textbf{题面模型。} 树上相邻节点不能同时选，每个节点需要偷和不偷两个状态。

\textbf{暴力瓶颈。} 后序汇总：偷当前则孩子不能偷，不偷当前则孩子可偷可不偷。

\textbf{边界反例。} 空树、负值是否可能、链式树、递归深度。

\textbf{变体迁移。} 可用显式栈后序保存节点状态，避免调用栈。

\subsubsection{LeetCode 347 前 K 个高频元素}

\textbf{题面模型。} 先统计频次，再只维护前 K 个候选。

\textbf{暴力瓶颈。} 小顶堆大小超过 K 就弹出最低频；桶排序利用频次上界。

\textbf{边界反例。} K 等于不同元素数、频次相同、结果顺序不要求。

\textbf{变体迁移。} 若要按频次和字典序排序，比较器要同时处理两个键。

\subsubsection{LeetCode 394 字符串解码}

\textbf{题面模型。} 嵌套结构需要保存上一层字符串和重复次数。

\textbf{暴力瓶颈。} 遇到左括号入栈当前状态，右括号弹出并拼接展开。

\textbf{边界反例。} 多位数字、嵌套、空片段、数字后必须跟括号。

\textbf{变体迁移。} 递归下降解析更通用，但迭代栈更符合工程栈控制。

\subsubsection{LeetCode 416 分割等和子集}

\textbf{题面模型。} 总和一半是背包容量，每个数只能使用一次。

\textbf{暴力瓶颈。} 0-1 背包可达性，一维容量必须倒序。

\textbf{边界反例。} 总和为奇数、目标为零、数值较大、复杂度依赖 target。

\textbf{变体迁移。} 负数会破坏普通容量模型，需要重新建模。

\subsubsection{LeetCode 417 太平洋大西洋水流问题}

\textbf{题面模型。} 从每个格子向海搜索很慢，反向从海边沿非递减高度搜索。

\textbf{暴力瓶颈。} 分别标记能到太平洋和大西洋的格子，交集就是答案。

\textbf{边界反例。} 单行单列、平台高度、边界重复入队、方向条件反向。

\textbf{变体迁移。} 反向搜索是很多可达性题的重要技巧。

\subsubsection{LeetCode 435 无重叠区间}

\textbf{题面模型。} 保留最多不重叠区间等价于删除最少区间。

\textbf{暴力瓶颈。} 按终点升序选择，结束越早给后续留下空间越多。

\textbf{边界反例。} 端点相接是否重叠、完全覆盖、相同终点、空列表。

\textbf{变体迁移。} 用交换论证说明最早结束的选择不劣于其他第一选择。

\subsubsection{LeetCode 438 找到字符串中所有字母异位词}

\textbf{题面模型。} 固定长度窗口内字符频次等于目标频次即为答案。

\textbf{暴力瓶颈。} 右进左出维护计数，窗口长度达到目标长度后检查差异。

\textbf{边界反例。} 目标比源串长、重复字符、字符集固定、答案重叠。

\textbf{变体迁移。} 维护差异计数可以避免每次比较整个数组。

\subsubsection{LeetCode 450 删除二叉搜索树中的节点}

\textbf{题面模型。} 删除节点后仍要保持 BST 中序有序。

\textbf{暴力瓶颈。} 若有两个孩子，用后继或前驱替换当前值，再删除那个后继节点。

\textbf{边界反例。} 删除叶子、一个孩子、两个孩子、删除根节点。

\textbf{变体迁移。} 工程实现要清楚是移动值还是重连节点。

\subsubsection{LeetCode 452 用最少数量的箭引爆气球}

\textbf{题面模型。} 一支箭的位置可以覆盖所有包含该位置的区间。

\textbf{暴力瓶颈。} 按右端点排序，当前箭放在最早结束气球的右端。

\textbf{边界反例。} 端点相同、负坐标、区间包含、闭区间边界。

\textbf{变体迁移。} 它和无重叠区间是同一类最早结束贪心。

\subsubsection{LeetCode 494 目标和}

\textbf{题面模型。} 正负号选择可代数转换为选若干数凑出特定正和。

\textbf{暴力瓶颈。} 若 S 加 target 为奇数或目标绝对值过大，直接无解。

\textbf{边界反例。} 零元素会让方案数翻倍、目标为负、容量为零。

\textbf{变体迁移。} 也可用哈希 DP 保存所有可能和，但背包转换更高效。

\subsubsection{LeetCode 496 下一个更大元素 I}

\textbf{题面模型。} 第二个数组提供每个值右侧第一个更大值的全局关系。

\textbf{暴力瓶颈。} 单调栈预处理值到下一个更大值映射，再回答第一个数组查询。

\textbf{边界反例。} 不存在更大值、递减数组、值唯一假设。

\textbf{变体迁移。} 若值不唯一，应按下标而不是值建映射。

\subsubsection{LeetCode 503 下一个更大元素 II}

\textbf{题面模型。} 环形数组可以通过遍历两倍下标模拟绕回。

\textbf{暴力瓶颈。} 第一遍入栈，第二遍只帮助未解决下标找到答案。

\textbf{边界反例。} 全递减、全相等、长度一、取模下标。

\textbf{变体迁移。} 不要在第二遍重复入栈，否则可能死循环或重复计算。

\subsubsection{LeetCode 518 零钱兑换 II}

\textbf{题面模型。} 硬币无限使用，问组合数而不是最少数量。

\textbf{暴力瓶颈。} 外层遍历硬币、内层金额正序，保证组合按固定硬币顺序统计。

\textbf{边界反例。} amount 为零、无硬币、组合数较大、循环顺序反例。

\textbf{变体迁移。} 若外层金额内层硬币，会统计排列数。

\subsubsection{LeetCode 543 二叉树直径}

\textbf{题面模型。} 直径经过某节点时等于左高加右高，但向父亲只能贡献单边高度。

\textbf{暴力瓶颈。} 后序遍历同时更新全局直径并返回高度。

\textbf{边界反例。} 空树、单节点、链式树、边数还是节点数定义。

\textbf{变体迁移。} 最大路径和是同型题，但贡献值和全局更新有负数处理。

\subsubsection{LeetCode 547 省份数量}

\textbf{题面模型。} 城市连接关系形成无向图，省份就是连通分量。

\textbf{暴力瓶颈。} 并查集合并所有直接连接的城市，最后统计根数量。

\textbf{边界反例。} 邻接矩阵对称、自连接、孤立城市、只扫描上三角。

\textbf{变体迁移。} BFS 也可做；并查集更突出动态合并。

\subsubsection{LeetCode 560 和为 K 的子数组}

\textbf{题面模型。} 当前前缀和减去目标值，就是需要匹配的历史前缀。

\textbf{暴力瓶颈。} 哈希表保存前缀和出现次数，先查询再更新。

\textbf{边界反例。} 负数、目标为零、从下标零开始的区间、重复前缀。

\textbf{变体迁移。} 若要求最长区间，哈希表应保存最早位置而不是次数。

\subsubsection{LeetCode 581 最短无序连续子数组}

\textbf{题面模型。} 要找破坏全局有序性的最左和最右位置。

\textbf{暴力瓶颈。} 从左维护历史最大值确定右边界，从右维护历史最小值确定左边界。

\textbf{边界反例。} 已经有序、逆序、重复值、局部乱序。

\textbf{变体迁移。} 排序比较法更直观但复杂度更高。

\subsubsection{LeetCode 621 任务调度器}

\textbf{题面模型。} 相同任务之间有冷却间隔，核心受最高频任务骨架限制。

\textbf{暴力瓶颈。} 可用大顶堆模拟，也可用最高频公式计算最短时间。

\textbf{边界反例。} 冷却为零、多个最高频、空闲被其他任务填满。

\textbf{变体迁移。} 若任务有不同释放时间或执行时长，就变成更一般的调度问题。

\subsubsection{LeetCode 647 回文子串}

\textbf{题面模型。} 每个回文都有唯一中心或中心间隙。

\textbf{暴力瓶颈。} 对每个中心向两边扩展，扩展成功一次就计数一次。

\textbf{边界反例。} 空串、全相同字符、偶数回文、奇数回文。

\textbf{变体迁移。} Manacher 算法可线性解决，但中心扩展更适合面试基础。

\subsubsection{LeetCode 684 冗余连接}

\textbf{题面模型。} 树加一条边会形成唯一环。

\textbf{暴力瓶颈。} 按顺序加边，若一条边两端已连通，则它就是冗余边。

\textbf{边界反例。} 节点编号、最后一条成环边、重复边、并查集初始化大小。

\textbf{变体迁移。} 有向冗余连接需要同时处理入度为二和环，复杂度更高。

\subsubsection{LeetCode 695 岛屿的最大面积}

\textbf{题面模型。} 面积就是一个连通块中陆地格子的数量。

\textbf{暴力瓶颈。} 每发现新陆地就搜索完整连通块并计数，取最大值。

\textbf{边界反例。} 全水、全陆、细长岛、是否允许修改输入。

\textbf{变体迁移。} 与岛屿数量相比，只是连通块聚合值不同。

\subsubsection{LeetCode 704 二分查找}

\textbf{题面模型。} 基础题的重点是固定区间语义。

\textbf{暴力瓶颈。} 左闭右开写法能统一 lower bound 和存在性检查。

\textbf{边界反例。} 空数组、长度一、目标在边界、目标不存在。

\textbf{变体迁移。} 所有复杂二分都应该能退回到这道题的边界不变量。

\subsubsection{LeetCode 714 买卖股票含手续费}

\textbf{题面模型。} 手续费改变交易收益，仍可用持有和不持有状态。

\textbf{暴力瓶颈。} 卖出时扣费或买入时扣费都可以，但不要重复扣。

\textbf{边界反例。} 手续费为零、价格震荡、小利润不足手续费。

\textbf{变体迁移。} 这是状态机比局部贪心更稳的例子。

\subsubsection{LeetCode 739 每日温度}

\textbf{题面模型。} 每一天等待右侧第一个更高温度。

\textbf{暴力瓶颈。} 单调递减栈保存尚未找到答案的下标。

\textbf{边界反例。} 没有更高温、相等温度、递增、递减。

\textbf{变体迁移。} 下一个更大元素与它同型，只是输出值或距离不同。

\subsubsection{LeetCode 743 网络延迟时间}

\textbf{题面模型。} 有向正权图从起点到所有节点的最短路最大值。

\textbf{暴力瓶颈。} Dijkstra 求每个节点最短距离，最后取最大，若有无穷则不可达。

\textbf{边界反例。} 节点编号从一开始、不可达、重边、过期堆元素。

\textbf{变体迁移。} 边权相同才可用普通 BFS。

\subsubsection{LeetCode 752 打开转盘锁}

\textbf{题面模型。} 四位密码是状态，每次转动一位形成边。

\textbf{暴力瓶颈。} BFS 从起点扩展，死亡状态不能入队，访问状态避免重复。

\textbf{边界反例。} 起点死亡、目标就是起点、数字环绕、死亡集合很大。

\textbf{变体迁移。} 双向 BFS 能明显减少状态展开。

\subsubsection{LeetCode 763 划分字母区间}

\textbf{题面模型。} 每个片段必须覆盖其中所有字符的最后出现位置。

\textbf{暴力瓶颈。} 扫描维护当前片段最远右端，到达右端即可切分。

\textbf{边界反例。} 单字符、所有字符相同、字符多次跨段。

\textbf{变体迁移。} 这是最早合法切分贪心，能最大化片段数量。

\subsubsection{LeetCode 875 爱吃香蕉的珂珂}

\textbf{题面模型。} 速度越大，总耗时越小，答案具有单调可行性。

\textbf{暴力瓶颈。} 二分速度，检查函数用向上取整累加小时数。

\textbf{边界反例。} 速度下界、堆很大、h 等于堆数、乘法加法溢出。

\textbf{变体迁移。} 船运包裹和分割数组都是同类最小可行上限。

\subsubsection{LeetCode 973 最接近原点的 K 个点}

\textbf{题面模型。} 距离只需比较平方，避免开方误差和成本。

\textbf{暴力瓶颈。} 大小为 K 的大顶堆保存当前最近 K 个点中最远者。

\textbf{边界反例。} K 为一、K 为全部、距离相同、坐标平方溢出。

\textbf{变体迁移。} 快速选择可平均线性，但堆更适合数据流。

\subsubsection{LeetCode 994 腐烂的橘子}

\textbf{题面模型。} 所有初始腐烂橘子同时作为 BFS 源点。

\textbf{暴力瓶颈。} 按层扩散，分钟数等于最后一个新鲜橘子被首次访问的层数。

\textbf{边界反例。} 没有新鲜橘子、没有腐烂源、隔离新鲜橘子、空格阻断。

\textbf{变体迁移。} 多源 BFS 也适用于最近零距离、地图扩散等题。

\subsubsection{LeetCode 1011 在 D 天内送达包裹}

\textbf{题面模型。} 容量越大，需要天数越少，存在最小可行容量。

\textbf{暴力瓶颈。} 下界是最大包裹，上界是总重量，检查函数按顺序贪心装船。

\textbf{边界反例。} D 为一、D 等于包裹数、单个包裹超过猜测容量。

\textbf{变体迁移。} 它和分割数组最大值都是连续分段最大和模型。

\subsubsection{LeetCode 1143 最长公共子序列}

\textbf{题面模型。} 状态表示两个前缀的最长公共子序列长度。

\textbf{暴力瓶颈。} 末尾相同取左上加一，不同取上方和左方最大。

\textbf{边界反例。} 空串、完全相同、无公共字符、子序列不是子串。

\textbf{变体迁移。} 若要求还原序列，需要从表右下角反向追踪选择。

\subsubsection{LeetCode 1288 删除被覆盖区间}

\textbf{题面模型。} 被覆盖要求另一区间起点不晚且终点不早。

\textbf{暴力瓶颈。} 按起点升序、终点降序排序，扫描维护最大右端。

\textbf{边界反例。} 同起点区间、完全覆盖、无覆盖、端点相等。

\textbf{变体迁移。} 终点降序是关键，否则同起点小区间会先出现造成误判。

\subsubsection{LeetCode 1631 最小体力消耗路径}

\textbf{题面模型。} 路径代价是沿途最大高度差，不是边权和。

\textbf{暴力瓶颈。} Dijkstra 的松弛改成取当前代价和新边代价的最大值。

\textbf{边界反例。} 单格、平坦网格、必须绕路、多个相同代价。

\textbf{变体迁移。} 也可二分体力上限后 BFS 检查可达。


\section{逐题推导与实验工作坊}

这一节把前面的题号索引继续展开。每道题都从自己的题面模型出发，先说明慢在哪里，再说明状态怎样保存、哪些候选被安全排除、哪些边界会打破错误实现。题型共性只作为背景，真正要读的是每道题自己的瓶颈和反例。

\subsection{数组与原地维护工作坊}

先写一个不会漏答案的枚举或辅助数组版本，再观察哪些位置已经被前缀、后缀、慢指针或边界确定。优化时把数组划成语义清楚的区域，每次循环只扩大已确认区域。

\subsubsection{LeetCode 5 最长回文子串}

先从位置关系入手。本题的单点模型是回文围绕中心对称，中心可以是字符也可以是两个字符之间。朴素版可以为每个位置重新寻找最终状态，或者用辅助数组把答案先做出来；它虽然慢，却能说明哪些位置必须被覆盖。本题真正浪费的部分是中心扩展避免枚举子串后再逐个检查；每次扩展只比较新增左右边界，后面的优化只能消掉这部分重复，不能改变输出语义。

优化时把数组切成几段：已经确认的前缀、未知区、尚未处理的后缀，或者左侧贡献和右侧贡献。本题落到代码上就是中心扩展避免枚举子串后再逐个检查；每次扩展只比较新增左右边界。循环中每改一个下标，都要能说出哪一段扩大了，哪一段缩小了。放到“数组与原地维护”这条主线里看，关键原则是：优化要把数组切成若干语义明确的区间：已经处理好的前缀、未知区、已经确定的后缀，或者左侧贡献和右侧贡献。双指针、前后缀数组、原地交换和稳定写入，本质都是让每个元素只进入少数几次状态转移。

正确性可以按区间不变量证明：已处理部分满足题意，未知部分还没承诺，循环每次只把一个元素移入正确区域。对本题来说，回文围绕中心对称，中心可以是字符也可以是两个字符之间给出了答案范围，中心扩展避免枚举子串后再逐个检查；每次扩展只比较新增左右边界给出了推进方式。这道题的边界不能留到最后补丁式处理，实验时把偶数长度回文、单字符、全相同字符、多个同长答案列成表，再打印关键下标和有效区间。若某个元素被写入后又被未来步骤破坏，说明区间不变量没有成立。

C++ 实现上，C++ 中优先使用 \code{std::vector} 和清楚的整型下标；涉及乘积、面积、和或距离时提前考虑 \code{long long}。如果题目要求原地修改，返回长度、容器容量和有效区间要分开讲。如果追问变成“若要求回文子序列，应转成区间 DP，不能用中心扩展”，先判断原来的状态是否仍然覆盖新答案集合，再决定是微调边界、改 value 语义，还是换成另一类结构。

\subsubsection{LeetCode 26 删除有序数组重复项}

先从位置关系入手。本题的单点模型是有序数组让重复元素相邻，慢指针表示下一个唯一值写入位置。朴素版可以为每个位置重新寻找最终状态，或者用辅助数组把答案先做出来；它虽然慢，却能说明哪些位置必须被覆盖。本题真正浪费的部分是快指针扫描，只有发现新值才写入慢指针并推进，后面的优化只能消掉这部分重复，不能改变输出语义。

优化时把数组切成几段：已经确认的前缀、未知区、尚未处理的后缀，或者左侧贡献和右侧贡献。本题落到代码上就是快指针扫描，只有发现新值才写入慢指针并推进。循环中每改一个下标，都要能说出哪一段扩大了，哪一段缩小了。放到“数组与原地维护”这条主线里看，关键原则是：优化要把数组切成若干语义明确的区间：已经处理好的前缀、未知区、已经确定的后缀，或者左侧贡献和右侧贡献。双指针、前后缀数组、原地交换和稳定写入，本质都是让每个元素只进入少数几次状态转移。

正确性可以按区间不变量证明：已处理部分满足题意，未知部分还没承诺，循环每次只把一个元素移入正确区域。对本题来说，有序数组让重复元素相邻，慢指针表示下一个唯一值写入位置给出了答案范围，快指针扫描，只有发现新值才写入慢指针并推进给出了推进方式。这道题的边界不能留到最后补丁式处理，实验时把空数组、全重复、无重复、返回长度不等于物理容量列成表，再打印关键下标和有效区间。若某个元素被写入后又被未来步骤破坏，说明区间不变量没有成立。

C++ 实现上，C++ 中优先使用 \code{std::vector} 和清楚的整型下标；涉及乘积、面积、和或距离时提前考虑 \code{long long}。如果题目要求原地修改，返回长度、容器容量和有效区间要分开讲。如果追问变成“若允许每个元素最多保留两次，只需改变写入条件为和前两个比较”，先判断原来的状态是否仍然覆盖新答案集合，再决定是微调边界、改 value 语义，还是换成另一类结构。

\subsubsection{LeetCode 27 移除元素}

先从位置关系入手。本题的单点模型是原地过滤，慢指针左侧始终是不等于目标值的稳定前缀。朴素版可以为每个位置重新寻找最终状态，或者用辅助数组把答案先做出来；它虽然慢，却能说明哪些位置必须被覆盖。本题真正浪费的部分是保持顺序用快慢指针；不保持顺序可用左右交换减少写入，后面的优化只能消掉这部分重复，不能改变输出语义。

优化时把数组切成几段：已经确认的前缀、未知区、尚未处理的后缀，或者左侧贡献和右侧贡献。本题落到代码上就是保持顺序用快慢指针；不保持顺序可用左右交换减少写入。循环中每改一个下标，都要能说出哪一段扩大了，哪一段缩小了。放到“数组与原地维护”这条主线里看，关键原则是：优化要把数组切成若干语义明确的区间：已经处理好的前缀、未知区、已经确定的后缀，或者左侧贡献和右侧贡献。双指针、前后缀数组、原地交换和稳定写入，本质都是让每个元素只进入少数几次状态转移。

正确性可以按区间不变量证明：已处理部分满足题意，未知部分还没承诺，循环每次只把一个元素移入正确区域。对本题来说，原地过滤，慢指针左侧始终是不等于目标值的稳定前缀给出了答案范围，保持顺序用快慢指针；不保持顺序可用左右交换减少写入给出了推进方式。这道题的边界不能留到最后补丁式处理，实验时把所有元素都被移除、没有元素被移除、目标值在尾部密集列成表，再打印关键下标和有效区间。若某个元素被写入后又被未来步骤破坏，说明区间不变量没有成立。

C++ 实现上，C++ 中优先使用 \code{std::vector} 和清楚的整型下标；涉及乘积、面积、和或距离时提前考虑 \code{long long}。如果题目要求原地修改，返回长度、容器容量和有效区间要分开讲。如果追问变成“工程里要先确认是否要求稳定顺序，不同要求对应不同写法”，先判断原来的状态是否仍然覆盖新答案集合，再决定是微调边界、改 value 语义，还是换成另一类结构。

\subsubsection{LeetCode 31 下一个排列}

先从位置关系入手。本题的单点模型是字典序结构由最长非递增后缀决定，枢轴是后缀前一个位置。朴素版可以为每个位置重新寻找最终状态，或者用辅助数组把答案先做出来；它虽然慢，却能说明哪些位置必须被覆盖。本题真正浪费的部分是从右找第一个上升点，再从右找刚好更大的元素交换，最后反转后缀，后面的优化只能消掉这部分重复，不能改变输出语义。

优化时把数组切成几段：已经确认的前缀、未知区、尚未处理的后缀，或者左侧贡献和右侧贡献。本题落到代码上就是从右找第一个上升点，再从右找刚好更大的元素交换，最后反转后缀。循环中每改一个下标，都要能说出哪一段扩大了，哪一段缩小了。放到“数组与原地维护”这条主线里看，关键原则是：优化要把数组切成若干语义明确的区间：已经处理好的前缀、未知区、已经确定的后缀，或者左侧贡献和右侧贡献。双指针、前后缀数组、原地交换和稳定写入，本质都是让每个元素只进入少数几次状态转移。

正确性可以按区间不变量证明：已处理部分满足题意，未知部分还没承诺，循环每次只把一个元素移入正确区域。对本题来说，字典序结构由最长非递增后缀决定，枢轴是后缀前一个位置给出了答案范围，从右找第一个上升点，再从右找刚好更大的元素交换，最后反转后缀给出了推进方式。这道题的边界不能留到最后补丁式处理，实验时把完全降序、重复数字、只有一个元素、交换后后缀必须最小化列成表，再打印关键下标和有效区间。若某个元素被写入后又被未来步骤破坏，说明区间不变量没有成立。

C++ 实现上，C++ 中优先使用 \code{std::vector} 和清楚的整型下标；涉及乘积、面积、和或距离时提前考虑 \code{long long}。如果题目要求原地修改，返回长度、容器容量和有效区间要分开讲。如果追问变成“若要求上一个排列，比较方向和后缀处理对称变化”，先判断原来的状态是否仍然覆盖新答案集合，再决定是微调边界、改 value 语义，还是换成另一类结构。

\subsubsection{LeetCode 54 螺旋矩阵}

先从位置关系入手。本题的单点模型是四个边界收缩模拟一圈一圈访问。朴素版可以为每个位置重新寻找最终状态，或者用辅助数组把答案先做出来；它虽然慢，却能说明哪些位置必须被覆盖。本题真正浪费的部分是每走完一条边就收缩对应边界，并在走反方向前检查是否仍合法，后面的优化只能消掉这部分重复，不能改变输出语义。

优化时把数组切成几段：已经确认的前缀、未知区、尚未处理的后缀，或者左侧贡献和右侧贡献。本题落到代码上就是每走完一条边就收缩对应边界，并在走反方向前检查是否仍合法。循环中每改一个下标，都要能说出哪一段扩大了，哪一段缩小了。放到“数组与原地维护”这条主线里看，关键原则是：优化要把数组切成若干语义明确的区间：已经处理好的前缀、未知区、已经确定的后缀，或者左侧贡献和右侧贡献。双指针、前后缀数组、原地交换和稳定写入，本质都是让每个元素只进入少数几次状态转移。

正确性可以按区间不变量证明：已处理部分满足题意，未知部分还没承诺，循环每次只把一个元素移入正确区域。对本题来说，四个边界收缩模拟一圈一圈访问给出了答案范围，每走完一条边就收缩对应边界，并在走反方向前检查是否仍合法给出了推进方式。这道题的边界不能留到最后补丁式处理，实验时把单行、单列、空矩阵、最后一圈只剩一个元素列成表，再打印关键下标和有效区间。若某个元素被写入后又被未来步骤破坏，说明区间不变量没有成立。

C++ 实现上，C++ 中优先使用 \code{std::vector} 和清楚的整型下标；涉及乘积、面积、和或距离时提前考虑 \code{long long}。如果题目要求原地修改，返回长度、容器容量和有效区间要分开讲。如果追问变成“若改成生成矩阵，边界逻辑相同，只是读变成写”，先判断原来的状态是否仍然覆盖新答案集合，再决定是微调边界、改 value 语义，还是换成另一类结构。

\subsubsection{LeetCode 75 颜色分类}

先从位置关系入手。本题的单点模型是三指针维护零区、一号区、未知区和二区。朴素版可以为每个位置重新寻找最终状态，或者用辅助数组把答案先做出来；它虽然慢，却能说明哪些位置必须被覆盖。本题真正浪费的部分是遇到二时只移动右边界，不移动扫描指针，因为换回来的元素未知，后面的优化只能消掉这部分重复，不能改变输出语义。

优化时把数组切成几段：已经确认的前缀、未知区、尚未处理的后缀，或者左侧贡献和右侧贡献。本题落到代码上就是遇到二时只移动右边界，不移动扫描指针，因为换回来的元素未知。循环中每改一个下标，都要能说出哪一段扩大了，哪一段缩小了。放到“数组与原地维护”这条主线里看，关键原则是：优化要把数组切成若干语义明确的区间：已经处理好的前缀、未知区、已经确定的后缀，或者左侧贡献和右侧贡献。双指针、前后缀数组、原地交换和稳定写入，本质都是让每个元素只进入少数几次状态转移。

正确性可以按区间不变量证明：已处理部分满足题意，未知部分还没承诺，循环每次只把一个元素移入正确区域。对本题来说，三指针维护零区、一号区、未知区和二区给出了答案范围，遇到二时只移动右边界，不移动扫描指针，因为换回来的元素未知给出了推进方式。这道题的边界不能留到最后补丁式处理，实验时把全零、全二、已经有序、逆序、交换后漏检查列成表，再打印关键下标和有效区间。若某个元素被写入后又被未来步骤破坏，说明区间不变量没有成立。

C++ 实现上，C++ 中优先使用 \code{std::vector} 和清楚的整型下标；涉及乘积、面积、和或距离时提前考虑 \code{long long}。如果题目要求原地修改，返回长度、容器容量和有效区间要分开讲。如果追问变成“这是荷兰国旗问题，能推广到三类分区的原地维护”，先判断原来的状态是否仍然覆盖新答案集合，再决定是微调边界、改 value 语义，还是换成另一类结构。

\subsubsection{LeetCode 88 合并两个有序数组}

先从位置关系入手。本题的单点模型是第一个数组尾部有空位，适合从后往前写入最终最大值。朴素版可以为每个位置重新寻找最终状态，或者用辅助数组把答案先做出来；它虽然慢，却能说明哪些位置必须被覆盖。本题真正浪费的部分是逆向填充避免覆盖尚未读取的 nums1 元素，后面的优化只能消掉这部分重复，不能改变输出语义。

优化时把数组切成几段：已经确认的前缀、未知区、尚未处理的后缀，或者左侧贡献和右侧贡献。本题落到代码上就是逆向填充避免覆盖尚未读取的 nums1 元素。循环中每改一个下标，都要能说出哪一段扩大了，哪一段缩小了。放到“数组与原地维护”这条主线里看，关键原则是：优化要把数组切成若干语义明确的区间：已经处理好的前缀、未知区、已经确定的后缀，或者左侧贡献和右侧贡献。双指针、前后缀数组、原地交换和稳定写入，本质都是让每个元素只进入少数几次状态转移。

正确性可以按区间不变量证明：已处理部分满足题意，未知部分还没承诺，循环每次只把一个元素移入正确区域。对本题来说，第一个数组尾部有空位，适合从后往前写入最终最大值给出了答案范围，逆向填充避免覆盖尚未读取的 nums1 元素给出了推进方式。这道题的边界不能留到最后补丁式处理，实验时把第二个数组为空、第一个有效元素为空、重复值、尾部剩余列成表，再打印关键下标和有效区间。若某个元素被写入后又被未来步骤破坏，说明区间不变量没有成立。

C++ 实现上，C++ 中优先使用 \code{std::vector} 和清楚的整型下标；涉及乘积、面积、和或距离时提前考虑 \code{long long}。如果题目要求原地修改，返回长度、容器容量和有效区间要分开讲。如果追问变成“若没有额外空位，只能新开数组或使用更复杂原地归并”，先判断原来的状态是否仍然覆盖新答案集合，再决定是微调边界、改 value 语义，还是换成另一类结构。

\subsubsection{LeetCode 189 轮转数组}

先从位置关系入手。本题的单点模型是右移 k 位等价于把数组切成两段后换序。朴素版可以为每个位置重新寻找最终状态，或者用辅助数组把答案先做出来；它虽然慢，却能说明哪些位置必须被覆盖。本题真正浪费的部分是三次反转能原地完成：整体、前 k、后 n 减 k，后面的优化只能消掉这部分重复，不能改变输出语义。

优化时把数组切成几段：已经确认的前缀、未知区、尚未处理的后缀，或者左侧贡献和右侧贡献。本题落到代码上就是三次反转能原地完成：整体、前 k、后 n 减 k。循环中每改一个下标，都要能说出哪一段扩大了，哪一段缩小了。放到“数组与原地维护”这条主线里看，关键原则是：优化要把数组切成若干语义明确的区间：已经处理好的前缀、未知区、已经确定的后缀，或者左侧贡献和右侧贡献。双指针、前后缀数组、原地交换和稳定写入，本质都是让每个元素只进入少数几次状态转移。

正确性可以按区间不变量证明：已处理部分满足题意，未知部分还没承诺，循环每次只把一个元素移入正确区域。对本题来说，右移 k 位等价于把数组切成两段后换序给出了答案范围，三次反转能原地完成：整体、前 k、后 n 减 k给出了推进方式。这道题的边界不能留到最后补丁式处理，实验时把k 为零、k 大于 n、空数组、单元素列成表，再打印关键下标和有效区间。若某个元素被写入后又被未来步骤破坏，说明区间不变量没有成立。

C++ 实现上，C++ 中优先使用 \code{std::vector} 和清楚的整型下标；涉及乘积、面积、和或距离时提前考虑 \code{long long}。如果题目要求原地修改，返回长度、容器容量和有效区间要分开讲。如果追问变成“环状替换也可做，但要处理最大公约数个循环”，先判断原来的状态是否仍然覆盖新答案集合，再决定是微调边界、改 value 语义，还是换成另一类结构。

\subsubsection{LeetCode 238 除自身以外数组的乘积}

先从位置关系入手。本题的单点模型是不能用除法时，答案由左侧乘积和右侧乘积组成。朴素版可以为每个位置重新寻找最终状态，或者用辅助数组把答案先做出来；它虽然慢，却能说明哪些位置必须被覆盖。本题真正浪费的部分是先写入左侧前缀积，再从右向左乘右侧后缀积，后面的优化只能消掉这部分重复，不能改变输出语义。

优化时把数组切成几段：已经确认的前缀、未知区、尚未处理的后缀，或者左侧贡献和右侧贡献。本题落到代码上就是先写入左侧前缀积，再从右向左乘右侧后缀积。循环中每改一个下标，都要能说出哪一段扩大了，哪一段缩小了。放到“数组与原地维护”这条主线里看，关键原则是：优化要把数组切成若干语义明确的区间：已经处理好的前缀、未知区、已经确定的后缀，或者左侧贡献和右侧贡献。双指针、前后缀数组、原地交换和稳定写入，本质都是让每个元素只进入少数几次状态转移。

正确性可以按区间不变量证明：已处理部分满足题意，未知部分还没承诺，循环每次只把一个元素移入正确区域。对本题来说，不能用除法时，答案由左侧乘积和右侧乘积组成给出了答案范围，先写入左侧前缀积，再从右向左乘右侧后缀积给出了推进方式。这道题的边界不能留到最后补丁式处理，实验时把一个零、多个零、负数、乘积溢出列成表，再打印关键下标和有效区间。若某个元素被写入后又被未来步骤破坏，说明区间不变量没有成立。

C++ 实现上，C++ 中优先使用 \code{std::vector} 和清楚的整型下标；涉及乘积、面积、和或距离时提前考虑 \code{long long}。如果题目要求原地修改，返回长度、容器容量和有效区间要分开讲。如果追问变成“这个题展示输出数组可作为状态存储的一部分”，先判断原来的状态是否仍然覆盖新答案集合，再决定是微调边界、改 value 语义，还是换成另一类结构。

\subsubsection{LeetCode 283 移动零}

先从位置关系入手。本题的单点模型是稳定保留非零元素顺序，把零集中到尾部。朴素版可以为每个位置重新寻找最终状态，或者用辅助数组把答案先做出来；它虽然慢，却能说明哪些位置必须被覆盖。本题真正浪费的部分是慢指针写入下一个非零位置，最后补零或用交换，后面的优化只能消掉这部分重复，不能改变输出语义。

优化时把数组切成几段：已经确认的前缀、未知区、尚未处理的后缀，或者左侧贡献和右侧贡献。本题落到代码上就是慢指针写入下一个非零位置，最后补零或用交换。循环中每改一个下标，都要能说出哪一段扩大了，哪一段缩小了。放到“数组与原地维护”这条主线里看，关键原则是：优化要把数组切成若干语义明确的区间：已经处理好的前缀、未知区、已经确定的后缀，或者左侧贡献和右侧贡献。双指针、前后缀数组、原地交换和稳定写入，本质都是让每个元素只进入少数几次状态转移。

正确性可以按区间不变量证明：已处理部分满足题意，未知部分还没承诺，循环每次只把一个元素移入正确区域。对本题来说，稳定保留非零元素顺序，把零集中到尾部给出了答案范围，慢指针写入下一个非零位置，最后补零或用交换给出了推进方式。这道题的边界不能留到最后补丁式处理，实验时把全零、无零、零在开头结尾、稳定顺序要求列成表，再打印关键下标和有效区间。若某个元素被写入后又被未来步骤破坏，说明区间不变量没有成立。

C++ 实现上，C++ 中优先使用 \code{std::vector} 和清楚的整型下标；涉及乘积、面积、和或距离时提前考虑 \code{long long}。如果题目要求原地修改，返回长度、容器容量和有效区间要分开讲。如果追问变成“若目标是移动指定值，模型与移除元素相同”，先判断原来的状态是否仍然覆盖新答案集合，再决定是微调边界、改 value 语义，还是换成另一类结构。

\subsubsection{LeetCode 581 最短无序连续子数组}

先从位置关系入手。本题的单点模型是要找破坏全局有序性的最左和最右位置。朴素版可以为每个位置重新寻找最终状态，或者用辅助数组把答案先做出来；它虽然慢，却能说明哪些位置必须被覆盖。本题真正浪费的部分是从左维护历史最大值确定右边界，从右维护历史最小值确定左边界，后面的优化只能消掉这部分重复，不能改变输出语义。

优化时把数组切成几段：已经确认的前缀、未知区、尚未处理的后缀，或者左侧贡献和右侧贡献。本题落到代码上就是从左维护历史最大值确定右边界，从右维护历史最小值确定左边界。循环中每改一个下标，都要能说出哪一段扩大了，哪一段缩小了。放到“数组与原地维护”这条主线里看，关键原则是：优化要把数组切成若干语义明确的区间：已经处理好的前缀、未知区、已经确定的后缀，或者左侧贡献和右侧贡献。双指针、前后缀数组、原地交换和稳定写入，本质都是让每个元素只进入少数几次状态转移。

正确性可以按区间不变量证明：已处理部分满足题意，未知部分还没承诺，循环每次只把一个元素移入正确区域。对本题来说，要找破坏全局有序性的最左和最右位置给出了答案范围，从左维护历史最大值确定右边界，从右维护历史最小值确定左边界给出了推进方式。这道题的边界不能留到最后补丁式处理，实验时把已经有序、逆序、重复值、局部乱序列成表，再打印关键下标和有效区间。若某个元素被写入后又被未来步骤破坏，说明区间不变量没有成立。

C++ 实现上，C++ 中优先使用 \code{std::vector} 和清楚的整型下标；涉及乘积、面积、和或距离时提前考虑 \code{long long}。如果题目要求原地修改，返回长度、容器容量和有效区间要分开讲。如果追问变成“排序比较法更直观但复杂度更高”，先判断原来的状态是否仍然覆盖新答案集合，再决定是微调边界、改 value 语义，还是换成另一类结构。

\subsubsection{LeetCode 647 回文子串}

先从位置关系入手。本题的单点模型是每个回文都有唯一中心或中心间隙。朴素版可以为每个位置重新寻找最终状态，或者用辅助数组把答案先做出来；它虽然慢，却能说明哪些位置必须被覆盖。本题真正浪费的部分是对每个中心向两边扩展，扩展成功一次就计数一次，后面的优化只能消掉这部分重复，不能改变输出语义。

优化时把数组切成几段：已经确认的前缀、未知区、尚未处理的后缀，或者左侧贡献和右侧贡献。本题落到代码上就是对每个中心向两边扩展，扩展成功一次就计数一次。循环中每改一个下标，都要能说出哪一段扩大了，哪一段缩小了。放到“数组与原地维护”这条主线里看，关键原则是：优化要把数组切成若干语义明确的区间：已经处理好的前缀、未知区、已经确定的后缀，或者左侧贡献和右侧贡献。双指针、前后缀数组、原地交换和稳定写入，本质都是让每个元素只进入少数几次状态转移。

正确性可以按区间不变量证明：已处理部分满足题意，未知部分还没承诺，循环每次只把一个元素移入正确区域。对本题来说，每个回文都有唯一中心或中心间隙给出了答案范围，对每个中心向两边扩展，扩展成功一次就计数一次给出了推进方式。这道题的边界不能留到最后补丁式处理，实验时把空串、全相同字符、偶数回文、奇数回文列成表，再打印关键下标和有效区间。若某个元素被写入后又被未来步骤破坏，说明区间不变量没有成立。

C++ 实现上，C++ 中优先使用 \code{std::vector} 和清楚的整型下标；涉及乘积、面积、和或距离时提前考虑 \code{long long}。如果题目要求原地修改，返回长度、容器容量和有效区间要分开讲。如果追问变成“Manacher 算法可线性解决，但中心扩展更适合面试基础”，先判断原来的状态是否仍然覆盖新答案集合，再决定是微调边界、改 value 语义，还是换成另一类结构。

\subsection{滑动窗口与双指针工作坊}

先枚举所有窗口，再确认右端扩展和左端收缩是否具有单调性。只有窗口状态可以增量维护，且移动左端不会漏掉未来更优答案时，滑动窗口才成立。

\subsubsection{LeetCode 3 无重复字符的最长子串}

先枚举所有连续窗口，明确左端、右端和窗口内容怎样决定答案。本题的模型是窗口保存当前无重复子串，右端每次加入一个字符；暴力版会反复统计相邻窗口中相同的内容。慢点具体落在用上次出现位置直接跳过无效左端，左边界只能向右不能回退，所以后面要证明左端或右端移动时不会漏掉可能答案。

窗口题只在条件有单调性或窗口长度固定时成立。本题用用上次出现位置直接跳过无效左端，左边界只能向右不能回退维护窗口，含义是右端只带来一个新字符或新元素，左端只移走一个旧元素；如果输入条件改变导致状态不再单调，就必须换方法。放到“滑动窗口与双指针”这条主线里看，关键原则是：滑动窗口成立必须有单调性或固定长度边界。右端扩展让某个条件单调变好或变坏，左端收缩可以恢复合法性或缩短答案。若数组含负数、状态会来回震荡，就不能硬套窗口。

证明重点是排除左端。若某个左端被移过，它对应的窗口已经被完整检查，或继续保留只会让窗口更长、更不合法。本题的窗口保存当前无重复子串，右端每次加入一个字符和用上次出现位置直接跳过无效左端，左边界只能向右不能回退必须一起说明这一点。这道题的边界不能留到最后补丁式处理，实验时重点跑空串、全相同字符、重复字符出现在窗口左侧之外、扩展字符集，并打印左右端、窗口计数和答案。若左端发生回退，或者窗口失去合法性后仍更新答案，就能很快定位错误。

C++ 实现上，窗口计数若字符集固定，可以用 \code{std::array<int, 128>} 或 26 长度数组；字符集不固定再用哈希表。左闭右闭和左闭右开选一种坚持到底，避免窗口长度和删除位置错位。如果追问变成“若要求恰好包含 k 种字符，窗口条件从无重复变成种类计数”，先判断原来的状态是否仍然覆盖新答案集合，再决定是微调边界、改 value 语义，还是换成另一类结构。

\subsubsection{LeetCode 11 盛最多水的容器}

先枚举所有连续窗口，明确左端、右端和窗口内容怎样决定答案。本题的模型是左右指针代表当前选取的两条边，面积由短板和宽度共同决定；暴力版会反复统计相邻窗口中相同的内容。慢点具体落在每次移动短板，因为移动长板只会缩小宽度且不会提高短板，所以后面要证明左端或右端移动时不会漏掉可能答案。

窗口题只在条件有单调性或窗口长度固定时成立。本题用每次移动短板，因为移动长板只会缩小宽度且不会提高短板维护窗口，含义是右端只带来一个新字符或新元素，左端只移走一个旧元素；如果输入条件改变导致状态不再单调，就必须换方法。放到“滑动窗口与双指针”这条主线里看，关键原则是：滑动窗口成立必须有单调性或固定长度边界。右端扩展让某个条件单调变好或变坏，左端收缩可以恢复合法性或缩短答案。若数组含负数、状态会来回震荡，就不能硬套窗口。

证明重点是排除左端。若某个左端被移过，它对应的窗口已经被完整检查，或继续保留只会让窗口更长、更不合法。本题的左右指针代表当前选取的两条边，面积由短板和宽度共同决定和每次移动短板，因为移动长板只会缩小宽度且不会提高短板必须一起说明这一点。这道题的边界不能留到最后补丁式处理，实验时重点跑两端高度相等、数组长度为二、面积乘法可能溢出，并打印左右端、窗口计数和答案。若左端发生回退，或者窗口失去合法性后仍更新答案，就能很快定位错误。

C++ 实现上，窗口计数若字符集固定，可以用 \code{std::array<int, 128>} 或 26 长度数组；字符集不固定再用哈希表。左闭右闭和左闭右开选一种坚持到底，避免窗口长度和删除位置错位。如果追问变成“接雨水计算每个位置贡献，不能把本题双指针证明直接搬过去”，先判断原来的状态是否仍然覆盖新答案集合，再决定是微调边界、改 value 语义，还是换成另一类结构。

\subsubsection{LeetCode 15 三数之和}

先枚举所有连续窗口，明确左端、右端和窗口内容怎样决定答案。本题的模型是排序后固定第一个数，剩余两数转成有序双指针；暴力版会反复统计相邻窗口中相同的内容。慢点具体落在和太小移动左指针，和太大移动右指针；固定值和左右值都要跳过重复，所以后面要证明左端或右端移动时不会漏掉可能答案。

窗口题只在条件有单调性或窗口长度固定时成立。本题用和太小移动左指针，和太大移动右指针；固定值和左右值都要跳过重复维护窗口，含义是右端只带来一个新字符或新元素，左端只移走一个旧元素；如果输入条件改变导致状态不再单调，就必须换方法。放到“滑动窗口与双指针”这条主线里看，关键原则是：滑动窗口成立必须有单调性或固定长度边界。右端扩展让某个条件单调变好或变坏，左端收缩可以恢复合法性或缩短答案。若数组含负数、状态会来回震荡，就不能硬套窗口。

证明重点是排除左端。若某个左端被移过，它对应的窗口已经被完整检查，或继续保留只会让窗口更长、更不合法。本题的排序后固定第一个数，剩余两数转成有序双指针和和太小移动左指针，和太大移动右指针；固定值和左右值都要跳过重复必须一起说明这一点。这道题的边界不能留到最后补丁式处理，实验时重点跑全零、重复元素很多、没有答案、结果顺序不要求但组合不能重复，并打印左右端、窗口计数和答案。若左端发生回退，或者窗口失去合法性后仍更新答案，就能很快定位错误。

C++ 实现上，窗口计数若字符集固定，可以用 \code{std::array<int, 128>} 或 26 长度数组；字符集不固定再用哈希表。左闭右闭和左闭右开选一种坚持到底，避免窗口长度和删除位置错位。如果追问变成“四数之和再外层固定一个数，并用 long long 防止目标和溢出”，先判断原来的状态是否仍然覆盖新答案集合，再决定是微调边界、改 value 语义，还是换成另一类结构。

\subsubsection{LeetCode 16 最接近的三数之和}

先枚举所有连续窗口，明确左端、右端和窗口内容怎样决定答案。本题的模型是排序后固定一数，双指针寻找最接近目标的两数和；暴力版会反复统计相邻窗口中相同的内容。慢点具体落在每次根据当前和与目标的大小移动指针，同时更新最小绝对差，所以后面要证明左端或右端移动时不会漏掉可能答案。

窗口题只在条件有单调性或窗口长度固定时成立。本题用每次根据当前和与目标的大小移动指针，同时更新最小绝对差维护窗口，含义是右端只带来一个新字符或新元素，左端只移走一个旧元素；如果输入条件改变导致状态不再单调，就必须换方法。放到“滑动窗口与双指针”这条主线里看，关键原则是：滑动窗口成立必须有单调性或固定长度边界。右端扩展让某个条件单调变好或变坏，左端收缩可以恢复合法性或缩短答案。若数组含负数、状态会来回震荡，就不能硬套窗口。

证明重点是排除左端。若某个左端被移过，它对应的窗口已经被完整检查，或继续保留只会让窗口更长、更不合法。本题的排序后固定一数，双指针寻找最接近目标的两数和和每次根据当前和与目标的大小移动指针，同时更新最小绝对差必须一起说明这一点。这道题的边界不能留到最后补丁式处理，实验时重点跑差值相等、负数、目标很大、三数和溢出，并打印左右端、窗口计数和答案。若左端发生回退，或者窗口失去合法性后仍更新答案，就能很快定位错误。

C++ 实现上，窗口计数若字符集固定，可以用 \code{std::array<int, 128>} 或 26 长度数组；字符集不固定再用哈希表。左闭右闭和左闭右开选一种坚持到底，避免窗口长度和删除位置错位。如果追问变成“若要求返回所有最接近组合，需要记录同差值并去重”，先判断原来的状态是否仍然覆盖新答案集合，再决定是微调边界、改 value 语义，还是换成另一类结构。

\subsubsection{LeetCode 18 四数之和}

先枚举所有连续窗口，明确左端、右端和窗口内容怎样决定答案。本题的模型是两层固定加一层双指针，关键是剪枝和去重；暴力版会反复统计相邻窗口中相同的内容。慢点具体落在排序提供单调性，外层可用最小可能和和最大可能和提前跳过，所以后面要证明左端或右端移动时不会漏掉可能答案。

窗口题只在条件有单调性或窗口长度固定时成立。本题用排序提供单调性，外层可用最小可能和和最大可能和提前跳过维护窗口，含义是右端只带来一个新字符或新元素，左端只移走一个旧元素；如果输入条件改变导致状态不再单调，就必须换方法。放到“滑动窗口与双指针”这条主线里看，关键原则是：滑动窗口成立必须有单调性或固定长度边界。右端扩展让某个条件单调变好或变坏，左端收缩可以恢复合法性或缩短答案。若数组含负数、状态会来回震荡，就不能硬套窗口。

证明重点是排除左端。若某个左端被移过，它对应的窗口已经被完整检查，或继续保留只会让窗口更长、更不合法。本题的两层固定加一层双指针，关键是剪枝和去重和排序提供单调性，外层可用最小可能和和最大可能和提前跳过必须一起说明这一点。这道题的边界不能留到最后补丁式处理，实验时重点跑重复元素、目标超出 int、答案很多导致输出空间主导复杂度，并打印左右端、窗口计数和答案。若左端发生回退，或者窗口失去合法性后仍更新答案，就能很快定位错误。

C++ 实现上，窗口计数若字符集固定，可以用 \code{std::array<int, 128>} 或 26 长度数组；字符集不固定再用哈希表。左闭右闭和左闭右开选一种坚持到底，避免窗口长度和删除位置错位。如果追问变成“k 数之和可以递归降维，但工程实现要控制重复和边界”，先判断原来的状态是否仍然覆盖新答案集合，再决定是微调边界、改 value 语义，还是换成另一类结构。

\subsubsection{LeetCode 76 最小覆盖子串}

先枚举所有连续窗口，明确左端、右端和窗口内容怎样决定答案。本题的模型是窗口内字符计数必须覆盖目标计数，满足后尽量收缩左端；暴力版会反复统计相邻窗口中相同的内容。慢点具体落在用 formed 记录满足需求的字符种类，避免每次全量比较，所以后面要证明左端或右端移动时不会漏掉可能答案。

窗口题只在条件有单调性或窗口长度固定时成立。本题用用 formed 记录满足需求的字符种类，避免每次全量比较维护窗口，含义是右端只带来一个新字符或新元素，左端只移走一个旧元素；如果输入条件改变导致状态不再单调，就必须换方法。放到“滑动窗口与双指针”这条主线里看，关键原则是：滑动窗口成立必须有单调性或固定长度边界。右端扩展让某个条件单调变好或变坏，左端收缩可以恢复合法性或缩短答案。若数组含负数、状态会来回震荡，就不能硬套窗口。

证明重点是排除左端。若某个左端被移过，它对应的窗口已经被完整检查，或继续保留只会让窗口更长、更不合法。本题的窗口内字符计数必须覆盖目标计数，满足后尽量收缩左端和用 formed 记录满足需求的字符种类，避免每次全量比较必须一起说明这一点。这道题的边界不能留到最后补丁式处理，实验时重点跑目标有重复字符、源串缺字符、最小窗口在开头或结尾、大小写，并打印左右端、窗口计数和答案。若左端发生回退，或者窗口失去合法性后仍更新答案，就能很快定位错误。

C++ 实现上，窗口计数若字符集固定，可以用 \code{std::array<int, 128>} 或 26 长度数组；字符集不固定再用哈希表。左闭右闭和左闭右开选一种坚持到底，避免窗口长度和删除位置错位。如果追问变成“若字符集固定，数组计数更快；若是 Unicode，哈希表更通用”，先判断原来的状态是否仍然覆盖新答案集合，再决定是微调边界、改 value 语义，还是换成另一类结构。

\subsubsection{LeetCode 167 两数之和 II 输入有序数组}

先枚举所有连续窗口，明确左端、右端和窗口内容怎样决定答案。本题的模型是有序数组中左右指针的和随指针移动单调变化；暴力版会反复统计相邻窗口中相同的内容。慢点具体落在当前和太小左指针右移，太大右指针左移，等于目标返回，所以后面要证明左端或右端移动时不会漏掉可能答案。

窗口题只在条件有单调性或窗口长度固定时成立。本题用当前和太小左指针右移，太大右指针左移，等于目标返回维护窗口，含义是右端只带来一个新字符或新元素，左端只移走一个旧元素；如果输入条件改变导致状态不再单调，就必须换方法。放到“滑动窗口与双指针”这条主线里看，关键原则是：滑动窗口成立必须有单调性或固定长度边界。右端扩展让某个条件单调变好或变坏，左端收缩可以恢复合法性或缩短答案。若数组含负数、状态会来回震荡，就不能硬套窗口。

证明重点是排除左端。若某个左端被移过，它对应的窗口已经被完整检查，或继续保留只会让窗口更长、更不合法。本题的有序数组中左右指针的和随指针移动单调变化和当前和太小左指针右移，太大右指针左移，等于目标返回必须一起说明这一点。这道题的边界不能留到最后补丁式处理，实验时重点跑两个数在边界、重复值、题目下标从一开始、保证有解，并打印左右端、窗口计数和答案。若左端发生回退，或者窗口失去合法性后仍更新答案，就能很快定位错误。

C++ 实现上，窗口计数若字符集固定，可以用 \code{std::array<int, 128>} 或 26 长度数组；字符集不固定再用哈希表。左闭右闭和左闭右开选一种坚持到底，避免窗口长度和删除位置错位。如果追问变成“无序版本通常用哈希表，不能直接套双指针”，先判断原来的状态是否仍然覆盖新答案集合，再决定是微调边界、改 value 语义，还是换成另一类结构。

\subsubsection{LeetCode 209 长度最小的子数组}

先枚举所有连续窗口，明确左端、右端和窗口内容怎样决定答案。本题的模型是正数数组让窗口和随右端增加、随左端减少；暴力版会反复统计相邻窗口中相同的内容。慢点具体落在右端扩展到满足目标后，不断左移缩短并更新答案，所以后面要证明左端或右端移动时不会漏掉可能答案。

窗口题只在条件有单调性或窗口长度固定时成立。本题用右端扩展到满足目标后，不断左移缩短并更新答案维护窗口，含义是右端只带来一个新字符或新元素，左端只移走一个旧元素；如果输入条件改变导致状态不再单调，就必须换方法。放到“滑动窗口与双指针”这条主线里看，关键原则是：滑动窗口成立必须有单调性或固定长度边界。右端扩展让某个条件单调变好或变坏，左端收缩可以恢复合法性或缩短答案。若数组含负数、状态会来回震荡，就不能硬套窗口。

证明重点是排除左端。若某个左端被移过，它对应的窗口已经被完整检查，或继续保留只会让窗口更长、更不合法。本题的正数数组让窗口和随右端增加、随左端减少和右端扩展到满足目标后，不断左移缩短并更新答案必须一起说明这一点。这道题的边界不能留到最后补丁式处理，实验时重点跑不存在答案、单元素达到目标、全正是必要条件，并打印左右端、窗口计数和答案。若左端发生回退，或者窗口失去合法性后仍更新答案，就能很快定位错误。

C++ 实现上，窗口计数若字符集固定，可以用 \code{std::array<int, 128>} 或 26 长度数组；字符集不固定再用哈希表。左闭右闭和左闭右开选一种坚持到底，避免窗口长度和删除位置错位。如果追问变成“若含负数，需要前缀和加单调队列”，先判断原来的状态是否仍然覆盖新答案集合，再决定是微调边界、改 value 语义，还是换成另一类结构。

\subsubsection{LeetCode 438 找到字符串中所有字母异位词}

先枚举所有连续窗口，明确左端、右端和窗口内容怎样决定答案。本题的模型是固定长度窗口内字符频次等于目标频次即为答案；暴力版会反复统计相邻窗口中相同的内容。慢点具体落在右进左出维护计数，窗口长度达到目标长度后检查差异，所以后面要证明左端或右端移动时不会漏掉可能答案。

窗口题只在条件有单调性或窗口长度固定时成立。本题用右进左出维护计数，窗口长度达到目标长度后检查差异维护窗口，含义是右端只带来一个新字符或新元素，左端只移走一个旧元素；如果输入条件改变导致状态不再单调，就必须换方法。放到“滑动窗口与双指针”这条主线里看，关键原则是：滑动窗口成立必须有单调性或固定长度边界。右端扩展让某个条件单调变好或变坏，左端收缩可以恢复合法性或缩短答案。若数组含负数、状态会来回震荡，就不能硬套窗口。

证明重点是排除左端。若某个左端被移过，它对应的窗口已经被完整检查，或继续保留只会让窗口更长、更不合法。本题的固定长度窗口内字符频次等于目标频次即为答案和右进左出维护计数，窗口长度达到目标长度后检查差异必须一起说明这一点。这道题的边界不能留到最后补丁式处理，实验时重点跑目标比源串长、重复字符、字符集固定、答案重叠，并打印左右端、窗口计数和答案。若左端发生回退，或者窗口失去合法性后仍更新答案，就能很快定位错误。

C++ 实现上，窗口计数若字符集固定，可以用 \code{std::array<int, 128>} 或 26 长度数组；字符集不固定再用哈希表。左闭右闭和左闭右开选一种坚持到底，避免窗口长度和删除位置错位。如果追问变成“维护差异计数可以避免每次比较整个数组”，先判断原来的状态是否仍然覆盖新答案集合，再决定是微调边界、改 value 语义，还是换成另一类结构。

\subsubsection{LeetCode 567 字符串的排列}

先枚举所有连续窗口，明确左端、右端和窗口内容怎样决定答案。本题的模型是目标串任意排列出现在源串中，等价于固定长度窗口频次相同；暴力版会反复统计相邻窗口中相同的内容。慢点具体落在窗口长度固定为目标长度，右进左出维护字符计数差异，所以后面要证明左端或右端移动时不会漏掉可能答案。

窗口题只在条件有单调性或窗口长度固定时成立。本题用窗口长度固定为目标长度，右进左出维护字符计数差异维护窗口，含义是右端只带来一个新字符或新元素，左端只移走一个旧元素；如果输入条件改变导致状态不再单调，就必须换方法。放到“滑动窗口与双指针”这条主线里看，关键原则是：滑动窗口成立必须有单调性或固定长度边界。右端扩展让某个条件单调变好或变坏，左端收缩可以恢复合法性或缩短答案。若数组含负数、状态会来回震荡，就不能硬套窗口。

证明重点是排除左端。若某个左端被移过，它对应的窗口已经被完整检查，或继续保留只会让窗口更长、更不合法。本题的目标串任意排列出现在源串中，等价于固定长度窗口频次相同和窗口长度固定为目标长度，右进左出维护字符计数差异必须一起说明这一点。这道题的边界不能留到最后补丁式处理，实验时重点跑目标比源串长、重复字符、字符集固定、窗口重叠，并打印左右端、窗口计数和答案。若左端发生回退，或者窗口失去合法性后仍更新答案，就能很快定位错误。

C++ 实现上，窗口计数若字符集固定，可以用 \code{std::array<int, 128>} 或 26 长度数组；字符集不固定再用哈希表。左闭右闭和左闭右开选一种坚持到底，避免窗口长度和删除位置错位。如果追问变成“与找到所有异位词同型，只是返回布尔值”，先判断原来的状态是否仍然覆盖新答案集合，再决定是微调边界、改 value 语义，还是换成另一类结构。

\subsubsection{LeetCode 862 和至少为 K 的最短子数组}

先枚举所有连续窗口，明确左端、右端和窗口内容怎样决定答案。本题的模型是数组允许负数，普通滑动窗口的和不再随左端移动单调变化；暴力版会反复统计相邻窗口中相同的内容。慢点具体落在在前缀和数组上用单调队列维护可能成为最优左端的下标，所以后面要证明左端或右端移动时不会漏掉可能答案。

窗口题只在条件有单调性或窗口长度固定时成立。本题用在前缀和数组上用单调队列维护可能成为最优左端的下标维护窗口，含义是右端只带来一个新字符或新元素，左端只移走一个旧元素；如果输入条件改变导致状态不再单调，就必须换方法。放到“滑动窗口与双指针”这条主线里看，关键原则是：滑动窗口成立必须有单调性或固定长度边界。右端扩展让某个条件单调变好或变坏，左端收缩可以恢复合法性或缩短答案。若数组含负数、状态会来回震荡，就不能硬套窗口。

证明重点是排除左端。若某个左端被移过，它对应的窗口已经被完整检查，或继续保留只会让窗口更长、更不合法。本题的数组允许负数，普通滑动窗口的和不再随左端移动单调变化和在前缀和数组上用单调队列维护可能成为最优左端的下标必须一起说明这一点。这道题的边界不能留到最后补丁式处理，实验时重点跑负数、答案不存在、K 为负或零、队尾支配关系，并打印左右端、窗口计数和答案。若左端发生回退，或者窗口失去合法性后仍更新答案，就能很快定位错误。

C++ 实现上，窗口计数若字符集固定，可以用 \code{std::array<int, 128>} 或 26 长度数组；字符集不固定再用哈希表。左闭右闭和左闭右开选一种坚持到底，避免窗口长度和删除位置错位。如果追问变成“这是从窗口题升级到前缀和加单调队列的代表”，先判断原来的状态是否仍然覆盖新答案集合，再决定是微调边界、改 value 语义，还是换成另一类结构。

\subsubsection{LeetCode 1438 绝对差不超过限制的最长连续子数组}

先枚举所有连续窗口，明确左端、右端和窗口内容怎样决定答案。本题的模型是窗口合法性取决于窗口内最大值和最小值之差；暴力版会反复统计相邻窗口中相同的内容。慢点具体落在两个单调队列分别维护最大和最小，超过限制时移动左端并清理过期下标，所以后面要证明左端或右端移动时不会漏掉可能答案。

窗口题只在条件有单调性或窗口长度固定时成立。本题用两个单调队列分别维护最大和最小，超过限制时移动左端并清理过期下标维护窗口，含义是右端只带来一个新字符或新元素，左端只移走一个旧元素；如果输入条件改变导致状态不再单调，就必须换方法。放到“滑动窗口与双指针”这条主线里看，关键原则是：滑动窗口成立必须有单调性或固定长度边界。右端扩展让某个条件单调变好或变坏，左端收缩可以恢复合法性或缩短答案。若数组含负数、状态会来回震荡，就不能硬套窗口。

证明重点是排除左端。若某个左端被移过，它对应的窗口已经被完整检查，或继续保留只会让窗口更长、更不合法。本题的窗口合法性取决于窗口内最大值和最小值之差和两个单调队列分别维护最大和最小，超过限制时移动左端并清理过期下标必须一起说明这一点。这道题的边界不能留到最后补丁式处理，实验时重点跑limit 为零、重复元素、最大最小同时过期、窗口收缩多次，并打印左右端、窗口计数和答案。若左端发生回退，或者窗口失去合法性后仍更新答案，就能很快定位错误。

C++ 实现上，窗口计数若字符集固定，可以用 \code{std::array<int, 128>} 或 26 长度数组；字符集不固定再用哈希表。左闭右闭和左闭右开选一种坚持到底，避免窗口长度和删除位置错位。如果追问变成“它说明窗口状态可以是动态极值，不只是频次或总和”，先判断原来的状态是否仍然覆盖新答案集合，再决定是微调边界、改 value 语义，还是换成另一类结构。

\subsection{前缀和与哈希表工作坊}

先用双层枚举区间，再把区间条件改写成两个前缀状态的关系。哈希表保存历史前缀的次数、最早位置或存在性，具体保存什么由题目目标决定。

\subsubsection{LeetCode 1 两数之和}

先用双层循环枚举区间或候选对，再把条件写成历史状态和当前状态的关系。本题的模型是当前元素只需要寻找唯一补数，历史值到下标的映射就是最小状态；暴力版的问题不是不会找，而是每个当前位置都回头扫太多历史。优化目标就是把先查再插入，避免同一个下标被用两次；若数组有序才考虑双指针变成一次查表或一次状态读取。

哈希表保存的是当前下标之前的历史，不是随手放进去的缓存。本题需要先查再插入，避免同一个下标被用两次；若数组有序才考虑双指针，所以要先决定 value 是次数、最早位置、最近位置还是布尔值。这个决定直接改变答案类型。放到“前缀和与哈希表”这条主线里看，关键原则是：前缀和把区间问题改写成两个前缀状态的差；哈希表把寻找匹配前缀改成查表。频次表、最早位置表、最近位置表对应不同目标：计数、最长区间、最短区间不能混用。

证明回到代数关系：任意合法答案都能对应到一个当前状态和一个历史状态；算法查到的历史状态也一定构成合法答案。本题的当前元素只需要寻找唯一补数，历史值到下标的映射就是最小状态说明要找什么，先查再插入，避免同一个下标被用两次；若数组有序才考虑双指针说明历史表怎样保存它。这道题的边界不能留到最后补丁式处理，实验时把重复数字、目标为偶数且补数等于当前值、没有答案时返回空结果加入对拍集合，用平方暴力和优化版比较。失败时先看初始历史状态、查询更新顺序和哈希表 value 类型。

C++ 实现上，哈希表查询优先用 \code{find}，只在确实要插入或累加时使用 \code{operator[]}。前缀和、区间和、计数乘法常需要 \code{long long}；模运算还要处理负余数归一化。如果追问变成“若改成返回所有不重复数对，要先排序或用频次表处理重复”，先判断原来的状态是否仍然覆盖新答案集合，再决定是微调边界、改 value 语义，还是换成另一类结构。

\subsubsection{LeetCode 49 字母异位词分组}

先用双层循环枚举区间或候选对，再把条件写成历史状态和当前状态的关系。本题的模型是异位词共享同一个规范表示，可以是排序字符串或字符计数；暴力版的问题不是不会找，而是每个当前位置都回头扫太多历史。优化目标就是把哈希表按规范键分桶，避免两两比较字符串变成一次查表或一次状态读取。

哈希表保存的是当前下标之前的历史，不是随手放进去的缓存。本题需要哈希表按规范键分桶，避免两两比较字符串，所以要先决定 value 是次数、最早位置、最近位置还是布尔值。这个决定直接改变答案类型。放到“前缀和与哈希表”这条主线里看，关键原则是：前缀和把区间问题改写成两个前缀状态的差；哈希表把寻找匹配前缀改成查表。频次表、最早位置表、最近位置表对应不同目标：计数、最长区间、最短区间不能混用。

证明回到代数关系：任意合法答案都能对应到一个当前状态和一个历史状态；算法查到的历史状态也一定构成合法答案。本题的异位词共享同一个规范表示，可以是排序字符串或字符计数说明要找什么，哈希表按规范键分桶，避免两两比较字符串说明历史表怎样保存它。这道题的边界不能留到最后补丁式处理，实验时把空字符串、重复单词、字符集不止小写字母、计数键必须无歧义加入对拍集合，用平方暴力和优化版比较。失败时先看初始历史状态、查询更新顺序和哈希表 value 类型。

C++ 实现上，哈希表查询优先用 \code{find}，只在确实要插入或累加时使用 \code{operator[]}。前缀和、区间和、计数乘法常需要 \code{long long}；模运算还要处理负余数归一化。如果追问变成“若字符串很长且字符集固定，计数键通常比排序键更稳定”，先判断原来的状态是否仍然覆盖新答案集合，再决定是微调边界、改 value 语义，还是换成另一类结构。

\subsubsection{LeetCode 128 最长连续序列}

先用双层循环枚举区间或候选对，再把条件写成历史状态和当前状态的关系。本题的模型是连续段只需要从没有前驱的数开始扩展；暴力版的问题不是不会找，而是每个当前位置都回头扫太多历史。优化目标就是把哈希集合判断 x 减一 是否存在，跳过非起点避免重复扫描变成一次查表或一次状态读取。

哈希表保存的是当前下标之前的历史，不是随手放进去的缓存。本题需要哈希集合判断 x 减一 是否存在，跳过非起点避免重复扫描，所以要先决定 value 是次数、最早位置、最近位置还是布尔值。这个决定直接改变答案类型。放到“前缀和与哈希表”这条主线里看，关键原则是：前缀和把区间问题改写成两个前缀状态的差；哈希表把寻找匹配前缀改成查表。频次表、最早位置表、最近位置表对应不同目标：计数、最长区间、最短区间不能混用。

证明回到代数关系：任意合法答案都能对应到一个当前状态和一个历史状态；算法查到的历史状态也一定构成合法答案。本题的连续段只需要从没有前驱的数开始扩展说明要找什么，哈希集合判断 x 减一 是否存在，跳过非起点避免重复扫描说明历史表怎样保存它。这道题的边界不能留到最后补丁式处理，实验时把重复元素、负数、空数组、整数边界加入对拍集合，用平方暴力和优化版比较。失败时先看初始历史状态、查询更新顺序和哈希表 value 类型。

C++ 实现上，哈希表查询优先用 \code{find}，只在确实要插入或累加时使用 \code{operator[]}。前缀和、区间和、计数乘法常需要 \code{long long}；模运算还要处理负余数归一化。如果追问变成“若数据流动态插入，可以用并查集或区间合并维护长度”，先判断原来的状态是否仍然覆盖新答案集合，再决定是微调边界、改 value 语义，还是换成另一类结构。

\subsubsection{LeetCode 202 快乐数}

先用双层循环枚举区间或候选对，再把条件写成历史状态和当前状态的关系。本题的模型是反复变换数字会进入一或循环，状态空间最终有限；暴力版的问题不是不会找，而是每个当前位置都回头扫太多历史。优化目标就是把用集合检测重复状态，或用快慢指针看成函数图变成一次查表或一次状态读取。

哈希表保存的是当前下标之前的历史，不是随手放进去的缓存。本题需要用集合检测重复状态，或用快慢指针看成函数图，所以要先决定 value 是次数、最早位置、最近位置还是布尔值。这个决定直接改变答案类型。放到“前缀和与哈希表”这条主线里看，关键原则是：前缀和把区间问题改写成两个前缀状态的差；哈希表把寻找匹配前缀改成查表。频次表、最早位置表、最近位置表对应不同目标：计数、最长区间、最短区间不能混用。

证明回到代数关系：任意合法答案都能对应到一个当前状态和一个历史状态；算法查到的历史状态也一定构成合法答案。本题的反复变换数字会进入一或循环，状态空间最终有限说明要找什么，用集合检测重复状态，或用快慢指针看成函数图说明历史表怎样保存它。这道题的边界不能留到最后补丁式处理，实验时把输入为一、进入循环、位平方和函数、负数不在题意内加入对拍集合，用平方暴力和优化版比较。失败时先看初始历史状态、查询更新顺序和哈希表 value 类型。

C++ 实现上，哈希表查询优先用 \code{find}，只在确实要插入或累加时使用 \code{operator[]}。前缀和、区间和、计数乘法常需要 \code{long long}；模运算还要处理负余数归一化。如果追问变成“快慢指针版本能训练把数值迭代看成链表”，先判断原来的状态是否仍然覆盖新答案集合，再决定是微调边界、改 value 语义，还是换成另一类结构。

\subsubsection{LeetCode 217 存在重复元素}

先用双层循环枚举区间或候选对，再把条件写成历史状态和当前状态的关系。本题的模型是判断是否出现过是哈希集合最直接的职责；暴力版的问题不是不会找，而是每个当前位置都回头扫太多历史。优化目标就是把扫描时若插入前已经存在则立即返回真变成一次查表或一次状态读取。

哈希表保存的是当前下标之前的历史，不是随手放进去的缓存。本题需要扫描时若插入前已经存在则立即返回真，所以要先决定 value 是次数、最早位置、最近位置还是布尔值。这个决定直接改变答案类型。放到“前缀和与哈希表”这条主线里看，关键原则是：前缀和把区间问题改写成两个前缀状态的差；哈希表把寻找匹配前缀改成查表。频次表、最早位置表、最近位置表对应不同目标：计数、最长区间、最短区间不能混用。

证明回到代数关系：任意合法答案都能对应到一个当前状态和一个历史状态；算法查到的历史状态也一定构成合法答案。本题的判断是否出现过是哈希集合最直接的职责说明要找什么，扫描时若插入前已经存在则立即返回真说明历史表怎样保存它。这道题的边界不能留到最后补丁式处理，实验时把空数组、全唯一、重复在末尾、值域很小加入对拍集合，用平方暴力和优化版比较。失败时先看初始历史状态、查询更新顺序和哈希表 value 类型。

C++ 实现上，哈希表查询优先用 \code{find}，只在确实要插入或累加时使用 \code{operator[]}。前缀和、区间和、计数乘法常需要 \code{long long}；模运算还要处理负余数归一化。如果追问变成“若允许修改输入，排序后看相邻可以降低额外空间”，先判断原来的状态是否仍然覆盖新答案集合，再决定是微调边界、改 value 语义，还是换成另一类结构。

\subsubsection{LeetCode 242 有效的字母异位词}

先用双层循环枚举区间或候选对，再把条件写成历史状态和当前状态的关系。本题的模型是两个字符串是异位词等价于每个字符出现次数相同；暴力版的问题不是不会找，而是每个当前位置都回头扫太多历史。优化目标就是把固定小写字母集时用定长计数数组，字符集未知时用哈希表变成一次查表或一次状态读取。

哈希表保存的是当前下标之前的历史，不是随手放进去的缓存。本题需要固定小写字母集时用定长计数数组，字符集未知时用哈希表，所以要先决定 value 是次数、最早位置、最近位置还是布尔值。这个决定直接改变答案类型。放到“前缀和与哈希表”这条主线里看，关键原则是：前缀和把区间问题改写成两个前缀状态的差；哈希表把寻找匹配前缀改成查表。频次表、最早位置表、最近位置表对应不同目标：计数、最长区间、最短区间不能混用。

证明回到代数关系：任意合法答案都能对应到一个当前状态和一个历史状态；算法查到的历史状态也一定构成合法答案。本题的两个字符串是异位词等价于每个字符出现次数相同说明要找什么，固定小写字母集时用定长计数数组，字符集未知时用哈希表说明历史表怎样保存它。这道题的边界不能留到最后补丁式处理，实验时把长度不同、空串、Unicode 字符、计数减到负数加入对拍集合，用平方暴力和优化版比较。失败时先看初始历史状态、查询更新顺序和哈希表 value 类型。

C++ 实现上，哈希表查询优先用 \code{find}，只在确实要插入或累加时使用 \code{operator[]}。前缀和、区间和、计数乘法常需要 \code{long long}；模运算还要处理负余数归一化。如果追问变成“异位词分组把这个二元比较扩展成规范 key 分桶”，先判断原来的状态是否仍然覆盖新答案集合，再决定是微调边界、改 value 语义，还是换成另一类结构。

\subsubsection{LeetCode 325 和等于 k 的最长子数组长度}

先用双层循环枚举区间或候选对，再把条件写成历史状态和当前状态的关系。本题的模型是最长区间要求在同一前缀差关系下尽量拉大左右端距离；暴力版的问题不是不会找，而是每个当前位置都回头扫太多历史。优化目标就是把哈希表保存每个前缀和最早出现下标，当前前缀查询 prefix 减 k变成一次查表或一次状态读取。

哈希表保存的是当前下标之前的历史，不是随手放进去的缓存。本题需要哈希表保存每个前缀和最早出现下标，当前前缀查询 prefix 减 k，所以要先决定 value 是次数、最早位置、最近位置还是布尔值。这个决定直接改变答案类型。放到“前缀和与哈希表”这条主线里看，关键原则是：前缀和把区间问题改写成两个前缀状态的差；哈希表把寻找匹配前缀改成查表。频次表、最早位置表、最近位置表对应不同目标：计数、最长区间、最短区间不能混用。

证明回到代数关系：任意合法答案都能对应到一个当前状态和一个历史状态；算法查到的历史状态也一定构成合法答案。本题的最长区间要求在同一前缀差关系下尽量拉大左右端距离说明要找什么，哈希表保存每个前缀和最早出现下标，当前前缀查询 prefix 减 k说明历史表怎样保存它。这道题的边界不能留到最后补丁式处理，实验时把从下标零开始、负数、前缀重复时不能覆盖最早位置、答案不存在加入对拍集合，用平方暴力和优化版比较。失败时先看初始历史状态、查询更新顺序和哈希表 value 类型。

C++ 实现上，哈希表查询优先用 \code{find}，只在确实要插入或累加时使用 \code{operator[]}。前缀和、区间和、计数乘法常需要 \code{long long}；模运算还要处理负余数归一化。如果追问变成“计数版本保存频次，最长版本保存最早位置，这是核心差别”，先判断原来的状态是否仍然覆盖新答案集合，再决定是微调边界、改 value 语义，还是换成另一类结构。

\subsubsection{LeetCode 454 四数相加 II}

先用双层循环枚举区间或候选对，再把条件写成历史状态和当前状态的关系。本题的模型是四个数组的和为零可以拆成两个二数组和互为相反数；暴力版的问题不是不会找，而是每个当前位置都回头扫太多历史。优化目标就是把先统计 A 和 B 的所有两数和频次，再枚举 C 和 D 查询相反数变成一次查表或一次状态读取。

哈希表保存的是当前下标之前的历史，不是随手放进去的缓存。本题需要先统计 A 和 B 的所有两数和频次，再枚举 C 和 D 查询相反数，所以要先决定 value 是次数、最早位置、最近位置还是布尔值。这个决定直接改变答案类型。放到“前缀和与哈希表”这条主线里看，关键原则是：前缀和把区间问题改写成两个前缀状态的差；哈希表把寻找匹配前缀改成查表。频次表、最早位置表、最近位置表对应不同目标：计数、最长区间、最短区间不能混用。

证明回到代数关系：任意合法答案都能对应到一个当前状态和一个历史状态；算法查到的历史状态也一定构成合法答案。本题的四个数组的和为零可以拆成两个二数组和互为相反数说明要找什么，先统计 A 和 B 的所有两数和频次，再枚举 C 和 D 查询相反数说明历史表怎样保存它。这道题的边界不能留到最后补丁式处理，实验时把重复值、负数、答案数量很大、两数和范围加入对拍集合，用平方暴力和优化版比较。失败时先看初始历史状态、查询更新顺序和哈希表 value 类型。

C++ 实现上，哈希表查询优先用 \code{find}，只在确实要插入或累加时使用 \code{operator[]}。前缀和、区间和、计数乘法常需要 \code{long long}；模运算还要处理负余数归一化。如果追问变成“这是 meet-in-the-middle 思想，比四层枚举少两个维度”，先判断原来的状态是否仍然覆盖新答案集合，再决定是微调边界、改 value 语义，还是换成另一类结构。

\subsubsection{LeetCode 523 连续的子数组和}

先用双层循环枚举区间或候选对，再把条件写成历史状态和当前状态的关系。本题的模型是长度至少二的子数组和为 k 的倍数，等价于两个前缀和同余；暴力版的问题不是不会找，而是每个当前位置都回头扫太多历史。优化目标就是把哈希表保存某个余数最早出现下标，当前余数若重复且距离至少二则成立变成一次查表或一次状态读取。

哈希表保存的是当前下标之前的历史，不是随手放进去的缓存。本题需要哈希表保存某个余数最早出现下标，当前余数若重复且距离至少二则成立，所以要先决定 value 是次数、最早位置、最近位置还是布尔值。这个决定直接改变答案类型。放到“前缀和与哈希表”这条主线里看，关键原则是：前缀和把区间问题改写成两个前缀状态的差；哈希表把寻找匹配前缀改成查表。频次表、最早位置表、最近位置表对应不同目标：计数、最长区间、最短区间不能混用。

证明回到代数关系：任意合法答案都能对应到一个当前状态和一个历史状态；算法查到的历史状态也一定构成合法答案。本题的长度至少二的子数组和为 k 的倍数，等价于两个前缀和同余说明要找什么，哈希表保存某个余数最早出现下标，当前余数若重复且距离至少二则成立说明历史表怎样保存它。这道题的边界不能留到最后补丁式处理，实验时把k 为零、负数余数、长度限制、从开头开始的区间加入对拍集合，用平方暴力和优化版比较。失败时先看初始历史状态、查询更新顺序和哈希表 value 类型。

C++ 实现上，哈希表查询优先用 \code{find}，只在确实要插入或累加时使用 \code{operator[]}。前缀和、区间和、计数乘法常需要 \code{long long}；模运算还要处理负余数归一化。如果追问变成“保存最早下标是为了满足长度条件并保留最长距离”，先判断原来的状态是否仍然覆盖新答案集合，再决定是微调边界、改 value 语义，还是换成另一类结构。

\subsubsection{LeetCode 525 连续数组}

先用双层循环枚举区间或候选对，再把条件写成历史状态和当前状态的关系。本题的模型是把零看成负一后，零和一数量相等等价于前缀和相同；暴力版的问题不是不会找，而是每个当前位置都回头扫太多历史。优化目标就是把哈希表保存每个前缀和最早下标，重复出现时更新最长长度变成一次查表或一次状态读取。

哈希表保存的是当前下标之前的历史，不是随手放进去的缓存。本题需要哈希表保存每个前缀和最早下标，重复出现时更新最长长度，所以要先决定 value 是次数、最早位置、最近位置还是布尔值。这个决定直接改变答案类型。放到“前缀和与哈希表”这条主线里看，关键原则是：前缀和把区间问题改写成两个前缀状态的差；哈希表把寻找匹配前缀改成查表。频次表、最早位置表、最近位置表对应不同目标：计数、最长区间、最短区间不能混用。

证明回到代数关系：任意合法答案都能对应到一个当前状态和一个历史状态；算法查到的历史状态也一定构成合法答案。本题的把零看成负一后，零和一数量相等等价于前缀和相同说明要找什么，哈希表保存每个前缀和最早下标，重复出现时更新最长长度说明历史表怎样保存它。这道题的边界不能留到最后补丁式处理，实验时把全零、全一、从下标零开始、不要覆盖最早位置加入对拍集合，用平方暴力和优化版比较。失败时先看初始历史状态、查询更新顺序和哈希表 value 类型。

C++ 实现上，哈希表查询优先用 \code{find}，只在确实要插入或累加时使用 \code{operator[]}。前缀和、区间和、计数乘法常需要 \code{long long}；模运算还要处理负余数归一化。如果追问变成“这是前缀和最早位置语义，不是频次语义”，先判断原来的状态是否仍然覆盖新答案集合，再决定是微调边界、改 value 语义，还是换成另一类结构。

\subsubsection{LeetCode 560 和为 K 的子数组}

先用双层循环枚举区间或候选对，再把条件写成历史状态和当前状态的关系。本题的模型是当前前缀和减去目标值，就是需要匹配的历史前缀；暴力版的问题不是不会找，而是每个当前位置都回头扫太多历史。优化目标就是把哈希表保存前缀和出现次数，先查询再更新变成一次查表或一次状态读取。

哈希表保存的是当前下标之前的历史，不是随手放进去的缓存。本题需要哈希表保存前缀和出现次数，先查询再更新，所以要先决定 value 是次数、最早位置、最近位置还是布尔值。这个决定直接改变答案类型。放到“前缀和与哈希表”这条主线里看，关键原则是：前缀和把区间问题改写成两个前缀状态的差；哈希表把寻找匹配前缀改成查表。频次表、最早位置表、最近位置表对应不同目标：计数、最长区间、最短区间不能混用。

证明回到代数关系：任意合法答案都能对应到一个当前状态和一个历史状态；算法查到的历史状态也一定构成合法答案。本题的当前前缀和减去目标值，就是需要匹配的历史前缀说明要找什么，哈希表保存前缀和出现次数，先查询再更新说明历史表怎样保存它。这道题的边界不能留到最后补丁式处理，实验时把负数、目标为零、从下标零开始的区间、重复前缀加入对拍集合，用平方暴力和优化版比较。失败时先看初始历史状态、查询更新顺序和哈希表 value 类型。

C++ 实现上，哈希表查询优先用 \code{find}，只在确实要插入或累加时使用 \code{operator[]}。前缀和、区间和、计数乘法常需要 \code{long long}；模运算还要处理负余数归一化。如果追问变成“若要求最长区间，哈希表应保存最早位置而不是次数”，先判断原来的状态是否仍然覆盖新答案集合，再决定是微调边界、改 value 语义，还是换成另一类结构。

\subsubsection{LeetCode 974 和可被 K 整除的子数组}

先用双层循环枚举区间或候选对，再把条件写成历史状态和当前状态的关系。本题的模型是区间和能被 K 整除等价于两个前缀和余数相同；暴力版的问题不是不会找，而是每个当前位置都回头扫太多历史。优化目标就是把统计每个余数出现次数，当前余数能和所有历史同余前缀组成答案变成一次查表或一次状态读取。

哈希表保存的是当前下标之前的历史，不是随手放进去的缓存。本题需要统计每个余数出现次数，当前余数能和所有历史同余前缀组成答案，所以要先决定 value 是次数、最早位置、最近位置还是布尔值。这个决定直接改变答案类型。放到“前缀和与哈希表”这条主线里看，关键原则是：前缀和把区间问题改写成两个前缀状态的差；哈希表把寻找匹配前缀改成查表。频次表、最早位置表、最近位置表对应不同目标：计数、最长区间、最短区间不能混用。

证明回到代数关系：任意合法答案都能对应到一个当前状态和一个历史状态；算法查到的历史状态也一定构成合法答案。本题的区间和能被 K 整除等价于两个前缀和余数相同说明要找什么，统计每个余数出现次数，当前余数能和所有历史同余前缀组成答案说明历史表怎样保存它。这道题的边界不能留到最后补丁式处理，实验时把负数余数归一化、K 为一、从开头开始、重复余数很多加入对拍集合，用平方暴力和优化版比较。失败时先看初始历史状态、查询更新顺序和哈希表 value 类型。

C++ 实现上，哈希表查询优先用 \code{find}，只在确实要插入或累加时使用 \code{operator[]}。前缀和、区间和、计数乘法常需要 \code{long long}；模运算还要处理负余数归一化。如果追问变成“如果问最长长度，保存最早位置而不是次数”，先判断原来的状态是否仍然覆盖新答案集合，再决定是微调边界、改 value 语义，还是换成另一类结构。

\subsubsection{LeetCode 1074 元素和为目标值的子矩阵数量}

先用双层循环枚举区间或候选对，再把条件写成历史状态和当前状态的关系。本题的模型是二维矩阵子区域和可以固定上下边界后转成一维子数组和；暴力版的问题不是不会找，而是每个当前位置都回头扫太多历史。优化目标就是把枚举行边界，压缩每列区间和，再用前缀和哈希统计目标和子数组变成一次查表或一次状态读取。

哈希表保存的是当前下标之前的历史，不是随手放进去的缓存。本题需要枚举行边界，压缩每列区间和，再用前缀和哈希统计目标和子数组，所以要先决定 value 是次数、最早位置、最近位置还是布尔值。这个决定直接改变答案类型。放到“前缀和与哈希表”这条主线里看，关键原则是：前缀和把区间问题改写成两个前缀状态的差；哈希表把寻找匹配前缀改成查表。频次表、最早位置表、最近位置表对应不同目标：计数、最长区间、最短区间不能混用。

证明回到代数关系：任意合法答案都能对应到一个当前状态和一个历史状态；算法查到的历史状态也一定构成合法答案。本题的二维矩阵子区域和可以固定上下边界后转成一维子数组和说明要找什么，枚举行边界，压缩每列区间和，再用前缀和哈希统计目标和子数组说明历史表怎样保存它。这道题的边界不能留到最后补丁式处理，实验时把单行、单列、负数、目标为零、矩阵较大加入对拍集合，用平方暴力和优化版比较。失败时先看初始历史状态、查询更新顺序和哈希表 value 类型。

C++ 实现上，哈希表查询优先用 \code{find}，只在确实要插入或累加时使用 \code{operator[]}。前缀和、区间和、计数乘法常需要 \code{long long}；模运算还要处理负余数归一化。如果追问变成“这是二维前缀思想和一维哈希计数的组合”，先判断原来的状态是否仍然覆盖新答案集合，再决定是微调边界、改 value 语义，还是换成另一类结构。

\subsubsection{LeetCode 1248 统计优美子数组}

先用双层循环枚举区间或候选对，再把条件写成历史状态和当前状态的关系。本题的模型是恰好 K 个奇数可转成前缀奇数计数差为 K；暴力版的问题不是不会找，而是每个当前位置都回头扫太多历史。优化目标就是把哈希表保存某个奇数计数出现次数，当前计数查询 count - K变成一次查表或一次状态读取。

哈希表保存的是当前下标之前的历史，不是随手放进去的缓存。本题需要哈希表保存某个奇数计数出现次数，当前计数查询 count - K，所以要先决定 value 是次数、最早位置、最近位置还是布尔值。这个决定直接改变答案类型。放到“前缀和与哈希表”这条主线里看，关键原则是：前缀和把区间问题改写成两个前缀状态的差；哈希表把寻找匹配前缀改成查表。频次表、最早位置表、最近位置表对应不同目标：计数、最长区间、最短区间不能混用。

证明回到代数关系：任意合法答案都能对应到一个当前状态和一个历史状态；算法查到的历史状态也一定构成合法答案。本题的恰好 K 个奇数可转成前缀奇数计数差为 K说明要找什么，哈希表保存某个奇数计数出现次数，当前计数查询 count - K说明历史表怎样保存它。这道题的边界不能留到最后补丁式处理，实验时把K 为零、全偶数、全奇数、从开头开始的区间加入对拍集合，用平方暴力和优化版比较。失败时先看初始历史状态、查询更新顺序和哈希表 value 类型。

C++ 实现上，哈希表查询优先用 \code{find}，只在确实要插入或累加时使用 \code{operator[]}。前缀和、区间和、计数乘法常需要 \code{long long}；模运算还要处理负余数归一化。如果追问变成“也可用至多 K 减至多 K-1 的滑动窗口”，先判断原来的状态是否仍然覆盖新答案集合，再决定是微调边界、改 value 语义，还是换成另一类结构。

\subsection{二分查找与答案空间工作坊}

先线性扫描或线性试答案，再找出能把候选一分为二的单调谓词。二分的核心不是 mid 写法，而是被排除的一半确实不可能包含答案。

\subsubsection{LeetCode 33 搜索旋转排序数组}

先用线性扫描或线性试答案保证正确，再找边界。本题的模型是旋转数组每次二分至少有一半保持有序；二分不是为了套循环，而是因为判断哪一半有序，再看目标是否落在有序半边范围内给出了一条能排除候选的边界线。若这条边界线说不清，宁愿先保留线性版本。

二分的状态不是数组本身，而是候选区间。本题用判断哪一半有序，再看目标是否落在有序半边范围内缩小区间；每次更新左右边界时，都要能说清被丢掉的一半为什么没有答案。边界不变量比 mid 公式更重要。放到“二分查找与答案空间”这条主线里看，关键原则是：二分前必须先定义谓词：某个位置是否已经满足边界条件，或者某个答案是否可行。数组二分找的是边界，答案二分找的是最小可行值或最大可行值。

证明二分时，不要只说复杂度是对数。要证明当判断哪一半有序，再看目标是否落在有序半边范围内成立或不成立时，被排除的一侧确实没有目标。本题的旋转数组每次二分至少有一半保持有序提供候选空间，循环只是在这个空间里保留答案。这道题的边界不能留到最后补丁式处理，实验时覆盖未旋转、长度一、目标不存在、边界等号、存在重复元素时退化，逐轮打印 left、mid、right 和谓词值。只要有一次被排除区间仍含答案，二分不变量就错了。

C++ 实现上，中点写成 \code{left + (right - left) / 2}；平方、乘法和总量比较避免溢出。答案二分最好把检查函数单独写出，让单调谓词可读。如果追问变成“若重复元素很多，需要在无法判断有序半边时收缩边界”，先判断原来的状态是否仍然覆盖新答案集合，再决定是微调边界、改 value 语义，还是换成另一类结构。

\subsubsection{LeetCode 34 排序数组中查找首尾位置}

先用线性扫描或线性试答案保证正确，再找边界。本题的模型是首尾位置可以拆成两个边界：第一个大于等于和第一个大于；二分不是为了套循环，而是因为不要找到一个目标后向两边线性扩展，否则重复元素很多时退化给出了一条能排除候选的边界线。若这条边界线说不清，宁愿先保留线性版本。

二分的状态不是数组本身，而是候选区间。本题用不要找到一个目标后向两边线性扩展，否则重复元素很多时退化缩小区间；每次更新左右边界时，都要能说清被丢掉的一半为什么没有答案。边界不变量比 mid 公式更重要。放到“二分查找与答案空间”这条主线里看，关键原则是：二分前必须先定义谓词：某个位置是否已经满足边界条件，或者某个答案是否可行。数组二分找的是边界，答案二分找的是最小可行值或最大可行值。

证明二分时，不要只说复杂度是对数。要证明当不要找到一个目标后向两边线性扩展，否则重复元素很多时退化成立或不成立时，被排除的一侧确实没有目标。本题的首尾位置可以拆成两个边界：第一个大于等于和第一个大于提供候选空间，循环只是在这个空间里保留答案。这道题的边界不能留到最后补丁式处理，实验时覆盖目标不存在、目标在开头或结尾、数组为空、全是目标，逐轮打印 left、mid、right 和谓词值。只要有一次被排除区间仍含答案，二分不变量就错了。

C++ 实现上，中点写成 \code{left + (right - left) / 2}；平方、乘法和总量比较避免溢出。答案二分最好把检查函数单独写出，让单调谓词可读。如果追问变成“这个模型能迁移到计数某值出现次数和插入区间边界”，先判断原来的状态是否仍然覆盖新答案集合，再决定是微调边界、改 value 语义，还是换成另一类结构。

\subsubsection{LeetCode 35 搜索插入位置}

先用线性扫描或线性试答案保证正确，再找边界。本题的模型是题目等价于寻找第一个大于等于目标的位置；二分不是为了套循环，而是因为统一用 lower bound 语义，不在相等时提前返回也能保持边界一致给出了一条能排除候选的边界线。若这条边界线说不清，宁愿先保留线性版本。

二分的状态不是数组本身，而是候选区间。本题用统一用 lower bound 语义，不在相等时提前返回也能保持边界一致缩小区间；每次更新左右边界时，都要能说清被丢掉的一半为什么没有答案。边界不变量比 mid 公式更重要。放到“二分查找与答案空间”这条主线里看，关键原则是：二分前必须先定义谓词：某个位置是否已经满足边界条件，或者某个答案是否可行。数组二分找的是边界，答案二分找的是最小可行值或最大可行值。

证明二分时，不要只说复杂度是对数。要证明当统一用 lower bound 语义，不在相等时提前返回也能保持边界一致成立或不成立时，被排除的一侧确实没有目标。本题的题目等价于寻找第一个大于等于目标的位置提供候选空间，循环只是在这个空间里保留答案。这道题的边界不能留到最后补丁式处理，实验时覆盖目标小于所有元素、目标大于所有元素、重复值、空数组，逐轮打印 left、mid、right 和谓词值。只要有一次被排除区间仍含答案，二分不变量就错了。

C++ 实现上，中点写成 \code{left + (right - left) / 2}；平方、乘法和总量比较避免溢出。答案二分最好把检查函数单独写出，让单调谓词可读。如果追问变成“手写二分时用左闭右开可直接返回 left”，先判断原来的状态是否仍然覆盖新答案集合，再决定是微调边界、改 value 语义，还是换成另一类结构。

\subsubsection{LeetCode 69 x 的平方根}

先用线性扫描或线性试答案保证正确，再找边界。本题的模型是答案是最大的整数 r，使得 r 的平方不超过 x；二分不是为了套循环，而是因为在整数范围上二分右边界，比较时用除法或 long long 避免平方溢出给出了一条能排除候选的边界线。若这条边界线说不清，宁愿先保留线性版本。

二分的状态不是数组本身，而是候选区间。本题用在整数范围上二分右边界，比较时用除法或 long long 避免平方溢出缩小区间；每次更新左右边界时，都要能说清被丢掉的一半为什么没有答案。边界不变量比 mid 公式更重要。放到“二分查找与答案空间”这条主线里看，关键原则是：二分前必须先定义谓词：某个位置是否已经满足边界条件，或者某个答案是否可行。数组二分找的是边界，答案二分找的是最小可行值或最大可行值。

证明二分时，不要只说复杂度是对数。要证明当在整数范围上二分右边界，比较时用除法或 long long 避免平方溢出成立或不成立时，被排除的一侧确实没有目标。本题的答案是最大的整数 r，使得 r 的平方不超过 x提供候选空间，循环只是在这个空间里保留答案。这道题的边界不能留到最后补丁式处理，实验时覆盖x 为零、一、非完全平方、接近 int 最大值，逐轮打印 left、mid、right 和谓词值。只要有一次被排除区间仍含答案，二分不变量就错了。

C++ 实现上，中点写成 \code{left + (right - left) / 2}；平方、乘法和总量比较避免溢出。答案二分最好把检查函数单独写出，让单调谓词可读。如果追问变成“若要求浮点精度，可以二分实数区间或用牛顿迭代”，先判断原来的状态是否仍然覆盖新答案集合，再决定是微调边界、改 value 语义，还是换成另一类结构。

\subsubsection{LeetCode 74 搜索二维矩阵}

先用线性扫描或线性试答案保证正确，再找边界。本题的模型是矩阵每行有序且下一行首元素大于上一行末元素，可视为一维有序数组；二分不是为了套循环，而是因为把下标 mid 映射到 row 和 col，再做标准二分给出了一条能排除候选的边界线。若这条边界线说不清，宁愿先保留线性版本。

二分的状态不是数组本身，而是候选区间。本题用把下标 mid 映射到 row 和 col，再做标准二分缩小区间；每次更新左右边界时，都要能说清被丢掉的一半为什么没有答案。边界不变量比 mid 公式更重要。放到“二分查找与答案空间”这条主线里看，关键原则是：二分前必须先定义谓词：某个位置是否已经满足边界条件，或者某个答案是否可行。数组二分找的是边界，答案二分找的是最小可行值或最大可行值。

证明二分时，不要只说复杂度是对数。要证明当把下标 mid 映射到 row 和 col，再做标准二分成立或不成立时，被排除的一侧确实没有目标。本题的矩阵每行有序且下一行首元素大于上一行末元素，可视为一维有序数组提供候选空间，循环只是在这个空间里保留答案。这道题的边界不能留到最后补丁式处理，实验时覆盖空矩阵、单行、单列、目标小于全部或大于全部，逐轮打印 left、mid、right 和谓词值。只要有一次被排除区间仍含答案，二分不变量就错了。

C++ 实现上，中点写成 \code{left + (right - left) / 2}；平方、乘法和总量比较避免溢出。答案二分最好把检查函数单独写出，让单调谓词可读。如果追问变成“若只是行列分别有序但行之间不连续，应改用右上角排除法”，先判断原来的状态是否仍然覆盖新答案集合，再决定是微调边界、改 value 语义，还是换成另一类结构。

\subsubsection{LeetCode 81 搜索旋转排序数组 II}

先用线性扫描或线性试答案保证正确，再找边界。本题的模型是旋转数组存在重复值时，二分可能无法判断哪一半有序；二分不是为了套循环，而是因为当左右和中点值相同导致信息不足时，收缩边界；否则仍按有序半边排除给出了一条能排除候选的边界线。若这条边界线说不清，宁愿先保留线性版本。

二分的状态不是数组本身，而是候选区间。本题用当左右和中点值相同导致信息不足时，收缩边界；否则仍按有序半边排除缩小区间；每次更新左右边界时，都要能说清被丢掉的一半为什么没有答案。边界不变量比 mid 公式更重要。放到“二分查找与答案空间”这条主线里看，关键原则是：二分前必须先定义谓词：某个位置是否已经满足边界条件，或者某个答案是否可行。数组二分找的是边界，答案二分找的是最小可行值或最大可行值。

证明二分时，不要只说复杂度是对数。要证明当当左右和中点值相同导致信息不足时，收缩边界；否则仍按有序半边排除成立或不成立时，被排除的一侧确实没有目标。本题的旋转数组存在重复值时，二分可能无法判断哪一半有序提供候选空间，循环只是在这个空间里保留答案。这道题的边界不能留到最后补丁式处理，实验时覆盖全相同元素、目标不存在、重复值跨过旋转点、退化到线性，逐轮打印 left、mid、right 和谓词值。只要有一次被排除区间仍含答案，二分不变量就错了。

C++ 实现上，中点写成 \code{left + (right - left) / 2}；平方、乘法和总量比较避免溢出。答案二分最好把检查函数单独写出，让单调谓词可读。如果追问变成“这题说明二分依赖信息量，重复值会削弱每次排除一半的能力”，先判断原来的状态是否仍然覆盖新答案集合，再决定是微调边界、改 value 语义，还是换成另一类结构。

\subsubsection{LeetCode 153 寻找旋转排序数组中的最小值}

先用线性扫描或线性试答案保证正确，再找边界。本题的模型是最小值是旋转断点，右端点可帮助判断中点在哪一段；二分不是为了套循环，而是因为若 mid 大于 right，最小值在右侧；否则在左侧含 mid给出了一条能排除候选的边界线。若这条边界线说不清，宁愿先保留线性版本。

二分的状态不是数组本身，而是候选区间。本题用若 mid 大于 right，最小值在右侧；否则在左侧含 mid缩小区间；每次更新左右边界时，都要能说清被丢掉的一半为什么没有答案。边界不变量比 mid 公式更重要。放到“二分查找与答案空间”这条主线里看，关键原则是：二分前必须先定义谓词：某个位置是否已经满足边界条件，或者某个答案是否可行。数组二分找的是边界，答案二分找的是最小可行值或最大可行值。

证明二分时，不要只说复杂度是对数。要证明当若 mid 大于 right，最小值在右侧；否则在左侧含 mid成立或不成立时，被排除的一侧确实没有目标。本题的最小值是旋转断点，右端点可帮助判断中点在哪一段提供候选空间，循环只是在这个空间里保留答案。这道题的边界不能留到最后补丁式处理，实验时覆盖未旋转、长度一、旋转一位、无重复条件，逐轮打印 left、mid、right 和谓词值。只要有一次被排除区间仍含答案，二分不变量就错了。

C++ 实现上，中点写成 \code{left + (right - left) / 2}；平方、乘法和总量比较避免溢出。答案二分最好把检查函数单独写出，让单调谓词可读。如果追问变成“有重复时遇到相等只能收缩 right，最坏退化”，先判断原来的状态是否仍然覆盖新答案集合，再决定是微调边界、改 value 语义，还是换成另一类结构。

\subsubsection{LeetCode 154 寻找旋转排序数组中的最小值 II}

先用线性扫描或线性试答案保证正确，再找边界。本题的模型是重复元素会让 mid 和 right 相等时无法判断最小值方向；二分不是为了套循环，而是因为相等时只能安全地丢掉一个 right，因为它不是唯一必要候选给出了一条能排除候选的边界线。若这条边界线说不清，宁愿先保留线性版本。

二分的状态不是数组本身，而是候选区间。本题用相等时只能安全地丢掉一个 right，因为它不是唯一必要候选缩小区间；每次更新左右边界时，都要能说清被丢掉的一半为什么没有答案。边界不变量比 mid 公式更重要。放到“二分查找与答案空间”这条主线里看，关键原则是：二分前必须先定义谓词：某个位置是否已经满足边界条件，或者某个答案是否可行。数组二分找的是边界，答案二分找的是最小可行值或最大可行值。

证明二分时，不要只说复杂度是对数。要证明当相等时只能安全地丢掉一个 right，因为它不是唯一必要候选成立或不成立时，被排除的一侧确实没有目标。本题的重复元素会让 mid 和 right 相等时无法判断最小值方向提供候选空间，循环只是在这个空间里保留答案。这道题的边界不能留到最后补丁式处理，实验时覆盖全相等、重复值跨断点、最小值出现多次，逐轮打印 left、mid、right 和谓词值。只要有一次被排除区间仍含答案，二分不变量就错了。

C++ 实现上，中点写成 \code{left + (right - left) / 2}；平方、乘法和总量比较避免溢出。答案二分最好把检查函数单独写出，让单调谓词可读。如果追问变成“这题说明二分依赖的信息不足时会退化，不是所有旋转题都严格对数”，先判断原来的状态是否仍然覆盖新答案集合，再决定是微调边界、改 value 语义，还是换成另一类结构。

\subsubsection{LeetCode 162 寻找峰值}

先用线性扫描或线性试答案保证正确，再找边界。本题的模型是比较 mid 和 mid 加一 的斜率能判断某侧一定存在峰值；二分不是为了套循环，而是因为上升则右侧有峰，下降则左侧含 mid 有峰给出了一条能排除候选的边界线。若这条边界线说不清，宁愿先保留线性版本。

二分的状态不是数组本身，而是候选区间。本题用上升则右侧有峰，下降则左侧含 mid 有峰缩小区间；每次更新左右边界时，都要能说清被丢掉的一半为什么没有答案。边界不变量比 mid 公式更重要。放到“二分查找与答案空间”这条主线里看，关键原则是：二分前必须先定义谓词：某个位置是否已经满足边界条件，或者某个答案是否可行。数组二分找的是边界，答案二分找的是最小可行值或最大可行值。

证明二分时，不要只说复杂度是对数。要证明当上升则右侧有峰，下降则左侧含 mid 有峰成立或不成立时，被排除的一侧确实没有目标。本题的比较 mid 和 mid 加一 的斜率能判断某侧一定存在峰值提供候选空间，循环只是在这个空间里保留答案。这道题的边界不能留到最后补丁式处理，实验时覆盖长度一、峰在边界、多个峰、相邻不相等假设，逐轮打印 left、mid、right 和谓词值。只要有一次被排除区间仍含答案，二分不变量就错了。

C++ 实现上，中点写成 \code{left + (right - left) / 2}；平方、乘法和总量比较避免溢出。答案二分最好把检查函数单独写出，让单调谓词可读。如果追问变成“二维峰值需要按行列最大值进行更复杂二分”，先判断原来的状态是否仍然覆盖新答案集合，再决定是微调边界、改 value 语义，还是换成另一类结构。

\subsubsection{LeetCode 240 搜索二维矩阵 II}

先用线性扫描或线性试答案保证正确，再找边界。本题的模型是右上角同时具有向左变小、向下变大的排除能力；二分不是为了套循环，而是因为当前值大于目标就左移，小于目标就下移给出了一条能排除候选的边界线。若这条边界线说不清，宁愿先保留线性版本。

二分的状态不是数组本身，而是候选区间。本题用当前值大于目标就左移，小于目标就下移缩小区间；每次更新左右边界时，都要能说清被丢掉的一半为什么没有答案。边界不变量比 mid 公式更重要。放到“二分查找与答案空间”这条主线里看，关键原则是：二分前必须先定义谓词：某个位置是否已经满足边界条件，或者某个答案是否可行。数组二分找的是边界，答案二分找的是最小可行值或最大可行值。

证明二分时，不要只说复杂度是对数。要证明当当前值大于目标就左移，小于目标就下移成立或不成立时，被排除的一侧确实没有目标。本题的右上角同时具有向左变小、向下变大的排除能力提供候选空间，循环只是在这个空间里保留答案。这道题的边界不能留到最后补丁式处理，实验时覆盖空矩阵、单行、单列、目标在角落、重复值，逐轮打印 left、mid、right 和谓词值。只要有一次被排除区间仍含答案，二分不变量就错了。

C++ 实现上，中点写成 \code{left + (right - left) / 2}；平方、乘法和总量比较避免溢出。答案二分最好把检查函数单独写出，让单调谓词可读。如果追问变成“从左上角无法一次排除一行或一列”，先判断原来的状态是否仍然覆盖新答案集合，再决定是微调边界、改 value 语义，还是换成另一类结构。

\subsubsection{LeetCode 410 分割数组的最大值}

先用线性扫描或线性试答案保证正确，再找边界。本题的模型是连续分成 m 段后最大段和越大，越容易完成分割；二分不是为了套循环，而是因为二分最大段和上限，检查函数贪心累计，超过上限就新开一段给出了一条能排除候选的边界线。若这条边界线说不清，宁愿先保留线性版本。

二分的状态不是数组本身，而是候选区间。本题用二分最大段和上限，检查函数贪心累计，超过上限就新开一段缩小区间；每次更新左右边界时，都要能说清被丢掉的一半为什么没有答案。边界不变量比 mid 公式更重要。放到“二分查找与答案空间”这条主线里看，关键原则是：二分前必须先定义谓词：某个位置是否已经满足边界条件，或者某个答案是否可行。数组二分找的是边界，答案二分找的是最小可行值或最大可行值。

证明二分时，不要只说复杂度是对数。要证明当二分最大段和上限，检查函数贪心累计，超过上限就新开一段成立或不成立时，被排除的一侧确实没有目标。本题的连续分成 m 段后最大段和越大，越容易完成分割提供候选空间，循环只是在这个空间里保留答案。这道题的边界不能留到最后补丁式处理，实验时覆盖m 为一、m 等于数组长度、单个元素超过猜测上限、总和溢出，逐轮打印 left、mid、right 和谓词值。只要有一次被排除区间仍含答案，二分不变量就错了。

C++ 实现上，中点写成 \code{left + (right - left) / 2}；平方、乘法和总量比较避免溢出。答案二分最好把检查函数单独写出，让单调谓词可读。如果追问变成“也可用 DP 求精确最优，但答案二分更适合大输入”，先判断原来的状态是否仍然覆盖新答案集合，再决定是微调边界、改 value 语义，还是换成另一类结构。

\subsubsection{LeetCode 540 有序数组中的单一元素}

先用线性扫描或线性试答案保证正确，再找边界。本题的模型是除一个元素外其余元素都成对出现，单一元素会改变配对下标的奇偶规律；二分不是为了套循环，而是因为把 mid 调整到偶数位，若 nums[mid] 等于 nums[mid+1]，单一元素在右侧，否则在左侧给出了一条能排除候选的边界线。若这条边界线说不清，宁愿先保留线性版本。

二分的状态不是数组本身，而是候选区间。本题用把 mid 调整到偶数位，若 nums[mid] 等于 nums[mid+1]，单一元素在右侧，否则在左侧缩小区间；每次更新左右边界时，都要能说清被丢掉的一半为什么没有答案。边界不变量比 mid 公式更重要。放到“二分查找与答案空间”这条主线里看，关键原则是：二分前必须先定义谓词：某个位置是否已经满足边界条件，或者某个答案是否可行。数组二分找的是边界，答案二分找的是最小可行值或最大可行值。

证明二分时，不要只说复杂度是对数。要证明当把 mid 调整到偶数位，若 nums[mid] 等于 nums[mid+1]，单一元素在右侧，否则在左侧成立或不成立时，被排除的一侧确实没有目标。本题的除一个元素外其余元素都成对出现，单一元素会改变配对下标的奇偶规律提供候选空间，循环只是在这个空间里保留答案。这道题的边界不能留到最后补丁式处理，实验时覆盖单一元素在开头或结尾、数组长度一、mid 加一越界，逐轮打印 left、mid、right 和谓词值。只要有一次被排除区间仍含答案，二分不变量就错了。

C++ 实现上，中点写成 \code{left + (right - left) / 2}；平方、乘法和总量比较避免溢出。答案二分最好把检查函数单独写出，让单调谓词可读。如果追问变成“异或也能求值，但二分利用有序性达到对数时间”，先判断原来的状态是否仍然覆盖新答案集合，再决定是微调边界、改 value 语义，还是换成另一类结构。

\subsubsection{LeetCode 704 二分查找}

先用线性扫描或线性试答案保证正确，再找边界。本题的模型是基础题的重点是固定区间语义；二分不是为了套循环，而是因为左闭右开写法能统一 lower bound 和存在性检查给出了一条能排除候选的边界线。若这条边界线说不清，宁愿先保留线性版本。

二分的状态不是数组本身，而是候选区间。本题用左闭右开写法能统一 lower bound 和存在性检查缩小区间；每次更新左右边界时，都要能说清被丢掉的一半为什么没有答案。边界不变量比 mid 公式更重要。放到“二分查找与答案空间”这条主线里看，关键原则是：二分前必须先定义谓词：某个位置是否已经满足边界条件，或者某个答案是否可行。数组二分找的是边界，答案二分找的是最小可行值或最大可行值。

证明二分时，不要只说复杂度是对数。要证明当左闭右开写法能统一 lower bound 和存在性检查成立或不成立时，被排除的一侧确实没有目标。本题的基础题的重点是固定区间语义提供候选空间，循环只是在这个空间里保留答案。这道题的边界不能留到最后补丁式处理，实验时覆盖空数组、长度一、目标在边界、目标不存在，逐轮打印 left、mid、right 和谓词值。只要有一次被排除区间仍含答案，二分不变量就错了。

C++ 实现上，中点写成 \code{left + (right - left) / 2}；平方、乘法和总量比较避免溢出。答案二分最好把检查函数单独写出，让单调谓词可读。如果追问变成“所有复杂二分都应该能退回到这道题的边界不变量”，先判断原来的状态是否仍然覆盖新答案集合，再决定是微调边界、改 value 语义，还是换成另一类结构。

\subsubsection{LeetCode 875 爱吃香蕉的珂珂}

先用线性扫描或线性试答案保证正确，再找边界。本题的模型是速度越大，总耗时越小，答案具有单调可行性；二分不是为了套循环，而是因为二分速度，检查函数用向上取整累加小时数给出了一条能排除候选的边界线。若这条边界线说不清，宁愿先保留线性版本。

二分的状态不是数组本身，而是候选区间。本题用二分速度，检查函数用向上取整累加小时数缩小区间；每次更新左右边界时，都要能说清被丢掉的一半为什么没有答案。边界不变量比 mid 公式更重要。放到“二分查找与答案空间”这条主线里看，关键原则是：二分前必须先定义谓词：某个位置是否已经满足边界条件，或者某个答案是否可行。数组二分找的是边界，答案二分找的是最小可行值或最大可行值。

证明二分时，不要只说复杂度是对数。要证明当二分速度，检查函数用向上取整累加小时数成立或不成立时，被排除的一侧确实没有目标。本题的速度越大，总耗时越小，答案具有单调可行性提供候选空间，循环只是在这个空间里保留答案。这道题的边界不能留到最后补丁式处理，实验时覆盖速度下界、堆很大、h 等于堆数、乘法加法溢出，逐轮打印 left、mid、right 和谓词值。只要有一次被排除区间仍含答案，二分不变量就错了。

C++ 实现上，中点写成 \code{left + (right - left) / 2}；平方、乘法和总量比较避免溢出。答案二分最好把检查函数单独写出，让单调谓词可读。如果追问变成“船运包裹和分割数组都是同类最小可行上限”，先判断原来的状态是否仍然覆盖新答案集合，再决定是微调边界、改 value 语义，还是换成另一类结构。

\subsubsection{LeetCode 1011 在 D 天内送达包裹}

先用线性扫描或线性试答案保证正确，再找边界。本题的模型是容量越大，需要天数越少，存在最小可行容量；二分不是为了套循环，而是因为下界是最大包裹，上界是总重量，检查函数按顺序贪心装船给出了一条能排除候选的边界线。若这条边界线说不清，宁愿先保留线性版本。

二分的状态不是数组本身，而是候选区间。本题用下界是最大包裹，上界是总重量，检查函数按顺序贪心装船缩小区间；每次更新左右边界时，都要能说清被丢掉的一半为什么没有答案。边界不变量比 mid 公式更重要。放到“二分查找与答案空间”这条主线里看，关键原则是：二分前必须先定义谓词：某个位置是否已经满足边界条件，或者某个答案是否可行。数组二分找的是边界，答案二分找的是最小可行值或最大可行值。

证明二分时，不要只说复杂度是对数。要证明当下界是最大包裹，上界是总重量，检查函数按顺序贪心装船成立或不成立时，被排除的一侧确实没有目标。本题的容量越大，需要天数越少，存在最小可行容量提供候选空间，循环只是在这个空间里保留答案。这道题的边界不能留到最后补丁式处理，实验时覆盖D 为一、D 等于包裹数、单个包裹超过猜测容量，逐轮打印 left、mid、right 和谓词值。只要有一次被排除区间仍含答案，二分不变量就错了。

C++ 实现上，中点写成 \code{left + (right - left) / 2}；平方、乘法和总量比较避免溢出。答案二分最好把检查函数单独写出，让单调谓词可读。如果追问变成“它和分割数组最大值都是连续分段最大和模型”，先判断原来的状态是否仍然覆盖新答案集合，再决定是微调边界、改 value 语义，还是换成另一类结构。

\subsection{栈、单调栈与表达式工作坊}

先对每个位置向左或向右扫描，再识别还没有被解决的候选。栈保存等待未来元素解决的对象，新元素到来时一次性结算被它支配的一批旧候选。

\subsubsection{LeetCode 20 有效的括号}

先让每个位置向左或向右暴力寻找答案，观察哪些候选一直等到未来才被解决。本题的模型是括号匹配要求最近打开的左括号先被关闭；慢点集中在栈保存尚未匹配的左括号，遇到右括号必须和栈顶类型一致。栈的作用不是保存所有历史，而是保存仍未被当前输入解决的那一小段历史。

栈里应保存还没确定答案的对象。本题用栈保存尚未匹配的左括号，遇到右括号必须和栈顶类型一致触发结算；弹出时要说清当前元素是右边界、左边界还是运算触发点。栈里存值还是下标，要由距离、宽度和过期判断决定。放到“栈、单调栈与表达式”这条主线里看，关键原则是：当新元素到来时，它会一次性解决栈顶一串被它支配的旧元素。被弹出的元素以后再也不会参与比较，因此总体线性。

正确性来自弹出时机。一旦栈保存尚未匹配的左括号，遇到右括号必须和栈顶类型一致触发，栈顶对象的答案已经由当前边界和新的栈顶共同确定；未来元素不会再改变它。本题的括号匹配要求最近打开的左括号先被关闭说明栈中对象等待的是什么。这道题的边界不能留到最后补丁式处理，实验时准备空串、以右括号开头、交叉匹配、结束后栈非空，打印每次入栈、弹栈和结算对象。重点观察相等元素、空栈、哨兵和最后一次结算。

C++ 实现上，栈题多数时候用 \code{std::vector<int>} 保存下标，比适配器更方便查看栈顶和计算宽度。表达式题要注意左右操作数顺序，柱状图题要注意哨兵和宽度。如果追问变成“若加入通配符星号，需要额外维护可行区间或双栈”，先判断原来的状态是否仍然覆盖新答案集合，再决定是微调边界、改 value 语义，还是换成另一类结构。

\subsubsection{LeetCode 32 最长有效括号}

先让每个位置向左或向右暴力寻找答案，观察哪些候选一直等到未来才被解决。本题的模型是有效括号子串需要知道每段非法边界之后的最早可连接位置；慢点集中在栈保存下标，先放入哨兵负一；匹配后用当前下标减栈顶得到长度。栈的作用不是保存所有历史，而是保存仍未被当前输入解决的那一小段历史。

栈里应保存还没确定答案的对象。本题用栈保存下标，先放入哨兵负一；匹配后用当前下标减栈顶得到长度触发结算；弹出时要说清当前元素是右边界、左边界还是运算触发点。栈里存值还是下标，要由距离、宽度和过期判断决定。放到“栈、单调栈与表达式”这条主线里看，关键原则是：当新元素到来时，它会一次性解决栈顶一串被它支配的旧元素。被弹出的元素以后再也不会参与比较，因此总体线性。

正确性来自弹出时机。一旦栈保存下标，先放入哨兵负一；匹配后用当前下标减栈顶得到长度触发，栈顶对象的答案已经由当前边界和新的栈顶共同确定；未来元素不会再改变它。本题的有效括号子串需要知道每段非法边界之后的最早可连接位置说明栈中对象等待的是什么。这道题的边界不能留到最后补丁式处理，实验时准备空串、全左括号、全右括号、多个断开的合法段，打印每次入栈、弹栈和结算对象。重点观察相等元素、空栈、哨兵和最后一次结算。

C++ 实现上，栈题多数时候用 \code{std::vector<int>} 保存下标，比适配器更方便查看栈顶和计算宽度。表达式题要注意左右操作数顺序，柱状图题要注意哨兵和宽度。如果追问变成“DP 写法也可做，但栈写法更直接表达最近非法边界”，先判断原来的状态是否仍然覆盖新答案集合，再决定是微调边界、改 value 语义，还是换成另一类结构。

\subsubsection{LeetCode 42 接雨水}

先让每个位置向左或向右暴力寻找答案，观察哪些候选一直等到未来才被解决。本题的模型是每个位置的水由左侧最高和右侧最高的较小者决定；慢点集中在前后缀、双指针、单调栈都是不同状态维护方式；要能说明各自结算对象。栈的作用不是保存所有历史，而是保存仍未被当前输入解决的那一小段历史。

栈里应保存还没确定答案的对象。本题用前后缀、双指针、单调栈都是不同状态维护方式；要能说明各自结算对象触发结算；弹出时要说清当前元素是右边界、左边界还是运算触发点。栈里存值还是下标，要由距离、宽度和过期判断决定。放到“栈、单调栈与表达式”这条主线里看，关键原则是：当新元素到来时，它会一次性解决栈顶一串被它支配的旧元素。被弹出的元素以后再也不会参与比较，因此总体线性。

正确性来自弹出时机。一旦前后缀、双指针、单调栈都是不同状态维护方式；要能说明各自结算对象触发，栈顶对象的答案已经由当前边界和新的栈顶共同确定；未来元素不会再改变它。本题的每个位置的水由左侧最高和右侧最高的较小者决定说明栈中对象等待的是什么。这道题的边界不能留到最后补丁式处理，实验时准备单调递增、单调递减、平台高度、多个凹槽、数组长度不足三，打印每次入栈、弹栈和结算对象。重点观察相等元素、空栈、哨兵和最后一次结算。

C++ 实现上，栈题多数时候用 \code{std::vector<int>} 保存下标，比适配器更方便查看栈顶和计算宽度。表达式题要注意左右操作数顺序，柱状图题要注意哨兵和宽度。如果追问变成“若柱子变成二维网格，需要最小堆从边界向内扩展”，先判断原来的状态是否仍然覆盖新答案集合，再决定是微调边界、改 value 语义，还是换成另一类结构。

\subsubsection{LeetCode 84 柱状图中最大的矩形}

先让每个位置向左或向右暴力寻找答案，观察哪些候选一直等到未来才被解决。本题的模型是每根柱子作为最矮高度时，需要知道左右第一个更矮位置；慢点集中在单调递增栈在遇到更矮柱时结算被弹出柱子的最大宽度。栈的作用不是保存所有历史，而是保存仍未被当前输入解决的那一小段历史。

栈里应保存还没确定答案的对象。本题用单调递增栈在遇到更矮柱时结算被弹出柱子的最大宽度触发结算；弹出时要说清当前元素是右边界、左边界还是运算触发点。栈里存值还是下标，要由距离、宽度和过期判断决定。放到“栈、单调栈与表达式”这条主线里看，关键原则是：当新元素到来时，它会一次性解决栈顶一串被它支配的旧元素。被弹出的元素以后再也不会参与比较，因此总体线性。

正确性来自弹出时机。一旦单调递增栈在遇到更矮柱时结算被弹出柱子的最大宽度触发，栈顶对象的答案已经由当前边界和新的栈顶共同确定；未来元素不会再改变它。本题的每根柱子作为最矮高度时，需要知道左右第一个更矮位置说明栈中对象等待的是什么。这道题的边界不能留到最后补丁式处理，实验时准备递增、递减、相等高度、哨兵零、宽度计算，打印每次入栈、弹栈和结算对象。重点观察相等元素、空栈、哨兵和最后一次结算。

C++ 实现上，栈题多数时候用 \code{std::vector<int>} 保存下标，比适配器更方便查看栈顶和计算宽度。表达式题要注意左右操作数顺序，柱状图题要注意哨兵和宽度。如果追问变成“最大矩形可以作为二维矩阵最大矩形的行压缩子问题”，先判断原来的状态是否仍然覆盖新答案集合，再决定是微调边界、改 value 语义，还是换成另一类结构。

\subsubsection{LeetCode 85 最大矩形}

先让每个位置向左或向右暴力寻找答案，观察哪些候选一直等到未来才被解决。本题的模型是每一行都可以把连续的一当成柱状图高度；慢点集中在逐行更新每列高度，然后把当前行转成柱状图最大矩形问题。栈的作用不是保存所有历史，而是保存仍未被当前输入解决的那一小段历史。

栈里应保存还没确定答案的对象。本题用逐行更新每列高度，然后把当前行转成柱状图最大矩形问题触发结算；弹出时要说清当前元素是右边界、左边界还是运算触发点。栈里存值还是下标，要由距离、宽度和过期判断决定。放到“栈、单调栈与表达式”这条主线里看，关键原则是：当新元素到来时，它会一次性解决栈顶一串被它支配的旧元素。被弹出的元素以后再也不会参与比较，因此总体线性。

正确性来自弹出时机。一旦逐行更新每列高度，然后把当前行转成柱状图最大矩形问题触发，栈顶对象的答案已经由当前边界和新的栈顶共同确定；未来元素不会再改变它。本题的每一行都可以把连续的一当成柱状图高度说明栈中对象等待的是什么。这道题的边界不能留到最后补丁式处理，实验时准备全零、全一、单行、单列、宽度和高度结算，打印每次入栈、弹栈和结算对象。重点观察相等元素、空栈、哨兵和最后一次结算。

C++ 实现上，栈题多数时候用 \code{std::vector<int>} 保存下标，比适配器更方便查看栈顶和计算宽度。表达式题要注意左右操作数顺序，柱状图题要注意哨兵和宽度。如果追问变成“这是二维问题压成一维单调栈的典型做法”，先判断原来的状态是否仍然覆盖新答案集合，再决定是微调边界、改 value 语义，还是换成另一类结构。

\subsubsection{LeetCode 150 逆波兰表达式求值}

先让每个位置向左或向右暴力寻找答案，观察哪些候选一直等到未来才被解决。本题的模型是后缀表达式把优先级和括号都编码进 token 顺序；慢点集中在遇到数字入栈，遇到运算符弹出右操作数和左操作数计算。栈的作用不是保存所有历史，而是保存仍未被当前输入解决的那一小段历史。

栈里应保存还没确定答案的对象。本题用遇到数字入栈，遇到运算符弹出右操作数和左操作数计算触发结算；弹出时要说清当前元素是右边界、左边界还是运算触发点。栈里存值还是下标，要由距离、宽度和过期判断决定。放到“栈、单调栈与表达式”这条主线里看，关键原则是：当新元素到来时，它会一次性解决栈顶一串被它支配的旧元素。被弹出的元素以后再也不会参与比较，因此总体线性。

正确性来自弹出时机。一旦遇到数字入栈，遇到运算符弹出右操作数和左操作数计算触发，栈顶对象的答案已经由当前边界和新的栈顶共同确定；未来元素不会再改变它。本题的后缀表达式把优先级和括号都编码进 token 顺序说明栈中对象等待的是什么。这道题的边界不能留到最后补丁式处理，实验时准备负数 token、多位数、除法向零截断、减法顺序，打印每次入栈、弹栈和结算对象。重点观察相等元素、空栈、哨兵和最后一次结算。

C++ 实现上，栈题多数时候用 \code{std::vector<int>} 保存下标，比适配器更方便查看栈顶和计算宽度。表达式题要注意左右操作数顺序，柱状图题要注意哨兵和宽度。如果追问变成“中缀表达式求值需要另一个运算符栈处理优先级”，先判断原来的状态是否仍然覆盖新答案集合，再决定是微调边界、改 value 语义，还是换成另一类结构。

\subsubsection{LeetCode 155 最小栈}

先让每个位置向左或向右暴力寻找答案，观察哪些候选一直等到未来才被解决。本题的模型是普通栈只能看到栈顶，不能常数时间知道历史最小值；慢点集中在辅助栈同步保存每个深度对应的最小值。栈的作用不是保存所有历史，而是保存仍未被当前输入解决的那一小段历史。

栈里应保存还没确定答案的对象。本题用辅助栈同步保存每个深度对应的最小值触发结算；弹出时要说清当前元素是右边界、左边界还是运算触发点。栈里存值还是下标，要由距离、宽度和过期判断决定。放到“栈、单调栈与表达式”这条主线里看，关键原则是：当新元素到来时，它会一次性解决栈顶一串被它支配的旧元素。被弹出的元素以后再也不会参与比较，因此总体线性。

正确性来自弹出时机。一旦辅助栈同步保存每个深度对应的最小值触发，栈顶对象的答案已经由当前边界和新的栈顶共同确定；未来元素不会再改变它。本题的普通栈只能看到栈顶，不能常数时间知道历史最小值说明栈中对象等待的是什么。这道题的边界不能留到最后补丁式处理，实验时准备重复最小值、弹出最小值、空栈操作约定，打印每次入栈、弹栈和结算对象。重点观察相等元素、空栈、哨兵和最后一次结算。

C++ 实现上，栈题多数时候用 \code{std::vector<int>} 保存下标，比适配器更方便查看栈顶和计算宽度。表达式题要注意左右操作数顺序，柱状图题要注意哨兵和宽度。如果追问变成“也可保存差值压缩空间，但可读性和溢出处理更复杂”，先判断原来的状态是否仍然覆盖新答案集合，再决定是微调边界、改 value 语义，还是换成另一类结构。

\subsubsection{LeetCode 224 基本计算器}

先让每个位置向左或向右暴力寻找答案，观察哪些候选一直等到未来才被解决。本题的模型是括号会开启新的表达式上下文，括号结束时要回到上一层；慢点集中在栈保存进入括号前的结果和符号，当前层算完后乘以上层符号并合并。栈的作用不是保存所有历史，而是保存仍未被当前输入解决的那一小段历史。

栈里应保存还没确定答案的对象。本题用栈保存进入括号前的结果和符号，当前层算完后乘以上层符号并合并触发结算；弹出时要说清当前元素是右边界、左边界还是运算触发点。栈里存值还是下标，要由距离、宽度和过期判断决定。放到“栈、单调栈与表达式”这条主线里看，关键原则是：当新元素到来时，它会一次性解决栈顶一串被它支配的旧元素。被弹出的元素以后再也不会参与比较，因此总体线性。

正确性来自弹出时机。一旦栈保存进入括号前的结果和符号，当前层算完后乘以上层符号并合并触发，栈顶对象的答案已经由当前边界和新的栈顶共同确定；未来元素不会再改变它。本题的括号会开启新的表达式上下文，括号结束时要回到上一层说明栈中对象等待的是什么。这道题的边界不能留到最后补丁式处理，实验时准备多位数、前导空格、一元负号、嵌套括号，打印每次入栈、弹栈和结算对象。重点观察相等元素、空栈、哨兵和最后一次结算。

C++ 实现上，栈题多数时候用 \code{std::vector<int>} 保存下标，比适配器更方便查看栈顶和计算宽度。表达式题要注意左右操作数顺序，柱状图题要注意哨兵和宽度。如果追问变成“基本计算器 II 没有括号但有乘除优先级，需要不同的栈语义”，先判断原来的状态是否仍然覆盖新答案集合，再决定是微调边界、改 value 语义，还是换成另一类结构。

\subsubsection{LeetCode 227 基本计算器 II}

先让每个位置向左或向右暴力寻找答案，观察哪些候选一直等到未来才被解决。本题的模型是乘除优先级高于加减，可以在读到下一个运算符时结算上一项；慢点集中在栈保存已经确定符号的项，乘除直接修改栈顶，加减压入正负项。栈的作用不是保存所有历史，而是保存仍未被当前输入解决的那一小段历史。

栈里应保存还没确定答案的对象。本题用栈保存已经确定符号的项，乘除直接修改栈顶，加减压入正负项触发结算；弹出时要说清当前元素是右边界、左边界还是运算触发点。栈里存值还是下标，要由距离、宽度和过期判断决定。放到“栈、单调栈与表达式”这条主线里看，关键原则是：当新元素到来时，它会一次性解决栈顶一串被它支配的旧元素。被弹出的元素以后再也不会参与比较，因此总体线性。

正确性来自弹出时机。一旦栈保存已经确定符号的项，乘除直接修改栈顶，加减压入正负项触发，栈顶对象的答案已经由当前边界和新的栈顶共同确定；未来元素不会再改变它。本题的乘除优先级高于加减，可以在读到下一个运算符时结算上一项说明栈中对象等待的是什么。这道题的边界不能留到最后补丁式处理，实验时准备多位数、空格、除法向零截断、最后一个数字需要结算，打印每次入栈、弹栈和结算对象。重点观察相等元素、空栈、哨兵和最后一次结算。

C++ 实现上，栈题多数时候用 \code{std::vector<int>} 保存下标，比适配器更方便查看栈顶和计算宽度。表达式题要注意左右操作数顺序，柱状图题要注意哨兵和宽度。如果追问变成“若加入括号，需要再引入上下文栈或递归下降解析”，先判断原来的状态是否仍然覆盖新答案集合，再决定是微调边界、改 value 语义，还是换成另一类结构。

\subsubsection{LeetCode 232 用栈实现队列}

先让每个位置向左或向右暴力寻找答案，观察哪些候选一直等到未来才被解决。本题的模型是两个栈分别负责输入顺序和输出顺序；慢点集中在只有输出栈为空时，才把输入栈整体倒过去，保证均摊常数。栈的作用不是保存所有历史，而是保存仍未被当前输入解决的那一小段历史。

栈里应保存还没确定答案的对象。本题用只有输出栈为空时，才把输入栈整体倒过去，保证均摊常数触发结算；弹出时要说清当前元素是右边界、左边界还是运算触发点。栈里存值还是下标，要由距离、宽度和过期判断决定。放到“栈、单调栈与表达式”这条主线里看，关键原则是：当新元素到来时，它会一次性解决栈顶一串被它支配的旧元素。被弹出的元素以后再也不会参与比较，因此总体线性。

正确性来自弹出时机。一旦只有输出栈为空时，才把输入栈整体倒过去，保证均摊常数触发，栈顶对象的答案已经由当前边界和新的栈顶共同确定；未来元素不会再改变它。本题的两个栈分别负责输入顺序和输出顺序说明栈中对象等待的是什么。这道题的边界不能留到最后补丁式处理，实验时准备连续 peek、空队列、交替 push pop，打印每次入栈、弹栈和结算对象。重点观察相等元素、空栈、哨兵和最后一次结算。

C++ 实现上，栈题多数时候用 \code{std::vector<int>} 保存下标，比适配器更方便查看栈顶和计算宽度。表达式题要注意左右操作数顺序，柱状图题要注意哨兵和宽度。如果追问变成“用队列实现栈则要通过轮转保持最近元素在队首”，先判断原来的状态是否仍然覆盖新答案集合，再决定是微调边界、改 value 语义，还是换成另一类结构。

\subsubsection{LeetCode 239 滑动窗口最大值}

先让每个位置向左或向右暴力寻找答案，观察哪些候选一直等到未来才被解决。本题的模型是每个窗口最大值来自仍未过期且没有被新元素支配的候选；慢点集中在单调队列保存下标，队尾小于等于新值的候选被弹出。栈的作用不是保存所有历史，而是保存仍未被当前输入解决的那一小段历史。

栈里应保存还没确定答案的对象。本题用单调队列保存下标，队尾小于等于新值的候选被弹出触发结算；弹出时要说清当前元素是右边界、左边界还是运算触发点。栈里存值还是下标，要由距离、宽度和过期判断决定。放到“栈、单调栈与表达式”这条主线里看，关键原则是：当新元素到来时，它会一次性解决栈顶一串被它支配的旧元素。被弹出的元素以后再也不会参与比较，因此总体线性。

正确性来自弹出时机。一旦单调队列保存下标，队尾小于等于新值的候选被弹出触发，栈顶对象的答案已经由当前边界和新的栈顶共同确定；未来元素不会再改变它。本题的每个窗口最大值来自仍未过期且没有被新元素支配的候选说明栈中对象等待的是什么。这道题的边界不能留到最后补丁式处理，实验时准备k 为一、k 等于数组长度、重复最大值、队首过期，打印每次入栈、弹栈和结算对象。重点观察相等元素、空栈、哨兵和最后一次结算。

C++ 实现上，栈题多数时候用 \code{std::vector<int>} 保存下标，比适配器更方便查看栈顶和计算宽度。表达式题要注意左右操作数顺序，柱状图题要注意哨兵和宽度。如果追问变成“受限子序列和使用同样队列维护最近 k 个 DP 最大值”，先判断原来的状态是否仍然覆盖新答案集合，再决定是微调边界、改 value 语义，还是换成另一类结构。

\subsubsection{LeetCode 394 字符串解码}

先让每个位置向左或向右暴力寻找答案，观察哪些候选一直等到未来才被解决。本题的模型是嵌套结构需要保存上一层字符串和重复次数；慢点集中在遇到左括号入栈当前状态，右括号弹出并拼接展开。栈的作用不是保存所有历史，而是保存仍未被当前输入解决的那一小段历史。

栈里应保存还没确定答案的对象。本题用遇到左括号入栈当前状态，右括号弹出并拼接展开触发结算；弹出时要说清当前元素是右边界、左边界还是运算触发点。栈里存值还是下标，要由距离、宽度和过期判断决定。放到“栈、单调栈与表达式”这条主线里看，关键原则是：当新元素到来时，它会一次性解决栈顶一串被它支配的旧元素。被弹出的元素以后再也不会参与比较，因此总体线性。

正确性来自弹出时机。一旦遇到左括号入栈当前状态，右括号弹出并拼接展开触发，栈顶对象的答案已经由当前边界和新的栈顶共同确定；未来元素不会再改变它。本题的嵌套结构需要保存上一层字符串和重复次数说明栈中对象等待的是什么。这道题的边界不能留到最后补丁式处理，实验时准备多位数字、嵌套、空片段、数字后必须跟括号，打印每次入栈、弹栈和结算对象。重点观察相等元素、空栈、哨兵和最后一次结算。

C++ 实现上，栈题多数时候用 \code{std::vector<int>} 保存下标，比适配器更方便查看栈顶和计算宽度。表达式题要注意左右操作数顺序，柱状图题要注意哨兵和宽度。如果追问变成“递归下降解析更通用，但迭代栈更符合工程栈控制”，先判断原来的状态是否仍然覆盖新答案集合，再决定是微调边界、改 value 语义，还是换成另一类结构。

\subsubsection{LeetCode 402 移掉 K 位数字}

先让每个位置向左或向右暴力寻找答案，观察哪些候选一直等到未来才被解决。本题的模型是要删除 k 位使剩余数字最小，本质是让高位尽量小；慢点集中在单调递增栈中遇到更小数字时，弹出前面较大的数字并消耗删除次数。栈的作用不是保存所有历史，而是保存仍未被当前输入解决的那一小段历史。

栈里应保存还没确定答案的对象。本题用单调递增栈中遇到更小数字时，弹出前面较大的数字并消耗删除次数触发结算；弹出时要说清当前元素是右边界、左边界还是运算触发点。栈里存值还是下标，要由距离、宽度和过期判断决定。放到“栈、单调栈与表达式”这条主线里看，关键原则是：当新元素到来时，它会一次性解决栈顶一串被它支配的旧元素。被弹出的元素以后再也不会参与比较，因此总体线性。

正确性来自弹出时机。一旦单调递增栈中遇到更小数字时，弹出前面较大的数字并消耗删除次数触发，栈顶对象的答案已经由当前边界和新的栈顶共同确定；未来元素不会再改变它。本题的要删除 k 位使剩余数字最小，本质是让高位尽量小说明栈中对象等待的是什么。这道题的边界不能留到最后补丁式处理，实验时准备前导零、k 未用完、k 等于长度、数字已经递增，打印每次入栈、弹栈和结算对象。重点观察相等元素、空栈、哨兵和最后一次结算。

C++ 实现上，栈题多数时候用 \code{std::vector<int>} 保存下标，比适配器更方便查看栈顶和计算宽度。表达式题要注意左右操作数顺序，柱状图题要注意哨兵和宽度。如果追问变成“去除重复字母也用单调栈，但还要保证每个字符保留一次”，先判断原来的状态是否仍然覆盖新答案集合，再决定是微调边界、改 value 语义，还是换成另一类结构。

\subsubsection{LeetCode 496 下一个更大元素 I}

先让每个位置向左或向右暴力寻找答案，观察哪些候选一直等到未来才被解决。本题的模型是第二个数组提供每个值右侧第一个更大值的全局关系；慢点集中在单调栈预处理值到下一个更大值映射，再回答第一个数组查询。栈的作用不是保存所有历史，而是保存仍未被当前输入解决的那一小段历史。

栈里应保存还没确定答案的对象。本题用单调栈预处理值到下一个更大值映射，再回答第一个数组查询触发结算；弹出时要说清当前元素是右边界、左边界还是运算触发点。栈里存值还是下标，要由距离、宽度和过期判断决定。放到“栈、单调栈与表达式”这条主线里看，关键原则是：当新元素到来时，它会一次性解决栈顶一串被它支配的旧元素。被弹出的元素以后再也不会参与比较，因此总体线性。

正确性来自弹出时机。一旦单调栈预处理值到下一个更大值映射，再回答第一个数组查询触发，栈顶对象的答案已经由当前边界和新的栈顶共同确定；未来元素不会再改变它。本题的第二个数组提供每个值右侧第一个更大值的全局关系说明栈中对象等待的是什么。这道题的边界不能留到最后补丁式处理，实验时准备不存在更大值、递减数组、值唯一假设，打印每次入栈、弹栈和结算对象。重点观察相等元素、空栈、哨兵和最后一次结算。

C++ 实现上，栈题多数时候用 \code{std::vector<int>} 保存下标，比适配器更方便查看栈顶和计算宽度。表达式题要注意左右操作数顺序，柱状图题要注意哨兵和宽度。如果追问变成“若值不唯一，应按下标而不是值建映射”，先判断原来的状态是否仍然覆盖新答案集合，再决定是微调边界、改 value 语义，还是换成另一类结构。

\subsubsection{LeetCode 503 下一个更大元素 II}

先让每个位置向左或向右暴力寻找答案，观察哪些候选一直等到未来才被解决。本题的模型是环形数组可以通过遍历两倍下标模拟绕回；慢点集中在第一遍入栈，第二遍只帮助未解决下标找到答案。栈的作用不是保存所有历史，而是保存仍未被当前输入解决的那一小段历史。

栈里应保存还没确定答案的对象。本题用第一遍入栈，第二遍只帮助未解决下标找到答案触发结算；弹出时要说清当前元素是右边界、左边界还是运算触发点。栈里存值还是下标，要由距离、宽度和过期判断决定。放到“栈、单调栈与表达式”这条主线里看，关键原则是：当新元素到来时，它会一次性解决栈顶一串被它支配的旧元素。被弹出的元素以后再也不会参与比较，因此总体线性。

正确性来自弹出时机。一旦第一遍入栈，第二遍只帮助未解决下标找到答案触发，栈顶对象的答案已经由当前边界和新的栈顶共同确定；未来元素不会再改变它。本题的环形数组可以通过遍历两倍下标模拟绕回说明栈中对象等待的是什么。这道题的边界不能留到最后补丁式处理，实验时准备全递减、全相等、长度一、取模下标，打印每次入栈、弹栈和结算对象。重点观察相等元素、空栈、哨兵和最后一次结算。

C++ 实现上，栈题多数时候用 \code{std::vector<int>} 保存下标，比适配器更方便查看栈顶和计算宽度。表达式题要注意左右操作数顺序，柱状图题要注意哨兵和宽度。如果追问变成“不要在第二遍重复入栈，否则可能死循环或重复计算”，先判断原来的状态是否仍然覆盖新答案集合，再决定是微调边界、改 value 语义，还是换成另一类结构。

\subsubsection{LeetCode 739 每日温度}

先让每个位置向左或向右暴力寻找答案，观察哪些候选一直等到未来才被解决。本题的模型是每一天等待右侧第一个更高温度；慢点集中在单调递减栈保存尚未找到答案的下标。栈的作用不是保存所有历史，而是保存仍未被当前输入解决的那一小段历史。

栈里应保存还没确定答案的对象。本题用单调递减栈保存尚未找到答案的下标触发结算；弹出时要说清当前元素是右边界、左边界还是运算触发点。栈里存值还是下标，要由距离、宽度和过期判断决定。放到“栈、单调栈与表达式”这条主线里看，关键原则是：当新元素到来时，它会一次性解决栈顶一串被它支配的旧元素。被弹出的元素以后再也不会参与比较，因此总体线性。

正确性来自弹出时机。一旦单调递减栈保存尚未找到答案的下标触发，栈顶对象的答案已经由当前边界和新的栈顶共同确定；未来元素不会再改变它。本题的每一天等待右侧第一个更高温度说明栈中对象等待的是什么。这道题的边界不能留到最后补丁式处理，实验时准备没有更高温、相等温度、递增、递减，打印每次入栈、弹栈和结算对象。重点观察相等元素、空栈、哨兵和最后一次结算。

C++ 实现上，栈题多数时候用 \code{std::vector<int>} 保存下标，比适配器更方便查看栈顶和计算宽度。表达式题要注意左右操作数顺序，柱状图题要注意哨兵和宽度。如果追问变成“下一个更大元素与它同型，只是输出值或距离不同”，先判断原来的状态是否仍然覆盖新答案集合，再决定是微调边界、改 value 语义，还是换成另一类结构。

\subsubsection{LeetCode 895 最大频率栈}

先让每个位置向左或向右暴力寻找答案，观察哪些候选一直等到未来才被解决。本题的模型是弹出时既要频率最高，又要在同频中最近入栈；慢点集中在哈希表记录值频率，频率到栈记录达到该频率的入栈顺序，维护当前最大频率。栈的作用不是保存所有历史，而是保存仍未被当前输入解决的那一小段历史。

栈里应保存还没确定答案的对象。本题用哈希表记录值频率，频率到栈记录达到该频率的入栈顺序，维护当前最大频率触发结算；弹出时要说清当前元素是右边界、左边界还是运算触发点。栈里存值还是下标，要由距离、宽度和过期判断决定。放到“栈、单调栈与表达式”这条主线里看，关键原则是：当新元素到来时，它会一次性解决栈顶一串被它支配的旧元素。被弹出的元素以后再也不会参与比较，因此总体线性。

正确性来自弹出时机。一旦哈希表记录值频率，频率到栈记录达到该频率的入栈顺序，维护当前最大频率触发，栈顶对象的答案已经由当前边界和新的栈顶共同确定；未来元素不会再改变它。本题的弹出时既要频率最高，又要在同频中最近入栈说明栈中对象等待的是什么。这道题的边界不能留到最后补丁式处理，实验时准备重复 push、pop 后最大频率下降、同频最近性、空栈不在题意内，打印每次入栈、弹栈和结算对象。重点观察相等元素、空栈、哨兵和最后一次结算。

C++ 实现上，栈题多数时候用 \code{std::vector<int>} 保存下标，比适配器更方便查看栈顶和计算宽度。表达式题要注意左右操作数顺序，柱状图题要注意哨兵和宽度。如果追问变成“它是栈和哈希表按频率维度组合的设计题”，先判断原来的状态是否仍然覆盖新答案集合，再决定是微调边界、改 value 语义，还是换成另一类结构。

\subsection{堆与动态极值工作坊}

先每轮扫描所有候选找极值，再把动态极值查询交给堆。堆只保证堆顶，不保证全局遍历有序；有过期元素时要在使用堆顶前清理。

\subsubsection{LeetCode 23 合并 K 个升序链表}

先每轮扫描所有候选找极值，确认答案选择的对象。本题的模型是每条链表当前头节点都是可能的下一个最小元素；暴力版慢在动态候选集合不断变化，却每次重新找最值。本题的优化点是小顶堆保存每条非空链的当前头，弹出最小节点后把它的后继入堆，因此堆中只能放足以比较下一步选择的轻量状态。

堆只保证堆顶，不保证内部全局有序。本题的小顶堆保存每条非空链的当前头，弹出最小节点后把它的后继入堆要转成堆顶可回答的问题；如果元素会过期，还要在真正使用堆顶前清理或跳过旧状态。放到“堆与动态极值”这条主线里看，关键原则是：堆把插入和弹出极值控制在对数复杂度。Top K 用固定大小堆维护边界，合并多路序列用堆保存每一路当前头部，数据流中位数用双堆维护两半。

证明堆题时，要说明堆覆盖了所有可能的下一步候选，并且堆顶过期时不会被使用。本题的小顶堆保存每条非空链的当前头，弹出最小节点后把它的后继入堆是选择顺序，每条链表当前头节点都是可能的下一个最小元素是候选集合来源。这道题的边界不能留到最后补丁式处理，实验时覆盖空链表数组、某条链为空、重复值、堆里保存节点指针还是值，记录堆顶是否过期、比较器是否让正确候选在堆顶、K 或窗口大小为边界值时是否稳定。

C++ 实现上，\code{std::priority_queue} 默认大顶堆，小顶堆需要比较器。堆中保存大对象时尽量保存下标或轻量记录；若不能删除任意元素，就用懒删除或过期检查。如果追问变成“也可以分治两两合并，堆更适合 K 较大且链长不均的场景”，先判断原来的状态是否仍然覆盖新答案集合，再决定是微调边界、改 value 语义，还是换成另一类结构。

\subsubsection{LeetCode 215 数组中的第 K 个最大元素}

先每轮扫描所有候选找极值，确认答案选择的对象。本题的模型是只关心第 K 大，不需要完整排序；暴力版慢在动态候选集合不断变化，却每次重新找最值。本题的优化点是大小为 K 的小顶堆保存当前前 K 大边界；快速选择可做到平均线性，因此堆中只能放足以比较下一步选择的轻量状态。

堆只保证堆顶，不保证内部全局有序。本题的大小为 K 的小顶堆保存当前前 K 大边界；快速选择可做到平均线性要转成堆顶可回答的问题；如果元素会过期，还要在真正使用堆顶前清理或跳过旧状态。放到“堆与动态极值”这条主线里看，关键原则是：堆把插入和弹出极值控制在对数复杂度。Top K 用固定大小堆维护边界，合并多路序列用堆保存每一路当前头部，数据流中位数用双堆维护两半。

证明堆题时，要说明堆覆盖了所有可能的下一步候选，并且堆顶过期时不会被使用。本题的大小为 K 的小顶堆保存当前前 K 大边界；快速选择可做到平均线性是选择顺序，只关心第 K 大，不需要完整排序是候选集合来源。这道题的边界不能留到最后补丁式处理，实验时覆盖K 为一、K 为 n、重复值、随机 pivot，记录堆顶是否过期、比较器是否让正确候选在堆顶、K 或窗口大小为边界值时是否稳定。

C++ 实现上，\code{std::priority_queue} 默认大顶堆，小顶堆需要比较器。堆中保存大对象时尽量保存下标或轻量记录；若不能删除任意元素，就用懒删除或过期检查。如果追问变成“若数据流持续到来，堆方案比一次性快速选择更自然”，先判断原来的状态是否仍然覆盖新答案集合，再决定是微调边界、改 value 语义，还是换成另一类结构。

\subsubsection{LeetCode 295 数据流的中位数}

先每轮扫描所有候选找极值，确认答案选择的对象。本题的模型是数据流需要动态维护较小一半和较大一半；暴力版慢在动态候选集合不断变化，却每次重新找最值。本题的优化点是大顶堆保存较小一半，小顶堆保存较大一半，大小差不超过一，因此堆中只能放足以比较下一步选择的轻量状态。

堆只保证堆顶，不保证内部全局有序。本题的大顶堆保存较小一半，小顶堆保存较大一半，大小差不超过一要转成堆顶可回答的问题；如果元素会过期，还要在真正使用堆顶前清理或跳过旧状态。放到“堆与动态极值”这条主线里看，关键原则是：堆把插入和弹出极值控制在对数复杂度。Top K 用固定大小堆维护边界，合并多路序列用堆保存每一路当前头部，数据流中位数用双堆维护两半。

证明堆题时，要说明堆覆盖了所有可能的下一步候选，并且堆顶过期时不会被使用。本题的大顶堆保存较小一半，小顶堆保存较大一半，大小差不超过一是选择顺序，数据流需要动态维护较小一半和较大一半是候选集合来源。这道题的边界不能留到最后补丁式处理，实验时覆盖偶数个元素、负数、重复值、平均值溢出，记录堆顶是否过期、比较器是否让正确候选在堆顶、K 或窗口大小为边界值时是否稳定。

C++ 实现上，\code{std::priority_queue} 默认大顶堆，小顶堆需要比较器。堆中保存大对象时尽量保存下标或轻量记录；若不能删除任意元素，就用懒删除或过期检查。如果追问变成“若还要删除任意元素，需要延迟删除哈希表”，先判断原来的状态是否仍然覆盖新答案集合，再决定是微调边界、改 value 语义，还是换成另一类结构。

\subsubsection{LeetCode 347 前 K 个高频元素}

先每轮扫描所有候选找极值，确认答案选择的对象。本题的模型是先统计频次，再只维护前 K 个候选；暴力版慢在动态候选集合不断变化，却每次重新找最值。本题的优化点是小顶堆大小超过 K 就弹出最低频；桶排序利用频次上界，因此堆中只能放足以比较下一步选择的轻量状态。

堆只保证堆顶，不保证内部全局有序。本题的小顶堆大小超过 K 就弹出最低频；桶排序利用频次上界要转成堆顶可回答的问题；如果元素会过期，还要在真正使用堆顶前清理或跳过旧状态。放到“堆与动态极值”这条主线里看，关键原则是：堆把插入和弹出极值控制在对数复杂度。Top K 用固定大小堆维护边界，合并多路序列用堆保存每一路当前头部，数据流中位数用双堆维护两半。

证明堆题时，要说明堆覆盖了所有可能的下一步候选，并且堆顶过期时不会被使用。本题的小顶堆大小超过 K 就弹出最低频；桶排序利用频次上界是选择顺序，先统计频次，再只维护前 K 个候选是候选集合来源。这道题的边界不能留到最后补丁式处理，实验时覆盖K 等于不同元素数、频次相同、结果顺序不要求，记录堆顶是否过期、比较器是否让正确候选在堆顶、K 或窗口大小为边界值时是否稳定。

C++ 实现上，\code{std::priority_queue} 默认大顶堆，小顶堆需要比较器。堆中保存大对象时尽量保存下标或轻量记录；若不能删除任意元素，就用懒删除或过期检查。如果追问变成“若要按频次和字典序排序，比较器要同时处理两个键”，先判断原来的状态是否仍然覆盖新答案集合，再决定是微调边界、改 value 语义，还是换成另一类结构。

\subsubsection{LeetCode 480 滑动窗口中位数}

先每轮扫描所有候选找极值，确认答案选择的对象。本题的模型是窗口移动时既要加入新元素，也要删除滑出的旧元素并查询中位数；暴力版慢在动态候选集合不断变化，却每次重新找最值。本题的优化点是双堆维护左右两半，延迟删除表记录已离窗但尚未到堆顶的元素，因此堆中只能放足以比较下一步选择的轻量状态。

堆只保证堆顶，不保证内部全局有序。本题的双堆维护左右两半，延迟删除表记录已离窗但尚未到堆顶的元素要转成堆顶可回答的问题；如果元素会过期，还要在真正使用堆顶前清理或跳过旧状态。放到“堆与动态极值”这条主线里看，关键原则是：堆把插入和弹出极值控制在对数复杂度。Top K 用固定大小堆维护边界，合并多路序列用堆保存每一路当前头部，数据流中位数用双堆维护两半。

证明堆题时，要说明堆覆盖了所有可能的下一步候选，并且堆顶过期时不会被使用。本题的双堆维护左右两半，延迟删除表记录已离窗但尚未到堆顶的元素是选择顺序，窗口移动时既要加入新元素，也要删除滑出的旧元素并查询中位数是候选集合来源。这道题的边界不能留到最后补丁式处理，实验时覆盖重复值、偶数窗口、平均值溢出、堆顶过期，记录堆顶是否过期、比较器是否让正确候选在堆顶、K 或窗口大小为边界值时是否稳定。

C++ 实现上，\code{std::priority_queue} 默认大顶堆，小顶堆需要比较器。堆中保存大对象时尽量保存下标或轻量记录；若不能删除任意元素，就用懒删除或过期检查。如果追问变成“普通 priority\_queue 不能任意删除，所以必须解释懒删除”，先判断原来的状态是否仍然覆盖新答案集合，再决定是微调边界、改 value 语义，还是换成另一类结构。

\subsubsection{LeetCode 621 任务调度器}

先每轮扫描所有候选找极值，确认答案选择的对象。本题的模型是相同任务之间有冷却间隔，核心受最高频任务骨架限制；暴力版慢在动态候选集合不断变化，却每次重新找最值。本题的优化点是可用大顶堆模拟，也可用最高频公式计算最短时间，因此堆中只能放足以比较下一步选择的轻量状态。

堆只保证堆顶，不保证内部全局有序。本题的可用大顶堆模拟，也可用最高频公式计算最短时间要转成堆顶可回答的问题；如果元素会过期，还要在真正使用堆顶前清理或跳过旧状态。放到“堆与动态极值”这条主线里看，关键原则是：堆把插入和弹出极值控制在对数复杂度。Top K 用固定大小堆维护边界，合并多路序列用堆保存每一路当前头部，数据流中位数用双堆维护两半。

证明堆题时，要说明堆覆盖了所有可能的下一步候选，并且堆顶过期时不会被使用。本题的可用大顶堆模拟，也可用最高频公式计算最短时间是选择顺序，相同任务之间有冷却间隔，核心受最高频任务骨架限制是候选集合来源。这道题的边界不能留到最后补丁式处理，实验时覆盖冷却为零、多个最高频、空闲被其他任务填满，记录堆顶是否过期、比较器是否让正确候选在堆顶、K 或窗口大小为边界值时是否稳定。

C++ 实现上，\code{std::priority_queue} 默认大顶堆，小顶堆需要比较器。堆中保存大对象时尽量保存下标或轻量记录；若不能删除任意元素，就用懒删除或过期检查。如果追问变成“若任务有不同释放时间或执行时长，就变成更一般的调度问题”，先判断原来的状态是否仍然覆盖新答案集合，再决定是微调边界、改 value 语义，还是换成另一类结构。

\subsubsection{LeetCode 632 最小区间覆盖 K 个列表}

先每轮扫描所有候选找极值，确认答案选择的对象。本题的模型是每个列表至少取一个元素时，当前区间由最小当前值和最大当前值决定；暴力版慢在动态候选集合不断变化，却每次重新找最值。本题的优化点是小顶堆保存各列表当前元素，同时维护当前最大值；弹出最小所在列表并推进，因此堆中只能放足以比较下一步选择的轻量状态。

堆只保证堆顶，不保证内部全局有序。本题的小顶堆保存各列表当前元素，同时维护当前最大值；弹出最小所在列表并推进要转成堆顶可回答的问题；如果元素会过期，还要在真正使用堆顶前清理或跳过旧状态。放到“堆与动态极值”这条主线里看，关键原则是：堆把插入和弹出极值控制在对数复杂度。Top K 用固定大小堆维护边界，合并多路序列用堆保存每一路当前头部，数据流中位数用双堆维护两半。

证明堆题时，要说明堆覆盖了所有可能的下一步候选，并且堆顶过期时不会被使用。本题的小顶堆保存各列表当前元素，同时维护当前最大值；弹出最小所在列表并推进是选择顺序，每个列表至少取一个元素时，当前区间由最小当前值和最大当前值决定是候选集合来源。这道题的边界不能留到最后补丁式处理，实验时覆盖某个列表为空、重复值、区间长度相同的 tie-breaker、列表耗尽，记录堆顶是否过期、比较器是否让正确候选在堆顶、K 或窗口大小为边界值时是否稳定。

C++ 实现上，\code{std::priority_queue} 默认大顶堆，小顶堆需要比较器。堆中保存大对象时尽量保存下标或轻量记录；若不能删除任意元素，就用懒删除或过期检查。如果追问变成“这是多路归并和覆盖窗口的结合”，先判断原来的状态是否仍然覆盖新答案集合，再决定是微调边界、改 value 语义，还是换成另一类结构。

\subsubsection{LeetCode 703 数据流中的第 K 大元素}

先每轮扫描所有候选找极值，确认答案选择的对象。本题的模型是数据不断到达，只关心当前第 K 大而非完整排序；暴力版慢在动态候选集合不断变化，却每次重新找最值。本题的优化点是维护大小为 K 的小顶堆，堆顶就是当前第 K 大元素，因此堆中只能放足以比较下一步选择的轻量状态。

堆只保证堆顶，不保证内部全局有序。本题的维护大小为 K 的小顶堆，堆顶就是当前第 K 大元素要转成堆顶可回答的问题；如果元素会过期，还要在真正使用堆顶前清理或跳过旧状态。放到“堆与动态极值”这条主线里看，关键原则是：堆把插入和弹出极值控制在对数复杂度。Top K 用固定大小堆维护边界，合并多路序列用堆保存每一路当前头部，数据流中位数用双堆维护两半。

证明堆题时，要说明堆覆盖了所有可能的下一步候选，并且堆顶过期时不会被使用。本题的维护大小为 K 的小顶堆，堆顶就是当前第 K 大元素是选择顺序，数据不断到达，只关心当前第 K 大而非完整排序是候选集合来源。这道题的边界不能留到最后补丁式处理，实验时覆盖初始元素少于 K、重复值、K 为一、连续 add，记录堆顶是否过期、比较器是否让正确候选在堆顶、K 或窗口大小为边界值时是否稳定。

C++ 实现上，\code{std::priority_queue} 默认大顶堆，小顶堆需要比较器。堆中保存大对象时尽量保存下标或轻量记录；若不能删除任意元素，就用懒删除或过期检查。如果追问变成“这题是 Top K 的在线版本，和一次性排序的适用场景不同”，先判断原来的状态是否仍然覆盖新答案集合，再决定是微调边界、改 value 语义，还是换成另一类结构。

\subsubsection{LeetCode 857 雇佣 K 名工人的最低成本}

先每轮扫描所有候选找极值，确认答案选择的对象。本题的模型是工资比例由当前队伍中最高工资质量比决定，总成本等于比例乘质量和；暴力版慢在动态候选集合不断变化，却每次重新找最值。本题的优化点是按比例升序扫描，用大顶堆维护当前 K 个工人的最小质量和，因此堆中只能放足以比较下一步选择的轻量状态。

堆只保证堆顶，不保证内部全局有序。本题的按比例升序扫描，用大顶堆维护当前 K 个工人的最小质量和要转成堆顶可回答的问题；如果元素会过期，还要在真正使用堆顶前清理或跳过旧状态。放到“堆与动态极值”这条主线里看，关键原则是：堆把插入和弹出极值控制在对数复杂度。Top K 用固定大小堆维护边界，合并多路序列用堆保存每一路当前头部，数据流中位数用双堆维护两半。

证明堆题时，要说明堆覆盖了所有可能的下一步候选，并且堆顶过期时不会被使用。本题的按比例升序扫描，用大顶堆维护当前 K 个工人的最小质量和是选择顺序，工资比例由当前队伍中最高工资质量比决定，总成本等于比例乘质量和是候选集合来源。这道题的边界不能留到最后补丁式处理，实验时覆盖K 为一、比例相同、浮点精度、质量和更新，记录堆顶是否过期、比较器是否让正确候选在堆顶、K 或窗口大小为边界值时是否稳定。

C++ 实现上，\code{std::priority_queue} 默认大顶堆，小顶堆需要比较器。堆中保存大对象时尽量保存下标或轻量记录；若不能删除任意元素，就用懒删除或过期检查。如果追问变成“这是排序确定价格比例，堆维护可替换候选的组合题”，先判断原来的状态是否仍然覆盖新答案集合，再决定是微调边界、改 value 语义，还是换成另一类结构。

\subsubsection{LeetCode 871 最低加油次数}

先每轮扫描所有候选找极值，确认答案选择的对象。本题的模型是行驶过程中，经过的加油站形成可反悔的燃料候选；暴力版慢在动态候选集合不断变化，却每次重新找最值。本题的优化点是按位置扫描，油量不足到达下一点时，从已路过站点中取最大燃料补充，因此堆中只能放足以比较下一步选择的轻量状态。

堆只保证堆顶，不保证内部全局有序。本题的按位置扫描，油量不足到达下一点时，从已路过站点中取最大燃料补充要转成堆顶可回答的问题；如果元素会过期，还要在真正使用堆顶前清理或跳过旧状态。放到“堆与动态极值”这条主线里看，关键原则是：堆把插入和弹出极值控制在对数复杂度。Top K 用固定大小堆维护边界，合并多路序列用堆保存每一路当前头部，数据流中位数用双堆维护两半。

证明堆题时，要说明堆覆盖了所有可能的下一步候选，并且堆顶过期时不会被使用。本题的按位置扫描，油量不足到达下一点时，从已路过站点中取最大燃料补充是选择顺序，行驶过程中，经过的加油站形成可反悔的燃料候选是候选集合来源。这道题的边界不能留到最后补丁式处理，实验时覆盖起点油足够、无法到达第一个站、多个站同位置、目标作为终点事件，记录堆顶是否过期、比较器是否让正确候选在堆顶、K 或窗口大小为边界值时是否稳定。

C++ 实现上，\code{std::priority_queue} 默认大顶堆，小顶堆需要比较器。堆中保存大对象时尽量保存下标或轻量记录；若不能删除任意元素，就用懒删除或过期检查。如果追问变成“这是反悔贪心，堆里保存的是已经经过但尚未使用的资源”，先判断原来的状态是否仍然覆盖新答案集合，再决定是微调边界、改 value 语义，还是换成另一类结构。

\subsubsection{LeetCode 973 最接近原点的 K 个点}

先每轮扫描所有候选找极值，确认答案选择的对象。本题的模型是距离只需比较平方，避免开方误差和成本；暴力版慢在动态候选集合不断变化，却每次重新找最值。本题的优化点是大小为 K 的大顶堆保存当前最近 K 个点中最远者，因此堆中只能放足以比较下一步选择的轻量状态。

堆只保证堆顶，不保证内部全局有序。本题的大小为 K 的大顶堆保存当前最近 K 个点中最远者要转成堆顶可回答的问题；如果元素会过期，还要在真正使用堆顶前清理或跳过旧状态。放到“堆与动态极值”这条主线里看，关键原则是：堆把插入和弹出极值控制在对数复杂度。Top K 用固定大小堆维护边界，合并多路序列用堆保存每一路当前头部，数据流中位数用双堆维护两半。

证明堆题时，要说明堆覆盖了所有可能的下一步候选，并且堆顶过期时不会被使用。本题的大小为 K 的大顶堆保存当前最近 K 个点中最远者是选择顺序，距离只需比较平方，避免开方误差和成本是候选集合来源。这道题的边界不能留到最后补丁式处理，实验时覆盖K 为一、K 为全部、距离相同、坐标平方溢出，记录堆顶是否过期、比较器是否让正确候选在堆顶、K 或窗口大小为边界值时是否稳定。

C++ 实现上，\code{std::priority_queue} 默认大顶堆，小顶堆需要比较器。堆中保存大对象时尽量保存下标或轻量记录；若不能删除任意元素，就用懒删除或过期检查。如果追问变成“快速选择可平均线性，但堆更适合数据流”，先判断原来的状态是否仍然覆盖新答案集合，再决定是微调边界、改 value 语义，还是换成另一类结构。

\subsubsection{LeetCode 1046 最后一块石头的重量}

先每轮扫描所有候选找极值，确认答案选择的对象。本题的模型是每次都要取当前最重的两块石头碰撞；暴力版慢在动态候选集合不断变化，却每次重新找最值。本题的优化点是大顶堆保存所有重量，弹出两个最大值，若不同则把差值放回，因此堆中只能放足以比较下一步选择的轻量状态。

堆只保证堆顶，不保证内部全局有序。本题的大顶堆保存所有重量，弹出两个最大值，若不同则把差值放回要转成堆顶可回答的问题；如果元素会过期，还要在真正使用堆顶前清理或跳过旧状态。放到“堆与动态极值”这条主线里看，关键原则是：堆把插入和弹出极值控制在对数复杂度。Top K 用固定大小堆维护边界，合并多路序列用堆保存每一路当前头部，数据流中位数用双堆维护两半。

证明堆题时，要说明堆覆盖了所有可能的下一步候选，并且堆顶过期时不会被使用。本题的大顶堆保存所有重量，弹出两个最大值，若不同则把差值放回是选择顺序，每次都要取当前最重的两块石头碰撞是候选集合来源。这道题的边界不能留到最后补丁式处理，实验时覆盖没有石头、一块石头、两块相等、多次产生相同重量，记录堆顶是否过期、比较器是否让正确候选在堆顶、K 或窗口大小为边界值时是否稳定。

C++ 实现上，\code{std::priority_queue} 默认大顶堆，小顶堆需要比较器。堆中保存大对象时尽量保存下标或轻量记录；若不能删除任意元素，就用懒删除或过期检查。如果追问变成“最后一块石头 II 则是分组差值最小化，模型转为 0-1 背包”，先判断原来的状态是否仍然覆盖新答案集合，再决定是微调边界、改 value 语义，还是换成另一类结构。

\subsubsection{LeetCode 1353 最多可以参加的会议数目}

先每轮扫描所有候选找极值，确认答案选择的对象。本题的模型是每天只能参加一个正在进行的会议，应优先选择最早结束的会议；暴力版慢在动态候选集合不断变化，却每次重新找最值。本题的优化点是按开始日排序扫描日期，把当天可参加会议的结束日入小顶堆，弹出最早结束者，因此堆中只能放足以比较下一步选择的轻量状态。

堆只保证堆顶，不保证内部全局有序。本题的按开始日排序扫描日期，把当天可参加会议的结束日入小顶堆，弹出最早结束者要转成堆顶可回答的问题；如果元素会过期，还要在真正使用堆顶前清理或跳过旧状态。放到“堆与动态极值”这条主线里看，关键原则是：堆把插入和弹出极值控制在对数复杂度。Top K 用固定大小堆维护边界，合并多路序列用堆保存每一路当前头部，数据流中位数用双堆维护两半。

证明堆题时，要说明堆覆盖了所有可能的下一步候选，并且堆顶过期时不会被使用。本题的按开始日排序扫描日期，把当天可参加会议的结束日入小顶堆，弹出最早结束者是选择顺序，每天只能参加一个正在进行的会议，应优先选择最早结束的会议是候选集合来源。这道题的边界不能留到最后补丁式处理，实验时覆盖同一天多个会议、过期会议清理、空闲日期跳跃、结束日相同，记录堆顶是否过期、比较器是否让正确候选在堆顶、K 或窗口大小为边界值时是否稳定。

C++ 实现上，\code{std::priority_queue} 默认大顶堆，小顶堆需要比较器。堆中保存大对象时尽量保存下标或轻量记录；若不能删除任意元素，就用懒删除或过期检查。如果追问变成“这是扫描线和堆结合的调度题”，先判断原来的状态是否仍然覆盖新答案集合，再决定是微调边界、改 value 语义，还是换成另一类结构。

\subsubsection{LeetCode 1675 数组的最小偏移量}

先每轮扫描所有候选找极值，确认答案选择的对象。本题的模型是奇数可以先乘二，偶数可以不断除二，目标是缩小当前最大最小差；暴力版慢在动态候选集合不断变化，却每次重新找最值。本题的优化点是把所有数规范到可降低的偶数形态，用大顶堆维护当前最大值，同时记录当前最小值，因此堆中只能放足以比较下一步选择的轻量状态。

堆只保证堆顶，不保证内部全局有序。本题的把所有数规范到可降低的偶数形态，用大顶堆维护当前最大值，同时记录当前最小值要转成堆顶可回答的问题；如果元素会过期，还要在真正使用堆顶前清理或跳过旧状态。放到“堆与动态极值”这条主线里看，关键原则是：堆把插入和弹出极值控制在对数复杂度。Top K 用固定大小堆维护边界，合并多路序列用堆保存每一路当前头部，数据流中位数用双堆维护两半。

证明堆题时，要说明堆覆盖了所有可能的下一步候选，并且堆顶过期时不会被使用。本题的把所有数规范到可降低的偶数形态，用大顶堆维护当前最大值，同时记录当前最小值是选择顺序，奇数可以先乘二，偶数可以不断除二，目标是缩小当前最大最小差是候选集合来源。这道题的边界不能留到最后补丁式处理，实验时覆盖全奇数、全偶数、最大值变奇数停止、数值范围，记录堆顶是否过期、比较器是否让正确候选在堆顶、K 或窗口大小为边界值时是否稳定。

C++ 实现上，\code{std::priority_queue} 默认大顶堆，小顶堆需要比较器。堆中保存大对象时尽量保存下标或轻量记录；若不能删除任意元素，就用懒删除或过期检查。如果追问变成“这是堆维护动态极值和状态空间规范化的结合”，先判断原来的状态是否仍然覆盖新答案集合，再决定是微调边界、改 value 语义，还是换成另一类结构。

\subsubsection{LeetCode 1851 包含每个查询的最小区间}

先每轮扫描所有候选找极值，确认答案选择的对象。本题的模型是每个查询点需要找覆盖它的最短区间，区间和查询都可离线排序；暴力版慢在动态候选集合不断变化，却每次重新找最值。本题的优化点是按查询值升序扫描，把左端不大于查询的区间入堆，堆顶过期区间弹出，因此堆中只能放足以比较下一步选择的轻量状态。

堆只保证堆顶，不保证内部全局有序。本题的按查询值升序扫描，把左端不大于查询的区间入堆，堆顶过期区间弹出要转成堆顶可回答的问题；如果元素会过期，还要在真正使用堆顶前清理或跳过旧状态。放到“堆与动态极值”这条主线里看，关键原则是：堆把插入和弹出极值控制在对数复杂度。Top K 用固定大小堆维护边界，合并多路序列用堆保存每一路当前头部，数据流中位数用双堆维护两半。

证明堆题时，要说明堆覆盖了所有可能的下一步候选，并且堆顶过期时不会被使用。本题的按查询值升序扫描，把左端不大于查询的区间入堆，堆顶过期区间弹出是选择顺序，每个查询点需要找覆盖它的最短区间，区间和查询都可离线排序是候选集合来源。这道题的边界不能留到最后补丁式处理，实验时覆盖没有覆盖区间、区间长度相同、查询重复、右端过期，记录堆顶是否过期、比较器是否让正确候选在堆顶、K 或窗口大小为边界值时是否稳定。

C++ 实现上，\code{std::priority_queue} 默认大顶堆，小顶堆需要比较器。堆中保存大对象时尽量保存下标或轻量记录；若不能删除任意元素，就用懒删除或过期检查。如果追问变成“这是离线扫描线、堆和区间覆盖的组合题”，先判断原来的状态是否仍然覆盖新答案集合，再决定是微调边界、改 value 语义，还是换成另一类结构。

\subsection{链表与指针重连工作坊}

先用数组或多次扫描得到直观答案，再把操作改成节点重连。链表题的关键是保存后继、明确前驱和当前段边界，不能靠值交换掩盖指针关系。

\subsubsection{LeetCode 2 两数相加}

先把链表节点拷到数组里得到一个容易验证的版本，再回到节点重连。本题的模型是链表按低位到高位保存数字，本质是竖式加法而不是整数转换；暴力版能帮助确认节点顺序或节点身份。优化时关注维护两个游标和进位，任一链未结束或进位非零都继续生成节点，不要用值交换掩盖链表结构要求。

链表优化要把指针角色固定下来。本题用维护两个游标和进位，任一链未结束或进位非零都继续生成节点推进，至少要画出前驱、当前节点、后继节点和已处理链尾。每次改 next 前先保存后继，防止链断掉。放到“链表与指针重连”这条主线里看，关键原则是：优化来自哨兵节点、快慢指针、前驱指针和局部反转。链表题的核心不是值交换，而是保持断开和接回的顺序正确。

链表题的正确性靠结构不变量：已经处理好的链仍然连通，未处理部分仍然可达，二者之间只有明确连接点。本题的维护两个游标和进位，任一链未结束或进位非零都继续生成节点每执行一次，都不能破坏这三个条件。这道题的边界不能留到最后补丁式处理，实验时覆盖两链长度不同、最后进位、输入为零、不能用内置整数承载超长数字，每步画出前驱、当前、后继和下一段头。链表题不要只看输出值，还要检查节点是否丢失或成环。

C++ 实现上，平台管理输入节点时不要随意 \code{delete}。使用虚拟头节点统一头部边界；每次重连前保存后继。比较交点时比较节点地址，不比较节点值。如果追问变成“若链表高位在前，可以用栈反向处理或先反转链表”，先判断原来的状态是否仍然覆盖新答案集合，再决定是微调边界、改 value 语义，还是换成另一类结构。

\subsubsection{LeetCode 19 删除链表的倒数第 N 个结点}

先把链表节点拷到数组里得到一个容易验证的版本，再回到节点重连。本题的模型是倒数位置可以转成两个指针之间固定 n 步的距离；暴力版能帮助确认节点顺序或节点身份。优化时关注快指针先走 n 步，再让快慢指针同步移动，慢指针停在待删节点前驱，不要用值交换掩盖链表结构要求。

链表优化要把指针角色固定下来。本题用快指针先走 n 步，再让快慢指针同步移动，慢指针停在待删节点前驱推进，至少要画出前驱、当前节点、后继节点和已处理链尾。每次改 next 前先保存后继，防止链断掉。放到“链表与指针重连”这条主线里看，关键原则是：优化来自哨兵节点、快慢指针、前驱指针和局部反转。链表题的核心不是值交换，而是保持断开和接回的顺序正确。

链表题的正确性靠结构不变量：已经处理好的链仍然连通，未处理部分仍然可达，二者之间只有明确连接点。本题的快指针先走 n 步，再让快慢指针同步移动，慢指针停在待删节点前驱每执行一次，都不能破坏这三个条件。这道题的边界不能留到最后补丁式处理，实验时覆盖删除头节点、n 等于链表长度、只有一个节点、题目保证 n 合法，每步画出前驱、当前、后继和下一段头。链表题不要只看输出值，还要检查节点是否丢失或成环。

C++ 实现上，平台管理输入节点时不要随意 \code{delete}。使用虚拟头节点统一头部边界；每次重连前保存后继。比较交点时比较节点地址，不比较节点值。如果追问变成“若 n 不保证合法，需要在快指针预走阶段检测长度不足”，先判断原来的状态是否仍然覆盖新答案集合，再决定是微调边界、改 value 语义，还是换成另一类结构。

\subsubsection{LeetCode 21 合并两个有序链表}

先把链表节点拷到数组里得到一个容易验证的版本，再回到节点重连。本题的模型是两个链表已经各自有序，下一节点只能来自两个当前头中较小者；暴力版能帮助确认节点顺序或节点身份。优化时关注虚拟头统一结果链表，尾指针不断接上较小节点并推进来源链，不要用值交换掩盖链表结构要求。

链表优化要把指针角色固定下来。本题用虚拟头统一结果链表，尾指针不断接上较小节点并推进来源链推进，至少要画出前驱、当前节点、后继节点和已处理链尾。每次改 next 前先保存后继，防止链断掉。放到“链表与指针重连”这条主线里看，关键原则是：优化来自哨兵节点、快慢指针、前驱指针和局部反转。链表题的核心不是值交换，而是保持断开和接回的顺序正确。

链表题的正确性靠结构不变量：已经处理好的链仍然连通，未处理部分仍然可达，二者之间只有明确连接点。本题的虚拟头统一结果链表，尾指针不断接上较小节点并推进来源链每执行一次，都不能破坏这三个条件。这道题的边界不能留到最后补丁式处理，实验时覆盖某一条链为空、重复值、剩余整段接回、是否复用原节点，每步画出前驱、当前、后继和下一段头。链表题不要只看输出值，还要检查节点是否丢失或成环。

C++ 实现上，平台管理输入节点时不要随意 \code{delete}。使用虚拟头节点统一头部边界；每次重连前保存后继。比较交点时比较节点地址，不比较节点值。如果追问变成“合并 K 条链表可用分治或堆，核心仍是维护每条链的当前头”，先判断原来的状态是否仍然覆盖新答案集合，再决定是微调边界、改 value 语义，还是换成另一类结构。

\subsubsection{LeetCode 24 两两交换链表中的节点}

先把链表节点拷到数组里得到一个容易验证的版本，再回到节点重连。本题的模型是链表按固定长度为二的小组做局部重连；暴力版能帮助确认节点顺序或节点身份。优化时关注使用组前驱、第一节点、第二节点和下一组头，按顺序改指针并推进前驱，不要用值交换掩盖链表结构要求。

链表优化要把指针角色固定下来。本题用使用组前驱、第一节点、第二节点和下一组头，按顺序改指针并推进前驱推进，至少要画出前驱、当前节点、后继节点和已处理链尾。每次改 next 前先保存后继，防止链断掉。放到“链表与指针重连”这条主线里看，关键原则是：优化来自哨兵节点、快慢指针、前驱指针和局部反转。链表题的核心不是值交换，而是保持断开和接回的顺序正确。

链表题的正确性靠结构不变量：已经处理好的链仍然连通，未处理部分仍然可达，二者之间只有明确连接点。本题的使用组前驱、第一节点、第二节点和下一组头，按顺序改指针并推进前驱每执行一次，都不能破坏这三个条件。这道题的边界不能留到最后补丁式处理，实验时覆盖空链、单节点、奇数长度、交换后前驱更新到组尾，每步画出前驱、当前、后继和下一段头。链表题不要只看输出值，还要检查节点是否丢失或成环。

C++ 实现上，平台管理输入节点时不要随意 \code{delete}。使用虚拟头节点统一头部边界；每次重连前保存后继。比较交点时比较节点地址，不比较节点值。如果追问变成“K 个一组翻转是同一思想的推广，只是组内反转更长”，先判断原来的状态是否仍然覆盖新答案集合，再决定是微调边界、改 value 语义，还是换成另一类结构。

\subsubsection{LeetCode 25 K 个一组翻转链表}

先把链表节点拷到数组里得到一个容易验证的版本，再回到节点重连。本题的模型是链表被切成长度为 K 的连续组，只有完整组才反转；暴力版能帮助确认节点顺序或节点身份。优化时关注每轮先探测 K 个节点确认完整，再保存下一组头并做局部反转接回，不要用值交换掩盖链表结构要求。

链表优化要把指针角色固定下来。本题用每轮先探测 K 个节点确认完整，再保存下一组头并做局部反转接回推进，至少要画出前驱、当前节点、后继节点和已处理链尾。每次改 next 前先保存后继，防止链断掉。放到“链表与指针重连”这条主线里看，关键原则是：优化来自哨兵节点、快慢指针、前驱指针和局部反转。链表题的核心不是值交换，而是保持断开和接回的顺序正确。

链表题的正确性靠结构不变量：已经处理好的链仍然连通，未处理部分仍然可达，二者之间只有明确连接点。本题的每轮先探测 K 个节点确认完整，再保存下一组头并做局部反转接回每执行一次，都不能破坏这三个条件。这道题的边界不能留到最后补丁式处理，实验时覆盖不足 K 个不反转、K 为一、刚好整除、反转后头尾接错，每步画出前驱、当前、后继和下一段头。链表题不要只看输出值，还要检查节点是否丢失或成环。

C++ 实现上，平台管理输入节点时不要随意 \code{delete}。使用虚拟头节点统一头部边界；每次重连前保存后继。比较交点时比较节点地址，不比较节点值。如果追问变成“递归写法短但可能有栈风险，迭代写法更接近工程实现”，先判断原来的状态是否仍然覆盖新答案集合，再决定是微调边界、改 value 语义，还是换成另一类结构。

\subsubsection{LeetCode 61 旋转链表}

先把链表节点拷到数组里得到一个容易验证的版本，再回到节点重连。本题的模型是右旋 k 位等价于把链表尾部一段切到头部；暴力版能帮助确认节点顺序或节点身份。优化时关注先求长度并让 k 对长度取模，再找到新尾节点并重连成新头，不要用值交换掩盖链表结构要求。

链表优化要把指针角色固定下来。本题用先求长度并让 k 对长度取模，再找到新尾节点并重连成新头推进，至少要画出前驱、当前节点、后继节点和已处理链尾。每次改 next 前先保存后继，防止链断掉。放到“链表与指针重连”这条主线里看，关键原则是：优化来自哨兵节点、快慢指针、前驱指针和局部反转。链表题的核心不是值交换，而是保持断开和接回的顺序正确。

链表题的正确性靠结构不变量：已经处理好的链仍然连通，未处理部分仍然可达，二者之间只有明确连接点。本题的先求长度并让 k 对长度取模，再找到新尾节点并重连成新头每执行一次，都不能破坏这三个条件。这道题的边界不能留到最后补丁式处理，实验时覆盖空链、单节点、k 为零、k 大于长度、刚好整圈，每步画出前驱、当前、后继和下一段头。链表题不要只看输出值，还要检查节点是否丢失或成环。

C++ 实现上，平台管理输入节点时不要随意 \code{delete}。使用虚拟头节点统一头部边界；每次重连前保存后继。比较交点时比较节点地址，不比较节点值。如果追问变成“若本来把链表接成环，最后必须断开新尾的 next”，先判断原来的状态是否仍然覆盖新答案集合，再决定是微调边界、改 value 语义，还是换成另一类结构。

\subsubsection{LeetCode 82 删除排序链表中的重复元素 II}

先把链表节点拷到数组里得到一个容易验证的版本，再回到节点重连。本题的模型是有序链表中重复值连续出现，题目要求把所有重复值整组删除；暴力版能帮助确认节点顺序或节点身份。优化时关注虚拟头保存结果前驱，遇到一段重复值时跳过整段，否则前驱前进，不要用值交换掩盖链表结构要求。

链表优化要把指针角色固定下来。本题用虚拟头保存结果前驱，遇到一段重复值时跳过整段，否则前驱前进推进，至少要画出前驱、当前节点、后继节点和已处理链尾。每次改 next 前先保存后继，防止链断掉。放到“链表与指针重连”这条主线里看，关键原则是：优化来自哨兵节点、快慢指针、前驱指针和局部反转。链表题的核心不是值交换，而是保持断开和接回的顺序正确。

链表题的正确性靠结构不变量：已经处理好的链仍然连通，未处理部分仍然可达，二者之间只有明确连接点。本题的虚拟头保存结果前驱，遇到一段重复值时跳过整段，否则前驱前进每执行一次，都不能破坏这三个条件。这道题的边界不能留到最后补丁式处理，实验时覆盖头部重复、尾部重复、全部重复、没有重复，每步画出前驱、当前、后继和下一段头。链表题不要只看输出值，还要检查节点是否丢失或成环。

C++ 实现上，平台管理输入节点时不要随意 \code{delete}。使用虚拟头节点统一头部边界；每次重连前保存后继。比较交点时比较节点地址，不比较节点值。如果追问变成“若只保留一个重复值，就是删除重复元素的较简单版本”，先判断原来的状态是否仍然覆盖新答案集合，再决定是微调边界、改 value 语义，还是换成另一类结构。

\subsubsection{LeetCode 83 删除排序链表中的重复元素}

先把链表节点拷到数组里得到一个容易验证的版本，再回到节点重连。本题的模型是有序链表让重复节点相邻，只需保留每段相同值的第一个节点；暴力版能帮助确认节点顺序或节点身份。优化时关注当前节点和后继值相同就跳过后继，否则当前节点前进，不要用值交换掩盖链表结构要求。

链表优化要把指针角色固定下来。本题用当前节点和后继值相同就跳过后继，否则当前节点前进推进，至少要画出前驱、当前节点、后继节点和已处理链尾。每次改 next 前先保存后继，防止链断掉。放到“链表与指针重连”这条主线里看，关键原则是：优化来自哨兵节点、快慢指针、前驱指针和局部反转。链表题的核心不是值交换，而是保持断开和接回的顺序正确。

链表题的正确性靠结构不变量：已经处理好的链仍然连通，未处理部分仍然可达，二者之间只有明确连接点。本题的当前节点和后继值相同就跳过后继，否则当前节点前进每执行一次，都不能破坏这三个条件。这道题的边界不能留到最后补丁式处理，实验时覆盖空链、单节点、全重复、重复段在尾部，每步画出前驱、当前、后继和下一段头。链表题不要只看输出值，还要检查节点是否丢失或成环。

C++ 实现上，平台管理输入节点时不要随意 \code{delete}。使用虚拟头节点统一头部边界；每次重连前保存后继。比较交点时比较节点地址，不比较节点值。如果追问变成“和删除重复元素 II 的区别是是否整段删除”，先判断原来的状态是否仍然覆盖新答案集合，再决定是微调边界、改 value 语义，还是换成另一类结构。

\subsubsection{LeetCode 92 反转链表 II}

先把链表节点拷到数组里得到一个容易验证的版本，再回到节点重连。本题的模型是只反转链表中一段连续区间，区间外节点必须保持原连接；暴力版能帮助确认节点顺序或节点身份。优化时关注用虚拟头找到区间前驱，再用头插法把区间内部节点逐个移到前面，不要用值交换掩盖链表结构要求。

链表优化要把指针角色固定下来。本题用用虚拟头找到区间前驱，再用头插法把区间内部节点逐个移到前面推进，至少要画出前驱、当前节点、后继节点和已处理链尾。每次改 next 前先保存后继，防止链断掉。放到“链表与指针重连”这条主线里看，关键原则是：优化来自哨兵节点、快慢指针、前驱指针和局部反转。链表题的核心不是值交换，而是保持断开和接回的顺序正确。

链表题的正确性靠结构不变量：已经处理好的链仍然连通，未处理部分仍然可达，二者之间只有明确连接点。本题的用虚拟头找到区间前驱，再用头插法把区间内部节点逐个移到前面每执行一次，都不能破坏这三个条件。这道题的边界不能留到最后补丁式处理，实验时覆盖从头开始反转、只反转一个节点、反转到尾部、区间边界，每步画出前驱、当前、后继和下一段头。链表题不要只看输出值，还要检查节点是否丢失或成环。

C++ 实现上，平台管理输入节点时不要随意 \code{delete}。使用虚拟头节点统一头部边界；每次重连前保存后继。比较交点时比较节点地址，不比较节点值。如果追问变成“K 个一组反转是在多个固定长度区间上重复这一类局部重连”，先判断原来的状态是否仍然覆盖新答案集合，再决定是微调边界、改 value 语义，还是换成另一类结构。

\subsubsection{LeetCode 141 环形链表}

先把链表节点拷到数组里得到一个容易验证的版本，再回到节点重连。本题的模型是快慢指针在有环时必然相遇，无环时快指针先到空；暴力版能帮助确认节点顺序或节点身份。优化时关注不用哈希表也能常数空间检测环，不要用值交换掩盖链表结构要求。

链表优化要把指针角色固定下来。本题用不用哈希表也能常数空间检测环推进，至少要画出前驱、当前节点、后继节点和已处理链尾。每次改 next 前先保存后继，防止链断掉。放到“链表与指针重连”这条主线里看，关键原则是：优化来自哨兵节点、快慢指针、前驱指针和局部反转。链表题的核心不是值交换，而是保持断开和接回的顺序正确。

链表题的正确性靠结构不变量：已经处理好的链仍然连通，未处理部分仍然可达，二者之间只有明确连接点。本题的不用哈希表也能常数空间检测环每执行一次，都不能破坏这三个条件。这道题的边界不能留到最后补丁式处理，实验时覆盖空链、单节点自环、两节点环、尾部无环，每步画出前驱、当前、后继和下一段头。链表题不要只看输出值，还要检查节点是否丢失或成环。

C++ 实现上，平台管理输入节点时不要随意 \code{delete}。使用虚拟头节点统一头部边界；每次重连前保存后继。比较交点时比较节点地址，不比较节点值。如果追问变成“若要求环入口，用相遇后双指针同步走”，先判断原来的状态是否仍然覆盖新答案集合，再决定是微调边界、改 value 语义，还是换成另一类结构。

\subsubsection{LeetCode 142 环形链表 II}

先把链表节点拷到数组里得到一个容易验证的版本，再回到节点重连。本题的模型是相遇点到入口的距离关系来自快慢指针速度差；暴力版能帮助确认节点顺序或节点身份。优化时关注相遇后一个指针回头，两个指针每次走一步，再次相遇就是入口，不要用值交换掩盖链表结构要求。

链表优化要把指针角色固定下来。本题用相遇后一个指针回头，两个指针每次走一步，再次相遇就是入口推进，至少要画出前驱、当前节点、后继节点和已处理链尾。每次改 next 前先保存后继，防止链断掉。放到“链表与指针重连”这条主线里看，关键原则是：优化来自哨兵节点、快慢指针、前驱指针和局部反转。链表题的核心不是值交换，而是保持断开和接回的顺序正确。

链表题的正确性靠结构不变量：已经处理好的链仍然连通，未处理部分仍然可达，二者之间只有明确连接点。本题的相遇后一个指针回头，两个指针每次走一步，再次相遇就是入口每执行一次，都不能破坏这三个条件。这道题的边界不能留到最后补丁式处理，实验时覆盖入口在头节点、长尾短环、无环、单节点环，每步画出前驱、当前、后继和下一段头。链表题不要只看输出值，还要检查节点是否丢失或成环。

C++ 实现上，平台管理输入节点时不要随意 \code{delete}。使用虚拟头节点统一头部边界；每次重连前保存后继。比较交点时比较节点地址，不比较节点值。如果追问变成“也可用哈希集合，但空间更高”，先判断原来的状态是否仍然覆盖新答案集合，再决定是微调边界、改 value 语义，还是换成另一类结构。

\subsubsection{LeetCode 146 LRU 缓存}

先把链表节点拷到数组里得到一个容易验证的版本，再回到节点重连。本题的模型是哈希表负责按 key 定位，双向链表负责维护新旧顺序；暴力版能帮助确认节点顺序或节点身份。优化时关注访问和更新都把节点移动到头部，容量满时删除尾部最旧节点，不要用值交换掩盖链表结构要求。

链表优化要把指针角色固定下来。本题用访问和更新都把节点移动到头部，容量满时删除尾部最旧节点推进，至少要画出前驱、当前节点、后继节点和已处理链尾。每次改 next 前先保存后继，防止链断掉。放到“链表与指针重连”这条主线里看，关键原则是：优化来自哨兵节点、快慢指针、前驱指针和局部反转。链表题的核心不是值交换，而是保持断开和接回的顺序正确。

链表题的正确性靠结构不变量：已经处理好的链仍然连通，未处理部分仍然可达，二者之间只有明确连接点。本题的访问和更新都把节点移动到头部，容量满时删除尾部最旧节点每执行一次，都不能破坏这三个条件。这道题的边界不能留到最后补丁式处理，实验时覆盖容量为一、重复 put、get 不存在、更新已有值，每步画出前驱、当前、后继和下一段头。链表题不要只看输出值，还要检查节点是否丢失或成环。

C++ 实现上，平台管理输入节点时不要随意 \code{delete}。使用虚拟头节点统一头部边界；每次重连前保存后继。比较交点时比较节点地址，不比较节点值。如果追问变成“C++ 可用 \code{std::list} 加迭代器，也可手写双向链表理解所有权”，先判断原来的状态是否仍然覆盖新答案集合，再决定是微调边界、改 value 语义，还是换成另一类结构。

\subsubsection{LeetCode 148 排序链表}

先把链表节点拷到数组里得到一个容易验证的版本，再回到节点重连。本题的模型是链表不能随机访问，适合归并排序而不是快速排序数组 partition；暴力版能帮助确认节点顺序或节点身份。优化时关注自底向上迭代归并能避免递归栈，同时保持 O(n log n)，不要用值交换掩盖链表结构要求。

链表优化要把指针角色固定下来。本题用自底向上迭代归并能避免递归栈，同时保持 O(n log n)推进，至少要画出前驱、当前节点、后继节点和已处理链尾。每次改 next 前先保存后继，防止链断掉。放到“链表与指针重连”这条主线里看，关键原则是：优化来自哨兵节点、快慢指针、前驱指针和局部反转。链表题的核心不是值交换，而是保持断开和接回的顺序正确。

链表题的正确性靠结构不变量：已经处理好的链仍然连通，未处理部分仍然可达，二者之间只有明确连接点。本题的自底向上迭代归并能避免递归栈，同时保持 O(n log n)每执行一次，都不能破坏这三个条件。这道题的边界不能留到最后补丁式处理，实验时覆盖空链、单节点、重复值、奇数长度、切断子链，每步画出前驱、当前、后继和下一段头。链表题不要只看输出值，还要检查节点是否丢失或成环。

C++ 实现上，平台管理输入节点时不要随意 \code{delete}。使用虚拟头节点统一头部边界；每次重连前保存后继。比较交点时比较节点地址，不比较节点值。如果追问变成“若把值复制到数组排序，会改变节点身份语义”，先判断原来的状态是否仍然覆盖新答案集合，再决定是微调边界、改 value 语义，还是换成另一类结构。

\subsubsection{LeetCode 160 相交链表}

先把链表节点拷到数组里得到一个容易验证的版本，再回到节点重连。本题的模型是相交是节点地址相同，不是节点值相同；暴力版能帮助确认节点顺序或节点身份。优化时关注两个指针走完各自链表后切换到另一条链，走过总长度相同后相遇，不要用值交换掩盖链表结构要求。

链表优化要把指针角色固定下来。本题用两个指针走完各自链表后切换到另一条链，走过总长度相同后相遇推进，至少要画出前驱、当前节点、后继节点和已处理链尾。每次改 next 前先保存后继，防止链断掉。放到“链表与指针重连”这条主线里看，关键原则是：优化来自哨兵节点、快慢指针、前驱指针和局部反转。链表题的核心不是值交换，而是保持断开和接回的顺序正确。

链表题的正确性靠结构不变量：已经处理好的链仍然连通，未处理部分仍然可达，二者之间只有明确连接点。本题的两个指针走完各自链表后切换到另一条链，走过总长度相同后相遇每执行一次，都不能破坏这三个条件。这道题的边界不能留到最后补丁式处理，实验时覆盖不相交、头节点相交、尾节点相交、长度差很大，每步画出前驱、当前、后继和下一段头。链表题不要只看输出值，还要检查节点是否丢失或成环。

C++ 实现上，平台管理输入节点时不要随意 \code{delete}。使用虚拟头节点统一头部边界；每次重连前保存后继。比较交点时比较节点地址，不比较节点值。如果追问变成“哈希集合更直观但空间更高”，先判断原来的状态是否仍然覆盖新答案集合，再决定是微调边界、改 value 语义，还是换成另一类结构。

\subsubsection{LeetCode 206 反转链表}

先把链表节点拷到数组里得到一个容易验证的版本，再回到节点重连。本题的模型是每次把当前节点的 next 指向已经反转好的前缀头；暴力版能帮助确认节点顺序或节点身份。优化时关注先保存后继，再改指针，再推进 previous 和 current，不要用值交换掩盖链表结构要求。

链表优化要把指针角色固定下来。本题用先保存后继，再改指针，再推进 previous 和 current推进，至少要画出前驱、当前节点、后继节点和已处理链尾。每次改 next 前先保存后继，防止链断掉。放到“链表与指针重连”这条主线里看，关键原则是：优化来自哨兵节点、快慢指针、前驱指针和局部反转。链表题的核心不是值交换，而是保持断开和接回的顺序正确。

链表题的正确性靠结构不变量：已经处理好的链仍然连通，未处理部分仍然可达，二者之间只有明确连接点。本题的先保存后继，再改指针，再推进 previous 和 current每执行一次，都不能破坏这三个条件。这道题的边界不能留到最后补丁式处理，实验时覆盖空链、单节点、两个节点、不要丢失剩余链，每步画出前驱、当前、后继和下一段头。链表题不要只看输出值，还要检查节点是否丢失或成环。

C++ 实现上，平台管理输入节点时不要随意 \code{delete}。使用虚拟头节点统一头部边界；每次重连前保存后继。比较交点时比较节点地址，不比较节点值。如果追问变成“K 组反转是在这段局部反转外增加分组边界”，先判断原来的状态是否仍然覆盖新答案集合，再决定是微调边界、改 value 语义，还是换成另一类结构。

\subsubsection{LeetCode 234 回文链表}

先把链表节点拷到数组里得到一个容易验证的版本，再回到节点重连。本题的模型是链表回文要求前半段和后半段逆序后逐一相等；暴力版能帮助确认节点顺序或节点身份。优化时关注快慢指针找中点，反转后半段，再和前半段同步比较，不要用值交换掩盖链表结构要求。

链表优化要把指针角色固定下来。本题用快慢指针找中点，反转后半段，再和前半段同步比较推进，至少要画出前驱、当前节点、后继节点和已处理链尾。每次改 next 前先保存后继，防止链断掉。放到“链表与指针重连”这条主线里看，关键原则是：优化来自哨兵节点、快慢指针、前驱指针和局部反转。链表题的核心不是值交换，而是保持断开和接回的顺序正确。

链表题的正确性靠结构不变量：已经处理好的链仍然连通，未处理部分仍然可达，二者之间只有明确连接点。本题的快慢指针找中点，反转后半段，再和前半段同步比较每执行一次，都不能破坏这三个条件。这道题的边界不能留到最后补丁式处理，实验时覆盖空链、单节点、奇数长度、偶数长度、是否恢复链表，每步画出前驱、当前、后继和下一段头。链表题不要只看输出值，还要检查节点是否丢失或成环。

C++ 实现上，平台管理输入节点时不要随意 \code{delete}。使用虚拟头节点统一头部边界；每次重连前保存后继。比较交点时比较节点地址，不比较节点值。如果追问变成“如果不允许修改链表，可用栈保存前半段但空间变为线性”，先判断原来的状态是否仍然覆盖新答案集合，再决定是微调边界、改 value 语义，还是换成另一类结构。

\subsubsection{LeetCode 430 扁平化多级双向链表}

先把链表节点拷到数组里得到一个容易验证的版本，再回到节点重连。本题的模型是多级链表按深度优先顺序展开成一条双向链表；暴力版能帮助确认节点顺序或节点身份。优化时关注显式栈保存后续 next，遇到 child 先接 child，再在子链结束后恢复原 next，不要用值交换掩盖链表结构要求。

链表优化要把指针角色固定下来。本题用显式栈保存后续 next，遇到 child 先接 child，再在子链结束后恢复原 next推进，至少要画出前驱、当前节点、后继节点和已处理链尾。每次改 next 前先保存后继，防止链断掉。放到“链表与指针重连”这条主线里看，关键原则是：优化来自哨兵节点、快慢指针、前驱指针和局部反转。链表题的核心不是值交换，而是保持断开和接回的顺序正确。

链表题的正确性靠结构不变量：已经处理好的链仍然连通，未处理部分仍然可达，二者之间只有明确连接点。本题的显式栈保存后续 next，遇到 child 先接 child，再在子链结束后恢复原 next每执行一次，都不能破坏这三个条件。这道题的边界不能留到最后补丁式处理，实验时覆盖child 为空、next 和 child 同时存在、多层嵌套、prev 指针恢复，每步画出前驱、当前、后继和下一段头。链表题不要只看输出值，还要检查节点是否丢失或成环。

C++ 实现上，平台管理输入节点时不要随意 \code{delete}。使用虚拟头节点统一头部边界；每次重连前保存后继。比较交点时比较节点地址，不比较节点值。如果追问变成“它和二叉树前序展开类似，核心是保存未处理分支”，先判断原来的状态是否仍然覆盖新答案集合，再决定是微调边界、改 value 语义，还是换成另一类结构。

\subsection{二叉树与树形状态工作坊}

先想每个节点重复扫描子树会得到什么信息，再把这些信息压成节点向父节点返回的状态。树题要区分自顶向下约束、后序汇总和按层传播。

\subsubsection{LeetCode 94 二叉树中序遍历}

先想最直接的树遍历会在每个节点重复求什么。本题的模型是中序访问顺序是左子树、根、右子树，BST 中会得到有序序列；如果每个节点都重新扫描子树，慢点就会非常明显。本题要压缩的状态是用显式栈模拟一路向左，再回退访问根并转向右子树，也就是节点和父子方向之间真正传递的信息。

树题要先判定状态方向。本题的用显式栈模拟一路向左，再回退访问根并转向右子树可能是自顶向下的约束、后序向上的汇总，或队列中的层号。返回值或栈元素不能只有节点指针，还要带上题目需要的信息。放到“二叉树与树形状态”这条主线里看，关键原则是：树形 DP 把每个节点看成子问题根；BFS 把层次边界显式放进队列；BST 题利用中序有序或上下界约束。工程实现优先考虑显式栈避免深树调用栈风险。

树题常用对子树归纳。假设子节点状态已经正确，父节点只按用显式栈模拟一路向左，再回退访问根并转向右子树合并或传递信息。本题的中序访问顺序是左子树、根、右子树，BST 中会得到有序序列决定空节点、叶子和极端链式树如何处理。这道题的边界不能留到最后补丁式处理，实验时覆盖空树、单节点、链式左树、链式右树，尤其关注空节点、单节点、链式树和全负值。递归版本和显式栈版本可以互相对拍。

C++ 实现上，极端深树可能造成递归栈风险，工程写法可以用 \code{std::vector<TreeNode*>} 或队列显式保存状态。BST 边界最好用更宽整数或可选前驱值。如果追问变成“Morris 遍历可做到常数额外空间，但会临时改树结构”，先判断原来的状态是否仍然覆盖新答案集合，再决定是微调边界、改 value 语义，还是换成另一类结构。

\subsubsection{LeetCode 98 验证二叉搜索树}

先想最直接的树遍历会在每个节点重复求什么。本题的模型是BST 约束是全局上下界，不是只比较父子节点；如果每个节点都重新扫描子树，慢点就会非常明显。本题要压缩的状态是中序严格递增或显式传递上下界都可以；中序版本要保存前一个访问值，也就是节点和父子方向之间真正传递的信息。

树题要先判定状态方向。本题的中序严格递增或显式传递上下界都可以；中序版本要保存前一个访问值可能是自顶向下的约束、后序向上的汇总，或队列中的层号。返回值或栈元素不能只有节点指针，还要带上题目需要的信息。放到“二叉树与树形状态”这条主线里看，关键原则是：树形 DP 把每个节点看成子问题根；BFS 把层次边界显式放进队列；BST 题利用中序有序或上下界约束。工程实现优先考虑显式栈避免深树调用栈风险。

树题常用对子树归纳。假设子节点状态已经正确，父节点只按中序严格递增或显式传递上下界都可以；中序版本要保存前一个访问值合并或传递信息。本题的BST 约束是全局上下界，不是只比较父子节点决定空节点、叶子和极端链式树如何处理。这道题的边界不能留到最后补丁式处理，实验时覆盖重复值、int 最小最大值、局部合法但全局非法，尤其关注空节点、单节点、链式树和全负值。递归版本和显式栈版本可以互相对拍。

C++ 实现上，极端深树可能造成递归栈风险，工程写法可以用 \code{std::vector<TreeNode*>} 或队列显式保存状态。BST 边界最好用更宽整数或可选前驱值。如果追问变成“若允许重复值，需要明确重复放左还是放右”，先判断原来的状态是否仍然覆盖新答案集合，再决定是微调边界、改 value 语义，还是换成另一类结构。

\subsubsection{LeetCode 102 二叉树层序遍历}

先想最直接的树遍历会在每个节点重复求什么。本题的模型是队列在每轮开始时保存当前层所有节点；如果每个节点都重新扫描子树，慢点就会非常明显。本题要压缩的状态是记录 level size 固定层边界，避免下一层节点混入当前层，也就是节点和父子方向之间真正传递的信息。

树题要先判定状态方向。本题的记录 level size 固定层边界，避免下一层节点混入当前层可能是自顶向下的约束、后序向上的汇总，或队列中的层号。返回值或栈元素不能只有节点指针，还要带上题目需要的信息。放到“二叉树与树形状态”这条主线里看，关键原则是：树形 DP 把每个节点看成子问题根；BFS 把层次边界显式放进队列；BST 题利用中序有序或上下界约束。工程实现优先考虑显式栈避免深树调用栈风险。

树题常用对子树归纳。假设子节点状态已经正确，父节点只按记录 level size 固定层边界，避免下一层节点混入当前层合并或传递信息。本题的队列在每轮开始时保存当前层所有节点决定空节点、叶子和极端链式树如何处理。这道题的边界不能留到最后补丁式处理，实验时覆盖空树、只有根、最后一层不满、左右顺序，尤其关注空节点、单节点、链式树和全负值。递归版本和显式栈版本可以互相对拍。

C++ 实现上，极端深树可能造成递归栈风险，工程写法可以用 \code{std::vector<TreeNode*>} 或队列显式保存状态。BST 边界最好用更宽整数或可选前驱值。如果追问变成“锯齿形层序只是在每层输出顺序上增加方向控制”，先判断原来的状态是否仍然覆盖新答案集合，再决定是微调边界、改 value 语义，还是换成另一类结构。

\subsubsection{LeetCode 103 二叉树锯齿形层序遍历}

先想最直接的树遍历会在每个节点重复求什么。本题的模型是BFS 层边界不变，变化的是当前层写入结果的方向；如果每个节点都重新扫描子树，慢点就会非常明显。本题要压缩的状态是可以层内反转，也可以预分配数组按方向写入目标下标，也就是节点和父子方向之间真正传递的信息。

树题要先判定状态方向。本题的可以层内反转，也可以预分配数组按方向写入目标下标可能是自顶向下的约束、后序向上的汇总，或队列中的层号。返回值或栈元素不能只有节点指针，还要带上题目需要的信息。放到“二叉树与树形状态”这条主线里看，关键原则是：树形 DP 把每个节点看成子问题根；BFS 把层次边界显式放进队列；BST 题利用中序有序或上下界约束。工程实现优先考虑显式栈避免深树调用栈风险。

树题常用对子树归纳。假设子节点状态已经正确，父节点只按可以层内反转，也可以预分配数组按方向写入目标下标合并或传递信息。本题的BFS 层边界不变，变化的是当前层写入结果的方向决定空节点、叶子和极端链式树如何处理。这道题的边界不能留到最后补丁式处理，实验时覆盖单层、偶数层、奇数层、空树，尤其关注空节点、单节点、链式树和全负值。递归版本和显式栈版本可以互相对拍。

C++ 实现上，极端深树可能造成递归栈风险，工程写法可以用 \code{std::vector<TreeNode*>} 或队列显式保存状态。BST 边界最好用更宽整数或可选前驱值。如果追问变成“若要求从底向上输出，收集后反转层列表即可”，先判断原来的状态是否仍然覆盖新答案集合，再决定是微调边界、改 value 语义，还是换成另一类结构。

\subsubsection{LeetCode 104 二叉树最大深度}

先想最直接的树遍历会在每个节点重复求什么。本题的模型是深度可以看成层数，也可以作为栈中携带的状态；如果每个节点都重新扫描子树，慢点就会非常明显。本题要压缩的状态是显式栈保存节点和深度，避免极端链式树递归栈过深，也就是节点和父子方向之间真正传递的信息。

树题要先判定状态方向。本题的显式栈保存节点和深度，避免极端链式树递归栈过深可能是自顶向下的约束、后序向上的汇总，或队列中的层号。返回值或栈元素不能只有节点指针，还要带上题目需要的信息。放到“二叉树与树形状态”这条主线里看，关键原则是：树形 DP 把每个节点看成子问题根；BFS 把层次边界显式放进队列；BST 题利用中序有序或上下界约束。工程实现优先考虑显式栈避免深树调用栈风险。

树题常用对子树归纳。假设子节点状态已经正确，父节点只按显式栈保存节点和深度，避免极端链式树递归栈过深合并或传递信息。本题的深度可以看成层数，也可以作为栈中携带的状态决定空节点、叶子和极端链式树如何处理。这道题的边界不能留到最后补丁式处理，实验时覆盖空树返回零、根深度是一、链式树，尤其关注空节点、单节点、链式树和全负值。递归版本和显式栈版本可以互相对拍。

C++ 实现上，极端深树可能造成递归栈风险，工程写法可以用 \code{std::vector<TreeNode*>} 或队列显式保存状态。BST 边界最好用更宽整数或可选前驱值。如果追问变成“最小深度不能简单取左右最小，因为空子树不构成叶子路径”，先判断原来的状态是否仍然覆盖新答案集合，再决定是微调边界、改 value 语义，还是换成另一类结构。

\subsubsection{LeetCode 105 前序中序构造二叉树}

先想最直接的树遍历会在每个节点重复求什么。本题的模型是前序给根，中序把左右子树切开；如果每个节点都重新扫描子树，慢点就会非常明显。本题要压缩的状态是用哈希表记录中序下标，避免每次线性寻找根位置，也就是节点和父子方向之间真正传递的信息。

树题要先判定状态方向。本题的用哈希表记录中序下标，避免每次线性寻找根位置可能是自顶向下的约束、后序向上的汇总，或队列中的层号。返回值或栈元素不能只有节点指针，还要带上题目需要的信息。放到“二叉树与树形状态”这条主线里看，关键原则是：树形 DP 把每个节点看成子问题根；BFS 把层次边界显式放进队列；BST 题利用中序有序或上下界约束。工程实现优先考虑显式栈避免深树调用栈风险。

树题常用对子树归纳。假设子节点状态已经正确，父节点只按用哈希表记录中序下标，避免每次线性寻找根位置合并或传递信息。本题的前序给根，中序把左右子树切开决定空节点、叶子和极端链式树如何处理。这道题的边界不能留到最后补丁式处理，实验时覆盖节点值唯一、空子树、左右子树长度计算、输入不合法，尤其关注空节点、单节点、链式树和全负值。递归版本和显式栈版本可以互相对拍。

C++ 实现上，极端深树可能造成递归栈风险，工程写法可以用 \code{std::vector<TreeNode*>} 或队列显式保存状态。BST 边界最好用更宽整数或可选前驱值。如果追问变成“后序中序构造类似，只是根来自后序末尾”，先判断原来的状态是否仍然覆盖新答案集合，再决定是微调边界、改 value 语义，还是换成另一类结构。

\subsubsection{LeetCode 106 中序后序构造二叉树}

先想最直接的树遍历会在每个节点重复求什么。本题的模型是后序末尾是根，中序位置确定左右子树规模；如果每个节点都重新扫描子树，慢点就会非常明显。本题要压缩的状态是构造时要小心后序左右区间边界，避免切分错位，也就是节点和父子方向之间真正传递的信息。

树题要先判定状态方向。本题的构造时要小心后序左右区间边界，避免切分错位可能是自顶向下的约束、后序向上的汇总，或队列中的层号。返回值或栈元素不能只有节点指针，还要带上题目需要的信息。放到“二叉树与树形状态”这条主线里看，关键原则是：树形 DP 把每个节点看成子问题根；BFS 把层次边界显式放进队列；BST 题利用中序有序或上下界约束。工程实现优先考虑显式栈避免深树调用栈风险。

树题常用对子树归纳。假设子节点状态已经正确，父节点只按构造时要小心后序左右区间边界，避免切分错位合并或传递信息。本题的后序末尾是根，中序位置确定左右子树规模决定空节点、叶子和极端链式树如何处理。这道题的边界不能留到最后补丁式处理，实验时覆盖空树、单节点、全左链、全右链，尤其关注空节点、单节点、链式树和全负值。递归版本和显式栈版本可以互相对拍。

C++ 实现上，极端深树可能造成递归栈风险，工程写法可以用 \code{std::vector<TreeNode*>} 或队列显式保存状态。BST 边界最好用更宽整数或可选前驱值。如果追问变成“若给前序和后序，二叉树通常不能唯一确定，除非额外限制”，先判断原来的状态是否仍然覆盖新答案集合，再决定是微调边界、改 value 语义，还是换成另一类结构。

\subsubsection{LeetCode 108 有序数组转二叉搜索树}

先想最直接的树遍历会在每个节点重复求什么。本题的模型是有序数组中点作为根能让左右规模接近；如果每个节点都重新扫描子树，慢点就会非常明显。本题要压缩的状态是每次选择区间中点，左右子区间构造左右子树，也就是节点和父子方向之间真正传递的信息。

树题要先判定状态方向。本题的每次选择区间中点，左右子区间构造左右子树可能是自顶向下的约束、后序向上的汇总，或队列中的层号。返回值或栈元素不能只有节点指针，还要带上题目需要的信息。放到“二叉树与树形状态”这条主线里看，关键原则是：树形 DP 把每个节点看成子问题根；BFS 把层次边界显式放进队列；BST 题利用中序有序或上下界约束。工程实现优先考虑显式栈避免深树调用栈风险。

树题常用对子树归纳。假设子节点状态已经正确，父节点只按每次选择区间中点，左右子区间构造左右子树合并或传递信息。本题的有序数组中点作为根能让左右规模接近决定空节点、叶子和极端链式树如何处理。这道题的边界不能留到最后补丁式处理，实验时覆盖空区间、偶数长度选左中还是右中、高度平衡定义，尤其关注空节点、单节点、链式树和全负值。递归版本和显式栈版本可以互相对拍。

C++ 实现上，极端深树可能造成递归栈风险，工程写法可以用 \code{std::vector<TreeNode*>} 或队列显式保存状态。BST 边界最好用更宽整数或可选前驱值。如果追问变成“工程上递归可读，若严格避免递归可用队列保存区间和父指针”，先判断原来的状态是否仍然覆盖新答案集合，再决定是微调边界、改 value 语义，还是换成另一类结构。

\subsubsection{LeetCode 110 平衡二叉树}

先想最直接的树遍历会在每个节点重复求什么。本题的模型是每个节点需要知道子树高度以及是否已经失衡；如果每个节点都重新扫描子树，慢点就会非常明显。本题要压缩的状态是后序汇总时一旦子树失衡就向上返回失败标记，避免重复算高度，也就是节点和父子方向之间真正传递的信息。

树题要先判定状态方向。本题的后序汇总时一旦子树失衡就向上返回失败标记，避免重复算高度可能是自顶向下的约束、后序向上的汇总，或队列中的层号。返回值或栈元素不能只有节点指针，还要带上题目需要的信息。放到“二叉树与树形状态”这条主线里看，关键原则是：树形 DP 把每个节点看成子问题根；BFS 把层次边界显式放进队列；BST 题利用中序有序或上下界约束。工程实现优先考虑显式栈避免深树调用栈风险。

树题常用对子树归纳。假设子节点状态已经正确，父节点只按后序汇总时一旦子树失衡就向上返回失败标记，避免重复算高度合并或传递信息。本题的每个节点需要知道子树高度以及是否已经失衡决定空节点、叶子和极端链式树如何处理。这道题的边界不能留到最后补丁式处理，实验时覆盖空树、单节点、左右高度差正好一、链式树，尤其关注空节点、单节点、链式树和全负值。递归版本和显式栈版本可以互相对拍。

C++ 实现上，极端深树可能造成递归栈风险，工程写法可以用 \code{std::vector<TreeNode*>} 或队列显式保存状态。BST 边界最好用更宽整数或可选前驱值。如果追问变成“可以把返回值设计成 optional 高度，空表示不平衡”，先判断原来的状态是否仍然覆盖新答案集合，再决定是微调边界、改 value 语义，还是换成另一类结构。

\subsubsection{LeetCode 111 二叉树最小深度}

先想最直接的树遍历会在每个节点重复求什么。本题的模型是最小深度是到最近叶子的节点数，不是左右深度简单取最小；如果每个节点都重新扫描子树，慢点就会非常明显。本题要压缩的状态是BFS 第一次遇到叶子就是最小深度，因为按层扩展，也就是节点和父子方向之间真正传递的信息。

树题要先判定状态方向。本题的BFS 第一次遇到叶子就是最小深度，因为按层扩展可能是自顶向下的约束、后序向上的汇总，或队列中的层号。返回值或栈元素不能只有节点指针，还要带上题目需要的信息。放到“二叉树与树形状态”这条主线里看，关键原则是：树形 DP 把每个节点看成子问题根；BFS 把层次边界显式放进队列；BST 题利用中序有序或上下界约束。工程实现优先考虑显式栈避免深树调用栈风险。

树题常用对子树归纳。假设子节点状态已经正确，父节点只按BFS 第一次遇到叶子就是最小深度，因为按层扩展合并或传递信息。本题的最小深度是到最近叶子的节点数，不是左右深度简单取最小决定空节点、叶子和极端链式树如何处理。这道题的边界不能留到最后补丁式处理，实验时覆盖只有左子树、只有右子树、根为空、根是叶子，尤其关注空节点、单节点、链式树和全负值。递归版本和显式栈版本可以互相对拍。

C++ 实现上，极端深树可能造成递归栈风险，工程写法可以用 \code{std::vector<TreeNode*>} 或队列显式保存状态。BST 边界最好用更宽整数或可选前驱值。如果追问变成“DFS 写法必须特殊处理空子树不能作为最小路径”，先判断原来的状态是否仍然覆盖新答案集合，再决定是微调边界、改 value 语义，还是换成另一类结构。

\subsubsection{LeetCode 112 路径总和}

先想最直接的树遍历会在每个节点重复求什么。本题的模型是路径必须从根到叶子，状态是剩余目标值；如果每个节点都重新扫描子树，慢点就会非常明显。本题要压缩的状态是沿路径减去当前节点值，到叶子时检查剩余是否等于叶子值，也就是节点和父子方向之间真正传递的信息。

树题要先判定状态方向。本题的沿路径减去当前节点值，到叶子时检查剩余是否等于叶子值可能是自顶向下的约束、后序向上的汇总，或队列中的层号。返回值或栈元素不能只有节点指针，还要带上题目需要的信息。放到“二叉树与树形状态”这条主线里看，关键原则是：树形 DP 把每个节点看成子问题根；BFS 把层次边界显式放进队列；BST 题利用中序有序或上下界约束。工程实现优先考虑显式栈避免深树调用栈风险。

树题常用对子树归纳。假设子节点状态已经正确，父节点只按沿路径减去当前节点值，到叶子时检查剩余是否等于叶子值合并或传递信息。本题的路径必须从根到叶子，状态是剩余目标值决定空节点、叶子和极端链式树如何处理。这道题的边界不能留到最后补丁式处理，实验时覆盖负数节点、目标为零、非叶子提前达到目标、空树，尤其关注空节点、单节点、链式树和全负值。递归版本和显式栈版本可以互相对拍。

C++ 实现上，极端深树可能造成递归栈风险，工程写法可以用 \code{std::vector<TreeNode*>} 或队列显式保存状态。BST 边界最好用更宽整数或可选前驱值。如果追问变成“若要求列出所有路径，需要保存路径数组并在回退时弹出”，先判断原来的状态是否仍然覆盖新答案集合，再决定是微调边界、改 value 语义，还是换成另一类结构。

\subsubsection{LeetCode 113 路径总和 II}

先想最直接的树遍历会在每个节点重复求什么。本题的模型是相比判定题，状态还要保存当前路径；如果每个节点都重新扫描子树，慢点就会非常明显。本题要压缩的状态是到叶子且和满足时复制路径；回溯时恢复路径，也就是节点和父子方向之间真正传递的信息。

树题要先判定状态方向。本题的到叶子且和满足时复制路径；回溯时恢复路径可能是自顶向下的约束、后序向上的汇总，或队列中的层号。返回值或栈元素不能只有节点指针，还要带上题目需要的信息。放到“二叉树与树形状态”这条主线里看，关键原则是：树形 DP 把每个节点看成子问题根；BFS 把层次边界显式放进队列；BST 题利用中序有序或上下界约束。工程实现优先考虑显式栈避免深树调用栈风险。

树题常用对子树归纳。假设子节点状态已经正确，父节点只按到叶子且和满足时复制路径；回溯时恢复路径合并或传递信息。本题的相比判定题，状态还要保存当前路径决定空节点、叶子和极端链式树如何处理。这道题的边界不能留到最后补丁式处理，实验时覆盖多个答案、负数、空树、路径数组复制时机，尤其关注空节点、单节点、链式树和全负值。递归版本和显式栈版本可以互相对拍。

C++ 实现上，极端深树可能造成递归栈风险，工程写法可以用 \code{std::vector<TreeNode*>} 或队列显式保存状态。BST 边界最好用更宽整数或可选前驱值。如果追问变成“若路径不要求从根到叶，模型会变成前缀和加哈希”，先判断原来的状态是否仍然覆盖新答案集合，再决定是微调边界、改 value 语义，还是换成另一类结构。

\subsubsection{LeetCode 114 二叉树展开为链表}

先想最直接的树遍历会在每个节点重复求什么。本题的模型是目标是按前序顺序把树原地改成右指针链；如果每个节点都重新扫描子树，慢点就会非常明显。本题要压缩的状态是显式栈按右后左先入顺序处理，逐个接到前驱右侧，也就是节点和父子方向之间真正传递的信息。

树题要先判定状态方向。本题的显式栈按右后左先入顺序处理，逐个接到前驱右侧可能是自顶向下的约束、后序向上的汇总，或队列中的层号。返回值或栈元素不能只有节点指针，还要带上题目需要的信息。放到“二叉树与树形状态”这条主线里看，关键原则是：树形 DP 把每个节点看成子问题根；BFS 把层次边界显式放进队列；BST 题利用中序有序或上下界约束。工程实现优先考虑显式栈避免深树调用栈风险。

树题常用对子树归纳。假设子节点状态已经正确，父节点只按显式栈按右后左先入顺序处理，逐个接到前驱右侧合并或传递信息。本题的目标是按前序顺序把树原地改成右指针链决定空节点、叶子和极端链式树如何处理。这道题的边界不能留到最后补丁式处理，实验时覆盖空树、单节点、原有右子树不能丢、左指针置空，尤其关注空节点、单节点、链式树和全负值。递归版本和显式栈版本可以互相对拍。

C++ 实现上，极端深树可能造成递归栈风险，工程写法可以用 \code{std::vector<TreeNode*>} 或队列显式保存状态。BST 边界最好用更宽整数或可选前驱值。如果追问变成“也可用寻找左子树最右节点的原地方法”，先判断原来的状态是否仍然覆盖新答案集合，再决定是微调边界、改 value 语义，还是换成另一类结构。

\subsubsection{LeetCode 124 二叉树最大路径和}

先想最直接的树遍历会在每个节点重复求什么。本题的模型是路径可以从任意节点到任意节点，但不能分叉向父节点贡献两边；如果每个节点都重新扫描子树，慢点就会非常明显。本题要压缩的状态是每个节点向父亲贡献单边最大增益，全局答案可用左右两边加当前节点更新，也就是节点和父子方向之间真正传递的信息。

树题要先判定状态方向。本题的每个节点向父亲贡献单边最大增益，全局答案可用左右两边加当前节点更新可能是自顶向下的约束、后序向上的汇总，或队列中的层号。返回值或栈元素不能只有节点指针，还要带上题目需要的信息。放到“二叉树与树形状态”这条主线里看，关键原则是：树形 DP 把每个节点看成子问题根；BFS 把层次边界显式放进队列；BST 题利用中序有序或上下界约束。工程实现优先考虑显式栈避免深树调用栈风险。

树题常用对子树归纳。假设子节点状态已经正确，父节点只按每个节点向父亲贡献单边最大增益，全局答案可用左右两边加当前节点更新合并或传递信息。本题的路径可以从任意节点到任意节点，但不能分叉向父节点贡献两边决定空节点、叶子和极端链式树如何处理。这道题的边界不能留到最后补丁式处理，实验时覆盖全负数、单节点、左右增益为负时取零、答案初始值，尤其关注空节点、单节点、链式树和全负值。递归版本和显式栈版本可以互相对拍。

C++ 实现上，极端深树可能造成递归栈风险，工程写法可以用 \code{std::vector<TreeNode*>} 或队列显式保存状态。BST 边界最好用更宽整数或可选前驱值。如果追问变成“若路径必须从根到叶，状态和答案位置完全不同”，先判断原来的状态是否仍然覆盖新答案集合，再决定是微调边界、改 value 语义，还是换成另一类结构。

\subsubsection{LeetCode 199 二叉树右视图}

先想最直接的树遍历会在每个节点重复求什么。本题的模型是每层从右侧能看到的就是该层最后访问或最先访问的右侧节点；如果每个节点都重新扫描子树，慢点就会非常明显。本题要压缩的状态是BFS 固定层边界，取每层最后一个；或 DFS 按右优先首次到达深度，也就是节点和父子方向之间真正传递的信息。

树题要先判定状态方向。本题的BFS 固定层边界，取每层最后一个；或 DFS 按右优先首次到达深度可能是自顶向下的约束、后序向上的汇总，或队列中的层号。返回值或栈元素不能只有节点指针，还要带上题目需要的信息。放到“二叉树与树形状态”这条主线里看，关键原则是：树形 DP 把每个节点看成子问题根；BFS 把层次边界显式放进队列；BST 题利用中序有序或上下界约束。工程实现优先考虑显式栈避免深树调用栈风险。

树题常用对子树归纳。假设子节点状态已经正确，父节点只按BFS 固定层边界，取每层最后一个；或 DFS 按右优先首次到达深度合并或传递信息。本题的每层从右侧能看到的就是该层最后访问或最先访问的右侧节点决定空节点、叶子和极端链式树如何处理。这道题的边界不能留到最后补丁式处理，实验时覆盖空树、左链、右链、左右交错，尤其关注空节点、单节点、链式树和全负值。递归版本和显式栈版本可以互相对拍。

C++ 实现上，极端深树可能造成递归栈风险，工程写法可以用 \code{std::vector<TreeNode*>} 或队列显式保存状态。BST 边界最好用更宽整数或可选前驱值。如果追问变成“若求左视图，只改变层内取值方向”，先判断原来的状态是否仍然覆盖新答案集合，再决定是微调边界、改 value 语义，还是换成另一类结构。

\subsubsection{LeetCode 208 实现 Trie}

先想最直接的树遍历会在每个节点重复求什么。本题的模型是前缀树把字符串公共前缀共享成路径；如果每个节点都重新扫描子树，慢点就会非常明显。本题要压缩的状态是插入和查询都逐字符沿边移动，不存在的边按需要创建，也就是节点和父子方向之间真正传递的信息。

树题要先判定状态方向。本题的插入和查询都逐字符沿边移动，不存在的边按需要创建可能是自顶向下的约束、后序向上的汇总，或队列中的层号。返回值或栈元素不能只有节点指针，还要带上题目需要的信息。放到“二叉树与树形状态”这条主线里看，关键原则是：树形 DP 把每个节点看成子问题根；BFS 把层次边界显式放进队列；BST 题利用中序有序或上下界约束。工程实现优先考虑显式栈避免深树调用栈风险。

树题常用对子树归纳。假设子节点状态已经正确，父节点只按插入和查询都逐字符沿边移动，不存在的边按需要创建合并或传递信息。本题的前缀树把字符串公共前缀共享成路径决定空节点、叶子和极端链式树如何处理。这道题的边界不能留到最后补丁式处理，实验时覆盖空字符串约定、单词结束标记、前缀存在不等于单词存在，尤其关注空节点、单节点、链式树和全负值。递归版本和显式栈版本可以互相对拍。

C++ 实现上，极端深树可能造成递归栈风险，工程写法可以用 \code{std::vector<TreeNode*>} 或队列显式保存状态。BST 边界最好用更宽整数或可选前驱值。如果追问变成“若字符集很大，用哈希子节点；小写字母用数组更快”，先判断原来的状态是否仍然覆盖新答案集合，再决定是微调边界、改 value 语义，还是换成另一类结构。

\subsubsection{LeetCode 211 添加与搜索单词}

先想最直接的树遍历会在每个节点重复求什么。本题的模型是点号通配符让查询可能分叉，但前缀树仍能剪掉不匹配前缀；如果每个节点都重新扫描子树，慢点就会非常明显。本题要压缩的状态是普通字符沿唯一边走，点号枚举当前节点所有孩子，也就是节点和父子方向之间真正传递的信息。

树题要先判定状态方向。本题的普通字符沿唯一边走，点号枚举当前节点所有孩子可能是自顶向下的约束、后序向上的汇总，或队列中的层号。返回值或栈元素不能只有节点指针，还要带上题目需要的信息。放到“二叉树与树形状态”这条主线里看，关键原则是：树形 DP 把每个节点看成子问题根；BFS 把层次边界显式放进队列；BST 题利用中序有序或上下界约束。工程实现优先考虑显式栈避免深树调用栈风险。

树题常用对子树归纳。假设子节点状态已经正确，父节点只按普通字符沿唯一边走，点号枚举当前节点所有孩子合并或传递信息。本题的点号通配符让查询可能分叉，但前缀树仍能剪掉不匹配前缀决定空节点、叶子和极端链式树如何处理。这道题的边界不能留到最后补丁式处理，实验时覆盖连续通配符、空结果、单词结束标记、字符集大小，尤其关注空节点、单节点、链式树和全负值。递归版本和显式栈版本可以互相对拍。

C++ 实现上，极端深树可能造成递归栈风险，工程写法可以用 \code{std::vector<TreeNode*>} 或队列显式保存状态。BST 边界最好用更宽整数或可选前驱值。如果追问变成“大量通配符会接近遍历整棵 Trie，需要规模约束”，先判断原来的状态是否仍然覆盖新答案集合，再决定是微调边界、改 value 语义，还是换成另一类结构。

\subsubsection{LeetCode 226 翻转二叉树}

先想最直接的树遍历会在每个节点重复求什么。本题的模型是每个节点交换左右孩子即可，遍历顺序不影响最终结构；如果每个节点都重新扫描子树，慢点就会非常明显。本题要压缩的状态是显式队列或栈逐个节点交换，避免递归深度，也就是节点和父子方向之间真正传递的信息。

树题要先判定状态方向。本题的显式队列或栈逐个节点交换，避免递归深度可能是自顶向下的约束、后序向上的汇总，或队列中的层号。返回值或栈元素不能只有节点指针，还要带上题目需要的信息。放到“二叉树与树形状态”这条主线里看，关键原则是：树形 DP 把每个节点看成子问题根；BFS 把层次边界显式放进队列；BST 题利用中序有序或上下界约束。工程实现优先考虑显式栈避免深树调用栈风险。

树题常用对子树归纳。假设子节点状态已经正确，父节点只按显式队列或栈逐个节点交换，避免递归深度合并或传递信息。本题的每个节点交换左右孩子即可，遍历顺序不影响最终结构决定空节点、叶子和极端链式树如何处理。这道题的边界不能留到最后补丁式处理，实验时覆盖空树、单节点、只有一侧子树、重复值，尤其关注空节点、单节点、链式树和全负值。递归版本和显式栈版本可以互相对拍。

C++ 实现上，极端深树可能造成递归栈风险，工程写法可以用 \code{std::vector<TreeNode*>} 或队列显式保存状态。BST 边界最好用更宽整数或可选前驱值。如果追问变成“这题简单但适合训练树节点指针修改”，先判断原来的状态是否仍然覆盖新答案集合，再决定是微调边界、改 value 语义，还是换成另一类结构。

\subsubsection{LeetCode 230 二叉搜索树中第 K 小}

先想最直接的树遍历会在每个节点重复求什么。本题的模型是BST 中序遍历按升序访问节点；如果每个节点都重新扫描子树，慢点就会非常明显。本题要压缩的状态是显式栈中序，访问第 k 个节点时返回，也就是节点和父子方向之间真正传递的信息。

树题要先判定状态方向。本题的显式栈中序，访问第 k 个节点时返回可能是自顶向下的约束、后序向上的汇总，或队列中的层号。返回值或栈元素不能只有节点指针，还要带上题目需要的信息。放到“二叉树与树形状态”这条主线里看，关键原则是：树形 DP 把每个节点看成子问题根；BFS 把层次边界显式放进队列；BST 题利用中序有序或上下界约束。工程实现优先考虑显式栈避免深树调用栈风险。

树题常用对子树归纳。假设子节点状态已经正确，父节点只按显式栈中序，访问第 k 个节点时返回合并或传递信息。本题的BST 中序遍历按升序访问节点决定空节点、叶子和极端链式树如何处理。这道题的边界不能留到最后补丁式处理，实验时覆盖k 为一、k 为节点数、树不平衡、是否有重复值，尤其关注空节点、单节点、链式树和全负值。递归版本和显式栈版本可以互相对拍。

C++ 实现上，极端深树可能造成递归栈风险，工程写法可以用 \code{std::vector<TreeNode*>} 或队列显式保存状态。BST 边界最好用更宽整数或可选前驱值。如果追问变成“若频繁查询，可在节点维护子树大小成为顺序统计树”，先判断原来的状态是否仍然覆盖新答案集合，再决定是微调边界、改 value 语义，还是换成另一类结构。

\subsubsection{LeetCode 235 二叉搜索树最近公共祖先}

先想最直接的树遍历会在每个节点重复求什么。本题的模型是BST 的值域关系能决定两个目标在当前节点的哪一侧；如果每个节点都重新扫描子树，慢点就会非常明显。本题要压缩的状态是两个值都小去左，都大去右，否则当前节点就是分叉点，也就是节点和父子方向之间真正传递的信息。

树题要先判定状态方向。本题的两个值都小去左，都大去右，否则当前节点就是分叉点可能是自顶向下的约束、后序向上的汇总，或队列中的层号。返回值或栈元素不能只有节点指针，还要带上题目需要的信息。放到“二叉树与树形状态”这条主线里看，关键原则是：树形 DP 把每个节点看成子问题根；BFS 把层次边界显式放进队列；BST 题利用中序有序或上下界约束。工程实现优先考虑显式栈避免深树调用栈风险。

树题常用对子树归纳。假设子节点状态已经正确，父节点只按两个值都小去左，都大去右，否则当前节点就是分叉点合并或传递信息。本题的BST 的值域关系能决定两个目标在当前节点的哪一侧决定空节点、叶子和极端链式树如何处理。这道题的边界不能留到最后补丁式处理，实验时覆盖一个节点是另一个祖先、目标顺序、重复值约定，尤其关注空节点、单节点、链式树和全负值。递归版本和显式栈版本可以互相对拍。

C++ 实现上，极端深树可能造成递归栈风险，工程写法可以用 \code{std::vector<TreeNode*>} 或队列显式保存状态。BST 边界最好用更宽整数或可选前驱值。如果追问变成“普通二叉树不能用值域剪枝，需要父指针或后序状态”，先判断原来的状态是否仍然覆盖新答案集合，再决定是微调边界、改 value 语义，还是换成另一类结构。

\subsubsection{LeetCode 236 二叉树最近公共祖先}

先想最直接的树遍历会在每个节点重复求什么。本题的模型是普通树没有有序性，只能依靠祖先关系；如果每个节点都重新扫描子树，慢点就会非常明显。本题要压缩的状态是父指针集合或后序返回目标命中状态都可行，也就是节点和父子方向之间真正传递的信息。

树题要先判定状态方向。本题的父指针集合或后序返回目标命中状态都可行可能是自顶向下的约束、后序向上的汇总，或队列中的层号。返回值或栈元素不能只有节点指针，还要带上题目需要的信息。放到“二叉树与树形状态”这条主线里看，关键原则是：树形 DP 把每个节点看成子问题根；BFS 把层次边界显式放进队列；BST 题利用中序有序或上下界约束。工程实现优先考虑显式栈避免深树调用栈风险。

树题常用对子树归纳。假设子节点状态已经正确，父节点只按父指针集合或后序返回目标命中状态都可行合并或传递信息。本题的普通树没有有序性，只能依靠祖先关系决定空节点、叶子和极端链式树如何处理。这道题的边界不能留到最后补丁式处理，实验时覆盖目标不存在、p 等于 q、一个是另一个祖先、比较节点地址，尤其关注空节点、单节点、链式树和全负值。递归版本和显式栈版本可以互相对拍。

C++ 实现上，极端深树可能造成递归栈风险，工程写法可以用 \code{std::vector<TreeNode*>} 或队列显式保存状态。BST 边界最好用更宽整数或可选前驱值。如果追问变成“多次 LCA 查询可预处理倍增祖先表”，先判断原来的状态是否仍然覆盖新答案集合，再决定是微调边界、改 value 语义，还是换成另一类结构。

\subsubsection{LeetCode 337 打家劫舍 III}

先想最直接的树遍历会在每个节点重复求什么。本题的模型是树上相邻节点不能同时选，每个节点需要偷和不偷两个状态；如果每个节点都重新扫描子树，慢点就会非常明显。本题要压缩的状态是后序汇总：偷当前则孩子不能偷，不偷当前则孩子可偷可不偷，也就是节点和父子方向之间真正传递的信息。

树题要先判定状态方向。本题的后序汇总：偷当前则孩子不能偷，不偷当前则孩子可偷可不偷可能是自顶向下的约束、后序向上的汇总，或队列中的层号。返回值或栈元素不能只有节点指针，还要带上题目需要的信息。放到“二叉树与树形状态”这条主线里看，关键原则是：树形 DP 把每个节点看成子问题根；BFS 把层次边界显式放进队列；BST 题利用中序有序或上下界约束。工程实现优先考虑显式栈避免深树调用栈风险。

树题常用对子树归纳。假设子节点状态已经正确，父节点只按后序汇总：偷当前则孩子不能偷，不偷当前则孩子可偷可不偷合并或传递信息。本题的树上相邻节点不能同时选，每个节点需要偷和不偷两个状态决定空节点、叶子和极端链式树如何处理。这道题的边界不能留到最后补丁式处理，实验时覆盖空树、负值是否可能、链式树、递归深度，尤其关注空节点、单节点、链式树和全负值。递归版本和显式栈版本可以互相对拍。

C++ 实现上，极端深树可能造成递归栈风险，工程写法可以用 \code{std::vector<TreeNode*>} 或队列显式保存状态。BST 边界最好用更宽整数或可选前驱值。如果追问变成“可用显式栈后序保存节点状态，避免调用栈”，先判断原来的状态是否仍然覆盖新答案集合，再决定是微调边界、改 value 语义，还是换成另一类结构。

\subsubsection{LeetCode 450 删除二叉搜索树中的节点}

先想最直接的树遍历会在每个节点重复求什么。本题的模型是删除节点后仍要保持 BST 中序有序；如果每个节点都重新扫描子树，慢点就会非常明显。本题要压缩的状态是若有两个孩子，用后继或前驱替换当前值，再删除那个后继节点，也就是节点和父子方向之间真正传递的信息。

树题要先判定状态方向。本题的若有两个孩子，用后继或前驱替换当前值，再删除那个后继节点可能是自顶向下的约束、后序向上的汇总，或队列中的层号。返回值或栈元素不能只有节点指针，还要带上题目需要的信息。放到“二叉树与树形状态”这条主线里看，关键原则是：树形 DP 把每个节点看成子问题根；BFS 把层次边界显式放进队列；BST 题利用中序有序或上下界约束。工程实现优先考虑显式栈避免深树调用栈风险。

树题常用对子树归纳。假设子节点状态已经正确，父节点只按若有两个孩子，用后继或前驱替换当前值，再删除那个后继节点合并或传递信息。本题的删除节点后仍要保持 BST 中序有序决定空节点、叶子和极端链式树如何处理。这道题的边界不能留到最后补丁式处理，实验时覆盖删除叶子、一个孩子、两个孩子、删除根节点，尤其关注空节点、单节点、链式树和全负值。递归版本和显式栈版本可以互相对拍。

C++ 实现上，极端深树可能造成递归栈风险，工程写法可以用 \code{std::vector<TreeNode*>} 或队列显式保存状态。BST 边界最好用更宽整数或可选前驱值。如果追问变成“工程实现要清楚是移动值还是重连节点”，先判断原来的状态是否仍然覆盖新答案集合，再决定是微调边界、改 value 语义，还是换成另一类结构。

\subsubsection{LeetCode 543 二叉树直径}

先想最直接的树遍历会在每个节点重复求什么。本题的模型是直径经过某节点时等于左高加右高，但向父亲只能贡献单边高度；如果每个节点都重新扫描子树，慢点就会非常明显。本题要压缩的状态是后序遍历同时更新全局直径并返回高度，也就是节点和父子方向之间真正传递的信息。

树题要先判定状态方向。本题的后序遍历同时更新全局直径并返回高度可能是自顶向下的约束、后序向上的汇总，或队列中的层号。返回值或栈元素不能只有节点指针，还要带上题目需要的信息。放到“二叉树与树形状态”这条主线里看，关键原则是：树形 DP 把每个节点看成子问题根；BFS 把层次边界显式放进队列；BST 题利用中序有序或上下界约束。工程实现优先考虑显式栈避免深树调用栈风险。

树题常用对子树归纳。假设子节点状态已经正确，父节点只按后序遍历同时更新全局直径并返回高度合并或传递信息。本题的直径经过某节点时等于左高加右高，但向父亲只能贡献单边高度决定空节点、叶子和极端链式树如何处理。这道题的边界不能留到最后补丁式处理，实验时覆盖空树、单节点、链式树、边数还是节点数定义，尤其关注空节点、单节点、链式树和全负值。递归版本和显式栈版本可以互相对拍。

C++ 实现上，极端深树可能造成递归栈风险，工程写法可以用 \code{std::vector<TreeNode*>} 或队列显式保存状态。BST 边界最好用更宽整数或可选前驱值。如果追问变成“最大路径和是同型题，但贡献值和全局更新有负数处理”，先判断原来的状态是否仍然覆盖新答案集合，再决定是微调边界、改 value 语义，还是换成另一类结构。

\subsection{图建模、BFS 与 DFS工作坊}

先定义节点、边和状态，再决定 BFS、DFS、拓扑或并查集。图题失败常常不是搜索模板错，而是状态维度不完整或邻居生成过慢。

\subsubsection{LeetCode 127 单词接龙}

先画图，不先选搜索模板。本题的模型是每个单词是节点，一次修改一个字符形成边；朴素做法通常从多个起点反复搜索，或者把状态维度压得太少。优化首先要围绕通配模式桶加速邻居查询，BFS 第一次到达终点就是最短转换长度明确节点、边和访问标记。

图状态由节点和附加资源共同组成。本题的通配模式桶加速邻居查询，BFS 第一次到达终点就是最短转换长度要决定访问标记何时设置、状态是否只按位置去重、邻居如何生成。建模错了，BFS 或 DFS 再标准也会错。放到“图建模、BFS 与 DFS”这条主线里看，关键原则是：无权最短步数用 BFS，连通块和状态检测用 DFS 或显式栈，有向无环依赖用拓扑排序。多源问题把所有源同时入队，状态带资源的问题不能只按位置去重。

图搜索证明要回到可达性或层次。一次搜索必须覆盖从起点按每个单词是节点，一次修改一个字符形成边可达的全部状态，同时不能越过非法边；通配模式桶加速邻居查询，BFS 第一次到达终点就是最短转换长度保证重复状态不会反复入队或入栈。这道题的边界不能留到最后补丁式处理，实验时覆盖终点不在词表、起点和终点相同、桶重复扫描、单词长度一致，并打印入队时的完整状态。若同一位置但资源不同的状态被合并，很多图题会出现隐蔽错误。

C++ 实现上，连续编号图用 \code{std::vector<std::vector<int>>}，字符串节点先编号或使用哈希表。BFS 通常入队时标记访问，方向数组用固定常量集中维护。如果追问变成“双向 BFS 可以进一步降低搜索空间”，先判断原来的状态是否仍然覆盖新答案集合，再决定是微调边界、改 value 语义，还是换成另一类结构。

\subsubsection{LeetCode 130 被围绕的区域}

先画图，不先选搜索模板。本题的模型是只有不连通到边界的 O 才会被翻转；朴素做法通常从多个起点反复搜索，或者把状态维度压得太少。优化首先要围绕从边界 O 出发标记安全区域，再翻转剩余 O明确节点、边和访问标记。

图状态由节点和附加资源共同组成。本题的从边界 O 出发标记安全区域，再翻转剩余 O要决定访问标记何时设置、状态是否只按位置去重、邻居如何生成。建模错了，BFS 或 DFS 再标准也会错。放到“图建模、BFS 与 DFS”这条主线里看，关键原则是：无权最短步数用 BFS，连通块和状态检测用 DFS 或显式栈，有向无环依赖用拓扑排序。多源问题把所有源同时入队，状态带资源的问题不能只按位置去重。

图搜索证明要回到可达性或层次。一次搜索必须覆盖从起点按只有不连通到边界的 O 才会被翻转可达的全部状态，同时不能越过非法边；从边界 O 出发标记安全区域，再翻转剩余 O保证重复状态不会反复入队或入栈。这道题的边界不能留到最后补丁式处理，实验时覆盖全边界、无 O、单行单列、内部岛屿，并打印入队时的完整状态。若同一位置但资源不同的状态被合并，很多图题会出现隐蔽错误。

C++ 实现上，连续编号图用 \code{std::vector<std::vector<int>>}，字符串节点先编号或使用哈希表。BFS 通常入队时标记访问，方向数组用固定常量集中维护。如果追问变成“这是反向思维：找不能被翻转的区域比直接判断每块区域更简单”，先判断原来的状态是否仍然覆盖新答案集合，再决定是微调边界、改 value 语义，还是换成另一类结构。

\subsubsection{LeetCode 133 克隆图}

先画图，不先选搜索模板。本题的模型是每个原节点需要唯一对应一个新节点，边按原图连接；朴素做法通常从多个起点反复搜索，或者把状态维度压得太少。优化首先要围绕哈希表保存原节点到克隆节点的映射，BFS 或 DFS 遍历图明确节点、边和访问标记。

图状态由节点和附加资源共同组成。本题的哈希表保存原节点到克隆节点的映射，BFS 或 DFS 遍历图要决定访问标记何时设置、状态是否只按位置去重、邻居如何生成。建模错了，BFS 或 DFS 再标准也会错。放到“图建模、BFS 与 DFS”这条主线里看，关键原则是：无权最短步数用 BFS，连通块和状态检测用 DFS 或显式栈，有向无环依赖用拓扑排序。多源问题把所有源同时入队，状态带资源的问题不能只按位置去重。

图搜索证明要回到可达性或层次。一次搜索必须覆盖从起点按每个原节点需要唯一对应一个新节点，边按原图连接可达的全部状态，同时不能越过非法边；哈希表保存原节点到克隆节点的映射，BFS 或 DFS 遍历图保证重复状态不会反复入队或入栈。这道题的边界不能留到最后补丁式处理，实验时覆盖空图、自环、重边、环结构、不能按节点值唯一假设，并打印入队时的完整状态。若同一位置但资源不同的状态被合并，很多图题会出现隐蔽错误。

C++ 实现上，连续编号图用 \code{std::vector<std::vector<int>>}，字符串节点先编号或使用哈希表。BFS 通常入队时标记访问，方向数组用固定常量集中维护。如果追问变成“若图节点值不唯一，必须用地址作为键”，先判断原来的状态是否仍然覆盖新答案集合，再决定是微调边界、改 value 语义，还是换成另一类结构。

\subsubsection{LeetCode 200 岛屿数量}

先画图，不先选搜索模板。本题的模型是陆地格子是节点，四方向相邻是边，岛屿是连通块；朴素做法通常从多个起点反复搜索，或者把状态维度压得太少。优化首先要围绕扫描遇到未访问陆地就启动搜索并标记整块明确节点、边和访问标记。

图状态由节点和附加资源共同组成。本题的扫描遇到未访问陆地就启动搜索并标记整块要决定访问标记何时设置、状态是否只按位置去重、邻居如何生成。建模错了，BFS 或 DFS 再标准也会错。放到“图建模、BFS 与 DFS”这条主线里看，关键原则是：无权最短步数用 BFS，连通块和状态检测用 DFS 或显式栈，有向无环依赖用拓扑排序。多源问题把所有源同时入队，状态带资源的问题不能只按位置去重。

图搜索证明要回到可达性或层次。一次搜索必须覆盖从起点按陆地格子是节点，四方向相邻是边，岛屿是连通块可达的全部状态，同时不能越过非法边；扫描遇到未访问陆地就启动搜索并标记整块保证重复状态不会反复入队或入栈。这道题的边界不能留到最后补丁式处理，实验时覆盖空网格、全水、全陆、斜对角不连通、是否允许修改输入，并打印入队时的完整状态。若同一位置但资源不同的状态被合并，很多图题会出现隐蔽错误。

C++ 实现上，连续编号图用 \code{std::vector<std::vector<int>>}，字符串节点先编号或使用哈希表。BFS 通常入队时标记访问，方向数组用固定常量集中维护。如果追问变成“也可用并查集合并相邻陆地，适合动态岛屿扩展”，先判断原来的状态是否仍然覆盖新答案集合，再决定是微调边界、改 value 语义，还是换成另一类结构。

\subsubsection{LeetCode 207 课程表}

先画图，不先选搜索模板。本题的模型是课程是节点，先修关系是有向边，问题是判断是否有环；朴素做法通常从多个起点反复搜索，或者把状态维度压得太少。优化首先要围绕入度拓扑排序不断学习当前无依赖课程明确节点、边和访问标记。

图状态由节点和附加资源共同组成。本题的入度拓扑排序不断学习当前无依赖课程要决定访问标记何时设置、状态是否只按位置去重、邻居如何生成。建模错了，BFS 或 DFS 再标准也会错。放到“图建模、BFS 与 DFS”这条主线里看，关键原则是：无权最短步数用 BFS，连通块和状态检测用 DFS 或显式栈，有向无环依赖用拓扑排序。多源问题把所有源同时入队，状态带资源的问题不能只按位置去重。

图搜索证明要回到可达性或层次。一次搜索必须覆盖从起点按课程是节点，先修关系是有向边，问题是判断是否有环可达的全部状态，同时不能越过非法边；入度拓扑排序不断学习当前无依赖课程保证重复状态不会反复入队或入栈。这道题的边界不能留到最后补丁式处理，实验时覆盖边方向、孤立课程、环、自依赖、重复边，并打印入队时的完整状态。若同一位置但资源不同的状态被合并，很多图题会出现隐蔽错误。

C++ 实现上，连续编号图用 \code{std::vector<std::vector<int>>}，字符串节点先编号或使用哈希表。BFS 通常入队时标记访问，方向数组用固定常量集中维护。如果追问变成“DFS 三色标记也能判环，但要明确访问状态”，先判断原来的状态是否仍然覆盖新答案集合，再决定是微调边界、改 value 语义，还是换成另一类结构。

\subsubsection{LeetCode 210 课程表 II}

先画图，不先选搜索模板。本题的模型是在判环基础上还要输出一个可行拓扑序；朴素做法通常从多个起点反复搜索，或者把状态维度压得太少。优化首先要围绕每弹出一个入度为零课程，就把它加入顺序并删除出边明确节点、边和访问标记。

图状态由节点和附加资源共同组成。本题的每弹出一个入度为零课程，就把它加入顺序并删除出边要决定访问标记何时设置、状态是否只按位置去重、邻居如何生成。建模错了，BFS 或 DFS 再标准也会错。放到“图建模、BFS 与 DFS”这条主线里看，关键原则是：无权最短步数用 BFS，连通块和状态检测用 DFS 或显式栈，有向无环依赖用拓扑排序。多源问题把所有源同时入队，状态带资源的问题不能只按位置去重。

图搜索证明要回到可达性或层次。一次搜索必须覆盖从起点按在判环基础上还要输出一个可行拓扑序可达的全部状态，同时不能越过非法边；每弹出一个入度为零课程，就把它加入顺序并删除出边保证重复状态不会反复入队或入栈。这道题的边界不能留到最后补丁式处理，实验时覆盖有环返回空数组、多个可行顺序、边方向，并打印入队时的完整状态。若同一位置但资源不同的状态被合并，很多图题会出现隐蔽错误。

C++ 实现上，连续编号图用 \code{std::vector<std::vector<int>>}，字符串节点先编号或使用哈希表。BFS 通常入队时标记访问，方向数组用固定常量集中维护。如果追问变成“若需要字典序最小拓扑序，把队列换成小顶堆”，先判断原来的状态是否仍然覆盖新答案集合，再决定是微调边界、改 value 语义，还是换成另一类结构。

\subsubsection{LeetCode 417 太平洋大西洋水流问题}

先画图，不先选搜索模板。本题的模型是从每个格子向海搜索很慢，反向从海边沿非递减高度搜索；朴素做法通常从多个起点反复搜索，或者把状态维度压得太少。优化首先要围绕分别标记能到太平洋和大西洋的格子，交集就是答案明确节点、边和访问标记。

图状态由节点和附加资源共同组成。本题的分别标记能到太平洋和大西洋的格子，交集就是答案要决定访问标记何时设置、状态是否只按位置去重、邻居如何生成。建模错了，BFS 或 DFS 再标准也会错。放到“图建模、BFS 与 DFS”这条主线里看，关键原则是：无权最短步数用 BFS，连通块和状态检测用 DFS 或显式栈，有向无环依赖用拓扑排序。多源问题把所有源同时入队，状态带资源的问题不能只按位置去重。

图搜索证明要回到可达性或层次。一次搜索必须覆盖从起点按从每个格子向海搜索很慢，反向从海边沿非递减高度搜索可达的全部状态，同时不能越过非法边；分别标记能到太平洋和大西洋的格子，交集就是答案保证重复状态不会反复入队或入栈。这道题的边界不能留到最后补丁式处理，实验时覆盖单行单列、平台高度、边界重复入队、方向条件反向，并打印入队时的完整状态。若同一位置但资源不同的状态被合并，很多图题会出现隐蔽错误。

C++ 实现上，连续编号图用 \code{std::vector<std::vector<int>>}，字符串节点先编号或使用哈希表。BFS 通常入队时标记访问，方向数组用固定常量集中维护。如果追问变成“反向搜索是很多可达性题的重要技巧”，先判断原来的状态是否仍然覆盖新答案集合，再决定是微调边界、改 value 语义，还是换成另一类结构。

\subsubsection{LeetCode 695 岛屿的最大面积}

先画图，不先选搜索模板。本题的模型是面积就是一个连通块中陆地格子的数量；朴素做法通常从多个起点反复搜索，或者把状态维度压得太少。优化首先要围绕每发现新陆地就搜索完整连通块并计数，取最大值明确节点、边和访问标记。

图状态由节点和附加资源共同组成。本题的每发现新陆地就搜索完整连通块并计数，取最大值要决定访问标记何时设置、状态是否只按位置去重、邻居如何生成。建模错了，BFS 或 DFS 再标准也会错。放到“图建模、BFS 与 DFS”这条主线里看，关键原则是：无权最短步数用 BFS，连通块和状态检测用 DFS 或显式栈，有向无环依赖用拓扑排序。多源问题把所有源同时入队，状态带资源的问题不能只按位置去重。

图搜索证明要回到可达性或层次。一次搜索必须覆盖从起点按面积就是一个连通块中陆地格子的数量可达的全部状态，同时不能越过非法边；每发现新陆地就搜索完整连通块并计数，取最大值保证重复状态不会反复入队或入栈。这道题的边界不能留到最后补丁式处理，实验时覆盖全水、全陆、细长岛、是否允许修改输入，并打印入队时的完整状态。若同一位置但资源不同的状态被合并，很多图题会出现隐蔽错误。

C++ 实现上，连续编号图用 \code{std::vector<std::vector<int>>}，字符串节点先编号或使用哈希表。BFS 通常入队时标记访问，方向数组用固定常量集中维护。如果追问变成“与岛屿数量相比，只是连通块聚合值不同”，先判断原来的状态是否仍然覆盖新答案集合，再决定是微调边界、改 value 语义，还是换成另一类结构。

\subsubsection{LeetCode 752 打开转盘锁}

先画图，不先选搜索模板。本题的模型是四位密码是状态，每次转动一位形成边；朴素做法通常从多个起点反复搜索，或者把状态维度压得太少。优化首先要围绕BFS 从起点扩展，死亡状态不能入队，访问状态避免重复明确节点、边和访问标记。

图状态由节点和附加资源共同组成。本题的BFS 从起点扩展，死亡状态不能入队，访问状态避免重复要决定访问标记何时设置、状态是否只按位置去重、邻居如何生成。建模错了，BFS 或 DFS 再标准也会错。放到“图建模、BFS 与 DFS”这条主线里看，关键原则是：无权最短步数用 BFS，连通块和状态检测用 DFS 或显式栈，有向无环依赖用拓扑排序。多源问题把所有源同时入队，状态带资源的问题不能只按位置去重。

图搜索证明要回到可达性或层次。一次搜索必须覆盖从起点按四位密码是状态，每次转动一位形成边可达的全部状态，同时不能越过非法边；BFS 从起点扩展，死亡状态不能入队，访问状态避免重复保证重复状态不会反复入队或入栈。这道题的边界不能留到最后补丁式处理，实验时覆盖起点死亡、目标就是起点、数字环绕、死亡集合很大，并打印入队时的完整状态。若同一位置但资源不同的状态被合并，很多图题会出现隐蔽错误。

C++ 实现上，连续编号图用 \code{std::vector<std::vector<int>>}，字符串节点先编号或使用哈希表。BFS 通常入队时标记访问，方向数组用固定常量集中维护。如果追问变成“双向 BFS 能明显减少状态展开”，先判断原来的状态是否仍然覆盖新答案集合，再决定是微调边界、改 value 语义，还是换成另一类结构。

\subsubsection{LeetCode 864 获取所有钥匙的最短路径}

先画图，不先选搜索模板。本题的模型是位置相同但钥匙集合不同，未来能开的门不同；朴素做法通常从多个起点反复搜索，或者把状态维度压得太少。优化首先要围绕BFS 状态包含行、列和钥匙掩码，访问标记也必须包含掩码明确节点、边和访问标记。

图状态由节点和附加资源共同组成。本题的BFS 状态包含行、列和钥匙掩码，访问标记也必须包含掩码要决定访问标记何时设置、状态是否只按位置去重、邻居如何生成。建模错了，BFS 或 DFS 再标准也会错。放到“图建模、BFS 与 DFS”这条主线里看，关键原则是：无权最短步数用 BFS，连通块和状态检测用 DFS 或显式栈，有向无环依赖用拓扑排序。多源问题把所有源同时入队，状态带资源的问题不能只按位置去重。

图搜索证明要回到可达性或层次。一次搜索必须覆盖从起点按位置相同但钥匙集合不同，未来能开的门不同可达的全部状态，同时不能越过非法边；BFS 状态包含行、列和钥匙掩码，访问标记也必须包含掩码保证重复状态不会反复入队或入栈。这道题的边界不能留到最后补丁式处理，实验时覆盖起点旁边有门、钥匙顺序、不可达、钥匙数量决定掩码大小，并打印入队时的完整状态。若同一位置但资源不同的状态被合并，很多图题会出现隐蔽错误。

C++ 实现上，连续编号图用 \code{std::vector<std::vector<int>>}，字符串节点先编号或使用哈希表。BFS 通常入队时标记访问，方向数组用固定常量集中维护。如果追问变成“这是状态图维度不能只按位置去重的典型题”，先判断原来的状态是否仍然覆盖新答案集合，再决定是微调边界、改 value 语义，还是换成另一类结构。

\subsubsection{LeetCode 994 腐烂的橘子}

先画图，不先选搜索模板。本题的模型是所有初始腐烂橘子同时作为 BFS 源点；朴素做法通常从多个起点反复搜索，或者把状态维度压得太少。优化首先要围绕按层扩散，分钟数等于最后一个新鲜橘子被首次访问的层数明确节点、边和访问标记。

图状态由节点和附加资源共同组成。本题的按层扩散，分钟数等于最后一个新鲜橘子被首次访问的层数要决定访问标记何时设置、状态是否只按位置去重、邻居如何生成。建模错了，BFS 或 DFS 再标准也会错。放到“图建模、BFS 与 DFS”这条主线里看，关键原则是：无权最短步数用 BFS，连通块和状态检测用 DFS 或显式栈，有向无环依赖用拓扑排序。多源问题把所有源同时入队，状态带资源的问题不能只按位置去重。

图搜索证明要回到可达性或层次。一次搜索必须覆盖从起点按所有初始腐烂橘子同时作为 BFS 源点可达的全部状态，同时不能越过非法边；按层扩散，分钟数等于最后一个新鲜橘子被首次访问的层数保证重复状态不会反复入队或入栈。这道题的边界不能留到最后补丁式处理，实验时覆盖没有新鲜橘子、没有腐烂源、隔离新鲜橘子、空格阻断，并打印入队时的完整状态。若同一位置但资源不同的状态被合并，很多图题会出现隐蔽错误。

C++ 实现上，连续编号图用 \code{std::vector<std::vector<int>>}，字符串节点先编号或使用哈希表。BFS 通常入队时标记访问，方向数组用固定常量集中维护。如果追问变成“多源 BFS 也适用于最近零距离、地图扩散等题”，先判断原来的状态是否仍然覆盖新答案集合，再决定是微调边界、改 value 语义，还是换成另一类结构。

\subsubsection{LeetCode 1091 二进制矩阵中的最短路径}

先画图，不先选搜索模板。本题的模型是八方向网格无权边，答案是从左上到右下的最少格子数；朴素做法通常从多个起点反复搜索，或者把状态维度压得太少。优化首先要围绕BFS 入队时标记访问，层数或距离随状态保存明确节点、边和访问标记。

图状态由节点和附加资源共同组成。本题的BFS 入队时标记访问，层数或距离随状态保存要决定访问标记何时设置、状态是否只按位置去重、邻居如何生成。建模错了，BFS 或 DFS 再标准也会错。放到“图建模、BFS 与 DFS”这条主线里看，关键原则是：无权最短步数用 BFS，连通块和状态检测用 DFS 或显式栈，有向无环依赖用拓扑排序。多源问题把所有源同时入队，状态带资源的问题不能只按位置去重。

图搜索证明要回到可达性或层次。一次搜索必须覆盖从起点按八方向网格无权边，答案是从左上到右下的最少格子数可达的全部状态，同时不能越过非法边；BFS 入队时标记访问，层数或距离随状态保存保证重复状态不会反复入队或入栈。这道题的边界不能留到最后补丁式处理，实验时覆盖起点或终点阻塞、单格、八方向越界、无路可达，并打印入队时的完整状态。若同一位置但资源不同的状态被合并，很多图题会出现隐蔽错误。

C++ 实现上，连续编号图用 \code{std::vector<std::vector<int>>}，字符串节点先编号或使用哈希表。BFS 通常入队时标记访问，方向数组用固定常量集中维护。如果追问变成“边权相同才能用 BFS，若每格有代价就要最短路”，先判断原来的状态是否仍然覆盖新答案集合，再决定是微调边界、改 value 语义，还是换成另一类结构。

\subsubsection{LeetCode 1293 网格中的最短路径}

先画图，不先选搜索模板。本题的模型是走到同一格时，剩余可消除障碍数量不同会影响未来可达性；朴素做法通常从多个起点反复搜索，或者把状态维度压得太少。优化首先要围绕BFS 状态包含位置和剩余消除次数，visited 记录该维度或记录到达该格的最大剩余次数明确节点、边和访问标记。

图状态由节点和附加资源共同组成。本题的BFS 状态包含位置和剩余消除次数，visited 记录该维度或记录到达该格的最大剩余次数要决定访问标记何时设置、状态是否只按位置去重、邻居如何生成。建模错了，BFS 或 DFS 再标准也会错。放到“图建模、BFS 与 DFS”这条主线里看，关键原则是：无权最短步数用 BFS，连通块和状态检测用 DFS 或显式栈，有向无环依赖用拓扑排序。多源问题把所有源同时入队，状态带资源的问题不能只按位置去重。

图搜索证明要回到可达性或层次。一次搜索必须覆盖从起点按走到同一格时，剩余可消除障碍数量不同会影响未来可达性可达的全部状态，同时不能越过非法边；BFS 状态包含位置和剩余消除次数，visited 记录该维度或记录到达该格的最大剩余次数保证重复状态不会反复入队或入栈。这道题的边界不能留到最后补丁式处理，实验时覆盖起点终点、k 很大、障碍密集、二维 visited 错误剪枝，并打印入队时的完整状态。若同一位置但资源不同的状态被合并，很多图题会出现隐蔽错误。

C++ 实现上，连续编号图用 \code{std::vector<std::vector<int>>}，字符串节点先编号或使用哈希表。BFS 通常入队时标记访问，方向数组用固定常量集中维护。如果追问变成“它和钥匙题一样说明状态维度不能丢”，先判断原来的状态是否仍然覆盖新答案集合，再决定是微调边界、改 value 语义，还是换成另一类结构。

\subsection{最短路与带权图工作坊}

先确认目标是步数最少还是代价最小，再检查边权是否非负。非负权用 Dijkstra，等权用 BFS，限制边数用分层松弛或 Bellman-Ford 思想。

\subsubsection{LeetCode 505 迷宫 II}

先区分步数和代价。本题的模型是小球一旦选择方向就会滚到墙前停止，边权是滚动经过的格子数；若暴力枚举路径，路径数量会爆炸。优化点邻居生成不是走一步，而是沿四个方向模拟滚动到停点，再用 Dijkstra 按累计距离松弛说明下一步扩展必须按某种代价顺序组织，而不是普通队列随便推进。

最短路状态至少包含当前节点和当前代价。本题的邻居生成不是走一步，而是沿四个方向模拟滚动到停点，再用 Dijkstra 按累计距离松弛决定优先队列按什么排序，以及旧状态如何被识别。非负边条件、代价合成方式和不可达处理都要写清楚。放到“最短路与带权图”这条主线里看，关键原则是：非负权图用 Dijkstra，小顶堆保存当前距离最小的候选；边权为一用 BFS；有负边要考虑 Bellman-Ford 或重新建模。路径代价若是最大边权，也可用 Dijkstra 或二分答案。

最短路证明依赖扩展顺序。若本题使用 Dijkstra，就要说明非负代价让弹出的未过期状态已经最优；若用二分可达性，就要说明阈值越大可用边只会增加。这道题的边界不能留到最后补丁式处理，实验时覆盖起点就是终点、目标无法停下、绕远路更短、同一停点多次被松弛，打印每次弹出的代价和是否过期。若普通 BFS 与 Dijkstra 结果不同，要能解释边权条件造成的差异。

C++ 实现上，Dijkstra 用邻接表保存 \code{(to, weight)}，距离数组用 \code{long long}。由于 \code{priority_queue} 没有 decrease-key，压入新状态后弹出时跳过旧状态。如果追问变成“若只问能否到达，可用 BFS 或 DFS；若问最少距离，就必须处理带权边”，先判断原来的状态是否仍然覆盖新答案集合，再决定是微调边界、改 value 语义，还是换成另一类结构。

\subsubsection{LeetCode 743 网络延迟时间}

先区分步数和代价。本题的模型是有向正权图从起点到所有节点的最短路最大值；若暴力枚举路径，路径数量会爆炸。优化点Dijkstra 求每个节点最短距离，最后取最大，若有无穷则不可达说明下一步扩展必须按某种代价顺序组织，而不是普通队列随便推进。

最短路状态至少包含当前节点和当前代价。本题的Dijkstra 求每个节点最短距离，最后取最大，若有无穷则不可达决定优先队列按什么排序，以及旧状态如何被识别。非负边条件、代价合成方式和不可达处理都要写清楚。放到“最短路与带权图”这条主线里看，关键原则是：非负权图用 Dijkstra，小顶堆保存当前距离最小的候选；边权为一用 BFS；有负边要考虑 Bellman-Ford 或重新建模。路径代价若是最大边权，也可用 Dijkstra 或二分答案。

最短路证明依赖扩展顺序。若本题使用 Dijkstra，就要说明非负代价让弹出的未过期状态已经最优；若用二分可达性，就要说明阈值越大可用边只会增加。这道题的边界不能留到最后补丁式处理，实验时覆盖节点编号从一开始、不可达、重边、过期堆元素，打印每次弹出的代价和是否过期。若普通 BFS 与 Dijkstra 结果不同，要能解释边权条件造成的差异。

C++ 实现上，Dijkstra 用邻接表保存 \code{(to, weight)}，距离数组用 \code{long long}。由于 \code{priority_queue} 没有 decrease-key，压入新状态后弹出时跳过旧状态。如果追问变成“边权相同才可用普通 BFS”，先判断原来的状态是否仍然覆盖新答案集合，再决定是微调边界、改 value 语义，还是换成另一类结构。

\subsubsection{LeetCode 778 水位上升的泳池中游泳}

先区分步数和代价。本题的模型是到达某格子的代价是路径上最大高度，而不是高度和；若暴力枚举路径，路径数量会爆炸。优化点用 Dijkstra 按当前路径最大高度扩展，松弛时取 max 当前代价和邻格高度说明下一步扩展必须按某种代价顺序组织，而不是普通队列随便推进。

最短路状态至少包含当前节点和当前代价。本题的用 Dijkstra 按当前路径最大高度扩展，松弛时取 max 当前代价和邻格高度决定优先队列按什么排序，以及旧状态如何被识别。非负边条件、代价合成方式和不可达处理都要写清楚。放到“最短路与带权图”这条主线里看，关键原则是：非负权图用 Dijkstra，小顶堆保存当前距离最小的候选；边权为一用 BFS；有负边要考虑 Bellman-Ford 或重新建模。路径代价若是最大边权，也可用 Dijkstra 或二分答案。

最短路证明依赖扩展顺序。若本题使用 Dijkstra，就要说明非负代价让弹出的未过期状态已经最优；若用二分可达性，就要说明阈值越大可用边只会增加。这道题的边界不能留到最后补丁式处理，实验时覆盖单格、起点高度很高、需要绕路、多个相同高度，打印每次弹出的代价和是否过期。若普通 BFS 与 Dijkstra 结果不同，要能解释边权条件造成的差异。

C++ 实现上，Dijkstra 用邻接表保存 \code{(to, weight)}，距离数组用 \code{long long}。由于 \code{priority_queue} 没有 decrease-key，压入新状态后弹出时跳过旧状态。如果追问变成“也可以按水位二分可达，或按高度排序并查集合并”，先判断原来的状态是否仍然覆盖新答案集合，再决定是微调边界、改 value 语义，还是换成另一类结构。

\subsubsection{LeetCode 787 K 站中转内最便宜的航班}

先区分步数和代价。本题的模型是路径代价受中转次数限制，不能只按费用最小弹出节点后永久确定；若暴力枚举路径，路径数量会爆炸。优化点用按边数分层的 Bellman-Ford 或带层数状态的优先队列，状态包含城市和已用边数说明下一步扩展必须按某种代价顺序组织，而不是普通队列随便推进。

最短路状态至少包含当前节点和当前代价。本题的用按边数分层的 Bellman-Ford 或带层数状态的优先队列，状态包含城市和已用边数决定优先队列按什么排序，以及旧状态如何被识别。非负边条件、代价合成方式和不可达处理都要写清楚。放到“最短路与带权图”这条主线里看，关键原则是：非负权图用 Dijkstra，小顶堆保存当前距离最小的候选；边权为一用 BFS；有负边要考虑 Bellman-Ford 或重新建模。路径代价若是最大边权，也可用 Dijkstra 或二分答案。

最短路证明依赖扩展顺序。若本题使用 Dijkstra，就要说明非负代价让弹出的未过期状态已经最优；若用二分可达性，就要说明阈值越大可用边只会增加。这道题的边界不能留到最后补丁式处理，实验时覆盖起点终点相同、不可达、K 为零、便宜路径中转过多，打印每次弹出的代价和是否过期。若普通 BFS 与 Dijkstra 结果不同，要能解释边权条件造成的差异。

C++ 实现上，Dijkstra 用邻接表保存 \code{(to, weight)}，距离数组用 \code{long long}。由于 \code{priority_queue} 没有 decrease-key，压入新状态后弹出时跳过旧状态。如果追问变成“若去掉中转限制，普通 Dijkstra 才能直接按费用组织扩展”，先判断原来的状态是否仍然覆盖新答案集合，再决定是微调边界、改 value 语义，还是换成另一类结构。

\subsubsection{LeetCode 1102 得分最高的路径}

先区分步数和代价。本题的模型是路径得分是沿途最小格子值，要最大化这个最小值；若暴力枚举路径，路径数量会爆炸。优化点用最大堆按当前路径得分扩展，进入邻格时取当前得分和邻格值的较小者说明下一步扩展必须按某种代价顺序组织，而不是普通队列随便推进。

最短路状态至少包含当前节点和当前代价。本题的用最大堆按当前路径得分扩展，进入邻格时取当前得分和邻格值的较小者决定优先队列按什么排序，以及旧状态如何被识别。非负边条件、代价合成方式和不可达处理都要写清楚。放到“最短路与带权图”这条主线里看，关键原则是：非负权图用 Dijkstra，小顶堆保存当前距离最小的候选；边权为一用 BFS；有负边要考虑 Bellman-Ford 或重新建模。路径代价若是最大边权，也可用 Dijkstra 或二分答案。

最短路证明依赖扩展顺序。若本题使用 Dijkstra，就要说明非负代价让弹出的未过期状态已经最优；若用二分可达性，就要说明阈值越大可用边只会增加。这道题的边界不能留到最后补丁式处理，实验时覆盖单格、起点或终点很小、必须绕路、多个相同得分路径，打印每次弹出的代价和是否过期。若普通 BFS 与 Dijkstra 结果不同，要能解释边权条件造成的差异。

C++ 实现上，Dijkstra 用邻接表保存 \code{(to, weight)}，距离数组用 \code{long long}。由于 \code{priority_queue} 没有 decrease-key，压入新状态后弹出时跳过旧状态。如果追问变成“也可把阈值二分后检查连通性，或按值从高到低并查集合并”，先判断原来的状态是否仍然覆盖新答案集合，再决定是微调边界、改 value 语义，还是换成另一类结构。

\subsubsection{LeetCode 1334 阈值距离内邻居最少的城市}

先区分步数和代价。本题的模型是每个城市都需要知道到其他城市的最短距离是否不超过阈值；若暴力枚举路径，路径数量会爆炸。优化点节点数较小时可用 Floyd，边稀疏时也可从每个点跑 Dijkstra说明下一步扩展必须按某种代价顺序组织，而不是普通队列随便推进。

最短路状态至少包含当前节点和当前代价。本题的节点数较小时可用 Floyd，边稀疏时也可从每个点跑 Dijkstra决定优先队列按什么排序，以及旧状态如何被识别。非负边条件、代价合成方式和不可达处理都要写清楚。放到“最短路与带权图”这条主线里看，关键原则是：非负权图用 Dijkstra，小顶堆保存当前距离最小的候选；边权为一用 BFS；有负边要考虑 Bellman-Ford 或重新建模。路径代价若是最大边权，也可用 Dijkstra 或二分答案。

最短路证明依赖扩展顺序。若本题使用 Dijkstra，就要说明非负代价让弹出的未过期状态已经最优；若用二分可达性，就要说明阈值越大可用边只会增加。这道题的边界不能留到最后补丁式处理，实验时覆盖距离等于阈值、并列时选编号最大、不可达、无向边，打印每次弹出的代价和是否过期。若普通 BFS 与 Dijkstra 结果不同，要能解释边权条件造成的差异。

C++ 实现上，Dijkstra 用邻接表保存 \code{(to, weight)}，距离数组用 \code{long long}。由于 \code{priority_queue} 没有 decrease-key，压入新状态后弹出时跳过旧状态。如果追问变成“这是多源最短路选择问题，算法取决于 n、边数和查询密度”，先判断原来的状态是否仍然覆盖新答案集合，再决定是微调边界、改 value 语义，还是换成另一类结构。

\subsubsection{LeetCode 1368 使网格图至少有一条有效路径的最小代价}

先区分步数和代价。本题的模型是沿着格子箭头走代价为零，改方向走代价为一；若暴力枚举路径，路径数量会爆炸。优化点边权只有零和一，用 0-1 BFS：零代价邻居入队首，一代价邻居入队尾说明下一步扩展必须按某种代价顺序组织，而不是普通队列随便推进。

最短路状态至少包含当前节点和当前代价。本题的边权只有零和一，用 0-1 BFS：零代价邻居入队首，一代价邻居入队尾决定优先队列按什么排序，以及旧状态如何被识别。非负边条件、代价合成方式和不可达处理都要写清楚。放到“最短路与带权图”这条主线里看，关键原则是：非负权图用 Dijkstra，小顶堆保存当前距离最小的候选；边权为一用 BFS；有负边要考虑 Bellman-Ford 或重新建模。路径代价若是最大边权，也可用 Dijkstra 或二分答案。

最短路证明依赖扩展顺序。若本题使用 Dijkstra，就要说明非负代价让弹出的未过期状态已经最优；若用二分可达性，就要说明阈值越大可用边只会增加。这道题的边界不能留到最后补丁式处理，实验时覆盖单格、箭头指向越界、多个零代价路径、过期状态，打印每次弹出的代价和是否过期。若普通 BFS 与 Dijkstra 结果不同，要能解释边权条件造成的差异。

C++ 实现上，Dijkstra 用邻接表保存 \code{(to, weight)}，距离数组用 \code{long long}。由于 \code{priority_queue} 没有 decrease-key，压入新状态后弹出时跳过旧状态。如果追问变成“普通 BFS 按步数扩展，不等于总代价最小”，先判断原来的状态是否仍然覆盖新答案集合，再决定是微调边界、改 value 语义，还是换成另一类结构。

\subsubsection{LeetCode 1514 概率最大的路径}

先区分步数和代价。本题的模型是路径概率是边概率乘积，目标是最大乘积路径；若暴力枚举路径，路径数量会爆炸。优化点用大顶堆每次扩展当前概率最高节点，乘以边概率得到候选概率说明下一步扩展必须按某种代价顺序组织，而不是普通队列随便推进。

最短路状态至少包含当前节点和当前代价。本题的用大顶堆每次扩展当前概率最高节点，乘以边概率得到候选概率决定优先队列按什么排序，以及旧状态如何被识别。非负边条件、代价合成方式和不可达处理都要写清楚。放到“最短路与带权图”这条主线里看，关键原则是：非负权图用 Dijkstra，小顶堆保存当前距离最小的候选；边权为一用 BFS；有负边要考虑 Bellman-Ford 或重新建模。路径代价若是最大边权，也可用 Dijkstra 或二分答案。

最短路证明依赖扩展顺序。若本题使用 Dijkstra，就要说明非负代价让弹出的未过期状态已经最优；若用二分可达性，就要说明阈值越大可用边只会增加。这道题的边界不能留到最后补丁式处理，实验时覆盖不可达、概率为零、起点等于终点、浮点比较，打印每次弹出的代价和是否过期。若普通 BFS 与 Dijkstra 结果不同，要能解释边权条件造成的差异。

C++ 实现上，Dijkstra 用邻接表保存 \code{(to, weight)}，距离数组用 \code{long long}。由于 \code{priority_queue} 没有 decrease-key，压入新状态后弹出时跳过旧状态。如果追问变成“取负对数后可转成最短路，但直接最大堆更直观”，先判断原来的状态是否仍然覆盖新答案集合，再决定是微调边界、改 value 语义，还是换成另一类结构。

\subsubsection{LeetCode 1631 最小体力消耗路径}

先区分步数和代价。本题的模型是路径代价是沿途最大高度差，不是边权和；若暴力枚举路径，路径数量会爆炸。优化点Dijkstra 的松弛改成取当前代价和新边代价的最大值说明下一步扩展必须按某种代价顺序组织，而不是普通队列随便推进。

最短路状态至少包含当前节点和当前代价。本题的Dijkstra 的松弛改成取当前代价和新边代价的最大值决定优先队列按什么排序，以及旧状态如何被识别。非负边条件、代价合成方式和不可达处理都要写清楚。放到“最短路与带权图”这条主线里看，关键原则是：非负权图用 Dijkstra，小顶堆保存当前距离最小的候选；边权为一用 BFS；有负边要考虑 Bellman-Ford 或重新建模。路径代价若是最大边权，也可用 Dijkstra 或二分答案。

最短路证明依赖扩展顺序。若本题使用 Dijkstra，就要说明非负代价让弹出的未过期状态已经最优；若用二分可达性，就要说明阈值越大可用边只会增加。这道题的边界不能留到最后补丁式处理，实验时覆盖单格、平坦网格、必须绕路、多个相同代价，打印每次弹出的代价和是否过期。若普通 BFS 与 Dijkstra 结果不同，要能解释边权条件造成的差异。

C++ 实现上，Dijkstra 用邻接表保存 \code{(to, weight)}，距离数组用 \code{long long}。由于 \code{priority_queue} 没有 decrease-key，压入新状态后弹出时跳过旧状态。如果追问变成“也可二分体力上限后 BFS 检查可达”，先判断原来的状态是否仍然覆盖新答案集合，再决定是微调边界、改 value 语义，还是换成另一类结构。

\subsubsection{LeetCode 1928 规定时间内到达终点的最小花费}

先区分步数和代价。本题的模型是路径同时有时间约束和费用目标，单一距离无法描述状态；若暴力枚举路径，路径数量会爆炸。优化点状态包含节点和已用时间，DP 或 Dijkstra 在时间维度内松弛最小费用说明下一步扩展必须按某种代价顺序组织，而不是普通队列随便推进。

最短路状态至少包含当前节点和当前代价。本题的状态包含节点和已用时间，DP 或 Dijkstra 在时间维度内松弛最小费用决定优先队列按什么排序，以及旧状态如何被识别。非负边条件、代价合成方式和不可达处理都要写清楚。放到“最短路与带权图”这条主线里看，关键原则是：非负权图用 Dijkstra，小顶堆保存当前距离最小的候选；边权为一用 BFS；有负边要考虑 Bellman-Ford 或重新建模。路径代价若是最大边权，也可用 Dijkstra 或二分答案。

最短路证明依赖扩展顺序。若本题使用 Dijkstra，就要说明非负代价让弹出的未过期状态已经最优；若用二分可达性，就要说明阈值越大可用边只会增加。这道题的边界不能留到最后补丁式处理，实验时覆盖时间刚好到达、费用高但时间短、不可达、状态数量由 maxTime 决定，打印每次弹出的代价和是否过期。若普通 BFS 与 Dijkstra 结果不同，要能解释边权条件造成的差异。

C++ 实现上，Dijkstra 用邻接表保存 \code{(to, weight)}，距离数组用 \code{long long}。由于 \code{priority_queue} 没有 decrease-key，压入新状态后弹出时跳过旧状态。如果追问变成“这是资源约束最短路，不能只保存每个节点一个最小费用”，先判断原来的状态是否仍然覆盖新答案集合，再决定是微调边界、改 value 语义，还是换成另一类结构。

\subsubsection{LeetCode 1976 到达目的地的方案数}

先区分步数和代价。本题的模型是需要在最短距离之外统计有多少条最短路径；若暴力枚举路径，路径数量会爆炸。优化点Dijkstra 松弛时若得到更短距离就重置方案数，若等于最短距离就累加方案数说明下一步扩展必须按某种代价顺序组织，而不是普通队列随便推进。

最短路状态至少包含当前节点和当前代价。本题的Dijkstra 松弛时若得到更短距离就重置方案数，若等于最短距离就累加方案数决定优先队列按什么排序，以及旧状态如何被识别。非负边条件、代价合成方式和不可达处理都要写清楚。放到“最短路与带权图”这条主线里看，关键原则是：非负权图用 Dijkstra，小顶堆保存当前距离最小的候选；边权为一用 BFS；有负边要考虑 Bellman-Ford 或重新建模。路径代价若是最大边权，也可用 Dijkstra 或二分答案。

最短路证明依赖扩展顺序。若本题使用 Dijkstra，就要说明非负代价让弹出的未过期状态已经最优；若用二分可达性，就要说明阈值越大可用边只会增加。这道题的边界不能留到最后补丁式处理，实验时覆盖多条等长路径、取模、旧堆状态、不可达，打印每次弹出的代价和是否过期。若普通 BFS 与 Dijkstra 结果不同，要能解释边权条件造成的差异。

C++ 实现上，Dijkstra 用邻接表保存 \code{(to, weight)}，距离数组用 \code{long long}。由于 \code{priority_queue} 没有 decrease-key，压入新状态后弹出时跳过旧状态。如果追问变成“这是最短路和计数 DP 的结合，更新顺序必须和距离松弛一致”，先判断原来的状态是否仍然覆盖新答案集合，再决定是微调边界、改 value 语义，还是换成另一类结构。

\subsubsection{LeetCode 2045 到达目的地的第二短时间}

先区分步数和代价。本题的模型是每个节点需要保留最短和次短到达步数，再考虑红绿灯等待时间；若暴力枚举路径，路径数量会爆炸。优化点BFS 式扩展边数或时间，只有候选时间不同且能进入前两名时才入队说明下一步扩展必须按某种代价顺序组织，而不是普通队列随便推进。

最短路状态至少包含当前节点和当前代价。本题的BFS 式扩展边数或时间，只有候选时间不同且能进入前两名时才入队决定优先队列按什么排序，以及旧状态如何被识别。非负边条件、代价合成方式和不可达处理都要写清楚。放到“最短路与带权图”这条主线里看，关键原则是：非负权图用 Dijkstra，小顶堆保存当前距离最小的候选；边权为一用 BFS；有负边要考虑 Bellman-Ford 或重新建模。路径代价若是最大边权，也可用 Dijkstra 或二分答案。

最短路证明依赖扩展顺序。若本题使用 Dijkstra，就要说明非负代价让弹出的未过期状态已经最优；若用二分可达性，就要说明阈值越大可用边只会增加。这道题的边界不能留到最后补丁式处理，实验时覆盖第二短不能等于最短、红灯等待、无向图往返、到达终点后继续找次短，打印每次弹出的代价和是否过期。若普通 BFS 与 Dijkstra 结果不同，要能解释边权条件造成的差异。

C++ 实现上，Dijkstra 用邻接表保存 \code{(to, weight)}，距离数组用 \code{long long}。由于 \code{priority_queue} 没有 decrease-key，压入新状态后弹出时跳过旧状态。如果追问变成“它说明最短路状态有时不是一个最优值，而是前两个不同值”，先判断原来的状态是否仍然覆盖新答案集合，再决定是微调边界、改 value 语义，还是换成另一类结构。

\subsection{并查集与动态连通性工作坊}

先用搜索回答连通性，再发现查询和合并交替出现。并查集只维护集合代表，如果题目还要路径、距离、比例或奇偶性，就必须扩展状态。

\subsubsection{LeetCode 305 岛屿数量 II}

先用搜索回答每次连通性查询，观察合并和查询如何交替出现。本题的模型是陆地逐步加入，每次加入只可能影响相邻陆地所在集合；慢点在每次都重新遍历已有连接。并查集只适合新增陆地先让岛屿数加一，再和四邻已存在陆地合并，合并成功则岛屿数减一这类集合代表问题，若题目要路径本身就不能硬套。

并查集状态是代表元和集合附加信息。本题的新增陆地先让岛屿数加一，再和四邻已存在陆地合并，合并成功则岛屿数减一只要求合并集合和查询同组关系；如果还要大小、边数、比例或奇偶性，就要在合并根时同步维护额外数组。放到“并查集与动态连通性”这条主线里看，关键原则是：路径压缩把查找路径上的节点直接挂到根，按大小或按秩合并控制树高。邮箱、等式、边、网格都可以映射成元素集合。

并查集正确性来自等价关系。合并两个代表后，两个集合中任意元素的代表都会一致；不合并时，集合关系保持不变。本题的新增陆地先让岛屿数加一，再和四邻已存在陆地合并，合并成功则岛屿数减一必须只依赖这种代表关系。这道题的边界不能留到最后补丁式处理，实验时覆盖重复加入同一格、边界位置、四邻属于同一岛、空操作列表，打印每个元素的代表元和集合大小。重复合并、已经连通的边和字符串编号最容易暴露实现问题。

C++ 实现上，并查集类中成员访问保持 \code{this->}，\code{find} 可以写成迭代路径压缩避免深递归。若维护集合大小、边数或权值差，合并根时同步更新。如果追问变成“这是动态连通性，比每次重新 BFS 更适合并查集”，先判断原来的状态是否仍然覆盖新答案集合，再决定是微调边界、改 value 语义，还是换成另一类结构。

\subsubsection{LeetCode 399 除法求值}

先用搜索回答每次连通性查询，观察合并和查询如何交替出现。本题的模型是变量之间的除法关系可以看成带权连通分量；慢点在每次都重新遍历已有连接。并查集只适合并查集维护节点到根的乘积权值，合并时根据等式调整根之间权重这类集合代表问题，若题目要路径本身就不能硬套。

并查集状态是代表元和集合附加信息。本题的并查集维护节点到根的乘积权值，合并时根据等式调整根之间权重只要求合并集合和查询同组关系；如果还要大小、边数、比例或奇偶性，就要在合并根时同步维护额外数组。放到“并查集与动态连通性”这条主线里看，关键原则是：路径压缩把查找路径上的节点直接挂到根，按大小或按秩合并控制树高。邮箱、等式、边、网格都可以映射成元素集合。

并查集正确性来自等价关系。合并两个代表后，两个集合中任意元素的代表都会一致；不合并时，集合关系保持不变。本题的并查集维护节点到根的乘积权值，合并时根据等式调整根之间权重必须只依赖这种代表关系。这道题的边界不能留到最后补丁式处理，实验时覆盖查询未知变量、自除、矛盾数据、路径压缩时权值同步更新，打印每个元素的代表元和集合大小。重复合并、已经连通的边和字符串编号最容易暴露实现问题。

C++ 实现上，并查集类中成员访问保持 \code{this->}，\code{find} 可以写成迭代路径压缩避免深递归。若维护集合大小、边数或权值差，合并根时同步更新。如果追问变成“也可建图 DFS，但多次查询时带权并查集更适合”，先判断原来的状态是否仍然覆盖新答案集合，再决定是微调边界、改 value 语义，还是换成另一类结构。

\subsubsection{LeetCode 547 省份数量}

先用搜索回答每次连通性查询，观察合并和查询如何交替出现。本题的模型是城市连接关系形成无向图，省份就是连通分量；慢点在每次都重新遍历已有连接。并查集只适合并查集合并所有直接连接的城市，最后统计根数量这类集合代表问题，若题目要路径本身就不能硬套。

并查集状态是代表元和集合附加信息。本题的并查集合并所有直接连接的城市，最后统计根数量只要求合并集合和查询同组关系；如果还要大小、边数、比例或奇偶性，就要在合并根时同步维护额外数组。放到“并查集与动态连通性”这条主线里看，关键原则是：路径压缩把查找路径上的节点直接挂到根，按大小或按秩合并控制树高。邮箱、等式、边、网格都可以映射成元素集合。

并查集正确性来自等价关系。合并两个代表后，两个集合中任意元素的代表都会一致；不合并时，集合关系保持不变。本题的并查集合并所有直接连接的城市，最后统计根数量必须只依赖这种代表关系。这道题的边界不能留到最后补丁式处理，实验时覆盖邻接矩阵对称、自连接、孤立城市、只扫描上三角，打印每个元素的代表元和集合大小。重复合并、已经连通的边和字符串编号最容易暴露实现问题。

C++ 实现上，并查集类中成员访问保持 \code{this->}，\code{find} 可以写成迭代路径压缩避免深递归。若维护集合大小、边数或权值差，合并根时同步更新。如果追问变成“BFS 也可做；并查集更突出动态合并”，先判断原来的状态是否仍然覆盖新答案集合，再决定是微调边界、改 value 语义，还是换成另一类结构。

\subsubsection{LeetCode 684 冗余连接}

先用搜索回答每次连通性查询，观察合并和查询如何交替出现。本题的模型是树加一条边会形成唯一环；慢点在每次都重新遍历已有连接。并查集只适合按顺序加边，若一条边两端已连通，则它就是冗余边这类集合代表问题，若题目要路径本身就不能硬套。

并查集状态是代表元和集合附加信息。本题的按顺序加边，若一条边两端已连通，则它就是冗余边只要求合并集合和查询同组关系；如果还要大小、边数、比例或奇偶性，就要在合并根时同步维护额外数组。放到“并查集与动态连通性”这条主线里看，关键原则是：路径压缩把查找路径上的节点直接挂到根，按大小或按秩合并控制树高。邮箱、等式、边、网格都可以映射成元素集合。

并查集正确性来自等价关系。合并两个代表后，两个集合中任意元素的代表都会一致；不合并时，集合关系保持不变。本题的按顺序加边，若一条边两端已连通，则它就是冗余边必须只依赖这种代表关系。这道题的边界不能留到最后补丁式处理，实验时覆盖节点编号、最后一条成环边、重复边、并查集初始化大小，打印每个元素的代表元和集合大小。重复合并、已经连通的边和字符串编号最容易暴露实现问题。

C++ 实现上，并查集类中成员访问保持 \code{this->}，\code{find} 可以写成迭代路径压缩避免深递归。若维护集合大小、边数或权值差，合并根时同步更新。如果追问变成“有向冗余连接需要同时处理入度为二和环，复杂度更高”，先判断原来的状态是否仍然覆盖新答案集合，再决定是微调边界、改 value 语义，还是换成另一类结构。

\subsubsection{LeetCode 685 冗余连接 II}

先用搜索回答每次连通性查询，观察合并和查询如何交替出现。本题的模型是有向图中额外边可能导致某节点入度为二，也可能导致环；慢点在每次都重新遍历已有连接。并查集只适合先记录入度为二的两条候选边，再用并查集排除其中一条测试是否成树这类集合代表问题，若题目要路径本身就不能硬套。

并查集状态是代表元和集合附加信息。本题的先记录入度为二的两条候选边，再用并查集排除其中一条测试是否成树只要求合并集合和查询同组关系；如果还要大小、边数、比例或奇偶性，就要在合并根时同步维护额外数组。放到“并查集与动态连通性”这条主线里看，关键原则是：路径压缩把查找路径上的节点直接挂到根，按大小或按秩合并控制树高。邮箱、等式、边、网格都可以映射成元素集合。

并查集正确性来自等价关系。合并两个代表后，两个集合中任意元素的代表都会一致；不合并时，集合关系保持不变。本题的先记录入度为二的两条候选边，再用并查集排除其中一条测试是否成树必须只依赖这种代表关系。这道题的边界不能留到最后补丁式处理，实验时覆盖只有环、只有入度二、两者同时存在、返回最后出现的边，打印每个元素的代表元和集合大小。重复合并、已经连通的边和字符串编号最容易暴露实现问题。

C++ 实现上，并查集类中成员访问保持 \code{this->}，\code{find} 可以写成迭代路径压缩避免深递归。若维护集合大小、边数或权值差，合并根时同步更新。如果追问变成“它比无向冗余连接多了有向树入度约束”，先判断原来的状态是否仍然覆盖新答案集合，再决定是微调边界、改 value 语义，还是换成另一类结构。

\subsubsection{LeetCode 721 账户合并}

先用搜索回答每次连通性查询，观察合并和查询如何交替出现。本题的模型是共享邮箱说明账户属于同一人，邮箱之间形成连通分量；慢点在每次都重新遍历已有连接。并查集只适合给邮箱编号并查集合并，同根邮箱收集后排序输出这类集合代表问题，若题目要路径本身就不能硬套。

并查集状态是代表元和集合附加信息。本题的给邮箱编号并查集合并，同根邮箱收集后排序输出只要求合并集合和查询同组关系；如果还要大小、边数、比例或奇偶性，就要在合并根时同步维护额外数组。放到“并查集与动态连通性”这条主线里看，关键原则是：路径压缩把查找路径上的节点直接挂到根，按大小或按秩合并控制树高。邮箱、等式、边、网格都可以映射成元素集合。

并查集正确性来自等价关系。合并两个代表后，两个集合中任意元素的代表都会一致；不合并时，集合关系保持不变。本题的给邮箱编号并查集合并，同根邮箱收集后排序输出必须只依赖这种代表关系。这道题的边界不能留到最后补丁式处理，实验时覆盖同名不同人、一个账户多个邮箱、重复邮箱、输出排序，打印每个元素的代表元和集合大小。重复合并、已经连通的边和字符串编号最容易暴露实现问题。

C++ 实现上，并查集类中成员访问保持 \code{this->}，\code{find} 可以写成迭代路径压缩避免深递归。若维护集合大小、边数或权值差，合并根时同步更新。如果追问变成“姓名不能作为合并依据，只能作为输出字段”，先判断原来的状态是否仍然覆盖新答案集合，再决定是微调边界、改 value 语义，还是换成另一类结构。

\subsubsection{LeetCode 947 移除最多的同行或同列石头}

先用搜索回答每次连通性查询，观察合并和查询如何交替出现。本题的模型是同一行或同一列的石头属于同一连通分量，每个分量至少留一块；慢点在每次都重新遍历已有连接。并查集只适合把行和列编号后合并，答案是石头数减去连通分量数这类集合代表问题，若题目要路径本身就不能硬套。

并查集状态是代表元和集合附加信息。本题的把行和列编号后合并，答案是石头数减去连通分量数只要求合并集合和查询同组关系；如果还要大小、边数、比例或奇偶性，就要在合并根时同步维护额外数组。放到“并查集与动态连通性”这条主线里看，关键原则是：路径压缩把查找路径上的节点直接挂到根，按大小或按秩合并控制树高。邮箱、等式、边、网格都可以映射成元素集合。

并查集正确性来自等价关系。合并两个代表后，两个集合中任意元素的代表都会一致；不合并时，集合关系保持不变。本题的把行和列编号后合并，答案是石头数减去连通分量数必须只依赖这种代表关系。这道题的边界不能留到最后补丁式处理，实验时覆盖单块石头、行列编号冲突、多个分量、重复坐标，打印每个元素的代表元和集合大小。重复合并、已经连通的边和字符串编号最容易暴露实现问题。

C++ 实现上，并查集类中成员访问保持 \code{this->}，\code{find} 可以写成迭代路径压缩避免深递归。若维护集合大小、边数或权值差，合并根时同步更新。如果追问变成“也可把石头当节点建图，但行列并查集更省边”，先判断原来的状态是否仍然覆盖新答案集合，再决定是微调边界、改 value 语义，还是换成另一类结构。

\subsubsection{LeetCode 959 由斜杠划分区域}

先用搜索回答每次连通性查询，观察合并和查询如何交替出现。本题的模型是每个网格可拆成四个小三角，斜杠决定内部哪些三角相连；慢点在每次都重新遍历已有连接。并查集只适合给每格四个区域编号，先按字符合并内部，再合并相邻格边界区域这类集合代表问题，若题目要路径本身就不能硬套。

并查集状态是代表元和集合附加信息。本题的给每格四个区域编号，先按字符合并内部，再合并相邻格边界区域只要求合并集合和查询同组关系；如果还要大小、边数、比例或奇偶性，就要在合并根时同步维护额外数组。放到“并查集与动态连通性”这条主线里看，关键原则是：路径压缩把查找路径上的节点直接挂到根，按大小或按秩合并控制树高。邮箱、等式、边、网格都可以映射成元素集合。

并查集正确性来自等价关系。合并两个代表后，两个集合中任意元素的代表都会一致；不合并时，集合关系保持不变。本题的给每格四个区域编号，先按字符合并内部，再合并相邻格边界区域必须只依赖这种代表关系。这道题的边界不能留到最后补丁式处理，实验时覆盖空格、正斜杠、反斜杠、边界合并、区域计数，打印每个元素的代表元和集合大小。重复合并、已经连通的边和字符串编号最容易暴露实现问题。

C++ 实现上，并查集类中成员访问保持 \code{this->}，\code{find} 可以写成迭代路径压缩避免深递归。若维护集合大小、边数或权值差，合并根时同步更新。如果追问变成“这是把几何连通性离散成并查集元素的典型题”，先判断原来的状态是否仍然覆盖新答案集合，再决定是微调边界、改 value 语义，还是换成另一类结构。

\subsubsection{LeetCode 990 等式方程的可满足性}

先用搜索回答每次连通性查询，观察合并和查询如何交替出现。本题的模型是相等关系形成集合，不等关系要求两端不在同一集合；慢点在每次都重新遍历已有连接。并查集只适合先合并所有等式，再检查每个不等式是否冲突这类集合代表问题，若题目要路径本身就不能硬套。

并查集状态是代表元和集合附加信息。本题的先合并所有等式，再检查每个不等式是否冲突只要求合并集合和查询同组关系；如果还要大小、边数、比例或奇偶性，就要在合并根时同步维护额外数组。放到“并查集与动态连通性”这条主线里看，关键原则是：路径压缩把查找路径上的节点直接挂到根，按大小或按秩合并控制树高。邮箱、等式、边、网格都可以映射成元素集合。

并查集正确性来自等价关系。合并两个代表后，两个集合中任意元素的代表都会一致；不合并时，集合关系保持不变。本题的先合并所有等式，再检查每个不等式是否冲突必须只依赖这种代表关系。这道题的边界不能留到最后补丁式处理，实验时覆盖自等式、自不等式、传递相等、多组变量，打印每个元素的代表元和集合大小。重复合并、已经连通的边和字符串编号最容易暴露实现问题。

C++ 实现上，并查集类中成员访问保持 \code{this->}，\code{find} 可以写成迭代路径压缩避免深递归。若维护集合大小、边数或权值差，合并根时同步更新。如果追问变成“先处理不等式会漏掉后续等式带来的传递冲突”，先判断原来的状态是否仍然覆盖新答案集合，再决定是微调边界、改 value 语义，还是换成另一类结构。

\subsubsection{LeetCode 1061 按字典序排列最小的等效字符串}

先用搜索回答每次连通性查询，观察合并和查询如何交替出现。本题的模型是等价字符形成集合，每个集合应选择字典序最小代表；慢点在每次都重新遍历已有连接。并查集只适合合并两个字符集合时让较小根成为代表，扫描目标串时替换为根这类集合代表问题，若题目要路径本身就不能硬套。

并查集状态是代表元和集合附加信息。本题的合并两个字符集合时让较小根成为代表，扫描目标串时替换为根只要求合并集合和查询同组关系；如果还要大小、边数、比例或奇偶性，就要在合并根时同步维护额外数组。放到“并查集与动态连通性”这条主线里看，关键原则是：路径压缩把查找路径上的节点直接挂到根，按大小或按秩合并控制树高。邮箱、等式、边、网格都可以映射成元素集合。

并查集正确性来自等价关系。合并两个代表后，两个集合中任意元素的代表都会一致；不合并时，集合关系保持不变。本题的合并两个字符集合时让较小根成为代表，扫描目标串时替换为根必须只依赖这种代表关系。这道题的边界不能留到最后补丁式处理，实验时覆盖传递等价、重复等价关系、自等价、未出现字符，打印每个元素的代表元和集合大小。重复合并、已经连通的边和字符串编号最容易暴露实现问题。

C++ 实现上，并查集类中成员访问保持 \code{this->}，\code{find} 可以写成迭代路径压缩避免深递归。若维护集合大小、边数或权值差，合并根时同步更新。如果追问变成“如果集合代表还要保存其他权值，合并规则需要同步维护”，先判断原来的状态是否仍然覆盖新答案集合，再决定是微调边界、改 value 语义，还是换成另一类结构。

\subsubsection{LeetCode 1202 交换字符串中的元素}

先用搜索回答每次连通性查询，观察合并和查询如何交替出现。本题的模型是可交换下标形成连通分量，同一分量内字符可以任意重排；慢点在每次都重新遍历已有连接。并查集只适合并查集合并所有可交换下标，对每个分量收集字符并按字典序分配回最小下标这类集合代表问题，若题目要路径本身就不能硬套。

并查集状态是代表元和集合附加信息。本题的并查集合并所有可交换下标，对每个分量收集字符并按字典序分配回最小下标只要求合并集合和查询同组关系；如果还要大小、边数、比例或奇偶性，就要在合并根时同步维护额外数组。放到“并查集与动态连通性”这条主线里看，关键原则是：路径压缩把查找路径上的节点直接挂到根，按大小或按秩合并控制树高。邮箱、等式、边、网格都可以映射成元素集合。

并查集正确性来自等价关系。合并两个代表后，两个集合中任意元素的代表都会一致；不合并时，集合关系保持不变。本题的并查集合并所有可交换下标，对每个分量收集字符并按字典序分配回最小下标必须只依赖这种代表关系。这道题的边界不能留到最后补丁式处理，实验时覆盖没有交换、多个分量、重复字符、分量下标排序，打印每个元素的代表元和集合大小。重复合并、已经连通的边和字符串编号最容易暴露实现问题。

C++ 实现上，并查集类中成员访问保持 \code{this->}，\code{find} 可以写成迭代路径压缩避免深递归。若维护集合大小、边数或权值差，合并根时同步更新。如果追问变成“它展示并查集不仅能回答连通性，还能把分量内资源重新组织”，先判断原来的状态是否仍然覆盖新答案集合，再决定是微调边界、改 value 语义，还是换成另一类结构。

\subsubsection{LeetCode 1319 连通网络的操作次数}

先用搜索回答每次连通性查询，观察合并和查询如何交替出现。本题的模型是多余网线可以挪用来连接不同连通分量；慢点在每次都重新遍历已有连接。并查集只适合先判断边数至少为 n 减一，再用并查集统计连通分量，答案是分量数减一这类集合代表问题，若题目要路径本身就不能硬套。

并查集状态是代表元和集合附加信息。本题的先判断边数至少为 n 减一，再用并查集统计连通分量，答案是分量数减一只要求合并集合和查询同组关系；如果还要大小、边数、比例或奇偶性，就要在合并根时同步维护额外数组。放到“并查集与动态连通性”这条主线里看，关键原则是：路径压缩把查找路径上的节点直接挂到根，按大小或按秩合并控制树高。邮箱、等式、边、网格都可以映射成元素集合。

并查集正确性来自等价关系。合并两个代表后，两个集合中任意元素的代表都会一致；不合并时，集合关系保持不变。本题的先判断边数至少为 n 减一，再用并查集统计连通分量，答案是分量数减一必须只依赖这种代表关系。这道题的边界不能留到最后补丁式处理，实验时覆盖边数不足、已经连通、重复边、多分量，打印每个元素的代表元和集合大小。重复合并、已经连通的边和字符串编号最容易暴露实现问题。

C++ 实现上，并查集类中成员访问保持 \code{this->}，\code{find} 可以写成迭代路径压缩避免深递归。若维护集合大小、边数或权值差，合并根时同步更新。如果追问变成“合并失败的边代表冗余资源，但只要总边数足够即可”，先判断原来的状态是否仍然覆盖新答案集合，再决定是微调边界、改 value 语义，还是换成另一类结构。

\subsubsection{LeetCode 1579 保证图可完全遍历}

先用搜索回答每次连通性查询，观察合并和查询如何交替出现。本题的模型是Alice 和 Bob 各自需要图连通，公共边应优先服务两人；慢点在每次都重新遍历已有连接。并查集只适合先处理类型三公共边同步合并两个并查集，再分别处理各自边，最后检查两图连通这类集合代表问题，若题目要路径本身就不能硬套。

并查集状态是代表元和集合附加信息。本题的先处理类型三公共边同步合并两个并查集，再分别处理各自边，最后检查两图连通只要求合并集合和查询同组关系；如果还要大小、边数、比例或奇偶性，就要在合并根时同步维护额外数组。放到“并查集与动态连通性”这条主线里看，关键原则是：路径压缩把查找路径上的节点直接挂到根，按大小或按秩合并控制树高。邮箱、等式、边、网格都可以映射成元素集合。

并查集正确性来自等价关系。合并两个代表后，两个集合中任意元素的代表都会一致；不合并时，集合关系保持不变。本题的先处理类型三公共边同步合并两个并查集，再分别处理各自边，最后检查两图连通必须只依赖这种代表关系。这道题的边界不能留到最后补丁式处理，实验时覆盖公共边冗余、某一方不连通、边数很多、节点从一开始编号，打印每个元素的代表元和集合大小。重复合并、已经连通的边和字符串编号最容易暴露实现问题。

C++ 实现上，并查集类中成员访问保持 \code{this->}，\code{find} 可以写成迭代路径压缩避免深递归。若维护集合大小、边数或权值差，合并根时同步更新。如果追问变成“这是双并查集和资源优先级的组合”，先判断原来的状态是否仍然覆盖新答案集合，再决定是微调边界、改 value 语义，还是换成另一类结构。

\subsection{回溯、枚举与剪枝工作坊}

先写搜索树：路径是什么、选择列表是什么、什么时候结束。然后用排序、剩余容量、同层去重或合法性预处理剪掉不可能分支。

\subsubsection{LeetCode 17 电话号码的字母组合}

先画前三层搜索树。本题的模型是每个数字对应若干字母，答案是所有按位选择形成的字符串；暴力搜索要明确路径、选择列表和终止条件。优化点是暴力就是完整笛卡尔积，优化重点是路径长度和下标同步推进，避免在每层重新拼接大量中间字符串，也就是哪些分支在展开前已经不可能产生合法答案。

回溯状态由路径和当前位置共同定义。本题的暴力就是完整笛卡尔积，优化重点是路径长度和下标同步推进，避免在每层重新拼接大量中间字符串要落成剪枝条件或同层去重条件。剪枝只能删掉不可能成功的分支，不能删掉只是暂时看起来不优的分支。放到“回溯、枚举与剪枝”这条主线里看，关键原则是：剪枝来自容量、剩余长度、排序后去重、合法性预检查和提前终止。组合题关注起点推进，排列题关注元素是否使用，分割题关注当前片段是否合法。

回溯证明分两半：搜索树覆盖所有合法答案，同层去重或剪枝不会删除合法答案。本题的每个数字对应若干字母，答案是所有按位选择形成的字符串给出选择来源，暴力就是完整笛卡尔积，优化重点是路径长度和下标同步推进，避免在每层重新拼接大量中间字符串给出提前停止的理由。这道题的边界不能留到最后补丁式处理，实验时覆盖空输入、包含 7 或 9、输入中不应出现 0 和 1、输出规模指数增长，统计搜索节点数量和剪枝次数。重复元素题要打印同一层跳过哪些值，确认没有误删合法路径。

C++ 实现上，路径容器修改要成对出现：加入、搜索、撤销。重复元素题先排序，再用同层去重条件；输出规模大时，复杂度要把答案数量算进去。如果追问变成“若字符映射来自键盘以外的字典，只需要替换候选表，搜索骨架不变”，先判断原来的状态是否仍然覆盖新答案集合，再决定是微调边界、改 value 语义，还是换成另一类结构。

\subsubsection{LeetCode 22 括号生成}

先画前三层搜索树。本题的模型是合法括号串的前缀中右括号数量不能超过左括号数量；暴力搜索要明确路径、选择列表和终止条件。优化点是状态保存已用左括号和右括号数量，只有满足约束的选择才继续递归，也就是哪些分支在展开前已经不可能产生合法答案。

回溯状态由路径和当前位置共同定义。本题的状态保存已用左括号和右括号数量，只有满足约束的选择才继续递归要落成剪枝条件或同层去重条件。剪枝只能删掉不可能成功的分支，不能删掉只是暂时看起来不优的分支。放到“回溯、枚举与剪枝”这条主线里看，关键原则是：剪枝来自容量、剩余长度、排序后去重、合法性预检查和提前终止。组合题关注起点推进，排列题关注元素是否使用，分割题关注当前片段是否合法。

回溯证明分两半：搜索树覆盖所有合法答案，同层去重或剪枝不会删除合法答案。本题的合法括号串的前缀中右括号数量不能超过左括号数量给出选择来源，状态保存已用左括号和右括号数量，只有满足约束的选择才继续递归给出提前停止的理由。这道题的边界不能留到最后补丁式处理，实验时覆盖n 为零或一、右括号提前超过左括号、输出规模是卡特兰数，统计搜索节点数量和剪枝次数。重复元素题要打印同一层跳过哪些值，确认没有误删合法路径。

C++ 实现上，路径容器修改要成对出现：加入、搜索、撤销。重复元素题先排序，再用同层去重条件；输出规模大时，复杂度要把答案数量算进去。如果追问变成“若有多种括号类型，状态还需要栈保存类型匹配关系”，先判断原来的状态是否仍然覆盖新答案集合，再决定是微调边界、改 value 语义，还是换成另一类结构。

\subsubsection{LeetCode 39 组合总和}

先画前三层搜索树。本题的模型是候选数字可重复使用，答案不关心排列顺序；暴力搜索要明确路径、选择列表和终止条件。优化点是排序后按起点枚举，递归时仍传当前下标表示可以继续使用同一数字，也就是哪些分支在展开前已经不可能产生合法答案。

回溯状态由路径和当前位置共同定义。本题的排序后按起点枚举，递归时仍传当前下标表示可以继续使用同一数字要落成剪枝条件或同层去重条件。剪枝只能删掉不可能成功的分支，不能删掉只是暂时看起来不优的分支。放到“回溯、枚举与剪枝”这条主线里看，关键原则是：剪枝来自容量、剩余长度、排序后去重、合法性预检查和提前终止。组合题关注起点推进，排列题关注元素是否使用，分割题关注当前片段是否合法。

回溯证明分两半：搜索树覆盖所有合法答案，同层去重或剪枝不会删除合法答案。本题的候选数字可重复使用，答案不关心排列顺序给出选择来源，排序后按起点枚举，递归时仍传当前下标表示可以继续使用同一数字给出提前停止的理由。这道题的边界不能留到最后补丁式处理，实验时覆盖目标小于最小数、刚好等于、候选有重复时题意如何处理、输出规模，统计搜索节点数量和剪枝次数。重复元素题要打印同一层跳过哪些值，确认没有误删合法路径。

C++ 实现上，路径容器修改要成对出现：加入、搜索、撤销。重复元素题先排序，再用同层去重条件；输出规模大时，复杂度要把答案数量算进去。如果追问变成“组合总和 II 每个数只能用一次，起点推进和去重条件都不同”，先判断原来的状态是否仍然覆盖新答案集合，再决定是微调边界、改 value 语义，还是换成另一类结构。

\subsubsection{LeetCode 40 组合总和 II}

先画前三层搜索树。本题的模型是每个候选只能使用一次，输入中可能存在重复值；暴力搜索要明确路径、选择列表和终止条件。优化点是排序后递归下一层从 i 加一开始，同层遇到相同值要跳过，也就是哪些分支在展开前已经不可能产生合法答案。

回溯状态由路径和当前位置共同定义。本题的排序后递归下一层从 i 加一开始，同层遇到相同值要跳过要落成剪枝条件或同层去重条件。剪枝只能删掉不可能成功的分支，不能删掉只是暂时看起来不优的分支。放到“回溯、枚举与剪枝”这条主线里看，关键原则是：剪枝来自容量、剩余长度、排序后去重、合法性预检查和提前终止。组合题关注起点推进，排列题关注元素是否使用，分割题关注当前片段是否合法。

回溯证明分两半：搜索树覆盖所有合法答案，同层去重或剪枝不会删除合法答案。本题的每个候选只能使用一次，输入中可能存在重复值给出选择来源，排序后递归下一层从 i 加一开始，同层遇到相同值要跳过给出提前停止的理由。这道题的边界不能留到最后补丁式处理，实验时覆盖全重复、目标为零、同层去重写成跨层去重、答案顺序不要求，统计搜索节点数量和剪枝次数。重复元素题要打印同一层跳过哪些值，确认没有误删合法路径。

C++ 实现上，路径容器修改要成对出现：加入、搜索、撤销。重复元素题先排序，再用同层去重条件；输出规模大时，复杂度要把答案数量算进去。如果追问变成“子集 II 使用同样的同层去重，只是每个路径节点都可能成为答案”，先判断原来的状态是否仍然覆盖新答案集合，再决定是微调边界、改 value 语义，还是换成另一类结构。

\subsubsection{LeetCode 46 全排列}

先画前三层搜索树。本题的模型是排列关心元素顺序，每个位置从所有未使用元素中选择一个；暴力搜索要明确路径、选择列表和终止条件。优化点是路径长度达到 n 时产生答案，used 数组保证同一元素不重复使用，也就是哪些分支在展开前已经不可能产生合法答案。

回溯状态由路径和当前位置共同定义。本题的路径长度达到 n 时产生答案，used 数组保证同一元素不重复使用要落成剪枝条件或同层去重条件。剪枝只能删掉不可能成功的分支，不能删掉只是暂时看起来不优的分支。放到“回溯、枚举与剪枝”这条主线里看，关键原则是：剪枝来自容量、剩余长度、排序后去重、合法性预检查和提前终止。组合题关注起点推进，排列题关注元素是否使用，分割题关注当前片段是否合法。

回溯证明分两半：搜索树覆盖所有合法答案，同层去重或剪枝不会删除合法答案。本题的排列关心元素顺序，每个位置从所有未使用元素中选择一个给出选择来源，路径长度达到 n 时产生答案，used 数组保证同一元素不重复使用给出提前停止的理由。这道题的边界不能留到最后补丁式处理，实验时覆盖空数组、单元素、输出规模为 n 阶乘、元素值唯一假设，统计搜索节点数量和剪枝次数。重复元素题要打印同一层跳过哪些值，确认没有误删合法路径。

C++ 实现上，路径容器修改要成对出现：加入、搜索、撤销。重复元素题先排序，再用同层去重条件；输出规模大时，复杂度要把答案数量算进去。如果追问变成“若有重复元素，需要排序并在同层跳过等价选择”，先判断原来的状态是否仍然覆盖新答案集合，再决定是微调边界、改 value 语义，还是换成另一类结构。

\subsubsection{LeetCode 47 全排列 II}

先画前三层搜索树。本题的模型是重复元素会让不同下标选择生成相同排列；暴力搜索要明确路径、选择列表和终止条件。优化点是排序后使用 used 数组，同层遇到相同值且前一个相同值未使用时跳过，也就是哪些分支在展开前已经不可能产生合法答案。

回溯状态由路径和当前位置共同定义。本题的排序后使用 used 数组，同层遇到相同值且前一个相同值未使用时跳过要落成剪枝条件或同层去重条件。剪枝只能删掉不可能成功的分支，不能删掉只是暂时看起来不优的分支。放到“回溯、枚举与剪枝”这条主线里看，关键原则是：剪枝来自容量、剩余长度、排序后去重、合法性预检查和提前终止。组合题关注起点推进，排列题关注元素是否使用，分割题关注当前片段是否合法。

回溯证明分两半：搜索树覆盖所有合法答案，同层去重或剪枝不会删除合法答案。本题的重复元素会让不同下标选择生成相同排列给出选择来源，排序后使用 used 数组，同层遇到相同值且前一个相同值未使用时跳过给出提前停止的理由。这道题的边界不能留到最后补丁式处理，实验时覆盖全重复、部分重复、同层去重条件写反、输出顺序不要求，统计搜索节点数量和剪枝次数。重复元素题要打印同一层跳过哪些值，确认没有误删合法路径。

C++ 实现上，路径容器修改要成对出现：加入、搜索、撤销。重复元素题先排序，再用同层去重条件；输出规模大时，复杂度要把答案数量算进去。如果追问变成“子集 II 和组合总和 II 也是同层去重，但起点推进方式不同”，先判断原来的状态是否仍然覆盖新答案集合，再决定是微调边界、改 value 语义，还是换成另一类结构。

\subsubsection{LeetCode 51 N 皇后}

先画前三层搜索树。本题的模型是每一行放一个皇后，列和两条对角线不能冲突；暴力搜索要明确路径、选择列表和终止条件。优化点是逐行搜索，用列集合、主对角线集合和副对角线集合快速判断合法性，也就是哪些分支在展开前已经不可能产生合法答案。

回溯状态由路径和当前位置共同定义。本题的逐行搜索，用列集合、主对角线集合和副对角线集合快速判断合法性要落成剪枝条件或同层去重条件。剪枝只能删掉不可能成功的分支，不能删掉只是暂时看起来不优的分支。放到“回溯、枚举与剪枝”这条主线里看，关键原则是：剪枝来自容量、剩余长度、排序后去重、合法性预检查和提前终止。组合题关注起点推进，排列题关注元素是否使用，分割题关注当前片段是否合法。

回溯证明分两半：搜索树覆盖所有合法答案，同层去重或剪枝不会删除合法答案。本题的每一行放一个皇后，列和两条对角线不能冲突给出选择来源，逐行搜索，用列集合、主对角线集合和副对角线集合快速判断合法性给出提前停止的理由。这道题的边界不能留到最后补丁式处理，实验时覆盖n 为一、n 为二或三无解、对角线编号、回退时状态恢复，统计搜索节点数量和剪枝次数。重复元素题要打印同一层跳过哪些值，确认没有误删合法路径。

C++ 实现上，路径容器修改要成对出现：加入、搜索、撤销。重复元素题先排序，再用同层去重条件；输出规模大时，复杂度要把答案数量算进去。如果追问变成“可用位掩码压缩三个约束集合，提高常数性能”，先判断原来的状态是否仍然覆盖新答案集合，再决定是微调边界、改 value 语义，还是换成另一类结构。

\subsubsection{LeetCode 78 子集}

先画前三层搜索树。本题的模型是每个元素都有选或不选两种选择，答案是所有路径节点；暴力搜索要明确路径、选择列表和终止条件。优化点是回溯按起点向后枚举，避免不同顺序生成同一集合，也就是哪些分支在展开前已经不可能产生合法答案。

回溯状态由路径和当前位置共同定义。本题的回溯按起点向后枚举，避免不同顺序生成同一集合要落成剪枝条件或同层去重条件。剪枝只能删掉不可能成功的分支，不能删掉只是暂时看起来不优的分支。放到“回溯、枚举与剪枝”这条主线里看，关键原则是：剪枝来自容量、剩余长度、排序后去重、合法性预检查和提前终止。组合题关注起点推进，排列题关注元素是否使用，分割题关注当前片段是否合法。

回溯证明分两半：搜索树覆盖所有合法答案，同层去重或剪枝不会删除合法答案。本题的每个元素都有选或不选两种选择，答案是所有路径节点给出选择来源，回溯按起点向后枚举，避免不同顺序生成同一集合给出提前停止的理由。这道题的边界不能留到最后补丁式处理，实验时覆盖空数组、单元素、输出数量为二的 n 次方、顺序不要求，统计搜索节点数量和剪枝次数。重复元素题要打印同一层跳过哪些值，确认没有误删合法路径。

C++ 实现上，路径容器修改要成对出现：加入、搜索、撤销。重复元素题先排序，再用同层去重条件；输出规模大时，复杂度要把答案数量算进去。如果追问变成“也可用位掩码枚举所有子集，适合 n 较小的场景”，先判断原来的状态是否仍然覆盖新答案集合，再决定是微调边界、改 value 语义，还是换成另一类结构。

\subsubsection{LeetCode 79 单词搜索}

先画前三层搜索树。本题的模型是网格路径不能重复使用同一格，状态包含位置、匹配下标和访问标记；暴力搜索要明确路径、选择列表和终止条件。优化点是剪枝包括首字符不匹配、剩余长度不足、字符频次不足和越界访问，也就是哪些分支在展开前已经不可能产生合法答案。

回溯状态由路径和当前位置共同定义。本题的剪枝包括首字符不匹配、剩余长度不足、字符频次不足和越界访问要落成剪枝条件或同层去重条件。剪枝只能删掉不可能成功的分支，不能删掉只是暂时看起来不优的分支。放到“回溯、枚举与剪枝”这条主线里看，关键原则是：剪枝来自容量、剩余长度、排序后去重、合法性预检查和提前终止。组合题关注起点推进，排列题关注元素是否使用，分割题关注当前片段是否合法。

回溯证明分两半：搜索树覆盖所有合法答案，同层去重或剪枝不会删除合法答案。本题的网格路径不能重复使用同一格，状态包含位置、匹配下标和访问标记给出选择来源，剪枝包括首字符不匹配、剩余长度不足、字符频次不足和越界访问给出提前停止的理由。这道题的边界不能留到最后补丁式处理，实验时覆盖单字符单词、路径需要回退、重复字母、单元格不能复用，统计搜索节点数量和剪枝次数。重复元素题要打印同一层跳过哪些值，确认没有误删合法路径。

C++ 实现上，路径容器修改要成对出现：加入、搜索、撤销。重复元素题先排序，再用同层去重条件；输出规模大时，复杂度要把答案数量算进去。如果追问变成“若要查很多单词，应使用 Trie 合并搜索前缀”，先判断原来的状态是否仍然覆盖新答案集合，再决定是微调边界、改 value 语义，还是换成另一类结构。

\subsubsection{LeetCode 90 子集 II}

先画前三层搜索树。本题的模型是重复元素会让不同下标路径生成相同集合；暴力搜索要明确路径、选择列表和终止条件。优化点是排序后在同一层跳过相同值，保留不同层使用重复值的可能，也就是哪些分支在展开前已经不可能产生合法答案。

回溯状态由路径和当前位置共同定义。本题的排序后在同一层跳过相同值，保留不同层使用重复值的可能要落成剪枝条件或同层去重条件。剪枝只能删掉不可能成功的分支，不能删掉只是暂时看起来不优的分支。放到“回溯、枚举与剪枝”这条主线里看，关键原则是：剪枝来自容量、剩余长度、排序后去重、合法性预检查和提前终止。组合题关注起点推进，排列题关注元素是否使用，分割题关注当前片段是否合法。

回溯证明分两半：搜索树覆盖所有合法答案，同层去重或剪枝不会删除合法答案。本题的重复元素会让不同下标路径生成相同集合给出选择来源，排序后在同一层跳过相同值，保留不同层使用重复值的可能给出提前停止的理由。这道题的边界不能留到最后补丁式处理，实验时覆盖全重复、空集、重复值相邻、去重条件写成同层而不是全局，统计搜索节点数量和剪枝次数。重复元素题要打印同一层跳过哪些值，确认没有误删合法路径。

C++ 实现上，路径容器修改要成对出现：加入、搜索、撤销。重复元素题先排序，再用同层去重条件；输出规模大时，复杂度要把答案数量算进去。如果追问变成“组合总和 II 使用同样的同层去重思想”，先判断原来的状态是否仍然覆盖新答案集合，再决定是微调边界、改 value 语义，还是换成另一类结构。

\subsubsection{LeetCode 93 复原 IP 地址}

先画前三层搜索树。本题的模型是字符串要切成四段，每段必须是 0 到 255 且无非法前导零；暴力搜索要明确路径、选择列表和终止条件。优化点是按段数和剩余长度剪枝，每次尝试长度一到三的片段，也就是哪些分支在展开前已经不可能产生合法答案。

回溯状态由路径和当前位置共同定义。本题的按段数和剩余长度剪枝，每次尝试长度一到三的片段要落成剪枝条件或同层去重条件。剪枝只能删掉不可能成功的分支，不能删掉只是暂时看起来不优的分支。放到“回溯、枚举与剪枝”这条主线里看，关键原则是：剪枝来自容量、剩余长度、排序后去重、合法性预检查和提前终止。组合题关注起点推进，排列题关注元素是否使用，分割题关注当前片段是否合法。

回溯证明分两半：搜索树覆盖所有合法答案，同层去重或剪枝不会删除合法答案。本题的字符串要切成四段，每段必须是 0 到 255 且无非法前导零给出选择来源，按段数和剩余长度剪枝，每次尝试长度一到三的片段给出提前停止的理由。这道题的边界不能留到最后补丁式处理，实验时覆盖前导零、数值超过 255、剩余字符太多或太少、刚好四段，统计搜索节点数量和剪枝次数。重复元素题要打印同一层跳过哪些值，确认没有误删合法路径。

C++ 实现上，路径容器修改要成对出现：加入、搜索、撤销。重复元素题先排序，再用同层去重条件；输出规模大时，复杂度要把答案数量算进去。如果追问变成“分割回文串也是前缀分割，但合法性检查不同”，先判断原来的状态是否仍然覆盖新答案集合，再决定是微调边界、改 value 语义，还是换成另一类结构。

\subsubsection{LeetCode 131 分割回文串}

先画前三层搜索树。本题的模型是每次选择一个从当前位置开始的回文前缀；暴力搜索要明确路径、选择列表和终止条件。优化点是预处理回文 DP 表，让回溯合法性检查为常数，也就是哪些分支在展开前已经不可能产生合法答案。

回溯状态由路径和当前位置共同定义。本题的预处理回文 DP 表，让回溯合法性检查为常数要落成剪枝条件或同层去重条件。剪枝只能删掉不可能成功的分支，不能删掉只是暂时看起来不优的分支。放到“回溯、枚举与剪枝”这条主线里看，关键原则是：剪枝来自容量、剩余长度、排序后去重、合法性预检查和提前终止。组合题关注起点推进，排列题关注元素是否使用，分割题关注当前片段是否合法。

回溯证明分两半：搜索树覆盖所有合法答案，同层去重或剪枝不会删除合法答案。本题的每次选择一个从当前位置开始的回文前缀给出选择来源，预处理回文 DP 表，让回溯合法性检查为常数给出提前停止的理由。这道题的边界不能留到最后补丁式处理，实验时覆盖空串、单字符、全相同字符、重复分割路径，统计搜索节点数量和剪枝次数。重复元素题要打印同一层跳过哪些值，确认没有误删合法路径。

C++ 实现上，路径容器修改要成对出现：加入、搜索、撤销。重复元素题先排序，再用同层去重条件；输出规模大时，复杂度要把答案数量算进去。如果追问变成“若只求最少切割次数，可改成 DP 最短切分”，先判断原来的状态是否仍然覆盖新答案集合，再决定是微调边界、改 value 语义，还是换成另一类结构。

\subsubsection{LeetCode 212 单词搜索 II}

先画前三层搜索树。本题的模型是多个单词在同一棋盘上搜索，公共前缀可以共享；暴力搜索要明确路径、选择列表和终止条件。优化点是先建 Trie，再从每个格子回溯；路径到达单词结束节点就收集答案，也就是哪些分支在展开前已经不可能产生合法答案。

回溯状态由路径和当前位置共同定义。本题的先建 Trie，再从每个格子回溯；路径到达单词结束节点就收集答案要落成剪枝条件或同层去重条件。剪枝只能删掉不可能成功的分支，不能删掉只是暂时看起来不优的分支。放到“回溯、枚举与剪枝”这条主线里看，关键原则是：剪枝来自容量、剩余长度、排序后去重、合法性预检查和提前终止。组合题关注起点推进，排列题关注元素是否使用，分割题关注当前片段是否合法。

回溯证明分两半：搜索树覆盖所有合法答案，同层去重或剪枝不会删除合法答案。本题的多个单词在同一棋盘上搜索，公共前缀可以共享给出选择来源，先建 Trie，再从每个格子回溯；路径到达单词结束节点就收集答案给出提前停止的理由。这道题的边界不能留到最后补丁式处理，实验时覆盖重复单词、同一格不能复用、前缀是完整单词、剪枝删除空子树，统计搜索节点数量和剪枝次数。重复元素题要打印同一层跳过哪些值，确认没有误删合法路径。

C++ 实现上，路径容器修改要成对出现：加入、搜索、撤销。重复元素题先排序，再用同层去重条件；输出规模大时，复杂度要把答案数量算进去。如果追问变成“这是单词搜索从单目标扩展到多目标后的结构升级”，先判断原来的状态是否仍然覆盖新答案集合，再决定是微调边界、改 value 语义，还是换成另一类结构。

\subsubsection{LeetCode 216 组合总和 III}

先画前三层搜索树。本题的模型是从一到九中选 k 个数和为 n，每个数最多用一次；暴力搜索要明确路径、选择列表和终止条件。优化点是按递增起点枚举，利用剩余个数和剩余和剪枝，也就是哪些分支在展开前已经不可能产生合法答案。

回溯状态由路径和当前位置共同定义。本题的按递增起点枚举，利用剩余个数和剩余和剪枝要落成剪枝条件或同层去重条件。剪枝只能删掉不可能成功的分支，不能删掉只是暂时看起来不优的分支。放到“回溯、枚举与剪枝”这条主线里看，关键原则是：剪枝来自容量、剩余长度、排序后去重、合法性预检查和提前终止。组合题关注起点推进，排列题关注元素是否使用，分割题关注当前片段是否合法。

回溯证明分两半：搜索树覆盖所有合法答案，同层去重或剪枝不会删除合法答案。本题的从一到九中选 k 个数和为 n，每个数最多用一次给出选择来源，按递增起点枚举，利用剩余个数和剩余和剪枝给出提前停止的理由。这道题的边界不能留到最后补丁式处理，实验时覆盖n 太小、n 太大、k 为零、路径长度已经超过 k，统计搜索节点数量和剪枝次数。重复元素题要打印同一层跳过哪些值，确认没有误删合法路径。

C++ 实现上，路径容器修改要成对出现：加入、搜索、撤销。重复元素题先排序，再用同层去重条件；输出规模大时，复杂度要把答案数量算进去。如果追问变成“可预估剩余最小和最大和进一步剪枝”，先判断原来的状态是否仍然覆盖新答案集合，再决定是微调边界、改 value 语义，还是换成另一类结构。

\subsubsection{LeetCode 698 划分为 k 个相等的子集}

先画前三层搜索树。本题的模型是每个元素要分配到 k 个桶中，所有桶目标和相同；暴力搜索要明确路径、选择列表和终止条件。优化点是先判断总和可整除并排序降序，再用桶剩余容量回溯放置元素，也就是哪些分支在展开前已经不可能产生合法答案。

回溯状态由路径和当前位置共同定义。本题的先判断总和可整除并排序降序，再用桶剩余容量回溯放置元素要落成剪枝条件或同层去重条件。剪枝只能删掉不可能成功的分支，不能删掉只是暂时看起来不优的分支。放到“回溯、枚举与剪枝”这条主线里看，关键原则是：剪枝来自容量、剩余长度、排序后去重、合法性预检查和提前终止。组合题关注起点推进，排列题关注元素是否使用，分割题关注当前片段是否合法。

回溯证明分两半：搜索树覆盖所有合法答案，同层去重或剪枝不会删除合法答案。本题的每个元素要分配到 k 个桶中，所有桶目标和相同给出选择来源，先判断总和可整除并排序降序，再用桶剩余容量回溯放置元素给出提前停止的理由。这道题的边界不能留到最后补丁式处理，实验时覆盖存在元素大于目标、重复元素、空桶对称、k 为一，统计搜索节点数量和剪枝次数。重复元素题要打印同一层跳过哪些值，确认没有误删合法路径。

C++ 实现上，路径容器修改要成对出现：加入、搜索、撤销。重复元素题先排序，再用同层去重条件；输出规模大时，复杂度要把答案数量算进去。如果追问变成“状态压缩 DP 也能做，适合用掩码表示已选元素集合”，先判断原来的状态是否仍然覆盖新答案集合，再决定是微调边界、改 value 语义，还是换成另一类结构。

\subsection{动态规划与状态定义工作坊}

先写递归搜索并找重复状态，再给状态命名。状态必须包含未来决策所需信息，遍历顺序必须保证依赖状态已经算完。

\subsubsection{LeetCode 53 最大子数组和}

先写递归或枚举，把重复子问题暴露出来。本题的模型是状态是必须以当前位置结尾的最大连续和；慢点不是递归形式，而是相同状态被反复求。本题需要命名的状态来自若前一个后缀为负，接上它只会拖累当前元素，因此可以重新开始，状态定义必须让未来决策不再依赖完整历史。

DP 状态必须足够描述未来。本题的若前一个后缀为负，接上它只会拖累当前元素，因此可以重新开始要写成数组下标、容量、前缀长度、持有状态或区间边界。初始化对应空状态，遍历顺序对应依赖方向。放到“动态规划与状态定义”这条主线里看，关键原则是：状态定义要包含足够信息但不能过度膨胀。线性 DP 看最后一步选择，二维 DP 看两个前缀或网格位置，背包看物品阶段和容量，区间 DP 看最后合并点。

DP 证明通常看最后一步。任意最优解或合法方案的最后一步必在转移枚举中；每个转移来源又已经按遍历顺序算完。本题的若前一个后缀为负，接上它只会拖累当前元素，因此可以重新开始要支撑这个依赖关系。这道题的边界不能留到最后补丁式处理，实验时覆盖全负数组、单元素、和溢出、答案不能初始化为零，先打印完整 DP 表，再比较空间压缩版本。若压缩版不同，优先检查遍历顺序和旧值保存。

C++ 实现上，DP 表初始化要对应状态语义，不可达哨兵不能溢出。空间压缩前先写完整表；压缩后要说清读到的是上一轮、本轮左侧还是左上角旧值。如果追问变成“若要求返回区间下标，同步记录重新开始位置和最优左右端”，先判断原来的状态是否仍然覆盖新答案集合，再决定是微调边界、改 value 语义，还是换成另一类结构。

\subsubsection{LeetCode 62 不同路径}

先写递归或枚举，把重复子问题暴露出来。本题的模型是网格状态由上方和左方两种最后一步转移而来；慢点不是递归形式，而是相同状态被反复求。本题需要命名的状态来自二维表可压缩为一维，因为当前行只依赖上一行同列和当前行左侧，状态定义必须让未来决策不再依赖完整历史。

DP 状态必须足够描述未来。本题的二维表可压缩为一维，因为当前行只依赖上一行同列和当前行左侧要写成数组下标、容量、前缀长度、持有状态或区间边界。初始化对应空状态，遍历顺序对应依赖方向。放到“动态规划与状态定义”这条主线里看，关键原则是：状态定义要包含足够信息但不能过度膨胀。线性 DP 看最后一步选择，二维 DP 看两个前缀或网格位置，背包看物品阶段和容量，区间 DP 看最后合并点。

DP 证明通常看最后一步。任意最优解或合法方案的最后一步必在转移枚举中；每个转移来源又已经按遍历顺序算完。本题的二维表可压缩为一维，因为当前行只依赖上一行同列和当前行左侧要支撑这个依赖关系。这道题的边界不能留到最后补丁式处理，实验时覆盖一行、一列、较大网格结果溢出、障碍物版本边界，先打印完整 DP 表，再比较空间压缩版本。若压缩版不同，优先检查遍历顺序和旧值保存。

C++ 实现上，DP 表初始化要对应状态语义，不可达哨兵不能溢出。空间压缩前先写完整表；压缩后要说清读到的是上一轮、本轮左侧还是左上角旧值。如果追问变成“带障碍物时遇到障碍状态清零，第一行第一列要按阻断处理”，先判断原来的状态是否仍然覆盖新答案集合，再决定是微调边界、改 value 语义，还是换成另一类结构。

\subsubsection{LeetCode 63 不同路径 II}

先写递归或枚举，把重复子问题暴露出来。本题的模型是障碍物让某些格子的路径数变为零；慢点不是递归形式，而是相同状态被反复求。本题需要命名的状态来自滚动数组中遇到障碍直接把当前位置清零，否则累加左侧，状态定义必须让未来决策不再依赖完整历史。

DP 状态必须足够描述未来。本题的滚动数组中遇到障碍直接把当前位置清零，否则累加左侧要写成数组下标、容量、前缀长度、持有状态或区间边界。初始化对应空状态，遍历顺序对应依赖方向。放到“动态规划与状态定义”这条主线里看，关键原则是：状态定义要包含足够信息但不能过度膨胀。线性 DP 看最后一步选择，二维 DP 看两个前缀或网格位置，背包看物品阶段和容量，区间 DP 看最后合并点。

DP 证明通常看最后一步。任意最优解或合法方案的最后一步必在转移枚举中；每个转移来源又已经按遍历顺序算完。本题的滚动数组中遇到障碍直接把当前位置清零，否则累加左侧要支撑这个依赖关系。这道题的边界不能留到最后补丁式处理，实验时覆盖起点或终点是障碍、整行被截断、单行障碍，先打印完整 DP 表，再比较空间压缩版本。若压缩版不同，优先检查遍历顺序和旧值保存。

C++ 实现上，DP 表初始化要对应状态语义，不可达哨兵不能溢出。空间压缩前先写完整表；压缩后要说清读到的是上一轮、本轮左侧还是左上角旧值。如果追问变成“若允许向上或向左走，图会有环，普通填表不再成立”，先判断原来的状态是否仍然覆盖新答案集合，再决定是微调边界、改 value 语义，还是换成另一类结构。

\subsubsection{LeetCode 64 最小路径和}

先写递归或枚举，把重复子问题暴露出来。本题的模型是状态是到达当前格子的最小代价，最后一步来自上方或左方；慢点不是递归形式，而是相同状态被反复求。本题需要命名的状态来自先初始化第一行和第一列，再处理内部，主转移保持简单，状态定义必须让未来决策不再依赖完整历史。

DP 状态必须足够描述未来。本题的先初始化第一行和第一列，再处理内部，主转移保持简单要写成数组下标、容量、前缀长度、持有状态或区间边界。初始化对应空状态，遍历顺序对应依赖方向。放到“动态规划与状态定义”这条主线里看，关键原则是：状态定义要包含足够信息但不能过度膨胀。线性 DP 看最后一步选择，二维 DP 看两个前缀或网格位置，背包看物品阶段和容量，区间 DP 看最后合并点。

DP 证明通常看最后一步。任意最优解或合法方案的最后一步必在转移枚举中；每个转移来源又已经按遍历顺序算完。本题的先初始化第一行和第一列，再处理内部，主转移保持简单要支撑这个依赖关系。这道题的边界不能留到最后补丁式处理，实验时覆盖单行、单列、权值为零、路径和溢出，先打印完整 DP 表，再比较空间压缩版本。若压缩版不同，优先检查遍历顺序和旧值保存。

C++ 实现上，DP 表初始化要对应状态语义，不可达哨兵不能溢出。空间压缩前先写完整表；压缩后要说清读到的是上一轮、本轮左侧还是左上角旧值。如果追问变成“若允许四方向移动，变成最短路而不是普通网格 DP”，先判断原来的状态是否仍然覆盖新答案集合，再决定是微调边界、改 value 语义，还是换成另一类结构。

\subsubsection{LeetCode 70 爬楼梯}

先写递归或枚举，把重复子问题暴露出来。本题的模型是最后一步来自前一阶或前两阶，是最小的重叠子问题模型；慢点不是递归形式，而是相同状态被反复求。本题需要命名的状态来自滚动变量保存最近两个状态即可，不需要完整数组，状态定义必须让未来决策不再依赖完整历史。

DP 状态必须足够描述未来。本题的滚动变量保存最近两个状态即可，不需要完整数组要写成数组下标、容量、前缀长度、持有状态或区间边界。初始化对应空状态，遍历顺序对应依赖方向。放到“动态规划与状态定义”这条主线里看，关键原则是：状态定义要包含足够信息但不能过度膨胀。线性 DP 看最后一步选择，二维 DP 看两个前缀或网格位置，背包看物品阶段和容量，区间 DP 看最后合并点。

DP 证明通常看最后一步。任意最优解或合法方案的最后一步必在转移枚举中；每个转移来源又已经按遍历顺序算完。本题的滚动变量保存最近两个状态即可，不需要完整数组要支撑这个依赖关系。这道题的边界不能留到最后补丁式处理，实验时覆盖n 为一、n 为二、题目是否允许零阶、结果范围，先打印完整 DP 表，再比较空间压缩版本。若压缩版不同，优先检查遍历顺序和旧值保存。

C++ 实现上，DP 表初始化要对应状态语义，不可达哨兵不能溢出。空间压缩前先写完整表；压缩后要说清读到的是上一轮、本轮左侧还是左上角旧值。如果追问变成“若每次可走一到 k 阶，转移变成固定宽度窗口求和”，先判断原来的状态是否仍然覆盖新答案集合，再决定是微调边界、改 value 语义，还是换成另一类结构。

\subsubsection{LeetCode 72 编辑距离}

先写递归或枚举，把重复子问题暴露出来。本题的模型是二维状态表示两个前缀之间的最少编辑操作；慢点不是递归形式，而是相同状态被反复求。本题需要命名的状态来自末尾字符相同继承左上；不同则枚举插入、删除、替换，状态定义必须让未来决策不再依赖完整历史。

DP 状态必须足够描述未来。本题的末尾字符相同继承左上；不同则枚举插入、删除、替换要写成数组下标、容量、前缀长度、持有状态或区间边界。初始化对应空状态，遍历顺序对应依赖方向。放到“动态规划与状态定义”这条主线里看，关键原则是：状态定义要包含足够信息但不能过度膨胀。线性 DP 看最后一步选择，二维 DP 看两个前缀或网格位置，背包看物品阶段和容量，区间 DP 看最后合并点。

DP 证明通常看最后一步。任意最优解或合法方案的最后一步必在转移枚举中；每个转移来源又已经按遍历顺序算完。本题的末尾字符相同继承左上；不同则枚举插入、删除、替换要支撑这个依赖关系。这道题的边界不能留到最后补丁式处理，实验时覆盖空串、完全相同、完全不同、操作代价不同，先打印完整 DP 表，再比较空间压缩版本。若压缩版不同，优先检查遍历顺序和旧值保存。

C++ 实现上，DP 表初始化要对应状态语义，不可达哨兵不能溢出。空间压缩前先写完整表；压缩后要说清读到的是上一轮、本轮左侧还是左上角旧值。如果追问变成“若只允许插入删除，问题退化到和最长公共子序列相关”，先判断原来的状态是否仍然覆盖新答案集合，再决定是微调边界、改 value 语义，还是换成另一类结构。

\subsubsection{LeetCode 121 买卖股票的最佳时机}

先写递归或枚举，把重复子问题暴露出来。本题的模型是卖出日固定时，最佳买入日是此前最低价格；慢点不是递归形式，而是相同状态被反复求。本题需要命名的状态来自扫描维护历史最低价和最大利润，一次交易无需复杂状态机，状态定义必须让未来决策不再依赖完整历史。

DP 状态必须足够描述未来。本题的扫描维护历史最低价和最大利润，一次交易无需复杂状态机要写成数组下标、容量、前缀长度、持有状态或区间边界。初始化对应空状态，遍历顺序对应依赖方向。放到“动态规划与状态定义”这条主线里看，关键原则是：状态定义要包含足够信息但不能过度膨胀。线性 DP 看最后一步选择，二维 DP 看两个前缀或网格位置，背包看物品阶段和容量，区间 DP 看最后合并点。

DP 证明通常看最后一步。任意最优解或合法方案的最后一步必在转移枚举中；每个转移来源又已经按遍历顺序算完。本题的扫描维护历史最低价和最大利润，一次交易无需复杂状态机要支撑这个依赖关系。这道题的边界不能留到最后补丁式处理，实验时覆盖单日价格、一直下跌、一直上涨、价格相同，先打印完整 DP 表，再比较空间压缩版本。若压缩版不同，优先检查遍历顺序和旧值保存。

C++ 实现上，DP 表初始化要对应状态语义，不可达哨兵不能溢出。空间压缩前先写完整表；压缩后要说清读到的是上一轮、本轮左侧还是左上角旧值。如果追问变成“多次交易、冷冻期和手续费版本都要扩展成状态机”，先判断原来的状态是否仍然覆盖新答案集合，再决定是微调边界、改 value 语义，还是换成另一类结构。

\subsubsection{LeetCode 123 买卖股票 III}

先写递归或枚举，把重复子问题暴露出来。本题的模型是最多两次交易，状态需要包含交易次数和持有状态；慢点不是递归形式，而是相同状态被反复求。本题需要命名的状态来自维护 buy1、sell1、buy2、sell2 四个状态，按天更新，状态定义必须让未来决策不再依赖完整历史。

DP 状态必须足够描述未来。本题的维护 buy1、sell1、buy2、sell2 四个状态，按天更新要写成数组下标、容量、前缀长度、持有状态或区间边界。初始化对应空状态，遍历顺序对应依赖方向。放到“动态规划与状态定义”这条主线里看，关键原则是：状态定义要包含足够信息但不能过度膨胀。线性 DP 看最后一步选择，二维 DP 看两个前缀或网格位置，背包看物品阶段和容量，区间 DP 看最后合并点。

DP 证明通常看最后一步。任意最优解或合法方案的最后一步必在转移枚举中；每个转移来源又已经按遍历顺序算完。本题的维护 buy1、sell1、buy2、sell2 四个状态，按天更新要支撑这个依赖关系。这道题的边界不能留到最后补丁式处理，实验时覆盖价格为空、只够一次交易、两次交易间不能重叠，先打印完整 DP 表，再比较空间压缩版本。若压缩版不同，优先检查遍历顺序和旧值保存。

C++ 实现上，DP 表初始化要对应状态语义，不可达哨兵不能溢出。空间压缩前先写完整表；压缩后要说清读到的是上一轮、本轮左侧还是左上角旧值。如果追问变成“最多 k 次交易是它的泛化，k 很大时退化为无限次交易”，先判断原来的状态是否仍然覆盖新答案集合，再决定是微调边界、改 value 语义，还是换成另一类结构。

\subsubsection{LeetCode 139 单词拆分}

先写递归或枚举，把重复子问题暴露出来。本题的模型是状态表示前缀能否由字典单词拼出；慢点不是递归形式，而是相同状态被反复求。本题需要命名的状态来自枚举最后一个单词的起点，左侧前缀可达且子串在字典中则当前可达，状态定义必须让未来决策不再依赖完整历史。

DP 状态必须足够描述未来。本题的枚举最后一个单词的起点，左侧前缀可达且子串在字典中则当前可达要写成数组下标、容量、前缀长度、持有状态或区间边界。初始化对应空状态，遍历顺序对应依赖方向。放到“动态规划与状态定义”这条主线里看，关键原则是：状态定义要包含足够信息但不能过度膨胀。线性 DP 看最后一步选择，二维 DP 看两个前缀或网格位置，背包看物品阶段和容量，区间 DP 看最后合并点。

DP 证明通常看最后一步。任意最优解或合法方案的最后一步必在转移枚举中；每个转移来源又已经按遍历顺序算完。本题的枚举最后一个单词的起点，左侧前缀可达且子串在字典中则当前可达要支撑这个依赖关系。这道题的边界不能留到最后补丁式处理，实验时覆盖空串、重复单词、长单词、substr 复制成本，先打印完整 DP 表，再比较空间压缩版本。若压缩版不同，优先检查遍历顺序和旧值保存。

C++ 实现上，DP 表初始化要对应状态语义，不可达哨兵不能溢出。空间压缩前先写完整表；压缩后要说清读到的是上一轮、本轮左侧还是左上角旧值。如果追问变成“若要输出所有句子，需要 DFS 加记忆化或前驱列表”，先判断原来的状态是否仍然覆盖新答案集合，再决定是微调边界、改 value 语义，还是换成另一类结构。

\subsubsection{LeetCode 152 乘积最大子数组}

先写递归或枚举，把重复子问题暴露出来。本题的模型是负数会让最大乘积和最小乘积互换；慢点不是递归形式，而是相同状态被反复求。本题需要命名的状态来自同时维护以当前位置结尾的最大和最小乘积，状态定义必须让未来决策不再依赖完整历史。

DP 状态必须足够描述未来。本题的同时维护以当前位置结尾的最大和最小乘积要写成数组下标、容量、前缀长度、持有状态或区间边界。初始化对应空状态，遍历顺序对应依赖方向。放到“动态规划与状态定义”这条主线里看，关键原则是：状态定义要包含足够信息但不能过度膨胀。线性 DP 看最后一步选择，二维 DP 看两个前缀或网格位置，背包看物品阶段和容量，区间 DP 看最后合并点。

DP 证明通常看最后一步。任意最优解或合法方案的最后一步必在转移枚举中；每个转移来源又已经按遍历顺序算完。本题的同时维护以当前位置结尾的最大和最小乘积要支撑这个依赖关系。这道题的边界不能留到最后补丁式处理，实验时覆盖零、负数个数奇偶、单元素、乘积溢出，先打印完整 DP 表，再比较空间压缩版本。若压缩版不同，优先检查遍历顺序和旧值保存。

C++ 实现上，DP 表初始化要对应状态语义，不可达哨兵不能溢出。空间压缩前先写完整表；压缩后要说清读到的是上一轮、本轮左侧还是左上角旧值。如果追问变成“如果要求返回区间，同步记录状态来源”，先判断原来的状态是否仍然覆盖新答案集合，再决定是微调边界、改 value 语义，还是换成另一类结构。

\subsubsection{LeetCode 198 打家劫舍}

先写递归或枚举，把重复子问题暴露出来。本题的模型是每间房只有偷和不偷两种结局，偷当前就不能偷前一间；慢点不是递归形式，而是相同状态被反复求。本题需要命名的状态来自滚动保存前一状态和前两状态，状态定义必须让未来决策不再依赖完整历史。

DP 状态必须足够描述未来。本题的滚动保存前一状态和前两状态要写成数组下标、容量、前缀长度、持有状态或区间边界。初始化对应空状态，遍历顺序对应依赖方向。放到“动态规划与状态定义”这条主线里看，关键原则是：状态定义要包含足够信息但不能过度膨胀。线性 DP 看最后一步选择，二维 DP 看两个前缀或网格位置，背包看物品阶段和容量，区间 DP 看最后合并点。

DP 证明通常看最后一步。任意最优解或合法方案的最后一步必在转移枚举中；每个转移来源又已经按遍历顺序算完。本题的滚动保存前一状态和前两状态要支撑这个依赖关系。这道题的边界不能留到最后补丁式处理，实验时覆盖空数组、单间、全零、金额很大，先打印完整 DP 表，再比较空间压缩版本。若压缩版不同，优先检查遍历顺序和旧值保存。

C++ 实现上，DP 表初始化要对应状态语义，不可达哨兵不能溢出。空间压缩前先写完整表；压缩后要说清读到的是上一轮、本轮左侧还是左上角旧值。如果追问变成“环形房屋要拆成不含首和不含尾两个线性问题”，先判断原来的状态是否仍然覆盖新答案集合，再决定是微调边界、改 value 语义，还是换成另一类结构。

\subsubsection{LeetCode 221 最大正方形}

先写递归或枚举，把重复子问题暴露出来。本题的模型是以当前位置为右下角的最大正方形由上、左、左上三个状态限制；慢点不是递归形式，而是相同状态被反复求。本题需要命名的状态来自当前位置为一时取三者最小加一，否则为零，状态定义必须让未来决策不再依赖完整历史。

DP 状态必须足够描述未来。本题的当前位置为一时取三者最小加一，否则为零要写成数组下标、容量、前缀长度、持有状态或区间边界。初始化对应空状态，遍历顺序对应依赖方向。放到“动态规划与状态定义”这条主线里看，关键原则是：状态定义要包含足够信息但不能过度膨胀。线性 DP 看最后一步选择，二维 DP 看两个前缀或网格位置，背包看物品阶段和容量，区间 DP 看最后合并点。

DP 证明通常看最后一步。任意最优解或合法方案的最后一步必在转移枚举中；每个转移来源又已经按遍历顺序算完。本题的当前位置为一时取三者最小加一，否则为零要支撑这个依赖关系。这道题的边界不能留到最后补丁式处理，实验时覆盖全零、全一、单行单列、字符一和数字一不要混淆，先打印完整 DP 表，再比较空间压缩版本。若压缩版不同，优先检查遍历顺序和旧值保存。

C++ 实现上，DP 表初始化要对应状态语义，不可达哨兵不能溢出。空间压缩前先写完整表；压缩后要说清读到的是上一轮、本轮左侧还是左上角旧值。如果追问变成“最大矩形则需要按行构造柱状图再用单调栈”，先判断原来的状态是否仍然覆盖新答案集合，再决定是微调边界、改 value 语义，还是换成另一类结构。

\subsubsection{LeetCode 300 最长递增子序列}

先写递归或枚举，把重复子问题暴露出来。本题的模型是二次 DP 固定结尾，二分优化维护每个长度的最小结尾；慢点不是递归形式，而是相同状态被反复求。本题需要命名的状态来自tails 不是实际答案序列，而是未来扩展潜力表，状态定义必须让未来决策不再依赖完整历史。

DP 状态必须足够描述未来。本题的tails 不是实际答案序列，而是未来扩展潜力表要写成数组下标、容量、前缀长度、持有状态或区间边界。初始化对应空状态，遍历顺序对应依赖方向。放到“动态规划与状态定义”这条主线里看，关键原则是：状态定义要包含足够信息但不能过度膨胀。线性 DP 看最后一步选择，二维 DP 看两个前缀或网格位置，背包看物品阶段和容量，区间 DP 看最后合并点。

DP 证明通常看最后一步。任意最优解或合法方案的最后一步必在转移枚举中；每个转移来源又已经按遍历顺序算完。本题的tails 不是实际答案序列，而是未来扩展潜力表要支撑这个依赖关系。这道题的边界不能留到最后补丁式处理，实验时覆盖重复元素、严格递增和非递减区别、空数组，先打印完整 DP 表，再比较空间压缩版本。若压缩版不同，优先检查遍历顺序和旧值保存。

C++ 实现上，DP 表初始化要对应状态语义，不可达哨兵不能溢出。空间压缩前先写完整表；压缩后要说清读到的是上一轮、本轮左侧还是左上角旧值。如果追问变成“若要还原路径，需要额外前驱数组，不能只看 tails”，先判断原来的状态是否仍然覆盖新答案集合，再决定是微调边界、改 value 语义，还是换成另一类结构。

\subsubsection{LeetCode 309 最佳买卖股票含冷冻期}

先写递归或枚举，把重复子问题暴露出来。本题的模型是冷冻期让买入来源受昨天卖出状态限制；慢点不是递归形式，而是相同状态被反复求。本题需要命名的状态来自状态机维护持有、刚卖出、休息三种日终状态，状态定义必须让未来决策不再依赖完整历史。

DP 状态必须足够描述未来。本题的状态机维护持有、刚卖出、休息三种日终状态要写成数组下标、容量、前缀长度、持有状态或区间边界。初始化对应空状态，遍历顺序对应依赖方向。放到“动态规划与状态定义”这条主线里看，关键原则是：状态定义要包含足够信息但不能过度膨胀。线性 DP 看最后一步选择，二维 DP 看两个前缀或网格位置，背包看物品阶段和容量，区间 DP 看最后合并点。

DP 证明通常看最后一步。任意最优解或合法方案的最后一步必在转移枚举中；每个转移来源又已经按遍历顺序算完。本题的状态机维护持有、刚卖出、休息三种日终状态要支撑这个依赖关系。这道题的边界不能留到最后补丁式处理，实验时覆盖空价格、一天、连续上涨、卖出后不能次日买入，先打印完整 DP 表，再比较空间压缩版本。若压缩版不同，优先检查遍历顺序和旧值保存。

C++ 实现上，DP 表初始化要对应状态语义，不可达哨兵不能溢出。空间压缩前先写完整表；压缩后要说清读到的是上一轮、本轮左侧还是左上角旧值。如果追问变成“手续费版本可以在买入或卖出转移中扣费”，先判断原来的状态是否仍然覆盖新答案集合，再决定是微调边界、改 value 语义，还是换成另一类结构。

\subsubsection{LeetCode 714 买卖股票含手续费}

先写递归或枚举，把重复子问题暴露出来。本题的模型是手续费改变交易收益，仍可用持有和不持有状态；慢点不是递归形式，而是相同状态被反复求。本题需要命名的状态来自卖出时扣费或买入时扣费都可以，但不要重复扣，状态定义必须让未来决策不再依赖完整历史。

DP 状态必须足够描述未来。本题的卖出时扣费或买入时扣费都可以，但不要重复扣要写成数组下标、容量、前缀长度、持有状态或区间边界。初始化对应空状态，遍历顺序对应依赖方向。放到“动态规划与状态定义”这条主线里看，关键原则是：状态定义要包含足够信息但不能过度膨胀。线性 DP 看最后一步选择，二维 DP 看两个前缀或网格位置，背包看物品阶段和容量，区间 DP 看最后合并点。

DP 证明通常看最后一步。任意最优解或合法方案的最后一步必在转移枚举中；每个转移来源又已经按遍历顺序算完。本题的卖出时扣费或买入时扣费都可以，但不要重复扣要支撑这个依赖关系。这道题的边界不能留到最后补丁式处理，实验时覆盖手续费为零、价格震荡、小利润不足手续费，先打印完整 DP 表，再比较空间压缩版本。若压缩版不同，优先检查遍历顺序和旧值保存。

C++ 实现上，DP 表初始化要对应状态语义，不可达哨兵不能溢出。空间压缩前先写完整表；压缩后要说清读到的是上一轮、本轮左侧还是左上角旧值。如果追问变成“这是状态机比局部贪心更稳的例子”，先判断原来的状态是否仍然覆盖新答案集合，再决定是微调边界、改 value 语义，还是换成另一类结构。

\subsubsection{LeetCode 1143 最长公共子序列}

先写递归或枚举，把重复子问题暴露出来。本题的模型是状态表示两个前缀的最长公共子序列长度；慢点不是递归形式，而是相同状态被反复求。本题需要命名的状态来自末尾相同取左上加一，不同取上方和左方最大，状态定义必须让未来决策不再依赖完整历史。

DP 状态必须足够描述未来。本题的末尾相同取左上加一，不同取上方和左方最大要写成数组下标、容量、前缀长度、持有状态或区间边界。初始化对应空状态，遍历顺序对应依赖方向。放到“动态规划与状态定义”这条主线里看，关键原则是：状态定义要包含足够信息但不能过度膨胀。线性 DP 看最后一步选择，二维 DP 看两个前缀或网格位置，背包看物品阶段和容量，区间 DP 看最后合并点。

DP 证明通常看最后一步。任意最优解或合法方案的最后一步必在转移枚举中；每个转移来源又已经按遍历顺序算完。本题的末尾相同取左上加一，不同取上方和左方最大要支撑这个依赖关系。这道题的边界不能留到最后补丁式处理，实验时覆盖空串、完全相同、无公共字符、子序列不是子串，先打印完整 DP 表，再比较空间压缩版本。若压缩版不同，优先检查遍历顺序和旧值保存。

C++ 实现上，DP 表初始化要对应状态语义，不可达哨兵不能溢出。空间压缩前先写完整表；压缩后要说清读到的是上一轮、本轮左侧还是左上角旧值。如果追问变成“若要求还原序列，需要从表右下角反向追踪选择”，先判断原来的状态是否仍然覆盖新答案集合，再决定是微调边界、改 value 语义，还是换成另一类结构。

\subsubsection{LeetCode 1235 规划兼职工作}

先写递归或枚举，把重复子问题暴露出来。本题的模型是每份工作有开始、结束和收益，选择不冲突工作使收益最大；慢点不是递归形式，而是相同状态被反复求。本题需要命名的状态来自按结束时间排序，dp[i] 在不选当前和选当前加上兼容前驱之间取最大，前驱用二分找，状态定义必须让未来决策不再依赖完整历史。

DP 状态必须足够描述未来。本题的按结束时间排序，dp[i] 在不选当前和选当前加上兼容前驱之间取最大，前驱用二分找要写成数组下标、容量、前缀长度、持有状态或区间边界。初始化对应空状态，遍历顺序对应依赖方向。放到“动态规划与状态定义”这条主线里看，关键原则是：状态定义要包含足够信息但不能过度膨胀。线性 DP 看最后一步选择，二维 DP 看两个前缀或网格位置，背包看物品阶段和容量，区间 DP 看最后合并点。

DP 证明通常看最后一步。任意最优解或合法方案的最后一步必在转移枚举中；每个转移来源又已经按遍历顺序算完。本题的按结束时间排序，dp[i] 在不选当前和选当前加上兼容前驱之间取最大，前驱用二分找要支撑这个依赖关系。这道题的边界不能留到最后补丁式处理，实验时覆盖结束等于开始是否兼容、收益相同、工作排序、空列表，先打印完整 DP 表，再比较空间压缩版本。若压缩版不同，优先检查遍历顺序和旧值保存。

C++ 实现上，DP 表初始化要对应状态语义，不可达哨兵不能溢出。空间压缩前先写完整表；压缩后要说清读到的是上一轮、本轮左侧还是左上角旧值。如果追问变成“这是排序、二分和 DP 的组合，不是按收益贪心”，先判断原来的状态是否仍然覆盖新答案集合，再决定是微调边界、改 value 语义，还是换成另一类结构。

\subsection{背包模型工作坊}

先枚举物品选择，再把剩余容量作为状态。0-1 背包倒序容量，完全背包正序容量，组合和排列的循环顺序不能混用。

\subsubsection{LeetCode 279 完全平方数}

先枚举每个物品选或不选、选几次。本题的模型是完全平方数可以看成无限使用的物品，目标是凑出 n 的最少数量；暴力慢在相同阶段和容量反复出现。优化点完全背包最少物品数，或把数字看成图节点做 BFS决定容量循环方向、状态值含义和是否允许当前物品在同一轮被再次使用。

背包状态围绕阶段和容量。本题的完全背包最少物品数，或把数字看成图节点做 BFS决定倒序还是正序、求最值还是计数、组合还是排列。循环顺序不是写法偏好，而是题意的一部分。放到“背包模型”这条主线里看，关键原则是：0-1 背包每个物品只能用一次，一维压缩时容量倒序；完全背包物品可无限使用，容量正序。计数组合和计数排列的循环顺序不同。

背包证明围绕当前物品使用次数。0-1 只有选和不选，完全背包允许同一物品继续转移。本题的完全背包最少物品数，或把数字看成图节点做 BFS要与循环方向一致，否则会把题意改掉。这道题的边界不能留到最后补丁式处理，实验时覆盖n 为零或一、完全平方数本身、较大 n，并故意把容量方向写反构造最小反例。这样能确认循环方向来自题意，不是背模板。

C++ 实现上，一维背包数组的容量方向要和题意一致。方案数用 \code{long long} 或题目要求的取模；容量来自数值大小时，复杂度是伪多项式。如果追问变成“数学四平方定理可给更强解法，但面试 DP 更通用”，先判断原来的状态是否仍然覆盖新答案集合，再决定是微调边界、改 value 语义，还是换成另一类结构。

\subsubsection{LeetCode 322 零钱兑换}

先枚举每个物品选或不选、选几次。本题的模型是金额是容量，硬币可无限使用，目标是最少硬币数；暴力慢在相同阶段和容量反复出现。优化点dp[value] 从 value 减 coin 的已知状态转移决定容量循环方向、状态值含义和是否允许当前物品在同一轮被再次使用。

背包状态围绕阶段和容量。本题的dp[value] 从 value 减 coin 的已知状态转移决定倒序还是正序、求最值还是计数、组合还是排列。循环顺序不是写法偏好，而是题意的一部分。放到“背包模型”这条主线里看，关键原则是：0-1 背包每个物品只能用一次，一维压缩时容量倒序；完全背包物品可无限使用，容量正序。计数组合和计数排列的循环顺序不同。

背包证明围绕当前物品使用次数。0-1 只有选和不选，完全背包允许同一物品继续转移。本题的dp[value] 从 value 减 coin 的已知状态转移要与循环方向一致，否则会把题意改掉。这道题的边界不能留到最后补丁式处理，实验时覆盖无法凑出、amount 为零、硬币面额大于目标、哨兵溢出，并故意把容量方向写反构造最小反例。这样能确认循环方向来自题意，不是背模板。

C++ 实现上，一维背包数组的容量方向要和题意一致。方案数用 \code{long long} 或题目要求的取模；容量来自数值大小时，复杂度是伪多项式。如果追问变成“零钱兑换 II 问组合数，循环语义不同”，先判断原来的状态是否仍然覆盖新答案集合，再决定是微调边界、改 value 语义，还是换成另一类结构。

\subsubsection{LeetCode 377 组合总和 IV}

先枚举每个物品选或不选、选几次。本题的模型是题目统计排列数，选择顺序不同算不同方案；暴力慢在相同阶段和容量反复出现。优化点外层遍历容量，内层遍历数字，让每个容量的方案按最后一个数累加决定容量循环方向、状态值含义和是否允许当前物品在同一轮被再次使用。

背包状态围绕阶段和容量。本题的外层遍历容量，内层遍历数字，让每个容量的方案按最后一个数累加决定倒序还是正序、求最值还是计数、组合还是排列。循环顺序不是写法偏好，而是题意的一部分。放到“背包模型”这条主线里看，关键原则是：0-1 背包每个物品只能用一次，一维压缩时容量倒序；完全背包物品可无限使用，容量正序。计数组合和计数排列的循环顺序不同。

背包证明围绕当前物品使用次数。0-1 只有选和不选，完全背包允许同一物品继续转移。本题的外层遍历容量，内层遍历数字，让每个容量的方案按最后一个数累加要与循环方向一致，否则会把题意改掉。这道题的边界不能留到最后补丁式处理，实验时覆盖target 为零、数字大于目标、方案数溢出、循环顺序，并故意把容量方向写反构造最小反例。这样能确认循环方向来自题意，不是背模板。

C++ 实现上，一维背包数组的容量方向要和题意一致。方案数用 \code{long long} 或题目要求的取模；容量来自数值大小时，复杂度是伪多项式。如果追问变成“若顺序不同不算新方案，就变成完全背包组合数，循环顺序要交换”，先判断原来的状态是否仍然覆盖新答案集合，再决定是微调边界、改 value 语义，还是换成另一类结构。

\subsubsection{LeetCode 416 分割等和子集}

先枚举每个物品选或不选、选几次。本题的模型是总和一半是背包容量，每个数只能使用一次；暴力慢在相同阶段和容量反复出现。优化点0-1 背包可达性，一维容量必须倒序决定容量循环方向、状态值含义和是否允许当前物品在同一轮被再次使用。

背包状态围绕阶段和容量。本题的0-1 背包可达性，一维容量必须倒序决定倒序还是正序、求最值还是计数、组合还是排列。循环顺序不是写法偏好，而是题意的一部分。放到“背包模型”这条主线里看，关键原则是：0-1 背包每个物品只能用一次，一维压缩时容量倒序；完全背包物品可无限使用，容量正序。计数组合和计数排列的循环顺序不同。

背包证明围绕当前物品使用次数。0-1 只有选和不选，完全背包允许同一物品继续转移。本题的0-1 背包可达性，一维容量必须倒序要与循环方向一致，否则会把题意改掉。这道题的边界不能留到最后补丁式处理，实验时覆盖总和为奇数、目标为零、数值较大、复杂度依赖 target，并故意把容量方向写反构造最小反例。这样能确认循环方向来自题意，不是背模板。

C++ 实现上，一维背包数组的容量方向要和题意一致。方案数用 \code{long long} 或题目要求的取模；容量来自数值大小时，复杂度是伪多项式。如果追问变成“负数会破坏普通容量模型，需要重新建模”，先判断原来的状态是否仍然覆盖新答案集合，再决定是微调边界、改 value 语义，还是换成另一类结构。

\subsubsection{LeetCode 474 一和零}

先枚举每个物品选或不选、选几次。本题的模型是每个字符串消耗零和一两种容量，每个字符串最多选一次；暴力慢在相同阶段和容量反复出现。优化点二维 0-1 背包，两个容量都要倒序遍历决定容量循环方向、状态值含义和是否允许当前物品在同一轮被再次使用。

背包状态围绕阶段和容量。本题的二维 0-1 背包，两个容量都要倒序遍历决定倒序还是正序、求最值还是计数、组合还是排列。循环顺序不是写法偏好，而是题意的一部分。放到“背包模型”这条主线里看，关键原则是：0-1 背包每个物品只能用一次，一维压缩时容量倒序；完全背包物品可无限使用，容量正序。计数组合和计数排列的循环顺序不同。

背包证明围绕当前物品使用次数。0-1 只有选和不选，完全背包允许同一物品继续转移。本题的二维 0-1 背包，两个容量都要倒序遍历要与循环方向一致，否则会把题意改掉。这道题的边界不能留到最后补丁式处理，实验时覆盖字符串全零或全一、容量为零、倒序方向、计数预处理，并故意把容量方向写反构造最小反例。这样能确认循环方向来自题意，不是背模板。

C++ 实现上，一维背包数组的容量方向要和题意一致。方案数用 \code{long long} 或题目要求的取模；容量来自数值大小时，复杂度是伪多项式。如果追问变成“这是背包从一维容量扩展到二维容量的代表”，先判断原来的状态是否仍然覆盖新答案集合，再决定是微调边界、改 value 语义，还是换成另一类结构。

\subsubsection{LeetCode 494 目标和}

先枚举每个物品选或不选、选几次。本题的模型是正负号选择可代数转换为选若干数凑出特定正和；暴力慢在相同阶段和容量反复出现。优化点若 S 加 target 为奇数或目标绝对值过大，直接无解决定容量循环方向、状态值含义和是否允许当前物品在同一轮被再次使用。

背包状态围绕阶段和容量。本题的若 S 加 target 为奇数或目标绝对值过大，直接无解决定倒序还是正序、求最值还是计数、组合还是排列。循环顺序不是写法偏好，而是题意的一部分。放到“背包模型”这条主线里看，关键原则是：0-1 背包每个物品只能用一次，一维压缩时容量倒序；完全背包物品可无限使用，容量正序。计数组合和计数排列的循环顺序不同。

背包证明围绕当前物品使用次数。0-1 只有选和不选，完全背包允许同一物品继续转移。本题的若 S 加 target 为奇数或目标绝对值过大，直接无解要与循环方向一致，否则会把题意改掉。这道题的边界不能留到最后补丁式处理，实验时覆盖零元素会让方案数翻倍、目标为负、容量为零，并故意把容量方向写反构造最小反例。这样能确认循环方向来自题意，不是背模板。

C++ 实现上，一维背包数组的容量方向要和题意一致。方案数用 \code{long long} 或题目要求的取模；容量来自数值大小时，复杂度是伪多项式。如果追问变成“也可用哈希 DP 保存所有可能和，但背包转换更高效”，先判断原来的状态是否仍然覆盖新答案集合，再决定是微调边界、改 value 语义，还是换成另一类结构。

\subsubsection{LeetCode 518 零钱兑换 II}

先枚举每个物品选或不选、选几次。本题的模型是硬币无限使用，问组合数而不是最少数量；暴力慢在相同阶段和容量反复出现。优化点外层遍历硬币、内层金额正序，保证组合按固定硬币顺序统计决定容量循环方向、状态值含义和是否允许当前物品在同一轮被再次使用。

背包状态围绕阶段和容量。本题的外层遍历硬币、内层金额正序，保证组合按固定硬币顺序统计决定倒序还是正序、求最值还是计数、组合还是排列。循环顺序不是写法偏好，而是题意的一部分。放到“背包模型”这条主线里看，关键原则是：0-1 背包每个物品只能用一次，一维压缩时容量倒序；完全背包物品可无限使用，容量正序。计数组合和计数排列的循环顺序不同。

背包证明围绕当前物品使用次数。0-1 只有选和不选，完全背包允许同一物品继续转移。本题的外层遍历硬币、内层金额正序，保证组合按固定硬币顺序统计要与循环方向一致，否则会把题意改掉。这道题的边界不能留到最后补丁式处理，实验时覆盖amount 为零、无硬币、组合数较大、循环顺序反例，并故意把容量方向写反构造最小反例。这样能确认循环方向来自题意，不是背模板。

C++ 实现上，一维背包数组的容量方向要和题意一致。方案数用 \code{long long} 或题目要求的取模；容量来自数值大小时，复杂度是伪多项式。如果追问变成“若外层金额内层硬币，会统计排列数”，先判断原来的状态是否仍然覆盖新答案集合，再决定是微调边界、改 value 语义，还是换成另一类结构。

\subsubsection{LeetCode 879 盈利计划}

先枚举每个物品选或不选、选几次。本题的模型是每个项目消耗人数并贡献利润，利润达到下限后可截断；暴力慢在相同阶段和容量反复出现。优化点二维 DP 用人数和已达到利润表示方案数，利润维度超过目标就压到目标决定容量循环方向、状态值含义和是否允许当前物品在同一轮被再次使用。

背包状态围绕阶段和容量。本题的二维 DP 用人数和已达到利润表示方案数，利润维度超过目标就压到目标决定倒序还是正序、求最值还是计数、组合还是排列。循环顺序不是写法偏好，而是题意的一部分。放到“背包模型”这条主线里看，关键原则是：0-1 背包每个物品只能用一次，一维压缩时容量倒序；完全背包物品可无限使用，容量正序。计数组合和计数排列的循环顺序不同。

背包证明围绕当前物品使用次数。0-1 只有选和不选，完全背包允许同一物品继续转移。本题的二维 DP 用人数和已达到利润表示方案数，利润维度超过目标就压到目标要与循环方向一致，否则会把题意改掉。这道题的边界不能留到最后补丁式处理，实验时覆盖minProfit 为零、人数不足、方案数取模、项目只能选一次，并故意把容量方向写反构造最小反例。这样能确认循环方向来自题意，不是背模板。

C++ 实现上，一维背包数组的容量方向要和题意一致。方案数用 \code{long long} 或题目要求的取模；容量来自数值大小时，复杂度是伪多项式。如果追问变成“这是 0-1 背包的计数扩展，目标维度需要饱和处理”，先判断原来的状态是否仍然覆盖新答案集合，再决定是微调边界、改 value 语义，还是换成另一类结构。

\subsubsection{LeetCode 956 最高的广告牌}

先枚举每个物品选或不选、选几次。本题的模型是两组钢筋高度相等时得到广告牌，高度差是关键状态；暴力慢在相同阶段和容量反复出现。优化点DP 保存某个高度差下较矮一侧的最大高度，新增钢筋可放高侧、低侧或不用决定容量循环方向、状态值含义和是否允许当前物品在同一轮被再次使用。

背包状态围绕阶段和容量。本题的DP 保存某个高度差下较矮一侧的最大高度，新增钢筋可放高侧、低侧或不用决定倒序还是正序、求最值还是计数、组合还是排列。循环顺序不是写法偏好，而是题意的一部分。放到“背包模型”这条主线里看，关键原则是：0-1 背包每个物品只能用一次，一维压缩时容量倒序；完全背包物品可无限使用，容量正序。计数组合和计数排列的循环顺序不同。

背包证明围绕当前物品使用次数。0-1 只有选和不选，完全背包允许同一物品继续转移。本题的DP 保存某个高度差下较矮一侧的最大高度，新增钢筋可放高侧、低侧或不用要与循环方向一致，否则会把题意改掉。这道题的边界不能留到最后补丁式处理，实验时覆盖空数组、无法平衡、重复长度、差值更新顺序，并故意把容量方向写反构造最小反例。这样能确认循环方向来自题意，不是背模板。

C++ 实现上，一维背包数组的容量方向要和题意一致。方案数用 \code{long long} 或题目要求的取模；容量来自数值大小时，复杂度是伪多项式。如果追问变成“这是背包从容量转成差值状态的代表”，先判断原来的状态是否仍然覆盖新答案集合，再决定是微调边界、改 value 语义，还是换成另一类结构。

\subsubsection{LeetCode 1049 最后一块石头的重量 II}

先枚举每个物品选或不选、选几次。本题的模型是把石头分成两堆，最终差值等于两堆重量差；暴力慢在相同阶段和容量反复出现。优化点0-1 背包找不超过总和一半的最大可达重量决定容量循环方向、状态值含义和是否允许当前物品在同一轮被再次使用。

背包状态围绕阶段和容量。本题的0-1 背包找不超过总和一半的最大可达重量决定倒序还是正序、求最值还是计数、组合还是排列。循环顺序不是写法偏好，而是题意的一部分。放到“背包模型”这条主线里看，关键原则是：0-1 背包每个物品只能用一次，一维压缩时容量倒序；完全背包物品可无限使用，容量正序。计数组合和计数排列的循环顺序不同。

背包证明围绕当前物品使用次数。0-1 只有选和不选，完全背包允许同一物品继续转移。本题的0-1 背包找不超过总和一半的最大可达重量要与循环方向一致，否则会把题意改掉。这道题的边界不能留到最后补丁式处理，实验时覆盖单块石头、总和为奇偶、多个相同重量、容量方向，并故意把容量方向写反构造最小反例。这样能确认循环方向来自题意，不是背模板。

C++ 实现上，一维背包数组的容量方向要和题意一致。方案数用 \code{long long} 或题目要求的取模；容量来自数值大小时，复杂度是伪多项式。如果追问变成“它和分割等和子集同源，只是目标从是否可达变成最小差”，先判断原来的状态是否仍然覆盖新答案集合，再决定是微调边界、改 value 语义，还是换成另一类结构。

\subsubsection{LeetCode 1155 掷骰子等于目标和的方法数}

先枚举每个物品选或不选、选几次。本题的模型是多个骰子每个只能取一到 faces 的点数，问目标和方案数；暴力慢在相同阶段和容量反复出现。优化点按骰子数量分阶段，容量为当前点数和，枚举当前骰子的点数转移决定容量循环方向、状态值含义和是否允许当前物品在同一轮被再次使用。

背包状态围绕阶段和容量。本题的按骰子数量分阶段，容量为当前点数和，枚举当前骰子的点数转移决定倒序还是正序、求最值还是计数、组合还是排列。循环顺序不是写法偏好，而是题意的一部分。放到“背包模型”这条主线里看，关键原则是：0-1 背包每个物品只能用一次，一维压缩时容量倒序；完全背包物品可无限使用，容量正序。计数组合和计数排列的循环顺序不同。

背包证明围绕当前物品使用次数。0-1 只有选和不选，完全背包允许同一物品继续转移。本题的按骰子数量分阶段，容量为当前点数和，枚举当前骰子的点数转移要与循环方向一致，否则会把题意改掉。这道题的边界不能留到最后补丁式处理，实验时覆盖target 太小或太大、取模、骰子数为一、滚动数组方向，并故意把容量方向写反构造最小反例。这样能确认循环方向来自题意，不是背模板。

C++ 实现上，一维背包数组的容量方向要和题意一致。方案数用 \code{long long} 或题目要求的取模；容量来自数值大小时，复杂度是伪多项式。如果追问变成“这是有界选择计数，和完全背包不能混”，先判断原来的状态是否仍然覆盖新答案集合，再决定是微调边界、改 value 语义，还是换成另一类结构。

\subsubsection{LeetCode 1449 数位成本和为目标值的最大数字}

先枚举每个物品选或不选、选几次。本题的模型是每个数字可重复选择，容量是成本，目标是组成字典序和长度最大的数字；暴力慢在相同阶段和容量反复出现。优化点完全背包先最大化位数，再按数字从大到小贪心还原答案决定容量循环方向、状态值含义和是否允许当前物品在同一轮被再次使用。

背包状态围绕阶段和容量。本题的完全背包先最大化位数，再按数字从大到小贪心还原答案决定倒序还是正序、求最值还是计数、组合还是排列。循环顺序不是写法偏好，而是题意的一部分。放到“背包模型”这条主线里看，关键原则是：0-1 背包每个物品只能用一次，一维压缩时容量倒序；完全背包物品可无限使用，容量正序。计数组合和计数排列的循环顺序不同。

背包证明围绕当前物品使用次数。0-1 只有选和不选，完全背包允许同一物品继续转移。本题的完全背包先最大化位数，再按数字从大到小贪心还原答案要与循环方向一致，否则会把题意改掉。这道题的边界不能留到最后补丁式处理，实验时覆盖无法组成目标、多个数字同成本、target 为零、字符串比较，并故意把容量方向写反构造最小反例。这样能确认循环方向来自题意，不是背模板。

C++ 实现上，一维背包数组的容量方向要和题意一致。方案数用 \code{long long} 或题目要求的取模；容量来自数值大小时，复杂度是伪多项式。如果追问变成“这题说明背包状态值不一定是数值收益，也可以是长度和字典序”，先判断原来的状态是否仍然覆盖新答案集合，再决定是微调边界、改 value 语义，还是换成另一类结构。

\subsection{贪心与交换证明工作坊}

先把选择空间完整枚举，再寻找可以交换到最优解前面的局部选择。贪心必须能证明当前选择不会让后续资源变差，而不是只看排序后代码短。

\subsubsection{LeetCode 45 跳跃游戏 II}

先把选择顺序完整枚举，哪怕这个版本只能处理很小输入。本题的模型是每一次跳跃可以看成 BFS 的一层，当前层内所有位置共享同一次跳数；真正要证明的是扫描当前层时维护下一层最远边界，到达当前层末尾才增加跳数为什么不会伤害后续选择。没有这一步证明，排序或局部选择都只是猜测。

贪心状态要表达当前方案为什么不落后。本题围绕扫描当前层时维护下一层最远边界，到达当前层末尾才增加跳数保留局部最优边界；排序只是把可比较的候选排到一起，证明仍要落在交换或领先关系上。放到“贪心与交换证明”这条主线里看，关键原则是：贪心需要找到可交换的局部最优：最早结束、最小右端、当前最远覆盖、当前最大收益或最小代价。排序只是手段，不是证明。

贪心证明要么用交换，要么用领先性。围绕扫描当前层时维护下一层最远边界，到达当前层末尾才增加跳数说明把任意最优解的第一步换成贪心选择不会变差，或说明当前方案在同等步数下始终不落后。这道题的边界不能留到最后补丁式处理，实验时除了数组长度一、刚好到达、存在多个可选跳点、不要在每个位置都加跳数，还要故意替换成一个看似合理但错误的排序规则，构造反例。能解释错误贪心为何失败，才算理解正确贪心。

C++ 实现上，比较器必须满足严格弱序，相同关键字的处理要服务于证明。涉及区间端点、收益和代价时，排序键要写成业务含义，不要只凭直觉写小于号。如果追问变成“若问能否到达，只需维护全局最远可达位置”，先判断原来的状态是否仍然覆盖新答案集合，再决定是微调边界、改 value 语义，还是换成另一类结构。

\subsubsection{LeetCode 55 跳跃游戏}

先把选择顺序完整枚举，哪怕这个版本只能处理很小输入。本题的模型是扫描过程中只关心当前能到达的最远位置；真正要证明的是如果下标超过最远可达，任何之前路径都不能到达它；否则用它扩展边界为什么不会伤害后续选择。没有这一步证明，排序或局部选择都只是猜测。

贪心状态要表达当前方案为什么不落后。本题围绕如果下标超过最远可达，任何之前路径都不能到达它；否则用它扩展边界保留局部最优边界；排序只是把可比较的候选排到一起，证明仍要落在交换或领先关系上。放到“贪心与交换证明”这条主线里看，关键原则是：贪心需要找到可交换的局部最优：最早结束、最小右端、当前最远覆盖、当前最大收益或最小代价。排序只是手段，不是证明。

贪心证明要么用交换，要么用领先性。围绕如果下标超过最远可达，任何之前路径都不能到达它；否则用它扩展边界说明把任意最优解的第一步换成贪心选择不会变差，或说明当前方案在同等步数下始终不落后。这道题的边界不能留到最后补丁式处理，实验时除了第一个元素为零、数组长度一、恰好到达最后、远超最后，还要故意替换成一个看似合理但错误的排序规则，构造反例。能解释错误贪心为何失败，才算理解正确贪心。

C++ 实现上，比较器必须满足严格弱序，相同关键字的处理要服务于证明。涉及区间端点、收益和代价时，排序键要写成业务含义，不要只凭直觉写小于号。如果追问变成“求最少跳数时把可达区间看成 BFS 层”，先判断原来的状态是否仍然覆盖新答案集合，再决定是微调边界、改 value 语义，还是换成另一类结构。

\subsubsection{LeetCode 122 买卖股票 II}

先把选择顺序完整枚举，哪怕这个版本只能处理很小输入。本题的模型是允许多次交易且不能同时持有，所有正收益差都可以收集；真正要证明的是把长上涨段拆成每天买卖收益相同，因此累加正差值最优为什么不会伤害后续选择。没有这一步证明，排序或局部选择都只是猜测。

贪心状态要表达当前方案为什么不落后。本题围绕把长上涨段拆成每天买卖收益相同，因此累加正差值最优保留局部最优边界；排序只是把可比较的候选排到一起，证明仍要落在交换或领先关系上。放到“贪心与交换证明”这条主线里看，关键原则是：贪心需要找到可交换的局部最优：最早结束、最小右端、当前最远覆盖、当前最大收益或最小代价。排序只是手段，不是证明。

贪心证明要么用交换，要么用领先性。围绕把长上涨段拆成每天买卖收益相同，因此累加正差值最优说明把任意最优解的第一步换成贪心选择不会变差，或说明当前方案在同等步数下始终不落后。这道题的边界不能留到最后补丁式处理，实验时除了下降段、平台价格、只有一天、手续费不存在，还要故意替换成一个看似合理但错误的排序规则，构造反例。能解释错误贪心为何失败，才算理解正确贪心。

C++ 实现上，比较器必须满足严格弱序，相同关键字的处理要服务于证明。涉及区间端点、收益和代价时，排序键要写成业务含义，不要只凭直觉写小于号。如果追问变成“有手续费时不能简单累加正差，需要状态机或调整贪心阈值”，先判断原来的状态是否仍然覆盖新答案集合，再决定是微调边界、改 value 语义，还是换成另一类结构。

\subsubsection{LeetCode 134 加油站}

先把选择顺序完整枚举，哪怕这个版本只能处理很小输入。本题的模型是绕行一圈是否可行取决于总油量差和局部剩余油量；真正要证明的是总和为负必不可行；扫描时若当前剩余为负，起点只能放到下一站为什么不会伤害后续选择。没有这一步证明，排序或局部选择都只是猜测。

贪心状态要表达当前方案为什么不落后。本题围绕总和为负必不可行；扫描时若当前剩余为负，起点只能放到下一站保留局部最优边界；排序只是把可比较的候选排到一起，证明仍要落在交换或领先关系上。放到“贪心与交换证明”这条主线里看，关键原则是：贪心需要找到可交换的局部最优：最早结束、最小右端、当前最远覆盖、当前最大收益或最小代价。排序只是手段，不是证明。

贪心证明要么用交换，要么用领先性。围绕总和为负必不可行；扫描时若当前剩余为负，起点只能放到下一站说明把任意最优解的第一步换成贪心选择不会变差，或说明当前方案在同等步数下始终不落后。这道题的边界不能留到最后补丁式处理，实验时除了总和刚好为零、多个可行起点、某站油量为零、消耗大于补给，还要故意替换成一个看似合理但错误的排序规则，构造反例。能解释错误贪心为何失败，才算理解正确贪心。

C++ 实现上，比较器必须满足严格弱序，相同关键字的处理要服务于证明。涉及区间端点、收益和代价时，排序键要写成业务含义，不要只凭直觉写小于号。如果追问变成“它的领先性证明是前一段任意站点都不能作为起点”，先判断原来的状态是否仍然覆盖新答案集合，再决定是微调边界、改 value 语义，还是换成另一类结构。

\subsubsection{LeetCode 135 分发糖果}

先把选择顺序完整枚举，哪怕这个版本只能处理很小输入。本题的模型是相邻评分约束需要同时满足左邻和右邻两个方向；真正要证明的是左到右满足递增关系，右到左满足反向递增关系，最后取两侧要求最大值为什么不会伤害后续选择。没有这一步证明，排序或局部选择都只是猜测。

贪心状态要表达当前方案为什么不落后。本题围绕左到右满足递增关系，右到左满足反向递增关系，最后取两侧要求最大值保留局部最优边界；排序只是把可比较的候选排到一起，证明仍要落在交换或领先关系上。放到“贪心与交换证明”这条主线里看，关键原则是：贪心需要找到可交换的局部最优：最早结束、最小右端、当前最远覆盖、当前最大收益或最小代价。排序只是手段，不是证明。

贪心证明要么用交换，要么用领先性。围绕左到右满足递增关系，右到左满足反向递增关系，最后取两侧要求最大值说明把任意最优解的第一步换成贪心选择不会变差，或说明当前方案在同等步数下始终不落后。这道题的边界不能留到最后补丁式处理，实验时除了评分相等、严格递增、严格递减、山峰和谷底，还要故意替换成一个看似合理但错误的排序规则，构造反例。能解释错误贪心为何失败，才算理解正确贪心。

C++ 实现上，比较器必须满足严格弱序，相同关键字的处理要服务于证明。涉及区间端点、收益和代价时，排序键要写成业务含义，不要只凭直觉写小于号。如果追问变成“这是局部约束两遍扫描，不能只用一个方向贪心”，先判断原来的状态是否仍然覆盖新答案集合，再决定是微调边界、改 value 语义，还是换成另一类结构。

\subsubsection{LeetCode 169 多数元素}

先把选择顺序完整枚举，哪怕这个版本只能处理很小输入。本题的模型是超过一半的元素可以和其他元素两两抵消后仍然剩下；真正要证明的是Boyer-Moore 投票维护候选和计数为什么不会伤害后续选择。没有这一步证明，排序或局部选择都只是猜测。

贪心状态要表达当前方案为什么不落后。本题围绕Boyer-Moore 投票维护候选和计数保留局部最优边界；排序只是把可比较的候选排到一起，证明仍要落在交换或领先关系上。放到“贪心与交换证明”这条主线里看，关键原则是：贪心需要找到可交换的局部最优：最早结束、最小右端、当前最远覆盖、当前最大收益或最小代价。排序只是手段，不是证明。

贪心证明要么用交换，要么用领先性。围绕Boyer-Moore 投票维护候选和计数说明把任意最优解的第一步换成贪心选择不会变差，或说明当前方案在同等步数下始终不落后。这道题的边界不能留到最后补丁式处理，实验时除了题目是否保证存在多数、计数归零、全相同，还要故意替换成一个看似合理但错误的排序规则，构造反例。能解释错误贪心为何失败，才算理解正确贪心。

C++ 实现上，比较器必须满足严格弱序，相同关键字的处理要服务于证明。涉及区间端点、收益和代价时，排序键要写成业务含义，不要只凭直觉写小于号。如果追问变成“若找超过三分之一的元素，需要维护两个候选”，先判断原来的状态是否仍然覆盖新答案集合，再决定是微调边界、改 value 语义，还是换成另一类结构。

\subsubsection{LeetCode 316 去除重复字母}

先把选择顺序完整枚举，哪怕这个版本只能处理很小输入。本题的模型是每个字符必须保留一次，同时结果字典序尽量小；真正要证明的是单调栈维护当前结果，若栈顶更大且后面还会出现，就可以弹出再等后面放回为什么不会伤害后续选择。没有这一步证明，排序或局部选择都只是猜测。

贪心状态要表达当前方案为什么不落后。本题围绕单调栈维护当前结果，若栈顶更大且后面还会出现，就可以弹出再等后面放回保留局部最优边界；排序只是把可比较的候选排到一起，证明仍要落在交换或领先关系上。放到“贪心与交换证明”这条主线里看，关键原则是：贪心需要找到可交换的局部最优：最早结束、最小右端、当前最远覆盖、当前最大收益或最小代价。排序只是手段，不是证明。

贪心证明要么用交换，要么用领先性。围绕单调栈维护当前结果，若栈顶更大且后面还会出现，就可以弹出再等后面放回说明把任意最优解的第一步换成贪心选择不会变差，或说明当前方案在同等步数下始终不落后。这道题的边界不能留到最后补丁式处理，实验时除了字符只出现一次、重复字符、弹出后 visited 恢复、字典序比较，还要故意替换成一个看似合理但错误的排序规则，构造反例。能解释错误贪心为何失败，才算理解正确贪心。

C++ 实现上，比较器必须满足严格弱序，相同关键字的处理要服务于证明。涉及区间端点、收益和代价时，排序键要写成业务含义，不要只凭直觉写小于号。如果追问变成“402 移掉 K 位数字也是单调栈删除，但删除预算和必须保留集合不同”，先判断原来的状态是否仍然覆盖新答案集合，再决定是微调边界、改 value 语义，还是换成另一类结构。

\subsubsection{LeetCode 406 根据身高重建队列}

先把选择顺序完整枚举，哪怕这个版本只能处理很小输入。本题的模型是高个子的位置先确定后，不会被矮个子影响其前方高个数量；真正要证明的是按身高降序、k 升序排序，再按 k 插入当前结果位置为什么不会伤害后续选择。没有这一步证明，排序或局部选择都只是猜测。

贪心状态要表达当前方案为什么不落后。本题围绕按身高降序、k 升序排序，再按 k 插入当前结果位置保留局部最优边界；排序只是把可比较的候选排到一起，证明仍要落在交换或领先关系上。放到“贪心与交换证明”这条主线里看，关键原则是：贪心需要找到可交换的局部最优：最早结束、最小右端、当前最远覆盖、当前最大收益或最小代价。排序只是手段，不是证明。

贪心证明要么用交换，要么用领先性。围绕按身高降序、k 升序排序，再按 k 插入当前结果位置说明把任意最优解的第一步换成贪心选择不会变差，或说明当前方案在同等步数下始终不落后。这道题的边界不能留到最后补丁式处理，实验时除了相同身高、k 为零、插入到末尾、结果稳定性，还要故意替换成一个看似合理但错误的排序规则，构造反例。能解释错误贪心为何失败，才算理解正确贪心。

C++ 实现上，比较器必须满足严格弱序，相同关键字的处理要服务于证明。涉及区间端点、收益和代价时，排序键要写成业务含义，不要只凭直觉写小于号。如果追问变成“这是排序后插入的构造型贪心，证明依赖高个子先固定”，先判断原来的状态是否仍然覆盖新答案集合，再决定是微调边界、改 value 语义，还是换成另一类结构。

\subsubsection{LeetCode 435 无重叠区间}

先把选择顺序完整枚举，哪怕这个版本只能处理很小输入。本题的模型是保留最多不重叠区间等价于删除最少区间；真正要证明的是按终点升序选择，结束越早给后续留下空间越多为什么不会伤害后续选择。没有这一步证明，排序或局部选择都只是猜测。

贪心状态要表达当前方案为什么不落后。本题围绕按终点升序选择，结束越早给后续留下空间越多保留局部最优边界；排序只是把可比较的候选排到一起，证明仍要落在交换或领先关系上。放到“贪心与交换证明”这条主线里看，关键原则是：贪心需要找到可交换的局部最优：最早结束、最小右端、当前最远覆盖、当前最大收益或最小代价。排序只是手段，不是证明。

贪心证明要么用交换，要么用领先性。围绕按终点升序选择，结束越早给后续留下空间越多说明把任意最优解的第一步换成贪心选择不会变差，或说明当前方案在同等步数下始终不落后。这道题的边界不能留到最后补丁式处理，实验时除了端点相接是否重叠、完全覆盖、相同终点、空列表，还要故意替换成一个看似合理但错误的排序规则，构造反例。能解释错误贪心为何失败，才算理解正确贪心。

C++ 实现上，比较器必须满足严格弱序，相同关键字的处理要服务于证明。涉及区间端点、收益和代价时，排序键要写成业务含义，不要只凭直觉写小于号。如果追问变成“用交换论证说明最早结束的选择不劣于其他第一选择”，先判断原来的状态是否仍然覆盖新答案集合，再决定是微调边界、改 value 语义，还是换成另一类结构。

\subsubsection{LeetCode 452 用最少数量的箭引爆气球}

先把选择顺序完整枚举，哪怕这个版本只能处理很小输入。本题的模型是一支箭的位置可以覆盖所有包含该位置的区间；真正要证明的是按右端点排序，当前箭放在最早结束气球的右端为什么不会伤害后续选择。没有这一步证明，排序或局部选择都只是猜测。

贪心状态要表达当前方案为什么不落后。本题围绕按右端点排序，当前箭放在最早结束气球的右端保留局部最优边界；排序只是把可比较的候选排到一起，证明仍要落在交换或领先关系上。放到“贪心与交换证明”这条主线里看，关键原则是：贪心需要找到可交换的局部最优：最早结束、最小右端、当前最远覆盖、当前最大收益或最小代价。排序只是手段，不是证明。

贪心证明要么用交换，要么用领先性。围绕按右端点排序，当前箭放在最早结束气球的右端说明把任意最优解的第一步换成贪心选择不会变差，或说明当前方案在同等步数下始终不落后。这道题的边界不能留到最后补丁式处理，实验时除了端点相同、负坐标、区间包含、闭区间边界，还要故意替换成一个看似合理但错误的排序规则，构造反例。能解释错误贪心为何失败，才算理解正确贪心。

C++ 实现上，比较器必须满足严格弱序，相同关键字的处理要服务于证明。涉及区间端点、收益和代价时，排序键要写成业务含义，不要只凭直觉写小于号。如果追问变成“它和无重叠区间是同一类最早结束贪心”，先判断原来的状态是否仍然覆盖新答案集合，再决定是微调边界、改 value 语义，还是换成另一类结构。

\subsubsection{LeetCode 630 课程表 III}

先把选择顺序完整枚举，哪怕这个版本只能处理很小输入。本题的模型是每门课有持续时间和截止日，已选课程总时长不能超过当前截止日；真正要证明的是按截止日排序，先选当前课；若超时，就用大顶堆丢掉已选中耗时最长的课为什么不会伤害后续选择。没有这一步证明，排序或局部选择都只是猜测。

贪心状态要表达当前方案为什么不落后。本题围绕按截止日排序，先选当前课；若超时，就用大顶堆丢掉已选中耗时最长的课保留局部最优边界；排序只是把可比较的候选排到一起，证明仍要落在交换或领先关系上。放到“贪心与交换证明”这条主线里看，关键原则是：贪心需要找到可交换的局部最优：最早结束、最小右端、当前最远覆盖、当前最大收益或最小代价。排序只是手段，不是证明。

贪心证明要么用交换，要么用领先性。围绕按截止日排序，先选当前课；若超时，就用大顶堆丢掉已选中耗时最长的课说明把任意最优解的第一步换成贪心选择不会变差，或说明当前方案在同等步数下始终不落后。这道题的边界不能留到最后补丁式处理，实验时除了课程无法单独完成、截止日相同、替换已选课程、空列表，还要故意替换成一个看似合理但错误的排序规则，构造反例。能解释错误贪心为何失败，才算理解正确贪心。

C++ 实现上，比较器必须满足严格弱序，相同关键字的处理要服务于证明。涉及区间端点、收益和代价时，排序键要写成业务含义，不要只凭直觉写小于号。如果追问变成“这是反悔贪心：接受后在约束被破坏时删最差候选”，先判断原来的状态是否仍然覆盖新答案集合，再决定是微调边界、改 value 语义，还是换成另一类结构。

\subsubsection{LeetCode 646 最长数对链}

先把选择顺序完整枚举，哪怕这个版本只能处理很小输入。本题的模型是每个数对像一个区间，选择链要求前一个右端小于后一个左端；真正要证明的是按右端升序选择，当前右端越小越容易接后续数对为什么不会伤害后续选择。没有这一步证明，排序或局部选择都只是猜测。

贪心状态要表达当前方案为什么不落后。本题围绕按右端升序选择，当前右端越小越容易接后续数对保留局部最优边界；排序只是把可比较的候选排到一起，证明仍要落在交换或领先关系上。放到“贪心与交换证明”这条主线里看，关键原则是：贪心需要找到可交换的局部最优：最早结束、最小右端、当前最远覆盖、当前最大收益或最小代价。排序只是手段，不是证明。

贪心证明要么用交换，要么用领先性。围绕按右端升序选择，当前右端越小越容易接后续数对说明把任意最优解的第一步换成贪心选择不会变差，或说明当前方案在同等步数下始终不落后。这道题的边界不能留到最后补丁式处理，实验时除了端点相等不允许连接、完全覆盖、负数端点、空数组，还要故意替换成一个看似合理但错误的排序规则，构造反例。能解释错误贪心为何失败，才算理解正确贪心。

C++ 实现上，比较器必须满足严格弱序，相同关键字的处理要服务于证明。涉及区间端点、收益和代价时，排序键要写成业务含义，不要只凭直觉写小于号。如果追问变成“也可以 DP，但贪心用最早结束证明更简洁”，先判断原来的状态是否仍然覆盖新答案集合，再决定是微调边界、改 value 语义，还是换成另一类结构。

\subsubsection{LeetCode 763 划分字母区间}

先把选择顺序完整枚举，哪怕这个版本只能处理很小输入。本题的模型是每个片段必须覆盖其中所有字符的最后出现位置；真正要证明的是扫描维护当前片段最远右端，到达右端即可切分为什么不会伤害后续选择。没有这一步证明，排序或局部选择都只是猜测。

贪心状态要表达当前方案为什么不落后。本题围绕扫描维护当前片段最远右端，到达右端即可切分保留局部最优边界；排序只是把可比较的候选排到一起，证明仍要落在交换或领先关系上。放到“贪心与交换证明”这条主线里看，关键原则是：贪心需要找到可交换的局部最优：最早结束、最小右端、当前最远覆盖、当前最大收益或最小代价。排序只是手段，不是证明。

贪心证明要么用交换，要么用领先性。围绕扫描维护当前片段最远右端，到达右端即可切分说明把任意最优解的第一步换成贪心选择不会变差，或说明当前方案在同等步数下始终不落后。这道题的边界不能留到最后补丁式处理，实验时除了单字符、所有字符相同、字符多次跨段，还要故意替换成一个看似合理但错误的排序规则，构造反例。能解释错误贪心为何失败，才算理解正确贪心。

C++ 实现上，比较器必须满足严格弱序，相同关键字的处理要服务于证明。涉及区间端点、收益和代价时，排序键要写成业务含义，不要只凭直觉写小于号。如果追问变成“这是最早合法切分贪心，能最大化片段数量”，先判断原来的状态是否仍然覆盖新答案集合，再决定是微调边界、改 value 语义，还是换成另一类结构。

\subsection{区间、排序与合并工作坊}

先两两比较或完整重排，再利用排序后相交区间连续出现的性质。动态区间题还要关注前驱、后继、端点闭开和覆盖关系。

\subsubsection{LeetCode 56 合并区间}

先两两比较区间，确认相交、覆盖或冲突的定义。本题的模型是排序后可能与当前区间相交的只会是结果尾部；暴力慢在每个区间都回头看太多无关区间。优化点按起点排序，扫描时维护结果最后区间的右端点通常依赖排序后相邻关系，必须先定义端点闭开。

区间结构的状态是当前右端、结果尾区间或动态有序表的相邻关系。本题的按起点排序，扫描时维护结果最后区间的右端点只允许丢掉已经不可能再冲突的区间；端点相等是否合并要先定义。放到“区间、排序与合并”这条主线里看，关键原则是：合并按起点排序，调度按终点排序，覆盖题常用起点升序且终点降序。扫描时维护当前最远右端、结果尾区间或已选区间终点。

区间题证明依赖排序后的邻接关系。若当前区间无法影响结果尾部或相邻区间，它就不会再影响更远的区间。本题的按起点排序，扫描时维护结果最后区间的右端点要说明哪些区间已经安全归档。这道题的边界不能留到最后补丁式处理，实验时覆盖端点相接、完全覆盖、无重叠、空列表、输入未排序，画出排序后的区间序列和当前结果尾部。动态版本要专门测前驱、后继和端点相接。

C++ 实现上，区间可用 \code{std::vector<std::array<int, 2>>} 或结构体表达端点，排序比较器要处理相同起点。动态区间需要 \code{std::map} 的前驱后继，而不是哈希表。如果追问变成“若区间是半开区间，端点相等不一定合并”，先判断原来的状态是否仍然覆盖新答案集合，再决定是微调边界、改 value 语义，还是换成另一类结构。

\subsubsection{LeetCode 57 插入区间}

先两两比较区间，确认相交、覆盖或冲突的定义。本题的模型是新区间只会影响与它相交的一段连续区间；暴力慢在每个区间都回头看太多无关区间。优化点先追加左侧完全不相交区间，再合并重叠段，最后追加右侧通常依赖排序后相邻关系，必须先定义端点闭开。

区间结构的状态是当前右端、结果尾区间或动态有序表的相邻关系。本题的先追加左侧完全不相交区间，再合并重叠段，最后追加右侧只允许丢掉已经不可能再冲突的区间；端点相等是否合并要先定义。放到“区间、排序与合并”这条主线里看，关键原则是：合并按起点排序，调度按终点排序，覆盖题常用起点升序且终点降序。扫描时维护当前最远右端、结果尾区间或已选区间终点。

区间题证明依赖排序后的邻接关系。若当前区间无法影响结果尾部或相邻区间，它就不会再影响更远的区间。本题的先追加左侧完全不相交区间，再合并重叠段，最后追加右侧要说明哪些区间已经安全归档。这道题的边界不能留到最后补丁式处理，实验时覆盖新区间在最前、最后、覆盖全部、刚好相邻、原列表为空，画出排序后的区间序列和当前结果尾部。动态版本要专门测前驱、后继和端点相接。

C++ 实现上，区间可用 \code{std::vector<std::array<int, 2>>} 或结构体表达端点，排序比较器要处理相同起点。动态区间需要 \code{std::map} 的前驱后继，而不是哈希表。如果追问变成“若输入不是有序且不重叠，必须先排序或换模型”，先判断原来的状态是否仍然覆盖新答案集合，再决定是微调边界、改 value 语义，还是换成另一类结构。

\subsubsection{LeetCode 252 会议室}

先两两比较区间，确认相交、覆盖或冲突的定义。本题的模型是每个会议占用一个时间区间，任意重叠都表示不能全部参加；暴力慢在每个区间都回头看太多无关区间。优化点按开始时间排序，只需检查相邻会议是否冲突通常依赖排序后相邻关系，必须先定义端点闭开。

区间结构的状态是当前右端、结果尾区间或动态有序表的相邻关系。本题的按开始时间排序，只需检查相邻会议是否冲突只允许丢掉已经不可能再冲突的区间；端点相等是否合并要先定义。放到“区间、排序与合并”这条主线里看，关键原则是：合并按起点排序，调度按终点排序，覆盖题常用起点升序且终点降序。扫描时维护当前最远右端、结果尾区间或已选区间终点。

区间题证明依赖排序后的邻接关系。若当前区间无法影响结果尾部或相邻区间，它就不会再影响更远的区间。本题的按开始时间排序，只需检查相邻会议是否冲突要说明哪些区间已经安全归档。这道题的边界不能留到最后补丁式处理，实验时覆盖端点相等是否冲突、空列表、单会议、完全重叠，画出排序后的区间序列和当前结果尾部。动态版本要专门测前驱、后继和端点相接。

C++ 实现上，区间可用 \code{std::vector<std::array<int, 2>>} 或结构体表达端点，排序比较器要处理相同起点。动态区间需要 \code{std::map} 的前驱后继，而不是哈希表。如果追问变成“若问最少会议室，需要维护同时进行的会议数量”，先判断原来的状态是否仍然覆盖新答案集合，再决定是微调边界、改 value 语义，还是换成另一类结构。

\subsubsection{LeetCode 253 会议室 II}

先两两比较区间，确认相交、覆盖或冲突的定义。本题的模型是同时进行的会议数量决定最少房间数；暴力慢在每个区间都回头看太多无关区间。优化点按开始时间扫描，用小顶堆维护当前会议的最早结束时间通常依赖排序后相邻关系，必须先定义端点闭开。

区间结构的状态是当前右端、结果尾区间或动态有序表的相邻关系。本题的按开始时间扫描，用小顶堆维护当前会议的最早结束时间只允许丢掉已经不可能再冲突的区间；端点相等是否合并要先定义。放到“区间、排序与合并”这条主线里看，关键原则是：合并按起点排序，调度按终点排序，覆盖题常用起点升序且终点降序。扫描时维护当前最远右端、结果尾区间或已选区间终点。

区间题证明依赖排序后的邻接关系。若当前区间无法影响结果尾部或相邻区间，它就不会再影响更远的区间。本题的按开始时间扫描，用小顶堆维护当前会议的最早结束时间要说明哪些区间已经安全归档。这道题的边界不能留到最后补丁式处理，实验时覆盖会议端点相接、多个同一时间开始、空列表、堆顶释放条件，画出排序后的区间序列和当前结果尾部。动态版本要专门测前驱、后继和端点相接。

C++ 实现上，区间可用 \code{std::vector<std::array<int, 2>>} 或结构体表达端点，排序比较器要处理相同起点。动态区间需要 \code{std::map} 的前驱后继，而不是哈希表。如果追问变成“也可以把开始和结束分别排序，用双指针统计并发数”，先判断原来的状态是否仍然覆盖新答案集合，再决定是微调边界、改 value 语义，还是换成另一类结构。

\subsubsection{LeetCode 352 将数据流变为多个不相交区间}

先两两比较区间，确认相交、覆盖或冲突的定义。本题的模型是数字逐个到达，需要动态维护已覆盖的有序不相交区间；暴力慢在每个区间都回头看太多无关区间。优化点用有序表找到新值附近的前驱和后继，判断是否连接、合并或新建区间通常依赖排序后相邻关系，必须先定义端点闭开。

区间结构的状态是当前右端、结果尾区间或动态有序表的相邻关系。本题的用有序表找到新值附近的前驱和后继，判断是否连接、合并或新建区间只允许丢掉已经不可能再冲突的区间；端点相等是否合并要先定义。放到“区间、排序与合并”这条主线里看，关键原则是：合并按起点排序，调度按终点排序，覆盖题常用起点升序且终点降序。扫描时维护当前最远右端、结果尾区间或已选区间终点。

区间题证明依赖排序后的邻接关系。若当前区间无法影响结果尾部或相邻区间，它就不会再影响更远的区间。本题的用有序表找到新值附近的前驱和后继，判断是否连接、合并或新建区间要说明哪些区间已经安全归档。这道题的边界不能留到最后补丁式处理，实验时覆盖重复插入、连接左右两个区间、插入最小最大值、相邻端点，画出排序后的区间序列和当前结果尾部。动态版本要专门测前驱、后继和端点相接。

C++ 实现上，区间可用 \code{std::vector<std::array<int, 2>>} 或结构体表达端点，排序比较器要处理相同起点。动态区间需要 \code{std::map} 的前驱后继，而不是哈希表。如果追问变成“这类动态区间不能用普通数组扫描长期维护”，先判断原来的状态是否仍然覆盖新答案集合，再决定是微调边界、改 value 语义，还是换成另一类结构。

\subsubsection{LeetCode 436 寻找右区间}

先两两比较区间，确认相交、覆盖或冲突的定义。本题的模型是每个区间需要寻找起点不小于当前终点的最小起点区间；暴力慢在每个区间都回头看太多无关区间。优化点把所有起点和原下标排序，对每个终点做 lower bound通常依赖排序后相邻关系，必须先定义端点闭开。

区间结构的状态是当前右端、结果尾区间或动态有序表的相邻关系。本题的把所有起点和原下标排序，对每个终点做 lower bound只允许丢掉已经不可能再冲突的区间；端点相等是否合并要先定义。放到“区间、排序与合并”这条主线里看，关键原则是：合并按起点排序，调度按终点排序，覆盖题常用起点升序且终点降序。扫描时维护当前最远右端、结果尾区间或已选区间终点。

区间题证明依赖排序后的邻接关系。若当前区间无法影响结果尾部或相邻区间，它就不会再影响更远的区间。本题的把所有起点和原下标排序，对每个终点做 lower bound要说明哪些区间已经安全归档。这道题的边界不能留到最后补丁式处理，实验时覆盖不存在右区间、起点相同、负端点、单区间，画出排序后的区间序列和当前结果尾部。动态版本要专门测前驱、后继和端点相接。

C++ 实现上，区间可用 \code{std::vector<std::array<int, 2>>} 或结构体表达端点，排序比较器要处理相同起点。动态区间需要 \code{std::map} 的前驱后继，而不是哈希表。如果追问变成“如果区间动态插入，则需要有序表而不是一次排序数组”，先判断原来的状态是否仍然覆盖新答案集合，再决定是微调边界、改 value 语义，还是换成另一类结构。

\subsubsection{LeetCode 729 我的日程安排表 I}

先两两比较区间，确认相交、覆盖或冲突的定义。本题的模型是每次插入新区间时，需要判断它是否和已有区间重叠；暴力慢在每个区间都回头看太多无关区间。优化点有序表按起点存区间，只检查新区间的前驱和后继通常依赖排序后相邻关系，必须先定义端点闭开。

区间结构的状态是当前右端、结果尾区间或动态有序表的相邻关系。本题的有序表按起点存区间，只检查新区间的前驱和后继只允许丢掉已经不可能再冲突的区间；端点相等是否合并要先定义。放到“区间、排序与合并”这条主线里看，关键原则是：合并按起点排序，调度按终点排序，覆盖题常用起点升序且终点降序。扫描时维护当前最远右端、结果尾区间或已选区间终点。

区间题证明依赖排序后的邻接关系。若当前区间无法影响结果尾部或相邻区间，它就不会再影响更远的区间。本题的有序表按起点存区间，只检查新区间的前驱和后继要说明哪些区间已经安全归档。这道题的边界不能留到最后补丁式处理，实验时覆盖端点相接、插入最前最后、完全覆盖已有区间、动态多次插入，画出排序后的区间序列和当前结果尾部。动态版本要专门测前驱、后继和端点相接。

C++ 实现上，区间可用 \code{std::vector<std::array<int, 2>>} 或结构体表达端点，排序比较器要处理相同起点。动态区间需要 \code{std::map} 的前驱后继，而不是哈希表。如果追问变成“哈希表无法回答前驱后继，动态区间要用有序结构”，先判断原来的状态是否仍然覆盖新答案集合，再决定是微调边界、改 value 语义，还是换成另一类结构。

\subsubsection{LeetCode 731 我的日程安排表 II}

先两两比较区间，确认相交、覆盖或冲突的定义。本题的模型是允许双重预订但不允许三重预订，插入会改变区间重叠层数；暴力慢在每个区间都回头看太多无关区间。优化点保存已有预订和已经双重预订的区间；新预订若碰到双重区间则失败，否则记录新产生的重叠通常依赖排序后相邻关系，必须先定义端点闭开。

区间结构的状态是当前右端、结果尾区间或动态有序表的相邻关系。本题的保存已有预订和已经双重预订的区间；新预订若碰到双重区间则失败，否则记录新产生的重叠只允许丢掉已经不可能再冲突的区间；端点相等是否合并要先定义。放到“区间、排序与合并”这条主线里看，关键原则是：合并按起点排序，调度按终点排序，覆盖题常用起点升序且终点降序。扫描时维护当前最远右端、结果尾区间或已选区间终点。

区间题证明依赖排序后的邻接关系。若当前区间无法影响结果尾部或相邻区间，它就不会再影响更远的区间。本题的保存已有预订和已经双重预订的区间；新预订若碰到双重区间则失败，否则记录新产生的重叠要说明哪些区间已经安全归档。这道题的边界不能留到最后补丁式处理，实验时覆盖端点相接、多个局部重叠、完全覆盖、回滚失败插入，画出排序后的区间序列和当前结果尾部。动态版本要专门测前驱、后继和端点相接。

C++ 实现上，区间可用 \code{std::vector<std::array<int, 2>>} 或结构体表达端点，排序比较器要处理相同起点。动态区间需要 \code{std::map} 的前驱后继，而不是哈希表。如果追问变成“扫描线差分也能做，但要处理失败时恢复状态”，先判断原来的状态是否仍然覆盖新答案集合，再决定是微调边界、改 value 语义，还是换成另一类结构。

\subsubsection{LeetCode 986 区间列表的交集}

先两两比较区间，确认相交、覆盖或冲突的定义。本题的模型是两个有序且不重叠的区间列表，交集只可能来自当前两个区间；暴力慢在每个区间都回头看太多无关区间。优化点每次比较两个当前区间的重叠部分，然后推进右端较小的区间通常依赖排序后相邻关系，必须先定义端点闭开。

区间结构的状态是当前右端、结果尾区间或动态有序表的相邻关系。本题的每次比较两个当前区间的重叠部分，然后推进右端较小的区间只允许丢掉已经不可能再冲突的区间；端点相等是否合并要先定义。放到“区间、排序与合并”这条主线里看，关键原则是：合并按起点排序，调度按终点排序，覆盖题常用起点升序且终点降序。扫描时维护当前最远右端、结果尾区间或已选区间终点。

区间题证明依赖排序后的邻接关系。若当前区间无法影响结果尾部或相邻区间，它就不会再影响更远的区间。本题的每次比较两个当前区间的重叠部分，然后推进右端较小的区间要说明哪些区间已经安全归档。这道题的边界不能留到最后补丁式处理，实验时覆盖空列表、端点相接、一个区间覆盖多个区间、无交集，画出排序后的区间序列和当前结果尾部。动态版本要专门测前驱、后继和端点相接。

C++ 实现上，区间可用 \code{std::vector<std::array<int, 2>>} 或结构体表达端点，排序比较器要处理相同起点。动态区间需要 \code{std::map} 的前驱后继，而不是哈希表。如果追问变成“这是区间双指针，不需要把两个列表合并排序”，先判断原来的状态是否仍然覆盖新答案集合，再决定是微调边界、改 value 语义，还是换成另一类结构。

\subsubsection{LeetCode 1094 拼车}

先两两比较区间，确认相交、覆盖或冲突的定义。本题的模型是乘客上下车事件构成沿路线位置变化的载客量；暴力慢在每个区间都回头看太多无关区间。优化点用差分数组或扫描线记录每个位置人数变化，累计人数不能超过容量通常依赖排序后相邻关系，必须先定义端点闭开。

区间结构的状态是当前右端、结果尾区间或动态有序表的相邻关系。本题的用差分数组或扫描线记录每个位置人数变化，累计人数不能超过容量只允许丢掉已经不可能再冲突的区间；端点相等是否合并要先定义。放到“区间、排序与合并”这条主线里看，关键原则是：合并按起点排序，调度按终点排序，覆盖题常用起点升序且终点降序。扫描时维护当前最远右端、结果尾区间或已选区间终点。

区间题证明依赖排序后的邻接关系。若当前区间无法影响结果尾部或相邻区间，它就不会再影响更远的区间。本题的用差分数组或扫描线记录每个位置人数变化，累计人数不能超过容量要说明哪些区间已经安全归档。这道题的边界不能留到最后补丁式处理，实验时覆盖同站上下车、容量刚好、站点范围、输入未排序，画出排序后的区间序列和当前结果尾部。动态版本要专门测前驱、后继和端点相接。

C++ 实现上，区间可用 \code{std::vector<std::array<int, 2>>} 或结构体表达端点，排序比较器要处理相同起点。动态区间需要 \code{std::map} 的前驱后继，而不是哈希表。如果追问变成“航班预订统计也是差分思想，只是只需要最终每段累加结果”，先判断原来的状态是否仍然覆盖新答案集合，再决定是微调边界、改 value 语义，还是换成另一类结构。

\subsubsection{LeetCode 1109 航班预订统计}

先两两比较区间，确认相交、覆盖或冲突的定义。本题的模型是每条预订给一个连续航班区间增加座位数；暴力慢在每个区间都回头看太多无关区间。优化点差分数组在区间起点加人数，终点后一位减人数，最后前缀还原每航班总数通常依赖排序后相邻关系，必须先定义端点闭开。

区间结构的状态是当前右端、结果尾区间或动态有序表的相邻关系。本题的差分数组在区间起点加人数，终点后一位减人数，最后前缀还原每航班总数只允许丢掉已经不可能再冲突的区间；端点相等是否合并要先定义。放到“区间、排序与合并”这条主线里看，关键原则是：合并按起点排序，调度按终点排序，覆盖题常用起点升序且终点降序。扫描时维护当前最远右端、结果尾区间或已选区间终点。

区间题证明依赖排序后的邻接关系。若当前区间无法影响结果尾部或相邻区间，它就不会再影响更远的区间。本题的差分数组在区间起点加人数，终点后一位减人数，最后前缀还原每航班总数要说明哪些区间已经安全归档。这道题的边界不能留到最后补丁式处理，实验时覆盖区间到最后一个航班、单点区间、多条重叠预订、下标从一开始，画出排序后的区间序列和当前结果尾部。动态版本要专门测前驱、后继和端点相接。

C++ 实现上，区间可用 \code{std::vector<std::array<int, 2>>} 或结构体表达端点，排序比较器要处理相同起点。动态区间需要 \code{std::map} 的前驱后继，而不是哈希表。如果追问变成“这是静态区间加法，不需要线段树”，先判断原来的状态是否仍然覆盖新答案集合，再决定是微调边界、改 value 语义，还是换成另一类结构。

\subsubsection{LeetCode 1229 安排会议日程}

先两两比较区间，确认相交、覆盖或冲突的定义。本题的模型是两个人的可用时间段都有序，目标是找最早长度足够的交集；暴力慢在每个区间都回头看太多无关区间。优化点双指针比较当前两个区间交集，若长度不足就推进结束更早的一侧通常依赖排序后相邻关系，必须先定义端点闭开。

区间结构的状态是当前右端、结果尾区间或动态有序表的相邻关系。本题的双指针比较当前两个区间交集，若长度不足就推进结束更早的一侧只允许丢掉已经不可能再冲突的区间；端点相等是否合并要先定义。放到“区间、排序与合并”这条主线里看，关键原则是：合并按起点排序，调度按终点排序，覆盖题常用起点升序且终点降序。扫描时维护当前最远右端、结果尾区间或已选区间终点。

区间题证明依赖排序后的邻接关系。若当前区间无法影响结果尾部或相邻区间，它就不会再影响更远的区间。本题的双指针比较当前两个区间交集，若长度不足就推进结束更早的一侧要说明哪些区间已经安全归档。这道题的边界不能留到最后补丁式处理，实验时覆盖无可行时间、刚好等于 duration、区间端点相接、输入需要先排序，画出排序后的区间序列和当前结果尾部。动态版本要专门测前驱、后继和端点相接。

C++ 实现上，区间可用 \code{std::vector<std::array<int, 2>>} 或结构体表达端点，排序比较器要处理相同起点。动态区间需要 \code{std::map} 的前驱后继，而不是哈希表。如果追问变成“它和区间列表交集相同，但额外带最早可行和长度约束”，先判断原来的状态是否仍然覆盖新答案集合，再决定是微调边界、改 value 语义，还是换成另一类结构。

\subsubsection{LeetCode 1288 删除被覆盖区间}

先两两比较区间，确认相交、覆盖或冲突的定义。本题的模型是被覆盖要求另一区间起点不晚且终点不早；暴力慢在每个区间都回头看太多无关区间。优化点按起点升序、终点降序排序，扫描维护最大右端通常依赖排序后相邻关系，必须先定义端点闭开。

区间结构的状态是当前右端、结果尾区间或动态有序表的相邻关系。本题的按起点升序、终点降序排序，扫描维护最大右端只允许丢掉已经不可能再冲突的区间；端点相等是否合并要先定义。放到“区间、排序与合并”这条主线里看，关键原则是：合并按起点排序，调度按终点排序，覆盖题常用起点升序且终点降序。扫描时维护当前最远右端、结果尾区间或已选区间终点。

区间题证明依赖排序后的邻接关系。若当前区间无法影响结果尾部或相邻区间，它就不会再影响更远的区间。本题的按起点升序、终点降序排序，扫描维护最大右端要说明哪些区间已经安全归档。这道题的边界不能留到最后补丁式处理，实验时覆盖同起点区间、完全覆盖、无覆盖、端点相等，画出排序后的区间序列和当前结果尾部。动态版本要专门测前驱、后继和端点相接。

C++ 实现上，区间可用 \code{std::vector<std::array<int, 2>>} 或结构体表达端点，排序比较器要处理相同起点。动态区间需要 \code{std::map} 的前驱后继，而不是哈希表。如果追问变成“终点降序是关键，否则同起点小区间会先出现造成误判”，先判断原来的状态是否仍然覆盖新答案集合，再决定是微调边界、改 value 语义，还是换成另一类结构。

\subsection{位运算与状态压缩工作坊}

先用计数或哈希保证正确，再寻找位运算的代数抵消、按位统计或函数图结构。位题的难点是证明被抵消的信息确实不影响答案。

\subsubsection{LeetCode 29 两数相除}

先用计数、集合或普通 DP 写出清楚版本。本题的模型是除法可以看成不断减去除数的倍增值；位运算只是把状态压进二进制表示。优化点把除数按二进制倍增，优先减去不超过被除数的最大倍数并累加商必须说明每一位或每次异或到底保留了什么信息。

位状态要给每一位固定含义。本题的把除数按二进制倍增，优先减去不超过被除数的最大倍数并累加商如果依赖异或抵消，就要说明成对元素为什么完全消失；如果依赖掩码，就要说明同一位置不同掩码为什么不能合并。放到“位运算与状态压缩”这条主线里看，关键原则是：异或解决成对抵消，按位计数处理出现三次，掩码表示访问集合或可用集合，低位提取用于枚举候选。

位运算证明要逐位说明。异或、按位计数或掩码转移都不能只靠直觉；本题的把除数按二进制倍增，优先减去不超过被除数的最大倍数并累加商必须保证被压掉的信息不会影响除法可以看成不断减去除数的倍增值要求的答案。这道题的边界不能留到最后补丁式处理，实验时覆盖除数为一、被除数为 int 最小值、结果溢出、符号处理，把数字写成二进制逐位检查。负数、最高位、空掩码和全集掩码都要单独测。

C++ 实现上，移位时优先使用无符号或足够宽的整数，位数超过 31 时考虑 \code{long long} 或 \code{std::bitset}。位运算表达式加括号，避免优先级误读。如果追问变成“它和快速幂类似，都利用二进制拆分减少重复加减”，先判断原来的状态是否仍然覆盖新答案集合，再决定是微调边界、改 value 语义，还是换成另一类结构。

\subsubsection{LeetCode 136 只出现一次的数字}

先用计数、集合或普通 DP 写出清楚版本。本题的模型是成对元素可以用异或抵消，剩下单个元素；位运算只是把状态压进二进制表示。优化点异或满足交换结合，扫描顺序不影响答案必须说明每一位或每次异或到底保留了什么信息。

位状态要给每一位固定含义。本题的异或满足交换结合，扫描顺序不影响答案如果依赖异或抵消，就要说明成对元素为什么完全消失；如果依赖掩码，就要说明同一位置不同掩码为什么不能合并。放到“位运算与状态压缩”这条主线里看，关键原则是：异或解决成对抵消，按位计数处理出现三次，掩码表示访问集合或可用集合，低位提取用于枚举候选。

位运算证明要逐位说明。异或、按位计数或掩码转移都不能只靠直觉；本题的异或满足交换结合，扫描顺序不影响答案必须保证被压掉的信息不会影响成对元素可以用异或抵消，剩下单个元素要求的答案。这道题的边界不能留到最后补丁式处理，实验时覆盖零、负数、唯一值在开头或结尾、所有其他值恰好两次，把数字写成二进制逐位检查。负数、最高位、空掩码和全集掩码都要单独测。

C++ 实现上，移位时优先使用无符号或足够宽的整数，位数超过 31 时考虑 \code{long long} 或 \code{std::bitset}。位运算表达式加括号，避免优先级误读。如果追问变成“若其他值出现三次，要改为按位计数取模”，先判断原来的状态是否仍然覆盖新答案集合，再决定是微调边界、改 value 语义，还是换成另一类结构。

\subsubsection{LeetCode 137 只出现一次的数字 II}

先用计数、集合或普通 DP 写出清楚版本。本题的模型是除一个数字外，其余数字都出现三次，单纯异或不能抵消；位运算只是把状态压进二进制表示。优化点逐位统计一的数量，对三取模后还原目标数字必须说明每一位或每次异或到底保留了什么信息。

位状态要给每一位固定含义。本题的逐位统计一的数量，对三取模后还原目标数字如果依赖异或抵消，就要说明成对元素为什么完全消失；如果依赖掩码，就要说明同一位置不同掩码为什么不能合并。放到“位运算与状态压缩”这条主线里看，关键原则是：异或解决成对抵消，按位计数处理出现三次，掩码表示访问集合或可用集合，低位提取用于枚举候选。

位运算证明要逐位说明。异或、按位计数或掩码转移都不能只靠直觉；本题的逐位统计一的数量，对三取模后还原目标数字必须保证被压掉的信息不会影响除一个数字外，其余数字都出现三次，单纯异或不能抵消要求的答案。这道题的边界不能留到最后补丁式处理，实验时覆盖负数、符号位、目标为零、所有其他数恰好三次，把数字写成二进制逐位检查。负数、最高位、空掩码和全集掩码都要单独测。

C++ 实现上，移位时优先使用无符号或足够宽的整数，位数超过 31 时考虑 \code{long long} 或 \code{std::bitset}。位运算表达式加括号，避免优先级误读。如果追问变成“若其他数出现 k 次，思路仍是按位计数取模”，先判断原来的状态是否仍然覆盖新答案集合，再决定是微调边界、改 value 语义，还是换成另一类结构。

\subsubsection{LeetCode 190 颠倒二进制位}

先用计数、集合或普通 DP 写出清楚版本。本题的模型是每一位的新位置由原位置镜像决定；位运算只是把状态压进二进制表示。优化点循环三十二次，把结果左移并接上当前最低位，再右移输入必须说明每一位或每次异或到底保留了什么信息。

位状态要给每一位固定含义。本题的循环三十二次，把结果左移并接上当前最低位，再右移输入如果依赖异或抵消，就要说明成对元素为什么完全消失；如果依赖掩码，就要说明同一位置不同掩码为什么不能合并。放到“位运算与状态压缩”这条主线里看，关键原则是：异或解决成对抵消，按位计数处理出现三次，掩码表示访问集合或可用集合，低位提取用于枚举候选。

位运算证明要逐位说明。异或、按位计数或掩码转移都不能只靠直觉；本题的循环三十二次，把结果左移并接上当前最低位，再右移输入必须保证被压掉的信息不会影响每一位的新位置由原位置镜像决定要求的答案。这道题的边界不能留到最后补丁式处理，实验时覆盖最高位、最低位、全零、全一、无符号语义，把数字写成二进制逐位检查。负数、最高位、空掩码和全集掩码都要单独测。

C++ 实现上，移位时优先使用无符号或足够宽的整数，位数超过 31 时考虑 \code{long long} 或 \code{std::bitset}。位运算表达式加括号，避免优先级误读。如果追问变成“若多次调用，可以按字节预处理反转表”，先判断原来的状态是否仍然覆盖新答案集合，再决定是微调边界、改 value 语义，还是换成另一类结构。

\subsubsection{LeetCode 191 位 1 的个数}

先用计数、集合或普通 DP 写出清楚版本。本题的模型是每次清掉最低位的一可以让循环次数等于一的个数；位运算只是把状态压进二进制表示。优化点使用 n \& (n - 1) 删除最低位一，比逐位检查更直接必须说明每一位或每次异或到底保留了什么信息。

位状态要给每一位固定含义。本题的使用 n \& (n - 1) 删除最低位一，比逐位检查更直接如果依赖异或抵消，就要说明成对元素为什么完全消失；如果依赖掩码，就要说明同一位置不同掩码为什么不能合并。放到“位运算与状态压缩”这条主线里看，关键原则是：异或解决成对抵消，按位计数处理出现三次，掩码表示访问集合或可用集合，低位提取用于枚举候选。

位运算证明要逐位说明。异或、按位计数或掩码转移都不能只靠直觉；本题的使用 n \& (n - 1) 删除最低位一，比逐位检查更直接必须保证被压掉的信息不会影响每次清掉最低位的一可以让循环次数等于一的个数要求的答案。这道题的边界不能留到最后补丁式处理，实验时覆盖零、只有最高位、一的个数很多、无符号输入，把数字写成二进制逐位检查。负数、最高位、空掩码和全集掩码都要单独测。

C++ 实现上，移位时优先使用无符号或足够宽的整数，位数超过 31 时考虑 \code{long long} 或 \code{std::bitset}。位运算表达式加括号，避免优先级误读。如果追问变成“比特位计数可以把这个过程 DP 化到一整段数”，先判断原来的状态是否仍然覆盖新答案集合，再决定是微调边界、改 value 语义，还是换成另一类结构。

\subsubsection{LeetCode 201 数字范围按位与}

先用计数、集合或普通 DP 写出清楚版本。本题的模型是区间内所有数按位与后，只保留左右端公共二进制前缀；位运算只是把状态压进二进制表示。优化点不断右移 left 和 right，直到二者相等，再把公共前缀移回原位必须说明每一位或每次异或到底保留了什么信息。

位状态要给每一位固定含义。本题的不断右移 left 和 right，直到二者相等，再把公共前缀移回原位如果依赖异或抵消，就要说明成对元素为什么完全消失；如果依赖掩码，就要说明同一位置不同掩码为什么不能合并。放到“位运算与状态压缩”这条主线里看，关键原则是：异或解决成对抵消，按位计数处理出现三次，掩码表示访问集合或可用集合，低位提取用于枚举候选。

位运算证明要逐位说明。异或、按位计数或掩码转移都不能只靠直觉；本题的不断右移 left 和 right，直到二者相等，再把公共前缀移回原位必须保证被压掉的信息不会影响区间内所有数按位与后，只保留左右端公共二进制前缀要求的答案。这道题的边界不能留到最后补丁式处理，实验时覆盖left 等于 right、区间跨过二的幂边界、结果为零，把数字写成二进制逐位检查。负数、最高位、空掩码和全集掩码都要单独测。

C++ 实现上，移位时优先使用无符号或足够宽的整数，位数超过 31 时考虑 \code{long long} 或 \code{std::bitset}。位运算表达式加括号，避免优先级误读。如果追问变成“也可以不断清除 right 的最低位一，直到 right 不大于 left”，先判断原来的状态是否仍然覆盖新答案集合，再决定是微调边界、改 value 语义，还是换成另一类结构。

\subsubsection{LeetCode 231 2 的幂}

先用计数、集合或普通 DP 写出清楚版本。本题的模型是二的幂在二进制中只有一个一；位运算只是把状态压进二进制表示。优化点正数 n 满足 n 与 n 减一按位与为零必须说明每一位或每次异或到底保留了什么信息。

位状态要给每一位固定含义。本题的正数 n 满足 n 与 n 减一按位与为零如果依赖异或抵消，就要说明成对元素为什么完全消失；如果依赖掩码，就要说明同一位置不同掩码为什么不能合并。放到“位运算与状态压缩”这条主线里看，关键原则是：异或解决成对抵消，按位计数处理出现三次，掩码表示访问集合或可用集合，低位提取用于枚举候选。

位运算证明要逐位说明。异或、按位计数或掩码转移都不能只靠直觉；本题的正数 n 满足 n 与 n 减一按位与为零必须保证被压掉的信息不会影响二的幂在二进制中只有一个一要求的答案。这道题的边界不能留到最后补丁式处理，实验时覆盖n 为零、负数、int 最小值、最高位，把数字写成二进制逐位检查。负数、最高位、空掩码和全集掩码都要单独测。

C++ 实现上，移位时优先使用无符号或足够宽的整数，位数超过 31 时考虑 \code{long long} 或 \code{std::bitset}。位运算表达式加括号，避免优先级误读。如果追问变成“判断四的幂还要额外限制一所在的位置为偶数位”，先判断原来的状态是否仍然覆盖新答案集合，再决定是微调边界、改 value 语义，还是换成另一类结构。

\subsubsection{LeetCode 260 只出现一次的数字 III}

先用计数、集合或普通 DP 写出清楚版本。本题的模型是两个只出现一次的数异或后至少有一位不同；位运算只是把状态压进二进制表示。优化点用最低不同位把数组分成两组，每组内部再异或抵消必须说明每一位或每次异或到底保留了什么信息。

位状态要给每一位固定含义。本题的用最低不同位把数组分成两组，每组内部再异或抵消如果依赖异或抵消，就要说明成对元素为什么完全消失；如果依赖掩码，就要说明同一位置不同掩码为什么不能合并。放到“位运算与状态压缩”这条主线里看，关键原则是：异或解决成对抵消，按位计数处理出现三次，掩码表示访问集合或可用集合，低位提取用于枚举候选。

位运算证明要逐位说明。异或、按位计数或掩码转移都不能只靠直觉；本题的用最低不同位把数组分成两组，每组内部再异或抵消必须保证被压掉的信息不会影响两个只出现一次的数异或后至少有一位不同要求的答案。这道题的边界不能留到最后补丁式处理，实验时覆盖两个答案一正一负、最低位分组、其他数恰好两次，把数字写成二进制逐位检查。负数、最高位、空掩码和全集掩码都要单独测。

C++ 实现上，移位时优先使用无符号或足够宽的整数，位数超过 31 时考虑 \code{long long} 或 \code{std::bitset}。位运算表达式加括号，避免优先级误读。如果追问变成“它是在异或抵消基础上增加分组位”，先判断原来的状态是否仍然覆盖新答案集合，再决定是微调边界、改 value 语义，还是换成另一类结构。

\subsubsection{LeetCode 268 丢失的数字}

先用计数、集合或普通 DP 写出清楚版本。本题的模型是下标零到 n 与数组值相比，缺失值会在异或抵消后剩下；位运算只是把状态压进二进制表示。优化点把所有下标和所有值一起异或，重复出现的数抵消必须说明每一位或每次异或到底保留了什么信息。

位状态要给每一位固定含义。本题的把所有下标和所有值一起异或，重复出现的数抵消如果依赖异或抵消，就要说明成对元素为什么完全消失；如果依赖掩码，就要说明同一位置不同掩码为什么不能合并。放到“位运算与状态压缩”这条主线里看，关键原则是：异或解决成对抵消，按位计数处理出现三次，掩码表示访问集合或可用集合，低位提取用于枚举候选。

位运算证明要逐位说明。异或、按位计数或掩码转移都不能只靠直觉；本题的把所有下标和所有值一起异或，重复出现的数抵消必须保证被压掉的信息不会影响下标零到 n 与数组值相比，缺失值会在异或抵消后剩下要求的答案。这道题的边界不能留到最后补丁式处理，实验时覆盖缺失零、缺失 n、数组长度一、也可用求和但要防溢出，把数字写成二进制逐位检查。负数、最高位、空掩码和全集掩码都要单独测。

C++ 实现上，移位时优先使用无符号或足够宽的整数，位数超过 31 时考虑 \code{long long} 或 \code{std::bitset}。位运算表达式加括号，避免优先级误读。如果追问变成“异或法适合展示值域完整且只有一个缺口的结构”，先判断原来的状态是否仍然覆盖新答案集合，再决定是微调边界、改 value 语义，还是换成另一类结构。

\subsubsection{LeetCode 287 寻找重复数}

先用计数、集合或普通 DP 写出清楚版本。本题的模型是数组值域使下标到值形成函数图，重复数是环入口；位运算只是把状态压进二进制表示。优化点快慢指针不修改数组且常数空间；值域二分用抽屉原理必须说明每一位或每次异或到底保留了什么信息。

位状态要给每一位固定含义。本题的快慢指针不修改数组且常数空间；值域二分用抽屉原理如果依赖异或抵消，就要说明成对元素为什么完全消失；如果依赖掩码，就要说明同一位置不同掩码为什么不能合并。放到“位运算与状态压缩”这条主线里看，关键原则是：异或解决成对抵消，按位计数处理出现三次，掩码表示访问集合或可用集合，低位提取用于枚举候选。

位运算证明要逐位说明。异或、按位计数或掩码转移都不能只靠直觉；本题的快慢指针不修改数组且常数空间；值域二分用抽屉原理必须保证被压掉的信息不会影响数组值域使下标到值形成函数图，重复数是环入口要求的答案。这道题的边界不能留到最后补丁式处理，实验时覆盖重复数出现多次、不能排序、不能改数组、长度为 n 加一，把数字写成二进制逐位检查。负数、最高位、空掩码和全集掩码都要单独测。

C++ 实现上，移位时优先使用无符号或足够宽的整数，位数超过 31 时考虑 \code{long long} 或 \code{std::bitset}。位运算表达式加括号，避免优先级误读。如果追问变成“快慢指针要求值都能作为合法下标”，先判断原来的状态是否仍然覆盖新答案集合，再决定是微调边界、改 value 语义，还是换成另一类结构。

\subsubsection{LeetCode 318 最大单词长度乘积}

先用计数、集合或普通 DP 写出清楚版本。本题的模型是两个单词没有公共字母时才可相乘，字母集合可压成位掩码；位运算只是把状态压进二进制表示。优化点每个单词生成 26 位掩码，两个掩码按位与为零表示无交集必须说明每一位或每次异或到底保留了什么信息。

位状态要给每一位固定含义。本题的每个单词生成 26 位掩码，两个掩码按位与为零表示无交集如果依赖异或抵消，就要说明成对元素为什么完全消失；如果依赖掩码，就要说明同一位置不同掩码为什么不能合并。放到“位运算与状态压缩”这条主线里看，关键原则是：异或解决成对抵消，按位计数处理出现三次，掩码表示访问集合或可用集合，低位提取用于枚举候选。

位运算证明要逐位说明。异或、按位计数或掩码转移都不能只靠直觉；本题的每个单词生成 26 位掩码，两个掩码按位与为零表示无交集必须保证被压掉的信息不会影响两个单词没有公共字母时才可相乘，字母集合可压成位掩码要求的答案。这道题的边界不能留到最后补丁式处理，实验时覆盖重复字母、空字符串、多个单词同掩码、乘积范围，把数字写成二进制逐位检查。负数、最高位、空掩码和全集掩码都要单独测。

C++ 实现上，移位时优先使用无符号或足够宽的整数，位数超过 31 时考虑 \code{long long} 或 \code{std::bitset}。位运算表达式加括号，避免优先级误读。如果追问变成“若字符集大于整数位数，需要 bitset 或哈希集合”，先判断原来的状态是否仍然覆盖新答案集合，再决定是微调边界、改 value 语义，还是换成另一类结构。

\subsubsection{LeetCode 338 比特位计数}

先用计数、集合或普通 DP 写出清楚版本。本题的模型是每个数的一的个数可由去掉最低位一后的较小数转移；位运算只是把状态压进二进制表示。优化点dp[x] = dp[x \& (x - 1)] + 1，或 dp[x] = dp[x >> 1] + (x \& 1)必须说明每一位或每次异或到底保留了什么信息。

位状态要给每一位固定含义。本题的dp[x] = dp[x \& (x - 1)] + 1，或 dp[x] = dp[x >> 1] + (x \& 1)如果依赖异或抵消，就要说明成对元素为什么完全消失；如果依赖掩码，就要说明同一位置不同掩码为什么不能合并。放到“位运算与状态压缩”这条主线里看，关键原则是：异或解决成对抵消，按位计数处理出现三次，掩码表示访问集合或可用集合，低位提取用于枚举候选。

位运算证明要逐位说明。异或、按位计数或掩码转移都不能只靠直觉；本题的dp[x] = dp[x \& (x - 1)] + 1，或 dp[x] = dp[x >> 1] + (x \& 1)必须保证被压掉的信息不会影响每个数的一的个数可由去掉最低位一后的较小数转移要求的答案。这道题的边界不能留到最后补丁式处理，实验时覆盖n 为零、二的幂、连续区间、表达式优先级，把数字写成二进制逐位检查。负数、最高位、空掩码和全集掩码都要单独测。

C++ 实现上，移位时优先使用无符号或足够宽的整数，位数超过 31 时考虑 \code{long long} 或 \code{std::bitset}。位运算表达式加括号，避免优先级误读。如果追问变成“它展示位运算和 DP 的结合，而不是逐个数循环数位”，先判断原来的状态是否仍然覆盖新答案集合，再决定是微调边界、改 value 语义，还是换成另一类结构。

\subsubsection{LeetCode 371 两整数之和}

先用计数、集合或普通 DP 写出清楚版本。本题的模型是不用加号时，加法可以拆成无进位和与进位；位运算只是把状态压进二进制表示。优化点异或得到无进位和，与运算后左移得到进位，迭代直到进位为零必须说明每一位或每次异或到底保留了什么信息。

位状态要给每一位固定含义。本题的异或得到无进位和，与运算后左移得到进位，迭代直到进位为零如果依赖异或抵消，就要说明成对元素为什么完全消失；如果依赖掩码，就要说明同一位置不同掩码为什么不能合并。放到“位运算与状态压缩”这条主线里看，关键原则是：异或解决成对抵消，按位计数处理出现三次，掩码表示访问集合或可用集合，低位提取用于枚举候选。

位运算证明要逐位说明。异或、按位计数或掩码转移都不能只靠直觉；本题的异或得到无进位和，与运算后左移得到进位，迭代直到进位为零必须保证被压掉的信息不会影响不用加号时，加法可以拆成无进位和与进位要求的答案。这道题的边界不能留到最后补丁式处理，实验时覆盖正负数、进位传播、有符号溢出、需要使用无符号中间值，把数字写成二进制逐位检查。负数、最高位、空掩码和全集掩码都要单独测。

C++ 实现上，移位时优先使用无符号或足够宽的整数，位数超过 31 时考虑 \code{long long} 或 \code{std::bitset}。位运算表达式加括号，避免优先级误读。如果追问变成“这是位级算术模拟，和异或抵消类题目的语义不同”，先判断原来的状态是否仍然覆盖新答案集合，再决定是微调边界、改 value 语义，还是换成另一类结构。

\subsubsection{LeetCode 393 UTF-8 编码验证}

先用计数、集合或普通 DP 写出清楚版本。本题的模型是每个字节的高位模式决定一个字符需要的后续字节数量；位运算只是把状态压进二进制表示。优化点用位掩码判断首字节类型，并维护还需要多少个 continuation 字节必须说明每一位或每次异或到底保留了什么信息。

位状态要给每一位固定含义。本题的用位掩码判断首字节类型，并维护还需要多少个 continuation 字节如果依赖异或抵消，就要说明成对元素为什么完全消失；如果依赖掩码，就要说明同一位置不同掩码为什么不能合并。放到“位运算与状态压缩”这条主线里看，关键原则是：异或解决成对抵消，按位计数处理出现三次，掩码表示访问集合或可用集合，低位提取用于枚举候选。

位运算证明要逐位说明。异或、按位计数或掩码转移都不能只靠直觉；本题的用位掩码判断首字节类型，并维护还需要多少个 continuation 字节必须保证被压掉的信息不会影响每个字节的高位模式决定一个字符需要的后续字节数量要求的答案。这道题的边界不能留到最后补丁式处理，实验时覆盖非法首字节、后续字节不足、后续字节格式错误、单字节字符，把数字写成二进制逐位检查。负数、最高位、空掩码和全集掩码都要单独测。

C++ 实现上，移位时优先使用无符号或足够宽的整数，位数超过 31 时考虑 \code{long long} 或 \code{std::bitset}。位运算表达式加括号，避免优先级误读。如果追问变成“这是位模式解析题，重点是协议规则和状态机”，先判断原来的状态是否仍然覆盖新答案集合，再决定是微调边界、改 value 语义，还是换成另一类结构。

\subsubsection{LeetCode 421 数组中两个数的最大异或值}

先用计数、集合或普通 DP 写出清楚版本。本题的模型是最大异或值由高位到低位尽量取一决定；位运算只是把状态压进二进制表示。优化点逐位尝试把当前前缀答案设为一，用哈希集合验证是否存在两个前缀配对必须说明每一位或每次异或到底保留了什么信息。

位状态要给每一位固定含义。本题的逐位尝试把当前前缀答案设为一，用哈希集合验证是否存在两个前缀配对如果依赖异或抵消，就要说明成对元素为什么完全消失；如果依赖掩码，就要说明同一位置不同掩码为什么不能合并。放到“位运算与状态压缩”这条主线里看，关键原则是：异或解决成对抵消，按位计数处理出现三次，掩码表示访问集合或可用集合，低位提取用于枚举候选。

位运算证明要逐位说明。异或、按位计数或掩码转移都不能只靠直觉；本题的逐位尝试把当前前缀答案设为一，用哈希集合验证是否存在两个前缀配对必须保证被压掉的信息不会影响最大异或值由高位到低位尽量取一决定要求的答案。这道题的边界不能留到最后补丁式处理，实验时覆盖零、重复数、最高位、负数语义，把数字写成二进制逐位检查。负数、最高位、空掩码和全集掩码都要单独测。

C++ 实现上，移位时优先使用无符号或足够宽的整数，位数超过 31 时考虑 \code{long long} 或 \code{std::bitset}。位运算表达式加括号，避免优先级误读。如果追问变成“也可以建二进制 Trie，每个数尽量走相反位”，先判断原来的状态是否仍然覆盖新答案集合，再决定是微调边界、改 value 语义，还是换成另一类结构。

\begin{principlebox}{工作坊使用方式}
不要一次性把本节当题解背完。建议每次选一个题型，先遮住优化落点，只看题面模型和暴力基线，自己写出慢在哪里；再看结构观察和正确性；最后用边界实验做对拍。能完成这个过程，才说明题目进入了自己的知识体系。
\end{principlebox}

\section{数据结构变种、工程应用和实验专题}

这一组专题补齐数据结构的变种、工程应用和实验要求。它不是题号清单，而是把同一类结构放到刷题、C++ 容器和系统源码之间比较，帮助读者理解为什么同一个复杂度在不同场景下表现差异很大。

\subsection{数组、字符串和连续内存的变种}

数组的优势不是语法简单，而是连续内存、随机访问和下标可计算。刷题里常见的前缀和、差分、双指针、二分、原地交换，都依赖下标能在常数时间跳转。工程里数组还带来缓存局部性：顺序扫描通常比链表快很多，即使二者同为线性复杂度。学习数组题时，要把每个下标变量写成区间语义，例如左闭右开窗口、已处理前缀、未知区、已确认后缀。只要区间语义清楚，代码中的加一减一就有根据。

数组变种很多。前缀和把区间求和变成两个端点相减，适合静态区间查询；差分把区间加法变成两个端点标记，适合批量更新后一次还原；二维前缀和把矩形查询变成四个角组合；二维差分适合批量矩形更新；坐标压缩把大值域映射到紧凑下标，适合树状数组、线段树和计数数组。每个变种都要问一个问题：它把哪类重复工作提前保存了，代价是多了什么空间和初始化步骤。

字符串本质上也是字符数组，但它多了编码、复制和匹配成本。LeetCode 基础字符串题常假设小写字母，此时 \code{std::array<int, 26>} 比哈希表更直接；若扩展到 ASCII，可以用 128 或 256 长度数组；若扩展到 Unicode，就不能简单按字节计数。字符串切片也要注意成本，\code{substr} 可能创建新字符串，热循环里更适合用下标区间、\code{string_view} 或 Trie。面试时如果能主动说明字符集假设，答案会更稳。

数组题的实验很直接。第一，给一个窗口题打印 \code{left}、\code{right}、窗口内容和计数表，观察每次移动只改变一个元素。第二，把二维矩阵随机生成后，用暴力矩形求和和二维前缀和对拍，确认四角公式。第三，把 \code{vector} 和 \code{list} 各存一百万个整数，分别顺序求和，比较时间，理解缓存局部性。第四，故意保存 \code{vector} 元素地址后继续 \code{push_back}，观察扩容后旧地址失效。

\subsection{链表、侵入式节点和稳定身份}

链表题的难点不是没有下标，而是节点身份和连接关系分离。数组元素的位置由下标决定，移动元素会改变位置；链表节点的位置由前后指针决定，重连节点不会改变节点自身地址。这个差异让链表适合 LRU、任务队列、空闲块链、内核对象列表和需要稳定迭代器的场景。刷题里的 \code{ListNode} 是最简模型，真实工程里一个业务对象可能带多个链表节点，分别挂进不同索引。

链表有多种变种。单链表空间少，适合反转、合并、快慢指针和环检测；双链表能从已知节点 O(1) 删除，适合 LRU 和双端队列；循环链表适合轮转和哨兵化边界；跳表在多层链表上提供期望对数查找；侵入式链表把链接字段放进业务对象，减少额外分配。每个变种都在回答访问模式：是否需要前驱，是否需要随机查找，是否需要稳定节点，是否会频繁移动到头尾。

链表面试要能讲副作用。反转链表会改变输入结构，回文链表若反转后半段，工程里最好恢复；合并链表可以复用节点，也可以新建节点，二者所有权不同；删除节点时如果平台管理内存，不应该随意释放输入节点；相交链表比较的是地址不是值；环形链表的快慢指针证明来自相对速度，而不是经验。只要这些边界说不清，代码短也不代表理解深。

链表实验建议做三组。第一，把反转链表每一步的 \code{prev}、\code{cur}、\code{next} 打印成边关系，确认没有丢失剩余链。第二，写 LRU 的 \code{get} 后移动到头部，统计哈希查询和链表移动次数。第三，写一个侵入式节点对象，同时挂入按时间排序链表和按 key 哈希桶，体会同一对象进入多个索引时删除顺序有多重要。

\subsection{栈、队列、双端队列和单调结构}

栈的核心是最近未完成结构。括号匹配保存最近未闭合左括号，表达式求值保存上一层上下文，单调栈保存还在等待未来元素解决的下标。不要把“用栈”当结论，必须说清栈顶代表什么，什么事件触发入栈，什么事件触发出栈，出栈时答案是否已经确定。柱状图最大矩形、每日温度、下一个更大元素、最长有效括号虽然都用栈，但弹出的语义完全不同。

队列的核心是按到达顺序处理。BFS 队列保证距离层次，任务队列保证提交顺序，滑动窗口队列保存候选下标。双端队列进一步允许两端进出，适合 0-1 BFS 和单调队列。单调队列不是普通队列加排序，而是在队尾删除被新元素支配的候选，在队首删除过期候选。每个元素最多进出一次，所以线性复杂度来自摊还分析。

单调结构最容易被写成模板。真正要问的是支配关系：一个旧候选为什么被新候选删除后永远不可能成为答案？滑动窗口最大值里，新值更大且更晚，旧小值在未来窗口中既不更大也不更持久；受限子序列和里，新 DP 值更大且下标更晚，也支配旧候选；柱状图中，遇到更矮柱时，被弹出柱子的右边界已经确定。支配关系说不清，单调结构就不可靠。

实验可以把队列状态打印出来。对滑动窗口最大值，记录每轮右端进入、队尾弹出、队首过期和答案读取。对 0-1 BFS，构造零权边和一权边混合的小图，观察零权边入队首如何保证距离顺序。对表达式求值，打印栈中数字和运算符，观察乘除为什么要立即处理，加减为什么可以延迟。

\subsection{哈希、前缀状态和等价类}

哈希表的本质是把对象归入等价类。两数之和把历史值映射到下标，字母异位词把字符串映射到规范签名，前缀和把区间条件映射成历史前缀，账户合并把邮箱映射到账户集合。哈希不是魔法，它只把等值查询变便宜；如果题目需要顺序、前驱、后继或范围最值，哈希表就不是合适结构。

前缀哈希题要先写代数关系。和为 K 的子数组来自 \code{prefix[j] - prefix[i] == k}；连续子数组和来自两个前缀同余；连续数组来自把 0 当成 -1 后找相同前缀；优美子数组来自奇数计数差。关系不同，哈希表保存内容也不同：计数题保存频次，最长题保存最早下标，判定题保存是否出现，最短题可能需要单调队列而不是普通哈希。

哈希 key 的设计决定正确性。异位词可以排序字符串做 key，也可以用字符计数做 key；计数 key 必须带分隔符，否则 \code{1,11} 和 \code{11,1} 可能混淆；结构化 key 可以用 pair 或自定义结构，但要提供哈希函数和相等比较；浮点数不适合直接做哈希 key，工程里要转成可比较的离散表示。面试中讲 key 比讲容器名字更重要。

实验建议做哈希对拍。第一，写和为 K 的子数组暴力版和哈希版，用包含负数、零和重复前缀的随机数组对拍。第二，把先查询后更新改成先更新后查询，观察 target 为 0 时如何出错。第三，给异位词分组设计没有分隔符的计数 key，构造冲突案例。第四，对 \code{unordered_map} 调用和不调用 \code{reserve}，比较大量插入时 rehash 次数。

\subsection{堆、优先队列和动态候选集}

堆解决的是动态候选集的极值查询。排序能给全局顺序，但每次插入删除后重新排序成本高；堆只维护足够的信息让堆顶是最大或最小。前 K 高频、数据流中位数、合并 K 个链表、任务调度、最小加油次数，都是把“下一步选谁”变成堆顶查询。要注意堆不能回答任意元素删除，也不能直接给有序遍历。

双堆是数据流中位数的核心。大顶堆保存较小一半，小顶堆保存较大一半，大小差不超过一，且大顶堆堆顶不大于小顶堆堆顶。两个不变量缺一不可。若要支持滑动窗口删除，就必须加延迟删除表，删除请求先记录，等元素到堆顶时再真正弹出。这和 Linux RCU 的延迟释放不是同一机制，但都体现“逻辑删除”和“物理清理”分离。

堆题要警惕比较器。C++ \code{priority_queue} 默认大顶堆，小顶堆要用 \code{greater} 或自定义比较器。结构体比较器应该明确谁的优先级低，而不是凭直觉写小于号。若堆中保存 pair，默认比较会先比第一维再比第二维；如果业务只看频次或距离，就要写清楚 tie-breaker。比较器一旦不满足严格弱序，结果可能不稳定。

堆实验可以从三个版本比较。第 K 大元素用排序、大小为 K 的小顶堆和快速选择分别实现，比较复杂度和适用场景。合并 K 个有序链表用每轮扫描 K 个头、堆和分治三种方式实现，观察 N 和 K 对性能的影响。任务调度器同时写模拟堆版本和最高频公式版本，解释公式为什么不需要逐步模拟。

\subsection{树、Trie、平衡树和区间树视角}

树题的关键是状态方向。自顶向下问题把父节点约束传给子节点，例如验证 BST 的上下界；自底向上问题让子树向父节点汇报信息，例如平衡二叉树、最大路径和、直径；按层问题用队列维护层边界，例如层序遍历和最小深度。把方向说清楚，递归或迭代只是实现选择。

BST 的价值是有序性。中序遍历得到升序序列，查找可以按大小排除半棵树，前驱后继可以沿树结构寻找。普通二叉树没有这些性质，不能把 BST 的值域剪枝搬过去。红黑树、AVL、Treap、跳表都是为了在动态更新时保持近似平衡；LeetCode 大多不要求实现平衡树，但 \code{map}、\code{set} 和 Linux rbtree 背后都依赖这一类思想。

Trie 把字符串集合按前缀共享。实现 Trie、单词搜索 II、搜索建议系统、异或最大值位 Trie 都是这个思想。数组子节点适合小固定字符集，哈希子节点适合稀疏大字符集；单词结束标记不能省，因为前缀存在不等于单词存在。Trie 的空间可能很大，工程里会考虑压缩 Trie、Radix Tree 或 DAWG 等变体。

树实验要同时做递归和迭代。用显式栈写中序遍历，比较它和递归版本的状态；用队列写层序遍历，打印每层 size；用 Trie 插入一批单词，统计节点数和共享前缀数；用 \code{std::map} 实现日程安排，理解前驱和后继为什么足以判断新区间冲突。

\subsection{图、状态图和搜索维度}

图题第一步永远是建模。网格题的节点可以是格子，也可以是格子加钥匙集合、剩余障碍次数或时间；课程表的节点是课程，边是依赖；单词接龙的节点是单词，边是一位字符变化；转盘锁的节点是四位状态，边是一次拨动。节点维度少了会漏解，节点维度多了会超时。

BFS 适合无权最短路，因为队列按层扩展，第一次到达就是最短。DFS 适合连通块、路径枚举和拓扑状态检测。拓扑排序适合有向依赖，入度为零代表当前可执行。并查集适合动态连通，但不保存路径。最短路算法选择取决于边权：等权 BFS，非负 Dijkstra，负边 Bellman-Ford 或 SPFA 风险评估，边权只有 0 和 1 则 0-1 BFS。

图的性能常败在邻居生成。单词接龙如果每次和所有单词比较，会很慢；通配模式桶把邻居查询降到共享模式。网格 BFS 如果出队才标记访问，可能被同层多个父节点重复入队；入队时标记能控制规模。课程表如果边方向建反，判环可能仍然对，但输出顺序会错。图题必须把邻接表、访问标记和状态编码讲清楚。

实验建议做状态维度对比。打开转盘锁写普通 BFS 和双向 BFS，统计扩展节点数。障碍消除最短路分别用二维 visited 和三维 visited，对比二维为何错误剪枝。课程表分别按两种边方向建图，观察输出顺序差异。最小体力路径写 Dijkstra、二分 BFS、并查集三版，比较它们利用的是哪种单调性。

\subsection{动态规划状态表和空间压缩}

DP 的起点不是公式，而是暴力搜索树。每个搜索节点描述处理到哪里、剩余资源是什么、已经做了什么选择；如果不同路径到达同一个节点后未来选择完全一样，就有重叠子问题。状态定义就是给这种节点命名。定义不稳定时，代码会偷偷依赖历史，边界也会变成补丁。

线性 DP 看最后一步，二维 DP 看两个前缀或网格位置，背包 DP 看物品阶段和容量，区间 DP 看最后合并点，树形 DP 看子树向父节点汇报的信息，状态压缩 DP 看集合中哪些元素已访问。每一类 DP 都要问：状态是否足够，转移是否覆盖全部选择，遍历顺序是否保证依赖已计算，初始化是否对应空状态。

空间压缩是高风险优化。0-1 背包倒序容量，是为了让当前物品只能用一次；完全背包正序容量，是为了允许当前物品重复使用；LCS 一维压缩需要保存左上角旧值；股票状态机需要先保存昨天状态，避免同一天污染。初学时先写清二维或完整状态，确认正确后再压缩。

DP 实验可以非常具体。打印编辑距离表，手动解释每个格子来自插入、删除还是替换。把 0-1 背包容量循环改成正序，用一个物品重复使用的反例验证错误。把零钱兑换 II 的循环顺序调换，观察组合数如何变成排列数。给股票冷冻期状态机打印每天 hold、sold、rest，观察买入来源为什么受限制。


\section{综合案例、实验与源码对照}

这一章专门补齐三个容易被刷题书忽略的环节：第一，经典题不能直接给最终答案，而要从暴力方法、瓶颈定位、结构观察和正确性证明一步步推出来；第二，C++ 容器不是语法糖，容器的内存模型、迭代器失效、比较器和分配成本会改变工程质量；第三，Linux 源码里的数据结构不是孤立技巧，而是访问模式、生命周期、并发和内存分配共同作用的结果。读本章时，不要把它当题号清单，而要把它当作一组可以反复做的实验。

\subsection{经典案例推导}

\noindent\textbf{经典题完整推导。}

\noindent\textbf{LeetCode 1 两数之和。}

这道题看似简单，却是“从暴力到结构”的最小样本。题面要求找两个下标，值相加等于目标。最直接的暴力做法是两层循环枚举 \code{i} 和 \code{j}，检查 \code{nums[i] + nums[j] == target}。这个版本一定正确，因为它覆盖了所有下标对；它的问题是同一个历史值会被未来很多位置反复扫描。把瓶颈说清楚后，优化就自然出现：当扫描到当前值 \code{x} 时，需要的另一个数只可能是 \code{target - x}，所以不必再遍历所有历史下标，只要查历史表即可。

这里的哈希表不是“模板”，而是把暴力中的历史扫描换成等值查询。表的 key 是已经出现过的数，value 是它的下标。顺序必须先查再插入当前数，因为题目不允许同一个下标使用两次。如果先插入再查，输入 \code{[3]}、目标 \code{6} 这种边界就会把当前元素和自己配对。重复数字也要靠这个顺序处理：输入 \code{[3,3]}、目标 \code{6} 时，扫描第二个 \code{3} 才能查到第一个 \code{3}。

正确性可以对照暴力证明。任意合法答案有两个下标 \code{i < j}。算法扫描到 \code{j} 时，\code{i} 的值已经插入历史表；由于 \code{nums[i] + nums[j] == target}，查 \code{target - nums[j]} 一定能找到 \code{i}。反过来，算法找到的历史下标一定早于当前下标，且值满足目标和，所以返回结果合法。复杂度从 \complexity{O(n^2)} 降到平均 \complexity{O(n)}，代价是 \complexity{O(n)} 额外空间。实验时要覆盖没有答案、重复值、负数、目标为零、补数等于当前值这些输入。

\noindent\textbf{LeetCode 3 无重复字符的最长子串。}

暴力版本可以枚举每个左端和右端，再检查子串内是否有重复字符。如果检查也用集合扫描，复杂度会很高。优化的第一步不是立刻说“滑动窗口”，而是观察连续子串的变化：右端向右扩展时，只新增一个字符；如果这个字符在当前窗口里重复，左端只需要移动到它上一次出现位置之后。左端永远不会向左回退，因为已经被排除的左端只会让窗口更长、更容易包含重复字符。

高质量实现通常保存字符上一次出现的位置。扫描右端 \code{right} 时，如果字符 \code{s[right]} 的上次位置在当前窗口内部，就把 \code{left} 更新为 \code{last[s[right]] + 1}；如果上次位置在窗口左侧之外，就不能把 \code{left} 拉回去。这里最常见的错误是直接写 \code{left = last[c] + 1}，导致输入 \code{abba} 时，处理最后一个 \code{a} 会把左端从 2 错误拉回 1。

这题的正确性来自窗口不变量：任意时刻，\code{[left,right]} 内没有重复字符，且 \code{left} 是在当前 \code{right} 下能保持无重复的最左合法边界。每次右端加入一个字符，只可能因为这个新字符破坏合法性；把左端移过它上次出现位置后，窗口重新合法。答案在每次窗口合法后更新。实验应打印 \code{left}、\code{right}、当前字符、上次位置和答案，特别观察 \code{abba}、空串、全相同字符、全部不同字符以及扩展字符集的情况。

\noindent\textbf{LeetCode 15 三数之和。}

暴力方法三层循环枚举三个下标，能够覆盖全部答案，但复杂度是 \complexity{O(n^3)}，并且重复组合处理很麻烦。优化从排序开始。排序的意义不是让数组变好看，而是制造单调性：固定第一个数以后，剩下两个数在有序数组中寻找目标和。如果左指针和右指针的和太小，左指针右移才能增大；如果和太大，右指针左移才能减小。这样一层固定加一层双指针就能覆盖所有二元组合。

去重是这题的核心边界。固定第一个数时，如果它和前一个固定值相同，就跳过，否则会生成重复三元组。找到一个三元组后，左右指针也要分别跳过相同值，避免同一组值被多个下标组合重复加入。注意，去重不是随便用集合过滤结果；集合过滤可以让答案看起来对，但会掩盖推导过程，也可能带来额外排序和哈希成本。

正确性可以从固定值证明。排序后，对于每个固定位置 \code{i}，双指针会在区间 \code{(i,n)} 中按单调规则排除不可能的配对。若当前和小于目标，所有使用当前左指针和更小右指针的组合只会更小，所以左指针可以安全右移；若当前和大于目标，所有使用当前右指针和更大左指针的组合只会更大，所以右指针可以安全左移。实验要覆盖全零、大量重复、无解、正负混合、答案很多以及目标和可能溢出的四数扩展。

\noindent\textbf{LeetCode 42 接雨水。}

这题最容易被误讲成“背双指针”。先从暴力出发：对每个位置 \code{i}，它能接的水量由左侧最高柱和右侧最高柱的较小者决定，再减去当前高度。暴力每个位置都向两边扫描，复杂度 \complexity{O(n^2)}。第一层优化是预处理左右最高数组，令 \code{left_max[i]} 表示 \code{i} 左侧含自身最高，\code{right_max[i]} 表示右侧含自身最高。这样每个位置的贡献能常数时间计算。

双指针进一步把左右最高数组压缩成两个变量。左指针和右指针从两端向中间走，同时维护 \code{left_max} 和 \code{right_max}。如果 \code{left_max <= right_max}，左侧当前位置的水位已经由 \code{left_max} 决定，因为右侧至少存在一个不低于 \code{left_max} 的边界；此时可以结算左指针并右移。反之结算右指针。这个证明比代码更重要：结算的一侧必须已经知道较小边界，否则会提前计算。

单调栈解法则从另一角度结算凹槽。栈中保存递减高度下标，遇到更高的右边界时，弹出槽底，用当前柱作为右边界、新栈顶作为左边界计算水量。三种解法都正确，但状态不同：前后缀数组保存每个位置的边界，双指针保存当前两侧边界，单调栈保存尚未找到右边界的柱子。实验时分别打印三个版本的状态，对比单调递增、单调递减、平台高度、多个凹槽和长度不足三的情况。

\noindent\textbf{LeetCode 76 最小覆盖子串。}

暴力方法枚举所有子串，再检查它是否覆盖目标串字符需求。它的问题有两层：枚举窗口多，检查窗口也重复统计。优化要抓住覆盖关系的单调性。右端扩展时，窗口字符只会增加，覆盖条件更容易满足；一旦满足覆盖，继续扩展只会让窗口更长，不可能改善当前左端下的答案，所以应该尽量收缩左端，直到刚好失去覆盖。

实现中需要两个计数结构：目标需求 \code{need} 和窗口计数 \code{window}。为了避免每次比较全部字符，可以维护 \code{formed}，表示有多少种字符已经达到需求次数。注意这里是“达到需求次数的字符种类”，不是窗口里匹配的字符总数。目标串包含重复字符时尤其容易错，例如 \code{AABC} 需要两个 \code{A}，窗口里一个 \code{A} 不算满足。

正确性来自最短窗口的收缩。每当右端使窗口满足覆盖，内层循环会尽可能移动左端，并在每个合法窗口上更新答案。对于某个固定右端，循环结束前已经检查了所有能覆盖目标的左端中最短的那个；对于所有右端扫描完成后，全局最短自然被覆盖。实验要覆盖目标字符缺失、目标含重复、最短窗口在开头、最短窗口在结尾、大小写混合和源串短于目标串。C++ 中若字符集固定，可用数组；若字符集不确定，用哈希表但要注意 \code{operator[]} 的默认插入。

\noindent\textbf{LeetCode 84 柱状图中最大的矩形。}

暴力方法枚举每根柱子作为矩形高度，再向左右扩展直到遇到更矮柱子。这个做法清楚地暴露了瓶颈：每根柱子都在寻找左右第一个更矮的位置。单调栈优化就是一次扫描求出“某根柱子作为最矮高度时的最大宽度”。栈中保存高度递增的下标；当遇到一个更矮柱子时，栈顶柱子的右边界已经确定为当前位置，左边界是弹出后新的栈顶。

宽度计算是这题最容易错的地方。弹出下标 \code{mid} 后，如果栈为空，说明它左侧没有更矮柱子，宽度是 \code{right}；如果栈不空，左边界是 \code{stack.top()}，宽度是 \code{right - stack.top() - 1}。为了统一结算，可以在末尾补一个高度为零的哨兵，让所有柱子最终都被弹出。相等高度如何处理也要固定：常见写法遇到小于当前栈顶才弹，或者遇到小于等于时弹，只要宽度语义和去重一致即可。

正确性来自结算时机。某根柱子被弹出，说明当前柱子是它右侧第一个更矮位置；弹出后新的栈顶是它左侧第一个更矮位置。两个边界之间所有柱子高度都不低于它，所以以它为最矮柱的最大矩形已经确定，未来再向右也会被当前更矮柱阻断。实验要覆盖递增、递减、全相等、单柱、空数组、哨兵和宽度计算。这个题也是 LeetCode 85 最大矩形的基础：二维矩阵每一行都能转成一组柱状图高度。

\noindent\textbf{LeetCode 146 LRU 缓存。}

LRU 不是“哈希表加链表”六个字。题目要求 \code{get} 和 \code{put} 都是常数时间，并且容量满时淘汰最久未使用项。先看暴力：用数组或向量保存访问顺序，每次访问都线性查找并移动到头部；定位和移动都会慢。哈希表能按 key 常数定位，但无法维护新旧顺序；双向链表能常数删除已知节点并移动到头部，但无法按 key 查找。两个结构职责互补，组合起来才满足要求。

高质量实现要先定义节点身份。链表节点里保存 key 和 value；哈希表从 key 映射到链表迭代器或节点指针。为什么节点里还要保存 key？因为淘汰尾部节点时，需要从哈希表删除对应 key。如果节点只保存 value，尾部删除后无法常数时间更新哈希表。访问已有 key 时，更新 value 后要把节点移动到头部；插入新 key 时，容量超过上限就删除尾部。

正确性来自两个不变量：哈希表中的每个 key 都指向链表中的一个节点，链表从头到尾按最近使用到最久未使用排序。每次 \code{get} 或成功更新都会把节点移到头部，所以尾部始终是最久未使用项。实验要覆盖容量为一、重复 \code{put}、访问不存在 key、更新已有值、连续淘汰和插入后立即访问。对照 Linux 源码时，可以进一步理解侵入式链表：业务对象自己带链表节点，哈希索引只保存对象位置，链表不拥有对象。

\noindent\textbf{LeetCode 200 岛屿数量。}

这题的第一步是建模：网格中的陆地格子是节点，四方向相邻陆地之间有边，岛屿就是连通分量。暴力错误做法常常从每个陆地都重新搜索，或者不标记访问导致重复计数。正确方法是扫描每个格子，遇到一个未访问陆地，就启动 BFS 或 DFS，把整个连通块标记为已访问，并把岛屿数量加一。

BFS 和 DFS 在这题里都可以。BFS 用队列保存待扩展格子，入队时标记访问，避免同一格子被多个邻居重复加入；DFS 可以递归，也可以用显式栈。工程上大网格可能导致递归栈风险，因此显式栈更稳。若允许修改输入，可以把访问过的 \code{1} 改成 \code{0}；若不允许修改，就使用单独访问数组。这个选择不是风格问题，而是题目副作用要求。

正确性来自连通块定义。一次搜索从某个未访问陆地出发，会沿合法边访问所有与它连通的陆地；搜索不会越过水或边界，所以不会把两个不同岛屿合并。外层扫描保证每个连通块至少会被第一次遇到的陆地启动一次，而访问标记保证同一连通块不会重复计数。实验覆盖全水、全陆、单个陆地、斜角相邻但不连通、细长岛、输入是否可修改以及大网格递归深度。

\noindent\textbf{LeetCode 207 课程表。}

课程表题不是普通遍历，而是有向依赖图判环。每门课是节点，先修关系是有向边。如果学习课程 \code{b} 后才能学 \code{a}，可以建边 \code{b -> a}，表示 \code{b} 完成后减少 \code{a} 的依赖。暴力思路是反复找没有先修要求的课程并删除它；这正是拓扑排序的朴素模型。优化做法用入度数组保存每门课尚未满足的依赖数量，用队列保存当前入度为零的课程。

拓扑排序的过程是：初始把所有入度为零的节点入队；每弹出一门课，就把它指向的后继课程入度减一；如果某个后继入度变成零，就入队。最后如果弹出的课程数等于总课程数，说明所有依赖都能被满足；否则剩余课程处在环里或依赖环。边方向要讲清楚，否则输出顺序会反。判能否完成时，边方向反了有时也能判环，但到了课程表 II 输出顺序就会错。

正确性来自入度语义。入度为零表示当前没有未完成先修，可以安全安排；安排它只会减少后继依赖，不会破坏其他课程合法性。若最终还有节点未输出，这些节点的入度无法降到零，说明它们之间存在循环依赖。实验要覆盖无依赖、单链依赖、简单环、多个连通分量、自环、重复边和边方向检查。DFS 三色标记也能做，但需要区分“当前路径中”和“已经完成”的节点。

\noindent\textbf{LeetCode 300 最长递增子序列。}

先写二次 DP。令 \code{dp[i]} 表示必须以 \code{nums[i]} 结尾的最长递增子序列长度。对每个 \code{i}，枚举所有 \code{j < i}，若 \code{nums[j] < nums[i]}，就可以从 \code{j} 转移到 \code{i}。这个状态非常清楚，复杂度是 \complexity{O(n^2)}。它的问题是每个位置都要扫描所有历史前驱，很多历史状态只是在比较“能否成为某个长度的更小结尾”。

二分优化维护 \code{tails[len]}：长度为 \code{len + 1} 的递增子序列，当前能达到的最小结尾值。这个数组不是某条真实子序列，而是未来扩展潜力表。扫描新数 \code{x} 时，用 \code{lower_bound} 找第一个大于等于 \code{x} 的位置并替换；如果位置在末尾，就把 \code{x} 追加，说明能形成更长序列。为什么替换安全？同样长度下，结尾越小，未来越容易接上更大的数，不会比原结尾差。

正确性要强调两个层次。第一，\code{tails} 的长度始终是已经扫描前缀中可形成的某个递增子序列长度；第二，\code{tails} 每个位置保存该长度的最小可能结尾，因此替换不会减少当前答案，只会改善未来扩展空间。重复元素时使用 \code{lower_bound} 保持严格递增；若题目改成非递减，就要用 \code{upper_bound}。若要还原路径，单靠 \code{tails} 不够，需要前驱数组和位置记录。

\noindent\textbf{LeetCode 322 零钱兑换。}

这题要求凑出金额的最少硬币数，硬币可以无限使用。暴力搜索会枚举每种硬币使用次数或递归尝试所有选择，很多剩余金额会重复出现。状态可以定义为 \code{dp[value]}：凑出金额 \code{value} 的最少硬币数。初始化 \code{dp[0] = 0}，其他为无穷。对每个金额，枚举硬币，如果 \code{value >= coin} 且 \code{dp[value - coin]} 可达，就尝试更新 \code{dp[value]}。

这里要区分“最少数量”和“组合数量”。零钱兑换 I 求最少硬币数，内外循环顺序有多种写法，只要转移语义正确即可；零钱兑换 II 求组合数，循环顺序决定是否把不同顺序算作不同方案。哨兵也要谨慎：如果无穷值设得太大，\code{dp[value - coin] + 1} 可能溢出。常见做法是把无穷设为 \code{amount + 1}，因为最坏情况下使用面额一的硬币也不会超过这个数量。

正确性来自最后一枚硬币。任意最优方案最后使用某枚 \code{coin}，去掉它后剩余金额是 \code{value - coin} 的最优方案；否则如果剩余部分不是最优，就能替换得到更少硬币。实验覆盖金额为零、无法凑出、硬币面额大于目标、重复硬币、存在面额一和不存在面额一。复杂度是 \complexity{O(amount \cdot coins)}，这是伪多项式复杂度，依赖金额数值而不是输入字符长度。

\noindent\textbf{LeetCode 560 和为 K 的子数组。}

暴力枚举所有左右端点并求和，若每次重新求和是 \complexity{O(n^3)}，用前缀和可以降到 \complexity{O(n^2)}。继续优化的关键是代数关系：区间 \code{(i,j]} 的和等于 \code{prefix[j] - prefix[i]}。当扫描到当前前缀 \code{prefix[j]} 时，需要寻找多少个历史前缀等于 \code{prefix[j] - k}。这就是哈希表保存历史前缀频次的理由。

先查询还是先更新是本题核心边界。我们要找的是当前下标之前的历史前缀，所以应该先用当前前缀查询 \code{prefix - k}，再把当前前缀加入表。初始表里要放 \code{0 -> 1}，表示从数组开头到当前下标的区间也可能满足条件。若漏掉这个初始值，输入 \code{[k]} 会错误返回零。

这题能处理负数，是它和正数滑动窗口的关键区别。滑动窗口依赖右端扩展和左端收缩带来的和单调变化；一旦有负数，窗口和可能来回变化，不能安全排除左端。前缀和加哈希不需要单调性，只需要代数等式。实验覆盖负数、零、目标为零、重复前缀、从下标零开始的答案和答案数量很多。若题目改成最长区间，哈希表应保存最早位置；若改成是否存在，保存出现与否即可。

\noindent\textbf{LeetCode 743 网络延迟时间。}

题目给有向带权图，要求从起点到所有节点的最短路最大值。暴力枚举路径会爆炸；普通 BFS 只适用于等权边，因为它按边数层次扩展，而题目边权可能不同。边权非负时，Dijkstra 是合适选择：维护每个节点当前已知最短距离，用小顶堆每次取距离最小的候选扩展。

C++ 的 \code{priority_queue} 不支持 decrease-key，所以常用“压入新状态、弹出时跳过过期状态”的写法。每次松弛边 \code{u -> v}，如果 \code{dist[u] + w < dist[v]}，就更新 \code{dist[v]} 并把新状态放入堆。堆里可能还保留旧距离；弹出时如果堆顶距离不等于当前 \code{dist[node]}，说明它已经过期，直接跳过。这个懒处理是容器限制下的工程实现，不是算法偷懒。

正确性依赖非负边。每次弹出的未过期节点拥有当前全图最小的临时距离，由于所有边权非负，从其他未确定节点绕回来不可能得到更短路径，因此它的距离可以确定。实验要覆盖不可达节点、重边、起点到自己、零权边、多条等长路径和堆中过期状态。最终若有节点距离仍为无穷，返回 -1；否则返回最大最短距离。

\noindent\textbf{LeetCode 862 和至少为 K 的最短子数组。}

如果数组全是正数，长度最小且和至少为 K 可以用滑动窗口；但本题允许负数，窗口和不再单调。一个负数进入窗口可能让和变小，一个负数离开窗口可能让和变大，所以单纯移动左右端无法安全排除候选。先写前缀和：需要找 \code{prefix[j] - prefix[i] >= k} 且 \code{j - i} 最短。对每个 \code{j}，我们希望历史前缀尽量小且下标尽量靠近。

单调队列维护候选前缀下标，队列中的前缀和严格递增。处理当前 \code{j} 时，若队首和当前前缀差至少 K，就可以更新答案并弹出队首，因为这个队首作为左端已经得到最短的当前右端，未来右端更远只会长度更长。然后维护队尾：如果队尾前缀和大于等于当前前缀，且当前下标更靠右，那么队尾被当前下标支配，未来用当前下标只会得到更大的区间和和更短长度。

正确性核心是支配关系。前缀和更大且下标更早的候选不如前缀和更小且下标更晚的候选；满足条件的队首在当前右端下已经结算完毕。实验要覆盖负数破坏普通窗口、答案从开头开始、无答案、多个候选同时满足、队尾连续弹出和 K 为负或零的边界。这个题是单调队列最值得精读的一题，因为它把前缀和、最短长度和支配关系放到一起。

\noindent\textbf{LeetCode 1631 最小体力消耗路径。}

这题的路径代价不是边权和，而是路径上最大高度差。暴力可以枚举路径，但路径数量巨大。第一种高质量建模是 Dijkstra 变体：节点是格子，边权是相邻格子高度差，路径代价是当前路径最大边权。从一个格子扩展到邻居时，新代价是 \code{max(current_cost, edge_cost)}。这个代价随着路径延长不会下降，因此仍然满足按当前最小代价优先扩展的单调性。

第二种建模是二分答案。猜一个体力上限 \code{limit}，只允许走高度差不超过 \code{limit} 的边，检查左上角是否能到右下角。上限越大，可用边越多，可达性单调，所以可以二分最小可行上限。第三种建模是并查集：把所有边按高度差排序，从小到大加入，直到起点和终点连通，此时加入的边权就是答案。

三种方法利用的是不同结构。Dijkstra 把路径代价作为动态极值扩展；二分 BFS 利用可行性单调；并查集利用从小到大开放边的连通性。面试中能比较三者，比只写一种代码更有说服力。实验覆盖单格、平坦网格、必须绕路、局部高度差很大但可绕开、多个同代价路径和大网格性能。注意不要把路径代价误写成边权和，否则就变成了另一道最短路题。

\subsection{容器实验和源码对照}

\noindent\textbf{C++ 容器底层实验。}

\noindent\textbf{实验一：\code{vector} 的容量和地址。}

很多刷题代码默认 \code{vector} 是“无限数组”，但工程上必须理解它的容量模型。写一个实验：从空 \code{vector<int>} 开始连续 \code{push_back}，每次当 \code{capacity} 改变时打印旧容量、新容量和 \code{data()} 地址。你会看到容量不是每次加一，而是按某种增长策略跳变；地址也可能变化。这解释了为什么保存元素指针、引用或迭代器后继续插入是危险的。

第二步比较 \code{reserve} 和 \code{resize}。\code{reserve(1000)} 只保证容量，不改变 \code{size}，不能直接访问 \code{v[999]}；\code{resize(1000)} 会构造一千个元素，\code{size} 变为一千。很多同学为了避免扩容使用 \code{resize}，结果把未写入元素也当成有效数据，答案被默认零污染。高质量题解要说明自己需要的是容量还是大小。

第三步比较顺序扫描 \code{vector} 和 \code{list}。两者理论上都是线性，但 \code{vector} 连续内存通常有更好的缓存局部性，\code{list} 每个节点分散分配，指针追踪成本高。这个实验能纠正一个常见误解：链表插入删除复杂度好，不代表所有场景都快。若题目主要是顺序扫描和随机访问，\code{vector} 往往是更好的默认选择。

\noindent\textbf{实验二：\code{unordered_map} 的查询语义。}

哈希表题最常见错误不是哈希函数，而是 API 语义。写三版前缀和计数：第一版用 \code{operator[]} 查询，第二版用 \code{find}，第三版用 \code{contains} 后再取值。观察不存在 key 时，\code{operator[]} 会插入默认值，导致表大小变化。计数累加时这很方便，但只想判断存在时可能污染状态，也会让调试输出误以为某些前缀已经真实出现过。

再做 \code{reserve} 实验。向 \code{unordered_map} 插入大量元素，记录桶数量和负载因子变化。没有 \code{reserve} 时，插入过程中可能多次 rehash，迭代器失效，性能也会波动。刷题中不一定必须 \code{reserve}，但当输入规模很大、哈希表是热路径时，主动预留容量可以降低常数成本。需要强调的是，\code{reserve} 不能修复错误 key 设计；如果 key 构造昂贵或容易冲突，复杂度仍然会差。

最后检查 key 设计。异位词分组若把计数直接拼成字符串而不加分隔符，\code{1,11} 和 \code{11,1} 可能混淆；二维坐标若用 \code{x * M + y} 编码，要确认不会溢出且 M 选择正确；自定义结构作为 key 时，哈希函数和相等比较必须表达同一个等价关系。哈希表只是查询工具，真正决定正确性的是 key 的语义。

\noindent\textbf{实验三：\code{priority_queue} 和过期状态。}

\code{priority_queue} 只支持访问和弹出堆顶，不支持删除任意元素，也不支持 decrease-key。Dijkstra、滑动窗口中位数、任务调度和 Top K 题都必须面对这个限制。写一个 Dijkstra 小实验：每当发现更短距离，就把新状态压入堆；弹出时如果距离和 \code{dist} 数组不一致，就统计一次过期弹出。这个实验能让你直观看到“堆中有旧状态”并不影响正确性，只要使用前检查。

再写滑动窗口最大值或中位数的懒删除实验。窗口滑动时，一些元素已经过期，但它们可能不在堆顶，无法立即删除。可以用哈希表记录待删除次数，等它们浮到堆顶时再真正弹出。这个思想和 RCU 的延迟释放相似但目的不同：堆懒删除是容器能力不足下的工程技巧，RCU 是并发读者安全要求。

比较器也要单独实验。C++ 默认 \code{priority_queue<int>} 是大顶堆；小顶堆要写 \code{greater<int>}。如果存的是结构体，比较器含义是“谁的优先级更低”，很多人会写反。若比较器对相等元素处理不满足严格弱序，排序和堆行为都可能不稳定。任何堆题都要先在纸上写出堆顶应该是谁，再写比较器。

\noindent\textbf{实验四：\code{map}、\code{set} 和动态区间。}

有序容器适合前驱、后继和范围边界。写一个“我的日程安排”实验：用 \code{map<int,int>} 保存不重叠区间，key 是起点，value 是终点。插入新区间 \code{[l,r)} 时，先找第一个起点不小于 \code{l} 的后继，再检查前驱和后继是否与新区间重叠。只要这两个邻居合法，其他区间就不可能冲突，因为已有区间按起点有序且互不重叠。

这个实验能解释为什么哈希表不适合动态区间。哈希表能查某个起点是否存在，却不能快速找到离 \code{l} 最近的前驱和后继。排序数组适合离线处理，但动态插入会移动大量元素。红黑树类结构牺牲一些常数，换来稳定的对数级顺序操作。Linux 的 rbtree 和 Maple Tree 也服务于类似需求，只是多了侵入式节点、并发和内存管理约束。

实验报告要记录闭区间和半开区间的差异。若使用 \code{[l,r)}，前一个区间终点等于新区间起点不算重叠；若使用闭区间，端点相等可能算重叠。很多区间题错误都来自没定义端点语义。写比较器时也要小心相同起点：覆盖类题常需要起点升序、终点降序；合并类题只按起点升序通常足够。

\noindent\textbf{Linux 源码数据结构对照。}

\noindent\textbf{\code{list_head} 和 LeetCode 146。}

Linux 的 \code{list_head} 和力扣 LRU 的双向链表有共同点：都需要已知节点常数时间删除、移动到头部和从尾部淘汰。区别在所有权。力扣常用 \code{std::list<Node>} 让容器拥有节点，哈希表保存迭代器；内核常把 \code{list_head} 嵌进业务对象，链表只保存链接关系，不负责对象分配和释放。这个侵入式设计让同一个对象可以同时挂进多个链表，例如一个页面对象可以进入 LRU 链表，也可以进入其他状态链表。

源码阅读时不要先追每个宏，而要先画对象布局：业务对象里有哪些 \code{list_head} 字段，每个字段对应哪条链，插入和删除由谁保护，删除后对象是否还能被其他索引找到。刷题中的 LRU 删除尾部只要维护哈希表和链表；内核里删除一个对象可能还要处理引用计数、锁、RCU 或延迟释放。相同点是“节点身份必须稳定”，不同点是“生命周期复杂得多”。

实验可以写两个版本的 LRU。第一版使用 \code{std::list} 和 \code{unordered_map}；第二版让 \code{Entry} 自己带 \code{prev}、\code{next} 指针，哈希表保存 \code{Entry*}。比较两版的删除顺序：先从链表摘除，再从哈希表删除，最后释放对象。如果先释放对象，链表或哈希表就会留下悬空指针。这个实验能把链表题、缓存题和内核侵入式结构连起来。

\noindent\textbf{\code{hlist} 和哈希冲突。}

哈希表的桶内冲突可以用链表处理。Linux 的 \code{hlist} 面向大量桶头做了空间优化：桶头较轻，节点保存足够信息以支持已知节点删除。刷题时 \code{unordered_map} 把桶数组、负载因子和冲突链隐藏起来，但理解桶结构能帮助你解释为什么哈希表平均常数不是无条件保证。

对照题包括两数之和、字母异位词分组、和为 K 的子数组、账户合并。每道题都不是“用了哈希就完了”，而是要回答 key 是什么、冲突时如何确认相等、value 保存次数还是下标、什么时候插入、什么时候删除。内核里这些问题更具体：对象节点是否已在桶中，删除时是否有读者正在遍历，key 对应的对象生命周期由谁保证。

实验可以写一个固定 16 个桶的小哈希表。第一版桶内使用普通单链表，删除已知节点时必须从桶头找前驱；第二版让节点保存前驱链接位置，删除时直接改指针。这个实验能解释 hlist 的设计动机：不是为了让算法更神秘，而是为了在桶很多、删除频繁的场景下减少空间和扫描成本。

\noindent\textbf{rbtree 和动态有序集合。}

Linux rbtree 与 C++ \code{map} 都是动态有序结构，但接口风格不同。C++ 容器把节点、比较器和内存管理封装起来；Linux rbtree 通常把树节点嵌入业务对象，比较逻辑由调用者在插入和查找时写出。这样做适合内核对象按不同字段进入不同索引，也避免通用容器隐藏分配和生命周期。

对照刷题时，可以把 rbtree 看作“需要前驱后继的动态集合”的底层代表。日程安排、区间冲突、滑动窗口有序统计、数据流排名，都不是哈希表擅长的问题。红黑树的旋转不会改变中序顺序，只是在保持有序语义的前提下修复平衡。理解这一点后，再看 \code{lower_bound}、前驱、后继和区间插入，会更容易说清为什么只检查相邻区间。

源码阅读建议选一条简单路径：从查找插入位置开始，看调用者如何比较 key；再看链接新节点；最后看插入修复如何通过旋转和变色维持平衡。不要第一天就背红黑树五条性质的所有修复分支。刷题阶段也一样：先证明有序性和相邻关系，再谈容器实现。

\noindent\textbf{XArray、Trie 和稀疏整数索引。}

XArray 面向整数索引到指针的映射，适合稀疏 ID 空间和内核对象管理。它和 Trie 有相似思想：把 key 分段，沿层级向下查找，只为实际存在的路径分配节点。与普通哈希相比，它更重视按索引组织、范围推进、标记和并发语义；与巨大数组相比，它避免为空洞分配大量空间。

刷题中可以用位 Trie、前缀树、坐标压缩来建立直觉。最大异或题把整数按二进制位分层；单词搜索 II 把字符串按字符分层；坐标压缩把稀疏大值域映射到紧凑下标。XArray 的工程环境更复杂，但核心问题仍是：key 空间很大，实际对象稀疏，等值查找之外还可能需要按序扫描和状态标记。

实验可以写一个简化的二进制 Trie 保存整数集合，再支持查找某个整数是否存在、查找不小于某个值的下一个元素。这个实验不会复刻 XArray，但能训练分层索引的感觉。报告中要比较数组、哈希表、平衡树和 Trie：它们分别适合连续小值域、稀疏等值查找、动态有序查找、可分段 key。

\noindent\textbf{Maple Tree 和动态区间。}

Maple Tree 服务于范围映射场景，典型直觉是虚拟地址区间管理。区间问题比单点 key 更难：查询一个地址需要找到覆盖它的区间；插入新区间需要检查前后邻居是否重叠；删除可能删除整段，也可能把原区间切成两段。力扣 56 和 57 是离线区间题，Maple Tree 面对的是动态范围集合和高频查询。

刷题阶段可以先用 \code{std::map} 写简化区间表。key 是区间起点，value 是终点和属性。点查询时找起点不大于地址的最后一个区间，再判断终点是否覆盖；插入时检查前驱和后继；删除中间一段时要拆分。这个小实验能帮助理解 Maple Tree 的问题动机，而不是把源码当成神秘结构背诵。

工程差异在于读写并发、内存分配和查询路径。一个简单 \code{map} 版本可以帮助学习区间不变量，却不能替代内核结构。源码阅读时要问：节点里保存哪些范围信息，读路径如何快速找到覆盖区间，更新如何保持相邻关系，旧节点什么时候释放。这些问题比记字段名更稳定。

\noindent\textbf{伙伴系统、SLUB 和分配成本。}

算法题通常把内存分配当成免费，但 Linux 内存管理会把分配成本放到中心。伙伴系统按 2 的幂次阶管理连续页，分配时从合适阶取块，不够就拆更高阶；释放时检查伙伴是否空闲，能合并就合并。它强调连续性和对齐，不只是总空闲量。一个系统可能总空闲页很多，却没有足够大的连续块。

SLUB 解决小对象频繁分配问题。同类型对象从缓存中分配，避免每次都走通用页分配路径，也提升局部性。回到刷题，链表节点、哈希节点、树节点都可能带来大量小对象分配；这就是为什么连续数组有时远快于节点容器。复杂度分析给出数量级，分配器和缓存局部性决定常数。

实验可以写简化伙伴系统：管理 \code{2^10} 个页框，维护 0 到 10 阶空闲链表。分配某阶时逐级拆分，释放时按伙伴编号尝试合并。再写一个对象池，为固定大小节点重复分配和释放，比较它与频繁 \code{new} 的差异。这个实验能把位运算、链表、区间合并和工程性能放在同一张图里。

\subsection{错误用例、复盘和训练路线}

\noindent\textbf{错误用例库。}

\noindent\textbf{数组和窗口。}

数组题的错误用例要专门打破区间语义。删除重复项要测空数组、全重复、无重复和只保留两次的变体；颜色分类要测全二、全零、逆序和交换后二次检查；接雨水要测单调递增、单调递减和平台高度。窗口题则要测试左端不能回退的情况，例如无重复子串的 \code{abba}；还要测试目标字符重复，例如最小覆盖子串中目标含两个相同字符。

设计反例时，不要只看最终答案，要让错误状态暴露出来。滑动窗口题可以打印左端、右端、计数和答案；双指针题可以打印每次为什么移动某一侧。一个好的反例能说明错误发生在初始化、循环条件、状态更新还是答案读取，而不是只告诉你“答案错了”。这也是把题解变成教材的关键。

\noindent\textbf{哈希和前缀和。}

哈希题的错误用例集中在查询和更新顺序、key 设计、保存次数还是下标。两数之和要测补数等于当前值；和为 K 的子数组要测从下标零开始的答案和目标为零；连续子数组和要测长度至少为二；连续数组要测重复前缀且要求最长。异位词分组要测空字符串、重复单词和计数 key 歧义。

前缀和题尤其适合对拍。写一个平方暴力版，再随机生成包含负数、零和重复值的数组，与哈希版比较。很多窗口解法在正数数组上正确，一加负数就失败；对拍能强迫你说明算法依赖的是单调性还是代数等式。若目标从“计数”改成“最长”，哈希表 value 也要从频次改成最早下标，这类变化必须进反例库。

\noindent\textbf{二分和答案空间。}

二分错误不是只差一，而是区间语义没有固定。测试要覆盖空数组、长度一、长度二、目标小于所有元素、目标大于所有元素、目标重复出现、答案在边界。答案二分还要测试刚好可行和刚好不可行的相邻值。每个测试都打印 \code{left}、\code{mid}、\code{right} 和谓词值，检查被排除的一半是否真的不含答案。

对于船运包裹、爱吃香蕉、分割数组这类题，检查函数本身就是贪心题。反例要针对检查函数，而不是只针对外层二分。比如容量小于最大包裹必须不可行；天数为一时容量必须至少是总和；天数等于包裹数时容量下界就是最大包裹。若检查函数没有这些边界，二分循环再标准也会返回错误答案。

\noindent\textbf{栈、堆和单调结构。}

单调栈要测试严格递增、严格递减、相等元素和哨兵。柱状图最大矩形中，递增数组能暴露末尾未结算问题；递减数组能暴露宽度计算问题；全相等能暴露相等元素处理策略。每日温度要测试没有更高温、相同温度不算更高、最后一个元素答案为零。

堆题要测试过期状态和比较器方向。Top K 高频要测试频次相同和 K 等于不同元素数；数据流中位数要测试偶数个元素平均溢出；Dijkstra 要测试重边和旧堆状态；滑动窗口中位数要测试要删除的元素不在堆顶。每个堆实验都应该打印堆顶是否真实可用，而不是默认堆顶一定有效。

\noindent\textbf{树、图和并查集。}

树题要测试空树、单节点、链式树、完全树、重复值、全负值和极端深度。验证 BST 的反例必须包含“局部父子合法但全局不合法”，例如某个右子树深处节点小于根。最大路径和要测试全负数，答案不能初始化为零。构造树题要测试节点值唯一性，如果值不唯一，单靠值到中序下标的哈希表就不成立。

图题要测试自环、重边、不连通、多个源点、状态资源不同但位置相同。打开转盘锁要测试起点死亡和目标就是起点；障碍消除最短路要测试同一格子剩余消除次数不同；课程表要测试自环和两个节点互相依赖。并查集要测试重复合并、已经连通的边、字符串编号和合并失败代表环，而不是路径本身。

\noindent\textbf{动态规划。}

DP 反例要打状态定义。最大子数组和要测试全负数；不同路径要测试一行一列；障碍物版本要测试起点或终点障碍；编辑距离要测试空串；背包要测试容量为零、总和奇偶、一个物品被重复使用的循环方向错误；零钱兑换 II 要测试组合和排列的差异。空间压缩题还要测试更新顺序，例如 0-1 背包正序会错误重复使用当前物品。

打印 DP 表是最有效的实验。编辑距离表能看到插入、删除、替换；LCS 表能看到不连续子序列和子串的区别；股票状态机能看到持有、卖出、休息之间的约束。若一个状态无法解释表中某个格子的来源，就说明状态定义还不够稳定。不要急着压缩空间，先让完整状态表正确。

\noindent\textbf{复盘报告模板。}

每道题复盘写一页即可，但必须有固定信息。第一行写题号和题名，例如 LeetCode 560 和为 K 的子数组。第二行写题面模型：对象、关系、输出。第三行写暴力基线：如何保证不漏候选，复杂度是多少。第四行写瓶颈：哪一步重复扫描、重复计算或无法快速查询。第五行写优化结构：保存什么，查询什么，更新什么，删除什么。第六行写正确性：覆盖和排除各一句。第七行写边界用例：至少三个能打破错误实现的输入。第八行写 C++ 注意点：容器、溢出、迭代器、所有权或副作用。

这份报告不需要文学化，越具体越好。比如“用哈希优化”不是合格描述，“哈希表保存当前下标之前每个前缀和出现次数，扫描到前缀 \code{s} 时先查询 \code{s-k} 再累加当前前缀”才是合格描述。比如“注意边界”也不是合格描述，“初始放入 \code{0 -> 1}，否则从下标零开始的区间会漏掉”才是合格描述。

最后要有实验记录。把暴力版和优化版同时保留，用随机小数据对拍；把边界用例写成单元测试；把大规模输入用于观察复杂度增长。对拍不是浪费时间，它能把“我觉得对”变成“我知道错在哪里”。当你能用一页报告解释一道题，再用相似报告解释变体，刷题才真正转化为面试能力和工程能力。

\noindent\textbf{高阶经典题补充推导。}

\noindent\textbf{LeetCode 4 寻找两个正序数组的中位数。}

这题经常被当成二分难题背答案，其实先从暴力看更清楚。两个数组已经有序，最直接的方法是归并到一半位置，复杂度 \complexity{O(m+n)} 或 \complexity{O((m+n)/2)}。它没有利用“只需要中位数而不是完整合并结果”这个条件。中位数的本质是把全部元素切成左右两半，左半元素数量等于右半或多一个，并且左半最大值不大于右半最小值。

优化做法是在较短数组上二分切分位置。假设第一个数组左边取 \code{i} 个元素，第二个数组左边就必须取 \code{half - i} 个元素。这样数量条件自动满足，只剩大小条件：第一个数组左侧最大不大于第二个数组右侧最小，第二个数组左侧最大不大于第一个数组右侧最小。如果第一个数组左侧太大，说明 \code{i} 取多了，要向左；如果第二个数组左侧太大，说明 \code{i} 取少了，要向右。

边界处理依赖哨兵。切分在数组开头时，左侧最大可视为负无穷；切分在数组末尾时，右侧最小可视为正无穷。这样不需要为每个边界写特殊分支。正确性来自切分条件：数量正确且左右有序时，中位数只可能由左侧最大和右侧最小决定。实验要覆盖一个数组为空、两个数组长度差很大、总长度奇偶、重复元素、所有元素都在一个数组左侧或右侧。讲这题时重点不是背代码，而是说明二分的对象是切分位置，不是某个具体值。

\noindent\textbf{LeetCode 10 正则表达式匹配。}

这题的困难在于星号。暴力递归从字符串下标和模式下标出发，遇到普通字符或点号就匹配一个字符，遇到星号就尝试匹配零次、一次、两次甚至更多次。暴力会在同一对下标上反复尝试，复杂度爆炸。DP 状态可以定义为 \code{dp[i][j]}：字符串前 \code{i} 个字符能否匹配模式前 \code{j} 个字符。这个定义把递归中的下标对固定下来，避免重复搜索。

转移要围绕模式最后一个字符。如果 \code{p[j-1]} 不是星号，那么只有当前字符能匹配时，才能从 \code{dp[i-1][j-1]} 转移。若 \code{p[j-1]} 是星号，它修饰的是 \code{p[j-2]}。星号匹配零次时，看 \code{dp[i][j-2]}；星号匹配一次或多次时，必须当前字符能被 \code{p[j-2]} 匹配，然后看 \code{dp[i-1][j]}，表示消耗一个字符串字符后仍由这个星号继续处理。

初始化是最容易漏的地方。空字符串可以被 \code{a*}、\code{a*b*} 这类模式匹配，所以第一行需要按星号零次匹配向前传播。正确性来自模式末尾分类：任意合法匹配的最后一步要么消耗一个普通模式字符，要么让星号匹配零次，要么让星号至少匹配一次，这三类覆盖全部可能。实验覆盖空串、只含星号组合的模式、点号、星号连续修饰的结构、完全不匹配和长重复字符串。

\noindent\textbf{LeetCode 25 K 个一组翻转链表。}

这题不是单纯反转链表，而是“先确定边界，再局部反转，再接回”。暴力直觉可以把链表节点放进数组，每 K 个一组反转数组片段，再重连节点。但原地链表解法要训练的是指针边界。每一轮从上一组尾部出发，向后找第 K 个节点；如果不足 K 个，剩余节点保持原样。找到组尾后，先保存下一组头，否则反转过程中会丢失后续链表。

指针关系可以固定成四个名字：\code{group_prev} 是上一组尾部，\code{group_head} 是当前组头，\code{group_tail} 是当前组尾，\code{next_group} 是下一组头。反转当前组后，原来的 \code{group_tail} 成为新头，原来的 \code{group_head} 成为新尾。接回时，\code{group_prev->next} 指向新头，新尾的 \code{next} 指向 \code{next_group}，然后 \code{group_prev} 移到新尾。

正确性来自链表结构不变量：已经处理的前缀是按 K 分组翻转后的合法链；当前组在反转前已经确认长度足够；未处理后缀仍然由 \code{next_group} 保存且可达。实验要覆盖空链表、长度小于 K、长度等于 K、长度不是 K 倍数、K 为一、连续多组和删除头部附近边界。C++ 中不要随意释放输入节点，除非题目明确要求；这里重连节点即可。

\noindent\textbf{LeetCode 32 最长有效括号。}

暴力方法枚举每个子串并检查括号是否合法，复杂度很高。栈解法的核心是保存边界。栈里放尚未匹配的左括号下标，同时保留最后一个无法匹配的右括号位置。一个常见写法是在栈中先放 \code{-1} 作为哨兵：遇到左括号入栈；遇到右括号先弹出一个左括号位置，如果弹完后栈空，说明当前右括号无法匹配，把当前下标作为新边界入栈；否则当前有效长度就是 \code{i - stack.top()}。

为什么这个长度公式正确？栈顶保存的是当前有效后缀左侧最近的非法边界，或者尚未匹配的左括号位置。当前右括号匹配成功后，从栈顶之后到当前下标之间都是一个有效括号串，因此长度是差值。如果没有哨兵，边界会被很多特殊分支打散。DP 解法也可做，\code{dp[i]} 表示以 \code{i} 结尾的最长有效括号长度，但转移需要回看匹配左括号位置，初学阶段栈更直观。

实验要覆盖空串、全左括号、全右括号、简单嵌套、相邻有效段、非法右括号打断后重新开始。调试时打印栈内容和当前答案，特别观察 \code{)()())} 这类样例。正确性不是“栈能匹配括号”这么简单，而是栈顶始终保存当前有效后缀左侧的边界。

\noindent\textbf{LeetCode 41 缺失的第一个正数。}

排序能做，哈希集合也能做，但题目要求线性时间和常数额外空间。关键观察是：长度为 n 的数组，答案只可能在 \code{1..n+1} 之间。小于一和大于 n 的数不会影响前 n 个正整数是否存在；每个值 \code{x} 如果在 \code{1..n} 内，它的理想位置是下标 \code{x-1}。于是可以把数组本身当作哈希表，通过交换把每个合法值放到自己的位置。

循环时，对每个下标 \code{i}，只要 \code{nums[i]} 在 \code{1..n} 内，并且目标位置上的值不是它，就交换 \code{nums[i]} 和 \code{nums[nums[i]-1]}。这里必须检查目标位置是否已经有相同值，否则重复数字会导致无限交换。所有能放到正确位置的数处理完后，再从左到右找第一个 \code{nums[i] != i+1} 的位置，答案就是 \code{i+1}；如果都正确，答案是 \code{n+1}。

正确性来自位置映射。任何在 \code{1..n} 内出现过的值，只要经过交换循环，就会被移动到对应位置，除非该位置已经有相同值。无关值被留在某些位置，不影响其他合法值继续归位。实验覆盖空数组、全负数、含零、重复正数、已经连续、缺一、答案为 n+1 和极端乱序。这个题训练的是原地哈希，不是普通排序。

\noindent\textbf{LeetCode 51 N 皇后。}

暴力可以在棋盘每个格子上选择放或不放皇后，但这会产生巨大搜索树。第一层优化是逐行放置，因为合法解中每行必须恰好一个皇后；这样每层只需要选择列。再进一步，用三个集合判断冲突：列集合、主对角线集合和副对角线集合。对于位置 \code{row,col}，主对角线可用 \code{row-col} 标识，副对角线可用 \code{row+col} 标识。

回溯状态包括当前行号、已经放置的列、两类对角线和棋盘路径。进入某个选择时，把对应列和对角线加入集合；递归到下一行；返回时撤销。剪枝正确性来自皇后的攻击规则完全由同列和两类对角线覆盖。不要用“遍历已有皇后逐个检查”作为最终写法，它虽然正确，但每次检查重复扫描；集合把冲突判断降成常数时间。

实验可以统计搜索节点数。先写无剪枝版本或只检查列的版本，再加入对角线集合，比较 n 增大时节点数差异。边界包括 n 为一、n 为二和三无解、n 为四有两个解。输出棋盘时要注意路径拷贝成本：每次找到答案再构造字符串，或维护可修改棋盘并在答案处复制。这个题的重点是搜索树建模和剪枝证明。

\noindent\textbf{LeetCode 72 编辑距离。}

编辑距离是二维 DP 的代表。暴力递归从两个字符串末尾开始，如果末尾相同，就处理两个更短前缀；如果不同，就尝试删除、插入、替换三种操作。不同路径会反复处理同一对前缀长度，于是定义 \code{dp[i][j]}：把 \code{word1} 前 \code{i} 个字符变成 \code{word2} 前 \code{j} 个字符的最少操作数。

初始化来自空串。把前 \code{i} 个字符变成空串，需要删除 \code{i} 次；把空串变成前 \code{j} 个字符，需要插入 \code{j} 次。转移时，如果两个末尾字符相同，\code{dp[i][j] = dp[i-1][j-1]}；否则三种候选分别是删除 \code{word1} 末尾、向 \code{word1} 插入 \code{word2} 末尾、替换末尾，对应 \code{dp[i-1][j] + 1}、\code{dp[i][j-1] + 1}、\code{dp[i-1][j-1] + 1}。

正确性来自最后一步分类。任意最优编辑序列的最后一步必然是三种操作之一，或者末尾字符相同不需要操作；每种操作前的子问题都由更短前缀定义。实验要打印 DP 表，手动解释每个格子来源。若操作代价不同，转移权重要改；若只允许插入删除，则问题和最长公共子序列相关。空间压缩可以做，但初学时先写二维表，避免左上角旧值被覆盖。

\noindent\textbf{LeetCode 212 单词搜索 II。}

如果对每个单词单独在棋盘上 DFS，会重复搜索大量相同前缀。优化来自 Trie：把所有单词的前缀合并成一棵树，网格 DFS 时同步沿 Trie 前进。如果当前路径在 Trie 中已经没有对应子节点，就可以立即剪枝；如果 Trie 节点标记了完整单词，就把它加入答案。这样搜索空间由“单词数乘以棋盘搜索”变成“棋盘路径和词典前缀共同约束”。

状态包括网格位置、当前 Trie 节点和访问标记。进入一个格子前，先看当前 Trie 节点是否有该字符子节点；没有就返回。有则进入子节点，标记格子已访问，向四周扩展，退出时恢复访问标记。找到单词后，为避免重复加入，可以把节点上的单词标记清空；进一步优化还可以在子树为空时剪掉 Trie 叶子，但要小心节点生命周期。

正确性来自前缀共享。任意答案单词对应 Trie 中从根到终止节点的一条路径，网格 DFS 会枚举所有不重复格子的相邻路径；只有当前缀存在时才继续，不会漏掉可能单词。实验覆盖重复单词、同一单词多条路径、棋盘字符频次不足、单字符单词、访问标记恢复和 Trie 剪枝。C++ 实现中，固定小写字母可以用 26 个指针或下标数组；若字符集更大，哈希子节点更通用但常数更高。

\noindent\textbf{LeetCode 312 戳气球。}

这题最重要的转折是从“第一个戳谁”改成“最后戳谁”。如果按第一个戳的气球思考，左右邻居会随着后续操作变化，状态不稳定。若考虑某个开区间内最后被戳的气球 \code{mid}，那么在它最后被戳时，区间左右边界气球仍然存在，收益可以稳定写成 \code{nums[left] * nums[mid] * nums[right]}，左右两侧已经独立处理完。

令 \code{dp[left][right]} 表示开区间 \code{(left,right)} 内所有气球戳完的最大收益。枚举最后戳的 \code{mid}，转移为左区间收益加右区间收益加最后一次收益。为了处理边界，在数组两端补一。遍历顺序按区间长度从小到大，因为大区间依赖更短区间。这个状态定义看起来绕，但它解决了邻居变化的问题。

正确性来自最后一步分解。任意最优方案在区间内都有最后一个被戳的气球；在它之前，左右两侧气球的操作互不影响，因为最后时左右边界固定。枚举所有 \code{mid} 就覆盖全部最优方案可能的最后一步。实验覆盖一个气球、两个气球、包含一、较大值在中间、不同最后戳选择。这个题是区间 DP 的核心样本：只要正向操作让状态变化复杂，就尝试倒过来看最后一步。

\noindent\textbf{LeetCode 460 LFU 缓存。}

LFU 比 LRU 多一个频次维度。题目要求淘汰访问频率最低的 key；若频率相同，淘汰最久未使用的 key。暴力可以每次淘汰时扫描所有 key 找最低频和最旧时间，复杂度太高。要做到常数时间，需要两个索引：key 到节点信息的哈希表，频次到双向链表的哈希表。每个频次链表内部按最近使用顺序排列。

访问一个 key 时，节点频次从 f 变成 f+1，要从旧频次链表删除，再插入新频次链表头部。如果旧频次链表为空，并且它正好是当前最小频次，就把最小频次加一。插入新 key 时，频次为一，放入频次一的链表头部，最小频次重置为一。容量满时，从最小频次链表尾部淘汰最旧节点，并从 key 哈希表删除。

正确性来自两个不变量：key 表能定位每个节点；频次表中，每个链表只保存该频次节点，且从头到尾是新到旧顺序；\code{min_freq} 始终指向当前存在节点的最低频次。实验覆盖容量为零、容量为一、重复访问同一 key、更新已有 key、频次相同按 LRU 淘汰、旧频次链表清空。这个题非常接近工程索引维护：多个结构必须同步更新，任何一步漏掉都会产生悬空节点或错误淘汰。

\noindent\textbf{Linux 专题阅读路线。}

\noindent\textbf{第一周：链表和对象身份。}

第一周只读链表，不要贪多。目标是理解 \code{list_head} 为什么嵌进对象，\code{container_of} 为什么能从成员地址回到外层对象，遍历时删除为什么需要 safe 版本。读源码时画三张图：空链表哨兵、自身在链表中的节点、删除当前节点前保存下一个节点。然后写一个小实验：同一个对象带两个链表节点，分别挂入 LRU 链和脏链。实验报告说明删除其中一条链的节点不会自动从另一条链删除。

这一周要和力扣链表题对照。反转链表训练保存后继，K 个一组翻转训练分组边界，LRU 训练已知节点移动和尾部淘汰。内核链表把这些小技巧放到真实生命周期里：对象可能被多个模块引用，删除链表节点不等于释放对象。读懂这一点，再看内核其他数据结构就不会只停留在指针操作。

\noindent\textbf{第二周：哈希链和有序树。}

第二周读 hlist 和 rbtree。hlist 关注桶头空间、冲突链和已知节点删除；rbtree 关注动态有序索引和调用者比较。读 hlist 时，写一个小哈希表，比较普通单链表删除和保存前驱链接位置的删除。读 rbtree 时，不要先背全部修复情况，而是先看调用者如何找到插入位置，再看节点如何链接到父节点，最后理解旋转保持中序顺序。

对应刷题题型是哈希与区间。两数之和、前缀和计数、异位词分组训练 key 设计；日程安排、合并区间动态版、滑动窗口有序统计训练前驱后继。你要能说出什么时候哈希表不够：只要需要“小于等于某值的最大 key”“大于等于某值的最小 key”或按顺序扫描，哈希表就不是合适结构。

\noindent\textbf{第三周：XArray、Maple Tree 和范围。}

第三周读整数索引和范围索引。XArray 可以先用 Trie 类比：把整数 key 分段，沿层级定位对象，适合稀疏空间。Maple Tree 可以先用 \code{std::map} 区间表类比：给地址找覆盖区间，插入区间检查邻居，删除区间可能拆分。源码细节很多，但第一遍只回答三个问题：key 或范围如何组织，查找路径如何走，更新后旧节点如何处理。

对应刷题题型包括 Trie、坐标压缩、区间合并、插入区间、删除被覆盖区间。区别在于内核还要处理并发和对象生命周期。不要把 Maple Tree 简化成“高级红黑树”，它面向的是范围映射和读路径效率。读完这一周，应能写出一个简化 VMA 表，支持点查询、插入不重叠区间、删除子区间，并解释它缺少哪些内核级能力。

\noindent\textbf{第四周：RCU 和内存分配。}

第四周读 RCU、伙伴系统和 SLUB。RCU 的关键词是读侧轻量、写侧复制或替换、旧对象延迟释放。把删除拆成逻辑删除和物理释放：先从索引摘除，让新读者找不到；再等旧读者离开；最后释放。这个思想能帮助你理解为什么源码里看似“已经删掉”的对象还不能马上归还内存。

伙伴系统和 SLUB 则回答分配成本。伙伴系统关心连续页和阶数，SLUB 关心固定大小对象复用。对应刷题里的对象池、链表节点分配、哈希表节点常数、数组缓存局部性。第四周实验可以写一个读多写少配置表：读者只读全局指针，写者复制新表并发布旧表到延迟列表；再写简化伙伴分配器，观察总空闲量足够但连续块不足的情况。

\noindent\textbf{专题实验清单。}

\noindent\textbf{实验一：暴力与优化对拍框架。}

为每类题保留一个暴力版。数组和窗口题用小规模枚举；图题用小图全搜索或慢 BFS；DP 题用记忆化递归和表格互相校验；哈希题用平方枚举校验。对拍框架随机生成小输入，比较暴力版和优化版输出，一旦不同就打印输入、两版结果和关键状态。这个实验能极大减少“看起来对”的错觉。

对拍不是为了替代证明。它的作用是暴露边界，证明的作用是解释为什么所有边界都被覆盖。报告里要写：随机输入如何生成，哪些约束被覆盖，暴力版为什么可信，优化版错在哪里。比如和为 K 的子数组，要让随机数组包含负数和零；滑动窗口正数题，则要分别生成全正和含负两组，观察算法适用边界。

\noindent\textbf{实验二：复杂度曲线。}

选三道题画复杂度曲线：两数之和的平方枚举和哈希版，柱状图最大矩形的暴力扩展和单调栈版，岛屿数量的重复搜索和访问标记版。每个版本运行多个输入规模，记录耗时中位数。不要把打印放进计时代码，不要只运行一次，不要在 Debug 编译下得出结论。实验结果不必精确符合数学曲线，但增长趋势应能说明优化替换了哪一步。

复杂度曲线能训练工程判断。有时 \complexity{O(n \log n)} 的排序在小数据上比复杂的线性算法更快；有时哈希表因为 key 构造成本高而慢；有时链表理论删除快，但缓存局部性差导致整体慢。教材里的复杂度是第一层判断，实验让你看到常数和内存布局。

\noindent\textbf{实验三：边界测试表。}

为每个题型建立边界表。二分题固定测试空、一、二、边界外、重复；窗口题固定测试空串、全重复、目标缺失、刚好覆盖、窗口多次收缩；图题固定测试不连通、自环、重边、多源、资源维度；DP 题固定测试空状态、全负、不可达、空间压缩顺序。每次写题先从表里选用例，不要等错了再补。

边界表的价值在于复用思维，而不是复用代码。比如“重复元素”在三数之和里是去重问题，在二分里是边界问题，在堆里是比较器 tie-breaker 问题，在哈希里是频次问题。表格要记录同一个输入特征在不同题型中的含义，这样复习时才能横向迁移。

\noindent\textbf{实验四：源码阅读报告。}

每次读源码只回答一个小问题。例如读 \code{list_head} 时，只回答删除节点改了哪几条边；读 rbtree 时，只回答插入前调用者如何比较 key；读 XArray 时，只回答索引如何分层；读 RCU 时，只回答旧对象什么时候释放。报告必须包含文件路径、接口名、对象结构、不变量、和一个刷题类比。

不要把源码报告写成字段翻译。字段名会随版本变，设计问题更稳定。一个合格报告应写出：这个结构服务的访问模式是什么，为什么普通数组、哈希表或链表不够，插入和删除维护什么关系，并发或生命周期带来什么额外约束。把这四个问题写清楚，比背十个宏更有价值。

\noindent\textbf{面试讲解脚本。}

\noindent\textbf{从题目到方案。}

高质量口述可以固定成六句。第一句复述题面模型，不加算法名；第二句给暴力解，说明它如何覆盖所有候选；第三句指出瓶颈，说明哪一步重复；第四句给关键观察，说明某个状态可以复用或某些候选可以排除；第五句讲数据结构职责，说明保存、查询、更新、删除；第六句讲正确性和复杂度。这个顺序能防止直接背答案，也能让面试官看到推导过程。

以 LeetCode 560 为例：题面模型是区间和计数；暴力枚举所有左右端点；瓶颈是每个右端要找很多历史左端；关键观察是 \code{prefix[j] - prefix[i] == k}；数据结构是哈希表保存历史前缀和频次；正确性来自任意合法区间都对应一对前缀。这个脚本比“用前缀和加哈希表”更完整。

\noindent\textbf{面对追问。}

追问通常改一个条件。数组从全正改成含负，滑动窗口可能失效；返回一个答案改成返回所有答案，输出规模和去重会变化；只判断存在改成计数，哈希表 value 要从布尔变成频次；求最短改成求最长，保存最早位置或单调队列；普通树改成 BST，可以利用中序和上下界。面对追问先问条件变了哪里，再判断原状态是否仍然足够。

不要用“这个也差不多”回答变体。三数之和和四数之和相似，但四数之和要注意外层剪枝和 \code{std::int64_t}；接雨水和盛水最多相似处很少，一个计算每个位置贡献，一个选择两条边；零钱兑换 I 和 II 都是背包，但一个求最少数量，一个求组合数，循环语义不同。能指出相似题的不同点，比会背更多题更重要。

\noindent\textbf{讲代码质量。}

面试代码要能说明工程取舍。为什么用 \code{vector} 而不是 \code{list}，为什么哈希表要先查再插入，为什么堆里保存下标而不是值，为什么距离用 \code{std::int64_t}，为什么比较器这样写，为什么递归可能有栈风险。代码短不是唯一目标，状态清楚、边界自然、复杂度可信才是高质量。

如果写类，例如 LRU、并查集、Trie，要说明成员职责。并查集的 \code{parent} 保存代表关系，\code{size} 或 \code{rank} 控制合并；LRU 的哈希表负责定位，链表负责顺序；Trie 节点负责前缀转移和单词结束。类方法内部访问成员时保持一致风格，复杂逻辑拆成小函数，如 \code{move_to_front}、\code{remove_node}、\code{link_after}。这不仅让代码好看，也让不变量有固定位置。

\noindent\textbf{最后的学习路线。}

\noindent\textbf{第一轮：题型地图。}

第一轮目标不是做难题，而是建立题型地图。数组、窗口、哈希、二分、栈、堆、树、图、回溯、DP、贪心、并查集都要各做一组代表题。每题写出暴力、瓶颈、优化结构、复杂度和两个边界。第一轮可以看题解，但看完必须用自己的话写一页复盘。不会证明也没关系，但要知道证明缺口在哪里。

\noindent\textbf{第二轮：证明和边界。}

第二轮不追求新题数量，而是补证明。二分题证明排除一半，窗口题证明左端不会漏，单调栈证明弹出时答案确定，贪心题证明局部选择不劣，DP 题证明最后一步覆盖全部可能，图题证明首次到达或最短距离。每个题型至少准备三个反例，能解释错误代码为什么会被打破。

\noindent\textbf{第三轮：工程化表达。}

第三轮把代码改成工程上也能接受的版本。统一变量语义，处理整数溢出，避免不必要复制，确认容器 API 副作用，给复杂状态写辅助函数。对照 C++ 容器实验和 Linux 源码专题，重新审视每道题的数据结构选择。最终目标是：你不仅能 AC，还能解释为什么这个结构适合这个访问模式，为什么边界不会漏，为什么复杂度可信。

\noindent\textbf{高频综合题继续精读。}

\noindent\textbf{LeetCode 85 最大矩形。}

最大矩形是柱状图最大矩形的二维化。暴力枚举左上角和右下角，再检查矩形内是否全是一，复杂度会非常高；即使用二维前缀和把检查降成常数，枚举矩形仍然是四次方级别。真正的观察是把每一行当作矩形底边：如果当前格子是 \code{1}，它上方连续一的数量可以累加成柱高；如果当前格子是 \code{0}，柱高清零。于是每一行都转化成一组柱状图高度，对这组高度运行 LeetCode 84 的单调栈算法。

这个转化有两个关键点。第一，任何全一矩形都有一个底边行，当扫描到这行时，它会在柱状图中表现为一段连续高度至少为矩形高度的柱子；单调栈会在某根柱子作为最矮高度时结算出它。第二，遇到零必须把高度清零，因为以当前行为底的矩形不能穿过零。不要把每行高度重新向上扫描，那会把复杂度又拉高；高度数组应当逐行增量维护。

实验时先在纸上画一个小矩阵，逐行写出高度数组，再运行柱状图栈。边界包括全零、全一、单行、单列、零把高度截断、最大矩形不在最后一行。C++ 中矩阵字符常是 \code{'0'} 和 \code{'1'}，不要和整数零一混淆。讲复杂度时要写 \complexity{O(mn)}，因为每一行的柱状图算法是线性的。

\noindent\textbf{LeetCode 97 交错字符串。}

交错字符串的暴力递归非常直观：目标串当前位置可以来自第一个字符串，也可以来自第二个字符串，只要字符相等就递归尝试。问题是同一对下标 \code{i,j} 会被不同路径反复到达。状态定义应当保留两个来源串各用了多少字符：\code{dp[i][j]} 表示 \code{s1} 前 \code{i} 个字符和 \code{s2} 前 \code{j} 个字符，能否交错组成 \code{s3} 前 \code{i+j} 个字符。

转移来自最后一个字符。若 \code{s1[i-1] == s3[i+j-1]}，并且 \code{dp[i-1][j]} 为真，则当前状态可达；若 \code{s2[j-1] == s3[i+j-1]}，并且 \code{dp[i][j-1]} 为真，也可达。两个来源是或关系。开始前必须先判断长度和，否则无论 DP 表如何都不可能匹配。第一行和第一列分别表示只使用一个来源串的情况，遇到第一个不匹配后后面都不能再真。

这题最容易错在把“交错”误解成“保持两串各自相对顺序即可，但不要求用完全部字符”。题目要求刚好用完两个字符串并组成整个 \code{s3}。实验覆盖空串、一个来源为空、两个来源字符大量相同、两条路径都可行、长度不等、最后一个字符只能来自某一侧。空间可以压缩成一维，但压缩后要小心 \code{dp[j]} 代表上一行还是当前行。

\noindent\textbf{LeetCode 115 不同的子序列。}

这题和最长公共子序列很像，但目标不是长度，而是方案数。暴力递归从 \code{s} 中选择或不选择字符，判断能否组成 \code{t}，搜索树会非常大。状态定义为 \code{dp[i][j]}：在 \code{s} 的前 \code{i} 个字符中，有多少个子序列等于 \code{t} 的前 \code{j} 个字符。这里 \code{i} 是候选来源长度，\code{j} 是目标匹配长度，两者角色不能交换。

转移分两种。无论当前 \code{s[i-1]} 是否使用，都可以继承 \code{dp[i-1][j]}，表示不用当前字符；如果 \code{s[i-1] == t[j-1]}，还可以用它匹配目标末尾，增加 \code{dp[i-1][j-1]} 种方案。初始化 \code{dp[i][0] = 1}，因为任何前缀都能用“什么都不选”组成空目标；而 \code{dp[0][j>0] = 0}。

正确性来自最后一个来源字符是否被使用的划分，两类互不重叠且覆盖全部方案。方案数可能很大，C++ 中要根据题目范围考虑 \code{std::int64_t} 或无符号长整型。实验覆盖目标为空、来源为空、来源比目标短、重复字符大量产生多方案、完全无匹配。空间压缩时，\code{j} 必须倒序，否则同一个来源字符会在一轮中被重复使用。

\noindent\textbf{LeetCode 126 单词接龙 II。}

单词接龙 II 不只要求最短长度，还要求输出所有最短路径。暴力 DFS 枚举所有变换路径会陷入指数爆炸，并且容易先走到较长路径。正确思路是先用 BFS 确定最短层级，再用前驱关系回溯所有最短路径。BFS 的层次性保证第一次到达某个单词时是最短距离；但本题必须保留同一层的多个前驱，否则会漏掉另一条同样短的路径。

实现时可以建立通配模式桶，例如 \code{hot} 生成 \code{*ot}、\code{h*t}、\code{ho*}，把共享模式的单词放在同一桶里。这样找邻居时不需要和所有单词逐个比较。BFS 时维护 \code{distance} 和 \code{parents}，当发现一个未访问单词，把距离设为当前距离加一并记录父亲；当发现一个已经在下一层的单词，也要追加当前父亲。不能因为它已经被本层访问就丢掉同层前驱。

BFS 完成后，从终点按 \code{parents} 反向回溯到起点，路径反转后加入答案。边界包括终点不在词典、起点等于终点、多个最短路径、词典中存在大量共享模式、同一层访问标记时机。报告中要说明为什么不能边 BFS 边 DFS 输出路径：那样会混合搜索层级和路径枚举，容易保留非最短路径或漏掉同层父亲。

\noindent\textbf{LeetCode 139 单词拆分。}

单词拆分的暴力递归会枚举每个切分点：如果前缀在字典中，就递归处理后缀。重复后缀会被多次处理，例如很多不同前缀都能走到同一个剩余位置。状态可以定义为 \code{dp[i]}：字符串前 \code{i} 个字符能否被字典拼出。转移枚举最后一个单词的起点 \code{j}，如果 \code{dp[j]} 为真且 \code{s[j:i]} 在字典中，则 \code{dp[i]} 为真。

优化细节有两个。第一，字典放进哈希集合做等值查询，但 \code{substr} 会创建新字符串，热循环里可能增加成本；可以先写清楚版本，再用最大单词长度限制 \code{j} 范围，或用 Trie 避免大量切片。第二，\code{dp[0]} 必须为真，表示空前缀可以被拼出，这是所有第一段单词的起点。如果漏掉这个初始化，任何以开头单词开始的方案都无法转移。

正确性来自最后一个单词。任意合法拆分都有最后一段 \code{s[j:i]}，其前缀 \code{s[0:j]} 必须也可拆；枚举所有 \code{j} 就覆盖全部可能。实验覆盖空串、字典为空、重复单词、长字符串无法拆、多个拆法和最大单词长度剪枝。若题目改成输出所有句子，就不能只保存布尔值，需要在可达性基础上 DFS 枚举路径，并考虑输出规模可能指数级。

\noindent\textbf{LeetCode 188 买卖股票 IV。}

股票 IV 的关键是交易次数。暴力搜索每天买、卖或不操作，会产生大量状态。状态必须包含天数、已完成交易次数和是否持有股票。常见压缩写法维护 \code{buy[t]} 和 \code{sell[t]}：\code{buy[t]} 表示完成至多 \code{t} 次交易后持有股票的最大收益，\code{sell[t]} 表示完成至多 \code{t} 次交易后不持有股票的最大收益。买入不增加完成交易次数，卖出才完成一次交易。

若 \code{k >= n/2}，交易次数限制失去意义，问题退化为无限次交易，可以把所有上涨差价累加，或用两状态 DP。否则每天更新所有 \code{t}。转移时 \code{buy[t] = max(buy[t], sell[t-1] - price)}，\code{sell[t] = max(sell[t], buy[t] + price)}。更新顺序要防止同一天买卖污染状态；一种稳妥方式是使用昨天状态的副本，或者按经过验证的循环顺序更新。

边界包括空价格、只有一天、k 为零、k 很大、连续上涨、连续下跌、价格震荡。正确性来自状态穷尽：每天结束时，要么持有，要么不持有；完成交易次数从零到 k；所有合法操作都在买入、卖出、休息中。面试里要主动说清“交易次数”是约束维度，不是最后统计能补出来的东西。

\noindent\textbf{LeetCode 295 数据流中位数。}

如果每次加入数字后完整排序，插入成本太高。中位数只关心较小一半的最大值和较大一半的最小值，所以用两个堆维护两半。大顶堆保存较小一半，小顶堆保存较大一半；大小差不超过一；并且大顶堆堆顶不大于小顶堆堆顶。两个不变量同时成立时，中位数可以由堆顶直接得到。

插入时可以先根据堆顶把元素放入某一边，再做重平衡；也可以统一先放大顶堆，再把大顶堆堆顶移到小顶堆，最后根据大小移回。无论写法如何，都必须同时维护大小和平序关系。若只维护大小，可能较大元素在左堆；若只维护顺序，大小差过大时堆顶不再代表中位数。

实验覆盖第一个元素、偶数个元素、负数、重复值、升序输入、降序输入、平均值溢出。求偶数中位数时两个 int 相加可能溢出，应先转成 \code{std::int64_t} 或分别除以浮点。若题目扩展为滑动窗口中位数，还要支持删除过期元素，普通堆不支持任意删除，需要懒删除表或有序多重集合。

\noindent\textbf{LeetCode 329 矩阵中的最长递增路径。}

暴力从每个格子 DFS 寻找递增路径，会重复计算同一个后续格子的最长长度。由于路径必须严格递增，状态图是有向无环图：从低值指向高值，不可能回到原值或更低值形成环。可以定义 \code{memo[r][c]} 为从该格子出发的最长递增路径长度。DFS 访问邻居时，只走高度更大的格子，并取邻居结果加一的最大值。

记忆化的正确性来自子问题稳定：从某个格子出发，未来可走路径只由它的位置和矩阵值决定，与之前如何到达它无关。因此计算一次后可以复用。也可以把所有格子按值排序，从大到小或小到大做 DAG DP；本质都是利用严格递增带来的无环性。如果题目允许不下降，就可能出现相等值环，原证明失效。

实验覆盖单格、全相等、严格递增矩阵、严格递减矩阵、多个局部峰值、长蛇形路径。递归实现要注意深度，极端矩阵可能路径很长；工程上可考虑排序 DP 或显式栈。调试时打印 \code{memo} 表，检查每个格子是否只计算一次。

\noindent\textbf{LeetCode 394 字符串解码。}

字符串解码是解析类问题。暴力可以遇到右括号时向前寻找匹配左括号并展开，但嵌套结构会让查找和字符串拼接复杂。更自然的模型是栈保存进入新括号前的上下文：上一层已经构造的字符串、当前重复次数。扫描字符时，数字累积成多位数；遇到左括号，把当前字符串和次数入栈，重置当前字符串；遇到右括号，弹出上一层，把当前片段重复后接回去。

这里的状态必须区分“当前层字符串”和“上一层字符串”。如果只用一个字符串，遇到嵌套时无法知道展开后接回哪里。数字也要累积，例如 \code{12[a]} 不是 \code{1} 和 \code{2} 两个次数。题目通常保证输入合法，但教材实现仍应知道非法输入会在哪里出错：括号不匹配、数字后没有括号、重复次数为空等。

实验覆盖多位数字、嵌套、相邻编码段、普通字符混合、空片段。C++ 中大量字符串重复拼接可能成本高，可以先写清晰版本，再讨论用 \code{reserve} 或输出缓冲优化。这个题和表达式求值同属“上下文栈”：遇到进入嵌套结构时保存旧上下文，退出时恢复并合并。

\noindent\textbf{LeetCode 416 分割等和子集。}

题面问能否把数组分成和相等的两部分。先做代数转换：总和若为奇数，直接不可能；否则问题变成能否选出若干数，使其和为 \code{sum/2}。每个数只能使用一次，所以这是 0-1 背包可达性问题。状态 \code{dp[cap]} 表示是否能用已处理数字凑出容量 \code{cap}。

一维压缩时容量必须倒序遍历。倒序保证当前数字只会基于上一轮状态转移，最多使用一次；如果正序遍历，当前数字刚更新出的状态会在同一轮再次被使用，变成完全背包。这个循环方向不是模板差异，而是物品使用次数的语义。初始化 \code{dp[0] = true}，表示什么都不选可以凑出零。

实验要专门构造循环方向反例，例如只有一个数却能在正序中被重复使用凑出更大容量。边界包括总和为奇数、目标为零、数组中有零、数值很大导致容量表很大。复杂度是伪多项式 \complexity{O(n \cdot target)}，如果数值范围很大，即使元素个数不多也可能慢。

\noindent\textbf{LeetCode 494 目标和。}

每个数前面放正号或负号，暴力枚举所有符号是指数级。设正号集合和为 P，负号集合和为 N，总和为 S，则有 \code{P - N = target}，且 \code{P + N = S}，推出 \code{P = (S + target) / 2}。于是问题变成从数组中选若干数凑出 P 的方案数。这一步转换要先检查 \code{S + target} 是否非负且为偶数，否则无解。

转成背包后，状态 \code{dp[cap]} 表示凑出容量 \code{cap} 的方案数。每个数只能选择一次，所以容量倒序遍历。零元素会自然让方案数翻倍：当 \code{num = 0} 时，\code{dp[cap] += dp[cap]}，表示给零加正号或负号都不改变和，但属于不同符号方案。很多手写特殊分支反而容易错，标准计数转移可以自然处理。

正确性来自代数一一对应。任意符号方案对应一个正号集合，正号集合和必须是 P；任意选出和为 P 的集合，也能给集合内元素正号、集合外元素负号得到目标。实验覆盖目标绝对值大于总和、奇偶不匹配、含多个零、全零、空数组语义。它和分割等和子集很像，但一个求可达性，一个求方案数。

\noindent\textbf{LeetCode 518 零钱兑换 II。}

零钱兑换 II 问组合数，硬币可以无限使用。如果外层遍历金额、内层遍历硬币，会把不同顺序算作不同方案；例如一加二和二加一会被重复统计。组合数要求每组硬币按固定顺序被考虑，因此外层应遍历硬币，内层正序遍历金额。状态 \code{dp[value]} 表示使用已经处理过的硬币种类，凑出金额 \code{value} 的组合数。

初始化 \code{dp[0] = 1}，表示凑出金额零有一种方案：什么都不选。处理某个硬币 \code{coin} 时，金额从 \code{coin} 到 \code{amount} 正序遍历，\code{dp[value] += dp[value - coin]}。正序表示当前硬币可以重复使用；外层硬币顺序固定表示组合不会因为排列不同重复计数。

实验要把循环顺序故意写反，观察组合数如何变成排列数。边界包括金额为零、硬币数组为空、硬币面额大于金额、组合数较大。和 LeetCode 322 的区别必须说清：322 求最少硬币数，518 求组合数量；一个是最值 DP，一个是计数 DP；循环语义和初始化都不同。

\noindent\textbf{LeetCode 621 任务调度器。}

任务调度器可以用堆模拟，也可以用公式。暴力模拟每个时间片，选择当前可执行且剩余次数最多的任务，不够时插入空闲。堆模拟比较直观，但公式能揭示结构：最短时间主要由最高频任务决定。设最高频为 \code{maxFreq}，有 \code{maxCount} 种任务达到最高频。前 \code{maxFreq - 1} 轮需要形成骨架，每轮长度至少为 \code{n + 1}，最后放所有最高频任务的尾部。

公式为 \code{max(tasks.size(), (maxFreq - 1) * (n + 1) + maxCount)}。为什么还要和任务总数取最大？因为其他任务足够多时，它们可以填满所有冷却空位，甚至让整体没有空闲，此时总时间就是任务数。若其他任务不够，骨架长度决定必须插入多少空闲。这个证明比记公式更重要。

实验覆盖冷却为零、只有一种任务、多个最高频任务、其他任务刚好填满空隙、其他任务超过空隙。堆模拟可以作为验证公式的对拍版本。若题目扩展为每个任务有不同执行时间、释放时间或不同冷却间隔，公式不再适用，需要更一般的调度模型。

\noindent\textbf{LeetCode 815 公交路线。}

公交路线题的建模很关键。如果把站点作为 BFS 节点，边表示同一公交路线内任意两站可达，那么一条路线内站点两两连边会非常多。更合适的模型是把“乘坐路线”作为 BFS 层级的一部分：从起点站找到所有可乘路线，乘上一条路线算一次公交；这条路线经过的所有站点都可以到达，并继续找到这些站点可换乘的其他路线。

实现通常建立站点到路线列表的哈希表。BFS 队列里可以放路线，也可以放站点加换乘次数；无论如何，都要分别标记访问过的路线和站点，避免重复展开。同一条路线一旦处理过，就不必再次扫描它的所有站点。目标站如果在当前路线中出现，就可以返回当前乘车次数。

正确性来自 BFS 层次：每一层代表乘坐公交次数增加一，第一次到达目标站就是最少乘车次数。实验覆盖起点等于终点、目标不可达、多条路线共享站点、路线很长、站点编号很大且稀疏。容器选择上，站点编号不一定连续，站点到路线列表适合哈希表；路线访问适合布尔数组，因为路线数量连续。

\noindent\textbf{LeetCode 895 最大频率栈。}

最大频率栈要求弹出出现频率最高的元素；频率相同，弹出最近压入的元素。暴力每次 pop 扫描所有元素找最高频和最近时间，太慢。需要两个维度的索引：value 到频次的哈希表，频次到栈的哈希表，以及当前最大频次 \code{maxFreq}。push 某个值时，先把它的频次加一，再压入对应频次的栈，并更新最大频次。

pop 时，从 \code{maxFreq} 对应的栈弹出栈顶值。由于每次达到某个频次时都按时间压入该频次栈，栈顶自然是该频次下最近的元素。弹出后把该值频次减一；如果最大频次栈为空，\code{maxFreq} 减一。注意，不需要从较低频次栈中删除旧记录，因为那些记录代表该元素历史上达到该频次的时刻；只有当它再次成为当前最大频次栈顶时才会被使用。

正确性来自频次分层和栈内时间顺序。哈希表给出当前真实频次，频次栈保存达到该频次的时间顺序，\code{maxFreq} 指向最高非空层。实验覆盖频次并列、同一元素连续 push、pop 后最大频次下降、多个历史记录存在。这个题能训练“按维度拆索引”：一个结构回答频次，一个结构回答同频最近。

\noindent\textbf{LeetCode 981 基于时间的键值存储。}

题目要求对每个 key 保存多个时间戳版本，并查询不超过给定时间戳的最新值。暴力可以把所有记录放到一个列表里，每次查询扫描；更好的模型是按 key 分组。每个 key 对应一个按时间递增的数组，元素是时间戳和值。查询时先用哈希表定位 key，再在该 key 的时间数组中二分找第一个大于目标时间的位置，答案是它前一个元素。

这个结构利用了题目常见条件：同一个 key 的 set 调用时间戳递增。如果没有这个条件，就需要插入时保持有序，可能使用有序容器或插入后二分位置。二分边界要处理目标时间小于该 key 最早时间，此时返回空字符串；目标时间大于所有记录，返回最后一个值。不要把所有 key 的记录混在一个全局有序数组里，否则查询某个 key 时还要过滤，复杂度和实现都变差。

实验覆盖 key 不存在、时间戳早于所有记录、等于某个记录、落在两个记录之间、晚于所有记录、多个 key 交错 set。C++ 中可以用 \code{unordered_map<string, vector<pair<int,string>>>}，二分时比较 pair 的第一维，或者手写二分避免构造哨兵 pair。这个题的重点是“先按 key 分组，再在组内二分”。

\noindent\textbf{LeetCode 1235 规划兼职工作。}

每个工作有开始时间、结束时间和收益，目标是选择互不重叠工作最大化收益。暴力枚举子集会指数爆炸。排序后可以做 DP：按结束时间排序，\code{dp[i]} 表示考虑前 \code{i} 个工作能获得的最大收益。第 \code{i} 个工作要么不选，收益 \code{dp[i-1]}；要么选，则需要找到最后一个结束时间不大于当前开始时间的工作 \code{p}，收益为 \code{dp[p] + profit[i]}。

二分用于寻找 \code{p}。因为工作按结束时间排序，所有结束时间不大于当前开始时间的工作构成前缀，可以用 \code{upper_bound} 找边界。这里的二分不是查目标值是否存在，而是找兼容工作的数量。若端点语义是前一个工作结束时间等于当前开始时间可兼容，比较条件应使用不大于。

正确性来自最后一个决策：最优解对当前工作只有选和不选两类，选当前工作后，之前能选择的工作只来自兼容前缀，且这部分最优值已经在 DP 中。实验覆盖工作完全不重叠、完全重叠、端点相接、相同结束时间、收益很高但阻塞多个小工作。不要用按收益贪心，局部最高收益可能破坏全局组合。

\noindent\textbf{LeetCode 1438 绝对差不超过限制的最长连续子数组。}

窗口合法性由最大值和最小值决定：窗口内最大值减最小值不超过限制。暴力每个窗口重新找最大最小会很慢；平衡树可以维护动态有序集合，复杂度 \complexity{O(n \log n)}；单调队列可以做到线性。维护两个队列：一个递减队列保存最大值候选，一个递增队列保存最小值候选，队列里存下标以便判断过期。

右端加入新元素时，最大队列弹出所有小于新值的尾部，因为它们既更小又更早，不可能再成为最大值；最小队列弹出所有大于新值的尾部。然后检查队首差值是否超过限制；若超过，就移动左端，并把过期下标从队首弹出。注意，左端可能需要连续移动，直到窗口重新合法。

正确性来自支配关系和窗口单调收缩。队首始终是当前窗口最大和最小；被队尾弹出的旧候选被新候选支配，不会在未来窗口中重新有用。实验覆盖限制为零、全相等、严格递增、严格递减、需要多次收缩、最大最小交替变化。这个题说明滑动窗口状态不一定是和或频次，也可以是动态极值。

\noindent\textbf{LeetCode 1514 概率最大的路径。}

这题是最短路思想的变体。路径概率是边概率乘积，目标是最大乘积。暴力枚举路径不可行；因为边概率在零到一之间，沿路径继续乘只会不增，所以可以用 Dijkstra 类似思想：每次扩展当前概率最大的节点。用大顶堆保存当前到达某节点的最大概率候选，若弹出的概率不是当前记录值，就跳过过期状态。

松弛边时，候选概率是 \code{prob[u] * edgeProb}。如果它大于 \code{prob[v]}，就更新并压入堆。也可以把概率取负对数，把乘积最大转成权重和最小；但直接最大堆更贴近题意。正确性来自单调性：当前堆顶概率最高，未来从其他更低概率路径继续乘边，不可能变成更高概率回头改写它。

实验覆盖起点等于终点、不可达、概率为零、多个同概率路径、浮点比较。C++ 中要使用 \code{double}，并避免用整数处理概率。这个题适合训练“框架迁移”：不是所有 Dijkstra 都叫最短距离，只要代价扩展满足合适单调性，优先队列扩展思想就可能成立。

\noindent\textbf{LeetCode 1675 数组的最小偏移量。}

每个数可以做有限变化：奇数可以乘二一次变成偶数，偶数可以反复除二。直接搜索所有状态组合会爆炸。关键是先把每个数规范化到它能达到的最大偶数形态：奇数乘二，偶数保持。之后每个数只允许不断减小。维护当前所有数中的最大值和最小值，每次尝试降低当前最大偶数，因为只有降低最大值才可能缩小偏移。

用大顶堆保存当前值，同时维护当前最小值。每次堆顶是最大值，先用最大值减最小值更新答案；如果最大值是偶数，可以除二后放回堆，并更新最小值；如果最大值是奇数，它已经不能再降低，过程结束。为什么只处理最大值？偏移由最大和最小决定，降低非最大值不会降低当前最大，反而可能降低最小使偏移变大。

正确性来自状态空间方向。规范化后，每个数的可达序列都是单调下降的；从当前最大值开始下降，枚举了所有可能改善最大边界的步骤。当最大值无法继续下降时，当前及未来都无法缩小最大边界。实验覆盖全奇数、全偶数、已经偏移为零、最大值多次下降、下降后成为新的最小值。注意值乘二可能溢出，要根据输入范围使用足够整数类型。

\noindent\textbf{跨题型迁移和混淆点。}

\noindent\textbf{盛水最多和接雨水不能混用。}

LeetCode 11 盛最多水的容器和 LeetCode 42 接雨水都出现左右指针和高度数组，但问题模型完全不同。盛水最多只选择两条边，面积由短板和宽度决定；每次移动短板，是因为移动长板无法提高短板且宽度必然变小。接雨水则计算每个位置的水量，位置贡献由左右最高边界的较小值决定。前者在候选对之间做排除，后者在每个位置上做边界结算。

如果把盛水最多的证明搬到接雨水，会漏掉中间位置的贡献；如果把接雨水的左右最高思路搬到盛水最多，也会计算不存在的“内部蓄水”。复盘时要写两张状态表：盛水最多保存的是当前左右边界和最大面积，接雨水保存的是左侧最高、右侧最高和当前位置贡献。它们的共同点只是都利用了高度边界，不能因为代码都有双指针就当成同一模板。

\noindent\textbf{最小覆盖子串和找到所有异位词。}

LeetCode 76 最小覆盖子串是可变长度窗口，窗口满足覆盖后要尽量收缩；LeetCode 438 找到字符串中所有字母异位词是固定长度窗口，只需要维护长度等于目标串长度的频次。两者都使用字符计数，但窗口边界语义不同。最小覆盖子串的左端移动由“覆盖条件仍然满足”驱动；异位词窗口的左端移动由“窗口长度固定”驱动。

混用这两题会出现典型错误。把固定窗口写成可变窗口，可能在频次超过需求时过度收缩，漏掉长度刚好的异位词；把可变窗口写成固定窗口，则无法找到更短覆盖。实验时分别打印窗口长度、字符差异和答案更新点。讲题时要先说输出目标：一个要求最短子串，一个要求所有起点；输出目标不同，状态读取位置也不同。

\noindent\textbf{普通 BFS、Dijkstra 和 0-1 BFS。}

看到“最短”不能直接写 BFS。普通 BFS 成立的条件是每条边代价相同，队列层数就是路径长度。Dijkstra 成立的条件是边权非负，每次扩展当前代价最小的未确定状态。0-1 BFS 则针对边权只有零和一的图，用双端队列把零权边放队首、一权边放队尾，从而保持距离顺序。三者都是按某种代价顺序扩展，但容器和证明不同。

网络延迟时间有不同正权边，普通 BFS 会按边数而不是时间扩展，可能错；打开转盘锁每次拨动代价相同，普通 BFS 正合适；带零一代价的状态图若用普通队列，会破坏距离顺序，若用普通 Dijkstra 又能做但常数更高。复盘图题时先写边权范围，再选队列、双端队列或堆。这个判断比背算法名字重要。

\noindent\textbf{并查集和 DFS 连通块。}

岛屿数量、朋友圈、省份数量、冗余连接都和连通性有关，但并查集和 DFS 的职责不同。DFS/BFS 适合从一个起点找完整连通块，能顺便统计面积、边界、路径等信息。并查集适合动态合并和快速回答两个元素是否同组，但它不保存具体路径，也不擅长枚举某个连通块的所有节点，除非额外维护集合成员。

冗余连接按边顺序加入，发现两个端点已经连通时，这条边形成环，因此并查集非常合适。岛屿最大面积需要统计每个连通块大小，BFS/DFS 更直接；当然也可以用并查集维护 size，但实现成本更高。选择结构时先问查询是什么：要遍历块，还是只问同组；要动态加边，还是静态扫描；要路径，还是只要集合代表。

\noindent\textbf{背包三题的状态区别。}

分割等和子集、目标和、零钱兑换 II 都能归到背包，但状态含义不同。分割等和是可达性，\code{dp[cap]} 是布尔值；目标和是方案数，\code{dp[cap]} 是计数；零钱兑换 II 是完全背包组合计数，硬币可无限使用且组合不区分顺序。循环方向和外层顺序都由这些语义决定。

若把分割等和的布尔状态搬到目标和，就只能知道是否存在，无法统计方案；若把零钱兑换 II 的正序容量搬到 0-1 背包，就会重复使用同一个物品；若把组合计数的循环顺序反过来，就会统计排列数。背包题不要背“倒序或正序”，而要先写一句：物品能用几次，答案是可达、最值还是计数，顺序是否有意义。

\noindent\textbf{源码思想迁移到题解表达。}

\noindent\textbf{对象生命周期。}

Linux 源码读多以后，刷题里也会更重视对象生命周期。链表节点是否由题目平台管理，返回的新链表是否复用输入节点，哈希表保存的是值、下标、迭代器还是指针，vector 扩容后旧引用是否还有效，这些都属于生命周期问题。很多面试题不会直接考内存释放，但会通过 LRU、链表反转、合并链表、复制随机链表等题考察你是否理解节点身份。

题解表达中可以主动说明副作用。例如“本题允许修改输入，所以我把访问过的陆地改成水”；“本题不应释放输入链表节点，只需要重连”；“LRU 的哈希表保存链表迭代器，链表节点必须稳定，所以不能用会移动元素的 vector 存节点”。这些话能把算法答案提升到工程答案。

\noindent\textbf{索引组合。}

内核对象常同时进入多个索引：一个对象可能在链表中表示顺序，在树中表示有序键，在哈希表中表示快速查找。刷题中的很多设计题也是索引组合。LRU 是 key 到节点的哈希表加新旧顺序链表；LFU 是 key 到节点、频次到链表、最小频次；时间键值存储是 key 到有序时间数组；账户合并是邮箱到集合编号再按根收集。

讲设计题时，不要只列容器，要说明每个索引回答什么问题。哈希表回答“这个 key 在哪里”，链表回答“谁最旧”，堆回答“当前极值是谁”，有序表回答“前驱后继是谁”，并查集回答“两个对象是否同组”。一个容器若回答不了所有操作，就需要组合；组合后每次更新都要维护所有索引的一致性。

\noindent\textbf{逻辑删除和物理删除。}

RCU、堆懒删除、滑动窗口过期元素、哈希表延迟清理都体现了一个思想：逻辑上不可用不等于物理上马上移除。Dijkstra 堆里的旧状态逻辑上过期，但直到弹出才物理删除；滑动窗口中位数里，过期元素可能埋在堆中，需要延迟清理；RCU 中，旧对象从新索引不可达后，还要等旧读者离开才能释放。

题解里要说清使用前检查。堆顶如果可能过期，就在读取答案前清理；哈希表如果保存历史频次，就要明确什么时候递减或删除；链表节点如果从一个索引摘除，还要确认其他索引是否仍引用它。逻辑删除能降低操作成本，但必须有可靠的清理时机，否则会读到脏状态。

\noindent\textbf{更多实验报告样例。}

\noindent\textbf{样例一：二分答案。}

题目选择 LeetCode 1011 在 D 天内送达包裹。暴力基线是从容量下界开始逐个尝试，检查能否在 D 天内完成。瓶颈是答案空间可能很大，线性试容量太慢。关键观察是容量越大，需要天数越少，可行性单调。检查函数按顺序装包，当前天能放就放，放不下再开新天。这个贪心正确，因为提前开新天只会让剩余天数更紧，不会减少所需天数。

边界测试包括 D 为一、D 等于包裹数、单个包裹等于最大值、容量小于最大包裹、所有包裹相同。实验记录 \code{left}、\code{right}、\code{mid} 和需要天数，检查每次排除是否符合“最小可行容量”语义。C++ 注意总重量可能较大，用足够整数类型。最终复杂度是 \complexity{O(n \log S)}，S 是容量范围。

\noindent\textbf{样例二：单调队列。}

题目选择 LeetCode 239 滑动窗口最大值。暴力基线是每个窗口扫描 K 个元素取最大，复杂度 \complexity{O(nk)}。优化观察是：若新元素比队尾候选更大且下标更晚，那么旧候选在未来任何同时包含二者的窗口中都不可能成为最大值；新元素既更大又更持久。单调队列保存下标，队首是当前窗口最大值。

边界测试包括 K 为一、K 等于数组长度、严格递增、严格递减、重复最大值。调试时打印右端进入、队尾弹出、队首过期和答案读取。证明分两步：队尾弹出不丢答案，因为旧候选被新候选支配；队首过期必须删除，因为它不再属于窗口。C++ 用 \code{deque<int>} 保存下标，而不是保存值，因为需要判断过期。

\noindent\textbf{样例三：树形 DP。}

题目选择 LeetCode 337 打家劫舍 III。暴力递归对每个节点尝试偷或不偷，会重复计算子树。状态应由子树向父节点汇报两个值：偷当前节点的最大收益、不偷当前节点的最大收益。若偷当前节点，孩子不能偷；若不偷当前节点，孩子可以偷也可以不偷，取各自最大值。

边界包括空树、单节点、链式树、左右子树收益差异大。正确性来自树形后序归纳：左右子树已经给出在偷或不偷根时的最优结果，父节点只按相邻不能同偷的约束合并。工程上递归清晰，但极端深树有栈风险；可用显式栈做后序遍历并用哈希表保存每个节点状态。这个实验把 DP 状态和树遍历顺序联系起来。

\noindent\textbf{样例四：区间贪心。}

题目选择 LeetCode 452 用最少数量的箭引爆气球。暴力可以尝试所有射箭位置组合，但选择空间巨大。贪心观察是：按右端点排序后，第一支箭放在最早结束气球的右端，不会比任何其他能射爆这个气球的位置更差，因为它给后续气球留下最大可能覆盖范围。然后跳过所有被这支箭覆盖的气球，继续处理下一个未覆盖气球。

正确性用交换论证。任意最优解中，必须有一支箭射爆最早结束的气球；把这支箭移动到该气球右端，仍能射爆原来被它覆盖且起点不大于该右端的气球，并且不会更早结束。因此存在一个最优解采用贪心第一步。边界包括端点相同、负坐标、区间包含、闭区间端点相接。比较器按右端升序，注意端点可能接近整数边界。

\noindent\textbf{全书复盘索引。}

\noindent\textbf{题型到能力。}

数组训练区间语义和缓存局部性；窗口训练增量维护和单调条件；哈希训练等价类和历史状态；二分训练边界谓词；栈训练最近未完成结构；堆训练动态极值；树训练状态方向；图训练建模和访问标记；DP 训练状态定义和依赖顺序；贪心训练交换证明；并查集训练动态连通；Linux 源码训练对象生命周期和索引组合。

复习时不要按题号孤立记忆，而要问每题训练哪种能力。LeetCode 560 训练前缀代数和哈希频次；LeetCode 862 训练前缀和上的单调队列；LeetCode 1631 训练同一问题的多模型；LeetCode 146 训练索引组合；LeetCode 312 训练区间 DP 的最后一步。把能力写在题号旁边，复习效率会明显高于只刷通过率。

\noindent\textbf{源码到能力。}

Linux 链表训练侵入式节点和稳定身份；hlist 训练桶和删除成本；rbtree 训练动态有序索引；XArray 训练稀疏整数映射；Maple Tree 训练范围查询；RCU 训练逻辑删除和物理释放分离；伙伴系统训练连续性、阶数和合并；SLUB 训练小对象分配成本。它们不是为了替代刷题算法，而是帮助你把数据结构选择讲到工程约束层面。

最终目标是形成统一语言：访问模式决定结构，结构维护不变量，不变量支撑正确性，容器实现带来工程成本，边界实验验证理解。无论面对力扣题、C++ 面试题还是 Linux 源码阅读，都按这条线分析，就不会退回“复制模板”的学习方式。

\noindent\textbf{代码审阅清单。}

\noindent\textbf{状态是否真的必要。}

每写完一道题，先检查状态是否最小且足够。最小是指没有保存未来不会再用的信息；足够是指未来决策所需的信息没有缺失。两数之和只需要历史值到下标，不需要保存所有数对；和为 K 的子数组计数需要历史前缀频次，不是只保存出现与否；最长区间需要最早位置，不是最后位置。状态一旦不匹配输出目标，代码可能在简单样例上通过，却在变体或边界上失败。

审阅时可以问四个问题：状态何时进入，何时更新，何时删除，何时读取答案。窗口题中，右端进入、左端删除和答案读取顺序非常关键；堆题中，堆顶读取前可能需要清理过期元素；DP 题中，状态更新顺序决定读到旧值还是新值。只要这四个问题有一个答不上来，就不应该把题标记为掌握。

\noindent\textbf{复杂度是否按真实维度说明。}

复杂度不能只写一个 \code{n}。图题要区分节点数和边数，字符串 DP 要区分两个字符串长度，背包题要说明容量来自数值大小，答案二分要说明对答案范围取对数，输出所有路径的题要把输出规模算进去。比如单词接龙 II 的 BFS 只是找到最短层级，最终输出所有路径可能本身就很大；N 皇后输出数量也是复杂度的一部分。

审阅复杂度时还要看容器操作。哈希表平均常数不等于最坏常数，\code{substr} 可能复制字符串，\code{priority_queue} 插入和弹出是对数，\code{map} 查询前驱后继也是对数。若代码在循环里构造大量临时字符串，表面 DP 复杂度可能被隐藏成本放大。复杂度说明要和实际代码一致。

\noindent\textbf{边界是否进入主逻辑。}

好的代码不是主体一段、后面一堆补丁分支。边界最好进入初始化和不变量。前缀和把 \code{0 -> 1} 放进哈希表，就自然处理从开头开始的区间；括号栈放 \code{-1} 哨兵，就自然处理第一个有效段；链表使用虚拟头节点，就自然处理删除头节点；二分使用左闭右开，就自然处理插入到末尾。审阅时要优先寻找能消除特殊分支的状态设计。

当然，有些入口边界应当直接返回，例如空数组、起点终点被阻塞、容量为零、目标长度不匹配。这些判断不是补丁，而是题意的不可行条件。区别在于：入口判断处理的是根本无解或无输入；不变量处理的是算法过程中的普通边界。把二者混在一起，会让代码难以解释。

\noindent\textbf{C++ 实现是否隐藏风险。}

C++ 审阅至少看五类问题。第一，整数溢出：面积、路径距离、前缀和、二分平方、概率转换都可能超出 \code{int}。第二，迭代器和引用失效：\code{vector} 扩容、\code{erase}、\code{unordered_map} rehash 都可能让保存的引用失效。第三，容器副作用：\code{operator[]} 会插入默认值。第四，比较器：排序和堆比较器必须满足严格弱序。第五，复制成本：循环中不要无意复制大 \code{vector}、\code{string} 或复杂节点。审阅必须具体。

设计类时还要看成员职责。并查集的 \code{find} 是否路径压缩，\code{unite} 是否按大小合并；LRU 的移动节点是否只改链表而不丢哈希；Trie 的析构和节点所有权是否清楚；图算法的邻接表是否按输入规模预分配。刷题代码不一定要写成工程库，但高质量教材必须告诉读者这些风险在哪里。

\noindent\textbf{十周训练安排。}

\noindent\textbf{第一、二周：线性结构。}

第一周集中数组、字符串、双指针和窗口。每天选择两道基础题和一道变体题，必须写暴力版和优化版。重点不是数量，而是把区间语义写清楚：左闭右开还是左闭右闭，慢指针左侧保存什么，窗口何时合法，答案何时更新。实验必须打印指针移动过程。推荐题包括两数之和、无重复字符最长子串、盛水最多、三数之和、最小覆盖子串、长度最小子数组、找到所有异位词。

第二周集中链表、栈、队列和单调结构。链表题每一步都画前驱、当前、后继；栈题说明栈顶代表什么；单调栈和单调队列说明支配关系。推荐题包括反转链表、K 个一组翻转链表、有效括号、最长有效括号、每日温度、柱状图最大矩形、滑动窗口最大值。周末写一个 LRU，把哈希表和双向链表的职责讲清楚。

\noindent\textbf{第三、四周：查找和有序结构。}

第三周集中哈希、前缀和和二分。哈希题必须写 key 和 value 的语义，前缀题必须写代数等式，二分题必须写谓词和区间不变量。推荐题包括字母异位词分组、和为 K 的子数组、连续子数组和、连续数组、搜索插入位置、查找首尾位置、爱吃香蕉、船运包裹、分割数组最大值。

第四周集中排序、区间、堆和有序容器。排序题说明排序换来了什么关系，区间题说明端点闭开，堆题说明堆顶是否可能过期，有序容器题说明前驱后继。推荐题包括合并区间、插入区间、无重叠区间、用箭引爆气球、前 K 高频、数据流中位数、最接近原点 K 个点、任务调度器。周末写 map 动态区间实验，对照 Maple Tree 思维。

\noindent\textbf{第五、六周：树和图。}

第五周集中二叉树、BST、Trie。每道树题先判断状态方向：自顶向下、后序汇总还是层序传播。递归写法要能说出迭代替代方案。推荐题包括层序遍历、验证 BST、最大深度、最近公共祖先、二叉树直径、打家劫舍 III、单词搜索 II、实现 Trie。实验打印显式栈和队列状态。

第六周集中图、拓扑、并查集和最短路。每道图题先写节点、边、边权和状态维度。推荐题包括岛屿数量、岛屿最大面积、腐烂橘子、打开转盘锁、课程表、省份数量、冗余连接、网络延迟时间、最小体力路径、概率最大路径。周末比较 BFS、Dijkstra、二分 BFS 和并查集在同一题上的建模差异。

\noindent\textbf{第七、八周：动态规划。}

第七周集中线性 DP、二维 DP 和字符串 DP。每题先写暴力递归，再圈出重复状态。推荐题包括爬楼梯、打家劫舍、最大子数组和、不同路径、最小路径和、最长公共子序列、编辑距离、交错字符串、不同子序列。每题至少打印一个小 DP 表，解释每个格子来源。

第八周集中背包、区间 DP、状态机和状态压缩。推荐题包括零钱兑换、零钱兑换 II、分割等和子集、目标和、买卖股票含冷冻期、买卖股票 IV、戳气球、最长递增子序列。重点检查循环方向、初始化、空间压缩旧值来源和伪多项式复杂度。周末把三道背包题放在一张表里比较状态含义。

\noindent\textbf{第九、十周：综合和源码。}

第九周做综合设计题和高频难题。推荐题包括 LFU 缓存、最大频率栈、时间键值存储、公交路线、最短子数组至少为 K、数组最小偏移量、规划兼职工作。每题都写索引组合图，说明每个结构回答什么查询，更新时如何保持一致。

第十周读源码和复盘。按链表、hlist、rbtree、XArray、Maple Tree、RCU、伙伴系统、SLUB 的顺序，每天只读一个小问题，写一页报告。最后选十道自己最容易混淆的题，重新讲一遍暴力、瓶颈、状态、不变量、边界和 C++ 风险。十周结束后，不要求记住所有题号，但要能看到题面就判断访问模式、状态维度和可疑边界。

\noindent\textbf{质量底线。}

\noindent\textbf{不能用模板替代理解。}

任何题解如果只有“这题用哈希”“这题用 DP”“这题用双指针”，都不合格。必须补上为什么：哈希保存什么等价类，DP 状态包含哪些未来信息，双指针排除了哪一批候选。模板只能作为复习索引，不能作为解释本身。读者应该在每道题里看到暴力版本如何演化，而不是看到一个突然出现的最终答案。

教材写作也一样。章节之间要有衔接：数组的区间语义过渡到窗口，窗口的增量维护过渡到单调队列，前缀和的代数关系过渡到哈希历史状态，树的后序汇总过渡到树形 DP，图的层次扩展过渡到最短路。只有这样，书才像教材，而不是题解堆。

\noindent\textbf{不能用字数掩盖空洞。}

字数目标只是底线，不是质量本身。有效内容应当增加背景、推导、证明、边界、实验或源码对照；无效内容只是换词重复同一句话。审阅时要删掉不能回答新问题的段落。一个段落如果不能说明“为什么这样做”“错在哪里”“边界如何处理”“和另一个结构有什么区别”，就应该重写。

本书后续继续扩展时，应优先补三类内容：第一，更多经典题的完整推导；第二，真实 C++ 实验和性能对拍；第三，Linux 源码阅读报告。它们能自然增加深度，不需要硬凑。最终读者拿到的应是一套可执行的学习系统：能刷题、能讲题、能写代码、能读源码。

\noindent\textbf{口述和代码审查样例。}

\noindent\textbf{样例一：讲清前缀和。}

以前缀和题为例，合格口述不能停在“我用前缀和”。应当这样讲：我先把任意区间和写成两个前缀的差，因此当前右端固定时，问题变成寻找某个历史前缀。若题目问数量，历史表保存频次；若题目问最长，历史表保存最早位置；若题目问是否存在，保存布尔即可。这个区分直接决定代码里的 value 类型。

审查代码时，第一看初始化是否包含前缀零；第二看先查询还是先更新；第三看是否使用足够整数类型；第四看负数是否影响算法。和为 K 的子数组不需要数组全正，因为它只依赖代数等式；长度最小且至少为 K 的问题若含负数，普通窗口失效，要换成前缀和加单调队列。能说清这两个题的差异，就说明前缀和不是死记公式。

\noindent\textbf{样例二：讲清二分。}

二分题的口述要先定义谓词。比如船运包裹，谓词是“容量为 cap 时，能否在 D 天内按顺序运完”。这个谓词随 cap 增大从假变真，因此要找第一个真。区间语义可以固定为左闭右开：答案始终在 \code{[left,right)} 内，若 \code{mid} 可行，就把右边界收到 \code{mid}；否则把左边界移到 \code{mid+1}。

代码审查时，不要只看循环是否会停。要看下界是否为最大包裹，上界是否为总重量，检查函数是否在容量小于某个包裹时正确失败，求和是否溢出。很多二分错题的根源不在二分循环，而在谓词写错。二分只是加速寻找边界，不能修复错误的可行性判断。

\noindent\textbf{样例三：讲清单调栈。}

单调栈题的口述要说明“栈里保存的是还没有被解决的候选”。每日温度中，栈里是还没找到更高温的日期；柱状图中，栈里是还没确定右侧第一个更矮位置的柱子；下一个更大元素中，栈里是等待未来更大值的元素。新元素到来时，如果它能解决栈顶，就不断弹出并结算。

代码审查时，第一看栈里存值还是下标。只要需要距离、宽度或过期判断，就必须存下标。第二看相等元素如何处理。第三看是否需要哨兵。第四看弹出时结算对象是谁。柱状图最大矩形里，当前更矮柱不是答案高度，它只是右边界；被弹出的柱子才是本次结算的最矮高度。

\noindent\textbf{样例四：讲清图建模。}

图题口述第一句必须定义节点和边。岛屿数量的节点是陆地格子，边是四方向相邻；课程表的节点是课程，边是先修关系；公交路线题更适合把路线或“乘坐路线”作为 BFS 层次，而不宜把同一路线内所有站点两两连边；最小体力路径的边权是高度差，路径代价是最大边权，不是和。

代码审查时，第一看访问标记时机。BFS 通常入队时标记，避免重复入队。第二看状态维度。位置相同但钥匙集合不同、剩余障碍次数不同，都不能合并。第三看边权。等权用 BFS，非负权用 Dijkstra，零一权用双端队列。第四看邻居生成成本。单词接龙如果每次扫描所有单词找邻居，复杂度会很差，应考虑通配桶。

\noindent\textbf{样例五：讲清 DP。}

DP 口述要从暴力搜索树开始。先说每个搜索节点需要哪些信息，再说哪些节点会重复出现，最后给状态命名。编辑距离的状态是两个前缀长度；背包的状态是物品阶段和容量；股票的状态是天数、持有状态和交易次数；树形 DP 的状态是子树向父节点返回的信息。状态不是为了填表而填表，而是为了让未来决策只依赖少量参数。

代码审查时，第一看初始化是否对应空状态；第二看转移是否覆盖所有最后一步；第三看遍历顺序是否满足依赖；第四看空间压缩是否读旧值。0-1 背包倒序容量，完全背包正序容量，LCS 一维压缩保存左上角旧值，股票状态机保存昨天状态。只要空间压缩解释不清，就先退回完整二维或多状态表。

\noindent\textbf{源码阅读样例报告。}

\noindent\textbf{报告一：链表删除。}

问题：删除一个已知链表节点需要维护什么不变量。对象：业务结构体中嵌入一个链表节点。索引：链表按某种顺序组织对象。插入时要让前驱、后继和新节点三者互相指向；删除时要让前驱直接指向后继，并让后继回指前驱。若遍历中删除当前节点，必须先保存下一个节点，否则删除后当前节点的链接可能被重置。

刷题类比：LRU 缓存移动节点到头部，K 个一组翻转链表接回前后段。工程差异：内核对象可能同时挂在多条链上，删除某条链的节点不代表对象释放。并发情况下，还要考虑锁或 RCU。实验结论：链表操作的难点不是复杂度，而是节点身份、链接顺序和生命周期。

\noindent\textbf{报告二：红黑树插入。}

问题：动态有序索引如何插入新对象。对象：业务结构体中嵌入树节点。索引：按某个 key 维护二叉搜索树顺序。插入前，调用者沿树比较 key，找到空孩子位置；链接新节点后，库函数通过旋转和变色修复平衡。旋转不改变中序顺序，所以查找语义保持不变。

刷题类比：日程安排用有序表找前驱后继，数据流排名维护动态有序集合。工程差异：C++ \code{map} 隐藏节点分配和比较器，Linux rbtree 让调用者负责比较和对象嵌入。实验结论：有序树解决的是顺序查询，不是比哈希表更高级；只要题目需要前驱后继，哈希表就不够。

\noindent\textbf{报告三：RCU 删除。}

问题：读多写少结构中，删除对象为什么不能立刻释放。对象：读者可能正在访问的共享节点。索引：链表、树或哈希表让新读者找到对象。删除第一步是从索引摘除，让后续读者不可达；第二步等待旧读者离开读侧临界区；第三步释放内存。若摘除后马上释放，旧读者可能访问悬空指针。

刷题类比：Dijkstra 堆中的旧状态、滑动窗口堆中的过期元素。相同点是逻辑不可用和物理清理分离；不同点是 RCU 为并发安全服务，堆懒删除多为容器能力限制服务。实验结论：删除语义要拆开看，不能把“从答案集合移除”和“销毁对象”混成一步。

\subsection{单点深挖和高频设计}

\noindent\textbf{单点深挖：把模板拆成问题。}

这一节专门回应一个写作底线：单点问题要单点分析，不能把同一套话粘到每道题下面。真正的教材应该让读者看到每个问题自己的结构。接雨水不是“又一个双指针”，它的难点是何时可以结算某个位置的水；柱状图不是“又一个单调栈”，它的难点是被弹出的柱子在那一刻同时知道左右边界；LRU 不是“哈希加链表”，它的难点是节点身份、两个索引同步和对象生命周期。下面每个小节只抓一个具体难点，把暴力、优化、反例、实验和工程迁移连成一条线。

\noindent\textbf{LeetCode 42 接雨水：先问结算对象。}

接雨水最常见的错误讲法是直接背“左右指针，谁小移谁”。这句话不解释为什么能算水，也不解释什么时候不能算。先从暴力出发：对位置 i，水位等于左侧最高柱和右侧最高柱的较小值，再减去当前位置高度。暴力每个位置都向两边扫描，所以慢在反复寻找左右边界。前后缀数组的优化很直观：左边界数组保存每个位置左侧最高，右边界数组保存每个位置右侧最高，最后逐点结算。这个版本空间高一些，但最容易证明，因为每个位置结算时两个边界都已经明确。

双指针版本必须继承这个结算条件。设左侧已经看到的最高柱是 left max，右侧已经看到的最高柱是 right max。如果 left max 不大于 right max，那么左指针位置的右侧至少存在一个不低于 left max 的边界，所以它的水位上限已经由 left max 决定，此时可以结算左侧并右移。反过来，如果 right max 更小，就结算右侧。注意，结算的是较低已知边界那一侧的位置，不是简单因为某个指针当前高度小就移动。代码里可以用当前高度更新边界，再根据边界大小结算，也可以先比较再更新，但不变量必须保持一致。

这道题的实验要同时跑三版：暴力版、前后缀版、双指针版。输入选择单调递增、单调递减、平台高度、多个凹槽、长度不足三、两端高而中间低。打印每一轮 left、right、left max、right max 和本轮新增水量。若某个实现提前结算了右侧边界未知的位置，平台高度和多凹槽会立刻暴露问题。把这组实验做完，读者才能明白双指针不是口诀，而是把“每个位置需要左右最高”压缩成两个移动边界。

单调栈解法是另一种结算对象。栈中保存还没有找到右边界的柱子，且高度保持单调。当遇到更高的当前柱时，栈顶作为槽底被弹出，当前柱是右边界，弹出后的新栈顶是左边界。此时槽底水量才确定。这个版本特别适合解释“弹出时机”：不是当前柱接水，而是被弹出的低柱终于知道了左右边界。把双指针和单调栈放在一起比较，能训练一个重要能力：同一题可以有不同状态，但每个状态都必须回答“现在能结算谁”。

\noindent\textbf{LeetCode 76 最小覆盖子串：覆盖不是匹配总数。}

最小覆盖子串的暴力做法是枚举所有子串，再检查是否覆盖目标字符需求。这个暴力版清楚地暴露两个重复：窗口候选太多，检查需求时又反复统计。滑动窗口能成立，是因为右端扩展只会增加窗口里的字符，覆盖条件更容易满足；当窗口已经覆盖目标时，继续扩展只会更长，所以应当收缩左端尝试变短。这个推理依赖“覆盖”是一个可增量维护的条件。

很多错误实现把覆盖理解成匹配字符总数。目标串如果是 AABC，只看到一个 A 不能算完成 A 的需求；窗口里多出来的无关字符也不能增加完成度。更稳的状态是 need 保存目标每个字符所需次数，window 保存当前窗口次数，formed 表示有多少种字符已经达到需求。一个字符从少于需求变成恰好达到需求时，formed 加一；从恰好达到需求被左端移走变成不足时，formed 减一。formed 等于 need 中字符种类数，才表示窗口覆盖目标。这里的“种类达标”比“总匹配数”更不容易被重复字符误导。

证明时按固定右端来看。当右端扩展到覆盖目标，内层循环会把左端一直右移到再移就不覆盖为止。在这个过程中，每个合法窗口都被检查，最短的那个被记录。对所有右端重复这个过程，全局最短一定被看见。被移过的左端不会再产生更短答案，因为未来右端只会更大，窗口只会更长。这个排除逻辑必须讲清楚，否则读者只能记住一段滑动窗口代码。

实验输入要包含目标字符重复、目标字符缺失、答案在开头、答案在结尾、大小写混合、源串短于目标串。调试时打印 left、right、当前字符、formed、need size 和当前最优区间。若用总匹配数写法，目标含重复字符的输入会暴露错误；若左端收缩后忘记更新 formed，答案在开头或结尾的用例会出错。C++ 里字符集固定时用数组即可；若题目不保证字符集，再用哈希表。容器选择只是实现细节，核心仍是“覆盖状态”到底表示什么。

\noindent\textbf{LeetCode 862 最短子数组至少为 K：负数为什么击穿窗口。}

很多人看到“连续子数组”和“至少为 K”就想用滑动窗口。若数组全为正数，右端扩展会让窗口和变大，左端收缩会让窗口和变小，所以窗口条件具有单调性。但一旦允许负数，右端加入一个负数可能让和变小，左端移走一个负数可能让和变大，左端不再能安全排除。用一个小反例就能说明问题：窗口当前不满足 K，不代表继续扩展一定更接近；窗口当前满足 K，移走左端也不一定破坏或改善。普通窗口失去依据。

这题正确的切入点是前缀和。区间和等于 prefix[j] 减 prefix[i]，目标是找到 j 固定时最靠右、同时 prefix[i] 不大于 prefix[j] 减 K 的历史下标 i。为了让长度最短，历史下标越大越好；为了让未来更容易满足条件，前缀和越小越好。单调队列保存候选历史下标，队列里的前缀和递增。扫描当前前缀时，若当前前缀减队首前缀至少为 K，就可以更新答案，并弹出队首，因为对未来来说，更靠后的当前 j 已经让这个队首不可能给出更短长度。随后，若队尾前缀和大于等于当前前缀，则队尾被当前下标支配：当前下标更靠后，前缀和还不大，未来任何右端选择当前下标都不差。

这道题的重点是两个弹出条件含义完全不同。队首弹出是“已经形成一个答案，为了找更短长度继续丢旧左端”；队尾弹出是“旧候选被新候选支配，以后没必要保留”。如果只背单调队列，很容易把两个条件混成一个。实验要准备含负数、全正、答案不存在、答案为一、前缀和反复下降、K 很大这些输入。打印 prefix、队列下标、队列前缀值和答案。看到队尾被支配的过程，读者才会理解单调队列不是排序结构，而是保留一组仍有可能成为最优左端的历史状态。

这题也适合作为“模板失效”的案例。LeetCode 209 长度最小的子数组可以用正数窗口，LeetCode 560 和为 K 的子数组可以用前缀和加哈希计数，LeetCode 862 要用前缀和加单调队列。三个题都在说连续区间，但保存的历史状态不同：正数窗口保存当前和，计数题保存历史前缀频次，最短且含负数保存单调的历史前缀下标。把它们放在一起，读者才能学会根据条件换结构。

\noindent\textbf{LeetCode 84 柱状图：弹栈时到底算谁。}

柱状图最大矩形的暴力模型很清楚：把每根柱子当作矩形最矮高度，向左右扩展到第一个更矮柱子之间，面积就是高度乘宽度。暴力慢在每根柱子都重复找左右第一个更矮位置。单调栈的目的就是在一次扫描中确定这些边界。栈中保存高度递增的下标，表示这些柱子还没遇到右侧第一个更矮位置。

当当前柱子高度小于栈顶柱子时，栈顶柱子的右边界已经确定为当前下标。弹出栈顶后，新的栈顶就是它左侧第一个更矮柱子。于是被弹出的柱子作为最矮高度时，最大宽度就是右边界减左边界再减一。如果弹出后栈为空，说明左侧没有更矮柱，宽度从零延伸到当前下标之前。这里最容易错的是把当前柱子当成答案高度。当前柱子只是触发结算的右边界；真正被结算的是被弹出的柱子。

相等高度也要有一致策略。可以遇到小于栈顶才弹，也可以遇到小于等于时弹，只要宽度语义和重复高度处理一致。教材阶段建议先用末尾哨兵零，让所有柱子最终都被结算，避免循环结束后再写一段补丁逻辑。实验输入必须包括严格递增、严格递减、全相等、单柱、空数组和中间凹槽。严格递增会暴露末尾未结算，严格递减会暴露宽度计算，全相等会暴露相等高度策略。

这道题还连接 LeetCode 85 最大矩形。二维矩阵中每一行都可以看成柱状图底部，连续一的高度向下累积，然后对每一行运行柱状图算法。这里的迁移不是复制代码，而是迁移“某个高度作为最矮高度时找左右边界”的模型。读者若能从柱状图自然推到最大矩形，就说明已经掌握了单调栈的结算对象，而不是只记住一段弹栈代码。

\noindent\textbf{LeetCode 15 和 18：排序、去重和输出规模。}

三数之和最容易被讲成“排序加双指针”，但排序为什么有用，去重为什么必须分层，往往没有讲透。暴力三层循环覆盖所有下标三元组，慢在组合数量太多，也慢在重复值会生成大量相同答案。排序后的第一层收益是制造单调性：固定第一个数后，剩下两个数的和随左指针右移而增大，随右指针左移而减小。这样可以安全排除一批配对，而不是逐个枚举。

第二层收益是去重变得局部可见。固定值如果和前一个固定值相同，跳过当前固定值，因为它会生成同一组数值答案。找到一组三元组后，左指针和右指针也要跳过相同值，避免相同数值组合重复进入结果。注意，去重不是最后把答案塞进集合。集合过滤能让结果看起来正确，但它掩盖了算法结构，也增加了额外成本。高质量题解应该在生成阶段就避免重复答案。

四数之和是三数之和的压力测试。外层多固定一个数，目标和可能超出 int，答案数量也可能很大。排序后可以用最小可能和与最大可能和提前剪枝：若当前固定前缀加上后面最小几个数已经大于目标，后续更大固定值也不必继续；若加上最大几个数仍小于目标，则当前固定值太小，可以跳过。这个剪枝必须建立在排序上，也必须注意溢出。C++ 中目标和、中间和最好用更宽整数。

实验要包含全零、大量重复、无答案、正负混合、目标极大、答案很多。还要专门测输出规模：当答案本身很多时，复杂度不能只写计算部分，结果数组也占空间和时间。面试追问 k 数之和时，不要直接写递归模板。先说明递归每层固定一个数，最后降到双指针；再说明每层去重和剪枝如何保持不重复、不漏解。这样讲，三数、四数和 k 数才是一条连续的教材线。

\noindent\textbf{答案二分：真正要审查的是谓词。}

二分答案题表面上在写循环，实际难点在检查函数。以 LeetCode 410 分割数组的最大值为例，猜一个最大段和上限 limit，检查能否把数组按原顺序分成不超过 k 段，使每段和不超过 limit。检查函数的贪心是尽量把当前数字放进当前段，放不下才开新段。为什么这样正确？因为顺序固定，提前开新段只会让剩余段数更紧，不会减少所需段数。这个证明比外层二分更重要。

LeetCode 875 爱吃香蕉也是一样。速度 speed 越大，吃完所需时间越少，因此可行性从假到真。检查函数里每堆耗时要向上取整，不能用普通除法。边界也由谓词决定：速度下界是一，上界可以取最大堆；船运包裹的容量下界应是最大包裹，上界是总重量；分割数组的 limit 下界是最大元素，上界是总和。若下界设错，检查函数可能处理本不该出现的非法状态。

二分实验要把外层循环和检查函数分开测。先写一个线性枚举答案的慢版本，再写二分版本，两者对拍。对每个 mid 打印 left、right、谓词值和检查函数返回的段数或小时数。若二分结果错，先怀疑谓词，不要立刻改加一减一。很多错误都是“谓词并不单调”或者“检查函数没有表达题意”，外层循环再标准也无济于事。

这类题还要说明复杂度维度。它通常不是 \complexity{O(n\log n)}，而是 \complexity{O(n\log M)}，其中 M 是答案范围。背包、金额、容量、速度、最大段和这类数值进入复杂度时，要主动说明它和输入长度不是同一个维度。面试中能讲清这个点，会明显区别于只背模板的人。

\noindent\textbf{LeetCode 146 LRU：两个索引和一个节点身份。}

LRU 缓存的暴力实现可以用数组保存访问顺序。每次 get 或 put 时线性查找 key，再把元素移动到头部；容量满时删除尾部。这个版本正确但慢，瓶颈在定位 key 和移动节点。哈希表能常数定位 key，但不能表达最近使用顺序；双向链表能常数删除已知节点并移动到头部，但不能按 key 查找。组合结构的必要性就在这里：哈希表负责定位，链表负责顺序，二者共享同一个节点身份。

节点里必须保存 key 和 value。很多人疑惑哈希表已经有 key，链表节点为什么还要保存 key。原因在淘汰尾部时，链表先告诉我们最久未使用的节点，但要从哈希表中删除对应条目，必须知道它的 key。若节点只保存 value，就无法常数时间同步删除哈希索引。更新已有 key 时也不能新建重复节点，否则哈希表和链表会出现两个身份不一致的对象。正确做法是更新节点值，并把同一个节点移动到头部。

这题可以直接对照 Linux 侵入式链表。刷题版常用 \code{std::list} 让容器拥有节点，哈希表保存迭代器。内核常把 \code{list_head} 嵌进业务对象，链表只组织对象，不拥有对象。共同点是节点身份必须稳定，删除必须同步多个索引；区别是内核还要处理引用计数、锁、RCU 和分配上下文。把这种差异讲出来，LRU 就不再是六个字，而是一个小型索引系统。

实验要跑操作序列而不是只测单次 get。容量为一、重复 put、get 不存在、更新已有值、get 后再 put、连续淘汰都要覆盖。每步打印链表从头到尾的 key 顺序和哈希表大小，检查两个不变量：哈希表每个 key 都能在链表中找到唯一节点，链表尾部就是最久未使用节点。若某个操作后链表有重复 key 或哈希表指向失效迭代器，说明节点身份没有维护好。

\noindent\textbf{并查集：代表元不是路径。}

并查集适合回答“两个元素是否在同一集合”以及“把两个集合合并”。它不保存两点之间的具体路径，也不适合回答最短距离。LeetCode 684 冗余连接中，按顺序加入无向边；若一条边的两个端点已经同属一个集合，再加入它就会形成环。这道题不需要知道环上所有边，也不需要给出路径，所以并查集正好匹配。暴力每加一条边都 DFS 检查连通性，会重复扫描已有结构。

并查集的正确性来自等价关系。每个集合有代表元，find 返回代表；unite 只把两个代表连起来。路径压缩改变的是树形实现，不改变集合语义；按大小合并降低树高，也不改变集合语义。初学者容易把 parent 数组看成图的真实父边，这是错误的。压缩后 parent 会直接指向根，原图路径信息已经不存在。如果题目要输出路径，不能只靠普通并查集。

账户合并是另一个代表性例子。邮箱是连接账户的证据，姓名通常不能作为合并依据。可以给每个邮箱编号，把同一账户中的邮箱合并；最后按根收集邮箱并排序。这里并查集解决的是动态连通，哈希表解决的是字符串到编号的映射，排序解决的是输出格式。三个结构各司其职。若题目变成等式除法或带奇偶约束，并查集就要扩展权值或颜色；普通代表元不够。

实验要打印合并前后每个元素的根、集合大小和 parent 树高度。重复合并、已经连通的边、孤立点、字符串编号映射都要测。还可以故意去掉按大小合并或路径压缩，观察 find 路径变长。这个实验的目的不是追求近乎常数的理论证明，而是让读者看见代表元如何把“许多连接关系”压缩成一个可查询的集合状态。

\noindent\textbf{Dijkstra：旧堆状态不是错误。}

带正权图不能用普通 BFS，因为 BFS 按边数扩展，不按总代价扩展。LeetCode 743 网络延迟时间中，从起点到各点的最短时间取决于边权和。Dijkstra 的核心状态是当前已知最短距离，优先队列每次弹出距离最小的候选。若边权非负，弹出的未过期状态已经不可能被其他路径再改小，因此可以扩展它的出边。

C++ 的 \code{priority_queue} 没有 decrease-key。发现更短距离时，常见写法不是修改堆中旧元素，而是把新距离再压入堆。这样堆里会同时存在同一个节点的旧状态和新状态。弹出时比较堆顶距离和 distance 数组，若不一致，说明它已经过期，跳过即可。旧堆状态不是错误，只要使用前检查，它只是容器接口限制下的工程实现。

这点可以和 RCU 或懒删除做概念对照，但不能混为一谈。Dijkstra 的旧状态是算法实现中的过期候选，跳过即可；RCU 的旧对象是并发读者可能仍在访问的内存，必须等宽限期后释放。共同点是逻辑状态和物理存在分离；不同点是安全目标完全不同。教材要把相似性和差异都说清楚，避免把源码概念泛化成空话。

实验要构造重边、零权边、不可达节点、多条等长路径和明显会产生旧堆状态的图。打印每次弹出的节点、堆中距离、当前 distance、是否过期。再用普通 BFS 跑同一张带权图，观察结果差异。最后让读者回答：如果边权为负，为什么这个证明失效？这个问题能把 Dijkstra 的条件边界讲清楚。

\noindent\textbf{DP：从暴力搜索树命名状态。}

动态规划最怕直接给公式。以编辑距离为例，暴力思路是从两个字符串末尾开始，如果末尾字符相同，就处理剩余前缀；如果不同，就尝试删除、插入、替换。这个递归会反复处理同一对前缀长度。于是状态自然命名为 dp[i][j]，表示第一个字符串前 i 个字符变成第二个字符串前 j 个字符的最少操作。公式不是凭空出现，而是暴力选择的缓存。

背包也是同样逻辑。暴力枚举每个物品选或不选，未来只依赖“处理到第几个物品”和“剩余容量”。0-1 背包每个物品只能用一次，所以一维压缩时容量倒序，防止当前物品在同一轮被重复使用；完全背包允许重复使用当前物品，所以容量正序。若把循环方向背成模板，不构造反例，就很容易在分割等和子集或零钱兑换之间写错。

最长递增子序列展示了另一类优化。二次 DP 的状态是必须以 i 结尾的最长长度，枚举所有前驱。二分优化维护 tails，表示某个长度的递增子序列能拥有的最小结尾。tails 不是答案序列，而是未来扩展潜力表。这个区别必须说清楚，否则读者会误以为 tails 中的值就是实际选择路径。若要还原路径，还要记录前驱和每个长度对应的位置。

DP 实验不要一上来压缩空间。先打印完整表，让读者解释每个格子来自哪里。编辑距离表能看到插入、删除、替换；LCS 表能看到不连续子序列；背包表能看到每个物品阶段如何影响容量；股票状态机能看到持有、卖出、休息的约束。只有完整表能讲清，再做滚动数组。空间压缩是优化，不是理解的起点。

\noindent\textbf{Trie 和单词搜索：前缀剪枝不等于路径合法。}

单词搜索 II 的暴力做法是对每个单词在棋盘上单独 DFS。慢点在于不同单词共享大量前缀，却被重复搜索。Trie 把所有单词的前缀合并，棋盘 DFS 时沿 Trie 节点前进；如果当前路径在 Trie 中没有对应前缀，就立即剪枝。这个优化节省的是“重复检查前缀”的工作，而不是改变网格路径规则。

路径合法性仍然由访问标记控制。同一条路径中，一个格子不能重复使用；进入格子要标记，递归返回要恢复。Trie 节点上存在单词结束标记，只表示当前路径构成一个词；它不允许你忽略网格访问约束。找到一个词后，可以清空该节点的单词标记避免重复加入答案，也可以在工程上做叶子删除减少搜索，但这些都不能破坏其他单词共享的前缀。

实验要准备重复字母棋盘、单字符单词、一个单词多条路径、多个单词共享前缀、路径必须回退的用例。打印当前坐标、路径字符串、Trie 深度和访问矩阵。若忘记恢复访问标记，其他分支会漏答案；若没有访问标记，同一格子会被重复使用；若找到词后错误删除共享节点，另一个共享前缀的词会丢失。这个单点能把回溯、Trie 和状态恢复三者联系起来。

C++ 实现时要考虑节点所有权。教学代码可以用数组下标池或智能指针；若用裸指针，必须说明释放责任。字符集固定为小写字母时，子节点数组更快；字符集稀疏时，哈希表更省空间。不要在 DFS 热路径里反复构造大字符串，可以用路径数组或直接在 Trie 节点保存完整单词。容器选择仍然服务于访问模式。

\noindent\textbf{动态区间和 Maple Tree 思维。}

离线区间题和动态区间题的难点不同。LeetCode 56 合并区间只需排序后线性扫描；所有区间一次性给出，排序把可能相交的区间放到相邻位置。日程安排、内存区间表或虚拟地址区域管理则是动态问题：区间会插入、删除、查询，不能每次都重新排序全部数据。此时需要有序结构维护相邻关系。

用 \code{std::map} 写简化动态区间表，key 是区间起点，value 是终点和属性。点查询时，找起点不大于目标地址的最后一个区间，再判断终点是否覆盖。插入新区间时，只需检查前驱和后继是否重叠；如果它们都不冲突，更远区间也不可能冲突。删除一段时更复杂：可能删除整个区间，可能截掉左侧或右侧，可能从中间挖掉一段并分裂成两个区间。

Maple Tree 的源码背景比这个实验复杂得多，但问题动机相通：大地址空间、稀疏区间、频繁查询、动态更新、读路径性能和并发安全。刷题阶段不需要复刻 Maple Tree，但可以通过 map 区间表理解范围映射的基本不变量。读源码时要问节点保存了哪些范围信息，查询如何找到覆盖区间，更新如何保持相邻关系，旧节点何时释放。

实验用例要包括插入到最前、插入到最后、端点相接、完全覆盖、跨多个区间、删除中间子段、删除不存在区间。报告要写清半开区间和闭区间的差别。半开区间下，前一个终点等于后一个起点不重叠；闭区间下，端点相等可能冲突。很多区间题的错误都来自端点语义没有先定下来。

\noindent\textbf{伙伴系统：总量够不代表连续块够。}

算法题常把内存分配当成免费，但系统里分配器本身就是算法。伙伴系统把连续页按 2 的幂次阶管理。分配某阶块时，若该阶没有空闲块，就向更高阶找，找到后逐层拆分；释放时，如果伙伴块同阶且空闲，就合并成高一阶块。这个机制强调连续性和对齐，而不是只看空闲页总量。

一个系统可能总空闲页很多，却没有足够大的连续块。比如很多小块分散在不同位置，总和超过请求大小，但无法合并成所需阶数。这就是碎片问题。刷题里可以把它类比为区间合并和二进制对齐，但不能简化成普通计数。每个空闲块必须只出现在它当前阶的链表里，拆分和合并都要维护这个不变量。

实验可以设总页数为 1024，维护从零阶到十阶的空闲链表。分配时打印从哪个高阶拆下，释放时打印伙伴编号和是否合并。构造一个碎片案例：先分配多个小块，再释放其中一部分，让总空闲量看似足够，但高阶请求失败。这个实验能把位运算、链表、区间、贪心分配和工程性能放到一起理解。

把这个视角带回 C++ 容器，读者会更容易理解为什么大量链表节点、哈希节点和树节点的分配成本不能忽略。复杂度都是线性或对数，连续数组可能因为缓存局部性和少分配而快很多；节点容器可能因为稳定身份和常数删除而必要。算法复杂度给出数量级，分配器和内存局部性决定真实表现。

\noindent\textbf{题解写作：每道题至少回答一个自己的问题。}

写题解时，不要让所有题都长成同一张表。可以有统一的复盘框架，但正文必须回答本题自己的问题。两数之和要回答为什么先查再插入；无重复子串要回答左端为什么不能回退；接雨水要回答何时能结算某个位置；柱状图要回答弹栈时算谁；LRU 要回答为什么节点里要保存 key；Dijkstra 要回答堆中旧状态为什么可以跳过；背包要回答容量循环方向从哪里来。

一个合格段落应当增加新信息。它可以解释一个错误实现为什么错，可以给一个最小反例，可以把暴力中的重复工作定位出来，可以说明一个容器的副作用，可以把源码结构和刷题模型做有限对照。如果一个段落只是在说“注意边界”“使用哈希优化”“复杂度降低”，但没有说明具体边界、具体 key、具体降低了哪一步，就应当重写。

题解还要有实验意识。每个高频题至少保留一个暴力对拍版本或小规模枚举器。前缀和题用随机数组对拍，二分题用线性枚举答案对拍，单调栈题打印入栈弹栈，DP 题打印小表，图题打印完整状态，链表题画节点连接。实验不是附属内容，它是把“我觉得对”变成“我知道为什么错不了”的过程。

最后是语言衔接。章节不是题号仓库，应该从一个问题过渡到下一个问题：连续内存引出双指针和窗口，窗口失效引出前缀和与单调队列，栈的最近未完成结构引出表达式和柱状图，动态候选集引出堆和最短路，子问题重复引出 DP，动态有序关系引出区间树和源码结构。这样的书读起来才像教材，而不是把许多答案贴在一起。

\noindent\textbf{源码与系统结构精读。}

刷题教材加入 Linux 源码，不是为了让读者背内核字段名，而是为了让读者看到同一类数据结构在真实系统中为何会变复杂。力扣题通常默认单线程、对象生命周期简单、内存分配随时可用；内核代码必须同时考虑并发读写、引用计数、缓存局部性、分配上下文和对象同时进入多个索引。下面这些小节仍然坚持单点分析，每节只追一个问题。

\noindent\textbf{list head：为什么节点嵌进对象。}

普通 C++ 初学者会把链表理解为“容器里存元素”。Linux 的侵入式链表不是这样。业务对象中嵌入链表节点，链表节点只负责前后指针，不拥有对象，也不知道对象类型。遍历时通过成员偏移从链表节点回到外层对象。这个设计看起来底层，但它直接解决三个工程问题：减少额外分配，让对象可以同时进入多条链表，保持对象地址稳定。

把它和 LRU 放在一起看会很清楚。刷题版 LRU 中，缓存条目既需要按 key 被哈希表找到，又需要按新旧顺序进入链表。如果链表节点和业务对象分离，删除和移动时要处理额外间接层；如果对象自己带链表节点，哈希表可以直接指向对象，链表只改变对象中的链接字段。内核里一个页面对象可能同时出现在 LRU 链、回写链或其他状态链上，所以对象中可以有多个链表节点字段，每个字段代表一条不同关系。

单点实验是写一个 Item，里面有 key、value、lru node、dirty node。先把三个 Item 挂入 LRU 链，再把其中两个挂入 dirty 链。删除一个对象的 dirty node，不应影响它在 LRU 链中的位置；释放对象前，必须先确认它已经从所有链和索引中摘除。这个实验比单纯反转链表更接近源码思维，因为它把“节点身份”和“对象生命周期”放在同一处观察。

\noindent\textbf{hlist：哈希桶为什么不一定用完整双链。}

哈希表有大量桶，很多桶可能为空。如果每个桶头都使用完整双向链表节点，桶数组本身就会消耗更多空间。hlist 的设计可以理解为桶头较轻，节点保存足够信息以支持已知节点删除。桶头只需要知道第一个节点，节点除了 next，还保存指向自己的前驱链接位置。这样删除已知节点时，可以直接修改前驱的 next 或桶头的 first，不必从桶头扫描找前驱。

刷题里哈希冲突通常被标准库隐藏，但很多题都在使用同一思想：key 先定位桶，桶内再比较实际 key。两数之和的 key 是历史值，字母异位词的 key 是规范签名，前缀和计数的 key 是前缀值，账户合并的 key 是邮箱。高质量题解要说清 key 和 value，不应只说“用 unordered map”。如果 key 设计错了，哈希表再快也会得到错误答案。

实验可以写两个固定桶数的小哈希表。第一版桶内普通单链表，删除已知节点必须找前驱；第二版节点保存前驱链接位置，删除时不用扫描。再设计一个遍历中删除的场景：如果读者正在走桶链，写者删除节点后什么时候能释放对象？单线程实验只需保存 next，内核并发场景则需要锁或 RCU。这个单点把哈希冲突、删除成本和并发生命周期自然连起来。

\noindent\textbf{rbtree：比较逻辑为什么交给调用者。}

C++ 的 map 把比较器和节点管理封装在容器里。Linux rbtree 更像一个底层平衡树工具：树节点嵌在业务对象里，插入和查找时，调用者自己沿树比较 key，找到位置后再调用库函数链接和修复平衡。这样设计的好处是业务对象可以按不同字段进入不同树，内存分配和对象生命周期也不被通用容器隐藏。

红黑树单点要抓住“旋转不改变中序顺序”。插入新节点后，树可能违反颜色和平衡约束，修复函数通过变色和旋转调整局部结构。旋转前后，中序遍历得到的 key 顺序相同，所以查找语义不变。很多教材一上来讲红黑树性质，读者很快迷路；源码阅读可以先问一个更小的问题：调用者如何比较 key，库函数如何在不破坏有序性的前提下修复平衡。

刷题对应的是动态有序集合。日程安排要找前驱后继，滑动窗口有序统计要插入删除并取极值，时间键值存储在每个 key 内部按时间二分，动态区间表需要按起点维护顺序。哈希表适合等值查找，不能回答前驱后继。实验用 map 实现日程安排，再手写一个简化 BST 插入，不要求实现红黑修复，只要求打印中序序列。等读者理解有序语义稳定后，再看平衡修复才不会把重点放错。

\noindent\textbf{XArray：稀疏整数索引不是普通哈希。}

XArray 面向整数索引到对象指针的映射，适合大而稀疏的索引空间。巨大数组会为空洞浪费空间，普通哈希表能做等值查找，但对按序推进、范围扫描、标记和并发读写支持不一定合适。XArray 这类结构把整数索引按位分层，只为实际存在的路径分配节点。它和 Trie 的直觉相似：字符串 Trie 每层消耗一个字符，整数分层索引每层消耗几位。

刷题里可以用三个模型建立直觉。位 Trie 用二进制位分层，常见于最大异或；坐标压缩把稀疏大值域映射到紧凑数组下标；有序 map 支持稀疏 key 的顺序遍历。XArray 和它们都不完全相同，但共享一个问题：key 空间很大，实际对象稀疏，需要在空间、查找、遍历和更新之间折中。

实验可以写一个简化二进制 Trie 保存整数集合，支持插入、查找、寻找不小于某个值的下一个元素。再和数组、哈希表、map 比较：数组适合小连续值域，哈希适合稀疏等值查找，map 适合动态有序，Trie 适合可按位分解的 key。读 XArray 时不要先背节点字段，先回答索引如何分层、空洞如何表示、读者什么时候能看到新对象。

\noindent\textbf{Maple Tree：范围映射的三个操作。}

范围映射比单点映射更容易出错。给定地址要找到覆盖它的区间，插入新区间要检查重叠，删除区间可能把原区间切开。Maple Tree 服务的正是这种范围索引场景。它比普通 map 复杂，是因为真实内核还要追求读路径性能、处理并发、减少分配和维护更复杂的节点布局。但教材阶段先抓三个操作：点查询、插入、删除。

用简化 map 区间表可以训练这些操作。点查询先找起点不大于地址的最后一个区间，再看终点是否覆盖；插入只需检查前驱和后继，因为已有区间不重叠且按起点有序；删除有四种情况：删除整段，截掉左侧，截掉右侧，从中间挖掉一段并分裂。每种情况都要画图，否则代码很容易在端点处错。

这和力扣 56、57、729、731 是一条线。合并区间是离线排序扫描，插入区间是一次性合并连续重叠段，我的日程安排是动态插入并检查冲突。Maple Tree 是系统级动态范围集合。读者不需要一开始理解全部源码，但要知道为什么一个普通哈希表无法回答“哪个区间覆盖这个地址”，也无法按地址顺序找下一个区域。

\noindent\textbf{RCU：删除要拆成摘除和释放。}

RCU 最适合从链表删除理解。单线程删除一个节点，可以改前后指针后立刻释放节点。读多写少并发场景不行：读者可能已经拿到旧指针，写者把节点从新索引中摘除后，旧读者仍可能继续访问它。于是删除必须拆成两个阶段：先从结构中摘除，让后来的读者找不到；再等待所有可能看到旧节点的读者离开；最后释放内存。

这个拆分和刷题里的懒删除有相似直觉，但目的不同。滑动窗口中位数的堆懒删除是因为 priority queue 不能删除任意元素；Dijkstra 的旧堆状态是因为没有 decrease-key；RCU 延迟释放是为了内存安全。相似点只是“逻辑不可用”和“物理存在”分离，不能把它们说成同一种机制。教材要防止概念类比过头。

实验可以模拟一个全局配置表。读者进入读侧临界区后读取当前指针；写者复制旧表，修改副本，发布新指针，把旧表放入延迟释放列表。只有当活跃读者计数为零时，才释放旧表。这个实验能让读者看到，读路径轻量的代价是写路径更复杂。回到算法题，读者也会更谨慎地使用“删除”这个词：从集合不可达、从答案中移除、从内存释放，三者不是同一件事。

\noindent\textbf{调度队列：为什么动态最小值不总是堆。}

调度器需要在可运行任务中选择下一个运行者。一个简化模型是每个任务有虚拟运行时间，每次选择虚拟运行时间最小的任务运行，运行后它的虚拟运行时间增加，再放回候选集合。直觉上堆可以取最小值，但真实调度还需要删除任意任务、处理睡眠唤醒、按顺序遍历、维护权重和调度类。堆只解决“取最小”，不解决全部操作。

刷题中类似问题很多。数据流中位数用双堆，因为只需要两半极值；日程安排用有序表，因为需要前驱后继；任务调度器可以用堆模拟，也可以用最高频公式；滑动窗口最大值用单调队列，因为候选只按窗口顺序过期。结构选择必须列出操作集合：插入、删除任意元素、取最值、找前驱后继、范围遍历、按时间过期。只看一个操作很容易选错结构。

实验写一个简化调度模拟器。第一版用 priority queue，每次任务运行后压入新状态，睡眠任务只能懒删除；第二版用 set，可以按迭代器删除任意任务，再重新插入更新后的任务。比较两版代码，读者会看到接口能力的差异。复杂度同为对数，维护语义却不同。这个单点能把堆、有序树和系统队列联系起来。

\noindent\textbf{页缓存和 LRU：真实淘汰不是一道缓存题。}

力扣 LRU 缓存的规则很干净：访问就移动到头部，容量满就淘汰尾部。页面回收里的 LRU 思想复杂得多。页面可能是文件页或匿名页，可能干净或脏，可能正在写回，可能被引用，可能属于不同内存区域。系统不能简单删除“最旧页面”，还要考虑 I/O 成本、回收收益和并发状态。

教材中引入这个差异，是为了防止读者把刷题结构绝对化。哈希表加双向链表是 LRU 的核心索引组合，但真实系统可能使用多条链表达状态分层，例如活跃和非活跃，文件和匿名。页面从一条链移动到另一条链，不仅是顺序变化，也是状态变化。删除页面也不是释放一个普通对象，还要处理引用、脏页写回和映射关系。

实验可以从简化两级 LRU 开始。把页面分成 active 和 inactive 两条链，访问 inactive 页面时提升到 active，周期性把 active 尾部降级，回收从 inactive 尾部开始。再加入 dirty 标记：脏页不能立即释放，需要先模拟写回。这个实验不追求复刻内核，只让读者看到“链表表达状态，哈希表达定位，标记表达附加约束”。再回到 LRU 题，读者会更理解为什么节点要稳定、为什么多个索引要同步。

\noindent\textbf{高频设计题精读。}

设计题比普通单题更容易露出模板化问题，因为它们通常不是一个数据结构就能完成，而是多个索引协同维护。每次读设计题，都要写出操作表：每个操作需要什么查询、什么更新、什么删除、什么顺序。下面几个题都按这个方式分析。

\noindent\textbf{LeetCode 460 LFU 缓存：频次和时间两个维度。}

LFU 缓存要求淘汰使用频次最低的项；如果频次相同，淘汰最久未使用的项。它比 LRU 多一个频次维度。暴力做法可以每次淘汰时扫描所有 key，找最低频次和最旧时间，复杂度太高。要做到常数时间，需要两个索引：key 到节点信息，frequency 到同频节点链表。还要维护当前最小频次 min freq。

节点信息至少包含 key、value、freq，以及它在对应频次链表中的位置。get 命中后，节点频次加一：先从旧频次链表删除，再加入新频次链表头部；若旧频次链表空且旧频次等于 min freq，min freq 加一。put 新 key 时，若容量满，从 min freq 对应链表尾部淘汰最旧节点。插入新节点的频次为一，因此 min freq 重置为一。

这题最容易错在只更新频次，不更新同频时间顺序；或者删除节点后忘记维护 min freq。实验序列要包含容量为零、容量为一、同频淘汰、get 后频次变化、put 已存在 key、低频链表清空。每步打印 min freq、每个频次链表顺序和 key 表大小。LFU 的核心不是“哈希加链表”四个字，而是频次索引和时间索引同时一致。

\noindent\textbf{LeetCode 895 最大频率栈：为什么 frequency 到 stack 足够。}

最大频率栈要求 pop 时返回出现频率最高的元素；若频率相同，返回最近压入的元素。暴力可以每次 pop 扫描所有元素，找最高频和最近时间。优化时观察两个维度：value 到当前频率，frequency 到达到该频率的元素栈。push 一个值时，它的频率从 f 变成 f+1，把这个值压入 group[f+1]，同时更新 max freq。

pop 时直接从 group[max freq] 弹出最近达到最高频的值，然后把它的当前频率减一。若 group[max freq] 变空，max freq 减一。为什么不需要从旧频率栈中删除这个值？因为一个值每次频率增加都会在新频率栈留下记录，这些记录对应历史时刻；pop 只会从当前最高频栈弹出合法的最近记录。频率减一后，低频栈中旧记录仍然代表它之前达到低频时的顺序，后续在那个频率层可能被使用。

实验序列要让同一个值多次 push，并和另一个值形成同频竞争。打印 value frequency、max freq、每个 group 栈。若实现使用单个堆，也能做，但要处理过期记录；frequency 到 stack 的结构更贴合题意，因为“同频最近”天然是栈。这个题训练的是把一个复杂排序规则拆成两级索引，而不是盲目使用优先队列。

\noindent\textbf{LeetCode 981 时间键值存储：按 key 分组再二分。}

时间键值存储支持 set(key,value,timestamp) 和 get(key,timestamp)，返回该 key 在不超过 timestamp 的最新值。暴力做法把所有记录放在一个列表里，每次 get 扫描。优化的第一步是按 key 分组，因为不同 key 的记录互不影响。对每个 key，保存按时间递增的 pair 列表。get 时在该 key 的列表中二分，找最后一个时间不超过目标的记录。

这题的关键条件是 set 的 timestamp 通常按递增顺序到来。如果不递增，就需要插入时保持有序，或者最后排序。二分查找的是 upper bound：第一个大于目标时间的位置，答案是它前一个。若位置为开头，说明没有可用记录。这里不能用哈希表直接按 timestamp 查，因为查询是“不超过”的前驱问题，不是等值问题。

实验要覆盖 key 不存在、时间早于第一条记录、等于某条记录、落在两条记录之间、晚于最后记录、多个 key 交错写入。报告里要写清为什么结构是 hash map 到 vector，而不是一个全局 map。这个题非常适合训练“先按独立维度分组，再在组内使用有序结构”的思想。

\noindent\textbf{LeetCode 632 最小区间：多路归并里的覆盖条件。}

最小区间覆盖 K 个列表，要求区间内至少包含每个列表一个数。暴力可以从所有元素中枚举左右端，再检查是否覆盖每个列表。优化从多路归并视角看：每个列表当前选一个元素，则当前 K 个元素天然覆盖所有列表，区间由其中最小值和最大值决定。为了缩小区间，应当推进当前最小值所在的列表，因为当前最大值不变时，只有提高最小值才可能缩短区间。

小顶堆保存每个列表当前元素，同时维护当前最大值。每次弹出最小元素，更新答案区间，然后把该元素所在列表的下一个元素压入堆，并更新最大值。如果某个列表耗尽，就无法再保持覆盖所有列表，算法结束。这个过程和合并 K 个有序链表相似，但多了当前最大值和覆盖区间的目标。

证明时看当前状态覆盖了所有列表。若不推进最小元素，而推进其他列表，当前最小值仍在，区间左端不变，最大值还可能变大，不可能得到更优区间。因此推进最小所在列表是安全的。实验覆盖某个列表只有一个元素、多个相同值、答案在开头、答案在中间、列表长度差异很大。C++ 中堆元素保存值、列表编号和下标，避免复制整个列表。

\noindent\textbf{LeetCode 815 公交路线：节点选错会爆炸。}

公交路线题要求从起点站到终点站最少乘坐几辆公交。若把站点作为节点，边表示两个站在同一路线上相邻，建图会非常大；若把同一路线内所有站两两连边，更会爆炸。更稳的建模是站点到路线列表的哈希表，BFS 的层次表示乘坐公交次数。起点站能乘坐的所有路线作为第一层路线，访问一条路线时，可以到达这条路线上的所有站，再从这些站扩展到其他路线。

这里的访问标记要分路线和站点两个层次。路线一旦乘坐过，就不需要再次乘坐；站点如果已经用于扩展过，也可以避免重复扫描它对应的路线列表。目标是到达终点站，不一定要构造完整的路线图。若起点等于终点，答案为零；若终点站不在任何可达路线中，返回失败。

实验要比较两种建模的扩展数量。构造多条长路线共享少量换乘站的输入，站点图会产生大量邻接关系，而路线 BFS 只在需要时通过站点找到可换乘路线。这个题说明图题第一步不是写 BFS，而是选择节点含义。节点选错，后面所有优化都只是在错误模型上补救。

\noindent\textbf{LeetCode 1235 规划兼职工作：排序、二分和 DP 合体。}

规划兼职工作给出若干工作，每个工作有开始时间、结束时间和收益，要求选择不重叠工作最大化收益。按收益贪心会失败，因为高收益工作可能阻塞多个组合后更高的工作。暴力枚举子集会指数爆炸。正确切入点是按结束时间排序，定义 dp[i] 表示考虑前 i 个按结束时间排序的工作时的最大收益。

对于第 i 个工作，有两种选择：不选它，收益是 dp[i-1]；选它，则要找到最后一个结束时间不超过当前开始时间的工作 j，收益是 dp[j] 加当前收益。因为工作已按结束时间排序，j 可以用二分找到。这个二分不是在答案空间上，而是在已经排序的结束时间数组上找前驱边界。

实验要覆盖完全不重叠、全部重叠、端点相接是否允许、同结束时间、同开始时间、收益很大的单个工作和多个小工作组合。打印排序后的工作、每个 i 找到的 j、选择和不选择的收益。这个题很好地展示了多技术衔接：排序制造可二分的结束时间，二分找到兼容前驱，DP 比较选或不选。任何一个环节单独背模板都不够。

\noindent\textbf{LeetCode 1675 最小偏移量：先规范化状态空间。}

数组最小偏移量的操作规则有点绕：偶数可以除以二，奇数可以乘以二。若直接搜索所有可能状态，分支很多。关键观察是每个奇数最多先乘二一次，变成偶数后就只能不断除以二；每个偶数也只能不断除以二。于是可以先把所有数规范化到“可往下减少的最大形态”：奇数乘二，偶数保持。之后每次只尝试降低当前最大偶数，同时维护当前最小值。

用大顶堆保存当前所有数中的最大值，另有变量保存当前最小值。每轮用最大值减最小值更新答案；如果最大值是偶数，就把它除以二放回堆，并更新最小值；如果最大值是奇数，就无法继续降低它，算法结束。为什么每次只动最大值？偏移量由最大和最小决定，降低非最大值不会降低当前最大值，反而可能降低最小值让偏移更大。

实验覆盖全奇数、全偶数、混合、已经最优、最大值很大、多个相同值。打印堆顶、当前最小值和答案。这个题的单点是“规范化状态空间”：先把所有数推到统一方向可操作的边界，再用堆做贪心推进。它和很多搜索题相通，先把可逆或双向操作改成单向可控过程，复杂度才降下来。

\noindent\textbf{实验报告样例库。}

下面给几份更具体的实验报告样例。它们不是为了增加形式，而是给读者一个可执行标准：每个实验都能复现一个算法为什么对，或者一个错误实现为什么错。

\noindent\textbf{报告样例：前缀和对拍。}

问题：验证和为 K 的子数组哈希版是否漏掉边界。暴力基线：两层枚举左右端点，内部用累加变量得到区间和。优化版本：扫描前缀和，哈希表保存历史前缀频次。输入生成：数组长度从零到二十，元素范围从负三到三，K 从负五到五随机选取。每组数据同时跑暴力和优化，结果不同就打印数组、K、当前前缀、哈希表内容和下标。

预期发现：若漏掉初始前缀零，所有从下标零开始的答案会少计；若先更新再查询，K 为零时可能多算空区间；若用滑动窗口替代哈希，含负数输入会失败。结论：本题依赖代数等式，不依赖单调性；哈希表 value 必须是频次，因为题目问数量；查询更新顺序是算法语义的一部分。

\noindent\textbf{报告样例：单调栈状态打印。}

问题：验证柱状图最大矩形宽度计算。暴力基线：每根柱子向左右找第一个更矮位置。优化版本：单调递增栈加末尾哨兵。输入集合：空数组、单柱、严格递增、严格递减、全相等、中间低谷。每次弹栈时打印当前下标、被弹出下标、高度、左边界、右边界、宽度和面积。

预期发现：若没有末尾哨兵，严格递增数组会漏算；若宽度写成 right-left，递减数组会面积过大；若把当前柱高度当作结算高度，中间低谷会错。结论：弹出时被弹出的柱子才是矩形高度，当前柱只提供右边界，弹出后的栈顶提供左边界。

\noindent\textbf{报告样例：二分谓词审查。}

问题：验证船运包裹容量二分。暴力基线：从最大包裹到总重量线性枚举容量，找到第一个可行容量。优化版本：左闭右开二分第一个可行容量。检查函数：按顺序装包，当前天放不下就开新天，统计天数是否不超过 D。输入集合：一天运完、每天一个包裹、最大包裹等于答案、多个相同重量、容量小于最大包裹的非法 mid。

预期发现：若下界不是最大包裹，检查函数必须额外处理单包超过容量；若天数统计初始值错，边界会偏一；若二分返回未验证的 left，也可能隐藏谓词错误。结论：答案二分的核心是可行性谓词单调，外层循环只是利用这个谓词找边界。

\noindent\textbf{报告样例：LRU 索引一致性。}

问题：验证 LRU 的哈希表和链表是否一致。暴力基线：用 vector 保存从新到旧的 key，每次操作线性移动。优化版本：哈希表加双向链表。操作序列：容量为二，依次 put 一、put 二、get 一、put 三、get 二、put 一更新、put 四。每步打印链表顺序、哈希表 key 集合、每个 key 指向的节点值。

预期发现：若 get 后不移动节点，淘汰顺序会错；若 put 已存在 key 时新建节点，链表会出现重复 key；若淘汰尾部时忘记删哈希表，后续 get 会访问失效节点。结论：LRU 的不变量有两个，哈希表到链表节点一一对应，链表顺序表达最近使用到最久未使用。

\noindent\textbf{报告样例：图状态维度。}

问题：验证障碍消除最短路是否错误合并状态。暴力基线：小网格上用 BFS 保存行、列、剩余消除次数。错误版本：visited 只保存行和列。输入集合：必须绕路保留消除次数的网格、起点终点相邻、K 为零、K 大于障碍数。每次入队时打印位置、剩余次数和距离。

预期发现：二维 visited 会把“到达同一格但剩余资源更多”的状态丢掉，导致错误不可达或路径过长。结论：BFS 第一次到达最短，只对完整状态成立。状态维度少了，访问标记就会错误剪枝。

\noindent\textbf{报告样例：背包循环方向。}

问题：验证 0-1 背包一维压缩方向。暴力基线：二维 dp，阶段是物品下标，容量是当前目标。错误版本：一维数组容量正序。输入集合：一个物品重量一、容量二；两个物品重量一和二、容量三；分割等和子集中的小数组。打印每轮物品处理后的 dp 数组。

预期发现：正序容量会让同一个物品在同一轮被重复使用。倒序容量读取的是上一轮阶段状态，符合每个物品只能用一次。结论：循环方向不是模板差异，而是题目对物品使用次数的约束。

\noindent\textbf{章节衔接重写原则。}

本书后续扩展时，章节应该按能力递进，而不是按题号堆叠。线性结构部分先讲连续内存、下标和缓存局部性，再讲双指针和窗口；窗口失效后，自然引出前缀和、哈希和单调队列。栈和队列部分先讲最近未完成结构，再讲单调栈、表达式解析和柱状图。堆部分先讲动态极值，再讲 Dijkstra、数据流中位数和多路归并。

树和图部分也要有顺序。树先讲遍历方向：自顶向下传约束、后序向上汇总、层序传播；再讲 BST 的有序性和 Trie 的前缀共享。图先讲节点和边的建模，再讲 BFS 的层次、DFS 的连通块、拓扑的依赖、并查集的集合代表，最后讲带权最短路。这样读者能看见一个概念如何扩展，而不是在相邻页里突然切换题型。

动态规划部分必须从暴力搜索树开始。线性 DP 解释最后一步，二维 DP 解释两个前缀，背包解释阶段和容量，区间 DP 解释最后合并点，状态机解释约束转移，状态压缩解释集合。每个小节都要有一个完整表格实验。空间压缩只能在完整状态讲清后出现，不能一开始就给滚动数组代码。

Linux 源码专题放在基础算法之后更合适。读者先掌握链表、哈希、有序树、区间、队列、堆、分配器这些算法结构，再去看内核里的侵入式节点、hlist、rbtree、XArray、Maple Tree、RCU、伙伴系统和 SLUB，才不会被字段名淹没。源码章节不是另一本操作系统书，而是把算法结构放进真实工程约束里重新审视。

目录也要克制。章标题保留大的学习路线，小节标题只在正文中服务阅读，不把每个微型点都塞进目录。一本教材的视觉舒适来自层级稳定、段落长度适中、代码块清晰、图示服务推理。章节过碎会让读者感觉像在翻题库；章节过长没有锚点又会迷路。后续排版应在这两者之间取平衡。

\noindent\textbf{新增专题候选清单。}

如果继续扩写，优先补这些不是模板、能提升质量的专题。第一，字符串算法专题：KMP 的失败函数为什么能跳转，Z 函数如何维护右边界，Manacher 如何把回文半径线性化，后缀数组和后缀自动机只作为进阶引导。每个算法都必须从暴力匹配出发，而不是直接给公式。

第二，数学和位运算专题：快速幂、欧几里得、质数筛、组合数、模逆元、同余、异或线性基。面试刷题不一定高频，但它们能训练代数不变量。写法上要强调适用条件，例如模逆元要求互质或模数为质数，位移要注意有符号整数。

第三，图高级专题：强连通分量、割点桥、最小生成树、二分图匹配、网络流只做求职向概览。重点不是铺算法名，而是说明它们分别回答什么问题：强连通回答有向互达，桥回答删边连通性，最小生成树回答连接全部点的最低总成本，匹配回答两侧资源配对。

第四，工程性能专题：缓存局部性、分配器、哈希退化、字符串拷贝、迭代器失效、递归栈、输入输出性能。每个点都要配小实验。算法书如果不讲这些，读者容易以为复杂度相同的代码真实表现也相同。

第五，面试表达专题：如何在五分钟内讲清一题，如何在追问中判断条件变化，如何承认不确定并回到暴力，如何用反例纠正错误贪心。这个专题能把刷题能力转成沟通能力，但必须用具体题演示，不能写成空泛建议。

\subsection{质量审查和错题复盘}

\noindent\textbf{反模板质量审查。}

这一节给后续写作和修订使用。算法教材最容易出现的低质量问题，不是少写几个题，而是把每道题都写成相同段落：先说暴力，再说优化，再说边界，再说 C++。这种结构本身可以作为复盘框架，但不能成为正文内容。正文必须回答本题自己的问题，否则读者读十页之后只会记住一套空话。

\noindent\textbf{审查一：每题是否有自己的核心疑问。}

每道题提交前，先写一句“本题真正要解释的问题是什么”。两数之和的问题是为什么历史表要先查再插入；无重复子串的问题是为什么左边界不能回退；接雨水的问题是何时能结算某一侧；柱状图的问题是弹栈时谁的左右边界确定；课程表的问题是边方向如何影响输出顺序；零钱兑换的问题是最少数量和组合数量的状态含义不同。若这句话写不出来，说明题解还停留在题型标签上。

核心疑问必须能被一个反例检验。比如左边界回退用 \code{abba} 检验，前缀和初始零用单元素等于目标检验，柱状图宽度用递减数组检验，背包循环方向用单物品容量二检验，图状态维度用剩余资源不同的同位置状态检验。没有反例的“注意边界”不是分析，只是提醒。

一个题可以有多个解法，但每个解法都要对应自己的疑问。接雨水的前后缀数组回答“每个位置的左右最高如何预处理”，双指针回答“何时能用较低已知边界结算”，单调栈回答“什么时候槽底知道左右边界”。如果把三种解法都写成“维护状态，线性扫描”，就是在抹掉差异。

\noindent\textbf{审查二：暴力版本是否真的服务推导。}

暴力版本不是形式。它应该说明候选空间是什么，检查条件是什么，复杂度慢在哪一步。三数之和暴力枚举三元组，慢在组合数量和重复答案；排序双指针正是利用单调性和相邻重复来减少这两类问题。编辑距离暴力尝试三种操作，慢在相同前缀对反复出现；DP 状态正是给前缀对命名。若暴力版本只写“暴力会超时”，它没有提供推导价值。

暴力也不一定要写代码，但必须能口头执行。读者应当能拿一个小输入，按暴力过程跑完，看到哪些操作被重复。窗口题可以枚举所有左右端；图题可以从每个起点重新搜索；区间题可以两两比较；设计题可以每次操作扫描全部元素。只有低效过程清楚，优化才不是突然出现。

审稿时可以问一句：优化到底替换了暴力中的哪一行工作？哈希替换历史扫描，前缀和替换重复区间求和，单调栈替换左右边界寻找，堆替换每轮找极值，DP 替换重复递归，并查集替换重复连通搜索。若回答不上来，就说明题解在堆算法名。

\noindent\textbf{审查三：状态是否有不变量。}

高质量题解必须写出状态含义，而不是只写变量名。slow 左侧保存什么，窗口内满足什么，哈希表保存当前下标之前还是包含当前下标，栈里是等待右边界的柱子还是等待更高温的日子，堆顶是否可能过期，dp[i][j] 表示哪个前缀，parent 数组是不是原图路径。这些问题都应在正文中说清。

不变量要能指导代码顺序。前缀和题先查再更新，是因为哈希表表示历史前缀；颜色分类遇到二不移动扫描指针，是因为换回来的元素仍属于未知区；LRU 更新已有 key 要移动节点，是因为访问也改变最近使用顺序；拓扑排序入度变零才入队，是因为所有先修已满足。若状态含义不能决定代码顺序，说明它还不够具体。

不变量还要能解释结束条件。二分结束时 left 为什么是答案，BFS 第一次到达为什么最短，单调队列队首为什么是当前窗口最大值，DP 表最后一个格子为什么是全局答案，回溯到叶子时为什么可以收集路径。很多错误代码不是主体循环错，而是答案读取时机错。

\noindent\textbf{审查四：C++ 内容是否和题目相关。}

C++ 提醒不能每题都写同一句“注意 \code{std::int64_t}”。只有涉及求和、乘积、面积、距离、时间和路径代价时，才重点讨论溢出。只有哈希查询语义会影响状态时，才强调 \code{find} 和 \code{operator[]} 的区别。只有保存节点身份时，才强调 \code{list} 迭代器稳定。只有比较器决定排序不变量时，才重点检查严格弱序。

容器选择也要跟访问模式绑定。vector 适合连续扫描和随机访问，deque 适合两端操作，unordered map 适合等值历史查询，map 适合前驱后继，priority queue 适合动态极值，list 适合已知节点常数移动，array 适合固定字符集计数。若题解只说“用某容器更方便”，没有说明它回答哪个查询，就应该重写。

代码块要少而精。一本教材不需要每道题都贴完整最终代码；更重要的是给关键骨架、状态更新和易错行。代码排版必须稳定：缩进清晰，命名表达状态，辅助函数拆出检查逻辑，复杂条件不要塞在一行。代码服务推理，不是用来填满页面。

\noindent\textbf{审查五：Linux 源码对照是否克制。}

源码对照只能在确实能解释访问模式时出现。LRU 对照 list head，动态区间对照 Maple Tree，哈希桶对照 hlist，动态有序集合对照 rbtree，读多写少删除对照 RCU，分配连续页对照伙伴系统。这些对照有价值，因为它们回答“真实系统为什么比题目复杂”。如果只是写“Linux 里也用了某结构”，没有说明对象生命周期、并发、分配或多索引，就没有教学价值。

源码字段名会随版本变化，教材不应把字段名当重点。重点应是接口语义、对象关系、热路径和不变量。读 list head 就问节点如何嵌入对象；读 rbtree 就问比较逻辑在哪里；读 XArray 就问稀疏整数索引如何分层；读 Maple Tree 就问范围查询如何保持相邻关系；读 RCU 就问删除为何不能立刻释放。每次只追一个问题，读者才不会迷路。

源码实验也要可执行。写最小侵入式链表，写固定桶数哈希表，写 map 区间表，写简化伙伴分配器，写读写配置表模拟 RCU。实验不复刻内核，而是复现一个不变量。报告里必须写明和刷题题目的相同点和差异点，防止类比过度。

\noindent\textbf{审查六：重复句子必须删除。}

修订时可以直接搜索一些危险句式：这题用某结构、注意边界、时间复杂度降低、维护状态、不要套模板、实验时对拍、C++ 要注意容器。这些句子不是绝对不能出现，但如果没有跟具体题目绑定，就应删除或改写。比如“注意边界”要改成“目标串含重复字符时 formed 表示达标字符种类，不是匹配总数”；“维护状态”要改成“哈希表保存当前下标之前每个前缀和的出现次数”。

同一章节内，相邻题目的开头也要避免同构。可以用题面冲突开头，可以用反例开头，可以用暴力慢点开头，可以用错误实现开头，可以用源码类比开头。读者需要节奏变化。若十道题都以“先从暴力开始”开头，即使内容不同，也会产生复制感。

真正必要的重复是术语稳定，而不是段落重复。比如“不变量”“候选”“历史状态”“前驱后继”“过期状态”这些词可以反复出现，因为它们构成学习框架；但每次出现都要落到不同对象上。稳定词汇让书有体系，重复段落会让书像拼接稿。

\noindent\textbf{审查七：章节是否能串起来。}

每章开头要告诉读者上一章解决了什么问题，本章接着解决什么问题。数组讲连续内存和下标，窗口讲连续区间的增量维护，前缀和讲窗口失效后如何用代数历史状态补救，单调结构讲如何保存未解决候选，堆讲动态极值，图讲状态扩展，DP 讲重复子问题，源码讲真实工程约束。这条线如果断了，读者就会觉得每章都在换话题。

章节结尾要给迁移任务，而不是只总结。学完前缀和，让读者比较 560、523、525、862；学完单调栈，让读者比较每日温度、柱状图、下一个更大元素；学完背包，让读者比较 416、494、518、322；学完图，让读者比较岛屿、课程表、网络延迟、公交路线。迁移任务能强迫读者看到条件变化，而不是只记题号。

目录层级也要服务这种串联。大章标题表达学习阶段，小节标题表达核心问题，过碎的小节不要全部进入目录。正文中的图示、表格和实验框可以帮助定位，但不应把每个提醒都变成标题。一本书的排版舒适，来自稳定层级和连续叙述，而不是标题数量。

\noindent\textbf{审查八：如何判断一段应该扩写还是删除。}

一段文字如果回答了为什么、如何证明、哪里会错、如何实验、如何迁移、如何和源码差异化，就值得扩写。比如“Dijkstra 堆里有旧状态”值得扩，因为它解释了 priority queue 的接口限制和正确性检查；“窗口要维护计数”只有在说明 formed、need、重复字符时才值得扩；“DP 可以空间压缩”只有在说明旧值来源时才值得扩。

一段文字如果只是换词重复前文，就应该删除。比如已经在题卡里讲过“哈希保存历史”，后面再写“哈希能快速查询历史”就没有新信息，除非加入具体 key、value、查询顺序或反例。已经在章节里讲过“数组连续内存缓存友好”，后面再重复这句话也没有意义，除非加入 vector 扩容、迭代器失效或和 list 的实验对比。

扩写和删除都要服务读者。字数目标是最低体量，不是写作目标。一本三十万字的教材如果每一万字都能教会一个新判断，就是扎实；如果只是同一句话重复三十遍，就是负担。后续修订时宁愿删除一千字模板，也要补一千字能解释一个具体错误的内容。

\noindent\textbf{审查九：最终交付前的抽样方法。}

最终交付前，不要只看总字数和编译成功。随机抽十道题，遮住标题之外的内容，只读正文第一段，判断是否能看出本题自己的问题。若第一段换到另一道题也能成立，就说明太模板。再抽五个代码块，检查缩进、字体、行宽、注释和变量名。再抽五个实验，确认它们能复现一个错误或证明一个不变量。

还要抽查跨章节引用。接雨水在数组窗口章节讲过，在单点深挖里又讲，两个地方应该互补，而不是重复。LRU 在链表、设计题和 Linux 对照里都出现，三个地方应分别强调链表移动、索引一致性和侵入式节点。Dijkstra 在图章节和设计实验里出现，应分别强调非负边证明和旧堆状态。重复题可以出现，但角度必须不同。

最后检查手机阅读体验。段落不能太碎，代码不能横向溢出，目录不要塞满细碎小标题，彩色强调不能影响正文黑色阅读，图示要解释状态流动。手机上读算法书最怕目录混乱和代码挤压，所以构建 PDF 和 EPUB 后都要实际打开检查。技术内容和排版体验必须一起达标。

抽查结果要写进修订记录。发现一个模板段，就回到对应题目重新补模型、反例和实验，而不是只换几个词。

\noindent\textbf{审查十：术语统一但例子必须变化。}

全书可以统一使用一些核心术语，例如候选、状态、不变量、历史信息、支配关系、边界、前驱、后继、过期状态、逻辑删除、物理释放。这些词能让读者在不同章节之间建立联系。但术语统一不等于例子重复。每次出现“不变量”，都要落到具体对象：窗口题是不重复窗口，数组题是已写入前缀，栈题是仍未结算的下标，DP 题是已经计算完的依赖状态，源码题是对象已经从索引摘除但还没释放。

同样，每次出现“支配关系”也要换成题目自己的语言。单调队列里，新下标前缀和更小且位置更靠后，所以支配旧下标；滑动窗口最大值里，新值更大且更晚，所以支配旧值；柱状图里，当前更矮柱不是支配旧柱，而是触发旧柱结算；贪心题里的支配常常要靠交换证明。若这些差异没有写出来，术语就会变成空壳。

最后抽查一句话是否合格，可以问它能否直接指导代码。比如“维护历史状态”不合格；“哈希表保存当前下标之前每个前缀和出现次数，扫描当前前缀时先查询再更新”合格。比如“注意链表指针”不合格；“改 current next 前先保存 next，反转后 previous 指向已处理前缀头，current 指向未处理第一节点”合格。教材中的每个关键句都应该尽量达到这种可执行程度。

\noindent\textbf{最终自测题。}

自测一：给一个含负数数组，问和至少为 K 的最短子数组。请说明为什么普通滑动窗口失效，为什么前缀和单调队列成立，队尾弹出条件是什么，队首弹出条件是什么。

自测二：给一个动态日程表，要求插入区间并拒绝冲突。请说明为什么哈希表不合适，为什么只检查前驱和后继，半开区间下端点相等是否冲突。

自测三：给一个带正权图，要求从起点到所有点的最短时间。请说明为什么 BFS 不够，Dijkstra 的过期堆状态如何处理，距离数组使用什么类型。

自测四：给一个链表设计题，要求 O(1) 移动节点。请说明节点由谁拥有，哈希表保存什么，删除节点时哪些索引要同步更新，是否能用 vector 保存节点。

自测五：给一个背包题，要求统计方案数。请说明物品能用几次，方案是否区分顺序，容量循环方向是什么，\code{dp[0]} 为什么初始化为一。

这些自测题不要求背出固定代码，但要求能完整说出模型、暴力、瓶颈、状态、不变量、边界和 C++ 风险。如果说不完整，就回到对应章节重写复盘报告。

\noindent\textbf{错题复盘样例。}

\noindent\textbf{错题一：窗口左端回退。}

错误现象：无重复字符最长子串在输入 \code{abba} 上返回三。错误原因通常是看到重复字符后直接把左端设为上次位置加一，没有检查上次位置是否仍在当前窗口内。处理最后一个 \code{a} 时，它上次出现在下标零，但当前窗口左端已经在下标二；如果把左端拉回一，就会让已经排除的字符重新进入窗口，破坏窗口无重复不变量。

复盘写法：窗口左端只能右移，不能回退。更新公式应当是取当前左端和上次位置后一位的较大值。这个错误说明我们没有把“窗口保存当前无重复子串”落实成不变量，只是在机械处理重复字符。修复后，再用 \code{abba}、\code{tmmzuxt}、全相同字符和全不同字符测试。报告中要记录每一步左端变化，而不是只写最终答案。

\noindent\textbf{错题二：前缀和漏掉开头区间。}

错误现象：和为 K 的子数组在输入 \code{[3]}、目标 \code{3} 时返回零。错误原因是哈希表没有预先放入前缀和零的一次出现。区间从下标零开始时，它对应的历史前缀是空前缀；如果不记录空前缀，就等于告诉算法“从开头开始的区间不存在”。

复盘写法：前缀和题必须解释历史状态的时间点。扫描到当前前缀 \code{s} 时，答案增加历史前缀 \code{s-k} 的出现次数；空前缀是合法历史状态，位置在数组开始之前。先查询再更新当前前缀，是为了避免把当前前缀当成历史前缀。修复后测试目标为零、负数、重复前缀和从开头开始的多段答案。

\noindent\textbf{错题三：二分检查函数错误。}

错误现象：船运包裹二分循环看似标准，但某些用例返回容量小于最大包裹。错误原因不是二分边界，而是检查函数没有在当前包裹超过容量时立即判不可行，或者下界没有设为最大包裹。二分只能在正确谓词上找边界，不能修复错误谓词。

复盘写法：先写谓词定义，再写二分。容量 \code{cap} 可行，含义是按原顺序把包裹分成不超过 D 段，每段和不超过 \code{cap}。如果某个包裹重量已经超过 \code{cap}，谓词必假。下界设为最大包裹，上界设为总重量，可以让检查函数更简单。每次失败用例都要打印 \code{cap} 和需要天数，确认单调性是否存在。

\noindent\textbf{错题四：柱状图宽度少一。}

错误现象：柱状图最大矩形在递减数组或含哨兵用例上面积偏小。错误原因常在弹栈后的宽度计算。弹出柱子后，当前下标是右侧第一个更矮位置；弹出后新的栈顶是左侧第一个更矮位置。有效宽度不包含左右两个更矮边界，因此是 \code{right - left - 1}。如果栈为空，说明左侧没有更矮边界，宽度是当前下标。

复盘写法：单调栈弹出时结算的是被弹出的柱子，不是当前柱子。当前柱子只是右边界。每一次弹出都要能说出左边界、右边界、高度和宽度。修复后测试严格递增、严格递减、全相等、单个柱子和末尾哨兵。这个错误说明我们背了栈模板，却没有理解弹出时机。

\noindent\textbf{错题五：背包循环方向写反。}

错误现象：分割等和子集在只有一个物品时却能凑出两倍容量。原因是 0-1 背包一维压缩时容量正序遍历，导致当前物品刚更新出的状态又被当前物品再次使用。这个错误在小样例上可能不明显，但一旦构造单物品反例就会暴露。

复盘写法：循环方向来自物品能使用几次。0-1 背包每个物品只能用一次，所以容量倒序，读取上一轮物品阶段的状态；完全背包物品可以无限使用，所以容量正序，允许当前物品在同一轮继续贡献。报告里不要写“背包模板”，要写“当前题每个数只能选一次，因此倒序容量”。再用零钱兑换 II 对比，说明组合计数为什么外层遍历硬币。

\noindent\textbf{错题六：图状态维度缺失。}

错误现象：带钥匙、障碍消除或剩余资源的网格最短路返回过大或错误不可达。原因是访问标记只按位置记录，把“同一位置但资源不同”的状态合并了。位置相同不代表未来能力相同；剩余障碍次数更多或钥匙集合不同，会改变后续可走边。

复盘写法：图题先定义状态，不只是节点位置。若未来决策依赖资源，就把资源放进状态和访问标记。障碍消除题的 visited 至少应包含行、列、剩余消除次数；钥匙迷宫要包含行、列、钥匙集合。BFS 的第一次到达最短，只对完整状态成立，不对被错误压缩的位置成立。测试要构造必须绕路保留资源的用例。

\noindent\textbf{错题七：LRU 更新已有 key。}

错误现象：LRU 在 \code{put} 已存在 key 后淘汰错误元素。原因是只更新了 value，没有把节点移动到最新位置，或者新建了一个重复节点而没有删除旧节点。LRU 的“使用”包括 get，也包括更新已有 key；只要访问或更新成功，它就应成为最近使用。

复盘写法：LRU 有两个不变量：哈希表中每个 key 对应链表中唯一节点，链表顺序从最新到最旧。更新已有 key 时，不能增加容量，不能产生重复节点，必须移动到头部。容量满淘汰时，要从链表尾找到 key，再从哈希表删除。测试容量为一、重复 put、get 后再 put、更新后淘汰等序列。

\noindent\textbf{错题八：Trie 搜索没有恢复访问标记。}

错误现象：单词搜索在某些路径上漏答案，或者同一格子被重复使用。原因通常是 DFS 进入格子后没有在返回时恢复访问标记，导致其他分支误以为该格不可用；或者根本没有访问标记，导致一条路径中重复使用同一格。Trie 剪枝只能判断前缀是否存在，不能替代网格路径的访问约束。

复盘写法：回溯代码必须成对出现：做选择、递归、撤销选择。访问标记属于路径状态，而不是全局永久状态。找到一个单词后可以清空 Trie 节点上的单词标记避免重复答案，但不能清空网格访问规则。测试棋盘有重复字母、同一单词多条路径、单字符单词和路径必须回退的用例。

\noindent\textbf{最终质量检查表。}

交付一本算法教材前，逐章检查以下问题。第一，是否每个核心算法都有暴力基线；第二，是否解释了优化替换了哪一步；第三，是否给出至少一个反例或边界实验；第四，是否说明 C++ 容器选择和风险；第五，是否有 Linux 或工程源码视角帮助读者迁移；第六，是否避免同一句话在大量题卡中机械重复。

对单题检查也用同样标准。题面模型是否一句话说清；输入规模是否决定复杂度预算；状态是否足够；不变量是否能反驳错误实现；边界是否进入测试表；代码是否有溢出、迭代器失效、默认插入、递归深度或比较器风险；变体是否重新证明。只有这些都通过，才算高质量题解。

最后检查阅读体验。章节标题不要过碎，小节应服务于连续讲解；目录只保留大结构，不把每个细碎点都塞进去；代码块使用稳定宽度和清晰缩进；正文以黑色为主，颜色用于提示、强调和代码；图示用于解释状态流动，而不是装饰。一本教材的质量不仅取决于内容量，也取决于读者能否顺着推理读下去。

\noindent\textbf{读者自查问答。}

\noindent\textbf{看到新题先问什么。}

第一问输入规模。规模决定暴力能否作为最终答案，也决定是否需要对数、线性、平方或伪多项式方法。第二问数据条件。数组是否有序，元素是否全正，字符集是否固定，图边权是否相同，树是否是搜索树，区间端点是闭区间还是半开区间。第三问输出目标。是判断、计数、最值、路径、所有方案还是修改后的结构。输出目标不同，保存的状态就不同。

第四问能否从暴力中看出重复工作。重复扫描历史值，可能用哈希；重复求区间和，可能用前缀；重复找极值，可能用堆或单调队列；重复计算子问题，可能用 DP；重复判断连通，可能用并查集。第五问优化依赖的条件是否稳定。如果条件变了，比如正数变成有负数，有序变成无序，等权变成带权，原算法就要重新证明。

\noindent\textbf{如何判断自己不是背模板。}

把答案遮住，只保留题面，尝试写出暴力版。如果暴力版说不清，说明还没理解候选空间。再尝试说出瓶颈是哪一步，而不是直接说某个算法名。然后说明优化结构维护什么状态，为什么这个状态足够回答输出目标。最后拿一个反例解释错误写法为什么错。能完成这四步，才说明不是背模板。

另一个检查方法是改条件。把数组加入负数，把返回一个答案改成返回所有答案，把判断存在改成统计数量，把静态输入改成动态更新，把普通树改成 BST 或把 BST 改成普通树。若你能判断原解法哪些部分保留、哪些部分失效，就说明理解的是结构。若只能说“这个题我见过”，说明还停留在题号记忆。

\noindent\textbf{如何把代码写得更像工程答案。}

工程答案要把风险前置。函数入口检查不可行条件，变量名表达状态语义，复杂判断拆成辅助函数，容器选择和访问模式一致。二分答案把检查函数单独写；图算法把邻居生成集中；链表操作把摘除和插入封装；DP 把初始化和转移分开。这样写不是为了增加行数，而是让不变量有清晰位置。

还要避免隐式成本。不要在热循环里反复构造大字符串；不要在 range-for 中复制复杂对象；不要把可能溢出的乘法留在 int；不要保存会被 vector 扩容失效的指针；不要用 \code{operator[]} 做纯查询；不要写不满足严格弱序的比较器。面试时间有限，也应主动说明这些风险，哪怕代码里只写最核心版本。

\noindent\textbf{如何读 Linux 源码不迷路。}

读源码前先写问题，不要漫无目的打开文件。比如“链表删除怎么保持前后指针”“红黑树插入时比较逻辑在哪里”“XArray 如何把整数索引分层”“RCU 为什么延迟释放”。每次只追一条路径，画出对象、索引、插入、查找、删除和释放。字段名可以暂时不背，但对象关系必须画清楚。

读完以后找一个刷题类比。list 对照 LRU，rbtree 对照动态区间，XArray 对照稀疏索引，RCU 对照逻辑删除和物理删除分离，伙伴系统对照分层空闲块。类比不是说二者完全相同，而是帮助理解访问模式。真正的源码能力来自差异：内核多了并发、分配上下文、引用计数和生命周期。

\noindent\textbf{如何长期复习。}

每周固定回看错题，而不是只做新题。错题报告按八项写：题面模型、暴力基线、瓶颈、优化结构、正确性、边界、C++ 风险、变体。一个月后随机抽十道题，只看题名重讲推导。如果讲不出暴力和瓶颈，就回到原章节；如果讲不出边界，就补实验；如果讲不出 C++ 风险，就补容器实验。

复习时把题按能力分组。前缀和组、单调结构组、答案二分组、拓扑和最短路组、背包组、区间 DP 组、设计题索引组合组。每组挑一题作为代表，再列两道变体。这样复习能建立迁移能力，而不是只增加题号数量。最终目标是看到新题时能自己搭桥：从题面条件到暴力，从暴力到瓶颈，从瓶颈到结构，从结构到证明。

